\begin{ex}%[Đề 12 - 2025 - Nguyen Son]%[2D4V1-1]
Cho hàm số $f(x)$ có đạo hàm liên tục trên $[1;4]$ thỏa mãn $f(1)=2$ và $f(x)=x^3+1-xf'(x)$. Tính $f(4)$ (Kết quả làm tròn đến hàng phần mười).
\shortans{$17{,}2$}
\loigiai{
Xét
\begin{eqnarray*}
f(x)=x^3+1-xf'(x)
&\Leftrightarrow&f(x)+xf'(x)=x^3+1\\
&\Leftrightarrow& [xf(x)]'=x^3+1 \Leftrightarrow xf(x)=\dfrac{x^4}{4}+x+C \quad(1)
\end{eqnarray*}
Thay $x=1$ vào (1), ta được $1\cdot f(1)=\dfrac{1}{4}+1+C \Rightarrow C=\dfrac{3}{4}$. \\
Vậy $xf(x)=\dfrac{x^4}{4}+x+\dfrac{3}{4} \Rightarrow f(4)=\dfrac{275}{16}\approx 17{,}2$.
}
\end{ex}

\begin{ex}%[Đề 12 - 2025 - Nguyen Son]%[2D4V1-1]
Cho hàm số $y=f(x)$ có đạo hàm liên tục trên $[1;2]$ thỏa mãn $f(1)=4$ và $f(x)=xf'(x)-2x^3-3x^2$. Tính $f(2)$.
\shortans{$20$}
\loigiai{
Với $x\in [1;2]$ ta có
\begin{eqnarray*}
f(x)=xf'(x)-2x^3-3x^2 &\Leftrightarrow& \dfrac{xf'(x)-f(x)}{x^2}=2x+3\\
&\Leftrightarrow& \left(\dfrac{f(x)}{x}\right)'=2x+3\\
&\Leftrightarrow& \dfrac{f(x)}{x}=x^2+3x+C.
\end{eqnarray*}
Do $f(1)=4$ nên $C=0 \Rightarrow f(x)=x^3+3x^2$. \\
Vậy $f(2)=2^3+3\cdot 2^2=20$.
}
\end{ex}

\begin{ex}%[2D4V1-1]%[2015_K12 Huyên Nguyễn]
Biết $\displaystyle\int   \dfrac{2x+3}{\mathrm{e}^x}\text{d}x= (ax+b)\mathrm{e}^{-x}+C$ (với $a$, $b\in \mathbb{R}$). Tính giá trị của $a-2b$.
\shortans{$8$}
\loigiai
{
Đặt $F(x)=(ax+b)\mathrm{e}^{-x}+C\Rightarrow F'(x)=\dfrac{2x+3}{\mathrm{e}^x}=(2x+3)\mathrm{e}^{-x}$.
\begin{eqnarray*}
F'(x)=(2x+3)\mathrm{e}^{-x} &\Leftrightarrow& -(ax+b)\mathrm{e}^{-x}+a  \mathrm{e}^{-x}=(2x+3)\mathrm{e}^{-x} \\
&\Leftrightarrow&(-ax+a-b)\mathrm{e}^{-x}=(2x+3)\mathrm{e}^{-x}\\
&\Leftrightarrow &\heva{&-a=2\\& a-b=3.}
\end{eqnarray*}
Suy ra $\heva{&a=-2\\&b=-5}\Rightarrow a-2b =8$.
}
\end{ex}

\begin{ex}%[2D4V1-2]%[2015_K12 Huyên Nguyễn]
Cho $F(x)$ là một nguyên hàm của hàm số $f(x)=\dfrac{1}{x}$ và thỏa mãn $F(1)=-1$. Có tất cả bao nhiêu giá trị của $x$ để $2F(x)=\dfrac{1}{F(x)+1}-1 $?
\shortans{$4$}
\loigiai{
Ta có $F(x)=\displaystyle\int\dfrac{1}{x}\mathrm{\,d}x=\ln|x|+C$.\\
Mà $F(1)=-1\Leftrightarrow \ln|1|+C=-1\Leftrightarrow C=-1$, suy ra $F(x)=\ln |x|-1$.\\
Do đó
\begin{eqnarray*}
2F(x)=\dfrac{1}{F(x)+1}-1&\Leftrightarrow& 2\ln |x|-1=\dfrac{1}{\ln|x|}\\
&\Leftrightarrow &2\ln^2|x|-\ln|x|-1=0\\
&\Leftrightarrow & \hoac{&\ln|x|=1\\&\ln|x|=-\dfrac{1}{2}}\\
&\Leftrightarrow & x\in \left\{\pm \mathrm{e};\pm \dfrac{1}{\sqrt{\mathrm{e}}}\right\}.
\end{eqnarray*}
Vậy có $4$ giá trị của $x$ để $2F(x)=\dfrac{1}{F(x)+1}-1 $.
}
\end{ex}

\begin{ex}%[2D4V1-2]
Cho hàm số $ f(x)$ có đạo hàm trên $\mathbb{R}$ thỏa mãn $\mathrm{e}^{f(x)}-\dfrac{x}{f'(x)}=0$, $\forall x\in\mathbb{R}$. Biết $f(1)=1$, tính $f\left(\mathrm{e}^2\right)$ (\textit{kết quả làm tròn đến hàng phần trăm}).
\shortans{$ 3{,}38$}
\loigiai{
Ta có
\begin{align*}
\mathrm{e}^{f(x)}-\dfrac{x}{f'(x)}=0&\Rightarrow f'(x)\mathrm{e}^{f(x)}=x\\&\Leftrightarrow \left(\mathrm{e}^{f(x)} \right)'=x\\&\Leftrightarrow
\mathrm{e}^{f(x)}=\displaystyle\int x\mathrm{\,d} x\\& \Leftrightarrow \mathrm{e}^{f(x)}=\dfrac{x^2}{2}+C.
\end{align*}
Mà 	$f(1)=1$ nên $\mathrm{e}=\dfrac{1}{2}+C\Rightarrow C=\mathrm{e}-\dfrac{1}{2}$.\\
Do đó $\mathrm{e}^{f(x)}=\dfrac{x^2}{2}+ \mathrm{e}-\dfrac{1}{2} \Rightarrow \mathrm{e}^{f\left(\mathrm{e}^2\right)}= \dfrac{\mathrm{e}^4}{2}+ \mathrm{e}-\dfrac{1}{2}\Rightarrow f\left(\mathrm{e}^2\right)=\ln \left(\dfrac{\mathrm{e}^4}{2}+ \mathrm{e}-\dfrac{1}{2}\right) \approx 3{,}38$.}
\end{ex}

\begin{ex}%[2D4V1-2]
Cho hàm số $ y=f(x)$ có đạo hàm liên tục trên $\mathbb{R}$ và thỏa mãn điều kiện $x^6\left( f'(x)\right) ^3+27\left[f(x)-1\right]^4=0$, $\forall x\in\mathbb{R}$ và $ f(1)=0$. Tính giá trị của $f(2)$.
\shortans{$-7$}
\loigiai{
Ta có
\begin{align*}
x^6\left( f'(x)\right)^3+27\left[f(x)-1\right]^4=0&\Leftrightarrow{x^6}{\left( f'(x)\right)^3}=-27\left( f(x)-1\right)^4\\&\Leftrightarrow\dfrac{\left( f'(x)\right) ^3}{\left( f(x)-1\right)^4}=-\dfrac{27}{x^6}\\
&\Leftrightarrow\dfrac{\left( f'(x)\right) ^3}{\left( f(x)-1\right) ^3\left( f(x)-1\right)}=-\dfrac{27}{x^6}\\&\Leftrightarrow\dfrac{f'(x)}{\left(f(x)-1\right)\sqrt[3]{f(x)-1}}=-\dfrac{3}{x^2}\\&\Leftrightarrow\dfrac{f'(x)}{-3\left(f(x)-1\right)\sqrt[3]{f(x)-1}}=\dfrac{1}{x^2}\\&\Leftrightarrow{\left[\dfrac{1}{\sqrt[3]{f(x)-1}}\right]'}=\dfrac{1}{x^2}
\end{align*}
Do đó $\displaystyle\int{\left( \dfrac{1}{\sqrt[3]{f(x)-1}}\right)'}\mathrm{\,d}x=\displaystyle\int{\dfrac{1}{x^2}\mathrm{\,d}x}=-\dfrac{1}{x}+C.$
\\
Suy ra $\dfrac{1}{\sqrt[3]{f(x)-1}}=-\dfrac{1}{x}+C$.\\
Ta có $ f(1)=0\Rightarrow C=0 \Rightarrow f(x)=1-x^3$.\\
Khi đó $ f(2)=-7$.}
\end{ex}

\begin{ex}%[2D4V1-2]
Cho hàm số $f(x)$ thỏa mãn $\left[x{f}'(x)\right]^2+1=x^2\left[1-f(x).f'' (x)\right]$ với mọi $x$ dương. Biết $f(1)=f'(1)=1$. Tính giá trị $f^2(2)$ (\textit{kết quả làm tròn đến hàng phần trăm}).
\shortans{$3{,}39$}
\loigiai{
Với mọi $x$ dương, ta có
\begin{align*}
\left[x{f}'(x)\right]^2+1=x^2\left[1-f(x)\cdot f'' (x)\right]; x>0&\Leftrightarrow{x^2}\cdot\left[f'(x)\right]^2+1=x^2\left[1-f(x)\cdot f'' (x)\right]\\
&\Leftrightarrow{\left[f'(x)\right]^2}+\dfrac{1}{x^2}=1-f(x)\cdot f'' (x)\\
&\Leftrightarrow{\left[f'(x)\right]^2}+f(x)\cdot f'' (x)=1-\dfrac{1}{x^2}\\
&\Leftrightarrow\left[f(x)\cdot f'(x)\right]'=1-\dfrac{1}{x^2}.
\end{align*}
Do đó $\displaystyle\int\left[f(x)\cdot f'(x)\right]'\mathrm{\, d}x=\displaystyle\int\left(1-\dfrac{1}{x^2}\right)\mathrm{\, d}x\Rightarrow f(x)\cdot f'(x)=x+\dfrac{1}{x}+C.$\\
Vì $ f(1)=f'(1)=1\Rightarrow 1=2+C\Leftrightarrow C=-1.$\\
Nên $\displaystyle\int f(x)\cdot f'(x)\mathrm{\, d}x=\displaystyle\int\left(x+\dfrac{1}{x}-1\right) \mathrm{\, d}x$ $\Leftrightarrow\displaystyle\int f(x)\mathrm{\, d}\left(f(x)\right)=\displaystyle\int{\left(x+\dfrac{1}{x}-1\right)}\mathrm{\, d}x$.\\
Suy ra				$\dfrac{f^2(x)}{2}=\dfrac{x^2}{2}+\ln x-x+C.$\\
Vì $ f(1)=1\Rightarrow\dfrac{1}{2}=\dfrac{1}{2}-1+C\Leftrightarrow C=1.$\\
Vậy $\dfrac{f^2(x)}{2}=\dfrac{x^2}{2}+\ln x-x+1\Rightarrow{f^2}(2)=2\ln 2+2\approx 3{,}39$.
}
\end{ex}

\begin{ex}%[2D4V1-2]
Biết rằng hàm số $f(x) = ax^2 + bx + c$ thỏa mãn $\displaystyle\int\limits_{0}^{1} f(x) \mathrm{d}x = -\dfrac{7}{2}$, $\displaystyle\int\limits_{0}^{2} f(x) \mathrm{d}x = -2$ và $\displaystyle\int\limits_{0}^{3} f(x) \mathrm{d}x = \dfrac{13}{2}$. Tính $P = a + b + c$. (\textit{làm tròn đến hàng phần trăm}).
\shortans{$-1{,}33$}
\loigiai{
Ta có
$	\displaystyle\int f(x) \mathrm{d}x = \displaystyle\int (ax^2 + bx + c) \mathrm{d}x = \dfrac{a}{3} x^3 + \dfrac{b}{2} x^2 + cx + C
$.\\
Lại có
$	\displaystyle\int\limits_{0}^{1} f(x) \mathrm{d}x = -\dfrac{7}{2} \Rightarrow \left( \dfrac{a}{3} x^3 + \dfrac{b}{2} x^2 + cx \right) \bigg|_{0}^{1} = -\dfrac{7}{2} \Rightarrow \dfrac{1}{3} a + \dfrac{1}{2} b + c = -\dfrac{7}{2} \quad (1)
$.
$	\displaystyle\int\limits_{0}^{2} f(x) \mathrm{d}x = -2 \Rightarrow \left( \dfrac{a}{3} x^3 + \dfrac{b}{2} x^2 + cx \right) \bigg|_{0}^{2} = -2 \Rightarrow \dfrac{8}{3} a + 2b + 2c = -2 \quad (2)
$.\\
$	\displaystyle\int\limits_{0}^{3} f(x) \mathrm{d}x = \dfrac{13}{2} \Rightarrow \left( \dfrac{a}{3} x^3 + \dfrac{b}{2} x^2 + cx \right) \bigg|_{0}^{3} = \dfrac{13}{2} \Rightarrow 9a + \dfrac{9}{2} b + 3c = \dfrac{13}{2} \quad (3)
$.\\
Từ (1), (2) và (3) ta có hệ phương trình:
\[\heva{&\dfrac{1}{3} a + \dfrac{1}{2} b + c = -\dfrac{7}{2} \\&
\dfrac{8}{3} a + 2b + 2c = -2 \\&
9a + \dfrac{9}{2} b + 3c = \dfrac{13}{2}}
\Rightarrow \heva{&	a = 1 \\&
b = 3 \\&
c = -\dfrac{16}{3}.}	\]
Vậy	$P = a + b + c = 1 + 3 + \left( -\dfrac{16}{3} \right) = -\dfrac{4}{3}\approx -1{,}33$.
}
\end{ex}

\begin{ex}%[2D4V1-2]
Cho $I = \displaystyle\int\limits_{0}^{1} (4x - 2m^2) \mathrm{d}x$. Có bao nhiêu giá trị nguyên của $m$ để $I + 6 > 0$?
\shortans{$3$}
\loigiai{
Theo định nghĩa tích phân ta có:
$	I = \displaystyle\int\limits_{0}^{1} (4x - 2m^2) \mathrm{d}x = \left( 2x^2 - 2m^2 x \right) \bigg|_{0}^{1} = -2m^2 + 2
$.\\
Khi đó $I + 6 > 0 \Leftrightarrow -2m^2 + 2 + 6 > 0 \Leftrightarrow -2m^2 + 8 > 0  \Leftrightarrow -2 < m < 2$.\\
Mà $m$ là số nguyên nên $m \in \{-1; 0; 1\}$.\\
Vậy có $3$ giá trị nguyên của $m$ thỏa mãn yêu cầu.
}
\end{ex}

\begin{ex}%[2D4V1-2]
Có bao nhiêu giá trị nguyên dương của $a$ để $\displaystyle\int\limits_{0}^{a} (2x - 3) \mathrm{d}x \leq 4$?
\shortans{$4$}
\loigiai{
Ta có
$
\displaystyle\int\limits_{0}^{a} (2x - 3) \mathrm{d}x = \left( x^2 - 3x \right) \bigg|_{0}^{a} = a^2 - 3a
$.\\
Khi đó
$
\displaystyle\int\limits_{0}^{a} (2x - 3) \mathrm{d}x \leq 4 \Leftrightarrow a^2 - 3a \leq 4 \Leftrightarrow -1 \leq a \leq 4
$.\\
Mà $a \in \mathbb{N}^*$ nên $a \in \{1; 2; 3; 4\}$.\\
Vậy có $4$ giá trị của $a$ thỏa đề bài.
}
\end{ex}

\begin{ex}%[2D4V1-2]
Có bao nhiêu số thực $b$ thuộc khoảng $(\pi; 3\pi)$ sao cho $\displaystyle\int\limits_{\pi}^{b} 4 \cos 2x \mathrm{d}x = 1$?
\shortans{$4$}
\loigiai{
Ta có
$
\displaystyle\int\limits_{\pi}^{b} 4 \cos 2x \mathrm{d}x = 1 \Leftrightarrow 2 \sin 2x \bigg|_{\pi}^{b} = 1 \Leftrightarrow \sin 2b - \sin 2\pi = \dfrac{1}{2} \Leftrightarrow \sin 2b = \dfrac{1}{2}
$.\\
$
\Rightarrow 2b = \dfrac{\pi}{6} + k2\pi \quad \text{ hoặc } \quad 2b = \dfrac{5\pi}{6} + k2\pi
$\\
$
\Rightarrow b = \dfrac{\pi}{12} + k\pi \quad \text{ hoặc } \quad b = \dfrac{5\pi}{12} + k\pi\qquad (k\in \mathbb{Z})
$.\\
Khi $	b = \dfrac{\pi}{12} + k\pi$, ta xét\\
\allowdisplaybreaks
\begin{eqnarray*}
&& \pi< \dfrac{\pi}{12} + k\pi<3\pi\\
&\Leftrightarrow& \dfrac{11}{12}<k<\dfrac{35}{12}\\
&\Leftrightarrow& k \in \{1;2\}.
\end{eqnarray*}
Khi $	b = \dfrac{5\pi}{12} + k\pi$, ta xét\\
\allowdisplaybreaks
\begin{eqnarray*}
&& \pi< \dfrac{5\pi}{12} + k\pi<3\pi\\
&\Leftrightarrow& \dfrac{7}{12}<k<\dfrac{31}{12}\\
&\Leftrightarrow& k \in \{1;2\}.
\end{eqnarray*}
Vậy có $4$ số thực $b$ thỏa mãn yêu cầu bài toán.
}
\end{ex}

\begin{ex}%[Dự án EX - TF - TLN 2024]%[Doan Hung]%[2D4V1-2]
Biết $F(x)$ là một họ nguyên hàm của $f(x)=\dfrac{x}{(x+1)^3}$ và $F(0)=\dfrac{1}{2}$. Khi đó $F(1)+F(2)$ bằng bao nhiêu?
\shortans{$2{,}25$}
\loigiai{
Ta có
\[F(x)=\displaystyle\int\limits \dfrac{x}{(x+1)^3}\mathrm{\,d}x=\displaystyle\int\limits \dfrac{x+1-1}{(x+1)^3}\mathrm{\,d}x=\displaystyle\int\limits \left(\dfrac{1}{(x+1)^2}-\dfrac{1}{(x+1)^3}\right)\mathrm{\,d}x=-\dfrac{1}{x+1}+\dfrac{1}{2(x+1)^2}+C.\]
Mà $F(0)=\dfrac{1}{2}$ nên suy ra \[-\dfrac{1}{2}+C=\dfrac{1}{2}\Leftrightarrow C=1.\]
Khi đó \[F(x)=-\dfrac{1}{x+1}+\dfrac{1}{2(x+1)^2}+1\Rightarrow F(1)+F(2)=\dfrac{5}{8}+\dfrac{13}{8}=\dfrac{9}{4}.\]
}
\end{ex}

\begin{ex}%[Dự án EX - TF - TLN 2024]%[Doan Hung]%[2D4V1-2]
Biết $\displaystyle\int\limits \dfrac{3x+1}{x^2+6x+9}\mathrm{\,d}x=a\ln|x+3|+\dfrac{2b}{x+3}+C$ với $a, b$ là các số nguyên. Tính giá trị của $T=2a^2+b$.
\shortans{$22$}
\loigiai{
Ta có \[f(x)=\dfrac{3x+1}{x^2+6x+9}=\dfrac{3x+1}{(x+3)^2}=\dfrac{3(x+3)-8}{(x+3)^2}=\dfrac{3}{x+3}-\dfrac{8}{(x+3)^2}.\]
Khi đó
\begin{align*}
\displaystyle\int\limits f(x)\mathrm{\,d}x=&\, \displaystyle\int\limits \dfrac{3}{x+3}\mathrm{\,d}x-\displaystyle\int\limits \dfrac{8}{(x+3)^2}\mathrm{\,d}x\\
=&\, 3\displaystyle\int\limits \dfrac{1}{x+3}\mathrm{\,d}(x+3)-8\displaystyle\int\limits (x+3)^{-2}\mathrm{\,d}(x+3)\\
=&\, 3\ln|x+3|+\dfrac{8}{x+3}+C.
\end{align*}
Suy ra $\heva{&a=3\\&2b=8}\Rightarrow\heva{&a=3\\&b=4}\Rightarrow T=2a^2+b=22$.
}
\end{ex}

\begin{ex}%[Dự án EX - TF - TLN 2024]%[Doan Hung]%[2D4V1-2]
Cho $f'(x)=\dfrac{2}{(2x-1)^2}-\dfrac{1}{(x-1)^2}$ thỏa mãn $f(2)=-\dfrac{1}{3}$. Biết phương trình có $f(x)=-1$ có nghiệm duy nhất $x=x_0$. Tính $T=2025^{x_0}$.
\shortans{$1$}
\loigiai{
Ta có $f(x)=\displaystyle\int f'(x)\mathrm{d}x=\displaystyle\int \left[\dfrac{2}{(2x-1)^2}-\dfrac{1}{(x-1)^2}\right]\mathrm{d}x=-\dfrac{1}{2x-1}+\dfrac{1}{x-1}+C$.\\
Khi đó $f(2)=-\dfrac{1}{3}\Leftrightarrow \dfrac{1}{1}-\dfrac{1}{3}+C=-\dfrac{1}{3}\Leftrightarrow C=-1\Rightarrow f(x)=-\dfrac{1}{2x-1}+\dfrac{1}{x-1}-1$\\
$f(x)=-1\Leftrightarrow -\dfrac{1}{2x-1}+\dfrac{1}{x-1}-1=-1\Leftrightarrow x=x_0=0\Rightarrow T=2025^0=1$.}
\end{ex}

\begin{ex}%[Dự án EX - TF - TLN 2024]%[Doan Hung]%[2D4V1-2]
Cho hàm số $y=f(x)$ xác định trên $\mathbb{R}\setminus \{1;2\}$ và thỏa mãn $f'(x)=|x-1|+|x-2|$, $f(0)+f\left(\dfrac{3}{2}\right)=1;f(4)=2$. Tính giá trị của biểu thức $f(-1)+f(3)+f\left(\dfrac{3}{2}\right)$.
\shortans{$-5$}
\loigiai{
$f'(x)=|x-1|+|x-2|=\heva{& 2x-3 \text{ khi } x>2 \\& 1 \text{ khi } 1<x<2 \\& 3-2x \text{ khi } x<1.} $\\
Suy ra $f(x)=\heva{& x^2-3x+c \text{ khi } x>2 \\& x+d \text{ khi } 1<x<2 \\& 3x-x^2+e \text{ khi } x<1.} $\\
Từ giả thiết $f(0)+f\left(\dfrac{3}{2}\right)=1;f(4)=2$ ta có
\[\heva{&e+\left(\dfrac{3}{2}+d\right)=1\\&c+4=2} \Rightarrow c+e+\left(\dfrac{3}{2}+d\right)=-1.\]
Suy ra $ f(-1)+f\left(\dfrac{3}{2}\right)+f(3)=(e-4)+\left(\dfrac{3}{2}+d\right)+c=-1-4=-5$.}
\end{ex}

\begin{ex}%[2D4V1-2]%[Dự án EX-TF-TLN lần 3 -Mui Doan]
Cho hàm số $y=f(x)$ thoả mãn $f(x)<0$, $\forall x>0$ và có đạo hàm $f'(x)$ liên tục trên $\left(0;+\infty\right)$ thoả mãn $f'(x)=\left(2x+1\right)f^2(x)$, $\forall x>0$ và $f(1)=-\dfrac{1}{2}$. Giá trị biểu thức $f(1)+f(2)+\cdots+f(2023)$ bằng $-\dfrac{a}{b}$ với $\dfrac{a}{b}$ là phân số tối giản. Tính $b-a$.
\shortans{$1$}
\loigiai{
Với $x\in\left(0;+\infty\right)$, do $f(x)<0$ nên $f'(x)=\left(2x+1\right)f^2(x)\Rightarrow \dfrac{f'(x)}{f^2(x)}=2x+1$, $\forall x\in\left(0;+\infty\right)$.\\
Suy ra $\displaystyle\int\dfrac{f'(x)}{f^2(x)}\mathrm{\,d}x=\displaystyle\int\left(2x+1\right)\mathrm{\,d}x\Leftrightarrow\displaystyle\int\dfrac{\mathrm{\,d}\left(f(x)\right)}{f^2(x)}=x^2+x+C\Leftrightarrow -\dfrac{1}{f(x)}=x^2+x+C$.\\
Cho $x=1$ suy ra $2=1+1+C\Rightarrow C=0$, hay $-\dfrac{1}{f(x)}=x^2+x\Rightarrow f(x)=-\dfrac{1}{x^2+x}$, $\forall x\in\left(0;+\infty\right)$.\\
Do đó $f(x)=-\dfrac{1}{x\left(x+1\right)}=-\dfrac{1}{x}+\dfrac{1}{x+1}$.\\
Vậy $f(1)+f(2)+\cdots+f(2023)=-1+\dfrac{1}{2}-\dfrac{1}{2}+\dfrac{1}{3}-\cdots-\dfrac{1}{2023}+\dfrac{1}{2024}=-\dfrac{2023}{2024}$.\\
Suy ra $a=2023$, $b=2024$'\\
Vậy $b-a=1$.
}
\end{ex}

\begin{ex}%[2D4V1-2]%[Dự án EX-TF-TLN lần 3 -Mui Doan]
Cho hàm số $y=f(x)$ xác định trên $\mathbb{R}\setminus\left\{\dfrac{1}{2}\right\}$ thoả mãn $f'(x)=\dfrac{2}{2x-1}$, $f(0)=1$, $f(1)=2$. Biết giá trị biểu thức $f(-1)+f(3)$ bằng $a+b\ln 3+c\ln 5$. Tính $a+b+c$.
\shortans{$5$}
\loigiai{
Ta có $f(x)=\displaystyle\int f'(x)\mathrm{\,d}x=\displaystyle\int\left(\dfrac{2}{2x-1}\right)\mathrm{\,d}x=\heva{&\ln\left(2x-1\right)+C_1,\,\text{khi}\, x>\dfrac{1}{2}\\& \ln\left(1-2x\right)+C_2,\,\text{khi}\, x<\dfrac{1}{2}.}$\\
Cho $x=1>\dfrac{1}{2}$ suy ra $2=\ln\left(2\cdot 1-1\right)+C_1\Rightarrow C_1=2$.\\
Cho $x=0<\dfrac{1}{2}$ suy ra $1=\ln\left(1-2\cdot 0\right)+C_2\Rightarrow C_2=1$.\\
Do đó $f(x)=\heva{&\ln\left(2x-1\right)+2,\,\text{khi}\, x>\dfrac{1}{2}\\& \ln\left(1-2x\right)+1,\,\text{khi}\, x<\dfrac{1}{2}.}$\\
Vậy $f(-1)+f(3)=\ln\left(1-2\cdot(-1)\right)+1+\ln\left(2\cdot 3-1\right)+2=3+\ln 3 +\ln 5$.\\
Suy ra $a=3$, $b=1$, $c=1$.\\
Vậy $a+b+c=3+1+1=5$.
}
\end{ex}

\begin{ex}%[2D4V1-2]%[Dự án EX-TF-TLN lần 3 -Mui Doan]
Cho hàm số $f(x)$ xác định trên $\mathbb{R}\setminus \{-2;1\}$ thoả mãn $f'(x)=\dfrac{1}{x^2+x-2}$; $f(0)=\dfrac{1}{3}$ và $f(-3)-f(3)=0$. Biết giá trị của biểu thức $T=f(-4)+f(-1)-f(4)=\dfrac{a}{b}\ln 2+\dfrac{c}{d}$ với $\dfrac{a}{b}$, $\dfrac{c}{d}$ là hai phân số tối giản. Tính $a+b-c-d$.
\shortans{$0$}
\loigiai{
Ta có
\begin{eqnarray*}
f(x)&=&\displaystyle\int f'(x) \mathrm{\,d}x=\displaystyle\int \dfrac{1}{x^2+x-2} \mathrm{\,d}x=\displaystyle\int \dfrac{1}{(x-1)(x+2)} \mathrm{\,d}x=\dfrac{1}{3}\displaystyle\int \left(\dfrac{1}{x-1}-\dfrac{1}{x+2}\right) \mathrm{\,d}x\\
&=&\heva{& \dfrac{1}{3}\ln \left|\dfrac{x-1}{x+2}\right|+C_1\text{ khi }x<-2 \\ & \dfrac{1}{3}\ln \left|\dfrac{x-1}{x+2}\right|+C_2\text{ khi }-2<x<1\\& \dfrac{1}{3}\ln \left|\dfrac{x-1}{x+2}\right|+C_3\text{ khi }x>1.}
\end{eqnarray*}
Với $x=0$ thì $f(x)=\dfrac{1}{3}\Leftrightarrow \dfrac{1}{3}\ln \dfrac{1}{2}+C_2=\dfrac{1}{3}\Leftrightarrow C_2=\dfrac{1}{3}+\dfrac{1}{3}\ln 2$.\\
Với $x=-3$ và $x=3$ thì
\begin{align*}
f(-3)-f(3)=0\Leftrightarrow \dfrac{1}{3}\ln 4+C_1-\dfrac{1}{3}\ln \dfrac{2}{5}-C_3=0\Leftrightarrow C_1-C_3=-\dfrac{1}{3}\ln 2-\dfrac{1}{3}\ln 5.
\end{align*}
Khi đó $T=f(-4)+f(-1)-f(4)=\dfrac{1}{3}\ln \dfrac{5}{2}+C_1+\dfrac{1}{3}\ln 2+C_2-\dfrac{1}{3}\ln \dfrac{1}{2}-C_3=\dfrac{1}{3}\ln 2+\dfrac{1}{3}$.\\
Suy ra $a=1$, $b=3$, $c=1$, $d=3$.\\
Vậy $a+b-c-d=1+3-1-3=0$.
}
\end{ex}

\begin{ex}%[2D4V1-2]
Gọi $F(x)$ là một nguyên hàm của hàm số $f(x)=\dfrac{(3x-1)^2}{x^2}$, biết đồ thị hàm số $y=F(x)$ đi qua điểm $M(1;-2)$. Tính $F\left(\mathrm{e}^2\right)$ (làm tròn kết quả đến hàng phần chục).
\shortans{44{,}4}
\loigiai{
Ta có
\allowdisplaybreaks
\begin{eqnarray*}
F(x)&=&\displaystyle\int \dfrac{(3x-1)^2}{x^2}\mathrm{\,d}x=\displaystyle\int \left(\dfrac{9x^2-6x+1}{x^2}\right) \mathrm{\,d}x\\
&=&\displaystyle\int \left(9-\dfrac{6}{x}+\dfrac{1}{x^2}\right) \mathrm{\,d}x\\
&=&9\displaystyle\int \mathrm{\,d}x-6\displaystyle\int\dfrac{1}{x} \mathrm{\,d}x+\displaystyle\int x^{-2} \mathrm{\,d}x\\
&=&9x-6\ln\left| x\right|-\dfrac{1}{x}+C.
\end{eqnarray*}
Theo giả thiết, đồ thị hàm số $y=F(x)$ đi qua điểm $M(1;-2)$ nên suy ra
\allowdisplaybreaks
\begin{eqnarray*}
F(1)=-2&\Rightarrow &9-6\ln 1-1+C=-2\\
&\Rightarrow & C=-10\Rightarrow F(x)=9x-6\ln\left|x\right|-\dfrac{1}{x}-10\\
&\Rightarrow& F\left(e^2\right)=9e^2-6\ln\left|\mathrm{e}^2\right|-\dfrac{1}{\mathrm{e}^2}-10\approx 44{,}4.
\end{eqnarray*}
}
\end{ex}

\begin{ex}%[12-EX-2-2025, Diamond]%[2D4V1-2]
Cho hàm số $F(x)$ là một nguyên hàm của hàm số $f(x) = 3x^2 - 4x + 1$ và $F(2)=2$. Tính $F(3)$.
\shortans{$12$}
\loigiai{
Ta có $F(x)=\displaystyle\int(3x^2-4x+1)\mathrm{\,d}x=x^3-2x^2+x+C$.\\
Từ $F(2)=2\Leftrightarrow 2^3-2\cdot 2^2+2+C=2\Leftrightarrow C=0$.\\
Do đó $F(x)=x^3-2x^2+x$, suy ra $F(3)=3^3-2\cdot 3^2+3=12$.
}
\end{ex}

\begin{ex}%[Hiếu Mai]%[2D4V1-2]
Cho hàm số $f\left(x \right)$ thỏa mãn $f''\left(x \right)=12x^2+6x-4$ và $f\left(0 \right)=1$, $f\left(1 \right)=3$. Tính $f\left(-1 \right).$
\shortans{$-3$}
\loigiai{
\begin{itemize}
\item  $f'(x)=\displaystyle\int f''(x)\mathrm{\,d}x=12\cdot \dfrac{x^3}{3}+6\cdot \dfrac{x^2}{2}-4x+C_1=4x^3+3x^2-4x+C_1$.
\item  $f(x)=\displaystyle\int f'(x)\mathrm{\,d}x=x^4+x^3-2x^2+C_1x+C_2$.
\item  Theo giả thiết $\heva{&f(0)=1\\&f(1)=3} \Leftrightarrow \heva{&C_2=1 \\ &C_1+C_2=3}\Leftrightarrow \heva{&C_2=1 \\ &C_1=2.}$
\item  Suy ra $f(x)=x^4+x^3-2x^2+2x+1$. Vậy $f(-1)=-3$.
\end{itemize}
}
\end{ex}

\begin{ex}%[HSA - đợt 1-25-26, Hector Tran]%[2D4V1-2]
Hàm số $f(x)$ xác định, liên tục trên $\mathbb{R}$ và có đạo hàm là $f'(x)=|x-1|$. Biết rằng $f(0)=3$. Tổng $f(2)+f(4)$ bằng bao nhiêu?
\par\shortans[oly]{12}
\loigiai{
Ta có $f'(x) = \heva{&x-1 & \text{khi } x \ge 1 \\&-(x-1) & \text{khi } x < 1.}$\\
Khi $x \ge 1$ thì $f(x) = \displaystyle\int (x-1)\mathrm{\,d}x = \dfrac{x^2}{2}-x+C_1$.\\
Khi $x < 1$ thì $f(x) = -\displaystyle\int (x-1)\mathrm{\,d}x = -\left(\dfrac{x^2}{2}-x\right)+C_2$.\\
Theo đề bài ta có $f(0)=3$ nên $C_2=3 \Rightarrow f(x) = -\left(\dfrac{x^2}{2}-x\right)+3$ khi $x<1$.\\
Mặt khác do hàm số $f(x)$ liên tục tại $x=1$ nên
\begin{eqnarray*}
&&\lim_{x \to 1^-} f(x) = \lim_{x \to 1^+} f(x) = f(1)\\
&\Leftrightarrow& \lim_{x \to 1^-} \left[ -\left(\dfrac{x^2}{2}-x\right)+3 \right] = \lim_{x \to 1^+} \left[ \left(\dfrac{x^2}{2}-x\right)+C_1 \right]\\
&\Leftrightarrow& -\left(\dfrac{1}{2}-1\right)+3 = \dfrac{1}{2}-1+C_1\\
&\Leftrightarrow& C_1=4.
\end{eqnarray*}
Vậy khi $x \ge 1$ thì $f(x) = \dfrac{x^2}{2}-x+4$.\\
$\Rightarrow f(2)+f(4)=12$.
}
\end{ex}

\begin{ex}%[HSA-đợt 1,Huy Trần- Đăng Tạ]%[2D4V1-2]
Cho hàm số $f(x)$ thỏa mãn $f(1)=4$ và $f(x)=x f'(x)-2 x^3-3 x^2$ với mọi $x>0$. Giá trị của $f(2)$ bằng bao nhiêu?
\shortans {20}
\loigiai{
Ta có\\
\begin{eqnarray*}
&&	f(x)=x f'(x)-2 x^3-3 x^2\\
&\Leftrightarrow& f(x)-x f'(x)=-2 x^3-3 x^2 \\
&\Leftrightarrow& \dfrac{1 \cdot  f(x)-x \cdot f'(x)}{x^2}=\frac{-2 x^3-3 x^2}{x^2}\\
& \Leftrightarrow& \left[\dfrac{f(x)}{x}\right]'=2 x+3.
\end{eqnarray*}
Suy ra $\dfrac{f(x)}{x}$ là một nguyên hàm của hàm số $g(x)=2 x+3$.\\
Ta có $\int(2 x+3) d x=x^2+3 x+C, C \in \mathbb{R}$.\\
Do đó, $\dfrac{f(x)}{x}=x^2+3 x+C_1$\quad (1) với $C_1 \in \mathbb{R}$.\\
Vì $f(1)=4$, thay $x=1$ vào hai vế của (1), suy ra $C_1=0$.\\
Vậy $f(x)=x^3+3 x^2$.\\
Do đó $f(2)=20$.
}
\end{ex}

\begin{ex}%[2D4V1-3]%[Tổ 1 - Đợt 17 - Chương 4 - - CD - Đề 2]%[Phùng Minh Phúc]
Nguyên hàm $F\left( x \right)$ của hàm số $F\left( x \right)=\text{si}{{\text{n}}^{2}}2x\cdot \text{co}{{\text{s}}^{3}}2x$ thỏa $F\left( \dfrac{\pi }{4} \right)=0$
\shortans[oly]{}
\loigiai{
Đặt $t=\sin2x\Rightarrow dt=2\cdot \cos2x\text{d}x\Rightarrow \dfrac{1}{2}dt=\cos2x\text{d}x$.\\
Ta có\\
$\begin{aligned}
F\left( x \right)&=\displaystyle\int\limits ~\text{si}{{\text{n}}^{2}}2x\cdot \text{co}{{\text{s}}^{3}}2x\text{d}x=\dfrac{1}{2}\displaystyle\int\limits ~{{t}^{2}}\cdot \left( 1-{{t}^{2}} \right)dt=\dfrac{1}{2}\displaystyle\int\limits ~\left( {{t}^{2}}-{{t}^{4}} \right)dt\\
&=\dfrac{1}{6}{{t}^{3}}-\dfrac{1}{10}{{t}^{5}}+C =\dfrac{1}{6}\text{si}{{\text{n}}^{3}}2x-\dfrac{1}{10}\text{si}{{\text{n}}^{5}}2x+C.\\
\end{aligned}$\\
$F\left( \dfrac{\pi }{4} \right)=0\Leftrightarrow \dfrac{1}{6}\text{si}{{\text{n}}^{3}}\dfrac{\pi }{2}-\dfrac{1}{10}\text{si}{{\text{n}}^{5}}\dfrac{\pi }{2}+C=0\Leftrightarrow C=-\dfrac{1}{15} $.\\
Vậy $F\left( x \right)=\dfrac{1}{6}\text{si}{{\text{n}}^{3}}2x-\dfrac{1}{10}\text{si}{{\text{n}}^{5}}2x-\dfrac{1}{15}$.
}
\end{ex}

\begin{ex}%[2D4V1-3]%[2015_K12 Huyên Nguyễn]
Cho $F(x)$ là một nguyên hàm của hàm số $f(x)=\cos 2x$ và thỏa mãn $F(\pi)=1$. Phương trình $F(x)=1$ có tất cả bao nhiêu nghiệm trong đoạn $[0 ; 3 \pi]$?
\shortans{$7$}
\loigiai{
Ta có $F(x)=\displaystyle\int\cos 2x\mathrm{\,d}x=\dfrac{1}{2}\sin 2x+C$.\\
$F(\pi)=1\Leftrightarrow C=1$, suy ra $F(x)=\dfrac{1}{2}\sin 2x +1$.\\
Phương trình
\begin{eqnarray*}
F(x)=1&\Leftrightarrow&\dfrac{1}{2}\sin 2x +1=1\\
&\Leftrightarrow&\sin 2x=0\\
&\Leftrightarrow& 2x=k\pi,\ k\in \mathbb{Z}\\
&\Leftrightarrow& x=\dfrac{k\pi}{2},\ k\in \mathbb{Z}
\end{eqnarray*}
Vậy trong đoạn $[0;3\pi]$ phương trình có tất cả $7$ nghiệm $x\in \left\{0;\dfrac{\pi}{2};\pi; \dfrac{3\pi}{2};2\pi; \dfrac{5\pi}{2};3\pi\right\}$.
}
\end{ex}

\begin{ex}%[2D4V1-3]
Cân nặng của một số lợn con mới sinh thuộc hai giống $A$ và $B$ được cho ở biểu đồ dưới đây (đơn vị: $\mathrm{kg}$).
\begin{center}
\begin{tikzpicture}[scale=1.2, font=\footnotesize, line join=round, line cap=round, >=stealth]
\coordinate [label= left: $0$](O) at (0,0) ;
\coordinate [label= left: (Số con)](y) at (0,6) ;
\coordinate [label= below: (kg)](x) at (8,0) ;
\foreach \i in {1,2,3,4,5}
{\draw[-,color=gray!40] (0,\i)--(8,\i); }
\draw (0,1)node[left]{$10$}(0,2)node[left]{$20$}(0,3)node[left]{$30$}(0,4)node[left]{$40$}(0,5)node[left]{$50$};
%--------------------------------------------
\foreach \i/\j/\a in
{0.4/0.8/8,1.2/1.3/13,2/2.8/28,2.8/1.4/14,3.6/3.2/32,4.4/2.4/24,	5.2/1.7/17,6/1.4/14}{
\draw[] ($(\i,\j)+(0,0.2)$) node[scale=0.8]{$\a$};
}
%-----------------------------------------------
\foreach \i/\j in
{0/0.8,1.6/2.8,3.2/3.2,4.8/1.7}{
\fill[blue!50] (\i,0) rectangle+ (0.8,\j);
\draw[thick] (\i,0) rectangle+ (0.8,\j);
}
%-----------------------------------------------
\foreach \i/\j in
{0.8/1.3,2.4/1.4,4/2.4,5.6/1.4}{
\fill[pattern=north east lines] (\i,0) rectangle+ (0.8,\j);
\draw[thick] (\i,0) rectangle+ (0.8,\j);
}
%-----------------------------------------------
\foreach \i/\j in
{0.4/[1{,}0;1{,}1),2/[1{,}1;1{,}2),3.6/[1{,}2;1{,}3),5.2/[1{,}3;1{,}4)}{
\draw($(\i,0)+(0.5,0)$) node[below]{\j};
}
\foreach \i/\j in {}{\draw (\i,\j) node[above]{\j};
}
\draw[->,very thick] (O)--(x);
\draw[->,very thick] (O)--(y);
\fill[blue!50,draw] (5.5,4.2) rectangle (6,4.8);
\fill[pattern=north east lines,draw] (5.5,3.2) rectangle (6,3.8);
\draw (6,4.5) node [right]{Giống $A$};
\draw (6,3.5) node [right]{Giống $B$};
\draw (4,5.5) node [above]{\textbf{Cân nặng của một số lợn con mới sinh}};
\end{tikzpicture}
\end{center}
Gọi $\Delta Q_1$, $\Delta Q_2$ lần lượt là giá trị khoảng tứ phân vị của lợn con giống $A$ và $B$. Tìm giá trị của $\Delta Q_2-\Delta Q_1$ (làm tròn kết quả đến hàng phần trăm).
\shortans[]{$0{,}03$}
\loigiai{
Cân nặng của lợn con giống $A$ và giống $B$ được thống kê như sau:
\begin{center}
\setlength{\extrarowheight}{2pt}
\newcommand{\PreBackslash}[1]{\let\temp=\\#1\let\\=\temp}
\let\PBS=\PreBackslash
\begin{tabular}{|>{\PBS\centering}m{4cm}|>{\PBS\centering}m{2cm}|>{\PBS\centering}m{2cm}| >{\PBS\centering}m{2cm}|>{\PBS\centering}m{2cm}|}\hline
\textbf{Cân nặng ($\mathrm{kg}$)} &$[1{,}0;1{,}1)$&$[1{,}1;1{,}2)$&$[1{,}2;1{,}3)$&$[1{,}3;1{,}4)$\\
\hline
\textbf{Giá trị đại diện}&$1{,}05$&$1{,}15$&$1{,}25$&$1{,}35$\\
\hline
\textbf{Số con giống $A$}&$8$&$28$&$32$&$17$\\
\hline
\textbf{Số con giống $B$}&$13$&$14$&$24$&$14$\\
\hline
\end{tabular}
\end{center}
\begin{enumerate}[$\bullet$]
\item Xét mẫu số liệu của giống lợn $A$.\\
Ta có $n=85$ nên tứ phân vị thứ nhất thuộc nhóm $[1{,}1;1{,}2)$.\\
Do đó $Q_1=1{,}1+\dfrac{\dfrac{85}{4}-8}{28}\cdot (1{,}2-1{,}1)=\dfrac{257}{224}$.\\
Tứ phân vị thứ ba thuộc nhóm $[1{,}2;1{,}3)$.\\
Do đó $Q_3=1{,}2+\dfrac{\dfrac{3\cdot 85}{4}-(8+28)}{32}\cdot (1{,}3-1{,}2)=\dfrac{1\,647}{1\,280}$.\\
Vậy $\Delta Q_1=Q_3-Q_1=\dfrac{1\,249}{8\,960}.$
\item Xét mẫu số liệu của giống lợn $B$.\\
Ta có $n=65$ nên tứ phân vị thứ nhất thuộc nhóm $[1{,}1;1{,}2)$.\\
Do đó $Q_1=1{,}1+\dfrac{\dfrac{65}{4}-13}{14}\cdot (1{,}2-1{,}1)=\dfrac{629}{560}$.\\
Tứ phân vị thứ ba thuộc nhóm $[1{,}2;1{,}3)$.\\
Do đó $Q_3=1{,}2+\dfrac{\dfrac{3\cdot 65}{4}-(13+14)}{24}\cdot (1{,}3-1{,}2)=\dfrac{413}{320}$.\\
Vậy $\Delta Q_2=Q_3-Q_1=\dfrac{75}{448}$.\\
Giá trị của $\Delta Q_2-\Delta Q_1\approx 0,03$.
\end{enumerate}
}
\end{ex}

\begin{ex}%[2D4V1-3]
Cho $F(x)$ là một nguyên hàm của $f(x)=2024-\sin^2 \dfrac{x}{2}$. Biết $F\left(\dfrac{\pi}{2}\right)=2025$. Tính $\sqrt{|F(0)|}$. (Làm tròn đến chữ số thập phân thứ nhất)
\shortans{$34$}
\loigiai{
Hàm số $f(x)=2024-\sin^2 \dfrac{x}{2}=2024-\dfrac{1-\cos x}{2}=\dfrac{4047+\cos x}{2}$.\\
Có $F(x)=\displaystyle \int f(x)\mathrm{\,d}x=\displaystyle \int\left(\dfrac{4047+\cos x}{2}\right)\mathrm{\,d}x=\dfrac{1}{2}(4047x+\sin x)+C$.\\
Do $F\left(\dfrac{\pi}{2}\right)=2025\Rightarrow 2025=\dfrac{1}{2}(4047\cdot \dfrac{\pi}{2} +\sin \dfrac{\pi}{2})+C\Rightarrow C=-\dfrac{4047}{4}\pi +\dfrac{4049}{2}$.\\
Suy ra $F(x)=\dfrac{1}{2}(4047x+\sin x)-\dfrac{4047}{4}\pi +\dfrac{4049}{2}$.\\
Vậy $\sqrt{|F(0)|}=34$.
}
\end{ex}

\begin{ex}%[2D4V1-3]
Cho $F(x)$ là một nguyên hàm của $f(x)=\sin^2 \dfrac{x}{4} \cdot \cos^2 \dfrac{x}{4}$. Biết $F\left(\dfrac{\pi}{3}\right)=0$. Tính giá trị của $F(\pi)$. (Làm tròn đến chữ số thập phân thứ hai)
\shortans{$0{,}37$}
\loigiai{
Hàm số $f(x)=\sin^2 \dfrac{x}{4} \cdot \cos^2 \dfrac{x}{4}=\dfrac{1}{8}(1-\cos x)$.\\
Có $F(x)=\displaystyle \int f(x)\mathrm{\,d}x=\displaystyle \int \dfrac{1}{8}(1-\cos x)\mathrm{\,d}x=\dfrac{1}{8}(x-\sin x)+C$.\\
Do $F\left(\dfrac{\pi}{3}\right)=0\Rightarrow 0=\dfrac{1}{8}\left(\dfrac{\pi}{3}-\sin \dfrac{\pi}{3}\right)+C\Rightarrow C=-\dfrac{\pi}{24}+\dfrac{\sqrt{3}}{16}$.\\
Suy ra $F(x)=\dfrac{1}{8}(x-\sin x)-\dfrac{\pi}{24}+\dfrac{\sqrt{3}}{16}$.\\
Vậy $F(\pi)=0{,}37$.
}
\end{ex}

\begin{ex}%[2D4V1-3]
Cho $F(x)$ là một nguyên hàm của hàm số $f(x)=\dfrac{1}{\cos^2 x}$. Biết $F\left(\dfrac{\pi}{4}+k\pi \right)=k$ với mọi $k\in \mathbb{Z}$. Tính giá trị của biểu thức $T=F(0)+F(\pi )+F(2\pi )+\cdots +F(10\pi )$.
\shortans{$44$}
\loigiai{
Ta có $\displaystyle \int f(x) \mathrm{\,d}x=\displaystyle \int \dfrac{\mathrm{\,d}x}{\cos^2 x}=\tan x+C$.\\
Suy ra
$F(x)=\heva{&\tan x+C_0, \quad x\in \left(-\dfrac{\pi}{2};\dfrac{\pi}{2}\right)\\ &\tan x+C_1, \quad x\in \left(\dfrac{\pi}{2};\dfrac{3\pi}{2}\right)\\ &\tan x+C_2, \quad x\in \left(\dfrac{3\pi}{2};\dfrac{5\pi}{2}\right)\\ &\cdots\\ &\tan x+C_9, \quad x\in \left(\dfrac{17\pi}{2};\dfrac{19\pi}{2}\right)\\ &\tan x+C_{10}, \quad x\in \left(\dfrac{19\pi}{2};\dfrac{21\pi}{2}\right)}\Rightarrow \heva{&F\left(\dfrac{\pi}{4}+0\pi\right)=1+C_0=0\Rightarrow C_0=-1\\ &F\left(\dfrac{\pi}{4}+\pi\right)=1+C_1=1\Rightarrow C_1=0\\ &F\left(\dfrac{\pi}{4}+2\pi\right)=1+C_2=2\Rightarrow C_2=1\\ &\cdots \\ &F\left(\dfrac{\pi}{4}+9\pi\right)=1+C_9=9\Rightarrow C_9=8\\ &F\left(\dfrac{\pi}{4}+10\pi\right)=1+C_{10}=0\Rightarrow C_{10}=9.}$\\
Vậy \begin{eqnarray*}
T&=&F(0)+F(\pi )+F(2\pi )+\cdots +F(10\pi )\\ &=&\tan 0-1+\tan \pi +\tan 2\pi +1+\cdots +\tan 10\pi +9\\ &=&44.
\end{eqnarray*}
}
\end{ex}

\begin{ex}%[2D4V1-3]
Hàm số $f(x)$ có đạo hàm liên tục trên $\mathbb{R}$ và $f'(x)=2024-2\sin^2 \dfrac{x}{2}$, $\forall x$;\hfill \break $f\left(\dfrac{\pi}{2}\right)=\dfrac{2023\pi}{2}$. Tính giá trị của $f(0)$.
\shortans{$-1$}
\loigiai{
$f(x)=\displaystyle \int \left(2024-2\sin^2 \dfrac{x}{2}\right)\mathrm{\,d}x=\displaystyle \int (2023+\cos x)\mathrm{\,d}x=2023x+\sin x+C$.\\
Tìm được $f(x)=2023x+\sin x+C$.\\
Do $f\left(\dfrac{\pi}{2}\right)=\dfrac{2023\pi}{2} \Leftrightarrow \dfrac{2023\pi}{2}=2023\cdot \dfrac{\pi}{2}+\sin \dfrac{\pi}{2}+C\Leftrightarrow C=-1$.\\
Vậy $f(x)=2023x+\sin x-1$.\\
Do đó $f(0)=-1$.
}
\end{ex}

\begin{ex}%[Hiếu Mai]%[2D4V1-3]
Gọi $F(x)$ là nguyên hàm của hàm số $f(x)=\sin^2x$ thỏa mãn $F(0)=0$. Tính $F(\pi)$ (làm tròn kết quả đến hàng phần mười).
\shortans{$1{,}6$}
\loigiai{
$F(x)=\displaystyle\int\limits f(x)\mathrm{\,d}x=\displaystyle\int\limits \sin^2x\mathrm{\,d}x=\displaystyle\int\limits \left( \dfrac{1}{2}-\dfrac{1}{2}\cos2x\right) \mathrm{\,d}x=\dfrac{x}{2}-\dfrac{\sin2x}{4}+C$.\\
Vì $F(0)=0$ nên $C=0$. Do đó, $F(x)=\dfrac{x}{2}-\dfrac{\sin2x}{4}$.\\
Khi đó, \[F(\pi)=\dfrac{\pi}{2}-\dfrac{\sin2\pi}{4}\approx 1{,}6\]
}
\end{ex}

\begin{ex}%[2D4V1-4]%[Tổ 1 - Đợt 17 - Chương 4 - - CD - Đề 2]%[Phùng Minh Phúc]
Cho hàm số $\text{y}=F\left( x \right)$ thỏa mãn $F\left( x \right)\left\langle 0,\forall x \right\rangle 0$ và có đạo hàm $f'\left( x \right)$ liên tục trên khoảng $\left( 0;+\infty  \right)$ thỏa mãn ${f}'\left( x \right)=\left( 2x+1 \right){{f}^{2}}\left( x \right),\forall x>0$ và $f\left( 1 \right)=-\dfrac{1}{2}$. Tính $f\left( 1 \right)+f\left( 2 \right)+\ldots +f\left( 2024 \right)$.
\shortans[oly]{}
\loigiai{
Ta có $f’\left( x \right)=\left( 2x+1 \right){{F}^{2}}\left( x \right)\Leftrightarrow \dfrac{f’\left( x \right)}{{{F}^{2}}\left( x \right)}=2x+1$.\\
$\Rightarrow \displaystyle\int\limits \dfrac{f’\left( x \right)}{{{F}^{2}}\left( x \right)}\text{d}x=\displaystyle\int\limits \left( 2x+1 \right)\text{d}x\Rightarrow -\dfrac{1}{F\left( x \right)}={{x}^{2}}+x+C$.\\
Mà $F\left( 1 \right)=-\dfrac{1}{2}\Rightarrow C=0\Rightarrow F\left( x \right)=\dfrac{-1}{{{x}^{2}}+x}=\dfrac{1}{x+1}-\dfrac{1}{x}$.\\
$\heva{
&F\left( 1 \right)=\dfrac{1}{2}-1  \\
&F\left( 2 \right)=\dfrac{1}{3}-\dfrac{1}{2}  \\
&F\left( 3 \right)=\dfrac{1}{4}-\dfrac{1}{3}  \\
&\vdots   \\
&F\left( 2024 \right)=\dfrac{1}{2025}-\dfrac{1}{2024}  \\
}$
$\Rightarrow F\left( 1 \right)+F\left( 2 \right)+\ldots +F\left( 2024 \right)=-1+\dfrac{1}{2025}=-\dfrac{2024}{2025}.$
}
\end{ex}

\begin{ex}%[2D4V1-4]
Cho hàm số $f(x)$ thỏa mãn $ f(0)=1-\ln 2$ và $\mathrm{e}^ xf'(x)=2^x\left[f(x)\right]^2$ với mọi $x\in\mathbb{R}$. Giá trị của $f(1)$ bằng bao nhiêu? (\textit{Kết quả làm tròn đến hàng phần trăm}).
\shortans{$0{,}42$}
\loigiai{
Từ giả thiết ta có $f'(x)=\dfrac{2^x}{\mathrm{e}^ x}{\left[f(x)\right]^2}$ với mọi $ x\in\left(1;2\right]$.\\
Do đó $ f(x)\ge f(1)=1>0$ với mọi $ x\in\left[1;2\right]$.\\
Xét với mọi $ x\in [1 ; 2]$ ta có
\begin{align*}
\mathrm{e}^ x{f}'(x)=2^x{\left[f(x)\right]^2}&\Rightarrow{f}'(x)=\dfrac{2^x}{\mathrm{e}^ x}{\left[f(x)\right]^2}\\&\Rightarrow\dfrac{f'(x)}{\left[f(x)\right]^2}=\left(\dfrac{2}{\mathrm{e}}\right)^x\\&\Rightarrow-\left(\dfrac{1}{f(x)}\right)'=\left(\dfrac{2}{\mathrm{e}}\right)^x \\&\Rightarrow{\left(\dfrac{1}{f(x)}\right)'}=-\left(\dfrac{2}{\mathrm{e}}\right)^x\\&\Rightarrow\dfrac{1}{f(x)}=-\displaystyle\int\left(\dfrac{2}{\mathrm{e}}\right)^x\mathrm{\,d} x\\&\Rightarrow\dfrac{1}{f(x)}=-\dfrac{\left(\dfrac{2}{\mathrm{e}}\right)^x}{\ln\dfrac{2}{\mathrm{e}}}+C\\&\Rightarrow\dfrac{1}{f(x)}=\dfrac{\left(\dfrac{2}{\mathrm{e}}\right)^x}{1-\ln 2}+C.
\end{align*}
Mà $ f(0)=1-\ln 2\Rightarrow C=0$. \\Do đó
$\dfrac{1}{f(x)}=\dfrac{\left(\dfrac{2}{\mathrm{e}}\right)^x}{1-\ln 2}$
$\Rightarrow f(x)=\dfrac{1-\ln 2}{\left(\dfrac{2}{\mathrm{e}}\right)^x}=\dfrac{(1-\ln 2)\mathrm{e}^x}{2^x}$.\\
Vậy $ f(1)=\dfrac{\mathrm{e}-\mathrm{e}\ln 2}{2}\approx 0{,}42$.}
\end{ex}

\begin{ex}%[2D4V1-4]
Cho hàm số $ y=f(x)$ đồng biến và có đạo hàm liên tục trên $\mathbb{R}$ thỏa mãn $\left(f'(x)\right)^2=f(x)\cdot\mathrm{e}^x$, $\forall x\in\mathbb{R}$ và $f(0)=2$. Tính $ f(2)$. (Kết quả làm tròn đến hàng phần trăm).
\shortans{$9{,}81$}
\loigiai{
Vì hàm số $ y=f(x)$ đồng biến và có đạo hàm liên tục trên $\mathbb{R}$ đồng thời $ f(0)=2$ nên $f'(x)\ge 0$ và $ f(x)>0$ với mọi $ x\in\left[0;+\infty\right)$.\\
Từ giả thiết $\left(f'(x)\right)^2=f(x)\cdot \mathrm{e}^x$, $\forall x\in\mathbb{R}$ suy ra $f'(x)=\sqrt{f(x)}\cdot\mathrm{e}^{\tfrac{x}{2}}$, $\forall x\in\left[0;+\infty\right).$\\
Do đó $\dfrac{f'(x)}{2\sqrt{f(x)}}=\dfrac{1}{2}{\mathrm{e}^{\tfrac{x}{2}}}$, $\forall x\in\left[0;+\infty\right).$\\
Lấy nguyên hàm hai vế, ta được $\sqrt{f(x)}=e^{\tfrac{x}{2}}+C$, $\forall x\in\left[0;+\infty\right)$ với $C$ là hằng số.\\
Kết hợp với $ f(0)=2$, ta được $C=\sqrt{2}-1$.\\
Suy ra $ f(2)=\left(\mathrm{e}+\sqrt{2}-1\right)^2\approx 9{,}81$.}
\end{ex}

\begin{ex}%[2D4V1-4]
Cho hàm số $ f(x)$ nhận giá trị dương và thỏa mãn $ f(0)=1$, $\left(f'(x)\right)^3=\mathrm{\mathrm{e}}^ x{\left(f(x)\right)^2}$, $\forall x\in\mathbb{R}$. Tính $ f(3)$ (\textit{kết quả làm tròn đến hàng phần mười}).
\shortans{$20{,}1$}
\loigiai{
Ta có

\begin{align*}
\left(f'(x)\right)^3=\mathrm{e}^x{\left(f(x)\right)^2},\,\forall x\in\mathbb{R}
&\Leftrightarrow{f}'(x)=\sqrt[3]{\mathrm{e}^x}\cdot \sqrt[3]{\left(f(x)\right)^2}\Leftrightarrow\dfrac{f'(x)}{\sqrt[3]{\left(f(x)\right)^2}}=\sqrt[3]{\mathrm{e}^x}\\
& \Leftrightarrow\dfrac{f'(x)}{\sqrt[3]{\left(f(x)\right)^2}}=\sqrt[3]{\mathrm{e}^x}\Leftrightarrow{f}'(x)\cdot \left(f(x)\right)^{-\tfrac{2}{3}}=\sqrt[3]{\mathrm{e}^x}\\&\Leftrightarrow 3\left[\left(f(x)\right)^{\tfrac{1}{3}}\right]'=\sqrt[3]{\mathrm{e}^x}\Leftrightarrow{\left[\left(f(x)\right)^{\tfrac{1}{3}}\right]'}=\dfrac{1}{3}\sqrt[3]{\mathrm{e}^x}\\&\Leftrightarrow{\left(f(x)\right)^{\tfrac{1}{3}}}=\dfrac{1}{3}\displaystyle\int{\sqrt[3]{\mathrm{e}^x}}\mathrm{\,d} x \Leftrightarrow{\left(f(x)\right)^{\tfrac{1}{3}}}=e^{\tfrac{x}{3}}+C.
\end{align*}
Vì	$f(0)=1$ nên $1=1+C\Rightarrow C=0\Rightarrow{\left(f(x)\right)^{\tfrac{1}{3}}}=e^{\tfrac{x}{3}}\Rightarrow f(x)=\mathrm{e}^x$.\\
Vậy	$f(3)=e^3\approx 20{,}1$.
}
\end{ex}

\begin{ex}%[2D4V1-4]
Gọi $F(x)$ là một nguyên hàm của hàm số $f(x)=2^x$, thỏa mãn $F(0)=\dfrac{1}{\ln2}$. Giá trị biểu thức $T=F(0)+F(1)+\cdots +F(2018)+F(2019)$ có dạng $\dfrac{{2^{2020}}+a}{\ln b}$. Giá trị của $\dfrac{a}{b}$ là
\shortans{$-0{,}5$}
\loigiai{
Ta có $\displaystyle \int f(x) \mathrm{\,d}x=\displaystyle \int 2^x \mathrm{\,d}x=\dfrac{2^x}{\ln2}+C$.\\
$F(x)$ là một nguyên hàm của hàm số $f(x)=2^x$, ta có $F(x)=\dfrac{2^x}{\ln2}+C$ mà $F(0)=\dfrac{1}{\ln2}$.\\
$\Rightarrow C=0\Rightarrow F(x)=\dfrac{2^x}{\ln2}$.\\
\begin{eqnarray*}
T&=&F(0)+F(1)+\cdots +F(2018)+F(2019)\\
&=&\dfrac{1}{\ln2}(1+2+2^2+\cdots +2^{2018}+2^{2019})\\
&=&\dfrac{1}{\ln2}\cdot \dfrac{{2^{2020}}-1}{2-1}\\
&=&\dfrac{{2^{2020}}-1}{\ln2}.\\
\end{eqnarray*}
Vậy $\dfrac{a}{b}=-\dfrac{1}{2}=-0{,}5$
}
\end{ex}

\begin{ex}%[2D4V1-4]%[Dự án EX-TF-TLN lần 3 -Mui Doan]
Gọi $F(x)=\dfrac{a}{b}\left(x^2+5\right)^c$ là một nguyên hàm của hàm số $f(x)=x\sqrt{x^2+5}$, trong đó $\dfrac{a}{b}$ tối giản và $a$, $b$ nguyên dương, $c$ là số hữu tỉ. Khi đó $a+b+c$ bằng bao nhiêu?
\shortans{$4{,}5$}
\loigiai{
Ta có $f(x)=F'(x)\Leftrightarrow x\sqrt{x^2+5}=\dfrac{a}{b}\cdot c\left(x^2+5\right)^{c-1}\cdot 2x\Leftrightarrow x\left(x^2+5\right)^{\frac{1}{2}}=\dfrac{2ac}{b}x\left(x^2+5\right)^{c-1}$.\\
Đồng nhất hệ số $2$ vế ta có hệ
\begin{align*}
\heva{& \dfrac{2ac}{b}=1 \\ & c-1=\dfrac{1}{2}}\Leftrightarrow \heva{& \dfrac{3a}{b}=1 \\ & c=\dfrac{3}{2}}\Leftrightarrow \heva{& \dfrac{a}{b}=\dfrac{1}{3} \\ & c=\dfrac{3}{2}}\Leftrightarrow \heva{& a=1 \\ & b=2\\& c=\dfrac{3}{2}.}
\end{align*}
Vậy $a+b+c=1+2+\dfrac{3}{2}=\dfrac{9}{2}=4{,}5$.
}
\end{ex}

\begin{ex}%[Mức độ 3]%[BG-12-New-4in1, Phạm Đức Thiệu]%[2D4V1-4]
Cho $F(x)$ là nguyên hàm của hàm số $f(x)=\dfrac{1}{\mathrm{e}^x+1}$ và $F(0)=-\ln(2\mathrm{e})$. Tìm nghiệm của phương trình $F(x)+\ln\left(\mathrm{e}^x+1\right)=2$.
\shortans{$3$}
\loigiai{
Ta có $F(x)=\displaystyle\int f(x)\mathrm{\,d}x
=\displaystyle\int\dfrac{1}{\mathrm{e}^x+1}\mathrm{\,d}x
=\displaystyle\int\left(1-\dfrac{\mathrm{e}^x}{\mathrm{e}^x+1}\right)\mathrm{d}x
=x-\ln\left(\mathrm{e}^x+1\right)+C$.\\
Theo bài ra $F(0)=-\ln (2\mathrm{e})\Leftrightarrow-\ln 2+C=-\ln (2\mathrm{e})
\Leftrightarrow-\ln 2+C=-\ln \mathrm{2}-1\Leftrightarrow C=-1$.\\
Do đó $F(x)=x-\ln\left(\mathrm{e}^x+1\right)-1$.\\
Khi đó, phương trình $F(x)+\ln\left(\mathrm{e}^x+1\right)=2
\Leftrightarrow x-\ln\left(\mathrm{e}^x+1\right)-1+\ln\left(\mathrm{e}^x+1\right)=2
\Leftrightarrow x-1=2
\Leftrightarrow x=3$.\\
Vậy nghiệm của phương trình là $x=3$.
}
\end{ex}

\begin{ex} %[Mức độ 3]%[BG-12-New-4in1, Phạm Đức Thiệu]%[2D4V1-4]
Cho hàm số $f(x)>0$ xác định trên $\mathbb{R}$ đồng thời thỏa mãn $f(0)=\dfrac{1}{2}$ và \break $f'(x)=-\mathrm{e}^x\cdot f^2(x)$,  $\forall x\in\mathbb{R}$. Tính giá trị của $f\left(\ln 2\right)$ (làm tròn kết quả đến hàng phần trăm).

\shortans{$0{,}33$}
\loigiai{
Ta có $f'(x)=-\mathrm{e}^x\cdot f^2(x)
\Leftrightarrow \dfrac{f'(x)}{f^2(x)}=-\mathrm{e}^x$ (do $f(x)>0$).\\
Lấy nguyên hàm 2 vế ta được
\begin{eqnarray*}
& & \displaystyle\int \dfrac{f'(x)}{f^2(x)} \mathrm{\,d}x=-\int \mathrm{e}^x \mathrm{\,d}x\\
&\Leftrightarrow& -\dfrac{1}{f(x)}=-\mathrm{e}^x+C\\
&\Leftrightarrow& f(x)=\dfrac{1}{\mathrm{e}^x-C}.
\end{eqnarray*}
Mà $f(0)=\dfrac{1}{2}
\Leftrightarrow \dfrac{1}{\mathrm{e}^0-C}=\dfrac{1}{2}
\Leftrightarrow 1-C=2
\Leftrightarrow C=-1$.\\
Do đó $f(x)=\dfrac{1}{\mathrm{e}^x+1}$.\\
Vậy $f(\ln 2)=\dfrac{1}{\mathrm{e}^{\ln 2}+1}
=\dfrac{1}{2+1}=\dfrac{1}{3}\simeq 0{,}33$.
}
\end{ex}

\begin{ex}%[2D4V1-4]
$F(x)$ là một nguyên hàm của hàm số $f(x)=2^x$, thỏa mãn $ F(0)=\dfrac{1}{\ln 2}$. Biểu thức $ F(0)+F(1)+F(2)+\ldots+F(2024)=\dfrac{a^b-c}{\ln a}$ $(a$, $b$, $c\in{N^*})$. Tính $ T=a+b-2c$.
\shortans{2025}
\loigiai{
Ta có $F(x)=\displaystyle\int 2^x \mathrm{\,d}x=\dfrac{2^x}{\ln 2}+C$.\\
Theo giả thiết $F(0)=\dfrac{1}{\ln 2}\Leftrightarrow\dfrac{2^0}{\ln 2}+C=\dfrac{1}{\ln 2}\Leftrightarrow C=0 \Rightarrow F(x)=\dfrac{2^x}{\ln 2}$.\\
Khi đó
\allowdisplaybreaks
\begin{eqnarray*}
F(0)+F(1)+F(2)+\ldots+F(2024)&=&\dfrac{2^0}{\ln 2}+\dfrac{2^1}{\ln 2}+\dfrac{2^2}{\ln 2}+\ldots+\dfrac{2^{2024}}{\ln 2}\\
&=&\dfrac{1}{\ln 2}(2^0+2^1+2^2+\ldots+2^{2024})\\
&=&\dfrac{1}{\ln 2}\cdot \dfrac{1(1-2^{2025})}{1-2}=\dfrac{2^{2025}-1}{\ln 2}.
\end{eqnarray*}
$\Rightarrow a=2$, $b=2025$, $c=1$.\\
Vậy $ T=a+b-2c=2025$.}
\end{ex}

\begin{ex}%[Hiếu Mai]%[2D4V1-5]
Gọi $F(x)$ là nguyên hàm của hàm số $f(x)=\dfrac{x}{\sqrt{8-x^2}}$ thỏa mãn $F(2)=0$. Tìm nghiệm của phương trình $F(x)=x$ (làm tròn kết quả đến hàng phần chục).
\shortans{$-0{,}7$}
\loigiai{
Ta có $\displaystyle\int{\dfrac{x}{\sqrt{8-x^2}}\mathrm{d}x}=-\dfrac{1}{2}\displaystyle\int{\left(8-x^2\right)^{-\tfrac{1}{2}}\mathrm{d}\left(8-x^2\right)}=-\sqrt{8-x^2}+C.$\\
Mặt khác $F(2)=0\Leftrightarrow-\sqrt{8-2^2}+C=0\Leftrightarrow C=2$.\\
Nên $F(x)=-\sqrt{8-x^2}+2.$\\
Suy ra
\begin{eqnarray*}
& &F(x)=x\Leftrightarrow-\sqrt{8-x^2}+2=x
\Leftrightarrow\sqrt{8-x^2}=2-x\\
&\Leftrightarrow& \heva{&2-x\ge 0\\&8-x^2=(2-x)^2}
\Leftrightarrow \heva{&x\le 2\\				&-2x^2+4x+4=0}
\Leftrightarrow \heva{&x\le 2\\
&\hoac{&x=1+\sqrt{3}\\&x=1-\sqrt{3}}}\\
&\Leftrightarrow& x=1-\sqrt{3}\simeq -0{,}7.
\end{eqnarray*}
Vậy nghiệm của phương trình đã cho là
$x \simeq -0{,}7$.
}
\end{ex}

\begin{ex}%[Hiếu Mai]%[2D4V1-5]
Cho hàm số $ f(x)$ xác định trên $\mathbb{R}\setminus\left\{\dfrac{1}{2}\right\}$ thỏa mãn $f'(x)=\dfrac{2}{2x-1}$; $f(0)=1$; $f(1)=2$. Tính giá trị của biểu thức $f(-1)+f(3)$ (làm tròn kết quả đến hàng phần trăm).
\shortans{$5{,}71$}
\loigiai{
Ta có $\displaystyle\int f'(x)\mathrm{\,d}x
=\displaystyle\int\dfrac{2}{2x-1}\mathrm{\,d}x
=\displaystyle\int\dfrac{1}{2x-1}\mathrm{\,d}(2x-1)
=\ln|2x-1|+C=f(x)$.\\
Theo bài ra $f(0)=1\Rightarrow C=1$ nên $f(-1)=1+\ln 3$ và $f(1)=2\Rightarrow C=2$ nên $f(3)=2+\ln 5$.\\
Do đó $f(-1)+f(3)=3+\ln 15\simeq 5{,}71$.
}
\end{ex}

\begin{ex}%[2D4V1-6]%[Tổ 19-Đợt 17-Chương 4-CD-Đề 1]%[Thinh Ngô]
Tại một lễ hội dân gian hàng năm, tốc độ thay đổi lượng khách tham dự được biểu diễn bằng hàm số $Q'(t)=8t^3-144t^2+576t$, trong đó $t$ tính bằng giờ $(0\leq t \leq 14)$, $Q'(t)$ tính bằng khách/giờ. Sau $1$ giờ đã có $300$ người có mặt. Hỏi số lượng khách tham dự đông nhất trong vòng $14$ là bao nhiêu?
\shortans{$2650$}
\loigiai{
Ta có $\displaystyle\int\limits Q'(t)\mathrm{\,d}t=\displaystyle\int\limits \left(8t^3-144t^2+576t\right)\mathrm{\,d}t=2t^4-48t^3+288t^2+C$.\\
Vậy số lượng khách tham quan được biểu diễn bằng hàm số $Q(t)=2t^4-48t^3+288t^2+C$, với $(0\leq t\leq 14)$.\\
Ta có $Q(1)=300 \Rightarrow C=58$, suy ra $Q(t)=2t^4-48t^3+288t^2+58$.\\
$Q'(t)=8t^3-144t^2+576t=0 \Leftrightarrow \hoac{& t=0 \\ & t=6\\&t=12.}$\\
Bảng biến thiên
\begin{center}
\begin{tikzpicture}
\tkzTabInit[espcl=2.5,lgt=1.5,nocadre]
{$t$/0.7,$Q'(t)$/0.7,$Q(t)$/2.1}
{$0$,$6$,$12$,$14$}
\tkzTabLine{,+,0,-,0,+,}
\tkzTabVar{-/$58$,+/$2650$,-/$58$,+/$1626$}
\end{tikzpicture}
\end{center}
Vậy số lượng khách tham dự đông nhất trong vòng $14$ giờ là $2650$ khách.
}
\end{ex}

\begin{ex}%[2D4V1-6]%[Tổ 19 - Đợt 17 - Chương 4 - - CD - Đề 2]%[Bình]
Một viên đạn được bắn lên trời với vận tốc là $72$~m/s bắt đầu từ độ cao $2$~m. Hãy xác định chiều cao của viên đạn sau thời gian $5$ giây kể từ lúc bắn biết gia tốc trọng trường là $9{,}8$~m/s$^2$. (Làm tròn đến hàng đơn vị)
\shortans{$240$}
\loigiai
{
Ta có vận tốc của viên đạn tại thời điểm $t$ là \[v(t)=\displaystyle\int -9{,}8\mathrm{\,d}t=-9{,}8t+C_1.\]
Do $v(0)=72$ nên $v(0)=-9{,}8\cdot 0+C_1=72\Leftrightarrow C_1=72\Rightarrow v(t)=-9{,}8t+72$.\\
Độ cao của viên đạn tại thời điểm $t$ là \[s(t)=\displaystyle\int v(t)\mathrm{\,d}t=\displaystyle\int (-9{,}8t+72)\mathrm{\,d}t=-4{,}9t^2+72t+C_2.\]
Vì $s(0)=2$ nên $s(0)=-4{,}9\cdot 0^2+72\cdot 0+C_2\Leftrightarrow C_2=2\Rightarrow s(t)=-4{,}9t^2+72t+2$.\\
Vậy sau khoảng thời gian $5$ giây kể từ lúc bắn, viên đạn ở độ cao\\ \centerline{$s(5)=-4{,}9\cdot 5^2+72\cdot 5+2=239{,}5 \approx 240$~m.}
}
\end{ex}

\begin{ex}%[2D4V1-6]%[Tổ 19 - Đợt 17 - Chương 4 - - CD - Đề 2]%[Bình]
Một bác thợ xây bơm nước vào bể chứa nước. Gọi $h(t)$ là thể tích nước bơm được sau $t$ giây. Cho $h'(t)=3at^2+bt$~(m$^3$/s) và ban đầu bể không có nước. Sau $5$ giây thì thể tích nước trong bể là $150$~m$^3$. Sau	$10$ giây thì thể tích nước trong bể là $1100$~m$^3$. Hỏi thể tích nước trong bể sau khi bơm được $20$ giây.
\shortans{$8400$}
\loigiai
{
Ta có $h'(t)=3at^2+bt$~(m$^3$/s).\\
$\Rightarrow h(t)=\displaystyle\int (3at^2+bt)\mathrm{\,d}t=at^3+\dfrac{1}{2}bt^2+C$.\\
Chọn $t=0\Rightarrow h(0)=0\Rightarrow C=0\Rightarrow h(t)=at^3+\dfrac{1}{2}bt^2$.\\
Sau $5$ giây thì thể tích nước trong bể là $150$~m$^3$, ta có\\ \centerline{$h(5)=150\Leftrightarrow 125a+\dfrac{25}{2}b=150.$}
Sau $10$ giây thì thể tích nước trong bể là $1100$~m$^3$, ta có\\
\centerline{$h(10)=1100\Leftrightarrow 1000a+50b=1100.$}
Ta có hệ $\heva{&125a+\dfrac{25}{2}b=150\\ &1000a+50b=1100}\Leftrightarrow \heva{&a=1\\ &b=2.}$\\
$\Rightarrow h(t)=t^3+t^2$.\\
Thể tích nước trong bể sau khi bơm được $20$ giây là $h(20)=20^3+20^2=8400$~m$^3$.
}
\end{ex}

\begin{ex}%[2D4V1-6]
\immini{Một vật chuyển động trong $4$ giờ với vận tốc $v$ (km / h) phụ thuộc vào thời gian $t$ (h) có đồ thị vận tốc là một đường parabol có đỉnh $I(3 ; 10)$ và trục đối xứng vuông góc với trục hoành như hình vẽ. Tính quãng đường vật di chuyển được trong nửa thời gian sau của chuyển động đó (kết quả làm tròn đến hàng phần chục và tính theo đơn vị km )}
{
\begin{tikzpicture}[scale=0.4,>=stealth, font=\footnotesize, line join=round, line cap=round]
\def\xmin{-3} \def\xmax{6}
\def\ymin{-1} \def\ymax{11}
\draw[->] (\xmin,0)--(\xmax,0) node [below]{$t$};
\draw[->] (0,\ymin)--(0,\ymax) node [left]{$v$};
\fill (0,0) circle (1pt) node[shift={(-135:2.5mm)}]{$O$};
\clip (\xmin+0.1,\ymin+0.1) rectangle (\xmax-0.1,\ymax-0.1);
\draw[smooth,red,samples=300,domain=(0:4)] plot(\x,{-(\x)^2+6*(\x)+1});
\foreach \x in {\xmin,...,\xmax}
\draw (\x,-0.05)--(\x,0.05);
\foreach \y in {\ymin,...,\ymax}
\draw (-0.05,\y)--(0.05,\y);
\node at (3,0)[below]{$3$};
\node at (0,1)[left]{$1$};
\node at (0,10)[left]{$10$};
\draw[dashed] (0,10)--(3,10)--(3,0);
%\node at (-1,1)[shift={(70:2mm)}]{$ I $};
%\draw[dashed]
%(-1,0)node[below]{$ -1 $}|-(0,1)node[right]{$1$}
%	;
\end{tikzpicture}
}
\shortans{$19{,}3$}

\loigiai{
Dựa vào đồ thị ta tìm được vận tốc $v(t)=-t^{2}+6 t+1, t \in[0 ; 4]$.\\
Quãng đường $s(t)$ vật di chuyển được trong thời gian $4$h là một nguyên hàm của $v(t)$, với $t \in[0 ; 4]$\\
Ta có $s(t)=\displaystyle\int\left(-t^{2}+6t+1\right) \mathrm{d}t=-\dfrac{t^{3}}{3}+3 t^{2}+t+C$.\\
Quãng đường vật di chuyển được trong $2$h đầu là $s(2)=\dfrac{34}{3}+C$ (km).\\
Quãng đường vật di chuyển được trong $4$h là $s(4)=\dfrac{92}{3}+C$ (km).\\
Quãng đường vật di chuyển được trong nửa thời gian sau của chuyển động là
\[s(4)-s(2)=\dfrac{58}{3} \text{~(km)} \approx 19{,}3\text{~(km)}.\]}

\end{ex}

\begin{ex}%[2D4V1-6]%[Tổ 1 - Đợt 17 - Chương 4 - - CD - Đề 2]%[Phùng Minh Phúc]
Tại một khu di tích vào ngày lễ hội, người ta tính được tốc độ thay đổi lượng khách tham quan được biểu diễn bằng hàm số $\text{{Q}'}\left( t \right)=4{{t}^{3}}-72{{t}^{2}}+288t$, trong đó $t$ tính bằng giờ $\left( 0\le t\le 13 \right),Q'\left( t \right)$ tính bằng khách/giờ. Sau $2$ giờ đã có 500 người có mặt.
\begin{enumerate}[a)]
\item Xác định hàm số $Q\left( t \right)$ biểu diễn lượng khách tham quan di tích.
\item Xác định thời điểm mà lượng khách tham quan lớn nhất.
\item Tìm thời điểm mà tốc độ thay đổi lượng khách tham quan là lớn nhất?
\end{enumerate}
\shortans[oly]{}
\loigiai{
\begin{enumerate}[a)]
\item Ta có $Q\left( t \right)=\displaystyle\int\limits Q'\left( t \right)\text{d}t=\displaystyle\int\limits \left( 4{{t}^{3}}-72{{t}^{2}}+288t \right)\text{d}t={{t}^{4}}-24{{t}^{3}}+144{{t}^{2}}+C$.\\
Mà sau $2$ giờ đã có $500$ người nên ta có $Q\left( 2 \right)=500$ suy ra $C=100$.\\
Vậy $Q\left( t \right)={{t}^{4}}-24{{t}^{3}}+144{{t}^{2}}+100$.
\item Ta tìm giá trị lớn nhất của hàm số $Q\left( t \right)$ trên đoạn $\left[ 0;13 \right]$.\\
Ta có ${Q}'\left( t \right)=0$ khi $t=0,t=6$ và $t=12$.\\
Mà $Q\left( 0 \right)=100,Q\left( 6 \right)=1\,396,Q\left( 12 \right)=100,Q\left( 13 \right)=269$.\\
Nên lượng khách tham quan lớn nhất là sau $6$ giờ, có $1\,396$ người.
\item Khảo sát hàm số ${Q}'\left( t \right)=4{{t}^{3}}-72{{t}^{2}}+288t$ trên đoạn $\left[ 0;13 \right]$.\\
Ta có ${Q}\rq\rq \left( t \right)=12{{t}^{2}}-144t+288$.\\
${Q}\rq\rq \left( t \right)=0\Leftrightarrow 12{{t}^{2}}-144t+288=0\Leftrightarrow t=6-2\sqrt{3}$ hoặc $t=6+2\sqrt{3}$.\\
Bảng biến thiên của hàm số ${Q}'\left( t \right)$ như sau:
\begin{center}
\begin{tikzpicture}
\tkzTabInit[lgt=2,espcl=4]
{$t$/0.7,$Q\rq\rq (t)$/0.7,$Q'(t)$/3}
{$0$,$6-2\sqrt{3}$,$6+2\sqrt{3}$,$13$}
\tkzTabLine{,+,0,-,0,+}
\tkzTabVar{-/$0$,+/$Q^{\prime}(6-2\sqrt{3})$,-/$Q^{\prime}(6+2\sqrt{3})$,+/$364$}
\end{tikzpicture}
\end{center}
Với ${Q}'\left( 6-2\sqrt{3} \right)\approx 332,6$ và ${Q}'\left( 6+2\sqrt{3} \right)\approx -332{,}6$.\\
Vậy tốc độ thay đổi lượng khách tham quan lớn nhất tại thời điểm $t=13$.
\end{enumerate}
}
\end{ex}

\begin{ex}%[Tổ 20 - Đợt 17 - Chương 4 - - CD - Đề 8]%[Huỳnh Xuân Tín]%[2D4V1-6]
Một vật chuyển động với gia tốc $a(t)=\dfrac{1}{t^2+3 t+2} \left(\mathrm{m} / \mathrm{s}^2\right)$, trong đó $t$ là khoảng thời gian tính từ thời điểm ban đầu. Vận tốc chuyển động của vật là $v(t)$. Vào thời điểm $t=11(\mathrm{~s})$ thì vận tốc của vật là $v(\mathrm{~m} / \mathrm{s})$, biết vận tốc ban đầu của vật là $v_0=3 \ln 2(\mathrm{~m} / \mathrm{s})$. Giá trị của $v$ (làm tròn đến hàng phần trăm) là
\shortans{$2{,}69$}
\loigiai
{Ta có $v(t)=\displaystyle\int\limits a(t) \mathrm{d}t=\displaystyle\int\limits\dfrac{1}{t^2+3 t+2} \mathrm{d}t=\displaystyle\int\limits\left(\dfrac{1}{t+1}-\dfrac{1}{t+2}\right) \mathrm{d}t=\ln \left|\dfrac{t+1}{t+2}\right|+C$.\\
$v(0)=\ln \left(\dfrac{1}{2}\right)+C=3 \ln 2 \Rightarrow C=4 \ln 2 \Rightarrow v(t)=\ln \left|\dfrac{t+1}{t+2}\right|+4 \ln 2$.\\
Ta có $v(11)=\ln \left(\dfrac{12}{13}\right)+4 \ln 2 \approx 2{,}69$.\\
Vậy $v \approx 2{,}69$.
}
\end{ex}

\begin{ex}%[Đề 12 - 2025 - Nguyen Son]%[2D4V1-6]
Một vật chuyển động có gia tốc là $a\left(t\right)=3 t^2+t$ (m/s$^2$). Biết rằng vận tốc ban đầu của vật là $2$ m/s. Vận tốc của vật đó sau $4$ giây là bao nhiêu (m/s)?
\shortans{$74$}
\loigiai{
Vận tốc của vật là $v(t)=\displaystyle \int \left(3 t^2+t\right)  \mathrm{\,d}t =t^3+\dfrac{1}{2}t^2+C$.\\
Vì vận tốc tốc ban đầu của vật là $2$ m/s nên $v(0)=2 \Leftrightarrow 0^3+\dfrac{1}{2}\cdot 0^2+C=2  \Leftrightarrow C=2$.\\
Suy ra $v(t)=t^3+\dfrac{1}{2}t^2+2$.\\
Vận tốc của vật đó sau $4$ giây là $v(4)=4^3+\dfrac{1}{2}\cdot 4^2+2=74$ (m/s).
}
\end{ex}

\begin{ex}%[2D4V1-6]
Một vật chuyển động có gia tốc là $a\left(t\right)=\dfrac{3}{t+1}$ (m/s$^2$). Biết rằng vận tốc ban đầu của vật là $6$ m/s. Vận tốc của vật đó sau $5$ giây là bao nhiêu m/s$^2$ (làm tròn kết quả đến hàng phần mười)?
\shortans{$13{,}2$}
\loigiai{
Vận tốc của vật là $v(t)=\displaystyle \int a\left(t\right)  \mathrm{\,d}t =\int \dfrac{3}{t+1}\mathrm{\,d}t=3\ln |t+1|+C$.\\
Vì vận tốc tốc ban đầu của vật là $6$ m/s nên $v(0)=6 \Leftrightarrow 3\ln 1+C=6  \Leftrightarrow C=6$.\\
Suy ra $v(t)=3\ln |t+1| +6$.\\
Vận tốc của vật đó sau $10$ giây là $v(10)=3\ln|10+1|+6\approx 13{,}2$ (m/s).
}
\end{ex}

\begin{ex}%[2D4V1-6]
Một vật được ném lên từ độ cao $300$ m với vận tốc được cho bởi công thức $v(t)=-9{,}81t+29{,}43$ (m/s) (\textit{Nguồn: R. Larson and B. Edwards, Calculus 10e, Cengane}). Gọi $h(t)$ (m) là độ cao của vật tại thời điểm $t$ (s). Sau bao lâu kể từ khi bắt đầu được ném lên thì vật đó chạm đất (làm tròn kết quả đến hàng phần chục giây)?
\shortans{$11{,}4$}\loigiai{
Ta có $h(t)=\displaystyle\int v(t) \mathrm{\,d} t=\displaystyle\int(-9{,}81 t+29{,}43) \mathrm{\,d} t=-\dfrac{9{,}81}{2} t^2+29{,}43 t+C$.\\
Vì vật được ném lên từ độ cao $300$ m nên $h(0)=300$. Suy ra $C=300$.\\
Vậy $h(t)=-\dfrac{9{,}81}{2} t^2+29{,}43 t+300$. Khi vật bắt đầu chạm đất ứng với $h(t)=0$.\\
Nên ta có $-\dfrac{9{,}81}{2} t^2+29{,}43 t+300=0 \Leftrightarrow t \approx 11{,}4$ hoặc $t \approx-5{,}4$.\\
Do $t>0$ nên $t \approx 11{,}4$ (s).
}
\end{ex}

\begin{ex}%[2D4V1-6]
Chủ một trung tâm thương mại muốn cho thuê một số gian hàng như nhau. Người đó muốn tăng giá cho thuê của mỗi gian hàng thêm $x$ (triệu đồng) $(x\ge 0)$. Tốc độ thay đổi doanh thu từ các gian hàng đó được biểu diễn bởi hàm số $T'(x)=-20x+300$, trong đó $T'(x)$ tính bằng triệu đồng (\textit{Nguồn: R. Larson and B. Edwards, Calculus ioe, Cengage}). Biết rằng nếu người đó tăng giá thuê cho mỗi gian hàng thêm $10$ triệu đồng thì doanh thu là $12\,000$ triệu đồng. Tìm giá trị của $x$ để người đó có doanh thu là cao nhất?
\shortans{$15$}\loigiai{
Ta có $T(x)=\displaystyle\int T'(x) \mathrm{\,d} x=\displaystyle\int(-20 x+300) \mathrm{\,d} x=-10 x^2+300 x+C$, $C \in \mathbb{R}$.\\
Khi người đó tăng giá cho thuê mỗi gian hàng thêm 10 triệu đồng thì doanh thu là $12\,000$ triệu đồng. \\
Nên ứng với $x=10$ ta có $T(10)=12\,000$ suy ra
\[12\,000=-10 \cdot 10^2+300 \cdot 10+C \Rightarrow C=10\,000.\]
Vậy $T(x)=-10 x^2+300 x+10\,000$. Ta có $T(x)$ là một hàm số bậc hai với hệ số $a<0$ và đồ thị hàm số có đỉnh là $I(15 ; 12\,250)$.\\
Vậy doanh thu cao nhất mà người đó có thể thu về là $12\,250$ triệu đồng và khi đó mỗi gian hàng đã tăng giá cho thuê thêm $15$ triệu đồng.
}
\end{ex}

\begin{ex}%[2D4V1-6]
Một xe ô tô đang chạy với tốc độ $90$ km/h thì người lái xe bất ngờ phát hiện chướng ngại vật trên đường cách đó $150$ m. Người lái xe phản ứng $2$ giây sau đó bằng cách đạp phanh cho xe chạy chậm hơn. Kể từ thời điểm này, ô tô chuyển động chậm dần đều với tốc độ $v(t)=-\dfrac{25}{4}t+25$(m/s), trong đó $t$ là thời gian tính bằng giây kể từ lúc đạp phanh. Quãng đường xe ô tô đã di chuyển kể từ lúc người lái xe phát hiện chướng ngại vật trên đường đến khi xe ô tô dừng hẳn là bao nhiêu mét?\\
\shortans{100}
\loigiai{
Gọi $s(t)$ là quãng đường xe ô tô đi được trong $t$ (giây) kể từ lúc đạp phanh.\\ Khi đó
$ s(t)=\displaystyle\int v(t) \mathrm{\,d}t=\displaystyle\int \left(-\dfrac{25}{4}t+25\right) \mathrm{\,d}t=-\dfrac{25}{8}{t^2}+25t+C$.\\
Do $s(0)=0$ nên $C=0$ . Suy ra $s(t)=-\dfrac{25}{8}{t^2}+25t$.\\
Xe ô tô dừng hẳn khi $ v(t)=0\Leftrightarrow-\dfrac{25}{4}t+25=0\Leftrightarrow t=4$.\\
Suy ra quãng đường xe ô tô còn di chuyển được kể từ lúc đạp phanh đến khi xe dừng hẳn là
$s(4)=-\dfrac{25}{8}{4^2}+25\cdot 4=50$ (m).\\
Ta có tốc độ $90$ km/h cũng là tốc độ $25$ m/s.\\
Do đó, quãng đường xe ô tô đã di chuyển kể từ lúc người lái xe phát hiện chướng ngại vật trên đường đến khi xe ô tô dừng hẳn là: $ 25\cdot 2+50=100$ (m).
}
\end{ex}

\begin{ex}%[2D4V1-6]
Một quần thể vi khuẩn ban đầu gồm $500$ vi khuẩn, sau đó bắt đầu tăng trưởng. Gọi $P(t)$ là số lượng vi khuẩn của quần thể đó tại thời điểm $t$, trong đó $ t$ tính theo ngày $(0\le t\le 10)$. Tốc độ tăng trưởng của quần thể vi khuẩn đó cho bởi hàm số $P'(t)=k\sqrt{t}$, trong đó k là hằng số. Sau 1 ngày, số lượng vi khuẩn của quần thể đó đã tăng lên thành 600 vi khuẩn (Nguồn: R. Larson and
Edwards, Calculus 10e, Cengage 2014). Tính số lượng vi khuẩn của quần thể đó sau $9$ ngày.

\shortans{3200}
\loigiai{
Ta có $P(t)=\displaystyle\int P'(t)\mathrm{\,d}t=\displaystyle\int k\sqrt{t} \mathrm{\,d}t=\displaystyle\int k\cdot t^{\tfrac{1}{2}}\mathrm{\,d}t=k\cdot\dfrac{2}{3}t\sqrt{t}+C$.\\
Từ giả thiết suy ra $\heva{&P(0)=500\\&P(1)=600
}\Rightarrow\heva{&k\cdot\dfrac{2}{3}\cdot0\sqrt 0+C=500\\&k\cdot\dfrac{2}{3}\cdot1\sqrt 1+C=600}\Rightarrow\heva{&C=500\\&
\dfrac{2}{3}k=100}\Rightarrow\heva{&C=500\\&k=150.}$\\
$\Rightarrow P(t)=100t\sqrt t+500$.\\
Do đó, số lượng vi khuẩn của quần thể đó sau $9$ ngày là $P(9)=100\cdot 9\sqrt{9}+500=3\,200$.}
\end{ex}

\begin{ex}%[2D4V1-6]
Cây cà chua khi trồng có chiều cao $5$ cm. Tốc độ tăng chiều cao của cây cà chua sau khi trồng được cho bởi hàm số $ v(t)=-0{,}1t^3+t^2$, trong đó $t$ tính theo tuần, $v(t)$ tính bằng cm/tuần. Gọi $h(t)$ (tính bằng centimét) là độ cao của cây cà chua ở tuần thứ $t$ (Nguồn:A. Bigalke et al., Mathematik, Grundkurs ma-I, Cornelsen 2016). Vào thời điểm cây cà chua đó phát triển nhanh nhất thì cây cà chua sẽ cao bao nhiêu? (làm tròn kết quả đến hàng phần chục).
\shortans{54{,}4}
\loigiai{
Ta có $h(t)=\displaystyle\int v(t) \mathrm{\,d}t=\displaystyle\int \left(-0{,}1t^3+t^2\right) \mathrm{\,d}t=-\dfrac{1}{40}{t^4}+\dfrac{t^3}{3}+C$.\\
Cây cà chua khi trồng có chiều cao $5$ cm nên $h(0)=5\Rightarrow C=5$.\\
Vậy độ cao của cây cà chua ở tuần thứ $t$ được cho bởi hàm số\\
\centerline {$h(t)=-\dfrac{1}{40}{t^4}+\dfrac{t^3}{3}+5$ $(t\ge 0)$.}\\
Ta tìm $ t$ $(t\ge 0)$ sao cho $v(t)$ đạt giá trị lớn nhất.\\
$v'(t)=-0{,}3t^2+2t$; $ v'(t)=0\Leftrightarrow-0{,}3t^2+2t=0\Leftrightarrow\hoac{&t=0\\&t=\dfrac{20}{3}.}$\\
Bảng biến thiên
\begin{center}
\begin{tikzpicture}
\tkzTabInit[espcl=2.5,lgt=1.5,nocadre]
{$x$/0.7,$y'$/0.7,$y$/2.1}
{$-\infty$,$0$,$\tfrac{20}{3}$,$+\infty$}
\tkzTabLine{,-,0,+,0,-,}
\tkzTabVar{+/$+\infty$,-/$0$,+/$\dfrac{400}{27}$,-/$-\infty$}
\end{tikzpicture}
\end{center}
Từ đó ta thấy $v(t)$ đạt giá trị lớn nhất tại $t=\dfrac{20}{3}$.\\
Khi đó, cây cà chua sẽ đạt chiều cao là $h\left(\dfrac{20}{3}\right)=\dfrac{4\,405}{81}\approx 54{,}4$ (cm).
}
\end{ex}

\begin{ex}%[12-EX-2-2025, Diamond]%[2D4V1-6]
Một ô tô chuyển động trên đường thẳng nằm ngang với gia tốc phụ thuộc thời gian $t$ (s) là $a(t) = 2t - 7$ (m/s$^2$). Biết vận tốc ban đầu bằng $10$ (m/s), hỏi sau bao lâu thì ô tô đạt vận tốc $18$ (m/s)?
\shortans{$8$}
\loigiai{
Ta có $v(t)=\displaystyle\int a(t)\mathrm{\,d}t=\int(2t-7)\mathrm{\,d}t=t^2-7t+C$ (m/s).\\
Với $v(0)=10\Leftrightarrow C=10$, suy ra $v(t)=t^2-7t+10$(m/s).\\
Với $v(t)=18\Leftrightarrow t^2-7t-8=0\Leftrightarrow\hoac{&t=-1\\&t=8.}$\\
Vậy sau $8$ giây thì vận tốc đạt $18$ m/s.
}
\end{ex}

\begin{ex}%[Dự án 12-EX-3-2026, Mui Doan]%[2D4V1-6]
Người ta truyền nhiệt cho một bình nuôi cấy vi sinh vật từ $1^\circ C$. Tốc độ tăng nhiệt độ của bình tại thời điểm $t$ phút ($0 \le t \le 5$) được cho bởi hàm số $f(t) = 3t^2$ ($^\circ C$/phút). Biết rằng nhiệt độ của bình đó tại thời điểm $t$ là một nguyên hàm của hàm số $f(t)$. Nhiệt độ của bình tại thời điểm $3$ phút kể từ khi truyền nhiệt là $a^\circ C$. Khi đó $a$ có giá trị bằng bao nhiêu?
\par\shortans[oly]{28}
\loigiai{
Ta có $F(t)=\displaystyle\int 3t^2\mathrm{\,d}t=t^3+C$.\\
Vì $F(0)=1$ nên $0^3+C=1\Leftrightarrow C=1$.\\
Suy ra $F(t)=t^3+1$.\\
Do đó  $F(3)=3^3+1=28^\circ C$.\\
Vậy $a=28$.
}
\end{ex}

\begin{ex}%[12-EX-2-2025, Đình Nguyên]%[2D4V1-6]
Khi nghiên cứu một quần thể vi khuẩn, người ta nhận thấy quần thể vi khuẩn đó ở ngày thứ $t$ có số lượng $N(t)$ con. Biết rằng tốc độ phát triển của quần thể đó là $N'(t) = \dfrac{8\,000}{t}$ và sau ngày thứ nhất ($t=1$) có $250$ nghìn con. Số lượng vi khuẩn sau $10$ ngày là bao nhiêu nghìn con (làm tròn đến hàng đơn vị).
\shortans{268}
\loigiai{
Ta có $N(t) = \displaystyle\int N'(t) \mathrm{\,d}t = \displaystyle\int \dfrac{8\,000}{t} \mathrm{\,d}t = 8\,000 \ln|t| + C$.\\
Vì tại thời điểm $t=1$, số lượng vi khuẩn là $250$ nghìn con nên \[N(1) = 250\,000 \Rightarrow 8\,000 \ln 1 + C = 250\,000 \Rightarrow C = 250\,000.\]
Suy ra hàm số biểu thị số lượng vi khuẩn là $N(t) = 8\,000 \ln|t| + 250\,000$.\\
Số lượng vi khuẩn sau $10$ ngày là $N(10) = 8\,000 \ln 10 + 250\,000 \approx 268\,421$ con $\approx 268$ nghìn con.
}
\end{ex}

\begin{ex}%[2D4V1-6]
Một người bơm nước vào một bể chứa nước. Gọi $h(t)$ là thể tích nước bơm được sau $t$ giây. Cho $h'(t)=6at^2+2bt$ và ban đầu bể không có nước. Sau $3$ giây thì thể tích nước trong bể là $90$ m$^3$, sau $6$ giây thì thể tích nước trong bể là $504$ m$^3$. Thể tích nước trong bể sau khi bơm được $9$ giây là bao nhiêu mét khối?
\par
\shortans[0]{1\,458}
\loigiai{
Ta có $\displaystyle\int\limits_{0}^{3} (6at^2+2bt) \,\mathrm{d}t = 90
\Leftrightarrow (2at^3+bt^2)\Bigg|_{0}^{3} = 90
\Leftrightarrow 54a+9b = 90. \hfill(1)$\\
Lại có $\displaystyle\int\limits_{0}^{6} (6at^2+2bt)\,\mathrm{d}t = 504
\Leftrightarrow (2at^3+bt^2)\Bigg|_{0}^{6} = 504
\Leftrightarrow 432a+36b = 504. \hfill(2)$\\
Từ $(1)$ và $(2)$ suy ra $a=\dfrac{2}{3}$ và $b=6$.\\
Sau khi bơm $9$ giây thì thể tích nước trong bể là
\[V = \displaystyle\int\limits_{0}^{9} (4t^2+12t)\,\mathrm{d}t = \left( \dfrac{4}{3}t^3 + 6t^2 \right)\Bigg|_{0}^{9} = 1\,458 \text{ (m}^3).\]
}
\end{ex}

\begin{ex}%[2D4V2-1]%[Tổ 20 - Đợt 17 - Chương 4 - - CD - Đề 2]%[Trần Thiên Đức]
Cho hàm số $f(x)$ liên tục trên $\mathbb{R}$. Gọi $F(x)$, $G(x)$, $H(x)$ là ba nguyên hàm của $f(x)$ trên $\mathbb{R}$ thỏa mãn $F(3)+G(3)+H(3)=4$ và $F(0)+G(0)+H(0)=1$. Tính $3\displaystyle\int\limits_0^1f(3x)\mathrm{\,d}x$.
\shortans[]{$1$}
\loigiai{
Ta có $I=\displaystyle\int\limits_0^1f(3x)\mathrm{\,d}x=\dfrac{1}{3}\displaystyle\int\limits_0^3f(t)\mathrm{\,d}t\Rightarrow3I=\displaystyle\int\limits_0^3f(x)\mathrm{\,d}x$.\\
Do $F(x)$, $G(x)$, $H(x)$ là ba nguyên hàm của $f(x)$ nên\\
$3I=F(3)-F(0)=G(3)-G(0)=H(3)-H(0)$.\\
$\Rightarrow9I=[F(3)+G(3)+H(3)]-[F(0)+G(0)+H(0)]=3$.\\
$\Rightarrow3I=1$.\\
Vậy $3\displaystyle\int\limits_0^1f(3x)\mathrm{\,d}x=1$
}
\end{ex}

\begin{ex}%[2D4V2-1]
Cho hàm số $f\left( x \right)$ liên tục trên $\mathbb{R}$. Gọi $F\left( x \right),G\left( x \right)$ là hai nguyên hàm của $f\left( x \right)$ trên $\mathbb{R}$ thỏa mãn $F\left( 2 \right)+G\left( 2 \right)=2$ và $F\left( 0 \right)+G\left( 0 \right)=-2$. Khi đó $
\displaystyle\int\limits_0^8f\left( \dfrac{x}{4} \right)\text{d}x$ bằng
\shortans[]{$8$}
\loigiai{
Ta có $G\left( x \right)=F\left( x \right)+C\Rightarrow \left\{ \begin{aligned}
& G\left( 2 \right)=F\left( 2 \right)+C \\
& G\left( 0 \right)=F\left( 0 \right)+C. \\
\end{aligned} \right.$ \\
Mà
$\left\{ \begin{aligned}
& F\left( 2 \right)+G\left( 2 \right)=2 \\
& F(0)+G(0)=-2 \\
\end{aligned} \right.\Leftrightarrow \left\{ \begin{aligned}
& 2F(2)+C=2 \\
& 2F(0)+C=-2 \\
\end{aligned} \right.\Leftrightarrow F(2)-F(0)=2.$\\
Vậy $\displaystyle\int\limits_0^8f\left(\dfrac{x}{4} \right)\text{d}x=4\displaystyle\int\limits_0^2{f(t)dt=4\left(F(2)-F(0) \right)=8.}$}
\end{ex}

\begin{ex}%[2D4V2-1]%[Dự án EX-TF-TLN lần 2 - Nguyễn Trần Phong]
Cho bất phương trình $\displaystyle \int\limits_{0}^{x} (3t^2 - 8t + 4)\, \mathrm{d}t \leq x, \quad (x > 0)$. Tính tổng các nghiệm nguyên của bất phương trình đã cho.
\shortans{$6$}
\loigiai{Ta có
\begin{eqnarray*}
\int\limits_{0}^{x} (3t^2 - 8t + 4)\, \mathrm{d}t \leq x & \Leftrightarrow & (t^3 - 4t^2 + 4t ) \Big|_0^x \leq x \\
& \Leftrightarrow & x^3 - 4x^2 + 4x \leq x \\
& \Leftrightarrow & x^2 - 4x + 4 \leq 1 \qquad (\text{vì } x > 0) \\
& \Leftrightarrow & x^2 - 4x + 3 \leq 0 \\
& \Leftrightarrow & 1 \leq x \leq 3.
\end{eqnarray*}
Vậy tập các nghiệm nguyên của bất phương trình đã cho là $S = \left\lbrace 1, 2, 3\right\rbrace$. Suy ra, tổng các nghiệm nguyên có giá trị bằng $6$.
}
\end{ex}

\begin{ex}%[Vanle Vo - DA tex hoa de cuoi ki 1-2425]%[2D4V2-2]
Bảng dưới đây thống kê cự li ném tạ của một vận động viên
\[\begin{array}{|c|c|c|c|c|c|}
\hline
\text{Cự li (m)} & [19;21) & [21;23) & [23;25) & [25;27) & [27;29) \\
\hline
\text{Tần số} & 13 & 45 & 24 & 12 & 6 \\
\hline
\end{array}\]
Hãy tính độ lệch chuẩn của mẫu số liệu ghép nhóm trên (kết quả được làm tròn đến hàng phần trăm).
\shortans[]{$2{,}07$.}
\loigiai{
Bảng dưới đây thống kê cự li ném tạ của một vận động viên
\[\begin{array}{|c|c|c|c|c|c|}
\hline
\text{Cự li đại diện (m)} & 20 & 22 & 24 & 26 & 28 \\
\hline
\text{Tần số} & 13 & 45 & 24 & 12 & 6 \\
\hline
\end{array}\]
Cỡ mẫu $n=100$.\\
Số trung bình
\[\overline{x}=\dfrac{13\cdot20+45\cdot22+24\cdot24+12\cdot26+6\cdot28}{100}=23{,}06.\]
Phương sai\\
$s^2=\left[13\cdot(20-23{,}06)^2 + 45\cdot(22-23{,}06)^2 + 24\cdot(24-23{,}06)^2 + 12\cdot(26-23{,}06)^2\right.\\
\text{}\qquad\left. + 6\cdot(28-23{,}06)^2\right]:100=4{,}28$.\\
Độ lệch chuẩn $\sigma = \sqrt{s^2} = \sqrt{4{,}28} \approx 2{,}07$.
}
\end{ex}

\begin{ex}%[2D4V2-2]%[To 20 - Dot 17 - Chuong 4 - Bai 3 - CD - De 3]%[Đình Nguyên]
Biết rằng hàm số $f(x)=ax^3+bx^2+cx+d$ thỏa mãn
$\displaystyle\int\limits_{0}^{1} f(x) \mathrm{\,d}x=-4$, $\displaystyle\int\limits_{0}^{2} f(x) \mathrm{\,d}x=-2$, $\displaystyle\int\limits_{0}^{3} f(x) \mathrm{\,d}x=18$, $\displaystyle\int\limits_{0}^{4} f(x) \mathrm{\,d}x=80$. Tính giá trị của biểu thức $P=2 a-3 b+4 c+5 d$.
\shortans{3}
\loigiai{
Ta có
\begin{eqnarray*}
\int\limits_{0}^{t} f(x) \mathrm{\,d}x&=&\int\limits_{0}^{t}\left(a x^{3}+b x^{2}+c x+d\right) dx\\
&=&\left.\left(\dfrac{a}{4} x^{4}+\dfrac{b}{3} x^{3}+\dfrac{c}{2} x^{2}+\mathrm{\,d}x\right)\right|_{0} ^{t}=\dfrac{a}{4} t^{4}+\dfrac{b}{3} t^{3}+\dfrac{c}{2} t^{2}+d t.
\end{eqnarray*}
Do đó
\begin{eqnarray*}
\heva{&\int\limits_{0}^{1} f(x) \mathrm{\,d}x=-4 \\ \displaystyle&\int\limits_{0}^{2} f(x) \mathrm{\,d}x=-2 \\ \displaystyle&\int\limits_{0}^{3} f(x) \mathrm{\,d}x=18 \\ \displaystyle&\int\limits_{0}^{4} f(x) \mathrm{\,d}x=80} &\Leftrightarrow& \heva{&\dfrac{a}{4} \cdot 1^{4}+\dfrac{b}{3} \cdot 1^{3}+\dfrac{c}{2} \cdot 1^{2}+d \cdot 1=-4 \\ &\dfrac{a}{4} \cdot 2^{4}+\dfrac{b}{3} \cdot 2^{3}+\dfrac{c}{2} \cdot 2^{2}+d \cdot 2=-2 \\ &\dfrac{a}{4} \cdot 3^{4}+\dfrac{b}{3} \cdot 3^{3}+\dfrac{c}{2} \cdot 3^{2}+d \cdot 3=18 \\ &\dfrac{a}{4} \cdot 4^{4}+\dfrac{b}{3} \cdot 4^{3}+\dfrac{c}{2} \cdot 4^{2}+d \cdot 4=80}\\
&\Leftrightarrow& \heva{&\dfrac { a } { 4 } + \dfrac { b } { 3 } + \dfrac { c } { 2 } + d = - 4  \\
& \dfrac { a } { 4 } \cdot 1 6 + \dfrac { b } { 3 } \cdot 8 + \dfrac { c } { 2 } \cdot 4 + d \cdot 2 = - 2  \\
& \dfrac { a } { 4 } \cdot 8 1 + \dfrac { b } { 3 } \cdot 2 7 + \dfrac { c } { 2 } \cdot 9 + d \cdot 3 = 1 8  \\
&\dfrac { a } { 4 } \cdot 2 5 6 + \dfrac { b } { 3 } \cdot 64 + \dfrac { c } { 2 } \cdot 1 6 + d \cdot 4 = 80 }\\
&\Leftrightarrow& \heva{&a=2 \\ &b=-3 \\ &c=5 \\ &d=-6.}
\end{eqnarray*}
Vậy $P=2a-3b+4c+5d=3$.
}
\end{ex}

\begin{ex}%[2D4V2-2]%[Tổ 20 - Đợt 17 - Chương 4 - - CD]%[Phạm Hà Giang]
Cho $\displaystyle\int\limits_0^1 \left( \dfrac{1}{x+1}-\dfrac{1}{x+2} \right) \mathrm{\,d}x=a\ln2 +b\ln3$ với $a, b$ là các số nguyên. Tính tích $ab$.
\shortans{$-2$}
\loigiai
{ Ta có $\displaystyle\int\limits_0^1 \dfrac{1}{x+1} \mathrm{\,d}x=\left. \ln \left| x+1 \right| \right| _0^1=\ln2$ và $\displaystyle\int\limits_0^1 \dfrac{1}{x+2} \mathrm{\,d}x=\left. \ln \left| x+2 \right| \right| _0^1=\ln3-\ln2$.\\
Do đó $\displaystyle\int\limits_0^1 \left( \dfrac{1}{x+1}-\dfrac{1}{x+2}\right) \mathrm{\,d}x=\ln2-(\ln3-\ln2)=2\ln2-\ln3 \Rightarrow a=2, b=-1$.\\
Vậy $ab=-2$.
}
\end{ex}

\begin{ex}%[2D4V2-2]
Cho $I=\displaystyle\int\limits_0^1 \dfrac{1}{3+2 x-x^2} \mathrm{\,d} x=(a-b) \ln 2+b \ln 3$. Tính $a+b$.
\shortans{$0{,}5$}
\loigiai{Ta có
$I=\displaystyle\int\limits_0^1 \dfrac{1}{3+2 x-x^2} \mathrm{\,d} x =\displaystyle\int\limits_0^1\left(\dfrac{\frac{1}{4}}{x+1}+\dfrac{\frac{1}{4}}{3-x}\right)\mathrm{\,d} x \\
=\left.\dfrac{1}{4}(\ln |x+1|-\ln |x-3|)\right|_0 ^1=\dfrac{1}{4} \ln 3 \Rightarrow a=b=\dfrac{1}{4} \Rightarrow a+b=\dfrac{1}{2}.
$}
\end{ex}

\begin{ex}%[2D4V2-2]
Tìm số thực dương $m$ để tích phân $\displaystyle\int\limits_0^m\left(x-x^2\right) \mathrm{\,d}x$ có giá trị lớn nhất.
\shortans{$1$}
\loigiai{\[
P=\displaystyle\int\limits_0^m\left(x-x^2\right) \mathrm{\,d}x=\left.\left(\dfrac{x^2}{2}-\dfrac{x^3}{3}\right)\right|_0 ^m=\dfrac{m^2}{2}-\dfrac{m^3}{3} \text {. }
\]
Đặt $f(m)=\dfrac{m^2}{2}-\dfrac{m^3}{3} \Rightarrow f^{\prime}(m)=m-m^2 \Rightarrow f^{\prime}(m)=0 \Leftrightarrow m=0$ hoặc $m=1$.\\
Lập bảng biến thiên
\begin{center}
\begin{tikzpicture}
\tkzTabInit[nocadre=false,lgt=1.5,espcl=2.5,deltacl=1]
{$m$ /0.7,$f'(m)$ /0.7,$f(m)$ /2}
{$0$,$1$,$+\infty$}
\tkzTabLine{,+,$0$,-,}
\tkzTabVar{-/$-\infty$,+/$\dfrac{1}{6}$,-/$-\infty$}
\end{tikzpicture}
\end{center}
Vậy $f(m)$ đạt GTLN tại $m=1$.
}
\end{ex}

\begin{ex}%[2D4V2-2]
Cho hàm số $ f(x)\ne 0$, liên tục trên đoạn $\left[1;2\right]$ và thỏa mãn $ f(1)=3$, \break $x^2\cdot f'(x)=f^2(x)$ với $\forall x\in\left[1;2\right]$. Tính $f(2)$.
\shortans{$-6$}
\loigiai{
Ta có
\begin{align*}
x^2\cdot f'(x)=f^2(x)& \Rightarrow\dfrac{f'(x)}{f^2(x)}=\dfrac{1}{x^2} \Rightarrow{\left(-\dfrac{1}{f(x)}\right)'}=\dfrac{1}{x^2}\\& \Rightarrow\displaystyle\int\limits_1^2\left(-\dfrac{1}{f(x)}\right)'\mathrm{\,d} x=\displaystyle\int\limits_1^2\dfrac{1}{x^2}\mathrm{\,d} x\\&\Rightarrow\left.\left(-\dfrac{1}{f(x)}\right)\right|_1^2=-\left.\dfrac{1}{x}\right|_1^2\\& \Rightarrow-\dfrac{1}{f(2)}+\dfrac{1}{f(1)}=-\dfrac{1}{2}+1\\& \Rightarrow-\dfrac{1}{f(2)}+\dfrac{1}{f(1)}=\dfrac{1}{2}.
\end{align*}
Vì $f(1)=3\Rightarrow-\dfrac{1}{f(2)}+\dfrac{1}{3}=\dfrac{1}{2}\Rightarrow f(2)=-6$.}
\end{ex}

\begin{ex}%[2D4V2-2]
Tính tích phân $I=\displaystyle\int\limits _{-3}^{3}\left|x^{2} -1\right|\mathrm{\,d}x $ (tính gần đúng đến hàng phần chục).
\shortans{$13{,}3$}
\loigiai{
\[I=\displaystyle\int\limits _{-3}^{3}\left|x^{2} -1\right|\mathrm{\,d}x.\]
Vì  $f(x)=x^{2} -1=0\to\hoac{&x=-1\\&x=1} \Rightarrow f(x)>0,~\forall x\in \left[-3;-1\right]\cup \left[1;3\right]$; $f(x)<0,~\forall x\in \left[-1;1\right]$.\\
Vậy
\begin{eqnarray*}
I&=&\displaystyle\int\limits _{-3}^{-1}\left(x^{2} -1\right)\mathrm{\,d}x+\displaystyle\int\limits _{-1}^{1}\left(1-x^{2} \right)\mathrm{\,d}x+\displaystyle\int\limits _{1}^{3}\left(x^{2} -1\right)\mathrm{\,d}x\\
&=&\left(\frac{1}{3} x^{3} -x\right)\Bigg|_{-3}^{-1}+\left(x-\frac{1}{3} x^{3} \right)\Bigg|_{-1}^1+\left(\frac{1}{3} x^{3} -x\right)\Bigg|_1^3\\
&=&\frac{20}{3} +\frac{4}{3} +\frac{16}{3} =\frac{40}{3}\approx 13{,}3.
\end{eqnarray*}
}
\end{ex}

\begin{ex}%[2D4V2-2]
Tính tích phân $I=\displaystyle\int\limits _{-1}^{2}\left|-x^{2} -2x+3\right|\mathrm{\,d}x $ (tính gần đúng đến hàng phần trăm).
\shortans{$7{,}67$}
\loigiai{
Vì  $f(x)=-x^{2} -2x+3=0\Rightarrow\hoac{&x=1\\&x=-3} \Rightarrow f(x)>0,~\forall x\in \left[-1;-1\right]$; $f(x)<0,~\forall x\in \left[1;2\right]$\\
Vậy
\begin{eqnarray*}
I&=&\displaystyle\int\limits _{-1}^{1}\left(-x^{2} -2x+3\right)\mathrm{\,d}x+\displaystyle\int\limits _{1}^{2}\left(x^{2}+2x-3 \right)\mathrm{\,d}x\\
&=&\left(-\dfrac{1}{3} x^{3}-x^2 +3x\right)\Bigg|_{-1}^{1}+\left(\dfrac{1}{3} x^{3}+x^2 -3x \right)\Bigg|_{1}^2\\
&=&-\dfrac{1}{3}-1+3-\dfrac{1}{3}+1+3 +\dfrac{8}{3}+4-6-\dfrac{1}{3}-1+3 \approx7{,}67.
\end{eqnarray*}
}
\end{ex}

\begin{ex}%[2D4V2-2]
Tính tích phân $I=\displaystyle\int\limits _{1}^{2}\left|\frac{x+1}{x} \right|\mathrm{\,d}x $ (tính gần đúng đến hàng phần trăm).
\shortans{$1{,}69$}
\loigiai{
Vì $\frac{x+1}{x}>0$, $\forall x\in [1;2]$ nên
\[I=\displaystyle\int\limits _{1}^{2}\left(\dfrac{x+1}{x}\right)\mathrm{\,d}x=\displaystyle\int\limits _{1}^{2}\left(1+\dfrac{1}{x}\right)\mathrm{\,d}x=\left(x+\ln x \right)\Bigg|1^2=2+\ln2-1=1+\ln2\approx1{,}69.  \]
}
\end{ex}

\begin{ex}%[2D4V2-2]
Tính tích phân $I=\displaystyle\int\limits _{2}^{6}\sqrt{x^{2} -8x+16} \mathrm{\,d}x $.
\shortans{$4$}
\loigiai{
Ta có $I=\displaystyle\int\limits _{2}^{6}\left| x-4\right|  \mathrm{\,d}x $.\\
Ta có $x-4\le 0$, $\forall x\in [2;4]$	và 	 $x-4\ge 0$, $\forall x\in [4;6]$. Khi đó
\[I=\displaystyle\int\limits _{2}^{4}\left( 4- x\right)   \mathrm{\,d}x+\displaystyle\int\limits _{4}^{6}\left( x- 4\right)   \mathrm{\,d}x=\left( 4 x-\dfrac{x^2}{2}\right)\Bigg|_{2}^{4}+\left( -4 x+\dfrac{x^2}{2}\right)\Bigg|_{4}^{6}=4.\]
}
\end{ex}

\begin{ex}%[2D4V2-2]
Tính tích phân $I=\displaystyle\int\limits _{-2}^{1}\sqrt{4x^{2} +6x+9} \mathrm{\,d}x $ (\textit{làm tròn đến hàng phần trăm}).
\shortans{$9{,}38$}
\loigiai{
Ta có $I=\displaystyle\int\limits _{-2}^{1}\sqrt{4x^{2} +6x+9} \mathrm{\,d}x=\displaystyle\int\limits _{-2}^{1}\left|2x+3 \right|  \mathrm{\,d}x$.\\
Ta có $2x+3\le 0$, $\forall x\in \left[-2;-\dfrac{3}{2} \right] $	và 	 $2x+3\ge 0$, $\forall x\in \left[-\dfrac{3}{2};1 \right]$. Khi đó
\begin{eqnarray*}
I&=&\displaystyle\int\limits _{-2}^{-\tfrac{3}{2}}\left( -2x-3 \right)   \mathrm{\,d}x+\displaystyle\int\limits _{-\tfrac{3}{2}}^{1}\left( 2x+3 \right)  \mathrm{\,d}x\\
&=&\left(-x^2-3x\right)\Bigg|_{-2}^{-\tfrac{3}{2}}+\left(x^2+3x \right)\Bigg|_{-\tfrac{3}{2}}^{1}\\
&\approx&9{,}38.
\end{eqnarray*}
}
\end{ex}

\begin{ex}%[2D4V2-2]
Tính tích phân $I=\displaystyle\int\limits _{0}^{1}\sqrt{9x^{2} -6x+1} \mathrm{\,d}x $ (\textit{làm tròn đến hàng phần trăm}).
\shortans{$0{,}83$}
\loigiai{
Ta có $I=\displaystyle\int\limits _{0}^{1}\sqrt{9x^{2} -6x+1} \mathrm{\,d}x =\displaystyle\int\limits _{0}^{1}\left| 3x-1\right| \mathrm{\,d}x $.\\
Ta có $3x+1\le 0$, $\forall x\in \left[1;\dfrac{1}{3} \right] $	và 	 $3x+1\ge 0$, $\forall x\in \left[\dfrac{1}{3};1 \right]$. Khi đó
\begin{eqnarray*}
I&=&\displaystyle\int\limits _{0}^{\tfrac{1}{3}}\left(-3x-1 \right)   \mathrm{\,d}x+\displaystyle\int\limits _{\tfrac{1}{3}}^{1}\left( 3x+1 \right)  \mathrm{\,d}x\\
&=&\left(-\dfrac{3x^2}{2}-x\right)\Bigg| _{0}^{\tfrac{1}{3}}+\left(\dfrac{3x^2}{2}+x\right)\Bigg|_{\tfrac{1}{3}}^{1}\\
&\approx&0{,}83.
\end{eqnarray*}
}
\end{ex}

\begin{ex}%[2D4V2-2]Biên soạn:Phùng Hoàng Em Phản biện:Nguyễn Đắc Kiên
Cho hàm số $f(x)=\heva{&x^2+1 & \text{ nếu } x\geq 1\\&2x & \text{ nếu } x<1}$. Tính tích phân $\displaystyle\int\limits_0^2 f(x)\mathrm{\,d}x$ (\textit{làm tròn đến hai chữ số thập phân}).
\shortans{$4{,}33$}
\loigiai
{Vì $\lim\limits_{x \to 1^+} f(x) = \lim\limits_{x \to 1^-} f(x)=f(1)=2$ nên hàm số $f(x)$ liên tục trên $\mathbb{R}$.\\
Vì vậy $\displaystyle\int\limits_0^2 f(x)\mathrm{\,d}x = \displaystyle\int\limits_0^1 2x\mathrm{\,d}x + \displaystyle\int\limits_1^2 \left(x^2+1\right)\mathrm{\,d}x =1+ \dfrac{10}{3} = \dfrac{13}{3}\approx 4{,}33$.}
\end{ex}

\begin{ex}%[2D4V2-2]
Thâm niên công tác của các công nhân hai nhà máy $A$ và $B$ được thống kê ở bảng dưới đây:
\begin{center}
\setlength{\extrarowheight}{2pt}
\newcommand{\PreBackslash}[1]{\let\temp=\\#1\let\\=\temp}
\let\PBS=\PreBackslash
\begin{tabular}{|>{\PBS\centering}m{6cm}|>{\PBS\centering}m{1.5cm}|>{\PBS\centering}m{1.5cm}| >{\PBS\centering}m{1.5cm}|>{\PBS\centering}m{1.5cm}|>{\PBS\centering}m{1.5cm}|}\hline
\textbf{Thâm niên công tác (năm)} &$[0;5)$&$[5;10)$&$[10;15)$&$[15;20)$&$[20;25)$\\
\hline
\textbf{Số công nhân nhà máy $A$}&$35$&$13$&$12$&$12$&$8$\\
\hline
\textbf{Số công nhân nhà máy $B$}&$14$&$26$&$24$&$11$&$5$\\
\hline
\end{tabular}
\end{center}
Gọi $S_A$, $S_B$ lần lượt là độ lệch chuẩn về thâm niên công tác của công nhân hai nhà máy $A$ và $B$. Tính giá trị của $S_A-S_B$ (làm tròn kết quả đến hàng phần trăm).
\shortans[]{$1{,}48$}
\loigiai{
Thâm niên công tác của công nhân hai nhà máy $A$ và $B$ được thống kê như sau:
\begin{center}
\setlength{\extrarowheight}{2pt}
\newcommand{\PreBackslash}[1]{\let\temp=\\#1\let\\=\temp}
\let\PBS=\PreBackslash
\begin{tabular}{|>{\PBS\centering}m{6cm}|>{\PBS\centering}m{1.5cm}|>{\PBS\centering}m{1.5cm}| >{\PBS\centering}m{1.5cm}|>{\PBS\centering}m{1.5cm}|>{\PBS\centering}m{1.5cm}|}
\hline
\textbf{Thâm niên công tác (năm)} &$[0;5)$&$[5;10)$&$[10;15)$&$[15;20)$&$[20;25)$\\
\hline
\textbf{Giá trị đại diện} &$2{,}5$&$7{,}5$&$12{,}5$&$17{,}5$&$22{,}5$\\
\hline
\textbf{Số công nhân nhà máy $A$}&$35$&$13$&$12$&$12$&$8$\\
\hline
\textbf{Số công nhân nhà máy $B$}&$14$&$26$&$24$&$11$&$5$\\
\hline
\end{tabular}
\end{center}
\begin{enumerate}[$\bullet$]
\item Xét mẫu số liệu thâm niên công tác của công nhân nhà máy $A$.\\
Thâm niên trung bình của công nhân nhà máy $A$ :
\[\overline{x}_A=\dfrac{1}{80}\left(2{,}5\cdot 35+7{,}5\cdot 13+12{,}5\cdot 12+17{,}5\cdot 12+22{,}5\cdot8\right)=\dfrac{145}{16}.\]
Phương sai về thâm niên của công nhân nhà máy $A$: \[S_A^2=\dfrac{1}{80}\left(2{,}5^2\cdot 35+7{,}5^2\cdot 13+12{,}5^2\cdot 12+17{,}5^2\cdot 12+22{,}5^2\cdot8\right)-\left(\dfrac{145}{16}\right)^2 =\dfrac{12\,735}{256}.\]
Độ lệch chuẩn  về thâm niên của công nhân nhà máy $A$: $S_A=\sqrt{S_A^2}\approx 7{,}05$.
\item Xét mẫu số liệu thâm niên công tác của công nhân nhà máy $B$.\\
Thâm niên trung bình của công nhân nhà máy $B$ :
\[\overline{x}_B=\dfrac{1}{80}\left(2{,}5\cdot 14+7{,}5\cdot 26+12{,}5\cdot 24+17{,}5\cdot 11+22{,}5\cdot5\right)=\dfrac{167}{16}.\]
Phương sai về thâm niên của công nhân nhà máy $B$: \[S_B^2=\dfrac{1}{80}\left(2{,}5^2\cdot 14+7{,}5^2\cdot 26+12{,}5^2\cdot 24+17{,}5^2\cdot 11+22{,}5^2\cdot 5\right)-\left(\dfrac{167}{16}\right)^2 =\dfrac{7\,951}{256}.\]
Độ lệch chuẩn  về thâm niên của công nhân nhà máy $B$: $S_B=\sqrt{S_B^2}\approx 5{,}57$.\\
Giá trị của $S_A-S_B\approx 7{,}05-5{,}57\approx 1{,}48$.
\end{enumerate}
}
\end{ex}

\begin{ex}%[2D4V2-2]
Biết $\displaystyle\int_1^3\dfrac{x+2}{x} \mathrm{\,d}x=a+b\ln c$, với $a,b,c\in \mathbb{Z},c < 9$. Tính tổng $S=a+b+c$.
\shortans[]{$7$}
\loigiai{
Ta có $\displaystyle\int_1^3\dfrac{x+2}{x} \mathrm{\,d}x=\displaystyle\int_1^3\left(1+\dfrac{2}{x} \right)\mathrm{\,d}x=\displaystyle\int_1^3\mathrm{\,d}x+\displaystyle\int_1^3\dfrac{2}{x} \mathrm{\,d}x=2+2\ln \left|x\right|\Bigr|_1^3=2+2\ln 3$.\\
Do đó $a=2$, $b=2$, $c=3$ $\Rightarrow S=7$.
}
\end{ex}

\begin{ex}%[2D4V2-2]
Cho $F(x)$ là một nguyên hàm của $f(x)=\dfrac{2}{x+2}$. Biết $F\left(-1\right)=0$. Tính $F(2)=a\ln b$ với $a$, $b$ là các số tự nhiên. Tổng $a+b$ bằng bao nhiêu?
\shortans[]{$6$}
\loigiai{
Ta có
\begin{eqnarray*}
&&\displaystyle\int \limits_{-1}^2f(x)\mathrm{\,d}x=F(2)-F\left(-1\right)\\
&\Leftrightarrow&\displaystyle\int\limits _{-1}^2\dfrac{2}{x+2}\mathrm{\,d}x=2\ln \left|x+2\right|\Bigr|_{-1}^2=2\ln 4-2\ln 1=2\ln 4\\
&\Leftrightarrow&F(2)-F(-1)=2\ln 4\\
&\Leftrightarrow& F(2)=2\ln 4 \quad (\text{do\,} F\left(-1\right)=0).
\end{eqnarray*}
Vậy $a=2$, $b=4$ nên tổng $a+b=6$.
}
\end{ex}

\begin{ex}%[2D4V2-2]
Có hai giá trị của số thực $a$ là $a_1$, $a_2$ ($0< a_1 < a_2$) thỏa mãn $\displaystyle\int\limits_1^a\left(2x-3\right)\mathrm{\,d}x=0$. Hãy tính $T=3^{a_1}+3^{a_2}+\log_2 \left(\dfrac{a_2}{a_1} \right)$.
\shortans[]{$13$}
\loigiai{
Ta có: $\displaystyle\int\limits_1^a\left(2x-3\right)\mathrm{\,d}x$ $=\left. \left(x^2-3x\right)\right|_1^a$ $=a^2-3a+2$.\\
Vì $\displaystyle\int\limits_1^a\left(2x-3\right)\mathrm{\,d}x=0$ nên $a^2-3a+2=0$, suy ra $\hoac{&a=1 \\&a=2}$.\\
Lại có $0< a_1 < a_2$ nên $a_1=1$; $a_2=2$.\\
Như vậy $T=3^{a_1}+3^{a_2}+\log_2 \left(\dfrac{a_2}{a_1} \right)$ $=3^1+3^2+\log_2 \left(\dfrac{2}{1} \right)$ $=13$.
}
\end{ex}

\begin{ex}%[2D4V2-2]
Biết rằng hàm số $f(x)=mx+n$ thỏa mãn $\displaystyle\int _0^1f(x) \mathrm{\,d}x=3$, $\displaystyle\int _0^2f(x) \mathrm{\,d}x=8$. Tính tổng $m+n$?
\shortans[]{$4$}
\loigiai{
Ta có $\displaystyle\int f(x) \mathrm{\,d}x=\displaystyle\int \left(mx+n\right) \mathrm{\,d}x$=$\dfrac{m}{2} x^2+nx+\mathbb{C}$.\\
Lại có $\displaystyle\int _0^1f(x) \mathrm{\,d}x=3$ $\Rightarrow \left(\dfrac{m}{2} x^2+nx\right)\Bigr|_0^1=3$ $\Leftrightarrow \dfrac{1}{2} m+n=3$. \quad(1)\\
$\displaystyle\int _0^2f(x) \mathrm{\,d}x=8\Rightarrow \left(\dfrac{m}{2} x^2+nx\right)\Bigr|_0^2=8\Leftrightarrow 2m+2n=8
$. \quad (2)\\
Từ $(1)$ và $(2)$ ta có hệ phương trình $\heva{&\dfrac{1}{2} m+n=3\\&2m+2n=8}\Leftrightarrow \heva{&{m=2} \\&{n=2.}}$\\
Vậy  $m+n=4$.
}
\end{ex}

\begin{ex}%[2D4V2-2]
Biết $\displaystyle\int_1^2\dfrac{\mathrm{\,d}x}{\left(x+1\right)\left(2x+1\right)}=a\ln 2+b\ln 3+c\ln 5$. Khi đó giá trị $a+b+c$ bằng
\shortans[]{$0$}
\loigiai{
Ta có
\[\displaystyle\int_1^2\dfrac{\mathrm{\,d}x}{\left(x+1\right)\left(2x+1\right)}=\displaystyle\int_1^2\left(\dfrac{2}{2x+1}-\dfrac{1}{x+1} \right)\mathrm{\,d}x=2\displaystyle\int_1^2\dfrac{1}{2x+1} \mathrm{\,d}x-\displaystyle\int_1^2\dfrac{1}{x+1} \mathrm{\,d}x
\]
\[=2\cdot\dfrac{1}{2} \ln \left|2x+1\right|\Bigr|_1^2-\ln \left|x+1\right|\Bigr|_1^2=\ln \left(2x+1\right)\Bigr|_1^2-\ln \left(x+1\right)\Bigr|_1^2=\ln 5-\ln 3-\left(\ln 3-\ln 2\right)
\]
\[=\ln 2-2\ln 3+\ln 5.
\]
Do đó $a=1$, $b=-2$, $c=1$.\\
Vậy $a+b+c=1+\left(-2\right)+1=0$.
}
\end{ex}

\begin{ex}%[DA-EX-TF-TLN2024 - Trần Xuân Hòa]%[2D4V2-2]
Biết $\displaystyle\int\limits_0^1\dfrac{x}{\sqrt{x+1}}\mathrm{\,d}x=\dfrac{a}{b}\left(-\sqrt{2}+c\right)$ với $\dfrac{a}{b}$ là phân số tối giản. Tính $a+b+c$.
\shortans{$7$}
\loigiai{
Ta có $\displaystyle\int\limits_0^1\dfrac{x}{\sqrt{x+1}}\mathrm{\,d}x=\displaystyle\int\limits_0^1\left(\sqrt{x+1}-\dfrac{2}{2\sqrt{x+1}}\right)\mathrm{\,d}x=\left.\left(\dfrac{2}{3}\sqrt{(x+1)^3}-2\sqrt{x+1}\right)\right|_0^1=\dfrac{2}{3}(-\sqrt{2}+2)$.\\
Từ đó suy ra $a=2$, $b=3$ và $c=2$.\\
Vậy $a+b+c=7$.
}
\end{ex}

\begin{ex}%[DA-EX-TF-TLN2024 - Trần Xuân Hòa]%[2D4V2-2]
Biết $I=\displaystyle\int\limits_3^4\dfrac{\mathrm{\,d}x}{x^2+x}=a\ln 2+b\ln 3+c\ln 5$ với $a$, $b$, $c$ là các số nguyên. Tính $S=a+b+c$.
\shortans{$2$}
\loigiai{
\begin{eqnarray*}
I&= & \displaystyle\int\limits_3^4\dfrac{\mathrm{\,d}x}{x^2+x}=\displaystyle\int\limits_3^4\dfrac{\mathrm{\,d}x}{x(x+1)}=\displaystyle\int\limits_3^4\dfrac{\mathrm{\,d}x}{x}
-\displaystyle\int\limits_3^4\dfrac{\mathrm{\,d}x}{x+1}=\ln {|x|}\bigg|_3^4-\ln {|x+1|}\bigg|_3^4\\
&= & \ln 4-\ln 3-\ln 5+\ln 4=4\ln 2-\ln 3-\ln 5.
\end{eqnarray*}
Suy ra $a=4$, $b=c=-1$ nên $S=2$.
}
\end{ex}

\begin{ex}%[2D4V2-2]
%	Cho $g(x)=\displaystyle\int\limits_0^x f(t)\mathrm{\,d} t$, $(0\le x\le 7)$ trong đó $f(t)$
\immini{	Cho $g(x)=\displaystyle\int\limits_0^x f(t)\mathrm{\,d}t$, $(0\le x\le 7)$ trong đó $f(t)$
là hàm số có đồ thị như \textit{Hình 8}. Tính $g(3)$.}
{ \begin{tikzpicture}[line join=round, line cap=round,>=stealth,thick,scale=.5]
\draw[step=1,gray!30,very thin] (-1,-3) grid (8,5);
\draw[->] (-1,0)--(8,0) node[below left] {$t$};
\foreach \x/\nx in {1/1,5/5}
\draw[thin] (\x,1pt)--(\x,-1pt) node [below right] {$\nx$};
\draw[->] (0,-3)--(0,5) node[below left] {$f$};
\foreach \y/\ny in {1/1}
\draw[thin] (-1pt,\y)--(1pt,\y) node [above left ] {$\ny$};
\draw (0,0) node [below left] {$O$} ;
\draw [very thick]  (0,2)--(1,2)--(2,4)--(3,0)--(5,-2)--(6,-2)--(7,0);
\draw  (3,-3) node [below]{Hình 8};
\draw (3.6,1.5)node{$f(t)$};
\end{tikzpicture}}
\shortans{$7$}
\loigiai{
Ta có
$f(t)=\heva{& 1 \quad\quad\,\, \text{ nếu } 0\le t\le 1,\\& 2t \quad\quad\,\text{ nếu } 1\le t\le 2,\\&12-4t \text{ nếu } 2\le t\le 3,\\& 3-t\quad\text{ nếu } 3\le t\le 5,\\&-2 \quad\,\,\,\text{ nếu } 5\le t\le 6, \\&35-5t\text{ nếu } 6\le t\le 7.}$\\
Do đó,
\begin{eqnarray*}
g(3)=\displaystyle\int\limits_0^3 f(t) \mathrm{\,d} t &=&\displaystyle\int\limits_0^1 f(t) \mathrm{\,d} t+\displaystyle\int\limits_1^2 f(t) \mathrm{\,d} t+\displaystyle\int\limits_2^3 f(t)\mathrm{\,d} t\\
& =&\displaystyle\int\limits_0^1 2 \mathrm{\,d} t+\displaystyle\int\limits_1^2 2 t \mathrm{\,d} t+\displaystyle\int\limits_2^3(12-4 t) \mathrm{\,d} t \\
& =&\left.2 t\right|_0 ^1+\left.t^2\right|_1 ^2+\left.\left(12 t-2 t^2\right)\right|_2 ^3=7.
\end{eqnarray*}
%Hinh 8
\textbf{Cách khác:}\\
\begin{center}
\begin{tikzpicture}[line join=round, line cap=round,>=stealth,thick,scale=.5]
\draw[step=1,gray!30,very thin] (-1,-3) grid (8,5);
\draw[->] (-1,0)--(8,0) node[below left] {$t$};
\foreach \x/\nx in {1/1,5/5}
%				\draw[thin] (\x,1pt)--(\x,-1pt) node [below right] {$\nx$};
\draw[->] (0,-3)--(0,5) node[below left] {$f$};
%				\foreach \y/\ny in {1/1}
%				\draw[thin] (-1pt,\y)--(1pt,\y) node [above left ] {$\ny$};
\draw (0,0) node [below left] {$O$} ;
\draw [very thick]  (0,2)--(1,2)--(2,4)--(3,0)--(5,-2)--(6,-2)--(7,0);
\draw  (3,-3) node [below]{Hình 8};
\draw (3.6,1.5)node{$f(t)$};
\draw[thick,teal]
(0,0)--(0,2)--(1,2)--(1,0)--cycle
(1,2)--(2,4)--(2,0)--(1,0)
(2,4)--(3,0)--(2,0)
;
\draw[fill=white]
(0,2) circle (0.05) node[left]{$A$}
(1,2) circle (0.05) node[above left]{$B$}
(2,4) circle (0.05) node[above]{$C$}
(3,0) circle (0.05) node[below]{$D$}
(2,0) circle (0.05) node[below]{$E$}
(1,0) circle (0.05) node[below]{$F$}
;
\end{tikzpicture}
\end{center}
Ta có $g(x)=\displaystyle \int _0^x f(t){\,d}t \Rightarrow g(0)=0.$\\
$\displaystyle \int _0^3 f(t){\,d}t =S_{OABCD} \Rightarrow g(3)-g(0)=S_{OABF}+S_{FBCE}+S_{CED} \Rightarrow g(3)=2+3+\dfrac{1}{2}.1.4=7$.
}
\end{ex}

\begin{ex}%[Mức độ 3]%[BG12, Nguyễn Kiều Nhã Tú]%[2D4V2-2]
Cho hàm số $f(x)=\heva{&2x+3 & \text {nếu} &\ x \le 2 \\& 3x^2+m & \text{nếu} &\ x>2}$. Biết rằng hàm số $f(x)$ liên tục trên $\mathbb{R}$. Hãy xác định giá trị của tích phân $\displaystyle\int\limits_{-1}^5f(x)\mathrm{\,d}x$.
\SA{$114$}
\loigiai{
Để hàm số $f(x)$ liên tục trên $\mathbb{R}$ thì hàm số $f(x)$ phải liên tục tại điểm $x_0=2$. Suy ra
\begin{eqnarray*}
&&\lim\limits _{x \to 2^-} f(x)=\lim\limits _{x \to 2^+} f(x)=f(2)\\
& \Leftrightarrow& \lim\limits _{x \to 2^-}(2x+3)=\lim\limits_{n \to +\infty}_{x \to 2^+}(3x^2+m)=7 \\
&\Leftrightarrow &7=12+m \Leftrightarrow m=-5.
\end{eqnarray*}
Khi đó ta có $I=\displaystyle\int\limits_{-1}^5f(x)\mathrm{\,d}x=\displaystyle\int\limits_{-1}^2 f(x)\mathrm{\,d}x+\displaystyle\int\limits_2^5 f(x)\mathrm{\,d}x=\displaystyle\int\limits_{-1}^2(2x+3)\mathrm{\,d}x+\displaystyle\int\limits_2^5(3x^2-5) \mathrm{\,d}x=114$.}
\end{ex}

\begin{ex}%[Mức độ 3]%[BG12, Nguyễn Kiều Nhã Tú]%[2D4V2-2]
Cho hàm số $f(x)=\heva{&x^2+x+a& \text{khi }x \geq 0\\&2+bx&\text{khi } x<0}$ có đạo hàm trên $\mathbb{R}$ ( với $a,b$ là các số thực). Nếu $\displaystyle\int\limits_{-1}^1 f(x)\mathrm{d}x=\dfrac{m}{n}$ (với $m,n \in \mathbb{Z}^+$ và phân số $\dfrac{m}{n}$ tối giản) thì $m +2n$ bằng bao nhiêu?
\SA{$19$}
\loigiai{Hàm số $f(x)$ có đạo hàm trên $\mathbb{R}$ nên $f(x)$ liên tục và có đạo hàm tại $x=0$.
\begin{itemize}
\item 	$\lim\limits_{x \to 0^-} f(x)=\lim\limits_{x \to 0^-}(2+bx)=2$.
\item 	$\lim\limits_{x \to 0^+} f(x)=\lim\limits_{x \to 0^+}(x^2+x+a)=a$.
\item $f(0)=a$.
\item Hàm số liên tục tại $x=0$ nên suy ra $a=2$.
\end{itemize}
Ta lại có
\begin{itemize}
\item $\lim\limits_{x \to 0^-}\dfrac{f(x)-f(0)}{x-0}=\lim\limits_{x \to 0^-}\dfrac{(2+bx)-2}{x-0}=b$.
\item $\lim\limits_{x \to 0^+}\dfrac{f(x)-f(0)}{x-0}=\lim\limits_{x \to 0^+}\dfrac{(x^2+x+2)-2}{x-0}=\lim\limits_{x \to 0^+}(x+1)=1$.
\item Hàm số có đạo hàm tại $x=0$ nên suy ra $b=1$.
\end{itemize}
Do đó $f(x)=\heva{&x^2+x+2& \text{khi }x \geq 0\\&2+x&\text{khi } x<0}$. \\
Ta có
\allowdisplaybreaks
\begin{eqnarray*}
\displaystyle\int\limits_{-1}^1 f(x)\mathrm{d}x=\displaystyle\int\limits_{-1}^0 f(x)\mathrm{d}x+\displaystyle\int\limits_0^1 f(x)\mathrm{d}x&=&\displaystyle\int\limits_{-1}^0 (x+1)\mathrm{d}x+\displaystyle\int\limits_0^1 (x^2+x+2)\mathrm{d}x\\
&=& \dfrac{3}{2}+\dfrac{17}{6}\\
&=& \dfrac{13}{3}.
\end{eqnarray*}
Suy ra $m=13$ và $n=3$. \\
Vậy $m+2n=13+2\cdot 3=19$.
}
\end{ex}

\begin{ex}%[Mức độ 3]%[BG12, Nguyễn Kiều Nhã Tú]%[2D4V2-2]
Cho hàm số $f(x)=\heva{&ax+1,\ x\ge 1\\&x^2+b,\ x<1}$, với $a,b$ là các số thực. Biết rằng $f(x)$ liên tục và có đạo hàm trên $\mathbb{R}$. Tính $\displaystyle\int\limits_{-1}^{2} {f(x)}\mathrm{\,d}x$. (Kết quả làm tròn đến hàng phần trăm).
\SA{$8{,}67$}
\loigiai{
Hàm số liên tục và có đạo hàm trên $\mathbb{R}\Rightarrow$ hàm số liên tục và có đạo hàm tại $x=1$\\
$\Leftrightarrow\heva{&\lim\limits_{x\rightarrow 1^+}f(x)=\lim\limits_{x\rightarrow 1^-}f(x)=f(1)\\&f'(1^-)=f'(1^+).}$\\
Hay ta có hệ $\heva{&a+1=1+b\\&a=2}\Leftrightarrow \heva{&a=2\\&b=2}\Rightarrow f(x)=\heva{&2x+1,\ x\ge 1\\&x^2+2,\ x<1.}$ \\
Do đó ta có
\begin{eqnarray*}
\displaystyle\int\limits_{-1}^{2} {f(x)}\mathrm{\,d}x & = &\displaystyle\int\limits_{-1}^{1} {f(x)}\mathrm{\,d}x+\displaystyle\int\limits_{1}^{2} {f(x)}\mathrm{\,d}x\\
& = & \displaystyle\int\limits_{-1}^{1} {\left(x^2+2\right)}\mathrm{\,d}x+\displaystyle\int\limits_{1}^{2} {(2x+1)}\mathrm{\,d}x\\
& = & \dfrac{26}{3}.
\end{eqnarray*}
Vậy $\displaystyle\int\limits_{-1}^{2} {f(x)}\mathrm{\,d}x=\dfrac{26}{3}$.
}
\end{ex}

\begin{ex}%[2D4V2-2]
Biết $\displaystyle I=\int\limits_1^3\left| \dfrac{x^2-6x+8}{x+1}\right|\mathrm{\,d}x=\dfrac{a}{b}+c\ln \dfrac{d}{e}$ biết $a$ là số nguyên âm và $b$, $c$, $d$, $e\in \mathbb{Z}^*$; $\left(a,b\right)=1$, $\left(d,e\right)=1$. Giá trị của $a+b+c+d+e$ bằng:
\shortans{$25$}
\loigiai{
Ta có $x+1>0$, $\forall x\in \left[1;3\right]$, suy ra $\displaystyle \int\limits_1^3\left| \dfrac{x^2-6x+8}{x+1}\right|\mathrm{\,d}x=\int\limits_1^3\dfrac{\left| x^2-6x+8\right|}{x+1}\mathrm{\,d}x$.\\
Lại có $x^2-6x+8\le 0\Leftrightarrow 2\le x\le 4$. Do đó
\allowdisplaybreaks
\begin{eqnarray*}
I&=&\int\limits_1^3\dfrac{\left| x^2-6x+8\right|}{x+1}\mathrm{\,d}x=\int\limits_1^2\dfrac{x^2-6x+8}{x+1}\mathrm{\,d}x-\int\limits_2^3\dfrac{x^2-6x+8}{x+1}\mathrm{\,d}x\\
&=&\int\limits_1^2\left(x-7+\dfrac{15}{x+1}\right)\mathrm{\,d}x-\int\limits_2^3\left(x-7+\dfrac{15}{x+1}\right)\mathrm{\,d}x\\
&=&\left(-\dfrac{11}{2}+15\ln \dfrac{3}{2}\right)-\left(-\dfrac{9}{2}+15\ln \dfrac{4}{3}\right)\\
&=&-\dfrac{9}{2}+15\ln \dfrac{9}{8}.
\end{eqnarray*}
Vậy $\heva{&a=-9 \\&b=2 \\&c=15 \\&d=9 \\&e=8}\Rightarrow a+b+c+d+e=25$.
}
\end{ex}

\begin{ex}%[2D4V2-2]
Cho hàm số $y=f(x)$ liên tục trên $\mathbb{R}$. Hàm số $y=f'(x)$ có đồ thị $(C)$ như hình vẽ, $(C)$ cắt trục $Ox$ tại ba điểm phân biệt có hoành độ $a<b<c$.
\begin{center}
\begin{tikzpicture}[line join = round, line cap = round,>=stealth,x = 1cm,y = .6cm]
%Vẽ hệ trục Oxy
\draw[->] (-2.5,0)--(0,0) node[below right]{$O$}--(4.5,0) node[below]{$x$};
\draw (-1.1,.3) node {$a$} (1.1,.3) node {$b$}  (3.1,-.3) node{$c$} (2.9,2) node[rotate=70]{$(C)$};
\draw[->] (0,-4.5)--(0,4.5) node[right]{$y$};
\draw[samples=200,domain=-1.42:3.42,smooth] plot (\x, {(\x)^3-3*(\x)^2-\x+3});
\end{tikzpicture}
\end{center}
Biết rằng diện tích hình phẳng giới hạn bởi $(C)\colon y=f'(x)$ và $O x$ bằng $15$, $f(a)=5$, $f(c)=6$. Tính $f(b)$.
\shortans{$13$}
\loigiai{
Diện tích hình phẳng giới hạn bởi $(C)\colon y=f'(x)$ và $Ox$ bằng $15$, do đó
\[
15=\displaystyle\int\limits_a^c\left|f'(x)\right| \mathrm{\,d} x=\displaystyle\int\limits_a^b\left|f'(x)\right| \mathrm{\,d} x+\displaystyle\int\limits_b^c\left|f'(x)\right| \mathrm{\,d} x=2 f(b)-f(a)-f(c) .
\]
Suy ra $2 f(b)=15+f(a)+f(c) \Rightarrow f(b)=13$.
}
\end{ex}

\begin{ex}%[Cau-1]%[2D4V2-2]
Cho $\displaystyle \int\limits_{1}^{4}\sqrt{\dfrac{1}{4x}+\dfrac{\sqrt{x}+ \mathrm{e}^x}{\sqrt{x}\cdot \mathrm{e}^{2x}}}\mathrm{\,d}x=a+\mathrm{e}^b-\mathrm{e}^c$ với $a$, $b$, $c$ là các số nguyên. Tính giá trị của biểu thức $S=a+b+c$.
\shortans{$-4$}
\loigiai{
Ta có
\begin{align*}
\displaystyle \int\limits_{1}^{4} \sqrt{\dfrac{1}{4x} + \dfrac{\sqrt{x} + \mathrm{e}^x}{\sqrt{x}\cdot \mathrm{e}^{2x}}} \mathrm{\,d}x
&= \displaystyle\int\limits_{1}^{4} \sqrt{\left(\dfrac{1}{2\sqrt{x}}\right)^2 + 2\dfrac{1}{2\sqrt{x}\cdot \mathrm{e}^x}+\left(\dfrac{1}{\mathrm{e}^x}\right)^2} \mathrm{\,d}x\\
&= \displaystyle\int\limits_{1}^{4}\sqrt{\left(\dfrac{1}{2\sqrt{x}} + \dfrac{1}{\mathrm{e}^x}\right)^2}\mathrm{d}x\\
&=\displaystyle\int\limits_{1}^{4}\left(\dfrac{1}{2\sqrt{x}}+\dfrac{1}{\mathrm{e}^x}\right) \mathrm{\,d}x\\
&=\left(\sqrt{x}-\mathrm{e}^{-x}\right)\bigg|_1^4\\
&=1-\mathrm{e}^{-4}+\mathrm{e}^{-1}\\
&=a+\mathrm{e}^b-\mathrm{e}^c.
\end{align*}
$\Rightarrow \heva{&a=1\\&b=-1\\&c=-4.}$\\
Vậy $a+b+c=1+(-1)+(-4)=-4$.
}
\end{ex}

\begin{ex}%[HSA - đợt 2, Nguyen The Tuan Vu]%[2D4V2-2]
Cho $\displaystyle\int\limits_{2}^{3} \left(\dfrac{x-1}{x}\right)^2 \mathrm{\,d}x=a+b\cdot\ln 2+c\cdot\ln 3$. Biết rằng tồn tại duy nhất bộ ba số hữu tỉ $a$, $b$, $c$ thỏa mãn đề bài. Tổng $6a+b+c$ có giá trị bằng bao nhiêu?
\shortans{$7$}
\loigiai{
Ta có $\begin{aligned}[t]
\displaystyle\int\limits_{2}^{3} \left(\dfrac{x-1}{x}\right)^2 \mathrm{\,d}x&=\int\limits_{2}^{3} \left(1-\dfrac{1}{x}\right)^2 \mathrm{\,d}x=\int\limits_{2}^{3} \left(1-\dfrac{2}{x}+\dfrac{1}{x^2}\right) \mathrm{\,d}x=\left(x-2\ln\left|x\right|-\dfrac{1}{x}\right)\Big|^3_2\\
&=\left(3-2\ln 3-\dfrac{1}{3}\right)-\left(2-2\ln 2-\dfrac{1}{2}\right)=\dfrac{7}{6}+2\ln 2-2\ln 3.
\end{aligned}$\\
$\Rightarrow a=\dfrac{7}{6}$, $b=2$ và $c=-2$.\\
$\Rightarrow 6a+b+c=7$.
}
\end{ex}

\begin{ex}%[Đề thi GHK2, THPT Nguyễn Hữu Huân, TPHCM năm học 2024-2025]%[BCTuan, dự án 12EX3-26]%[2D4V2-2]
\immini
{Cho hàm số $y=f(x)$ liên tục trên đoạn $[-3;3]$ có đồ thị như hình vẽ. Giá trị của tích phân $\displaystyle\int\limits_{-3}^{2}{f(x)}\mathrm{\,d}x$ bằng bao nhiêu?  (làm tròn kết quả đến hàng phần mười)}
{\begin{tikzpicture}[line cap=round,line join=round,>=stealth,font=\footnotesize,scale=.7]
\tikzset{declare function={xmin=-4;xmax=4;ymin=-5;ymax=2;},
smooth,samples=450
}
\draw[->] (xmin,0)--(xmax,0) node[shift={(-100:7pt)}]{$ x $};
\draw[->] (0,ymin)--(0,ymax) node[shift={(190:7pt)}]{$ y $};
\fill (0,0)circle (1pt) node[shift={(225:10pt)}]{$ O $};
\draw[dashed](-3,0)node[above]{$-3$}|-(0,-4)node[right]{$-4$}
(-1,0)node[below]{$-1$}|-(0,1)node[right]{$1$}
(1,0)node[above]{$1$}|-(0,-1)node[left]{$-1$}
(2,0)circle (1pt)node[above]{$2$}
(3,0)node[above]{$3$}|-(0,-3)node[left]{$-3$};
\draw (-3,-4) circle (1pt) 	(-1,1) circle (1pt)
(1,-1) circle (1pt)	(3,-3) circle (1pt);
\draw (-3,-4)--(-1,1)--(1,-1)--(2,0)--(3,-3);
\clip (xmin,ymin) rectangle (xmax,ymax);
\end{tikzpicture}}
\shortans{$-3{,}5$}
\loigiai{
\begin{itemize}
\item Trên đoạn $[-3;-1]$ đồ thị $f(x)$ là đường thẳng đi qua hai điểm $(-3,-4)$ và $(-1;1)$ nên có phương trình $f(x)=\dfrac{5}{2}x+\dfrac{7}{2}$.
\item Trên đoạn $[-1;1]$ đồ thị $f(x)$ là đường thẳng đi qua hai điểm $(-1,1)$ và $(1;-1)$ nên có phương trình $f(x)=-x$.
\item Trên đoạn $[1;2]$ đồ thị $f(x)$ là đường thẳng đi qua hai điểm $(1,-1)$ và $(2;0)$ nên có phương trình $f(x)=x-2$.
\end{itemize}
Do đó
\begin{align*}
\displaystyle\int\limits_{-3}^{2}{f(x)}\mathrm{\,d}x=\displaystyle\int\limits_{-3}^{-1}{\left(\dfrac{5}{2}x+\dfrac{7}{2}\right)}\mathrm{\,d}x+\displaystyle\int\limits_{-1}^{1}{(-x)}\mathrm{\,d}x+\displaystyle\int\limits_{1}^{2}{(x-2)}\mathrm{\,d}x=-\dfrac{7}{2}=-3{,}5.
\end{align*}

}
\end{ex}

\begin{ex}%[Đề kiểm tra giữa kì 2 - THPT Thanh Miện, Hải Dương, Nguyễn Anh Tuấn, Dự án 12EX-3-2026]%[2D4V2-2]
Cho hàm số $f(x) = \heva{&2x^2 - 1 \quad \text{khi } x < 0 \\&x - 1 \quad \text{khi } 0 \le x \le 2 \\& 5 - 2x \quad \text{khi } x > 2}$ và $I = \displaystyle \int\limits_{-5}^{9} \dfrac{1}{7} f(x) \mathrm{\,d} x$.\\ Biết $I = \dfrac{a}{b}$, ($a, b \in \mathbb{N}^*$; $(a,b)=1$). Tính $a - b$.
\shortans[oly]{88}
\loigiai{
Ta có $I = \dfrac{1}{7} \displaystyle \int\limits_{-5}^{9} f(x) \mathrm{\,d} x = \dfrac{1}{7} \left( \int\limits_{-5}^{0} f(x) \mathrm{\,d} x + \int\limits_{0}^{2} f(x) \mathrm{\,d} x + \int\limits_{2}^{9} f(x) \mathrm{\,d} x \right)$. \\
\begin{itemize}
\item $I_1 = \displaystyle \int\limits_{-5}^{0} (2x^2 - 1) \mathrm{\,d} x = \left. \left( \dfrac{2x^3}{3} - x \right) \right|_{-5}^{0} = 0 - \left( \dfrac{2(-5)^3}{3} - (-5) \right) = \dfrac{250}{3} - 5 = \dfrac{235}{3}$.
\item $I_2 = \displaystyle \int\limits_{0}^{2} (x - 1) \mathrm{\,d} x = \left. \left( \dfrac{x^2}{2} - x \right) \right|_{0}^{2} = (2 - 2) - 0 = 0$.
\item $I_3 = \displaystyle \int\limits_{2}^{9} (5 - 2x) \mathrm{\,d} x = \left. (5x - x^2) \right|_{2}^{9} = (45 - 81) - (10 - 4) = -36 - 6 = -42$.
\end{itemize}
Khi đó $I = \dfrac{1}{7} \left( \dfrac{235}{3} + 0 - 42 \right) = \dfrac{1}{7} \cdot \dfrac{109}{3} = \dfrac{109}{21}$. \\
Vì $I = \dfrac{a}{b}$ với $a, b \in \mathbb{N}^*$ và $(a, b) = 1$ nên $a = 109$, $b = 21$. \\
Vậy $a - b = 109 - 21 = 88$.
}
\end{ex}

\begin{ex}%[2D4V2-2]
Cho hàm số $g(x)=x^3\ln x$. Giá trị của $\displaystyle\int\limits_1^3\left(g'(x) + \dfrac{x^2-6}{x} \right)\mathrm{\,d}x$ là bao nhiêu? (kết quả làm tròn đến hàng đơn vị)
\shortans{27}
\loigiai{
Ta có $g(3)-g(1)=27\ln3$.\\
Mặt khác $\displaystyle\int\limits_1^3\left(\dfrac{x^2-6}{x}\right)\mathrm{\,d}x=\displaystyle\int\limits_1^3\left( x-\dfrac{6}{x} \right)\mathrm{\,d}x = \left( \dfrac{x^2}{2}-6\ln|x|\right) \bigg|_1^3=4-6\ln3$.\\
Vậy $\displaystyle\int\limits_1^3\left(g'(x) + \dfrac{x^2-6}{x}\right) \mathrm{\,d}x = g(x)\bigg|_1^3 + \displaystyle\int\limits_1^3\left( \dfrac{x^2-6}{x} \right)\mathrm{\,d}x=27\ln 3 + 4 - 6\ln 3 = 4 + 21\ln 3 \approx 27$.
}
\end{ex}

\begin{ex}%[HSA-đợt 1, Nguyễn Cường]%[2D4V2-2]
Biết $\displaystyle \int\limits_1^3\left|\dfrac{x^2-6x+8}{x+1}\right|\mathrm{\,d}x=\dfrac{a}{b}+c\ln\dfrac{d}{e}$ biết $a$ là số nguyên âm và $b$, $c$, $d$, $e\in\mathbb{Z}^+$ và $(a,b)=1$, $(d,e)=1$, đồng thời $d$, $e<10$. Giá trị của $a+b+c+d+e$ bằng bao nhiêu?
\shortans{32}
\loigiai{
Xét $\dfrac{x^2-6x+8}{x+1}=0\Leftrightarrow x^2-6x+8=0\Leftrightarrow\hoac{&x=2\\&x=4.}$\\
Khi đó
\allowdisplaybreaks
\begin{align*}
\displaystyle \int\limits_1^3\left|\dfrac{x^2-6x+8}{x+1}\right|\mathrm{\,d}x&=\displaystyle \int\limits_1^2\dfrac{x^2-6x+8}{x+1}\mathrm{\,d}x-\displaystyle \int\limits_2^3\dfrac{x^2-6x+8}{x+1}\mathrm{\,d}x\\
&=\displaystyle \int\limits_1^2\dfrac{x^2-6x+8}{x+1}\mathrm{\,d}x-\displaystyle \int\limits_2^3\dfrac{x^2-6x+8}{x+1}\mathrm{\,d}x\\
&=\displaystyle \int\limits_1^2\left(x-7+\dfrac{15}{x+1}\right)\mathrm{\,d}x-\displaystyle \int\limits_2^3\left(x-7+\dfrac{15}{x+1}\right)\mathrm{\,d}x\\
&=\left(\dfrac{x^2}{2}-7x+15\ln(x+1)\right)\Bigg|_1^2-\left(\dfrac{x^2}{2}-7x+15\ln(x+1)\right)\Bigg|_2^3\\
&=\left(-12+15\ln 3+\dfrac{13}{2}-15\ln 2 \right)-\left( -\dfrac{33}{2}+30\ln 2+12-15\ln 3 \right)\\
&=-1+30\ln 3-45\ln 2\\
&=-1+15\ln 9-15\ln 8\\
&=-1+15\ln\dfrac{9}{8}.
\end{align*}
Khi đó, $a=-1$, $b=1$, $c=15$, $d=9$ và $e=8$.\\
Vậy $a+b+c+d+e=-1+1+15+9+8=32$.
}
\end{ex}

\begin{ex}%[HSA-đợt 1, Nguyễn Cường]%[2D4V2-2]
Cho $\displaystyle\int\limits_0^1\dfrac{1}{x^2+3x+2}\mathrm{\,d}x=a\ln 2+b\ln 3$, với $a$, $b$ là các số nguyên. Khi đó $a+b$ bằng bao nhiêu?
\shortans{1}
\loigiai{
Ta có
\allowdisplaybreaks
\begin{eqnarray*}
\displaystyle\int\limits_0^1\dfrac{1}{x^2+3x+2}\mathrm{\,d}x&=\displaystyle\int\limits_0^1\dfrac{1}{(x+1)(x+2)}\mathrm{\,d}x\\
&=\displaystyle\int\limits_0^1\left(\dfrac{1}{x+1}-\dfrac{1}{x+2}\right)\mathrm{\,d}x\\
&=\ln\left(\dfrac{x+1}{x+2}\right)\Bigg|_0^1\\
&=\ln\dfrac{2}{3}-\ln\dfrac{1}{2}\\
&=2\ln 2-\ln 3.
\end{eqnarray*}
Do đó $a=2$ và $b=1$.\\
Vậy $a+b=2-1=1$.
}
\end{ex}

\begin{ex}%[2D4V2-2]
Biết $\displaystyle\int\limits_0^1\dfrac{x^2-2}{x+1}\mathrm{\,d}x=-\dfrac{1}{m}+n\cdot\ln2$, Với $m$, $n$ là các số nguyên. Tính $m-n$.
\par\shortans[]{$3$}
\loigiai{
Phương pháp giải: Tính tích phân.\\
$\displaystyle\int\limits_0^1\dfrac{x^2-2}{x+1}\mathrm{\,d}x=\displaystyle\int\limits_0^1(x-1)\mathrm{\,d}x-\displaystyle\int\limits_0^1\dfrac{\mathrm{\,d}x}{x+1}=\left.\dfrac{(x-1)^2}{2}\right|_0^1-\left.\ln|x+1|\right|_0^1=\dfrac{-1}{2}-\ln2$.\\
Suy ra $m=2$, $n=-1\Rightarrow m-n=3$.
}
\end{ex}

\begin{ex}%[2D4V2-3]%[To 20 - Dot 17 - Chuong 4 - Bai 3 - CD - De 1 - TLN]%[Lê Quân]
Biết rằng $I=\displaystyle\int\limits_0^{\tfrac{\pi}{2}}\dfrac{-4 \sin x+7 \cos x}{2 \sin x+3 \cos x}\mathrm{\,d} x=\dfrac{\pi}{a}+2 \ln \dfrac{b}{c}$ với $a>0 ; b, c \in \mathbb{N}^*$ và phân số $\dfrac{b}{c}$ tối giản. Hãy tính giá trị biểu thức $P=a+b+c$.
\shortans[3]{$7$}
\loigiai{
Xét đồng nhất thức
\begin{eqnarray}
-4 \sin x+7 \cos x&=&A(2 \sin x+3 \cos x)+B(2 \cos x-3 \sin x)\\
&=&(2 A-3 B) \sin x+(3 A+2 B) \cos x.
\end{eqnarray}
Do đó $\heva{&2 A-3 B=-4 \\& 3 A+2 B=7} \Rightarrow\heva{& A=1 \\& B=2.}$\\ Ta có
\begin{eqnarray*}
I&=&\displaystyle\int\limits_0^{\dfrac{\pi}{2}}\dfrac{-4 \sin x+7 \cos x}{2 \sin x+3 \cos x}\mathrm{\,d} x\\
&=&\displaystyle\int\limits_0^{\dfrac{\pi}{2}}\left(1+\dfrac{2(2 \sin x+3 \cos x)^{\prime}}{2 \sin x+3 \cos x}\right) \mathrm{\,d} x \\
&=&\left(x+2 \ln |2 \sin x+3 \cos x|\right)\Big|_0 ^{\tfrac{\pi}{2}}\\
&=&\dfrac{\pi}{2}+2 \ln \dfrac{2}{3}\\
&\Rightarrow& a=2, b=2, c=3.
\end{eqnarray*}
Vậy $P=a+b+c=7$.}
\end{ex}

\begin{ex}%[2D4V2-3]%[Tổ 20 - Đợt 17 - Chương 4 - - CD]%[Nguyễn Hồng Thạch]
Tích phân $\displaystyle\int\limits_{0}^{\frac{\pi}{3}}3\tan^2 x\mathrm{\,d}x=a\sqrt{b}+c\pi$. Tính $a+b+c$.
\shortans{$5$}
\loigiai
{Ta có $\displaystyle\int\limits_{0}^{\frac{\pi}{3}}3\tan^2 x\mathrm{\,d}x=3\displaystyle\int\limits_{0}^{\frac{\pi}{3}}\left(\dfrac{1}{\cos^2x}-1\right)\mathrm{\,d}x=3\left(\tan x-x\right)\Big|_{0}^{\frac{\pi}{3}}=3\sqrt{3}-\pi$.\\
Suy ra $a=3$, $b=3$, $c=-1$. Vậy $a+b+c=5$.
}
\end{ex}

\begin{ex}%[2D4V2-3]%[Tổ 20 - Đợt 17 - Chương 4 - - CD]%[Nguyễn Hồng Thạch]
Tích phân $\displaystyle\int\limits_{0}^{\frac{\pi}{3}}\left|\cos 2x\right|\mathrm{\,d}x=a-\dfrac{\sqrt{b}}{c}$. Tính $a+b+c$.
\shortans{$8$}
\loigiai
{Ta có \begin{eqnarray*}
\displaystyle\int\limits_{0}^{\frac{\pi}{3}}\left|\cos 2x\right|\mathrm{\,d}x&=&\displaystyle\int\limits_{0}^{\frac{\pi}{4}}\cos 2x\mathrm{\,d}x+\displaystyle\int\limits_{\frac{\pi}{4}}^{\frac{\pi}{3}}-\cos 2x\mathrm{\,d}x\\
&=&\dfrac{1}{2}\sin 2x\Big|_{0}^{\frac{\pi}{4}}-\dfrac{1}{2}\sin 2x\Big|_{\frac{\pi}{4}}^{\frac{\pi}{3}}\\
&=&\dfrac{1}{2}-\dfrac{1}{2}\left(\dfrac{\sqrt{3}}{2}-1\right)=1-\dfrac{\sqrt{3}}{4}
\end{eqnarray*}
Suy ra $a=1$, $b=3$, $c=4$. Vậy $a+b+c=8$.
}
\end{ex}

\begin{ex}%[2D4V2-3]
Tính tích phân $I=\displaystyle\int\limits _{0}^{2\pi }\sqrt{1+\cos 2x} \mathrm{\,d}x $ (\textit{làm tròn đến hàng phần trăm}).
\shortans{$5{,}66$}
\loigiai{
Ta có  $I=\displaystyle\int\limits _{0}^{2\pi }\sqrt{1+\cos 2x} \mathrm{\,d}x =\sqrt{2}\displaystyle\int\limits _{0}^{2\pi }|\cos x|\mathrm{\,d}x $.\\
Ta có $\cos x\ge 0, \forall x\in \left[0;\dfrac{\pi}{2} \right]\cup\left[\dfrac{3\pi}{2};2\pi \right] $ và $\cos x\le 0, \forall x\in \left[\dfrac{\pi}{2};\dfrac{3\pi}{2} \right] $. Khi đó
\begin{eqnarray*}
I&=&\sqrt{2}\displaystyle\int\limits _{0}^{\tfrac{\pi}{2} }\cos x\mathrm{\,d}x-\sqrt{2}\displaystyle\int\limits _{\tfrac{\pi}{2} }^{\tfrac{3\pi}{2} }\cos x\mathrm{\,d}x+\sqrt{2}\displaystyle\int\limits _{\tfrac{3\pi}{2} }^{2\pi}\cos x\mathrm{\,d}x\\
&=&\sqrt{2}\sin x\Bigg|_{0}^{\tfrac{\pi}{2} } -\sqrt{2}\sin x\Bigg|_{\tfrac{\pi}{2} }^{\tfrac{3\pi}{2} }+\sqrt{2}\sin x\Bigg|_{\tfrac{3\pi}{2} }^{2\pi}\\
&=&4\sqrt{2}\approx5{,}66.
\end{eqnarray*}		}
\end{ex}

\begin{ex}%[2D4V2-3]
Tính tích phân $I=\displaystyle\int\limits _{0}^{2\pi }\sqrt{1-\sin 2x} \mathrm{\,d}x $, (\textit{làm tròn đến hàng phần trăm}).
\shortans{$0{,}31$}
\loigiai{
Ta có  $I=\displaystyle\int\limits _{0}^{2\pi }\sqrt{1-\sin 2x} \mathrm{\,d}x =\displaystyle\int\limits _{0}^{2\pi }|\sin  x-\cos x|\mathrm{\,d}x $.\\
Ta có $\sin x-\cos x\le 0, \forall x\in \left[0;\dfrac{\pi}{4} \right]\cup\left[\dfrac{5\pi}{4};2\pi \right] $ và $\sin x-\cos x\ge 0, \forall x\in \left[\dfrac{\pi}{4};\dfrac{5\pi}{4} \right]$. Khi đó
\begin{eqnarray*}
I&=&\displaystyle\int\limits _{0}^{\tfrac{\pi}{4} }\left( \cos x-\sin x\right) \mathrm{\,d}x+\displaystyle\int\limits _{\tfrac{\pi}{4} }^{\tfrac{5\pi}{4} }\left( \sin x-\cos x\right) \mathrm{\,d}x+\displaystyle\int\limits _{\tfrac{5\pi}{4} }^{2\pi}\left( \cos x-\sin x\right) \mathrm{\,d}x\\
&=&\left( \sin x+\cos x\right) \Bigg|_{0}^{\tfrac{\pi}{4} }+\left( -\cos x-\sin x\right) \Bigg|_{\tfrac{\pi}{4} }^{\tfrac{5\pi}{4} }+\left( \sin x+\cos x\right) \Bigg|_{\tfrac{5\pi}{4} }^{2\pi}\\
&=&4\sqrt{2}\approx5{,}66.
\end{eqnarray*}
}
\end{ex}

\begin{ex}%[2D4V2-3]%[EX-TF-TLN-2024-Huyên Nguyễn]
Biết giá trị của tích phân $\displaystyle \int\limits_0^{\tfrac{\pi }{2}} \sqrt{1 - \sin 2x }\mathrm{\,d} x =a\sqrt{2}+b$. Tính $a-2b$.
\shortans{$6$}
\loigiai{Ta có
\begin{eqnarray*}
\displaystyle \int\limits_0^{\tfrac{\pi }{2}} \sqrt{1 - \sin 2x }\mathrm{\,d} x &=& \int\limits_0^{\tfrac{\pi }{2}} \sqrt{\sin^2 x - 2\sin x\cos x+\cos^2 x }\mathrm{\,d} x \\
&=&\int\limits_0^{\tfrac{\pi }{2}} \sqrt{\left(\sin x-\cos x\right)^2 }\mathrm{\,d} x \\
&=& \displaystyle \int\limits_0^{\tfrac{\pi }{2}}  \left| \sin x - \cos x \right|\mathrm{\,d}x\\
& = &\displaystyle \int\limits_0^{\tfrac{\pi}{4}} \left( \cos x  - \sin x \right) \mathrm{\,d} x  + \displaystyle \int\limits_{\tfrac{\pi }{4}}^{\tfrac{\pi }{2}} \left(\sin x - \cos x \right) \mathrm{\,d} x \\
&=&\left( \sin x+\cos x\right)\big|_0^{\tfrac{\pi }{4}}+\left( -\cos x-\sin x \right)\big|_{\tfrac{\pi }{4}}^{\tfrac{\pi }{2}}\\
& = &2\sqrt{2} -2.
\end{eqnarray*}
Vậy $a=2$ và $b=-2$ nên $a-2b =6$.}
\end{ex}

\begin{ex}%[2D4V2-3]%[EX-TF-TLN-2024-Huyên Nguyễn]
Một vật chuyển động với gia tốc được cho bởi hàm số $a(t)=5\cos t$ (m/s$^2$). Lúc bắt đầu chuyển động vật có vận tốc $0$ m/s. Tính quãng đường của vật đã di chuyển (đơn vị: m) tính đến thời điểm  $t=\dfrac{\pi}{2}$ (s).
\shortans{$5$}
\loigiai{
Vận tốc của vật được biểu diễn bởi hàm số $v(t)=\displaystyle\int a(t) \mathrm{\,d}t=\displaystyle\int 5\cos t \mathrm{\,d}t=5\sin t+C$.\\
Lúc bắt đầu vật có vận tốc $v(0)=0$ nên $C=0$ và $v(t)=5\sin 5t$ (m/s).\\
Quãng đường của vật đã di chuyển tính đến thời điểm  $t=\dfrac{\pi}{2}$ (s) là \[S\left(\dfrac{\pi}{2}\right)=S\left(\dfrac{\pi}{2}\right)-S(0)=\displaystyle\int\limits_0^{\tfrac{\pi}{2}} v(t) \mathrm{\,d}t=\displaystyle\int\limits_0^{\tfrac{\pi}{2}} 5\sin t \mathrm{\,d}t= \left(-5\cos t\right)\big|_{0}^{\tfrac{\pi }{2}}=
-5\cos \left(\dfrac{\pi}{2}\right)+5\cos 0=5\text{ (m)}.\]
}
\end{ex}

\begin{ex}%[MĐ3]%[2D4V2-3]
Biết $\displaystyle\int\limits _0^{\tfrac{\pi}{4}}\dfrac{1}{\sin ^2x\cdot \cos ^2x} \mathrm{d}x=\dfrac{a\sqrt{3}}{b} (a,b\in \mathbb{Z})$. Tính $P=\dfrac{a-2b}{b}$ (làm tròn kết quả đến hàng phần trăm)
\shortans{$1{,}33$}
\loigiai{
Ta có
\begin{eqnarray*}
& & \displaystyle\int\limits _{\dfrac{\pi}{4}}^{\dfrac{\pi}{3}}\dfrac{1}{\sin ^2x\cdot \cos ^2x} \mathrm{d}x=\displaystyle\int\limits _{\dfrac{\pi}{4}}^{\dfrac{\pi}{3}}\dfrac{\sin^2 x+\cos^2 x}{\sin^2 x\cdot \cos^2 x} \mathrm{d}x\\
&=& \displaystyle\int\limits _{\tfrac{\pi}{4}}^{\tfrac{\pi}{3}}\left(\dfrac{1}{\cos ^2x}+\dfrac{1}{\sin ^2x} \right)\mathrm{d}x\\
&= &(\tan x-\cot x)\bigg|_{\tfrac{\pi}{4}}^{\tfrac{\pi}{3}}=\dfrac{2\sqrt{3}}{3}.
\end{eqnarray*}
Vậy $P=\dfrac{2-2\cdot 3}{3}=-\dfrac{4}{3} \approx 1{,}33$.
}
\end{ex}

\begin{ex}%[MĐ3]%[2D4V2-3]
Có bao nhiêu số thực $b$ thuộc khoảng $(\pi; 3\pi)$ sao cho $\displaystyle\int\limits _{\pi}^b 4\cos 2x\mathrm{d}x=1$?
\shortans{$1$}
\loigiai{
Ta có $\displaystyle\int\limits _{\pi}^b4\cos 2x\mathrm{d}x=1\Leftrightarrow 2\sin 2x\bigg|_{\pi}^b=1
\Leftrightarrow \sin 2b=\dfrac{1}{2} \Leftrightarrow \hoac{&2b=\dfrac{\pi}{6}+k2\pi \\ &2b=\dfrac{5\pi}{6}+k2\pi} \Leftrightarrow \hoac{&b=\dfrac{\pi}{12}+k\pi \\ &b=\dfrac{5\pi}{12}+k\pi}, (k\in \mathbb{Z})$.
\begin{itemize}
\item[•] Với $b=\dfrac{\pi}{12}+k\pi$:
Mà $b\in (\pi; 3\pi)\Rightarrow 1< \dfrac{1}{12}+k < 3,(k\in \mathbb{Z})\Rightarrow k=1;2\Rightarrow b=\dfrac{13\pi}{12};\dfrac{25\pi}{12}$.
\item[•] Với $b=\dfrac{5\pi}{12}+k\pi$:
Mà $b\in (\pi; 3\pi)\Rightarrow 1< \dfrac{5}{12}+k < 3,(k\in \mathbb{Z})\Rightarrow k=1;2\Rightarrow b=\dfrac{17\pi}{12};\dfrac{29\pi}{12}$.
\end{itemize}
Vậy có $4$ số thực $b$ thỏa mãn yêu cầu bài toán.
}
\end{ex}

\begin{ex}%[Mức độ 3]%[Dat Thai, Dự án giảng 12 New]%[2D4V2-3]
Tính $\displaystyle\int\limits_0^{\tfrac{\pi}{6}}\sin 4x\cos x\mathrm{\,d}x$.
\shortans{$0{,}35$}
\loigiai{
%$
%\begin{aligned}
%\displaystyle\int\limits_0^{\tfrac{\pi}{6}}\sin 4x\cos x\mathrm{\,d}x&=\dfrac{1}{2}\displaystyle\int\limits_0^{\tfrac{\pi}{6}}[\sin(3x)+\sin(5x)]\mathrm{\,d}x=-\dfrac{1}{2}\left(\dfrac{\cos 3x}{3}+\dfrac{\cos 5x}{5}\right)\bigg|_0^\tfrac{\pi}{6}\\
%&=-\dfrac{1}{2}\left(\dfrac{0}{3}-\dfrac{\sqrt{3}}{10}\right)+\dfrac{1}{2}\left(\dfrac{1}{3}+\dfrac{1}{5}\right)=\dfrac{16+3\sqrt{3}}{60} =\approx 0{,}35.
%\end{aligned}
%$
}
\end{ex}

\begin{ex}%[Mức độ 3]%[Dat Thai, Dự án giảng 12 New]%[2D4V2-3]
Cho tích phân $\displaystyle\int\limits_{0}^{1}\left(\cos \dfrac{\pi}{2} x+\dfrac{x}{(x+1)^{3}}\right) \mathrm{\,d}x=\dfrac{a}{\pi}+\dfrac{b}{c}$; vơi $a, b, c$ là những số nguyên dương và phân số $\dfrac{b}{c}$ tối giản. Tính giá trị của biểu thức $T=a+b+c$.
\shortans{$11$}
\loigiai{
Ta có: $\displaystyle\dfrac{x}{(x+1)^3}=\dfrac{x+1}{(x+1)^{3}}-\dfrac{1}{(x+1)^{3}}=\dfrac{1}{(x+1)^{2}}-\dfrac{1}{(x+1)^{3}}$.
\\Suy ra $\displaystyle I=\int\limits_{0}^{1}\left(\dfrac{x}{(x+1)^{3}}\right) d x=\int\limits_{0}^{1}\left(\dfrac{1}{(x+1)^{2}}-\dfrac{1}{(x+1)^{3}}\right) d x=\left.\left(-\dfrac{1}{x+1}+\dfrac{1}{2(x+1)^{2}}\right)\right|_{0} ^{1}=\dfrac{1}{8}$
\\Và $J=\displaystyle\int\limits_{0}^{1} \cos \dfrac{\pi}{2} x d x=\left.\dfrac{2}{\pi} \sin \dfrac{\pi}{2} x\right|_{0} ^{1}=\dfrac{2}{\pi}$.
\\Do đó $\displaystyle\int\limits_{0}^{1}\left(\cos \dfrac{\pi}{2} x+\dfrac{x}{(x+1)^{3}}\right) \mathrm{\,d}x=\dfrac{2}{\pi}+\dfrac{1}{8}=\dfrac{a}{\pi}+\dfrac{b}{c}$
\\Suy ra $a=2 ; b=1 ; c=8 \Rightarrow T=a+b+c=11$.
}
\end{ex}

\begin{ex}%[Cau-4]%[2D4V2-3]
Cho $M$, $N$ là các số thực, xét hàm số $f(x)=M\cdot \sin \pi x + N\cdot \cos \pi x$ thỏa mãn $f(1)=3$ và $\displaystyle \int\limits_{0}^{\tfrac{1}{2}} f(x) \mathrm{\,d}x= -\dfrac{1}{\pi}$. Tính $f'\left(\dfrac{1}{4}\right)$. (Kết quả làm tròn đến hàng phần mười).
\shortans{$11{,}1$}
\loigiai{
Ta có $f(1)=3\Leftrightarrow M\cdot \sin \pi+N\cdot \cos \pi=3\Leftrightarrow N=-3$.\\
Mặt khác \\
\begin{align*}
\displaystyle \int\limits_{0}^{\tfrac{1}{2}} f(x) \mathrm{\,d}x = -\dfrac{1}{\pi} &\Leftrightarrow \int\limits_{0}^{\tfrac{1}{2}} \left(M\cdot \sin \pi x + N \cdot \cos \pi x\right) \mathrm{d}x = -\dfrac{1}{\pi}\\
&\Leftrightarrow \left(-\dfrac{M}{\pi}\cos \pi x - \dfrac{3}{\pi} \sin \pi x\right)\bigg|_0^{\tfrac{1}{2}} = -\dfrac{1}{\pi}\\
&\Leftrightarrow -\dfrac{3}{\pi} + \dfrac{M}{\pi} = -\dfrac{1}{\pi}\\
&\Leftrightarrow M = 2.
\end{align*}
Ta được $f(x) = 2\cdot \sin \pi x - 3\cdot \cos \pi x$ nên $f'(x) = 2\pi \cos \pi x + 3\pi \sin \pi x$.\\
Vậy $f'\left(\dfrac{1}{4}\right) = \dfrac{5\pi \sqrt{2}}{2} \approx 11{,}1$.
}
\end{ex}

\begin{ex}%[HSA-ĐỢT 2-Nguyễn Mạnh Trung Anh]%[2D4V2-3]
Biết $\displaystyle\int\limits _0^{\tfrac{\pi}{4}}\dfrac{1}{\sin ^2x\cdot \cos ^2x} \mathrm{d}x=\dfrac{a\sqrt{3}}{b} (a,b\in \mathbb{Z})$. Tính $P=\dfrac{a-2b}{b}$. (Làm tròn kết quả đến hàng phần trăm)
\par\shortans{1{,}33}
\loigiai{
Ta có
\begin{eqnarray*}
& & \displaystyle\int\limits _{\tfrac{\pi}{4}}^{\tfrac{\pi}{3}}\dfrac{1}{\sin ^2x\cdot \cos ^2x} \mathrm{d}x=\displaystyle\int\limits _{\tfrac{\pi}{4}}^{\tfrac{\pi}{3}}\dfrac{\sin^2 x+\cos^2 x}{\sin^2 x\cdot \cos^2 x} \mathrm{d}x\\
&=& \displaystyle\int\limits _{\tfrac{\pi}{4}}^{\tfrac{\pi}{3}}\left(\dfrac{1}{\cos ^2x}+\dfrac{1}{\sin ^2x} \right)\mathrm{d}x\\
&= &(\tan x-\cot x)\bigg|_{\tfrac{\pi}{4}}^{\tfrac{\pi}{3}}=\dfrac{2\sqrt{3}}{3}.
\end{eqnarray*}
Vậy $P=\dfrac{2-2\cdot 3}{3}=-\dfrac{4}{3} \approx 1{,}33$.
}
\end{ex}

\begin{ex}%[2D4V2-4]%[Anh Duy]
Tích phân  $\displaystyle\int\limits_{0}^{1}\dfrac{\mathrm{e}^{2x}}{\sqrt{\mathrm{e}^{x}}}\mathrm{\,d}x=\dfrac{a}{b}\left(\mathrm{e}^{\tfrac{3}{2}}-1\right)$. Tính $b-a$.
\shortans{$1$}
\loigiai{ Ta có \[\displaystyle\int\limits_{0}^{1}\dfrac{\mathrm{e}^{2x}}{\sqrt{\mathrm{e}^{x}}}\mathrm{\,d}x=\displaystyle\int\limits_{0}^{1}\dfrac{\mathrm{e}^{2x}}{\mathrm{e}^{\tfrac{x}{2}}}\mathrm{\,d}x=\displaystyle\int\limits_{0}^{1}\mathrm{e}^{\tfrac{3x}{2}}\mathrm{\,d}x=\dfrac{2}{3}\cdot \mathrm{e}^{\tfrac{3x}{2}}\bigg|_{0}^{1}=\dfrac{2}{3}\left(\mathrm{e}^{\tfrac{3}{2}}-1\right).\]
Do đó $a=2$, $b=3$ nên $b-a=1$.}
\end{ex}

\begin{ex}%[2D4V2-4]%[Anh Duy]
Tích phân $\displaystyle\int\limits_{0}^{\ln 4}\dfrac{\mathrm{e}^{2x}}{\sqrt{\mathrm{e}^{2x}+9}}\mathrm{\,d}x=a-\sqrt{b}$. Tính $a-b$.
\shortans{$-5$}
\loigiai{
Ta có \[\displaystyle\int\limits_{0}^{\ln 4}\dfrac{\mathrm{e}^{2x}}{\sqrt{\mathrm{e}^{2x}+9}}\mathrm{\,d}x=\dfrac{1}{2} \displaystyle\int\limits_{0}^{\ln 4}\dfrac{\mathrm{d}\left(\mathrm{e}^{2x}+9\right)}{\sqrt{\mathrm{e}^{2x}+9}}=\sqrt{\mathrm{e}^{2x}+9}\bigg|_{0}^{\ln 4}=5-\sqrt{10}.\]
Khi đó $a=5$, $b=10$ nên $a-b=-5$.}
\end{ex}

\begin{ex}%[MĐ3]%[2D4V2-4]
Biết $I=\displaystyle\int_1^2\dfrac{x\bigl(1+\mathrm{e}^x \bigr)+\ln x+1}{\bigl(x\ln x+\mathrm{e}^x \bigr)^2} \mathrm{d}x$ $=\dfrac{a}{b\ln 2+\mathrm{e}^c}+\dfrac{2}{\mathrm{e}}$ với $a$, $b$, $c$ là các số nguyên. Tổng $P=a+b+c$ là
\shortans{$1$}
\loigiai{
Ta có
\begin{eqnarray*}
I &=&\displaystyle\int_1^2\dfrac{x\bigl(1+\mathrm{e}^x \bigr)+\ln x+1}{\bigl(x\ln x+\mathrm{e}^x \bigr)^2} \mathrm{d}x=\displaystyle\int_1^2\dfrac{x\bigl(\ln x+1+\mathrm{e}^x \bigr)-\bigl(x\ln x+\mathrm{e}^x \bigr)+\bigl(\ln x+1+\mathrm{e}^x \bigr)}{\bigl(x\ln x+\mathrm{e}^x \bigr)^2} \mathrm{d}x\\
&=& -\displaystyle\int_1^2\dfrac{(x+1)'\bigl(x\ln x+\mathrm{e}^x \bigr)-(x+1)\bigl(x\ln x+\mathrm{e}^x \bigr)'}{\bigl(x\ln x+\mathrm{e}^x \bigr)^2} \mathrm{d}x \\
&=&-\displaystyle\int_1^2\left(\dfrac{x+1}{x\ln x+\mathrm{e}^x} \right)' \mathrm{d}x=-\dfrac{x+1}{x\ln x+\mathrm{e}^x} \bigg|_1^2=\dfrac{-3}{2\ln 2+\mathrm{e}^2}+\dfrac{2}{\mathrm{e}}\\
&\Rightarrow& a=-3;\,b=2;\,c=2.
\end{eqnarray*}
Vậy $P=a+b+c=1$.
}
\end{ex}

\begin{ex}%[MĐ3]%[2D4V2-4]
Biết $\displaystyle\int_1^4\sqrt{\dfrac{1}{4x}+\dfrac{\sqrt{x}+\mathrm{e}^x}{\sqrt{x} e^{2x}}} \mathrm{d}x=a+\mathrm{e}^b-\mathrm{e}^c$ với $a$, $b$, $c$ là các số nguyên. Tổng $T=a+b+c$ là
\shortans{$-4$}
\loigiai{
Ta có $\dfrac{1}{4x}+\dfrac{\sqrt{x}+\mathrm{e}^x}{\sqrt{x} \mathrm{e}^{2x}}=\left(\dfrac{1}{2\sqrt{x}}+\dfrac{1}{\mathrm{e}^x} \right)^2$ nên
\[\displaystyle\int_1^4\sqrt{\dfrac{1}{4x}+\dfrac{\sqrt{x}+\mathrm{e}^x}{\sqrt{x} \mathrm{e}^{2x}}} \mathrm{d}x=\displaystyle\int_1^4\left(\dfrac{1}{2\sqrt{x}} +\dfrac{1}{\mathrm{e}^x} \right)\mathrm{d}x=\bigl(\sqrt{x}-e^{-x} \bigr)\bigg|_1^4=1+\mathrm{e}^{-1}-\mathrm{e}^{-4}.
\]
Vậy $a=1$, $b=-1$, $c=-4$. Suy ra $T=-4$.
}
\end{ex}

\begin{ex}%[MĐ3]%[2D4V2-4]
Cho hàm số $f(x)=\heva{&2\mathrm{e}^{2x}+4x+1 &\qquad \text{khi}\; x\ge 0 \\ &3x^2+2x+3 &\qquad \text{khi} x < 0}$. Giả sử $F$ là nguyên hàm của $f$ trên $\mathbb{R}$ thoả mãn $F(-1)=1$. Biết rằng $2F(2)-F(-3)=a\mathrm{e}^4+b$ (trong đó $a, b$ là các số hữu tỉ). Khi đó $a+b$ bằng
\shortans{$51$}
\loigiai{
Vì $F$ là nguyên hàm của $f$ trên $\mathbb{R}$ nên $F(x)=\heva{&\mathrm{e}^{2x}+2x^2+x+C_1 &\qquad \text{khi}\; x\ge 0\\ &x^3+x^2+3x+C_2 &\qquad \text{khi} \;x < 0.}$\\
Ta có $F(-1)=1\Leftrightarrow-3+C_2=1$ $\Leftrightarrow C_2=4$.\\
Nhận xét: Hàm số $f(x)$ xác định và liên tục trên mỗi khoảng $(-\infty; 0)$ và $(0;+\infty)$.\\
$\lim\limits_{x\to 0^{+}} f(x)=\lim\limits_{x\to 0^{-}} f(x)=f(0)=3$ nên hàm số $f(x)$ liên tục tại $x=0$.\\
Suy ra hàm số $f(x)$ liên tục trên $\mathbb{R}$.\\
Do đó hàm số $F(x)$ liên tục trên $\mathbb{R}$ nên hàm số $F(x)$ liên tục tại $x=0$.\\
Suy ra ${\mathop{\lim}\limits_{x\to 0^{+}}} F(x)={\mathop{\lim}\limits_{x\to 0^{-}}} F(x)=F(0)$ $\Leftrightarrow 1+C_1=C_2$, mà $C_2=4$ nên $C_1=3$.\\
Vậy $F(x)=\heva{&\mathrm{e}^{2x}+2x^2+x+3 &\qquad \text{khi} \, x\ge 0 \\ &x^3+x^2+3x+4&\qquad \text{khi}\, x < 0.}$\\
Ta có $2F(2)-F(-3)=\bigl(2\mathrm{e}^4+26\bigr)-(-23)=2\mathrm{e}^4+49$. \\
Suy ra $a=2; b=49$. Vậy $a+b=51$.
}
\end{ex}

\begin{ex}%[2D4V2-4]
Giả sử hàm số $y=f(x)$ liên tục và thỏa mãn: $f(1)=1$ và $f'(x)\sqrt[3]{x^{-1}}=1$, với mọi $x>0$. Tính $4f(8)$.
\shortans{$47$}
\loigiai{
Ta có $f'(x)=\dfrac{1}{\sqrt[3]{x^{-1}}}=\dfrac{1}{x^{-\tfrac{1}{3}}}=x^{\tfrac{1}{3}}$\\
$\Rightarrow F(x)=\displaystyle\int f'(x)\mathrm{d}x= \int x^{\tfrac{1}{3}}\mathrm{d}x=\dfrac{3}{4}x^{\frac{4}{3}}+C=\frac{3}{4}\sqrt[3]{x^4}+C$.\\
$f(1)=1\Rightarrow \dfrac{3}{4}+C=1\Rightarrow C=-\dfrac{1}{4}.\\
\Rightarrow f(x)=\dfrac{3}{4}\sqrt[3]{x^4}-\dfrac{1}{4}$\\
$\Rightarrow 4f(8)=47$.}
\end{ex}

\begin{ex}%[2D4V2-4]
Gọi $F(x)$ là một nguyên hàm của hàm số $f(x)=3^{2x+1} 2^{1+3x}$, biết $F(0)=\dfrac{8}{\ln 72}$. Tính $F(-2)$. (làm tròn kết quả đến hàng phần trăm).
\shortans{$0{,}47$}
\loigiai{
Ta có:\\
$F(x)=\displaystyle \int{\left(3^{2x+1}\cdot 2^{1+3x} \right)}\mathrm{d}x=\int\left(3\cdot3^{2x}\cdot2\cdot2^{3x}\right)\mathrm{d}x=\int\left(6\cdot9^x\cdot8^x \right)\mathrm{d}x\\
=6\int 72^x\mathrm{d}x=6\cdot\dfrac{72^x}{\ln 72}+C$.\\
Theo giả thiết, $F(0)=\dfrac{8}{\ln 72}\Rightarrow 6\cdot\dfrac{72^0}{\ln 72}+C=\dfrac{8}{\ln 72}\Rightarrow C=\dfrac{2}{\ln 72}$\\
$\Rightarrow F(x)=6\cdot\dfrac{{{72}^{x}}}{\ln 72}+\dfrac{2}{\ln 72}\Rightarrow F\left( -2 \right)=6\cdot \dfrac{72^{-2}}{\ln 72}+\dfrac{2}{\ln 72}\approx 0{,}47$.
}
\end{ex}

\begin{ex}%[EX-3-2026, GHK2-12 THPT Nguyễn Quốc TRinh, Hà Nội]%[Huỳnh Xuân Tín]%[2D4V2-4]
Gọi $a$, $b$ là các số nguyên sao cho $\displaystyle\int\limits_0^2\sqrt{\mathrm{e}^{x+2}} \mathrm{~d} x=2a e^2+b e$. Giá trị của $a^2+b^2$ bằng
\shortans[0]{5}
\loigiai{Ta có $\displaystyle\int\limits_0^2\sqrt{\mathrm{e}^{x+2}} \mathrm{~d} x=\displaystyle\int\limits_0^2\mathrm{e}^{\tfrac{x}2+1} \mathrm{~d} x=2\mathrm{e}^{\tfrac{x}2+1}\Bigg|_0^2=2\mathrm{e}^2-2\mathrm{e}=2a e^2+b e$.\\
Khi đó $a=1$, $b=-2$. Suy ta $a^2+b^2=5$.}
\end{ex}

\begin{ex}%[2D4V2-5]%[To 20 - Dot 17 - Chuong 4 - Bai 3 - CD - De 1 - TLN]%[Lê Quân]
Cho $I=\displaystyle\int\limits_5^{12}\dfrac{\mathrm{\,d} x}{x \sqrt{x+4}}=\dfrac{1}{a}\ln \dfrac{b}{c}$ với $a, b, c$ là các số nguyên dương và phân số $\dfrac{b}{c}$ là phân số tối giản. Tính $P=a-b+c$.
\shortans[3]{$0$}
\loigiai{
Đặt $t=\sqrt{x+4}\Rightarrow t^2=x+4 \Rightarrow 2 t \mathrm{\,d}t=\mathrm{d}x$. \\
Đổi cận $x=5 \Rightarrow t=3 ; x=12 \Rightarrow t=4$.\\
Khi đó
\begin{eqnarray*}
I&=&\displaystyle\int\limits_5^{12}\dfrac{\mathrm{\,d}x}{x \sqrt{x+4}}\\
&=&\displaystyle\int\limits_3^4 \dfrac{2 t \mathrm{\,d}t}{t\left(t^2-4\right)}\\
&=&2 \displaystyle\int\limits_3^4 \dfrac{\mathrm{\,d}t}{t^2-4}\\
&=&\dfrac{1}{2}\displaystyle\int\limits_3^4\left(\dfrac{1}{t-2}-\dfrac{1}{t+2}\right) \mathrm{\,d}t\\
&=&\left.\dfrac{1}{2}\ln \left|\dfrac{t-2}{t+2}\right|\right|_3 ^4\\
&=&\dfrac{1}{2}\ln \dfrac{5}{3}.
\end{eqnarray*}
Vậy $a=2 ; b=5 ; c=3 \Rightarrow P=a-b+c=0$.}
\end{ex}

\begin{ex}%[2D4V2-5]%[Tổ 20 - Đợt 17 - Chương 4 - - CD - Đề 2]%[Trần Thiên Đức]
Cho biết $I=\displaystyle\int\limits_1^{\sqrt{5}}\dfrac{x^3}{\sqrt{x^2-1}+1}\mathrm{\,d}x=\dfrac{a}{b}-c\ln d$ với $a$, $b$, $c$, $d\in\mathbb{Z}^{+}$, $\dfrac{a}{b}$ tối giản, $d<9$. Tính giá trị của $T=a-b-c+d$.
\shortans[]{$12$}
\loigiai{
Đặt: $t=\sqrt{x^2-1}$.\\
$\Rightarrow t^2=x^2-1$.\\
$\Rightarrow x^2=t^2+1$.\\
$\Rightarrow t\mathrm{\,d}t=x\mathrm{\,d}x$.\\
Đổi cận
\begin{center}
\renewcommand{\arraystretch}{1.2}
\begin{longtable}{|c|c|c|}
\hline
$x$ & $1$ & $\sqrt{5}$ \\
\hline
$t$ & $0$ & $2$ \\
\hline
\end{longtable}
\end{center}
Ta được\\
$\displaystyle\int\limits_1^{\sqrt{5}}\dfrac{x^3}{\sqrt{x^2-1}+1}\mathrm{\,d}x=\displaystyle\int\limits_0^2\left(\dfrac{t^2+1}{t+1}\right)t\mathrm{\,d}t=\displaystyle\int\limits_0^2\left(t^2-t+2-\dfrac{2}{t+1}\right)\mathrm{\,d}t$.\\
$=\left(\dfrac{t^3}{3}-\dfrac{t^2}{2}+2t-2\ln|t+1|\right)\Bigg|_0^2=\dfrac{14}{3}-2\ln3$.\\
Vậy: $T=14-3-2+3=12$.
}
\end{ex}

\begin{ex}%[2D4V2-5]
Cho hàm số $f(x)$ có đạo hàm liên tục trên $[1; 2]$, có $f\left(2\right)=14$ và $\displaystyle\int\limits_1^2{\dfrac{f(x)}{x}\mathrm{\,d}x=6}$. Khi đó $I=\displaystyle\int\limits_1^2{f'(x)}\ln x\mathrm{\,d}x=a\ln b - c$. Tính $T=a+b+c$.
\shortans[]{$22$}
\loigiai{
Đặt $\heva{ & u=\ln x \\   & v'=f'(x) } \Rightarrow \heva{  & u'=\dfrac{1}{x} \\   & v=f(x) } $. \\
Khi đó $I=f(x)\ln x\bigg|_1^2 - \displaystyle\int\limits_1^2{\dfrac{f(x)}{x}\mathrm{\,d}x=f\left(2\right)\cdot \ln 2 - 0 - 6=14\ln 2 - 6}$.\\
Do đó $a=14$, $b=2$, $c=6$.\\
Suy ra $T=a+b+c=22$.
}
\end{ex}

\begin{ex}%[DA-EX-TF-TLN2024 - Trần Xuân Hòa]%[2D4V2-5]
Biết $\displaystyle \int \limits_1^5 \dfrac{1}{x\sqrt{3x+1}} \mathrm{\,d}x =a\ln 3+b\ln 5$ với $a$, $b\in \mathbb{Z}$. Giá trị của $S=a^2+ab+3b^2$ bằng
\shortans{$5$}
\loigiai{
Đặt $t=\sqrt[3]{3x+1} \Rightarrow t^2=3x+1 \Rightarrow 2t\mathrm{d}t=3\mathrm{d}x \Rightarrow \dfrac{2t}{3}\mathrm{d}t=\mathrm{d}x$.\\
Đổi cận: $x=1 \Rightarrow t=2$, $x=5 \Rightarrow t=4$.\\
Khi đó
\begin{eqnarray*}
I&=&\displaystyle \int \limits_1^5 \dfrac{1}{x\sqrt{3x+1}} \mathrm{\,d}x
=\displaystyle \int \limits_2^4 \dfrac{3}{(t^2-1)\cdot t}\cdot \dfrac{2}{3}t \mathrm{\,d}t\\
&=& \displaystyle \int \limits_2^4 \left(\dfrac{1}{t-1} -\dfrac{1}{t+1}\right) \mathrm{\,d}t
= \ln \left.\left| \dfrac{t-1}{t+1} \right|\right|_2^4
= 2\ln2 -\ln 5.
\end{eqnarray*}
Suy ra $a=2$ và $b=-1$. Vậy $S=a^2+ab+3b^2=5$.
}
\end{ex}

\begin{ex}%[DA-EX-TF-TLN2024 - Trần Xuân Hòa]%[2D4V2-5]
Cho tích phân $\displaystyle \int \limits_1^2 \sqrt{\dfrac{1}{x^8}+\dfrac{1}{x^6}} \mathrm{d}x=a\sqrt{2}-b\sqrt{5}$ với $a,b$ là các số hữu tỷ. Giá trị của biểu thức $a+b$ bằng (làm tròn đến số thập phân thứ nhất).
\shortans{$0{,}9$}
\loigiai{Ta có $I=\displaystyle \int \limits_1^2 \sqrt{\dfrac{1}{x^8}+\dfrac{1}{x^6}} \mathrm{d}x=\displaystyle \int \limits_1^2\sqrt{\dfrac{1}{x^6}\left(1+\dfrac{1}{x^2} \right)}=\displaystyle \int \limits_1^2 \dfrac{1}{x^3} \sqrt{1+\dfrac{1}{x^2}} \mathrm{d}x$.\\
Đặt $t=\dfrac{1}{x}$. Khi đó\\
+)  với $x=1 \Rightarrow t=1$, với $x=2 \Rightarrow t=\dfrac{1}{2}$.\\
+) $\mathrm{d}t=-\dfrac{\mathrm{d}x}{x^2}$.\\
Do đó  $I=-\displaystyle \int \limits_1^{\frac{1}{2}} t\cdot \sqrt{1+t^2} \mathrm{d}t=\displaystyle \int \limits_{\frac{1}{2}}^1 t \sqrt{t^2+1} \mathrm{d}t$.\\
Đặt $u=t^2+1$. Khi đó \\
+) với $t=\dfrac{1}{3} \Rightarrow u=\dfrac{5}{4}$, với $t=1 \Rightarrow u=2$.\\
+) $\mathrm{d}u=2t\mathrm{d}t$.\\
Do đó  $I=\dfrac{1}{2}\displaystyle \int \limits_{\frac{5}{4}}^2 \sqrt{u} \mathrm{d}u=\dfrac{1}{3}\left(\sqrt{2^3}-\sqrt{(\dfrac{5}{4})^3}\right)=\dfrac{2}{3}\sqrt{2}-\dfrac{5}{24}\sqrt{5}$.\\
Suy ra $a=\dfrac{2}{3}$, $b=\dfrac{5}{24} \Rightarrow a+b=\dfrac{7}{8}\approx 0{,}9$.

}
\end{ex}

\begin{ex}%[2D4V2-5]
Cho hàm số $y=f(x)=\heva{&\mathrm{e}^x+1&\text{ khi }x\ge 0 \\&x^2-2x+2&\text{ khi }x<0}$. Biết giá trị của $\displaystyle I=\int\limits_{1/e}^{\mathrm{e}^2}\dfrac{f(\ln x-1)}{x}\mathrm{\,d}x=\dfrac{a}{b}+ce$ với $a$, $b$, $c\in \mathbb{Z}$, $b\ne 0$ và $\left(a,b\right)=1$ bằng. Giá trị của $a+b+c$.
\shortans{$36$}
\loigiai{
Xét $\displaystyle I=\int\limits_{1/e}^{\mathrm{e}^2}\dfrac{f(\ln x-1)}{x}\mathrm{\,d}x$.\\
Đặt $u=\ln x-1\Rightarrow \mathrm{\,d}u=\dfrac{1}{x}\mathrm{\,d}x$. Đổi cận $\heva{&x=\dfrac{1}{e}\Rightarrow u=-2 \\&x=\mathrm{e}^2\Rightarrow u=1}$.\\
Khi đó $\displaystyle I=\int\limits_{-2}^1f(u)\mathrm{\,d}u=\int\limits_{-2}^1f(x)\mathrm{\,d}x=\int\limits_{-2}^0f(x)\mathrm{\,d}x+\int\limits_0^1f(x)\mathrm{\,d}x$
$\displaystyle =\int\limits_{-2}^0(x^2-2x+2)\mathrm{\,d}x+\int\limits_0^1(\mathrm{e}^x+1)\mathrm{\,d}x$ $=\left. \left(\dfrac{1}{3}x^3-x^2+2x\right)\right|_{-2}^0+\left. (\mathrm{e}^x+x)\right|_0^1=\dfrac{32}{3}+e$.\\
Vậy $\heva{&a=32 \\&b=3 \\&c=1}\Rightarrow a+b+c=36$.}
\end{ex}

\begin{ex}%[2D4V2-5]
Cho $\displaystyle\int\limits_{0}^{\frac{1}{2}}\sqrt{1-x^2}\mathrm{\,d}x=a\pi+b\sqrt{3}\,(*)$. Biết rằng tồn tại duy nhất bộ hai số hữu tỉ $a$, $b$ thỏa mãn $(*)$. Tính tổng $6a+4b$.
\shortans{1}
\loigiai{
Đặt $x=\sin t$ với $t\in\left[-\dfrac{\pi}{2} ; \dfrac{\pi}{2}\right]$, suy ra $\mathrm{d}x=\cos t\mathrm{\,d}t$.\\
Với $x=0$ thì $t=0$, với $x=\dfrac{1}{2}$ thì $t=\dfrac{\pi}{6}$.\\
Khi đó
$
\displaystyle\int\limits_{0}^{\frac{1}{2}}\sqrt{1-x^2}\mathrm{\,d}
= \displaystyle\int\limits_{0}^{\frac{\pi}{6}}\sqrt{1-\sin^2t}\cdot \cos t\mathrm{\,d}t
= \displaystyle\int\limits_{0}^{\frac{\pi}{6}}\cos^2t\mathrm{\,d}t
= \dfrac{1}{2}\displaystyle\int\limits_{0}^{\frac{\pi}{6}}(1+\cos 2t)\mathrm{\,d}t=\dfrac{\pi}{12}+\dfrac{\sqrt{3}}{8}$.\\
Suy ra $a=\dfrac{1}{12}$, $b=\dfrac{1}{8}$ và $6a+4b=1$
}
\end{ex}

\begin{ex}%[Diamond, dự án HSA đợt 2]%[2D4V2-5]
Cho hàm số $f(x) = x^3\cdot \mathrm{e}^x$. Giá trị (làm tròn đến hàng đơn vị) của $\displaystyle\int_2^3 \left(f'(x) + \dfrac{x-5}{x^2}\right) \mathrm{\,d}x$ là
\shortans{$483$}
\loigiai{
Ta có $f'(x)=3x^2\mathrm{e}^x+x^3\mathrm{e}^x=(x^3+3x^2)\mathrm{e}^x$.\\
Suy ra
\begin{eqnarray*}
\displaystyle\int_2^3 \left(f'(x) + \dfrac{x-5}{x^2}\right) \mathrm{\,d}x&=&\displaystyle\int_2^3 \left((x^3+3x^2)\mathrm{e}^x + \dfrac{1}{x}-\dfrac{5}{x^2}\right) \mathrm{\,d}x\\
&=&\left((x^3+3x^2-3x^2-6x+6x+6-6)\mathrm{e}^x+\ln|x|+\dfrac{5}{x}\right)\bigg|_2^3\\
&=&\left(x^3\mathrm{e}^x+\ln|x|+\dfrac{5}{x}\right)\bigg|_2^3\\
&=&3^3\mathrm{e}^3-2^3\mathrm{e}^2+\ln3-\ln2+\dfrac{5}{3}-\dfrac{5}{2}\\
&\approx&483.
\end{eqnarray*}
}
\end{ex}

\begin{ex}%[2D4V2-5]
Cho hàm số $y=f(x)$ liên tục, đồng biến và nhận giá trị dương trên khoảng $(0;+\infty)$ đồng thời thỏa mãn $f(3)=\dfrac{4}{9}$, $\left[f'(x)\right]^2=(x+1) \cdot f(x)$. Tính $f(8)$.
\shortans{49}
\loigiai{
Ta có $\left[f'(x)\right]^2=(x+1) \cdot f(x)$ nên $\dfrac{f'(x)}{\sqrt{f(x)}}=\sqrt{x+1}$. Khi đó
\[\begin{aligned}[t]
& \displaystyle\int\limits_3^8 \dfrac{f'(x)}{\sqrt{f(x)}}\mathrm{\,d}x=\displaystyle\int\limits_3^8 \sqrt{x+1}\mathrm{\,d}x\\
& \Rightarrow \left(2\sqrt{f(x)}\right)\bigg|_3^8=\dfrac{38}{3} \Leftrightarrow \sqrt{f(8)}-\sqrt{f(3)}=\dfrac{19}{3} \Leftrightarrow \sqrt{f(8)}=7 \Leftrightarrow f(8)=49.
\end{aligned}\]
}
\end{ex}

\begin{ex}%[2D4V2-5]
Cho hàm số $f(x)$ liên tục trên $[0; +\infty)$ và thỏa mãn $f\left(x^2+3x+1\right)=6x+1,\,\, \forall x\ge 0$. Tính $\displaystyle\int\limits_{1}^{5}f(x)\mathrm{\, d}x$.
\shortans{$17$}
\loigiai{
Ta có
$f\left(x^2+3x+1\right)=6x+1,\,\,\forall x\ge 0\Rightarrow (2x+3)f\left(x^2+3x+1\right)=(6x+1)(2x+3),\,\,\forall x\ge 0$\\
$\Rightarrow\displaystyle\int\limits_{0}^{1}(2x+3)f\left(x^2+3x+1\right)\mathrm{\, d}x=\displaystyle\int\limits_{0}^{1}(6x+1)(2x+3)\mathrm{\, d}x. \quad(1)$.\\
Đặt $t=x^2+3x+1$ suy ra $\mathrm{d}t=(2x+3)\mathrm{d}x$, khi đó từ $(1)$ ta được $\displaystyle\int\limits_{1}^{5}f(t)\mathrm{\, d}t=17$.\\
Vậy $\displaystyle\int\limits_{1}^{5}f(x)\mathrm{\, d}x=17$.
}
\end{ex}

\begin{ex}%[2D4V2-6]%[To 20 - Dot 17 - Chuong 4 - Bai 3 - CD - De 1 - TLN]%[Lê Quân]
\immini{ Cho đồ thị biểu diễn vận tốc của hai xe $A$ và $B$ khởi hành cùng một lúc và cùng vạch xuất phát, đi cùng chiều trên một con đường. Biết đồ thị biểu diễn vận tốc của xe $A$ là một đường parabol và đồ thị biểu diễn vận tốc của xe $B$ là một đường thẳng như hình vẽ bên. Hỏi sau $5$ giây kể từ lúc xuất phát thì khoảng cách giữa hai xe là bao nhiêu mét? (Làm tròn đến hàng phần chục và biết rằng xe $A$ sẽ dừng lại khi vận tốc bằng $0$).}
{\begin{tikzpicture}[line join=round, line cap=round,>=stealth,thick]
\tikzset{every node/.style={scale=0.9}}
\draw[->] (-1.1,0)--(5.1,0) node[below left] {$x$};
\draw[->] (0,-1.1)--(0,5.1) node[below left] {$y$};
\coordinate (B) at (3,3);
\coordinate (O) at (0,0);
\coordinate (M) at ($(O)!1.5!(B)$);
\draw (O) node [below left] {$O$};
\draw (O)--(M) node[below right ] {$v_B$};
\foreach \x/\nx in {3/3,4/4}
\draw[thin] (\x,1pt)--(\x,-1pt) node [below] {$\nx$};
\foreach \y/\ny in {3/60}
\draw[thin] (1pt,\y)--(-1pt,\y) node [left] {$\ny$};
\draw[dashed,thin](3,0)--(3,3)--(0,3);
\begin{scope}
\clip (-1,-1) rectangle (5,5);
\draw[samples=200,domain=0:4,smooth,variable=\x] plot (\x,{-1*(\x)^2+4*(\x)+0});
\end{scope}
\end{tikzpicture}}
\shortans[3]{$36{,}7$}
\loigiai{
Đồ thị biểu diễn vận tốc của xe $A$ là Parabol $(P)$ có phương trình $v_A=a t^2+b t+c(a \neq 0)$.
Vì $(P)$  đi qua các điểm có tọa độ
$(0 ; 0)$; $(3 ; 60)$; $(4 ; 0)$ nên \[ v_A=-20 t^2+80 t.\]
Đồ thị biểu diễn vận tốc của xe $B$ là đường thẳng $\Delta$ có phương trình  $v_B=m t+n(m \neq 0)$.\\
Vì $\Delta$  đi qua các điểm có tọa độ
$(0 ; 0)$; $(3 ; 60)$ nên  $ v_B=20 t.$\\
Ta có $v_A=-20 t^2+80 t=0 \Rightarrow t=4$ nên xe $A$ dừng lại sau giây thứ $4$.\\
Do đó quãng đường xe $A$ đi được sau $4$ giây là \[S_A=\displaystyle\int\limits_0^4\left(-20 t^2+80 t\right) \mathrm{\,d} t=\dfrac{640}{3}(\mathrm{m}).\]
Quãng đường xe $B$ đi được sau $5$ giây đầu là \[S_B=\displaystyle\int\limits_0^5 20 t \mathrm{\,d} t=250\,\,(\mathrm{m}).\]
Khoảng cách giữa hai xe sau $5$ giây kể từ lúc xuất phát là \[\Delta S=\left|S_A-S_B\right|=\dfrac{110}{3}=36{,}7\,\,(\mathrm{m}).\]
}
\end{ex}

\begin{ex}%[2D4V2-6]%[Tổ 20 - Đợt 17 - Chương 4 - - CD - Đề 2]%[Trần Thiên Đức]
\immini{
Một vật chuyển động trong $3$ giờ với vận tốc $v$ (km/h) phụ thuộc vào thời gian $t$ (h) có đồ thị vận tốc như hình bên. Tính quãng đường mà vật di chuyển được trong $3$ giờ (làm tròn đến hàng phần chục).
\shortans[]{$10{,}7$}
}
{
\begin{tikzpicture}[scale=1, font=\footnotesize, line join=round, line cap=round,>=stealth]%<DTools>
%Gán số liệu.
\def\xmin{-.5};\def\ymin{-.5};\def\xmax{5};\def\ymax{6};
%Gán tọa độ.
\coordinate (O) at (0,0);
%Trục Oxy.
\draw[->] (\xmin,0)--(\xmax,0) node[below]{$t$ (h)};
\draw[->] (0,\ymin)--(0,\ymax) node[right]{$v(t)$ (km/h)};
\fill (O) node[below left]{$O$} circle(1pt);
%Giới hạn đồ thị.
\clip ({\xmin-0.1},{\ymin-0.1}) rectangle ({\xmax+0.1},{\ymax+0.1});
\foreach \x in {1,2,3,4}{
\fill (\x,0) node[below]{$\x$} circle(1pt);
}
\foreach \y in {1,2,3,4,5}{
\fill (0,\y) node[left]{$\y$} circle(1pt);
}
%Vẽ đồ thị.
\draw[thick,samples=100] plot[domain=-0:4.24](\x,{-(\x)^2+4*\x+1});
\draw[dashed] (0,5)-|(2,0);
\end{tikzpicture}
}
\loigiai{
Phương trình chuyển động của vật theo đường parabol $v(t)=at^2+bt+c$ $(\mathrm{km/h})$.\\
Parabol có đỉnh $I(2;5)$ và đi qua điểm $(0;1)$ có phương trình $v(t)=at^2+bt+c$.\\
Ta có $\heva{
&c=1\\
&4a+2b+c=5\\
&\dfrac{-b}{2a}=2\\}\Leftrightarrow\heva{
&a=-1\\
&b=4\\
&c=1\\}\Rightarrow v(t)=-t^2+4t+1$.\\
Quãng đường vật đi được trong $1$ giờ đầu là\\
$S_1=\displaystyle\int\limits_0^1\left(-t^2+4t+1\right)\mathrm{\,d}t=\dfrac{8}{3}$.\\
Trong ba giờ vật đi được quãng đường là $S=S_1+S_2=\dfrac{8}{3}+8=\dfrac{32}{3}\approx10{,}7$ $(\mathrm{km})$.
}
\end{ex}

\begin{ex}%[2D4V2-6]%[Tổ 20 - Đợt 17 - Chương 4 - - CD - Đề 2]%[Trần Thiên Đức]
Một ô tô chuyển với vận tốc $v(t)=2t$ $(\mathrm{m/s})$ $(t=0)$ là lúc bắt đầu chuyển động) sau $20$ giây thì gặp chướng ngại vật nên đạp phanh và chuyển động với gia tốc $a(t)=-5$ $\left(\mathrm{m/s}^2\right)$. Tính quãng đường xe ô tô đi được từ khi bắt đầu đi đến khi dừng hẳn.
\shortans[]{$560$}
\loigiai{
Vận tốc của xe sau $20$ giây là $v=40$ $(\mathrm{m/s})$.\\
Gọi $v_1(t)$ là vận tốc của xe sau khi đạp phanh ($t=0$ là thời điểm đạp phanh).\\
Ta có $v_1(t)=\displaystyle\int a(t)\mathrm{\,d}t=-5t+C$.\\
$t=0$ thì $v_1=40\Leftrightarrow C=40$. Suy ra $v_1(t)=-5t+40$.\\
Khi xe dừng hẳn $v_1(t)=0\Leftrightarrow t=8$.\\
Vậy quãng đường xe ô tô đi được từ khi bắt đầu đi đến khi dừng hẳn là\\
$S=\displaystyle\int\limits_0^{20}v(t)\mathrm{\,d}t+\displaystyle\int\limits_0^8v_1(t)\mathrm{\,d}t=\displaystyle\int\limits_0^{20}(2t)\mathrm{\,d}t+\displaystyle\int\limits_0^8(-5t+40)\mathrm{\,d}t=560$ $(\mathrm{m})$.
}
\end{ex}

\begin{ex}%[2D4V2-6]%[Tổ 20 - Đợt 17 - Chương 4 - - CD - Đề 2]%[Trần Thiên Đức]
Một ôtô đang chuyển động đều với vận tốc $v_0$ (m/s) thì bắt đầu tăng tốc với gia tốc $a(t)=v_0t+t^2$ (m/s$^2$), trong đó $t$ là khoảng thời gian tính bằng giây kể từ thời điểm vật bắt đầu tăng tốc. Biết quãng đường ôtô đi được trong $3$ giây kể từ lúc bắt đầu tăng tốc là $100$ m. Tính vận tốc $v_0$ của ôtô? (Làm tròn đến hàng phần chục)
\shortans[]{$12{,}4$}
\loigiai{
Ta có $v(t)=\displaystyle\int a(t)\mathrm{\,d}t=\displaystyle\int\left(v_0t+t^2\right)\mathrm{\,d}t=\dfrac{t^3}{3}+v_0\dfrac{t^2}{2}+C$.\\
Tại $t=0\Rightarrow v(0)=v_0=C$.\\
Vì quãng đường ôtô đi được trong $3$ giây kể từ lúc bắt đầu tăng tốc là $100$ m nên
$100=\displaystyle\int\limits_0^3\left(\dfrac{t^3}{3}+v_0\dfrac{t^2}{2}+v_0\right)\mathrm{\,d}t=\dfrac{3^4}{12}+\dfrac{3^3}{6}v_0+3v_0$.\\
Suy ra $v_0=12{,}4$ (km/h).
}
\end{ex}

\begin{ex}%[2D4V2-6]%[Tổ 20 - Đợt 17 - Chương 4 - - CD - Đề 2]%[Trần Thiên Đức]
Một xe chuyển động với vận tốc thay đổi là $v(t)=3at^2+bt$. Gọi $S(t)$ là quãng đường đi được sau $t$ giây. Biết rằng sau $5$ giây thì quãng đường đi được là $150$ m, sau $10$ giây thì quãng đường đi được là $1\,100$ m. Tính quãng đường xe đi được sau $20$ giây.
\shortans[]{$8\,400$}
\loigiai{
Quãng đường đi được sau $5$ giây là\\
$S_1=\displaystyle\int\limits_0^5v(t)\mathrm{\,d}t=\displaystyle\int\limits_0^5\left(3at^2+bt\right)\mathrm{\,d}t=125a+\dfrac{25}{2}b$.\\
Quãng đường đi được sau $10$ giây là\\
$S_2=\displaystyle\int\limits_0^{10}v(t)\mathrm{\,d}t=\displaystyle\int\limits_0^{10}\left(3at^2+bt\right)\mathrm{\,d}t=1\,000a+50b$.\\
Theo đề bài, ta có\\
$\heva{
&125a+\dfrac{25}{2}b=150\\
&1\,000a+50b=1\,100\\}\Leftrightarrow\heva{
&a=1\\
&b=2.\\}$\\
Suy ra $v(t)=3t^2+2t$.\\
Quãng đường xe đi được sau $20$ giây là\\
$S=\displaystyle\int\limits_0^{20}v(t)\mathrm{\,d}t=\displaystyle\int\limits_0^{20}\left(3t^2+2t\right)\mathrm{\,d}t=8\,000+400=8\,400$ (m).
}
\end{ex}

\begin{ex}%[Tổ 20 - Đợt 17 - Chương 4 - - CD - Đề 8]%[Huỳnh Xuân Tín]%[2D4V2-6]
Một ô tô đang chạy thì người lái đạp phanh. Từ thời điểm đó, ô tô chuyển động chậm dần đều với vận tốc $v(t)=-12 t+36$ (m/s) trong đó $t$ là khoảng thời gian tính bằng giây, kể từ lúc bắt đầu đạp phanh. Từ lúc đạp phanh đến khi dừng hẳn, ô tô di chuyển được quãng đường là $s$ mét. Giá trị của $s$ là
\shortans{$54$}
\loigiai
{Thời gian từ lúc xe đạp phanh đến lúc dừng hẳn là $v(t)=0 \Leftrightarrow-12 t+36=0 \Leftrightarrow t=3(s)$.\\
Từ lúc đạp phanh đến khi dừng hẳn, ô tô di chuyển được quãng đường là \[s=\displaystyle\int\limits_0^3(-12 t+36) \mathrm{\,d}t=54 (\text{m}).\]
}
\end{ex}

\begin{ex}%[Tổ 20 - Đợt 17 - Chương 4 - - CD - Đề 8]%[Huỳnh Xuân Tín]%[2D4V2-6]
Để đảm bảo an toàn khi lưu thông trên đường, các xe ô tô khi dừng đèn đỏ phải cách nhau tối thiểu $1$ (m). Một ô tô $A$ đang chạy với vận tốc $15$ m/s bỗng gặp ô tô $B$ đang đứng chờ đèn đỏ nên ô tô $A$ hãm phanh và chuyển động chậm dần đều bởi vận tốc được biểu thị bởi công thức $v_A(t)=15-5 t$ (m/s). Để hai ô tô $A$ và $B$ đạt khoảng cách an toàn khi dừng lại thì ô tô $A$ phải hãm phanh khi cách ô tô $B$ một khoảng ít nhất là $s$ mét. Giá trị của biểu thức $s$ là
\shortans{$23{,}5$}
\loigiai
{Khi ô tô dừng lại $v_A(t)=0 \Leftrightarrow t=3$.\\
Quãng đường ô tô $A$ đi được từ lúc ô tô $A$ đạp phanh đến khi dừng hẳn là
\[\displaystyle\int\limits_0^3(15-5 t)\displaystyle\int\limits t=22,5 (\mathrm{~m}).\]
Vậy để đảm bảo an toàn thì ô tô $A$ phải hãm phanh khi cách ô tô $B$ một khoảng ít nhất là $23{,}5$ m.\\
Do đó $s=23{,}5$.
}
\end{ex}

\begin{ex}%[2D4V2-6]
Tại một nơi không có gió, một chiếc khí cầu đang đứng yên ở độ cao 162 (m) so với mặt đất đã được phi công cài đặt cho nó chế độ chuyển động đi xuống. Biết rằng, khí cầu đã chuyển động theo phương thẳng đứng với vận tốc tuân theo quy luật $v(t)=10 t-t^2$, trong đó $t$ (phút) là thời gian tính từ lúc bắt đầu chuyển động, $v(t)$ được tính theo đơn vị mét/phút (m/p). Nếu như vậy thì khi bắt đầu tiếp đất vận tốc $v$ của khí cầu bằng bao nhiêu m/p?
\shortans{$9$}
\loigiai{
Gọi thời điểm khí cầu bắt đầu chuyển động là $t=0$, thời điểm khinh khí cầu bắt đầu tiếp đất là $t_1$.\\
Quãng đường khí cầu đi được từ thời điểm $t=0$ đến thời điểm khinh khí cầu bắt đầu tiếp đất $t=t_1$ là
\[
\begin{aligned}
& \displaystyle\int\limits_0^{t_1}\left(10 t-t^2\right) \mathrm{\,d} t=5 t_1^2-\dfrac{t_1^3}{3}=162 \\
& \Leftrightarrow t \approx-4{,}93 \vee t \approx 10{,}93 \vee t=9.
\end{aligned}
\]
Do vận tốc chuyển động $v(t) \geq 0 \Leftrightarrow 0 \leq t \leq 10$ nên chọn $t_1=9$.\\
Vậy khi bắt đầu tiếp đất vận tốc $v$ của khí cầu là $v(9)=10\cdot9-9^2=9$ m/p.
}
\end{ex}

\begin{ex}%[2D4V2-6]
Gọi $h(t)$ cm là mức nước trong bồn chứa sau khi bơm được $t$ giây. Biết rằng tốc độ tăng giảm của mực nước là $h^{\prime}(t)=\dfrac{1}{5} \sqrt[3]{t+8}$ (cm/s) và lúc đầu bồn không có nước. Tìm mức nước ở bồn (đơn vị: cm) sau khi bơm nước được $6$ giây (làm tròn đến chữ số hàng phần trăm).
\shortans{$2{,}66$}
\loigiai{Hàm $h(t)=\displaystyle\int\limits \dfrac{1}{5} \sqrt[3]{t+8} \mathrm{\,d} t=\dfrac{3}{20}(t+8) \sqrt[3]{t+8}+C$.\\
Lúc $t=0$, bồn không chứa nước. Suy ra $h(0)=0 \Rightarrow \dfrac{12}{5}+C=0 \Leftrightarrow C=-\dfrac{12}{5}$.\\
Vậy, hàm $h(t)=\dfrac{3}{20}(t+8) \sqrt[3]{t+8}-\dfrac{12}{5}$.\\
Mức nước trong bồn sau $6$ giây là $h(6) \simeq 2{,}66$ cm.
}
\end{ex}

\begin{ex}%[2D4V2-6]
Một chất điểm $A$ xuất phát từ $O$, chuyển động thẳng với vận tốc biến thiên theo thời gian bởi quy luật $v \left(t\right) = \dfrac{1}{100}t^2 + \dfrac{13}{30}t$ (m/s), trong đó $t$ (giây) là khoảng thời gian tính từ lúc $A$ bắt đầu chuyển động. Từ trạng thái nghỉ, một chất điểm $B$ cũng xuất phát từ $O$, chuyển động thẳng cùng hướng với $A$ nhưng chậm hơn $10$ giây so với $A$ và có gia tốc bằng $a$ (m/s$^2$ ) ( $a$ là hằng số). Sau khi $B$ xuất phát được $15$ giây thì đuổi kịp $A$. Vận tốc của $B$ tại thời điểm đuổi kịp $A$ bằng bao nhiêu m/s?
\shortans{$25$}
\loigiai{
Ta có $v_{B}(t) = \displaystyle\int a \cdot \mathrm{\,d}t = at + C$, $v_{B} (0) = 0 \Rightarrow C = 0 \Rightarrow v_{B} \left(t\right) = at$.\\
Quãng đường chất điểm $A$ đi được trong $25$ giây là
\[S_{A} = \displaystyle\int\limits_0^{25} \left(\dfrac{1}{100}t^2 + \dfrac{13}{30}t \right) \mathrm{\,d}t = \left(\dfrac{1}{300}t^3 + \dfrac{13}{60}t^2 \right) \Big|_0^{25} = \dfrac{375}{2}.\]
Quãng đường chất điểm $B$ đi được trong $15$ giây là
\[S_{B} = \displaystyle\int\limits_0^{15} at \cdot \mathrm{\,d}t = \dfrac{at^2}{2} \Big|_0^{15} = \dfrac{225a}{2}.\]
Ta có $\dfrac{375}{2} = \dfrac{225a}{2} \Leftrightarrow a = \dfrac{5}{3}$.\\
Vận tốc của $B$ tại thời điểm đuổi kịp $A$ là $v_{B} \left(15\right) = \dfrac{5}{3} \cdot 15 = 25$ (m/s).
}
\end{ex}

\begin{ex}Cau22D%[2D4V2-6]
Một vật chuyển động trong $4$ giờ với vận tốc $v$ (km/h) phụ thuộc thời gian $t$ (h) có đồ thị của vận tốc như hình bên. Trong khoảng thời gian $3$ giờ kể từ khi bắt đầu chuyển động, đồ thị đó là một phần của đường parabol có đỉnh $I \left(2;9\right)$ với trục đối xứng song song với trục tung, khoảng thời gian còn lại đồ thị là một đoạn thẳng song song với trục hoành. Tính quãng đường $s$ mà vật di chuyển được trong $4$ giờ đó (đơn vị tính bằng km).
\begin{center}
\begin{tikzpicture}[>=stealth,scale=0.45]
% Vẽ 2 trục, điền các số lên trục
\draw[->] (-0.5,0)--(0,0) node[below left]{$O$}--(5,0) node[above]{$t$};
\foreach \x in {2,3,4}
\draw[shift={(\x,0)},color=black] (0pt,2pt)--(0pt,-2pt)
node[below] { $\x$};
\draw[->,color=black] (0,-0.5)--(0,10) node[right]{$v$};
\foreach \y in {9}
\draw[shift={(0,\y)},color=black] (2pt,0pt) -- (-2pt,0pt)
node[left] {$\y$};
\clip(-1,-1) rectangle (5,10); %vùng đồ thị
%\draw[gray!50,thin,opacity=.5] (-1,-1) grid (4,10); %ô vuông
%Vẽ đồ thị
\draw[smooth,samples=100,domain=0:3, font=\footnotesize, line join=round, line cap=round, thick, smooth]
plot(\x,{(-9/4)*(\x)^2+9*(\x)});
\draw[smooth,samples=100, font=\footnotesize, line join=round, line cap=round, thick, smooth,domain=3:4]
plot(\x,{27/4});
% Vẽ thêm mấy cái râu ria
\draw[dashed] (3,0)--(3,27/4) circle(1.5pt);  \draw[dashed] (2,0)--(2,9) circle(1.5pt) node[above]{$I$}--(0,9) circle(1.5pt); \draw[dashed] (4,0)--(4,27/4) circle(1.5pt);
%Vẽ dấu chấm tròn
\fill (0cm,0cm) circle (1.5pt);
\end{tikzpicture}
\end{center}
\shortans{$27$}
\loigiai{
Gọi $\left(P\right) \colon y = ax^2+bx+c$.\\
Vì $\left(P\right)$ qua $O\left(0;0\right)$ và có đỉnh $I\left(2;9\right)$ nên dễ tìm được phương trình là $y = \dfrac{-9}{4}x^2 + 9x$.\\
Ngoài ra tại $x=3$ ta có $y = \dfrac{27}{4}$.\\
Vậy quãng đường cần tìm là: $S = \displaystyle\int\limits_0^3 \left(\dfrac{-9}{4}x^2 +9x \right) \mathrm{\,d}x + \displaystyle\int\limits_3^4 \dfrac{27}{4} \mathrm{\,d}x = 27$ (km).
}
\end{ex}

\begin{ex}Cau23D%[2D4V2-6]
Một vật chuyển động trong $6$ giờ với vận tốc $v$ (km/h) phụ thuộc vào thời gian $t$ (h) có đồ thị như hình bên dưới. Trong khoảng thời gian $2$ giờ từ khi bắt đầu chuyển động, đồ thị là một phần đường Parabol có đỉnh $I\left(3;9\right)$ và có trục đối xứng song song với trục tung. Khoảng thời gian còn lại, đồ thị vận tốc là một đường thẳng có hệ số góc bằng $\dfrac{1}{4}$. Tính quảng đường $s$ mà vật di chuyển được trong $6$ giờ? (đơn vị tính bằng km, làm tròn đến chữ số thập phân thứ nhất).
\begin{center}
\begin{tikzpicture}[>=stealth,scale=0.5]
% Vẽ 2 trục, điền các số lên trục
\draw[->] (-0.5,0)--(0,0) node[below left]{$O$}--(7,0) node[above]{$t$}; %định dạng trục Ox
\foreach \x in {2,3,6}
\draw[shift={(\x,0)},color=black] (0pt,2pt)--(0pt,-2pt)
node[below] { $\x$};
\draw[->,color=black] (0,-0.5)--(0,10) node[right]{$v$};  %định dạng trục Oy
\foreach \y in {8,9}
\draw[shift={(0,\y)},color=black] (2pt,0pt) -- (-2pt,0pt)
node[left] {$\y$};
\clip(-1,-1) rectangle (7,10); %vùng đồ thị
%\draw[gray!50,thin,opacity=.5] (-1,-1) grid (4,10); %ô vuông
%Vẽ đồ thị
\draw[smooth,samples=100,domain=0:2,font=\footnotesize, line join=round, line cap=round, thick]
plot(\x,{(-1)*(\x)^2+6*(\x)});
\draw[smooth,domain=2:6, line join=round, line cap=round,dashed]
plot(\x,{(-1)*(\x)^2+6*(\x)});
\draw[smooth,samples=100,domain=2:6,font=\footnotesize, line join=round, line cap=round, thick]
plot(\x,{(1/4)*(\x)+15/2});
% Vẽ thêm mấy cái râu ria
\draw[dashed] (3,0)--(3,9) circle(1.5pt) node[above]{$I$}--(0,9) circle(1.5pt);
\draw[dashed] (2,0)--(2,8) circle(1.5pt) --(0,8);
\draw[dashed] (6,0)--(6,9) circle(1.5pt) --(0,9);
%Vẽ dấu chấm tròn
\fill (0cm,0cm) circle (1.5pt);
\end{tikzpicture}
\end{center}
\shortans{$43{,}3$}
\loigiai{
Vì Parabol đi qua $O\left(0;0\right)$ và có tọa độ đỉnh $I\left(3;9\right)$ nên thiết lập được phương trình Parabol là $\left(P\right) \colon y = v\left(t\right) = -t^2+6t$; $\forall t \in \left[0;2\right]$.\\
Sau $2$ giờ đầu thì hàm vận tốc có dạng là hàm bậc nhất $y = \dfrac{1}{4}t + m$, dựa trên đồ thị ta thấy đi qua điểm có tọa độ $\left(6;9\right)$ nên thế vào hàm số và tìm được $m = \dfrac{15}{2}$.\\
Nên hàm vận tốc từ giờ thứ $2$ đến giờ thứ $6$ là: $y = \dfrac{1}{4}t + \dfrac{15}{2},\forall t \in \left[2;6\right]$.\\
Quảng đường vật đi được bằng tổng đoạn đường $2$ giờ đầu và đoạn đường $4$ giờ sau.
\[S = S_1 +S_2 = \displaystyle\int\limits_0^2 \left(-t^2+6t\right) \mathrm{\,d}t + \displaystyle\int\limits_2^6 \left(\dfrac{1}{4}t+\dfrac{15}{2} \right) \mathrm{\,d}t = \dfrac{130}{3} \approx 43{,}3 \left(\text{km}\right).\]
}
\end{ex}

\begin{ex}Cau26D%[2D4V2-6]
Chất điểm chuyển động theo quy luật vận tốc $v\left(t\right)$ (m/s) có dạng đường Parapol khi $0 \leq t\leq 5$ (s) và $v\left(t\right)$ có dạng đường thẳng khi $5 \leq t \leq 10$ (s). Cho đỉnh Parapol là $I\left(2;3\right)$. Hỏi quãng đường đi được chất điểm trong thời gian $0 \leq t \leq 10$ (s) là bao nhiêu mét? (làm tròn đến hàng đơn vị)
\begin{center}
\begin{tikzpicture}[>=stealth,scale=0.3]
% Vẽ 2 trục, điền các số lên trục
\draw[->] (-0.5,0)--(0,0) node[below left]{$O$}--(13,0) node[above]{$t$};
\foreach \x in {2,5,10}
\draw[shift={(\x,0)},color=black] (0pt,2pt)--(0pt,-2pt)
node[below] { $\x$};
\draw[->,color=black] (0,-0.5)--(0,22) node[right]{$v$};
\foreach \y in {3,11}
\draw[shift={(0,\y)},color=black] (2pt,0pt) -- (-2pt,0pt)
node[left] {$\y$};
\clip(-1,-1) rectangle (11,23); %vùng đồ thị
\draw[smooth,samples=100,domain=0:5, font=\footnotesize, line join=round, line cap=round, thick]
plot(\x,{2*(\x)^2-8*(\x)+11});
\draw[smooth,samples=100,domain=5:10, font=\footnotesize, line join=round, line cap=round, thick]
plot(\x,{(-21/5)*(\x)+42});
% Vẽ thêm mấy cái râu ria
\draw[dashed] (5,0)--(5,21) circle(1.5pt);  \draw[dashed] (2,0)--(2,3) circle(1.5pt)--(0,3);
%Vẽ dấu chấm tròn
\fill (0cm,0cm) circle (1.5pt);
\end{tikzpicture}
\end{center}
\shortans{$91$}
\loigiai{
Gọi Parapol $\left(P\right) \colon y = ax^2+bx+c$ khi $0 \leq t \leq 5 \left(s\right)$.\\
Do $\left(P\right) \colon y = ax^2+bx+c$ đi qua $I\left(3;2\right)$; $A \left(0;11\right)$ nên
\[\heva{&4a+2b+c=3\\&c=11\\&4a+b=0} \Rightarrow \heva{&a=2\\&b=-8\\&c=11.}\]
Khi đó quãng đường vật di chuyển trong khoảng thời gian từ $0 \leq t \leq 5$ (s) là
\[S_1 = \displaystyle\int\limits_0^5 \left(2x^2 - 8x +11\right) \mathrm{\,d}x = \dfrac{115}{3} \left(\text{m}\right).\]
Ta có $f\left(5\right) = 21$.\\
Gọi $d \colon y = ax+b$ khi $5 \leq t \leq 10$ (s), do $d$ đi qua điểm $B\left(5;21\right)$ và $C\left(10;0\right)$ nên
\[\heva{&5a+b=11\\&10a+b=0} \Rightarrow \heva{&a=-\dfrac{21}{5}\\&b=42.}\]
Khi đó quãng đường vật di chuyển trong khoảng thời gian từ $5 \leq t \leq 10$ (s) là:
\[S_2 = \displaystyle\int\limits_5^{10} \left(-\dfrac{21}{5}x+42\right) \mathrm{\,d}x = \dfrac{105}{2} \left(\text{m}\right).\]
Quãng đường đi được chất điểm trong thời gian  $0 \leq t \leq 10$ (s) là:
\[S = \dfrac{115}{3}+\dfrac{105}{2}= \dfrac{545}{6} \approx 91 \left(\text{m}\right).\]
}
\end{ex}

\begin{ex}%[2D4V2-6]Biên soạn:Phùng Hoàng Em Phản biện:Nguyễn Đắc Kiên
Giá trị trung bình của hàm số liên tục $f(x)$ trên đoạn $[a ; b]$ được định nghĩa là
\[
\dfrac{1}{b-a} \int_a^b f(x) \mathrm{\,d}x .
\]
Giả sử nhiệt độ ngoài trời ở một thành phố vào các thời điểm khác nhau trong ngày được mô phỏng bởi công thức $	h(t)=29+3 \sin \left[\dfrac{\pi}{12}(t-9)\right]$, với $h$ tính bằng độ $\mathrm{C}$ và $t$ là thời gian trong ngày tính bằng giờ.	Tìm nhiệt độ trung bình (đơn vị: độ) của ngày trong khoảng thời gian từ $6$ giờ sáng đến $12$ giờ trưa.
\shortans{$29$}
\loigiai{
Nhiệt độ trung bình của ngày trong khoảng thời gian $6\le t \le 12$ được tính bởi công thức sau
\[\dfrac{1}{12-6} \int_6^{12} h(t) \mathrm{\,d}t=\dfrac{1}{6} \int_6^{12} \left( 29+3 \sin \left[\dfrac{\pi}{12}(t-9)\right]\right) \mathrm{\,d}t=29.\]
}
\end{ex}

\begin{ex}%[2D4V2-6]
Một vật chuyển động với gia tốc được cho bởi hàm số $a(t)=5\cos t$ (m/s$^2$). Lúc bắt đầu chuyển động vật có vận tốc $2{,}5$ m/s. Tính gia tốc của vật (đơn vị m/s$^2$) tại thời điểm vận tốc đạt giá trị lớn nhất trong $\pi$ (s) đầu tiên.
\shortans{$0$}
\loigiai{
Vận tốc của vật được biểu diễn bởi hàm số $v(t)=\displaystyle\int a(t) \mathrm{\,d}t=\displaystyle\int 5\cos t \mathrm{\,d}t=5\sin t+C$.\\
Khi bắt đầu chuyển động, vật có vận tốc $2{,}5$ m/s nên ta có
\[
v(0)=2{,}5 \Leftrightarrow 5 \sin 0+C=2{,}5 \Leftrightarrow C=2{,}5
\]
Suy ra $v(t)=5\sin t+2{,}5$. \\
Ta có $\sin t \le 1 \Rightarrow 5\sin t+2{,}5\leq 7{,}5$. \\
Vận tốc đạt giá trị lớn nhất khi $\sin t=1 \Rightarrow t=\dfrac{\pi}{2}$. \\
Khi đó, gia tốc của vật tại thời điểm $t=\dfrac{\pi}{2}$ là $a\left(\dfrac{\pi}{2}\right)=5\cos \left(\dfrac{\pi}{2}\right)=0$ m/s$^2$}
\end{ex}

\begin{ex}%[2D4V2-6]
Một vật chuyển động với gia tốc được cho bởi hàm số $a(t)=5 \cos t\left(\mathrm{~m} / \mathrm{s}^2\right)$. Lúc bắt đầu chuyển động vật có vận tốc $2,5 \mathrm{~m} / \mathrm{s}$. Tính gia tốc của vật tại thời điểm vận tốc đạt giá trị lớn nhất trong $\pi(s)$ đầu tiên.
\shortans{$0$}
\loigiai{
Vận tốc của vật được biểu diễn bởi hàm số $v(t)=\displaystyle\int a(t) \mathrm{d} t=\displaystyle\int 5 \cos t \mathrm{~d} t=5 \sin t+C$.
Khi bắt đầu chuyển động, vật có vận tốc $2,5 \mathrm{~m} / \mathrm{s}$ nên ta có:
\[
v(0)=2,5 \Leftrightarrow 5 \sin 0+C=2,5 \Leftrightarrow C=2,5 .
\]

Suy ra $v(t)=5 \sin t+2,5$. Mà $5 \sin t+2,5 \leq 7,5$. Vậy vận tốc đạt giá trị lớn nhất $7,5$ tại $t=\dfrac{\pi}{2}$. Khi đó, gia tốc của vật tại thời điểm $t=\dfrac{\pi}{2}$ là $a\left(\dfrac{\pi}{2}\right)=5 \cos \left(\dfrac{\pi}{2}\right)=0\left(\mathrm{~m} / \mathrm{s}^2\right)$.
}
\end{ex}

\begin{ex}%[2D4V2-6]
Một ô tô đang chạy với vận tốc $10$(m/s) thì người lái xe đạp phanh. Từ thời điểm đó, ô tô chuyển động chậm dần đều với vận tốc $v(t)=10-2t$ (m/s), trong đó $t$ là khoảng thời gian tính bằng giây kể từ lúc đạp phanh. Tính quãng đường ô tô di chuyển được trong $8$ giây cuối cùng.
\shortans{$55$}
\loigiai{
Chọn mốc thời gian và gốc tọa độ lúc ô tô bắt đầu đạp phanh. Suy ra $t=0;\,s=0$.\\
$s(t)=\displaystyle \int v(t)\mathrm{d}t=\int (10-2t)\mathrm{d}t=10t-t^2+C$.\\
$s(0)=0\Rightarrow C=0 \Rightarrow s(t)=10t-t^2$.\\
Ô tô dừng hẳn khi $v(t)=0\Leftrightarrow 10-2t=0\Leftrightarrow t=5$.\\
Trong $8$ giây cuối:
\begin{itemize}
\item ô tô chuyển động đều với vận tốc $10$(m/s) trong $3$ giây đầu.
\item ô tô chuyển động chậm dần đều trong $5$ giây cuối.
\end{itemize}
Quãng đường ô tô di chuyển là: $s=3\cdot 10+10\cdot 5-5^2=55$ m.}

\end{ex}

\begin{ex}%[2D4V2-6]
Một viên đạn được bắn thẳng đứng lên từ độ cao $1{,}5$ mét so với mặt đất. Giả sử tại thời điểm $t$ giây (coi $t=0$ là thời điểm viên đạn được bắn lên), vận tốc của nó được cho bởi $v(t)=170-9{,}8\,t\,\left( \text{m/s} \right)$. Tìm độ cao lớn nhất của viên đạn (làm tròn kết quả đến hàng đơn vị).
\shortans{$1476$}
\loigiai{
Gọi $h(t)$ là độ cao của viên đạn tại thời điểm $t$ giây sau khi bắn. Ta có:\\
$h(t)=\displaystyle \int v(t)\mathrm{d}t=\int{(170-9{,}8t)}\mathrm{d}t=170t-4{,}9t^2+C$.\\
Từ giả thiết suy ra: $h\left( 0 \right)=1,5\Rightarrow C=1{,}5\Rightarrow h(t)=170t-4{,}9t^2+1,5$.\\
Viên đạn đạt độ cao lớn nhất khi $v(t)=0\Leftrightarrow 170-9,8\,t\,=0\Leftrightarrow t=\dfrac{850}{49}$.\\
Khi đó, độ cao lớn nhất của viên đạn là:\\
$h\left(\dfrac{850}{49}\right)=170 \cdot\dfrac{850}{49}-4{,}9\left( \dfrac{850}{49} \right)^2+1{,}5=\dfrac{144647}{98}\approx 1476$ (m).}
\end{ex}

\begin{ex}%[2D4V2-6]
Một chiếc cốc chứa nước ở $95^\circ$ C được đặt trong phòng có nhiệt độ ${{20}^{0}}C$. Theo định luật làm mát của Newton, nhiệt độ của nước trong cốc sau $t$ phút (xem $t=0$ là thời điểm nước ở $95^\circ$ C là một hàm số $(t)$. Tốc độ giảm nhiệt độ của nước trong cốc tại thời điểm t phút được xác định bởi $T'(t)=\left(-\dfrac{3}{2} \mathrm\mathrm{e}^{-\tfrac{t}{50}}\right)^\circ$ C/phút). Tính nhiệt độ của nước tại thời điểm $t=40$ phút (làm tròn kết quả đến hàng phần chục).
\shortans{$53{,}7$}
\loigiai{
Ta có:\\
$\displaystyle T(t)=\int T'(t)\mathrm{d}t
=\int\left( -\frac{3}{2}\mathrm{e}^{-\tfrac{t}{50}} \right)\mathrm{d}t
=-\frac{3}{2}\int\left({\mathrm{e}^{-\tfrac{1}{50}}} \right)^t\mathrm{d}t\\
=-\frac{3}{2}\cdot\frac{\left(\mathrm{e}^{-\tfrac{1}{50}} \right)^t}{\ln \left(\mathrm{e}^{-\tfrac{1}{50}}\right)}+C
=75\left(\mathrm{e}^{-\frac{1}{50}}\right)^t+C$.\\
Vì $t=0$ là thời điểm nước ở $95^\circ$ C nên $T(0)=95\Rightarrow 75\left(\mathrm{e}^{-\tfrac{1}{50}} \right)^\circ+C=95\Rightarrow C=20$.\\
Suy ra $T(t)=75\left(\mathrm{e}^{-\frac{1}{50}} \right)^t+20$.\\
Do đó, nhiệt độ của nước tại thời điểm $t=40$ phút là: \\
$T(40)=75\left(\mathrm{e}^{-\tfrac{1}{50}} \right)^{40}+20\approx 53{,}7 ^\circ$ C.}
\end{ex}

\begin{ex}%[2D4V2-6]
Doanh thu bán hàng của một công ty khi bán một loại sản phẩm là số tiền $R(x)$ (triệu đồng) thu được khi $x$ đơn vị sản phẩm được bán ra. Tốc độ biến động (thay đổi) của doanh thu khi $x$ đơn vị sản phẩm đã được bán là hàm số $M_R(x)=R'(x)$. Một công ty công nghệ cho biết, tốc độ biến đổi của doanh thu khi bán một loại con chip của hãng được cho bởi $M_R(x)=40-0{,}1x$, ở đó $x$ là số lượng chip đã bán. Hỏi doanh thu của công ty khi đã bán 500 con chip bằng bao nhiêu tỉ đồng?
\shortans{$7{,}5$}
\loigiai{
Vì $R'(x)=M_R(x)$ nên doanh thu $R(x)$ là một nguyên hàm của $M_R(x)$.\\
Ta có: $R(x)=\displaystyle \int M_R(x) \mathrm{d}x=\int{(40-0{,}1x)}\mathrm{d}x=40x-0{,}05 x^2+C$.\\
Khi $x=0$, tức là chưa bán chip nào thì doanh thu sẽ bằng $0$ (triệu đồng), nên $R\left( 0 \right)=0\Rightarrow C=0$.\\
Suy ra $R(x)=40x-0{,}05 x^2$.\\
Do đó, doanh thu của công ty khi đã bán 500 con chip là:\\
$R(500)=40\cdot 500-0{,}05\cdot 500^2=7500$ (triệu đồng) $=7{,}5$ (tỉ đồng).
}
\end{ex}

\begin{ex}%[2D4V2-6]
Một ô tô chuyển động nhanh dần đều với vận tốc $v(t)=7 t$ m/s. Đi được $5$ s người lái xe phát hiện chướng ngại vật và phanh gấp, ô tô tiếp tục chuyển động chậm dần đều với gia tốc $a=-35 \mathrm{~m} / \mathrm{s}^2$. Tính quãng đường của ô tô đi được từ lúc bắt đầu chuyển bánh cho đến khi dừng hẳn? (quãng đường tính theo đơn vị m).
\shortans{$105$}
\loigiai{
Quãng đường ô tô đi được trong $5$ s đầu là $s_1=\displaystyle\int\limits_0^5 7 t \mathrm{\,d} t=7 \dfrac{t^2}{2}\,\bigg|_0 ^5=87{,}5$.\\
Phương trình vận tốc của ô tô khi người lái xe phát hiện chướng ngại vật là $v_2(t)=35-35 t$.\\
Khi xe dừng lại hẳn thì $v_2(t)=0 \Leftrightarrow 35-35 t=0 \Leftrightarrow t=1$.\\
Quãng đường ô tô đi được từ khi phanh gấp đến khi dừng lại hẳn là
\[
s_2=\displaystyle\int\limits_0^1\left(35-35 t\right) \mathrm{d} t=\left(35 t-35 \dfrac{t^2}{2}\right)\bigg|_0 ^1=17{,}5.\]
Do đó quãng đường của ô tô đi được từ lúc bắt đầu chuyển bánh cho đến khi dừng hẳn là \[s=s_1+s_2=87{,}5+17{,}5=105.\]
}
\end{ex}

\begin{ex}%[Cau-2]%[2D4V2-6]
Tốc độ chuyển động của thang máy từ tầng $1$ lên tầng cao nhất theo thời gian $t$ (giây) được cho bởi công thức
\[v(t) = \heva{&t&\text{khi}&\ 0 \le t \le 2\\&2 &\text{khi}&\ 2 < t \le 20\\&12-0{,}5t &\text{khi}&\ 20 < t \le 24.}\]
Tính vận tốc trung bình của thang máy.
\shortans{$1{,}75$}
\loigiai{
Quãng đường chuyển động của thang máy là\\
$s=\displaystyle \int\limits_{0}^{24} v(t)\mathrm{\,d}t = \displaystyle\int\limits_{0}^{2} v(t)\mathrm{\,d}t + \displaystyle\int\limits_{2}^{20} v(t)\mathrm{\,d}t + \displaystyle\int\limits_{20}^{24} v(t)\mathrm{\,d}t = \displaystyle\int\limits_{0}^{2} t\mathrm{\,d}t +\displaystyle\int\limits_{2}^{20} 2 \mathrm{\,d}t+\displaystyle\int\limits_{20}^{24} (12-0{,}5t)\mathrm{d}t = 42$.\\
Tốc độ trung bình của thang máy là $v_{tb}=\dfrac{s}{t}=\dfrac{42}{24}=1{,}75$ (m/s).
}
\end{ex}

\begin{ex}%[Cau-5]%[2D4V2-6]
Ba Tí muốn làm cửa sắt được thiết kế như hình bên dưới. Vòm cổng có hình dạng là một Parabol. Giá $1$ m$^2$ cửa sắt là $660\,000$ đồng. Cửa sắt có giá (nghìn đồng) là bao nhiêu?
\begin{center}
\begin{tikzpicture}[smooth,samples=300,scale=1,>=stealth, font=\footnotesize]
%\draw[->] (-3.0,0)--(4.0,0) node[below]{$x$};
%\draw[->] (0,-1.8)--(0,1.5) node[right]{$y$};
\path
(-2.5,0) coordinate (A)
(2.5,0) coordinate (B)
(0,0.5) coordinate (C)
(2.5,-1.5) coordinate (E)
(-2.5,-1.5) coordinate (D)
(0,0) coordinate (O)
;
\draw[thick, magenta,domain=-2.5:2.5] plot(\x,{-0.08*(\x)^2+1/2});
%\draw[thick, magenta,domain=-2.5:2.5] plot(\x,{-0.04*(\x)^2+2-0.2});
\draw[dashed] (D)--($(D)+(0,-0.5)$)--($(E)+(0,-0.5)$)--(E);
\draw[dashed] (C)--(3,0.5)--(3,-1.5)--(0,-1.5);
%\draw[dashed] (A)--(B);
\draw (A)--(D)--(E)--(B) (0,-1.5)--(0,0.5);
\node[rotate=90] at (3.2,-0.8) {$2$ m};
\node[rotate=90] at (2.7,-0.8) {$1{,}5$ m};
\node[rotate=0] at (0.3,-2.2) {$5$ m};
\end{tikzpicture}
\end{center}
\shortans{$6050$}
\loigiai{
\begin{center}
\begin{tikzpicture}[smooth,samples=300,scale=1,>=stealth, font=\footnotesize]
\draw[->] (-3.0,0)--(4.0,0) node[below]{$x$};
\draw[->] (0,-1.8)--(0,1.5) node[right]{$y$};
\path
(-2.5,0) coordinate (A)
(2.5,0) coordinate (B)
(0,0.5) coordinate (C)
(2.5,-1.5) coordinate (E)
(-2.5,-1.5) coordinate (D)
(0,0) coordinate (O)
;
\draw[thick, magenta,domain=-2.5:2.5] plot(\x,{-0.08*(\x)^2+1/2});
%\draw[thick, magenta,domain=-2.5:2.5] plot(\x,{-0.04*(\x)^2+2-0.2});
\draw[dashed] (D)--($(D)+(0,-0.5)$)--($(E)+(0,-0.5)$)--(E);
\draw[dashed] (C)--(3,0.5)--(3,-1.5)--(0,-1.5);
\draw[dashed] (A)--(B);
\draw (A)--(D)--(E)--(B);
\node[rotate=90] at (3.2,-0.8) {$2$ m};
\node[rotate=90] at (2.7,-0.8) {$1{,}5$ m};
\node[rotate=0] at (0.3,-2.2) {$5$ m};
\draw[fill=black] (O) circle(1pt) node[below right]{$O$};
\draw[fill=black] (A) circle(1pt) node[below left]{$A$};
\draw[fill=black] (B) circle(1pt) ($(B)+(60:0.3)$) node {$B$};
\draw[fill=black] (C) circle(1pt) node[above right]{$C$};
\draw[fill=black] (D) circle(1pt) node[left]{$D$};
\draw[fill=black] (E) circle(1pt) node[below right]{$E$};
\draw[->] (-3.5,0.25) node[left]{Phần 1}--(-1,0.25);
\draw[->] (-3.5,-1) node[left]{Phần 2}--(-2,-1);
\end{tikzpicture}
\end{center}
Từ hình vẽ ta chia cửa rào sắt ra thành $2$ phần như trên.\\
Khi đó $S = S_1 + S_2 = S_1 + 5\cdot 1{,}5 = S_1 + 7{,}5$.\\
Để tính $S_1$ ta vận dụng kiến thức tính diện tích hình phẳng của tích phân.\\
Gắn hệ trục $Oxy$ trong đó $O$ trung với trung điểm của $AB$, $OB\subset Ox$, $OC \subset Oy$.\\
Theo đề bài ta có đường cong có dạng hình Parabol. Giả sử $(P)\colon y=ax^2+bx+c$.\\
Khi đó $\heva{&A\left(-\dfrac{5}{2}; 0\right)\in (P)\\&B\left(\dfrac{5}{2}; 0\right)\in (P)\\&C\left(0;\dfrac{1}{2}\right)\in (P)} \Leftrightarrow\heva{&\dfrac{25}{4}a-\dfrac{5}{2}b+c=0\\&\dfrac{25}{4}a+ \dfrac{5}{2}b+c=0\\&c=\dfrac{1}{2}} \Leftrightarrow \heva{&a= -\dfrac{2}{25}\\ &b = 0\\&c=\dfrac{1}{2}.}$\\
$\Rightarrow (P) \colon y = -\dfrac{2}{25} x^2 + \dfrac{1}{2}$.\\
Diện tích $S_2 = 2\displaystyle\int\limits_{0}^{2{,}5} \left(-\dfrac{2}{25} x^2 + \dfrac{1}{2}\right) \mathrm{d}x = \dfrac{10}{6}\ (\mathrm{m}^2)$.\\
$\Rightarrow S = \dfrac{55}{6}$ (m$^2$).\\
Vậy  giá tiền cửa sắt là $\dfrac{55}{6} \cdot 660\,000= 6050$ (nghìn đồng).
}
\end{ex}

\begin{ex}%[Cau-6]%[2D4V2-6]
Một ô tô đang chạy đều với vận tốc $15$ (m/s) thì phía trước xuất hiện chướng ngại vật nên người lái đạp phanh gấp. Kể từ thời điểm đó, ô tô chuyển động chậm dần đều với gia tốc $-a$ (m/s$^2$). Tìm giá trị của $a$ biết ô tô chuyển động thêm được $20$ (m) thì dừng hẳn. (Kết quả làm tròn đến hàng phần trăm).
\shortans{$5{,}63$}
\loigiai{
Gọi $x(t)$ là hàm biểu diễn quãng đường, $v(t)$ là hàm vận tốc.\\
Ta có $\displaystyle \int\limits_{0}^{t} (-a) \mathrm{\,d}x = -at \Rightarrow v(t) = -at + 15$.\\
Mặt khác $x(t)-x(0)=\displaystyle \int\limits_{0}^{t}v(t)\mathrm{\,d}x =\displaystyle \int\limits_{0}^{t}(-at+15)\mathrm{d}x=-\dfrac{1}{2}at^2 + 15t$.\\
$\Rightarrow x(t)=-\dfrac{1}{2}at^2 + 15t$.\\
Ta có $\heva{&v(t)=0\\&x(t)=20}\Leftrightarrow\heva{&-at+15=0\\&-\dfrac{1}{2}at^2 + 15t = 20}\Rightarrow -\dfrac{15}{2}t+15t=20 \Rightarrow t=\dfrac{8}{3}$.\\
$\Rightarrow a=\dfrac{45}{8}\approx 5{,}63$.
}
\end{ex}

\begin{ex}%[2D4V2-6]
Một ô tô đang chạy với vận tốc $20$ (m/s) thì người lái xe đạp phanh. Sau khi đạp phanh, ô tô chuyển động chậm dần đều với vận tốc $v(t)=-4t+20$ (m/s), trong đó $t$ là khoảng thời gian tính bằng giây kể từ lúc bắt đầu đạp phanh. Hỏi từ lúc đạp phanh đến khi dừng hẳn, ô tô còn di chuyển được bao nhiêu mét?

\shortans{50}
\loigiai{
Đặt $t_0=0$ là thời điểm người lái xe ô tô bắt đầu đạp phanh.\\
Khi ô tô dừng hẳn thì vận tốc triệt tiêu nên $-4t+20=0\Leftrightarrow t=5$.\\
Từ lúc đạp phanh đến khi dừng hẳn, ô tô còn di chuyển được quãng đường là
\[\displaystyle\int\limits_0^5\left(-4t+20\right)\mathrm{d}t=50 \,\text{mét}.\]}
\end{ex}

\begin{ex}%[2D4V2-6]
Bác Năm làm một cái cửa nhà hình parabol có chiều cao từ mặt đất đến đỉnh là $2{,}25$ mét, chiều rộng tiếp giáp với mặt đất là $3$ mét. Nếu chi phí làm mỗi mét vuông cửa là $1{,}5$ triệu đồng thì số tiền bác Năm phải trả là bao nhiêu? (đơn vị triệu đồng)

\shortans{6,75}
\loigiai{
Gọi phương trình parabol $(P)\colon y=a{x^2}+bx+c$. Do tính đối xứng của parabol nên ta có thể chọn hệ trục tọa độ $Oxy$ sao cho $(P)$ có đỉnh $I\in Oy$ (như hình vẽ).\begin{center}
\begin{tikzpicture}[scale=1,>=stealth, font=\footnotesize, line join=round, line cap=round]
\def\a{-1} \def\b{0} \def\c{9/4} % Hệ số
\def\xmin{-2} \def\xmax{2.5}
\def\ymin{-1} \def\ymax{3}
\draw[->] (\xmin,0)--(\xmax,0) node [below]{$x$};
\draw[->] (0,\ymin)--(0,\ymax) node [left]{$y$};
\node at (0,0) [below left]{$O$};
\clip (\xmin+0.1,\ymin+0.1) rectangle (\xmax-0.5,\ymax-0.1);
\draw[smooth,samples=300,domain=-1.5:1.5] plot(\x,{\a*(\x)^2+\b*(\x)+\c});
\fill[pattern=north east lines,opacity=0.8] plot[domain=-1.5:1.5](\x,{\a*(\x)^2+\b*(\x)+\c})--cycle;
\draw[fill=black](0,2.35)node[right,scale=0.8]{$I\left(0,\tfrac94\right)$}
(0,2.25)circle(1pt)
(-1.5,0)circle(1pt)
(1.5,0)circle(1pt)
(-1.35,0)node[below,scale=0.8]{$A\left(-\tfrac32,0\right)$}
(1.3,0)node[below,scale=0.8]{$B\left(\tfrac32,0\right)$}
;
\end{tikzpicture}
\end{center}
Ta có hệ phương trình: $\heva{
&\dfrac{9}{4}=c&&\left(I\in(P)\right)\\
&\dfrac{9}{4}a-\dfrac{3}{2}b+c=0&&\left(A\in(P)\right)\\
&\dfrac{9}{4}a+\dfrac{3}{2}b+c=0&&\left(B\in(P)\right)} \Leftrightarrow\heva{
& c=\dfrac{9}{4}\\
& a=-1\\
& b=0.}$\\
Vậy $(P)\colon y=-x^2+\dfrac{9}{4}$.\\
Dựa vào đồ thị, diện tích cửa parabol là
\[ S=\displaystyle\int\limits_{\tfrac{-3}{2}}^{\tfrac{3}{2}}{\left(-x^2+\dfrac{9}{4}\right)\mathrm{d}x}=2\displaystyle\int\limits_0^{\tfrac{3}{2}}{\left(-x^2+\dfrac{9}{4}\right)\mathrm{d}x}= 2\left(\dfrac{-x^3}{3}+\dfrac{9}{4}x\right)\bigg|_0^{\tfrac{9}{4}}=\dfrac{9}{2}\,{\text{m}^2}.\]
Số tiền phải trả là $\dfrac{9}{2}\cdot 1{,}5=6{,}75$ (triệu đồng).}
\end{ex}

\begin{ex}%[2D4V2-6]
\immini{Một vật chuyển động trong $3$ giờ với vận tốc $v$ (km/h) là một hàm số phụ thuộc vào thời gian $t$ (h) và có đồ thị là một phần của đường parabol như hình vẽ.
Tính quãng đường $S$ mà vật di chuyển được trong $3$ giờ đó.
}
{
\begin{tikzpicture}[scale=0.8,yscale=0.7,xscale=0.7,>=stealth, font=\footnotesize, line join=round, line cap=round]
\def\a{1} \def\b{-2} \def\c{2} % Hệ số
\def\xmin{-1} \def\xmax{4}
\def\ymin{-1} \def\ymax{6}
\draw[->] (\xmin,0)--(\xmax,0) node [below]{$t$};
\draw[->] (0,\ymin)--(0,\ymax) node [left]{$v$};
\node at (0,0) [below left]{$O$};
\clip (\xmin+0.1,\ymin+0.1) rectangle (\xmax-0.5,\ymax-0.1);
\draw[smooth,samples=300,domain=0:3] plot(\x,{\a*(\x)^2+\b*(\x)+\c});
\draw[dashed](1,0)|-(0,1)
(3,0)|-(0,5);
\draw[fill=black]
(0,2)node[left]{$2$}circle(1pt)
(0,1)node[left]{$1$}circle(1pt)
(1,0)node[below]{$1$}circle(1pt)
(3,0)node[below]{$3$}circle(1pt)
(0,5)node[left]{$5$}circle(1pt)
;
\end{tikzpicture}
}
\shortans{6}
\loigiai{
Gọi $v(t)=at^2+bt+c$. Vì $v(t)$ đi qua các điểm có tọa độ $(0;2)$; $(1;1)$; $(3;5)$ nên ta có hệ phương trình:
\[\heva{
& a\cdot0+b\cdot0+c=2\\
& a\cdot1+b\cdot1+c=1\\
& a\cdot9+b\cdot3+c=5}\Leftrightarrow\heva{
& a=1\\
& b=-2\\
& c=2.}\]
Do đó $v(t)=t^2-2t+2$.\\
Vậy quãng đường vật di chuyển trong $3$ giờ là $S=\displaystyle\int\limits_0^3\left(t^2-2t+2\right)\mathrm{d}t=6$ (km).
}
\end{ex}

\begin{ex}%[2D4V2-6]
Một ô tô bắt đầu chuyển động nhanh dần đều với vận tốc $v_1(t)=5t+4$ (m/s). Đi được $6$ (s) người lái xe phát hiện chướng ngại vật và phanh gấp, ô tô tiếp tục chuyển động chậm dần đều với gia tốc $a=-34$ (m/s$^2$). Tính quãng đường $S$ (m) đi được của ô tô từ lúc bắt đầu chuyển bánh cho đến khi dừng hẳn.

\shortans{131}
\loigiai{
Vận tốc của ô tô sau khi phanh gấp là $v_2(t)=\displaystyle\int\left(-34\right)\mathrm{d}t=-34t+C$.\\
Khi người lái xe bắt đầu đạp phanh, ô tô chuyển động với vận tốc \[v_2(0)=v_1(6)=34\ (\text{m/s})\Leftrightarrow 34=C.\]
Vậy $v_2(t)=-34t+34$. \\
Khi ô tô dừng hẳn $v_2(t)=0\Leftrightarrow-34t+34=0\Leftrightarrow t=1$ (s).\\
Quãng đường ô tô chuyển động được trong $6$ giây đầu tiên là \[S_1=\displaystyle\int\limits_0^6v_1(t)\mathrm{d}t=\displaystyle\int\limits_0^6(5t+4)\mathrm{d}t=114\, (\text{m}).\]
Quãng đường ô tô chuyển động được kể từ lúc bắt đầu đạp phanh đến khi dừng hẳn là \[S_2=\displaystyle\int\limits_0^1v_2(t)\mathrm{d}t=\displaystyle\int\limits_0^1\left(-34t+34\right)\mathrm{d}t=17 (\text{m}).\]
Vậy $S=S_1+S_2=114+17=131$ (m).}
\end{ex}

\begin{ex}%[2D4V2-6]
Cho hàm số $f(x)$. Biết $f(0)=4$ và $f'(x)=2\cos^2x+1,\,\forall x \in \mathbb{R}$, khi đó $\displaystyle\int\limits_0^{\tfrac{\pi}{4}} f(x)\mathrm{\,d}x$ bằng bao nhiêu? (làm tròn đến hàng phần trăm)
\par \shortans{4,01}
\loigiai{
Ta có $f(x)=\displaystyle\int f'(x)\mathrm{\,d}x=\displaystyle\int \left(2\cos^2x+1\right)\mathrm{d}x=\displaystyle\int \left(2+\cos 2x\right)\mathrm{d}x=\dfrac{1}{2}\sin 2x+2x+C$. \\
Vì $f(0)=4 \Rightarrow C=4 \Rightarrow f(x)=\dfrac{1}{2}\sin 2x+2x+4$.\\
Vậy $\displaystyle\int\limits_0^{\tfrac{\pi}{4}} f(x)\mathrm{\,d}x=\displaystyle\int\limits_0^{\tfrac{\pi}{4}} \left(\dfrac{1}{2}\sin 2x+2x+4\right)\mathrm{d}x=\left(-\dfrac{1}{4}cos2x+x^2+4x\right)\bigg|_0^{\tfrac{\pi}{4}}=\dfrac{\pi^2+16\pi +4}{16} \approx 4{,}01$.}
\end{ex}

\begin{ex}%[2D4V2-6]
Một ly trà sữa dạng hình nón cụt, có đường kính đáy ly $6$ cm, đường kính miệng ly $9$ cm, chiều cao $13{,}4$ cm, ở miệng ly có sử dụng một nắp đậy có hình dạng nửa mặt cầu và ở đỉnh của nửa mặt cầu này có một hình tròn có đường kính $2$ cm để cắm ống hút, mặt phẳng chứa hình tròn này song song với mặt phẳng chứa miệng ly (tham khảo hình vẽ sau).
\begin{center}
% Tạo bảng với 2 cột canh giữa (c c)
\begin{tabular}{c c}
% --- CỘT 1: HÌNH BÊN TRÁI ---
\begin{tikzpicture}[scale=0.35]
\def\Rtop{4.5}
\def\Rbottom{3}
\def\Rhole{1}
\def\hcone{9.9}
\def\Rcap{4.5}

\draw (-\Rbottom,0) arc (180:360:\Rbottom cm and 0.6cm);
\draw[dashed] (-\Rbottom,0) arc (180:-360:\Rbottom cm and 0.6cm);
\draw (-\Rbottom,0) -- (-\Rtop,\hcone);
\draw (\Rbottom,0) -- (\Rtop,\hcone);
\draw (-\Rtop,\hcone) arc (180:360:\Rtop cm and 0.9cm);
\draw [dashed](-\Rtop,\hcone) arc (180:0:\Rtop cm and 0.9cm);
\draw (-\Rtop,\hcone) arc (180:104:\Rtop);
\draw (\Rtop,\hcone) arc (0:78:\Rtop);
\path (-1,\hcone+\Rcap-0.1) coordinate (top);
\draw ([xshift=-\Rhole]top) arc (180:360:\Rhole cm and 0.25cm);
\path (-1,\hcone+\Rcap-0.1) coordinate (top);
\draw ([xshift=-\Rhole]top) arc (180:-360:\Rhole cm and 0.25cm);
\draw[<->] (\Rtop+0.3,0) -- (\Rtop+0.3,\hcone) node[midway,right] {13,4 cm};
\draw[<->] (-\Rbottom,-0.7) -- (\Rbottom,-0.7) node[midway,below] {6 cm};
\draw[<->] (-\Rtop,\hcone) -- (\Rtop,\hcone) node[midway,above] {9 cm};
\draw[<->] (-\Rhole,\hcone+\Rcap+0.5) -- (\Rhole,\hcone+\Rcap+0.5) node[midway,above] {2 cm};
\end{tikzpicture}
& % Dấu ngăn cách cột trong tabular
% --- CỘT 2: HÌNH BÊN PHẢI ---
\begin{tikzpicture}[scale=0.35, rotate=-90]
% Trục Ox (nằm ngang theo thân ly)
% Trục Oy (vuông góc thân ly)
\draw[->] (5,9.9) -- (-5.5,9.9) node[left] {\(y\)};
\draw[->] (0,-1) -- (0,16) node[below right] {\(x\)};
\draw[dashed] (0,0) node[below]{$A$} -- (-3,0) node[above]{$B$};
\draw[dashed] (0,9.9) node[below left]{$O$} -- (-4.5,9.9) node[above right]{$C$};
\draw[dashed] (-1,14.4) node[right]{$D$} -- (0,14.4);

% Thông số
\def\Rtop{4.5}
\def\Rbottom{3}
\def\Rhole{1}
\def\hcone{9.9}
\def\Rcap{4.5}

% Vẽ đáy dưới
\draw (-\Rbottom,0) arc (180:360:\Rbottom cm and 0.6cm);
\draw[dashed] (-\Rbottom,0) arc (180:-360:\Rbottom cm and 0.6cm);

% Vẽ thân cốc
\draw (-\Rbottom,0) -- (-\Rtop,\hcone);
\draw (\Rbottom,0) -- (\Rtop,\hcone);

% Miệng cốc
\draw (-\Rtop,\hcone) arc (180:360:\Rtop cm and 0.9cm);
\draw [dashed](-\Rtop,\hcone) arc (180:0:\Rtop cm and 0.9cm);

% Nửa cầu (nắp)
\draw (-\Rtop,\hcone) arc (180:104:\Rtop);
\draw (\Rtop,\hcone) arc (0:78:\Rtop);

% Lỗ nhỏ trên nắp
\path (-1,\hcone+\Rcap-0.1) coordinate (top);
\draw ([xshift=-\Rhole]top) arc (180:360:\Rhole cm and 0.25cm);
\path (-1,\hcone+\Rcap-0.1) coordinate (top);
\draw ([xshift=-\Rhole]top) arc (180:-360:\Rhole cm and 0.25cm);
\end{tikzpicture}
\end{tabular}
\end{center}
Chọn hệ trục $Oxy$ (đơn vị trên trục là centimet) với trục $Ox$ đi qua tâm của $2$ đáy hình nón cụt và gốc tọa độ $O$ trùng với tâm của đáy lớn như hình vẽ trên. Tính thể tích bên trong của ly bao gồm cả thể tích của nắp \textit{(Kết quả làm tròn kết quả đến hàng đơn vị)}.
\par
\shortans[oly]{790}
\loigiai{
Thể tích bên trong của ly không bao gồm nắp là $V_1=\pi \displaystyle\int \limits_{-13{,}4}^0 \left(\dfrac{1{,}5x+60{,}3}{13{,}4}\right)^2\mathrm{d}x\approx 600$ (cm$^3$).\\
Thể tích bên trong của ly bao gồm cả thể tích của nắp là $V=V_1+(V_2-V_3)$ trong đó $V_2$ là nửa thể tích của khối cầu và $V_3$ là thể tích của chỏm cầu.\\
Nửa thể tích khối cầu là $V_2=\dfrac{1}{2}\cdot \dfrac{4}{3}\cdot \pi \cdot R^3=\dfrac{243\pi }{4}$ (cm$^3$).\\
Thể tích chỏm cầu
\[V_3=\pi \displaystyle\int \limits_{R-h}^R \left(\sqrt{R^2-x^2}\right)^2\mathrm{d}x
\Leftrightarrow V_3=\pi \displaystyle\int \limits_{R-h}^R (R^2-x^2)\mathrm{d}x
\Leftrightarrow V_3=\pi \left(R^2x-\dfrac{x^3}{3}\right)\bigg|_{R-h}^R.\]
Khi đó $V_3=\pi \left[\left(R^3-\dfrac{R^3}{3}\right)-(R^2\left(R-h)-\dfrac{(R-h)^3}{3}\right)\right]
\Leftrightarrow V_3=\pi h^2\left(R-\dfrac{h}{3}\right)$.\\
Thay số ta có $h=4{,}5-\dfrac{\sqrt{77}}{2}$; $R=4{,}5 \Rightarrow V_3=\dfrac{(729-83\sqrt{77})\pi }{12}$ (cm$^3$).\\
Thể tích bên trong của ly bao gồm cả thể tích của nắp là
\[V=600+\left[\dfrac{243\pi }{4}-\dfrac{(729-83\sqrt{77})\pi }{12}\right]\approx 790\, (\mathrm{ml}).\]
}
\end{ex}

\begin{ex}%[2-GHK2-1-ChuyenViThanh-HauGiang-25]%[12-EX-3-2026, Tran Quoc]%[2D4V2-6]
Một khinh khí cầu bay với độ cao (so với mực nước biển) tại thời điểm $t$ là $h(t)$, trong đó $t$ tính bằng phút, $h(t)$ tính bằng mét. Tốc độ bay của khinh khí cầu được cho bởi hàm số $v(t)=-0{,}12t^2+1{,}2t$, với $t$ tính bằng phút, $v(t)$ tính bằng mét/phút. Tại thời điểm xuất phát ($t=0$), khinh khí cầu ở độ cao $520\,\text{m}$. Tính độ cao tối đa của khinh khí cầu khi bay.
\par\shortans{524}
\loigiai{
Ta có $h'(t) = v(t) = -0{,}12t^2 + 1{,}2t$. \\
Do đó, hàm độ cao $h(t)$ là một nguyên hàm của hàm vận tốc $v(t)$.\\
Ta có
\[\displaystyle \int v(t) \mathrm{\, d}t = \int (-0{,}12t^2 + 1{,}2t) \mathrm{\, d}t = -0{,}04t^3 + 0{,}6t^2 + C.\]
Theo giả thiết, tại thời điểm $t=0$, ta có $h(0) = 520$ nên $C = 520$. \\
Độ cao của khinh khí cầu tại thời điểm $t$ là $h(t) =  -0{,}04t^3 + 0{,}6t^2 + 520$.\\
Phương trình $h'(t) = 0 \Leftrightarrow v(t)=0 \Leftrightarrow  -0{,}12t^2 + 1{,}2t = 0 \Leftrightarrow \hoac{&t=0 \\ &t=10.}$ \\
Bảng biến thiên của hàm số $h(t)$ trên khoảng $(0; +\infty)$ như sau:
\begin{center}
\begin{tikzpicture}
\tkzTabInit[nocadre=false,lgt=1.3,espcl=3,deltacl=0.6]
{$t$/0.7, $h'(t)$/0.7, $h(t)$/2}{$0$, $10$, $+\infty$}
\tkzTabLine{0, +, 0, -, }
\tkzTabVar{-/ 520, +/ 524, -/ $-\infty$}
\end{tikzpicture}
\end{center}
Dựa vào bảng biến thiên, độ cao tối đa đạt được tại $t=10$ là
\[h(10) = -0{,}04 \cdot 10^3 + 0{,}6 \cdot 10^2 + 520 = -40 + 60 + 520 = 524\,\text{(m)}.\]
}
\end{ex}

\begin{ex}%[Đề giữa HK2 Nguyễn Tất Thành - dự án 12EX3]%[Nguyễn Văn Nay]%[2D4V2-6]
Một vật chuyển động trong $6$ giờ với vận tốc $v(km/h)$ phụ thuộc vào thời gian $t(h)$ có đồ thị như hình bên dưới. Trong khoảng thời gian $2$ giờ từ khi bắt đầu chuyển động, đồ thị là một phần đường Parabol có đỉnh $I(3;9)$ và có trục đối xứng song song với trục tung. Khoảng thời gian còn lại, đồ thị vận tốc là một đường thẳng có hệ số góc bằng $\dfrac{1}{3}$. Quãng đường $s$ mà vật đi chuyển được trong $6$ giờ là $\dfrac{a}{b}$, ($a,b\in\mathbb{Z}$), $\dfrac{a}{b}$ là phân số tối giản. Khi đó $a-b$ bằng bao nhiêu?
\begin{center}
\begin{tikzpicture}[scale=.6, font=\footnotesize, line join=round, line cap=round, >=stealth]
\tikzset{label style/.style={font=\footnotesize}}
%Nhập giới hạn đồ thị và hàm số cần vẽ
\def \xmin{-1}
\def \xmax{7}
\def \ymin{-1}
\def \ymax{10}
\def \hamso{sin((2*(\x))*180/pi)}
%\def \tiemcanxien{\x+1}
%Tự động
\draw[->] (\xmin,0)--(\xmax,0) node[below left] {$x$};
\draw[->] (0,\ymin)--(0,\ymax) node[below left] {$y$};
\draw[fill=black] (0,0) circle(1pt) node [below right] {$O$};
\draw[dashed,thin,opacity=.5]
(\xmin,\ymin) grid (\xmax,\ymax);
%Vẽ các điểm trên 2 hệ trục
\foreach \x in {2,3}
\fill[black] (\x,0) circle(1pt) node [below] {$\x$};
\foreach \y in {8,9}
\fill[black] (0,\y) circle(1pt) node [left] {$\y$};
%Vẽ thêm mấy cái râu ria
\draw[line width=1.5](2,8)--(6,9);
%\draw[dashed,thin](2,0)--(2,2)--(-2,2)--(-2,0);
%Tự động
\begin{scope}
\clip (\xmin+0.01,\ymin+0.01) rectangle (\xmax-0.01,\ymax-0.01);
\draw[samples=350,domain=0:2,smooth,variable=\x,line width=1.5] plot (\x,{-(\x)^2+6*(\x)});
\draw[samples=350,domain=2:6,smooth,variable=\x,dashed] plot (\x,{-(\x)^2+6*(\x)});
\end{scope}
\end{tikzpicture}
\end{center}
\shortans[]{125}
\loigiai{
Ta có Parabol có dạng $y=ax^2+bx+c$.\\
Parabol đi qua điểm $0(0;0)$, có đỉnh $I(3;9)$ nên ta được
\[\heva{
&c=0\\
&9a+3b=9\\
& 6a+b=0
}
\Leftrightarrow \heva{
& c=0\\
& a=-1\\
& b=6.
}
\]
Vậy parabol được xác định là $y=-x^2+6x$.\\
Từ đồ thị đã cho ta thấy đường thẳng $d\colon y=\dfrac{1}{3}x+b$ đi qua điểm có tọa độ $(3;8)$ nên ta được $8=\dfrac{1}{3}\cdot 3+b
\Rightarrow b=7$.\\
Vậy $d\colon y=\dfrac{1}{3}x+7$.\\
Quãng đường $s$ mà vật đi chuyển được trong $6$ giờ là $\displaystyle\int\limits_0^2\left(-x^2+6x\right)\mathrm{\,d}x+\displaystyle\int\limits_2^6\left(\dfrac{1}{3}x+7\right)\mathrm{\,d}x
=\dfrac{128}{3}$.\\
Vậy $a-b=128-3=125$.\\
}
\end{ex}

\begin{ex}%[Dự án 12EX3-2526]%[Đề GHK2, Trường THPT Thực Hành SP HCM 25, Võ Hoàng Nghĩa]%[2D4V2-6]
Một chất điểm $A$ xuất phát từ $O$, chuyển động thẳng với vận tốc biến thiên theo thời gian bởi quy luật $v(t)=\dfrac{1}{150}t^2+\dfrac{59}{75}t$ m/s, trong đó $t$ (s) là khoảng thời gian tính từ lúc $A$ bắt đầu chuyển động. Từ trạng thái nghỉ, một chất điểm $B$ cũng xuất phát từ $O$, chuyển động thẳng cùng hướng với $A$ nhưng chậm hơn $5$ giây so với $A$ và có gia tốc bằng $a$ m/s$^2$ ($a$ là hằng số). Sau khi $B$ xuất phát được $10$ giây thì đuổi kịp $A$. Vận tốc của $B$ tại thời điểm đuổi kịp $A$ bằng bao nhiêu m/s?
\shortans[]{$16$}
\loigiai{
Tính từ lúc chất điểm $A$ bắt đầu chuyển động cho đến khi bị chất điểm $B$ bắt kịp thì $A$ đi được $15$ giây, $B$ đi được $12$ giây.\\
Biểu thức vận tốc của chất điểm $B$ có dạng $v_B(t)=\displaystyle\int a\mathrm{\,d}t=at+C$ mà $v_B(0)=0$ nên $v_B(t)=at$.\\
Từ lúc chất điểm $A$ bắt đầu chuyển động cho đến khi bị chất điểm $B$ bắt kịp thì quãng đường hai chất điểm đi được là bằng nhau. Do đó
\[\displaystyle\int\limits_0^{15}\left(\dfrac{1}{150} t^2+\dfrac{59}{75} t\right)\mathrm{\,d}t=\displaystyle\int\limits_0^{12}at\mathrm{\,d}t \Leftrightarrow  96=72a \Leftrightarrow  a=\dfrac{4}{3}.
\]
Do đó vận tốc của $B$ tại thời điểm đuổi kịp $A$ bằng $v_B(12)=\dfrac{4}{3}\cdot12=16$ m/s.
}
\end{ex}

\begin{ex}%[Dự án 12-EX-3-2026, Mui Doan]%[2D4V2-6]
Một vật chuyển động với vận tốc được cho bởi đồ thị ở hình vẽ dưới đây:
\begin{center}
\begin{tikzpicture}[scale=1, font=\footnotesize,line join=round, line cap=round, >=stealth]
\tikzset{declare function={xmin=-1;xmax=3;ymin=-1;ymax=3;
f(\x)=2*(\x);g(\x)=2;
},
smooth,samples=450
}
\draw[->] (xmin,0)--(xmax,0) node[shift={(-100:7pt)},font=\normalsize]{$ t(s) $};
\draw[->] (0,ymin)--(0,ymax) node[shift={(90:7pt)},font=\normalsize]{$ v$ (m/s)};
\fill (0,0) node[shift={(225:7pt)},font=\normalsize]{$ O $};
\foreach \x in { 1, 2}{
\draw (\x,2pt)--(\x,-2pt) +(0,-9pt) node[font=\scriptsize,fill=white,inner sep=1pt]{$\x$};
}
\draw (2pt,2)--(-2pt,2) +(-3pt,0) node[font=\scriptsize,anchor=east,fill=white,inner sep=1pt]{$2$};
\draw[dashed] (2,2)--(0,2) (2,0)--(2,2) (1,0)--(1,2);
\clip (xmin,ymin) rectangle (xmax,ymax);
\draw  plot[domain=0:1] (\x, {f(\x)});
\draw  plot[domain=1:2] (\x, {g(\x)});
\end{tikzpicture}
\end{center}
Quãng đường vật di chuyển được trong $2$ giây đầu tiên là $s$ mét. Khi đó $s$ có giá trị bằng bao nhiêu?
\par\shortans[oly]{3}
\loigiai{
Ta có $S=\displaystyle\int\limits_0^2v(t)\mathrm{\,d}t=\displaystyle\int\limits_0^1v(t)\mathrm{\,d}t+\displaystyle\int\limits_1^2v(t)\mathrm{\,d}t=\dfrac{1\cdot 2}{2}+1\cdot 2=3$ m.
}
\end{ex}

\begin{ex}%[12-EX-3-2026, Nguyễn Cường]%[2D4V2-6]
\immini{
Một chất điểm chuyển động theo quy luật vận tốc $v(t)$ (m/s) có đồ thị như hình vẽ, theo cung đường từ $A$ đến $B$ với vận tốc đầu bằng $9$\,(m/s), đến $C$ và dừng lại tại $D$. Tính quãng đường chất điểm đi được từ khi bắt đầu chuyển động đến khi dừng lại theo đơn vị mét. Kết quả làm tròn đến chữ số thập phân thứ nhất.
}
{
\begin{tikzpicture}[scale=.5,>=stealth, font=\footnotesize, line join=round, line cap=round]
\path
(0,0)coordinate(O)
(0,9)coordinate(A)
(2,9)coordinate(B)
(7,4)coordinate(C)
(10,0)coordinate(D)
;
\draw
(A)--(B) (C)--(D)
;
\draw [->](-0.5,0)--(11,0)node[below]{$t$};
\draw [->](0,-0.5)--(0,10)node[left]{$v$};
\draw[smooth] plot[domain=2:7] (\x,{(\x)^2-10*(\x)+25})
;
\foreach \p/\r in {A/45,B/60,C/90,D/45,O/45}
\fill (\p) circle (1pt) node[shift={(\r:3mm)}]{$\p$};
\path
(2,9)node[right]{$v_1=9$}
(6,5)node[above]{$v_2(t)=t^2-10t+25$}
(9,2)node[right]{$v_3(t)=-\dfrac{4}{3}t+\dfrac{40}{3}$}
;
\end{tikzpicture}
}
\shortans[]{$35{,}7$}
\loigiai{
Tại vị trí $B$ thì $t^2-10t+25=9\Leftrightarrow t^2-10t+16=0\Leftrightarrow \hoac{&t=2\\&t=8~\text{(Loại)}.}$\\
Tại vị trí $C$ thì $t^2-10t+25=-\dfrac{4}{3}t+\dfrac{40}{3}\Leftrightarrow t^2-\dfrac{26}{3}t+\dfrac{35}{3}=0\Leftrightarrow\hoac{&t=7\\&t=\dfrac{5}{3}~\text{(Loại)}.}$\\
Tại vị trí $D$ thì $-\dfrac{4}{3}t+\dfrac{40}{3}=0\Leftrightarrow t=10$.\\
Quãng đường chất điểm đi được là
\[\displaystyle\int\limits_0^29\mathrm{\,d}t+\displaystyle\int\limits_2^7\left(t^2-10t+25\right)\mathrm{\,d}t+\displaystyle\int\limits_7^{10}\left(-\dfrac{4}{3}t+\dfrac{40}{3}\right)\mathrm{\,d}t=\dfrac{107}{3}\approx 35{,}7.\]
Vậy quãng đường chất điểm đi được là $35{,}7$ m.
}
\end{ex}

\begin{ex}%[12-EX-3-2026, Nguyễn Thành Nhân]%[2D4V2-6]
\immini{Một chiếc xe đua F1 đạt tới vận tốc lớn nhất là $360$ km/h. Đồ thị bên biểu thị vận tốc $v$ của xe trong $5$ giây đầu tiên kể từ lúc xuất phát. Đồ thị trong $2$ giây đầu là một phần của một parabol đỉnh tại gốc tọa độ $O$, giây tiếp theo là đoạn thẳng và sau đúng ba giây thì xe đạt vận tốc lớn nhất. Biết rằng mỗi đơn vị trục hoành biểu thị $1$ giây, mỗi đơn vị trục tung biểu thị $10$ m/s và trong $5$ giây đầu xe chuyển động theo đường thẳng. Hỏi trong $5$ giây đó xe đã đi được quãng đường là bao nhiêu mét?
\shortans[]{$350$}}
{\begin{tikzpicture}[font=\footnotesize, line join=round, line cap=round, >=stealth, scale=1,x=0.7cm, y=0.166666666cm]
\draw[smooth] plot[domain=0:2] (\x,{1.5*(\x)^2});
\draw (2,6)--(3,30)--(5,30);
\draw[dashed] (2,0)|- (0,6) (3,0)|-(0,30) (5,0)--(5,30);
\draw[->] (-0.5,0)--(6,0)node[below]{$t$ (s)};
\draw (0,0)node[below left]{$O$};
\draw[->] (0,-1)--(0,35)node[right]{$v$ (m/s)};
\foreach \x/\g/\n in {(2,0)/below/2, (3,0)/below/3, (5,0)/below/5, (0,6)/left/6, (0,30)/left/, (2,6)/above/\, , (3,30)/above/\, , (5,30)/above/\,}{
\fill \x circle (1pt)node[\g]{$\n$};
}
\end{tikzpicture}}

\loigiai{Ta có $360$ km/h $=$ $100$ m/s.\\
Gọi $(P)\colon v_1=at^2+bt+c$ (m/s) là vận tốc xe đi được trong $2$ giây đầu.\\
Vì $(P)$ đi qua điểm $O(0;0)$ nên $b=0$ và $c=0$.\\
Ta có $A(2;60)\in (P)$ nên $60=4a\Leftrightarrow a=\dfrac{3}{2}$. Do đó ta được $v_1=15t^2$.\\
Gọi $d\colon v_2=mt+n$ (m/s) là vận tốc xe đi được từ giây thứ $2$ đến giây thứ $3$.\\
Ta có $A(2;60)\in d$ và $B(3;100)\in d$ ta có hệ phương trình\\
\[\heva{&2m+n=60\\&3m+n=100}\Leftrightarrow\heva{&m=40\\&n=-20.}\]
Do đó ta được $v_2=40t-20$.\\
Từ giây thứ $3$ đến giây thứ $5$ xe chuyển động đều với vận tốc $v=100$.\\
Quãng đường đi được sau $5$ giây là
\begin{eqnarray*}
&& \displaystyle\int\limits_0^5 v(t) \mathrm{\,d}t=\displaystyle\int\limits_0^2 15t^2 \mathrm{\,d}t + \displaystyle\int\limits_2^3 (40t-20) \mathrm{\,d}t + \displaystyle\int\limits_3^5 100 \mathrm{\,d}x =320 \text{ (m).}
\end{eqnarray*}

}
\end{ex}

\begin{ex}%[12-EX-3-2026, Nguyễn Thành Nhân]%[2D4V2-6]
\immini{Một người có miếng tôn hình tròn có bán kính bằng $5$ m. Người này tính trang trí sơn vẽ trên tấm tôn đó, biết mỗi mét vuông sơn hết $100$ nghìn đồng. Tuy nhiên cần có một khoảng trống để treo tấm tôn nên người này bớt lại một phần tấm tôn nhỏ không trang trí (phần màu trắng như hình vẽ), trong đó $AB=6$ m. Hỏi khi trang trí xong người này hết bao nhiêu tiền chi phí (làm tròn đến hàng đơn vị nghìn
đồng)?}
{\begin{tikzpicture}[font=\footnotesize, line join=round, line cap=round,scale=.4, >=stealth]
\path 	(0:0) coordinate (O)
(4,3) coordinate (A)
(4,-3) coordinate (B);
\filldraw[fill=gray!30!white, draw=gray!30!white] (A)--(O)--(B)-- cycle;
\filldraw[fill=gray!30!white, draw=gray!30!white] (36.87:5) arc (36.87:323.13:5) -- (0,0) -- cycle;
\draw (A)--(B);
\draw (O) circle (5);
\foreach \x /\goc in {A/30,B/-30}
\fill[black] (\x) circle (1pt)
($(\x)+(\goc:4mm)$) node {$\x$};
\end{tikzpicture}}
\shortans[oly]{$7445$}
\loigiai{
\immini{Gắn hệ trục toạ độ như hình vẽ.\\
Khi đó đường tròn có phương trình $x^2+y^2=25$.\\
Suy ra $y=\pm \sqrt{25-x^2}$.\\
Do $AB=6$ nên $AH=3$.\\
Suy ra $OH=\sqrt{OA^2-AH^2}=\sqrt{5^2-3^2}=4$.\\
Khi đó khoảng trống để treo tấm tôn chính là phần hình phẳng giới hạn bởi hai đồ thị đồ thị $y=\pm \sqrt{25-x^2}$ và hai đường thẳng $x=4;x=5$.\\
Nên diện tích khoảng trống treo tấm tôn là \[S_1=\displaystyle\int\limits_{4}^{5}{\left( \sqrt{25-x^2}-\left(-\sqrt{25-x^2}\right) \right)\mathrm{\,d}x}\approx 4{,}09\,\, (\mathrm{m}^2).\]
}
{\begin{tikzpicture}[font=\footnotesize, line join=round, line cap=round,scale=.5, >=stealth]
\path 	(0:0) coordinate (O)
(4,3) coordinate (A)
(4,-3) coordinate (B)
(4,0) coordinate (H);
\draw[->] (-6,0)--(6,0) node[below] {$x$};
\draw[->] (0,-6)--(0,6) node[left] {$y$};
\draw (5,0) node[above right] {$5$};
\draw (4.2,1.5) node {$3$};
\draw (O) circle (5);
\draw (O)--(A)--(B)--(O);
\foreach \x /\goc in {A/30,B/-30,O/-135,H/45}
\fill[black] (\x) circle (1pt)
($(\x)+(\goc:4mm)$) node {$\x$};
\end{tikzpicture}}
\noindent	Diện tích cả tấm tôn hình tròn là $S=\pi\cdot5^2=25\pi$ (m$^2$).\\
Diện tích phần trang trí là $S_2=S-S_1\approx 74{,}45$ (m$^2$).\\
Vậy số tiền chi phí là $100\cdot74{,}45=7445$ (nghìn đồng).
}
\end{ex}

\begin{ex}%[12-EX-2-2025, Đình Nguyên]%[2D4V2-6]
\immini{Một bồn chứa nước có dạng hình trụ với chiều cao $2$ m và bán kính đáy $0{,}5$ m. Lúc đầu bồn chứa đầy nước. Người ta tiến hành vặn van ở đáy bồn để xả nước. Kể từ khi bắt đầu xả nước, tốc độ thay đổi chiều cao của mực nước trong bồn theo thời gian $t$ là $h'(t) = \dfrac{t}{25} - \dfrac{2}{5}$ (m/phút). Sau khi xả $5$ phút, trong bồn còn bao nhiêu lít nước (kết quả làm tròn đến hàng đơn vị)?
\shortans{393}
}
{\begin{tikzpicture}[scale=0.7, font=\footnotesize, line join=round, line cap=round, >=stealth]
% Định nghĩa các tham số cơ bản
\pgfmathsetmacro{\R}{1.5} % Bán kính
\pgfmathsetmacro{\H}{4} % Chiều cao tổng thể của bình
\pgfmathsetmacro{\h}{2} % Chiều cao của nước
\pgfmathsetmacro{\waterLevel}{\h}
\pgfmathsetmacro{\nozzleY}{-0.2} % Vị trí vòi nước (trục y)
\pgfmathsetmacro{\nozzleX}{\R} % Vị trí vòi nước (trục x)

% Nước
% Đáy nước
\draw[fill=gray!50, opacity=0.8] (0, 0) ellipse ({\R} and {0.5*\R});

% Thân nước
\draw[fill=gray!50, opacity=0.6] ({-\R}, 0) -- ({-\R}, \waterLevel) arc (180:360:{\R} and {0.5*\R}) -- ({\R}, \waterLevel) -- ({\R}, 0);

% MẶT NƯỚC TRÊN (Phần được tô màu)
\draw[gray!70, fill=gray!30, opacity=0.8, thick] (0, \waterLevel) ellipse ({\R} and {0.5*\R});

% Thân bình (phần trên mực nước)
\draw[thick] ({-\R}, \waterLevel) -- ({-\R}, \H) arc (180:360:{\R} and {0.5*\R}); % Thân bình phía trước
\draw[thick] ({\R}, \waterLevel) -- ({\R}, \H);

% Miệng bình
\draw[thick] (0, \H) ellipse ({\R} and {0.5*\R});

% Ký hiệu chiều cao h
%			\draw[<->, line width=1pt] ({-\R-0.5}, 0) -- node[pos=0.5,fill=white] {$h$} ({-\R-0.5}, \h);
%			\draw[dashed] ({-\R-0.5}, 0)--(-\R,0) ({-\R-0.5}, \h)--(-\R,\h) ;
\end{tikzpicture}}
\loigiai{
Chiều cao mực nước giảm đi sau $5$ phút xả là:
$\Delta h = \left| \displaystyle\int\limits_0^5 \left( \dfrac{t}{25} - \dfrac{2}{5} \right) \mathrm{\,d}t \right| = \left| \left. \left( \dfrac{t^2}{50} - \dfrac{2t}{5} \right) \right|_0^5 \right| = |-1{,}5| = 1{,}5$ m.\\
Chiều cao mực nước còn lại trong bồn là $h_{cl} = 2 - 1{,}5 = 0{,}5$ m.\\
Thể tích nước còn lại là $V = \pi \cdot R^2 \cdot h_{cl} = \pi \cdot (0{,}5)^2 \cdot 0{,}5 = \dfrac{\pi}{8} \approx 0{,}3927$ $\text{m}^3$.\\
Đổi ra lít: $0{,}3927 \cdot 1\,000 \approx 393$ lít.
}
\end{ex}

\begin{ex}%[12-EX-3-2026, Trương Quan Kía]%[2D4V2-6]
Tại một nhà máy, $C(x)$ (triệu đồng) là tổng chi phí để sản xuất $x$ tấn sản phẩm B trong tháng.
Chi phí cận biên $C'(x)=8-0{,}13x+0{,}00022x^2$ với $0\le x\le 88$. Biết $C(0)=24$.
Tính $C(82)$ (làm tròn đến hàng đơn vị).
\shortans[]{$283$}
\loigiai{
\begin{eqnarray*}
&&C(82)=C(0)+\int_0^{82}C'(x)\,\mathrm{d}x
=24+\int_0^{82}\left(8-0{,}13x+0{,}00022x^2\right)\mathrm{d}x.\\
&&	C(82)=24+\left(8x-0{,}065x^2+\dfrac{0{,}00022}{3}x^3\right)\bigg|_0^{82}.
\end{eqnarray*}
Tính được
\[C(82)=24+\Bigl(656-437{,}06+40{,}361\Bigr)\approx 283{,}301\approx 283\ \text{(triệu đồng)}.\]
}
\end{ex}

\begin{ex}%[2D4V2-6]
\immini{ Sân vận động Sport Hub (Singapore) là sân có mái vòm kỳ vĩ nhất thế giới. Đây là nơi diễn ra lễ khai mạc Đại hội thể thao Đông Nam Á được tổ chức tại Singapore năm 2015.
Nền sân là một elip $(E)$ có trục lớn dài $150$ (m), trục bé dài $90$ (m) (hình vẽ). Nếu cắt sân vận động theo một mặt phẳng vuông góc với trục lớn của $(E)$ và cắt elip ở $M$, $N$ (hình vẽ) thì ta được thiết diện luôn là một phần của hình tròn có tâm $I$ (phần tô đậm trong hình vẽ) với MN là một dây cung và góc $\widehat{MIN}=90^{\circ}$. Để lắp máy điều hòa không khí thì các kỹ sư cần tính thể tích phần không gian bên dưới mái che và bên trên mặt sân, coi như mặt sân là một mặt phẳng và thể tích vật liệu là mái không đáng kể. Thể tích tính được là bao nhiêu (đơn vị nghìn m$^3$, \textit{làm tròn kết quả đến hàng đơn vị}).}
{
\begin{tikzpicture}[scale=1, font=\footnotesize,line join=round, line cap=round, >=stealth]
\path
(0,0) coordinate (I)
(45:2) coordinate (M)
(135:2) coordinate (N);
\foreach \i/\g in{I/-90,M/60,N/120}
\fill (\i) circle (1.5pt) node[shift={(\g:3mm)}]{$\i$};
\fill[blue!30] (M) -- (N) arc[start angle=135, end angle=45, radius=2] -- cycle;
\draw[thick] (I) circle (2cm) (M)--(N);
\draw (M)--(I)--(N);
\foreach \x/\y/\z in {M/I/N} {\draw pic[angle radius=1.5mm,draw=black] {right angle = \x--\y--\z};}
\begin{scope}[shift={(0,-4cm)}] % Di chuyển hình elip 4 cm qua phải
\path
(-3,0) arc (180:60:{3cm} and {1.5cm}) coordinate (M)
(-3,0) arc (-180:-60:{3cm} and {1.5cm}) coordinate (N)
;
\draw[thick]
(0,0) ellipse (3cm and 1.5cm)
(-3,0)--(3,0);
\fill
(-3,0) circle (1.5pt)
(3,0) circle (1.5pt);
\foreach \i/\g in{M/90,N/-90}
\fill (\i) circle (1.5pt) node[shift={(\g:3mm)}]{$\i$};
\draw (M)--(N);
\end{scope}
\end{tikzpicture}
}
\par	\shortans[0]{116}
\loigiai
{
\immini{Chọn hệ trục $Oxy$ như hình vẽ.\\
Phương trình chính tắc của elip đáy\\
$(E)\colon \dfrac{x^{2}}{75^{2}}+\dfrac{y^{2}}{45^{2}}=1$.\\
Gọi $M(x;y) \Rightarrow MN=2y=2 \sqrt{45^{2}\left(1-\dfrac{x^{2}}{75^{2}}\right)}=90 \sqrt{1-\dfrac{x^{2}}{75^{2}}}$.\\
$\Rightarrow R=\dfrac{MN}{\sqrt{2}}=\dfrac{90}{\sqrt{2}} \cdot \sqrt{1-\dfrac{x^{2}}{75^{2}}}$.\\
}{
\begin{tikzpicture}[scale=.8, font=\footnotesize,line join=round, line cap=round, >=stealth]
\path
(0,0) coordinate (O)
(-3,0) arc (180:60:{3cm} and {1.5cm}) coordinate (M)
(-3,0) arc (-180:-60:{3cm} and {1.5cm}) coordinate (N)
;
\draw[thick]
(0,0) ellipse (3cm and 1.5cm);
\draw[->]
(-4,0)--(4,0) node[right]{$x$};
\draw[->]
(0,-2)--(0,2) node[above]{$y$};
\fill
(-3,0) circle (1.5pt)
(3,0) circle (1.5pt);
\foreach \i/\g in{O/-120,M/90,N/-90}
\fill (\i) circle (1.5pt) node[shift={(\g:3mm)}]{$\i$};
\end{tikzpicture}
}
\immini{Ta có công thức diện tích hình quạt
$S_{\text{quạt}}=\dfrac{\pi R^{2}\cdot\alpha}{360^{\circ}}$,
với $\alpha$ (theo đơn vị độ) là số đo góc ở tâm chắn cung tương ứng.\\
Gọi $S(x)$ là diện tích thiết diện, ta có
\allowdisplaybreaks
\begin{eqnarray*}
S(x)&=&S_{\text{quạt}}-S_{\triangle IMN}\\
&=&\dfrac{1}{4} \pi R^{2}-\dfrac{1}{2} R^{2}\\
&=&\left(\dfrac{1}{4} \pi-\dfrac{1}{2}\right)R^{2}\\
&=&(\pi-2)\dfrac{2\,025}{2}\cdot\left(1-\dfrac{x^{2}}{75^{2}}\right).
\end{eqnarray*}
}{			\begin{tikzpicture}[scale=1, font=\footnotesize,line join=round, line cap=round, >=stealth]
% Vẽ hình tròn
\path
(0,0) coordinate (I)
(45:2) coordinate (M)
(135:2) coordinate (N)
;
\foreach \i/\g in{I/-90,M/60,N/120}
\fill (\i) circle (1.5pt) node[shift={(\g:3mm)}]{$\i$};
\fill[blue!30] (M) -- (N) arc[start angle=135, end angle=45, radius=2] -- cycle;
\draw[thick] (I) circle (2cm);
\draw[thick] (M)--(I)node[pos=.5,right]{$R$}--(N)--cycle;
\end{tikzpicture}
}

Thể tích khoảng không gian sân vận động là
\[V=\displaystyle\int\limits_{-75}^{75}(\pi-2) \dfrac{2\,025}{2} \cdot\left(1-\dfrac{x^{2}}{75^{2}}\right)\,\mathrm{d}x = 115\,586{,}2562 \, \mathrm{m}^{3} \approx 1{,}16 \cdot 10^5 \, (\mathrm{m}^{3}).\]
Vậy thể tích cần tìm là $116$ nghìn m$^3$.
}
\end{ex}

\begin{ex}%[2D4V2-6]
Một vật chuyển động với gia tốc được cho bởi hàm số $a(t)=5 \cos t$ (m/s$^2$). Lúc bắt đầu chuyển động vật có vận tốc $2{,}5$ m/s. Tính gia tốc của vật tại thời điểm vận tốc đạt giá trị lớn nhất trong $\pi$ (s) đầu tiên.
\par
\shortans[]{$0$}
\loigiai{
Vận tốc của vật được biểu diễn bởi hàm số \[v(t)=\displaystyle\int a(t) \,\mathrm{d} t=\displaystyle\int 5 \cos t \,\mathrm{d} t=5 \sin t+C.\]
Khi bắt đầu chuyển động, vật có vận tốc $2{,}5$ m/s nên ta có
\[v(0)=2{,}5\Leftrightarrow 5\sin 0+C=2{,}5\Leftrightarrow C=2{,}5.\]
Suy ra $v(t)=5\sin t+2{,}5$.\\
Mà $5\sin t+2{,}5\le 7{,}5$. \\
Vậy vận tốc đạt giá trị lớn nhất tại $t=\dfrac{\pi}{2}$. \\
Khi đó, gia tốc của vật tại thời điểm $t=\dfrac{\pi}{2}$ là $a\left(\dfrac{\pi}{2}\right)=5 \cos \left(\dfrac{\pi}{2}\right)=0$ (m/s$^2$).
}
\end{ex}

\begin{ex}%[DuAnC-dot3, Đoàn Hùng]%[2D4V2-6]
\immini[thm]{Kiến trúc sư thiết kế một khu sinh hoạt cộng đồng có dạng hình chữ nhật với chiều rộng và chiều dài lần lượt là $60$m và $80$m. Trong đó, phần được tô màu đậm là sân chơi, phần còn lại để trồng hoa. Mỗi phần trồng hoa có đường biên cong là một phần của parabol với đỉnh thuộc một trục đối xứng của hình chữ nhật và khoảng cách từ đỉnh đó đến trung điểm cạnh tương ứng của hình chữ nhật bằng $20$m (xem hình minh họa). Diện tích của phần sân chơi là bao nhiêu mét vuông?}{
\begin{tikzpicture}[scale=0.75, font=\footnotesize, line join=round, line cap=round, >=stealth]
\fill[gray,smooth] (-3,0)--plot[domain=-3:3](\x,{2/9*(\x)^2+6})--(3,0)--cycle;
\fill[white,smooth] (-3,0)--plot[domain=-3:3](\x,{-2/9*(\x)^2+2})--(3,0)--cycle;
\draw[<->, line width=0.8pt](0,6)--(0,8);
\draw[<->, line width=0.8pt](0,0)--(0,2);
\node at (0, 0) [below]{$60$\,m};
\node at (0, 7) [right]{$20$\,m};
\node at (0, 1) [right]{$20$\,m};
\node at (3, 4) [right]{$80$\,m};
\draw (-3,0)--(-3,8)--(3,8)--(3,0)--(-3,0);
\clip(-3,0) rectangle(3.05,8);
\draw[smooth,samples=300] plot(\x,{2/9*(\x)^2+6});
\draw[smooth,samples=300] plot(\x,{-2/9*(\x)^2+2});
\end{tikzpicture}
}
\par
\shortans[]{3200}
\loigiai{
\immini{	Diện tích của phần sân chơi bằng diện tích của hình chữ nhật trừ đi diện tích của hai phần trồng hoa.\\
Diện tích của hình chữ nhật là $60\cdot 80=4\,800$\,m$^2$.\\
Diện tích của mỗi phần trồng hoa bằng diện tích của một parabol.\\
Chọn hệ trục $Oxy$ như hình vẽ.\\
Phương trình parabol bên dưới có dạng:\\ $y=ax^2+b$ (vì đỉnh parabol thuộc trục tung).\\
Ta có $\heva{&y(0)=20\\&y(30)=0} \Leftrightarrow \heva{&b=20\\&a\cdot 30^2+20=0} \Rightarrow \heva{&a=-\dfrac{1}{45}\\&b=20.} $\\
Phương trình Parabol dưới là $y=-\dfrac{1}{45}x^2+20$.\\
Diện tích của mỗi phần trồng hoa bằng
\[\displaystyle\int \limits_{-30}^{30} \left(-\dfrac{1}{45}x^2+20\right)\mathrm{d}x=800 (\mathrm{m^2}).\]
Diện tích của phần sân chơi là: $4800-2\cdot 800=3200 (\mathrm{m^2}).$}{\begin{tikzpicture}[scale=.7, font=\footnotesize, line join=round, line cap=round, >=stealth]
\def\xt{-3.5} \def\xp{4} \def\yt{9} \def\yd{-1}
\fill[gray,smooth] (-3,0)--plot[domain=-3:3](\x,{2/9*(\x)^2+6})--(3,0)--cycle;
\fill[white,smooth] (-3,0)--plot[domain=-3:3](\x,{-2/9*(\x)^2+2})--(3,0)--cycle;
\draw[->, line width=0.8pt](\xt, 0)--(\xp, 0) node[below]{$x$};
\draw[->, line width=0.8pt](0,\yd)--(0,\yt) node[left]{$y$};
\node at (0, 0) [below left]{$O$};
\node at (-1.5, 8) [above]{$60$\,m};
\node at (0, 1) [left]{$20$\,m};
\node at (3, 4) [right]{$80$\,m};
\clip(-3,0) rectangle(3.05,8);
\draw[smooth,samples=300] plot(\x,{2/9*(\x)^2+6});
\draw[smooth,samples=300] plot(\x,{-2/9*(\x)^2+2});
\draw (-3,0)--(-3,8)--(3,8)--(3,0);
\end{tikzpicture}}

}
\end{ex}

\begin{ex}%[2D4V2-6] %[Thanh Phát]
Trong kinh tế học, người ta tính toán thu nhập của một khoản đầu tư thông qua công thức
$FV = \displaystyle\int\limits_0^T f(t) \mathrm{e}^{-rt} \mathrm{\,d}t$ trong đó với khoảng thời gian $t$ thì $f(t)$ là tốc độ sinh lời, $T$ là thời gian đầu tư và $r$ là lãi suất hàng kỳ. Bạn Nam đang cố gắng để cân nhắc giữa hai khoản đầu tư như sau
\begin{itemize}
\item Khoản đầu tư thứ nhất có giá $1\,000\,$USD và dự kiến sẽ tạo ra một dòng thu nhập liên tục với tốc độ $f_1(t) = 3\,000\mathrm{e}^{0{,}03t}$ (USD) mỗi năm.
\item Khoản đầu tư thứ hai có giá $4\,000\,$USD và được ước tính sẽ tạo ra thu nhập với tốc độ không đổi là $f_2(t) = 4\,000\mathrm{e}^{0{,}03t}$ (USD) mỗi năm.
\end{itemize}
Giả sử lãi suất hằng năm hiện hành vẫn cố định ở mức $5\%$ mỗi năm và được tính lãi kép liên tục. Em hãy giúp Nam chọn khoản đầu tư lời hơn bằng cách tính số Dollar chênh lệch giữa hai khoản đầu tư sau $5$ năm. (Lấy giá trị dương và làm tròn kết quả đến hàng đơn vị)
\par
\shortans[0]{1758}
\loigiai{
Ta biết rằng trong thực tế: Lợi nhuận $=$ Thu nhập $-$ Chi phí.
\begin{itemize}
\item Tính giá trị từ dòng thu nhập của khoản đầu tư thứ nhất.\\
Khoản đầu tư thứ nhất có giá $ 1\,000$ USD và tốc độ sinh lời $f_1(t) = 3\,000e^{0{,}03t}$.
\begin{eqnarray*}
FV_1 &=& \displaystyle\int_{0}^{5} 3\,000e^{0{,}03t} \cdot e^{-0,05t}\mathrm{\,d}t\\
&=&  3\,000 \displaystyle\int_{0}^{5} e^{-0{,}02t}\,\mathrm{d}t\\
&=& 3\,000\cdot \dfrac{e^{-0{,}02t}}{-0{,}02}\Bigg|_0^5\\
&\approx & 14\,274{,}39.
\end{eqnarray*}
Lợi nhuận ròng của khoản thứ nhất là
\[P_1 = FV_1 - 1\,000 = 14\,274{,}39 - 1\,000 = 13\,274{,}39.\]
\item Tính giá trị từ dòng thu nhập của khoản đầu tư thứ hai.\\
Khoản đầu tư thứ hai có giá $4\,000$ USD và tốc độ sinh lời $f_2(t) = 4\,000e^{0{,}03t}$.
\begin{eqnarray*}
FV_2 &=& \displaystyle\int_{0}^{5} 4\,000e^{0{,}03t} \cdot e^{-0{,}05t}\,\mathrm{d}t\\
&=& 4\,000 \displaystyle\int_{0}^{5} e^{-0{,}02t} \,\mathrm{d}t\\
&=& 4\,000 \cdot\dfrac{e^{-0{,}02t}}{-0{,}02}\Bigg|_0^5\\
&\approx & 19\,032{,}52.
\end{eqnarray*}
Lợi nhuận ròng của khoản đầu tư thứ hai là
\[P_2 = FV_2 - 4\,000 = 19\,032{,}52 - 4\,000 = 15\,032{,}52.\]
\end{itemize}
Số tiền chênh lệch giữa hai khoản đầu tư là
\[\Delta P = \left|P_2 - P_1\right| = \left|15\,032{,}52 - 13\,274{,}39\right| = 1\,758{,}13 \text{ USD}.\]
Làm tròn đến hàng đơn vị: $1\,758$ USD.
}
\end{ex}

\begin{ex}%[2D4V2-6]
Hệ thống lọc nước bể bơi vô cùng quan trọng khi tiến hành xây dựng công trình bơi lội để nguồn nước được làm sạch thường xuyên và giữ vệ sinh cho người bơi. Trong quá trình vận hành lọc nước thì lượng nước trong bể sẽ thay đổi theo thời gian. Lượng nước trong bể giảm nếu hệ thống đang xả nước bẩn ra khỏi bể và tăng nếu hệ thống đang cấp thêm nước sạch cho bể. Biết rằng $1$ gallon gần bằng $3{,}785$ lít, dung tích của bể là $1\,000$ gallon và thời điểm $6$ giờ sáng bể chứa $250$ gallon nước. Hàm số $f(t)$ biểu thị cho tốc độ thay đổi lượng nước trong bể theo thời gian $t$ giờ, từ thời điểm $6$ giờ sáng đến $6$ giờ chiều được cho bởi công thức $f(t) = \heva{&100t &\text{ khi } 0 \leq t \leq 3\\&900-200t &\text{ khi } 3 \leq t \leq 6\\&100t - 900 &\text{ khi } 6 \leq t \leq 12}$ với mốc thời gian $t= 0$ tại thời điểm $6$ giờ sáng. Hỏi ở thời điểm $6$ giờ chiều thì trong bể chứa bao nhiêu gallon nước?
\begin{center}
\definecolor{waterTop}{HTML}{4FC3F7} % Màu mặt nước
\definecolor{waterSide}{HTML}{29B6F6} % Màu hông khối nước
\definecolor{tileColor}{HTML}{E1F5FE} % Màu gạch lát
\definecolor{poolWall}{HTML}{B0BEC5} % Màu tường ngoài
\definecolor{deckColor}{HTML}{D7CCC8} % Màu sàn gỗ
\definecolor{filterBody}{HTML}{F5F5F5} % Màu bình lọc
\definecolor{pumpBody}{HTML}{37474F} % Màu máy bơm
\begin{tikzpicture}[
x={(0.9cm, -0.3cm)}, % Trục x
y={(0.6cm, 0.5cm)}, % Trục y
z={(0cm, 1cm)}, % Trục z
scale=1.2,
line join=round, line cap=round
]
\fill[deckColor] (-1, -1, 0) -- (6, -1, 0) -- (6, 5, 0) -- (-1, 5, 0) -- cycle;
\foreach \i in {-0.5, 0, ..., 4.5} {
\draw[deckColor!80!black, thin] (-1, \i, 0) -- (6, \i, 0);
}
\def\L{4} % Dài
\def\W{3} % Rộng
\def\H{1.5} % Cao
\def\Water{1.2} % Mực nước
\fill[tileColor!80!gray] (0, \W, 0) -- (\L, \W, 0) -- (\L, \W, \H) -- (0, \W, \H) -- cycle;
\fill[tileColor!90!gray] (0, 0, 0) -- (0, \W, 0) -- (0, \W, \H) -- (0, 0, \H) -- cycle;
\fill[waterSide!80!black] (0,0,0) -- (\L,0,0) -- (\L,\W,0) -- (0,\W,0) -- cycle;
\fill[waterTop, opacity=0.8] (0,0,\Water) -- (\L,0,\Water) -- (\L,\W,\Water) -- (0,\W,\Water) -- cycle;
\draw[white, opacity=0.5, thick] (0.5, 0.5, \Water) .. controls (1, 0.7, \Water) and (1.5, 0.3, \Water) .. (2, 0.5, \Water);
\draw[white, opacity=0.5, thick] (2.5, 2.5, \Water) .. controls (3, 2.7, \Water) and (3.5, 2.3, \Water) .. (\L-0.2, 2.5, \Water);
\fill[poolWall] (0,0,0) -- (\L,0,0) -- (\L,0,\H) -- (0,0,\H) -- cycle; % Trước
\fill[poolWall!80!black] (\L,0,0) -- (\L,\W,0) -- (\L,\W,\H) -- (\L,0,\H) -- cycle; % Phải
\draw[white, thick] (0,0,\H) -- (\L,0,\H) -- (\L,\W,\H);
\draw[white, thick] (\L,0,\H) -- (\L,0,0);
\begin{scope}[shift={(\L+0.8, 1, 0)}]
\fill[gray!50] (-0.5, -0.5, 0) rectangle (2, 2.5);
\node[cylinder, shape border rotate=90, draw=gray, fill=filterBody,
minimum height=1.8cm, minimum width=1.0cm, aspect=0.6]
(tank) at (1.0, 1.0, 0.8) {};
\fill[gray!80] (1.0, 1.0, 1.7) ellipse (0.5cm and 0.15cm);
\fill[black] (1.0, 1.0, 1.75) circle (0.1); % Đồng hồ áp suất
\fill[pumpBody] (0, 0.5, 0) -- (0.6, 0.5, 0) -- (0.6, 1.5, 0) -- (0, 1.5, 0) -- cycle;
\fill[pumpBody!80] (0, 0.5, 0.5) -- (0.6, 0.5, 0.5) -- (0.6, 1.5, 0.5) -- (0, 1.5, 0.5) -- cycle;
\fill[pumpBody!70] (0.6, 0.5, 0) -- (0.6, 0.5, 0.5) -- (0.6, 1.5, 0.5) -- (0.6, 1.5, 0) -- cycle;
\draw[gray, double=red!10, double distance=2pt, line width=0.5pt]
(-0.8, 1.0, 0.5) -- (0, 1.0, 0.5);
\fill[gray] (-0.8, 1.0, 0.5) circle (2pt); % Điểm nối vào tường hồ
\draw[gray, double=blue!10, double distance=2pt, line width=0.5pt]
(tank.east) -- ++(0.3, 0, 0) -- ++(0, -1.5, 0) -- ++(-1.5, 0, 0) -- (-0.8, -0.5, 0.8);
\fill[gray] (-0.8, -0.5, 0.8) circle (2pt); % Điểm nối vào tường hồ (trả nước)
\end{scope}
\begin{scope}[shift={(1.5, 0, \H)}]
\foreach \x in {0, 0.6} {
\draw[gray!30, double=gray, double distance=1.5pt] (\x,0,0) -- (\x,-0.3,0) -- (\x,-0.3,-1.2);
}
\foreach \z in {-0.3, -0.7, -1.1} {
\draw[gray!30, double=gray, double distance=1.5pt] (0,-0.3,\z) -- (0.6,-0.3,\z);
}
\end{scope}
\tikzstyle{lbl}=[fill=white, inner sep=1.5pt, rounded corners=2pt, opacity=0.9, text opacity=1, font=\tiny\sffamily, align=center, draw=gray!20]
\node[lbl] at (\L/2, \W/2, \Water) {BỂ BƠI\\(1000 gallon)};
\draw[->, red, thick] (\L, 1.5, 0.5) -- (\L+0.5, 1.5, 0.5) node[right, lbl] {Nước ra};
\draw[<-, blue, thick] (\L, 0.5, 0.8) -- (\L+0.5, 0.5, 0.8) node[right, lbl] {Nước vào};
\end{tikzpicture}
\end{center}
\par\shortans[]{$700$}
\loigiai{
Ở thời điểm $6$ giờ chiều thì trong bể chứa số gallon nước là \begin{align*}
V = 250 + \displaystyle\int\limits_0^3 100t\,\mathrm{d}t
+ \displaystyle\int\limits_3^6 \left(900 - 2t \right) \,\mathrm{d}t + \displaystyle\int\limits_6^{12} \left(100t-900 \right) \,\mathrm{d}t = 700.  \end{align*}
}
\end{ex}

\begin{ex}%[Dự án C THPTQG 2025]%[Võ Hoàng Nghĩa]%[2D4V2-6]
\immini[thm]
{Nhà bác Ba có tất cả $8$ cánh cửa sắt hình chữ nhật với chiều dài $2$ (m) và chiều rộng $1$ (m). Hai mặt của mỗi cánh cửa được thiết kế như hình vẽ dưới đây. Trong đó, phần được tô đậm được sơn màu xanh, phần còn lại được sơn màu trắng. Mỗi phần sơn màu trắng có đường biên cong là một phần của parabol có đỉnh nằm trên cạnh của hình chữ nhật. Biết rằng chi phí để sơn màu xanh là $120$ (nghìn đồng/m$^2$) và chi phí sơn màu trắng là $110$ (nghìn đồng/m$^2$). Hỏi, để sơn toàn bộ số cửa sắt trên, bác Ba phải tốn bao nhiêu triệu đồng \textit{(làm tròn kết quả đến hàng phần trăm theo đơn vị triệu đồng)}?
\par\shortans[0]{3,63}
}
{\begin{tikzpicture}[line cap=round,line join=round,font=\footnotesize,>=stealth,x=2cm,scale=0.6]
\draw[pattern=north east lines](-1,0)--plot[domain=-1:0] (\x,{4*(\x+1)^2})--plot[domain=0:1] (\x,{4*(\x-1)^2})--(1,0)--plot[domain=1:0] (\x,{-4*(\x-1)^2})--plot[domain=0:-1] (\x,{-4*(\x+1)^2})--(-1,0);
\draw(-1,4)rectangle(1,-4) (0,0)node[below left]{$O$};
\end{tikzpicture}
}
\loigiai{
\immini
{Xét diện tích phần son màu xanh của một phần tư cánh cửa.\\
Chọn hệ trục tọa độ như hình vẽ. Đường biên cong là parabol có dạng $y=ax^2(a\neq0)$.\\
Do parabol đi qua điểm $\left(\dfrac{1}{2};1\right)$ nên ta có $1=\dfrac{1}{4}a\Rightarrow a=4$.\\
Vậy parabol là $y=4x^2$.\\
Diện tích hình phẳng giới hạn bởi parabol $y=4x^2$, trục hoành và hai đường thẳng $x=0,x=\dfrac{1}{2}$ là $\displaystyle\int\limits_0^{\tfrac{1}{2}}4x^2\mathrm{\,d}x=\left.\dfrac{4}{3}x^3\right|_0^{\tfrac{1}{2}}=\dfrac{1}{6}$.
}
{\begin{tikzpicture}[line cap=round,line join=round,font=\footnotesize,>=stealth,x=2cm]
\draw plot[domain=-1:1] (\x,{4*(\x)^2});
\draw[->](-1,0)--(1,0)node[below]{$x$};
\draw[->](0,-1)--(0,4)node[left]{$y$};
\draw[dashed](0.5,0)node[below]{$0{,}5$}--(0.5,1)--(0,1);
\end{tikzpicture}}
\noindent
Khi đó, phần được sơn màu xanh của 1 mặt cửa có diện tích là $4\cdot \dfrac{1}{6}=\dfrac{2}{3}$ (m$^2$), còn phần diện tích được sơn màu trắng là $2\cdot 1-\dfrac{2}{3}=\dfrac{4}{3}$ (m$^2$).\\
Chi phí để sơn một cánh cửa là $2\cdot \left(120\cdot \dfrac{2}{3}+110\cdot \dfrac{4}{3}\right)=\dfrac{1\,360}{3}$ (nghìn đồng).\\
Chi phí để sơn 8 cánh cửa là $8\cdot \dfrac{1\,360}{3}=\dfrac{10\,880}{3}$ (nghìn đồng) $\approx3{,}63$ (triệu đồng).
}
\end{ex}

\begin{ex}%[2D4V2-6]
Vận tốc (dặm/giờ) của một máy bay khi bay ngược chiều gió được cho bởi công thức $v(t) = 30\left(16-t^2 \right)$ với $0\leq t \leq 3$. Khi vận tốc tức thời đạt $400$ dặm/giờ thì máy bay đã đi qua quãng đường bao xa kể từ thời điểm bay ngược chiều gió (kết qua làm tròn đến hàng đơn vị của dặm)?
\shortans[]{740}
\loigiai{
Khi vận tốc tức thời đạt $400$ dặm/giờ tức là $30\left(16 - t^2 \right) = 400 \Leftrightarrow t = \dfrac{\sqrt{24}}{3}$ (giờ).\\
Khi đó máy bay đã bay được quãng đường là $s(t) = \displaystyle\int^{\dfrac{\sqrt{24}}{3}}_{0} 30\left(16-t^2 \right)\,\mathrm{d}t \approx 740$ (dặm).
}
\end{ex}

\begin{ex}%[Hiếu Mai]%[2D4V2-6]
Một vật được ném lên độ cao $300$ m với vận tốc được cho bởi công thức $v(t)=-9{,}81t+29{,}43$ (m/s) (\textit{Nguồn: R. Larson and B. Edwards, Calculus 10e, Cengage}). Gọi $h(t)$ (m) là độ cao của vật tại thời điểm $t$ (s). Sau bao lâu kể từ khi bắt đầu được ném lên thì vật đó chạm đất (làm tròn kết quả đến hàng đơn vị của mét)?
\shortans[oly]{$11$}
\loigiai
{
Ta có $h(t)=\displaystyle\int\limits v(t)\mathrm{\,d}t=\displaystyle\int\limits \left(-9{,}81t+29{,}43\right)\mathrm{\,d}t=-\dfrac{9{,}81}{2}t^2+29{,}43t+C$. \\
Vì vật được ném lên ở độ cao $300$ m nên $h(0)=300$ suy ra $C=300$. \\
Vậy $h(t)=-\dfrac{9{,}81}{2}t^2+29{,}43t +300$. \\
Khi vật bắt đầu chạm đất thì $h(t)=0$ nên $-\dfrac{9{,}81}{2}t^2+29{,}43t+300=0 \Leftrightarrow \hoac{ & t\approx11 \\ & t\approx-5.}$ \\
Vì $t>0$ nên $t\approx11$ (s).
}
\end{ex}

\begin{ex}%[Hiếu Mai]%[2D4V2-6]
Chủ một trung tâm thương mại muốn cho thuê một số gian hàng khác nhau. Người đó muốn tăng giá cho thuê của mỗi gian hàng thêm $x$ (triệu đồng) $(x\geqslant0)$. Tốc độ thay đổi doanh thu từ các gian hàng đó được biểu diễn bởi hàm số $T'(x)=-20x+300$, trong đó $T'(x)$ tính bằng triệu đồng (\textit{Nguồn: R.Larson and B.Edwards, Calculus 10e, Cengage}). Biết rằng nếu người đó tăng giá thuê cho mỗi gian hàng thêm $10$ triệu đồng thì doanh thu là $12\;000$ triệu đồng. Tìm giá trị của $x$ để người đó có doanh thu cao nhất.
\shortans[oly]{$15$}
\loigiai
{
Ta có $T(x)=\displaystyle\int\limits T'(x)\mathrm{\,d}x=\displaystyle\int\limits (-20x+300)\mathrm{\,d}x=-10x^2+300x+C$ với $C\in\mathbb{R}$. \\
Khi người đó tăng giá cho thuê mỗi gian hàng thêm $10$ triệu đồng thì doanh thu là $12\;000$ triệu đồng. \\
Do đó ứng với $x=10$ ta có $T(10)=12\;000$. \\
Suy ra $12\;000=-10\cdot10^2+300\cdot10+C \Rightarrow C=10\;000$. \\
Suy ra $T(x)=-10x^2+300x+10\;000$. \\
Ta có $T(x)$ là một hàm số bậc hai với hệ số $a<0$ và đồ thị hàm số có đỉnh là $I(15;12\;250)$. \\
Vậy doanh thu cao nhất mà người đó có thể thu về là $12\;250$ triệu đồng và khi đó mỗi gian hàng đã tăng giá cho thuê thêm $15$ triệu đồng.
}
\end{ex}

\begin{ex}%[HSA-đợt 1, Nguyễn Cường]%[2D4V2-6]
Một ô tô bắt đầu chuyển động nhanh dần đều với vận tốc $v_1(t)=7t$ (m/s). Đi được $5$ s, người lái xe phát hiện chướng ngại vật và phanh gấp, ô tô tiếp tục chuyển động chậm dần đều với gia tốc $a=-70$ (m/s$^2$). Quãng đường đi được của ô tô từ lúc bắt đầu chuyển bánh cho đến khi dừng hẳn dài bao nhiêu mét? (Kết quả làm tròn đến hàng phần mười).
\shortans{96{,}3}
\loigiai{
Quãng đường xe đi được trong $5$ giây đầu tiên là
\[\displaystyle\int\limits_0^57t\mathrm{\,d}t=87{,}5\,(\mathrm{m}). \]
Khi đó, vận tốc của xe là $v=7\cdot 5=35$ (m/s).\\
Gọi vận tốc khi xe chuyển động chậm dần đều là $v_2(t)$.\\
Ta có $v_2(t)=\displaystyle\int -70\mathrm{\,d}t=-70t+C$.\\
Khi bắt đầu giảm tốc là $t=0$ thì vận tốc xe là $35$ m/s nên $C=35$.\\
Suy ra $v_2(t)=-70t+35$.\\
Xe dừng hẳn khi $v_2(t)=0\Leftrightarrow -70t+35=0\Leftrightarrow t=0{,}5$\,(s).\\
Khi đó xe sẽ đi thêm được $\displaystyle\int\limits_0^{0{,}5}(-70t+35)\mathrm{\,d}t=8{,}75$.\\
Vậy xe di chuyển được tổng quãng đường là $96{,}25\approx 96{,}3$\,m.
}
\end{ex}

\begin{ex}%[HSA-L1-2025]%[Hiep Nguyen Quang]%[2D4V2-6]
Một ô tô bắt đầu chuyển động nhanh dần đều với vận tốc $v(t)=2t$ (\text{m/s}). Đi được $12$ giây người đó phát hiện chướng ngại vật rồi phanh gấp, ô tô tiếp tục chuyển động chậm dần đều với gia tốc $a=-12$ (\text{m/s}$^2$). Quãng đường ô tô đi được từ lúc bắt đầu chuyển bánh đến khi dừng hẳn là bao nhiêu?
\shortans[oly]{168}
\loigiai{
Quãng đường ô tô đi được từ lúc lăn bánh đến khi được phanh
\[S_1=\displaystyle\int\limits_0^{12}v_1(t)\,\mathrm{d}t=\displaystyle\int\limits_0^{12}2t\,\mathrm{d}t=144\,(\text{m}). \]
Vận tốc $v_2(t)$ (m/s) của ô tô từ lúc được phanh đến khi dừng hẳn thỏa mãn
\[v_2(t)=\displaystyle\int-12\,\mathrm{d}t=-12t+C. \]
Ta có $v_2(12)=v_1(12)=24\Rightarrow C=168\Rightarrow v_2(t)=-12t+168$ (\text{m/s}).\\
Thời điểm xe dừng hẳn tương ứng với thỏa mãn $v_2(t)=0\Leftrightarrow t=14$ (\text{s}).\\
Quãng đường ô tô đi được từ lúc xe được phanh đến khi dừng hẳn là
\[S_2=\displaystyle\int\limits_{12}^{14}v_2(t)\,\mathrm{d}t=\displaystyle\int\limits_{12}^{14}(-12t+168)\,\mathrm{d}t=24\,(\text{m}). \]
Quãng đường cần tính $s=s_1+s_2=144+24=168\,(\text{m})$.
}
\end{ex}

\begin{ex}%[2D4V2-6]
Một chất điểm bắt đầu chuyển động thẳng đều với vận tốc $v_0$, sau $6$ giây chuyển động thì gặp chướng ngại vật nên bắt đầu giảm tốc độ với vận tốc chuyển động $v(t)=-\dfrac{5}{2}t+a$\,(m/s), $(t\geq6)$ cho đến khi dừng hẳn, biết rằng kể từ lúc chuyển động đến lúc dừng thì chất điểm đi được quãng đường là $80$\,m. Tìm $v_0$.
\par\shortans[]{$10$}
\loigiai{
Phương pháp giải: Công thức tích phân.\\
Ta có $v(6)=v_0\Leftrightarrow a=v_0+15$ suy ra $v(t)=-\dfrac{5}{2}t+v_0+15$.
Gọi $n$ là thời điểm vật dừng hẳn, khi đó ta có
\[ v(n)=0\Leftrightarrow n=\dfrac{2}{5}\left(v_0+15\right)\Leftrightarrow n=\dfrac{2v_0}{5}+6. \]
Khi đó ta có phương trình tổng quãng đường vật đi được là
\begin{eqnarray*}
&&80=6\cdot v_0+\displaystyle\int\limits_6^n\left(-\dfrac{5}{2}t+v_0+15\right)\mathrm{\,d}t.\\
&\Leftrightarrow& 80=6\cdot v_0-\dfrac{5}{4}\left(n^2-6^2\right)+v_0(n-6)+15(n-6).\\
&\Leftrightarrow& 80=6\cdot v_0-\dfrac{5}{4}\left(\dfrac{4v_0^2}{25}+\dfrac{24v_0}{25}\right)+v_0\dfrac{2v_0}{5}+15\dfrac{2v_0}{5}\\
&\Leftrightarrow& v_0^2+36v_0-400=0
\Leftrightarrow v_0=10.
\end{eqnarray*}
}
\end{ex}

\begin{ex}%[Tổ 20 - Đợt 17 - Chương 4 - - CD - Đề 8]%[Huỳnh Xuân Tín]%[2D4V3-1]
Cho hàm số $f(x)=\heva{&7-4x^2&\text{khi}& 0\le x\le 1\\&4-x^2&\text{khi}&  x> 1}$. Diện tích hình phẳng giới hạn bởi hàm số $f(x)$ và các đường thẳng $x=0$, $x=3$, $y=0$ bằng
\shortans{$10$}
\loigiai
{Ta có với $x\in [0;1]$ hàm số $y=7-4x^2>0$ và hàm số $y=4-x^2$ cắt trục hoàng tại $x=2$ trong đoạn $[0;3]$. Khi đó diện tích cần tìm là
\[
\begin{aligned}
& S=\displaystyle\int\limits_0^1 \left( 7-4 x^3\right)  \mathrm{\,d}x+\displaystyle\int\limits_1^2\left(  4-x^2\right)  \mathrm{\,d}x+\displaystyle\int\limits_2^3 \left( x^2-4\right)  \mathrm{\,d}x \\
& =7 x-x^4\Bigg|_0 ^1+\left(4 x-\dfrac{x^3}{3}\right)\Bigg|_1 ^2+\left(\dfrac{x^3}{3}-4 x\right)\Bigg|_2 ^3=6+4-\dfrac{7}{3}-3-\dfrac{8}{3}+8=10.
\end{aligned}
\]
}
\end{ex}

\begin{ex}%[2D4V3-1]
\immini{Cho hai hàm số $f(x)=x^3+ax^2+bx+c$ và $g(x)=\mathrm{\,d}x+e$, $(a,b,c,d,e\in \mathbb{R})$. Biết rằng đồ thị của hàm số $y=f(x)$ và $y=g(x)$ cắt nhau tại ba điểm $A$, $B$, $C$ sao cho $BC=2AB$, với phần diện tích $S_1$, $S_2$ như hình vẽ. Khi đó $\dfrac{S_1}{S_2}$ bằng bao nhiêu? (Làm tròn kết quả đến chữ số thập phân thứ hai).
}
{
\begin{tikzpicture}[scale=1, font=\footnotesize, line join=round, line cap=round, >=stealth]
\tikzset{label style/.style={font=\footnotesize}}
%Nhập giới hạn đồ thị và hàm số cần vẽ
\def \xmin{-2.5}
\def \xmax{1.5}
\def \ymin{-0.5}
\def \ymax{3.5}
\def \hamso{sin((2*(\x))*180/pi)}
%\def \tiemcanxien{\x+1}
%Tự động
\draw[->] (\xmin,0)--(\xmax,0) node[below left] {$x$};
\draw[->] (0,\ymin)--(0,\ymax) node[below left] {$y$};
\draw[fill=black] (0,0) circle(1pt) node [below left] {$O$};
%Vẽ các điểm trên 2 hệ trục
%\foreach \x in {-1,2}
%\fill[black] (\x,0) circle(1pt) node [above left] {$\x$};
%Tự động
\begin{scope}
\clip (\xmin+0.01,\ymin+0.01) rectangle (\xmax-0.01,\ymax-0.01);
\draw[samples=350,domain=\xmin+0.01:\xmax-0.01,smooth,variable=\x] plot (\x,{(\x+2)*(\x)^2});
\draw[samples=350,domain=\xmin+0.01:\xmax-0.01,smooth,variable=\x] plot (\x,{\x+2});
\end{scope}
\draw[pattern = north east lines] (-2,0)--plot[domain=-2:1] (\x,{(\x+2)*(\x)^2})--plot[domain=1:-2](\x,{\x+2});
\end{tikzpicture}
}

\shortans[]{$5{,}16$}
\loigiai{
Ta có $y'=-f'(4-x)$.\\
Vì $BC=2AB$ nên ta chọn các nghiệm của phương trình hoành độ giao điểm là $-2$; $-1$; $1$.\\
Ta có $f(x)-g(x)=(x+2)(x+1)(x-1)$.\\
Ta được: $S_1=\displaystyle \int\limits_{-2}^{-1} \left| (x+2)(x+1)(x-1) \right|\mathrm{\,d}x
=\dfrac{5}{12}$.\\
$S_2=\displaystyle \int\limits_{-1}^1 \left| (x+2)(x+1)(x-1) \right|\mathrm{\,d}x=\dfrac{8}{3}$.\\
Vậy $\dfrac{S_1}{S_2}=\dfrac{5}{32}\approx 5{,}16$.
}
\end{ex}

\begin{ex}%[ST12-dot1-Trần xuân Hòa]%[2D4V3-1]
\immini{Trong mặt phẳng tọa độ $Oxy$, cho hình chữ nhật $(H)$ có một cạnh nằm trên trục hoành và có hai đỉnh trên một đường chéo là $A(-1;0)$ và $C\left( a;\sqrt{a}\right)$, với $a>0$. Biết rằng đồ thị của hàm số $y=\sqrt{x}$ chia hình $(H)$ thành hai phần có diện tích bằng nhau. Tìm $a$.
}{\begin{tikzpicture}[line join=round, line cap=round,>=stealth,scale=1.5]
\tikzset{label style/.style={font=\footnotesize}}
\draw[->] (-1.6,0)--(4,0) node[below left] {$x$};
\draw[->] (0,-0.4)--(0,2.5) node[below left] {$y$};
\draw (0,0) node [below left] {$O$};
\foreach \x in {-1,1,2,3}
\draw[thin] (\x,1pt)--(\x,-1pt) node [below] {$\x$};
\foreach \y in {1,2}
\draw[thin] (1pt,\y)--(-1pt,\y) node [below left] {$\y$};
\draw[samples=200,domain=0:3.5,smooth,variable=\x] plot (\x,{sqrt(\x)});
\path
(-1,0) coordinate (A)
(-1,1.73) coordinate (D)
(3,0) coordinate (B)
(3,1.73) coordinate (C)
;
\foreach \p/\q in {A/135,B/45,C/80,D/135}
\fill[black] (\p) circle (1.0pt) ($(\p)+(\q:3mm)$) node{$\p$};
\draw[pattern = north east lines, line width = 1.2pt,draw=none] (0,0)%
plot[domain=0:3] (\x,{sqrt(\x)})--(3,1.73) -- (3,0)--cycle;
\draw (A)--(D)--(C)--(B);
\end{tikzpicture}}
\shortans{$3$}
\loigiai{
\begin{center}
\begin{tikzpicture}[line join=round, line cap=round,>=stealth,scale=2]
\tikzset{label style/.style={font=\footnotesize}}
\draw[->] (-1.6,0)--(4,0) node[below left] {$x$};
\draw[->] (0,-0.4)--(0,2.5) node[below left] {$y$};
\draw (0,0) node [below left] {$O$};
\foreach \x in {-1,1,2,3}
\draw[thin] (\x,1pt)--(\x,-1pt) node [below] {$\x$};
\foreach \y in {1,2}
\draw[thin] (1pt,\y)--(-1pt,\y) node [below left] {$\y$};
\draw[samples=200,domain=0:3.5,smooth,variable=\x] plot (\x,{sqrt(\x)});
\path
(-1,0) coordinate (A)
(-1,1.73) coordinate (D)
(3,0) coordinate (B)
(3,1.73) coordinate (C)
;
\foreach \p/\q in {A/135,B/45,C/80,D/135}
\fill[black] (\p) circle (1.0pt) ($(\p)+(\q:3mm)$) node{$\p$};
\draw[pattern = north east lines, line width = 1.2pt,draw=none] (0,0)%
plot[domain=0:3] (\x,{sqrt(\x)})--(3,1.73) -- (3,0)--cycle;
\draw (A)--(D)--(C)--(B);
\end{tikzpicture}
\end{center}
Gọi $ABCD$ là hình chữ nhật $(H)$ với $AB$ nằm trên trục $Ox$.\\
Nhận thấy đồ thị hàm số $y=\sqrt{x}$ cắt trục hoành tại điểm có hoành độ bằng $0$ và đi qua $C(a;\sqrt{a}).$\\
Do đó nó chia hình chữ nhật $ABCD$ ra làm 2 phần có diện tích lần lượt là $S_1$, $S_2$.\\
Gọi $S_1$ là phần diện tích hình phẳng giới hạn bởi các đường $y=\sqrt{x}$, trục $Ox$, $x=0$, $x=a$ và $S_2$ là phần diện tích còn lại.\\
Ta có $S_1=\displaystyle\int\limits_{0}^{a} \sqrt{x} \mathrm{\,d}x$.\\
Đặt $t=\sqrt{x}\Rightarrow t^2=x\Rightarrow 2t\mathrm{\,d}t =\mathrm{\,d}x$.\\
Khi $x=0\Rightarrow t=0$, $x=a\Rightarrow t=\sqrt{a}$.\\
Do đó $S_1=\displaystyle\int\limits_{0}^{\sqrt{a}} 2t^2 \mathrm{\,d}t =\left(\dfrac{2}{3}t^3\right)\bigg |_0^{\sqrt{a}}=\dfrac{2a\sqrt{a}}{3}$.\\
Hình chữ nhật $ABCD$ có $AB=a+1$, $AD=\sqrt{a}$.\\
Khi đó $S_2=S_{ABCD}-S_1=\sqrt{a}(a+1)-\dfrac{2a\sqrt{a}}{3}=\dfrac{a\sqrt{a}}{3}+\sqrt{a}$.\\
Theo đề bài ta có $S_1=S_2\Leftrightarrow \dfrac{2a\sqrt{a}}{3}=\dfrac{a\sqrt{a}}{3}+\sqrt{a}\Leftrightarrow a\sqrt{a}=3\sqrt{a}\Leftrightarrow a=3$ (do $a>0$).
}
\end{ex}

\begin{ex}%[2D4V3-1]
\immini[thm]{Cho hàm số $y=f(x)$ có đạo hàm liên tục trên $[-2;3]$ và $f'(x)$ như hình vẽ. Biết $\displaystyle \int_{-2}^{1}f'(x)\mathrm{\,d}x=3$ và diện tích $S=\dfrac{5}{3}$. Giá trị $f(3)-f(-2)=\dfrac{a}{b}$. Tích $a\cdot b$ bằng bao nhiêu, biết $\dfrac{a}{b}$ là phân số tối giản?}
{
\begin{tikzpicture}[scale=1, font=\footnotesize, line join=round, line cap=round,>=stealth,x=0.6cm,y=0.6cm]
\def \xmin{-2.7};
\def \xmax{4.1};
\def \ymin{-1.5};
\def \ymax{3.5};
\draw[->] (\xmin, 0.) -- (\xmax,0.) node[anchor=north] {$x$};
\draw[->] (0.,\ymin) -- (0.,\ymax) node[anchor=west] {$y$};
\clip(\xmin+0.1,\ymin+0.1) rectangle (\xmax-0.1,\ymax-0.1);
\draw[smooth,samples=100,domain=\xmin:\xmax] plot(\x,{-0.01*(\x)^4+0.36*(\x)^3-0.54*(\x)^2-1.6*(\x)+1.8});
\fill[pattern=north west lines, opacity=0.5,smooth] (-1,0)--plot[domain=1:2.65](\x,{-0.01*(\x)^4+0.36*(\x)^3-0.54*(\x)^2-1.6*(\x)+1.8})--(2.65,0)--cycle;
\foreach \x/\y/\l/\g  in {0/0/O/-45, -1.95/0/-2/160, 1/0/1/30, 2.65/0/3/-45}{\fill[black] (\x,\y) circle(1pt)+(\g:0.4) node{$\l$};}
\coordinate (A) at (1,0);
\coordinate (B) at (2.65,0);
\draw[-latex](2,-0.4)--(2.5,-1.1);%2 mũi tên
\path (2.6,-1.2)node{$S$};%nhãn điểm
\end{tikzpicture}
}
\shortans[]{$12$}
\loigiai{
Diện tích $S=\displaystyle\int_1^3|f(x)| \mathrm{\,d} x=\dfrac{5}{3} \Rightarrow S=-\displaystyle\int_1^3 f(x) \mathrm{\,d} x=\dfrac{5}{3} \Rightarrow \displaystyle\int_1^3 f(x) \mathrm{\,d} x=-\dfrac{5}{3}$.\\
Ta có $\displaystyle\int_{-2}^3 f'(x) \mathrm{\,d} x=\displaystyle\int_{-2}^1 f'(x) \mathrm{\,d} x+\displaystyle\int_1^3 f'(x) \mathrm{\,d} x \Rightarrow f(3)-f(-2)=3-\dfrac{5}{3}=\dfrac{4}{3}=\dfrac{a}{b}$. \\
Vậy tích $a\cdot b=4\cdot 3=12$.
}
\end{ex}

\begin{ex}%[2D4V3-1]
\immini{Cho hình phẳng $(\mathrm{H})$ giới hạn bởi các đường $y=x^2, y=0, x=0, x=4$. Đường thẳng $y=k(0<k<16)$ chia hình $(\mathrm{H})$ thành hai phần có diện tích $S_1, S_2$ (hình vẽ). Tìm $k$ để $S_1=S_2$.}
{\begin{tikzpicture}
\draw[-> ] (-3,0) -- (4,0)node[below] {$x$};
\draw[-> ] (0,-1.5) -- (0,5)node[left] {$y$};
\draw(0,0) node[below right=-0.1] { $O$};
\draw[smooth,samples=100,domain=-2.25:2.25]  plot (\x,{(\x)^2});
\draw[fill=black](2,-1)--(2,5);
\draw[fill=black](-3,1)--(4,1);
\draw[name path = c1](1,1)--(2,1);
\draw[name path = c4](0,0)--(2,0);
\draw[name path = c5](1,1)--(2,1);
\draw[name path = c2 ,domain=1:2]  plot (\x,{(\x)^2});
\draw[name path = c3 ,domain=0:2]  plot (\x,{(\x)^2});
\tikzfillbetween[
of=c2 and c1,split,
every even segment/.style={fill=yellow}
] {fill = black};
\node at (3,1)[above]{ \footnotesize $ y=k $};
\node at (2,0)[below right]{ \footnotesize $ x=4 $};
\node at (1.75,2)[below]{ \footnotesize $ S_1$};
\node at (1,0.5)[right]{ \footnotesize $ S_2$};
\end{tikzpicture}}
\shortans[]{$4$}
\loigiai{
Ta có hình $(\mathrm{H})$ giới hạn bởi các đường $x=\sqrt{y}, x=4, y=0, y=16$, khi đó diện tích hình $(\mathrm{H})$ là $S=\displaystyle\int\limits_{0}^{16}\left(4-\sqrt{y}\right)\mathrm{\, d}y=\left(4y-\dfrac{2}{3}\sqrt{y^3}\right)\Bigr|_0^{16}=\dfrac{64}{3}$.\\
Gọi $(\mathrm{H}_1)$ là hình giới hạn bởi các đường $x=\sqrt{y}, x=4, y=0, y=k$, khi đó diện tích hình $(\mathrm{H}_1)$ là
$S_1=\displaystyle\int\limits_{0}^{k}\left(4-\sqrt{y}\right)\mathrm{\, d}y=\left(4y-\dfrac{2}{3}\sqrt{y^3}\right)\Bigr|_0^{k}=4k-\dfrac{2}{3}\sqrt{k^3}.$
\begin{eqnarray*}
& &S_1=S_2=\dfrac{S}{2} \Leftrightarrow 4k-\dfrac{2}{3}\sqrt{k^3}=\dfrac{32}{3}
\Leftrightarrow-\dfrac{2}{3}\left( \sqrt{k}\right) ^3+4\left( \sqrt{k}\right) ^2-\dfrac{32}{3}=0\\
&\Leftrightarrow&\hoac{&\sqrt{k}=2+2\sqrt{3}\\&\sqrt{k}=2-2\sqrt{3}<0\,(\text{loại})\\&\sqrt{k}=2}\Leftrightarrow \hoac{&k=16+8\sqrt{3}\\&k=4.}
\end{eqnarray*}
Kết hợp với điều kiện $0<k<16$ ta được $k=4$.
}
\end{ex}

\begin{ex}%[2D4V3-1]
\immini{Cho hàm số bậc hai $y=f(x)$ có đồ thị $(P)$ và nửa đường tròn $(C)$ cắt $(P)$ tại hai điểm như trong hình bên. Biết rằng hình phẳng giới hạn bởi $(P)$ và $(C)$ có diện tích $S=\dfrac{\pi}{m}+\dfrac{1}{n}$ ($m,n$ nguyên dương). Giá trị biểu thức $P=m^2+n^2$ bằng bao nhiêu?}
{\begin{tikzpicture}[scale=1, font=\footnotesize, line join=round, line cap=round,>=stealth]
\def\a{1} \def\b{-2} \def\c{1} % Hệ số
\def\xmin{-1} \def\xmax{3}
\def\ymin{-1} \def\ymax{3}
%		\draw[color=gray!50,dashed] (\xmin,\ymin) grid (\xmax,\ymax);
\draw[->] (\xmin,0)--(\xmax,0) node [above]{$x$};
\draw[->] (0,\ymin)--(0,\ymax) node [left]{$y$};
\node at (0,0) [below right]{$O$};
\draw[smooth,samples=300,domain=-0.6:2.7] plot(\x,{\a*(\x)^2+\b*(\x)+\c});
\draw[smooth,samples=300,domain=-0.414213:2.414213] plot(\x,{sqrt(-1*(\x)^2+2*(\x)+1)});
\fill[pattern=north west lines,smooth] plot[domain=0:2](\x,{sqrt(-1*(\x)^2+2*(\x)+1)})--plot[domain=2:0](\x,{\a*(\x)^2+\b*(\x)+\c})--cycle;
\path (-0.414,0) coordinate (A);
\draw (A) arc (180:0:1.414);
\node at (-0.414,0) [below left]{$1-\sqrt{2}$};
\node at (1,0) [below]{$1$};
\node at (2.414,0) [below]{$1+\sqrt{2}$};
\node at (0,1) [left]{$1$};
\foreach\x/\y in{-0.414/0,1/0,2.414/0,0/1,0/0}
\fill(\x,\y) circle(1.0pt);
\end{tikzpicture}}
\shortans[]{$13$}
\loigiai{
Từ đồ thị ta có	$(P)\colon y=x^2-2x+1$ và đường tròn $(C)$ có phương trình \[(x-1)^2+y^2=2\Rightarrow y=\sqrt{2-(x-1)^2}.\]
Phương trình hoành độ giao điểm của $(P)$ và $(C)$ là \[\sqrt{2-(x-1)^2}=x^2-2x+1\Rightarrow\hoac{&x=0\\&x=2.}\]
Hình phẳng giới hạn bởi $(P)$ và $(C)$ có diện tích
\begin{eqnarray*}
S&=&\displaystyle\int\limits_0^2\left[\sqrt{2-(x-1)^2}-\left(x^2-2x+1\right)\right]\mathrm{\,d}x\\
&=&\displaystyle\int\limits_0^2\sqrt{2-(x-1)^2}\mathrm{\,d}x-\displaystyle\int\limits_0^2\left(x^2-2x+1\right)\mathrm{\,d}x\\
&=&I-\dfrac{2}{3}
\end{eqnarray*}
Với $I=\displaystyle\int\limits_0^2\sqrt{2-(x-1)^2}\mathrm{\,d}x$, đặt $x-1=\sqrt{2}\sin t\Rightarrow\mathrm{\,d}x=\sqrt{2}\cos t\mathrm{\,d}t$.
\begin{eqnarray*}
I&=&\displaystyle\int\limits_{\tfrac{-\pi}{4}}^{\tfrac{\pi}{4}}\sqrt{2-2\sin^2t}\cdot\sqrt{2}\cos t\mathrm{\,d}t\\
&=&\displaystyle\int\limits_{\tfrac{-\pi}{4}}^{\tfrac{\pi}{4}} 2\cos^2t\mathrm{\,d}t\\
&=&\displaystyle\int\limits_{\tfrac{-\pi}{4}}^{\tfrac{\pi}{4}} (1+\cos 2t)\mathrm{\,d}t\\
&=&\left(t+\dfrac{1}{2}\sin 2t\right)\bigg|_{\tfrac{-\pi}{4}}^{\tfrac{\pi}{4}}\\
&=&\dfrac{\pi}{2}+1.
\end{eqnarray*}
Suy ra $S=\dfrac{\pi}{2}+1-\dfrac{2}{3}=\dfrac{\pi}{2}+\dfrac{1}{3}$.\\
Vậy $P=m^2+n^2=13$.
}
\end{ex}

\begin{ex}%[2D4V3-1]
\immini{
Gọi $H$ là hình phẳng giới hạn bởi đồ thị hàm số $y=-x^2+4x$ và trục hoành. Hai đường thẳng $y=m$ và $y=n$ chia $(H)$ thành ba phần có diện tích bằng nhau (tham khảo hình vẽ). Giá trị của biểu thức $T=(4-m)^3+(4-n)^3$ bằng bao nhiêu? (Kết quả làm tròn đến hàng phần mười)
}{
\begin{tikzpicture}[thick,>=stealth,x=1cm,y=1cm,scale=.8]
\draw[->] (-1,0) -- (5,0) node[below] {\small $x$};
\draw[->] (0,-1) -- (0,5) node[right] {\small $y$};
\draw [fill=white,draw=black] (0,0) circle (1pt)node[above left] {\footnotesize $O$};
\clip(-1,-1) rectangle (5,5);
\draw[thick,smooth,samples=100,domain=-1:5] plot(\x,{-(\x)^2+4*(\x)});
\draw (-1,1.2)--(5.1,1.3)node[above left]{$y=n$};
\draw (-1,2.3)--(5.1,2.4)node[above left]{$y=m$};
\end{tikzpicture}}
\shortans[]{$35{,}6$}
\loigiai{
\immini{
Gọi $S$ là diện tích hình phẳng giới hạn bởi đồ thị hàm số $y=-x^2+4x$ và trục $Ox$ và hai đường thẳng $x=0$, $x=2$.\\
Khi đó $S=\displaystyle\int\limits_{0}^{2} (-x^2+4x)\mathrm{\,d}x =\dfrac{16}{3}$.\\
Đường thẳng $y=m$ và $y=n$ chia $S$ thành ba phần bằng nhau có diện tích theo thứ tự từ trên xuống là $S_1$; $S_2$; $S_3$.\\
Gọi hoành độ các giao điểm của parabol với hai đường thẳng như hình bên.\\
Ta có
\begin{eqnarray*}
&& S_1=2\displaystyle\int\limits_{a}^{2} (-x^2+4x-m)\mathrm{\,d}x =\dfrac{1}{3}S\\
&\Leftrightarrow & \left(-\dfrac{x^3}{3}+2x^2-mx\right)\Big|_{a}^{2}=\dfrac{1}{3}\cdot\dfrac{16}{3}\\
&\Leftrightarrow & \left(\dfrac{16}{3}-2m\right)-\left(-\dfrac{a^3}{3}+2a^2-ma\right)=\dfrac{16}{9}\quad(1).
\end{eqnarray*}
Mà $x=a$ là nghiệm của phương trình $-x^2+4x=m$ nên ta có $-a^2+4a=m\quad(2)$.\\
Thay $(2)$ vào $(1)$ ta được $-\dfrac{2a^3}{3}+4a^2-8a+\dfrac{32}{9}=0\Leftrightarrow a\approx 0{,}613277$.\\
Suy ra $m=-a^2+4a\approx 2{,}077$.\\
Tương tự ta có
\begin{eqnarray*}
&& S_1+S_2=\dfrac{2}{3}S\\
&\Rightarrow & 2\displaystyle\int\limits_{b}^{2} (-x^2+4x-n)\mathrm{\,d}x =\dfrac{2}{3}\cdot 2\cdot\displaystyle\int\limits_{0}^{2} (-x^2+4x)\mathrm{\,d}x\\
&\Leftrightarrow & -\dfrac{2}{3}b^3+4b^2-8b+\dfrac{16}{9}=0\\
&\Leftrightarrow & b\approx 0{,}252839\Rightarrow n=-b^2+4b=0{,}947428.
\end{eqnarray*}
Khi đó $T=(4-m)^3+(4-n)^3=\dfrac{320}{9}\approx35{,}6$.
}{
\begin{tikzpicture}[thick,>=stealth,x=1cm,y=1cm,scale=.8]
\draw[->] (-1,0) -- (5,0) node[below] {\small $x$};
\draw[->] (0,-1) -- (0,5) node[right] {\small $y$};
\draw [fill=white,draw=black] (0,0) circle (1pt)node[above left] {\footnotesize $O$};
\clip(-1,-1) rectangle (5,5);
\draw[thick,smooth,samples=100,domain=-1:5] plot(\x,{-(\x)^2+4*(\x)});
\draw (-1,1.2)--(5,1.2)node[above left]{$y=n$};
\draw (-1,2.3)--(5,2.3)node[above left]{$y=m$};
\draw[dashed](0.33,0)node[below]{$b$}--(0.33,1.2) (0.7,0)node[below]{$a$}--(0.7,2.3) (2,0)node[below]{$2$}--(2,4);
\draw[dashed] (0,4) node[below left]{$4$}--(2,4);
\end{tikzpicture}
}
}
\end{ex}

\begin{ex}%[2D4V3-1]
\immini{
Cho hàm số bậc ba $y=f(x)$ có đồ thị là đường cong trong hình bên dưới
Biết hàm số $f(x)$ đạt cực trị tại các điểm $x_1$, $x_2$ sao cho $x_2=x_1+2$ và $f'' (2)=0$. Gọi $S_1$ và $S_2$ là diện tích hai hình phẳng được gạch trong hình bên. Tỉ số $\dfrac{S_1}{S_2}$ bằng bao nhiêu? (Kết quả viết dưới dạng số thập phân làm tròn tới hàng phần trăm).}
{
\begin{tikzpicture}[scale=1, font=\footnotesize, line join=round, line cap=round, >=stealth]
\draw[->] (-0.5,0) --(4.5,0)node[below left] {$x$};
\draw[->] (0,-2.25) -- (0,2.5)node[right] {$y$};
\draw [fill=black] (0,0) circle(1pt) node[below left] {\footnotesize $O$};
\begin{scope}
\clip (-1,-2) rectangle (4.5,2.2);
\draw[smooth,black,thick,samples=100,domain=-1:4.5] plot(\x,{(\x)^3-6*(\x)^2+9*(\x)-2});
\filldraw [pattern=north west lines, opacity=0.5]
plot[domain=0:1] (\x,{(\x)^3-6*(\x)^2+9*(\x)-2})--(1,0)--(0,0) ;
\end{scope}
\draw [fill=black] (0,-2) circle(1pt)
(1,2) circle(1pt) (3,-2) circle(1pt)
(1,0) circle(1pt) node[below] {\footnotesize $x_1$}
(3,0) circle(1pt) node[above] {\footnotesize $x_2$}
;
\draw [dashed] (3,0)--(3,-2);
\draw[-latex](.4,-.55)--(0.05,-.15);\draw[-latex](1.2,.7)--(.68,.7);%2 mũi tên
\path (.55,-.55)node{$S_1$}(1.35,.7)node{$S_2$};%nhãn điểm
\end{tikzpicture}
}
\shortans[]{$0{,}25$}
\loigiai{
Đặt $f(x)=ax^3+bx^2+cx+d,\,(a > 0) \Rightarrow f'(x)=3ax^2+2bx+c\Rightarrow f'' (x)=6ax+2b$.\\
Ta có $f'' (x)=0\Leftrightarrow 6a x+2b=0\Leftrightarrow x=\dfrac{-b}{3a}$.\\
Vì $f'' (2)=0$ nên $\dfrac{-b}{3a}=2\Leftrightarrow b=-6a$.\\
Mặt khác, theo định lý Vi-ét thì $x_1+x_2=\dfrac{-2b}{3a}=4$, kết hợp với $x_2-x_1=2$ ta suy ra $x_1=1$, $x_2=3$.
Do đó $x_1 \cdot x_2=\dfrac{c}{3a}=3\Leftrightarrow c=9a$.\\
Từ đó ta có $f(x)=a x^3-6a x^2+9a x+d$.\\
Từ đồ thị hàm số ta suy ra $f(2)=0\Leftrightarrow 8a-24a+18a+d=0\Leftrightarrow d=-2a$.\\
Suy ra $f(x)=a\left(x^3-6x^2+9x-2\right)$.\\
Xét phương trình \[f(x)=0\Leftrightarrow x^3-6x^2+9x-2=0\Leftrightarrow\hoac{&x=2\\&x=2-\sqrt{3} \\&x=2+\sqrt{3}.}\]
Từ đây ta tính được	\[S_1=-a \displaystyle\int\limits_0^{2-\sqrt{3}}\left(x^3-6x^2+9x-2\right) \mathrm{d} x=\dfrac{a}{4}, \quad S_2=a \displaystyle\int\limits_{2-\sqrt{3}}^1\left(x^3-6x^2+9x-2\right)\mathrm{d}x=a.\]
Vậy $\dfrac{S_1}{S_2}=\dfrac{1}{4}=0{,}25$.
}
\end{ex}

\begin{ex}%[2D4V3-1]
Tính diện tích hình phẳng giới hạn bời đồ thị hàm số $y=x^2+x-1$, $y=x^4+x-1$, $x=-1$, $x=1$.
\shortans{$0,27$}
\loigiai{
Diện tích hình phẳng giới hạn bởi đồ thị hàm số $y=x^2+x-1, y=x^4+x-1, x=-1, x=1$ là
\begin{eqnarray*}
S&=&\displaystyle\int\limits_{-1}^1\left|x^2-x^4\right| \mathrm{d}x\\
&=&\displaystyle\int\limits_{-1}^0\left|x^2-x^4\right| \mathrm{d}x+\displaystyle\int\limits_0^1\left|x^2-x^4\right| \mathrm{d}x\\
&=&\left|\displaystyle\int\limits_{-1}^0\left(x^2-x^4\right) \mathrm{d}x\right|+\left|\displaystyle\int\limits_0^1\left(x^2-x^4\right) \mathrm{d}x\right|\\
&=&\left|\left(\dfrac{x^3}{3}-\dfrac{x^5}{5}\right)\right| 0|+|\left(\dfrac{x^3}{3}-\dfrac{x^5}{5}\right)|0|\\
&=&\dfrac{2}{15}+\dfrac{2}{15}=\dfrac{4}{15}\approx0,27.
\end{eqnarray*}
}
\end{ex}

\begin{ex}%[2D4V3-1]
Cho hình thang cong $(H)$ giới hạn bởi các đường $y=\mathrm{e}^{x}$, $y=0$, $x=0$, $x=\ln 4$. Đường thẳng $x=k$, $(0<k<\ln 4)$ chia $(H)$ thành hai phần có diện tích là $S_1$ và $S_2$ như hình vẽ bên. Tìm $k$ để $S_1=2 S_2$ (làm tròn kết quả đến hàng phần chục).
\begin{center}
\begin{tikzpicture}[font=\footnotesize, line join=round, line cap=round, >=stealth, scale =1]
\draw[->] (-1,0) --(0,0) node[below left]{$O$}--(2.5,0) node[below]{$x$};
\draw[->] (0,-.7) --(0,4.6) node[right]{$y$};
\draw[fill = black] (.8,0) node[below left]{$k$} circle (1pt);
\draw[fill = black] (1.4,0) node[below right]{$\ln 4$} circle (1pt);
\draw[fill = black] (0,1) node[above left]{$1$} circle (1pt);
\draw[fill = black] (0,0) circle (1pt);
\draw (.8,-.7)--(.8,4.4) (1.4,-.7)--(1.4,4.4);
\draw [samples=100, domain=-.9:1.5] plot (\x, {e^(\x)});
\fill[color=gray!10!black,shading=axis,opacity=0.2] (0,0) -- plot[smooth,samples=100,domain=0:1.4] (\x, {e^(\x)}) -- (1.4,0) -- cycle;
\draw (0.2,.3) node[right]{$ S_1 $} (1,1) node[above]{$ S_2 $};
\end{tikzpicture}
\end{center}
\shortans{$1,1$}
\loigiai{
Diện tích hình thang cong $(H)$ giới hạn bởi các đường $y=\mathrm{e}^{x}$, $y=0$, $x=0$, $x=\ln 4$ là
\[S=\displaystyle\int\limits_{0}^{\ln 4} \mathrm{e}^{x} \mathrm{\,d} x = \mathrm{e}^{x}\Bigg|_{0}^{\ln 4}=\mathrm{e}^{\ln 4}-\mathrm{e}^{0}=4-1=3.\]\\
Ta có $S=S_1+S_2=S_1+\dfrac{1}{2} S_1=\dfrac{3}{2} S_1$. Suy ra $S_1=\dfrac{2 S}{3}=\dfrac{2 \cdot 3}{3}=2$.\\
Vì $S_1$ là phần diện tích được giới hạn bởi các đường $y=\mathrm{e}^{x}$, $y=0$, $x=0$, $x=k$ nên\\
$
2=S_1=\displaystyle\int\limits_{0}^{k} \mathrm{e}^{x} \mathrm{\,d} x=\mathrm{e}^{x}\Bigg|_{0}^{k}=\mathrm{e}^{k}-\mathrm{e}^{0}=\mathrm{e}^{k}-1$.\\
Do đó $\mathrm{e}^{k}=3 \Leftrightarrow k=\ln 3 \approx 1{,}1$.}
\end{ex}

\begin{ex}%[2D4V3-1]
Cho hình phẳng $(H)$ giới hạn bởi các đường $y=\left|x^{2}-1\right|$ và $y=k$, với $0<k<1$. Tìm $k$ để diện tích hình phẳng $(H)$ gấp hai lần diện tích hình phẳng được kẻ sọc ở hình vẽ bên (làm tròn kết quả đến hàng phần trăm).
\begin{center}
\begin{tikzpicture}[font=\footnotesize, line join=round, line cap=round, >=stealth, scale =1]
\draw[->] (-2,0) --(0,0) node[below left]{$O$}--(2,0) node[below]{$x$};
\draw[->] (0,-3) --(0,3.3) node[right]{$y$};
\draw[fill = black] (1,0) node[below left]{$1$} circle (1pt);
\draw[fill = black] (0,1) node[above left]{$1$} circle (1pt);
\draw[fill = black] (0,0) circle (1pt);
\draw (-2,.4)--(2,.4) node[above right]{$ y=k $};
\draw [samples=100, domain=-1:1] plot (\x, {-(\x)^2+1});
\draw [samples=100, domain=-1:-1.8] plot (\x, {(\x)^2-1});
\draw [samples=100, domain=1:1.8] plot (\x, {(\x)^2-1});
\draw[dashed,samples=100,domain=-1:-1.8] plot (\x, {-(\x)^2+1});
\draw[dashed,samples=100,domain=1:1.8] plot (\x, {-(\x)^2+1});
\fill[pattern=north west lines] (-.77,.4) -- plot[smooth,samples=100,domain=-.77:.77] (\x, {-(\x)^2+1}) -- (.77,.4) -- cycle;

\end{tikzpicture}
\end{center}
\shortans{$ 0{,}59 $}
\loigiai{\begin{center}
\begin{tikzpicture}[font=\footnotesize, line join=round, line cap=round, >=stealth, scale =1]
\draw[->] (-2,0) --(0,0) node[below left]{$O$}--(2,0) node[below]{$x$};
\draw[->] (0,-3) --(0,3.3) node[right]{$y$};
\draw[fill = black] (1,0) node[below left]{$1$} circle (1pt);
\draw[fill = black] (0,1) node[above left]{$1$} circle (1pt);
\draw[fill = black] (0,0) circle (1pt);
\draw (-2,.4)--(2,.4) node[above right]{$ y=k $};
\draw [samples=100, domain=-1:1] plot (\x, {-(\x)^2+1});
\draw [samples=100, domain=-1:-1.8] plot (\x, {(\x)^2-1});
\draw [samples=100, domain=1:1.8] plot (\x, {(\x)^2-1});
\draw[dashed,samples=100,domain=-1:-1.8] plot (\x, {-(\x)^2+1});
\draw[dashed,samples=100,domain=1:1.8] plot (\x, {-(\x)^2+1});
\fill[pattern=dots] (0,.4) -- plot[smooth,samples=100,domain=0:.77] (\x, {-(\x)^2+1}) -- (.77,.4) -- cycle;
\fill[pattern=north west lines] (.77,.4)--plot[smooth,samples=100,domain=.77:1] (\x, {-(\x)^2+1}) -- plot[smooth,samples=100,domain=1:1.18] (\x, {(\x)^2-1}) -- (1.18,.4) -- cycle;
\draw (0.77,.4) node[above]{$ A $} circle (1pt) (1.18,.4) node[above right]{$ B $} circle (1pt);
\end{tikzpicture}
\end{center}
Gọi $S$ là diện tích hình phẳng $(H)$. Lúc dó $S=2 S_1+2 S_2$, trong đó $S_1$ là diện tích phần chấm bi và $S_2$ là diện tích phần gạch sọc trong hình vẽ bên.\\
Gọi $A$, $ B$ là các giao điếm có hoành độ dương của đường thẳng $y=k$ và đồ thị hàm số $y=\left|x^{2}-1\right|$, trong đó $A(\sqrt{1-k}; k)$ và $B(\sqrt{1+k}; k)$.\\
Theo yêu cầu bài toán
\begin{eqnarray*}
& & S=2 \cdot 2 S_1 \\
& \Leftrightarrow & S_1=S_2 \\
& \Leftrightarrow &\displaystyle\int\limits_{0}^{\sqrt{1-k}} \left(1-x^{2}-k\right) \mathrm{\,d} x =\displaystyle\int\limits_{\sqrt{1-k}}^{1}\left(k-1+x^{2}\right) \mathrm{\,d} x+\displaystyle\int\limits_{1}^{\sqrt{1+k}}\left(k-x^{2}+1\right) \mathrm{\,d} x \\
& \Leftrightarrow & (1-k) \sqrt{1-k}-\dfrac{1}{3}(1-k) \sqrt{1-k}=\dfrac{1}{3}-(1-k)-\dfrac{1}{3}(1-k) \sqrt{1-k}+(1-k) \sqrt{1-k}+(1+k) \sqrt{1+k}-\dfrac{1}{3}(1+k) \sqrt{1+k}-(1+k)+\dfrac{1}{3} \\
& \Leftrightarrow &\dfrac{2}{3}(1+k) \sqrt{1+k}=\dfrac{4}{3} \\
& \Leftrightarrow & (\sqrt{1+k})^{3}=2\\
& \Leftrightarrow & k=\sqrt[3]{4}-1\approx 0{,}59.
\end{eqnarray*}}

\end{ex}

\begin{ex}%[2D4V3-1]
Kí hiệu $S(t)$ là diện tích của hình phẳng giới hạn bởi các đường $y=2 x+1$, $y=0$, $x=1$, $x=t\left(t>1\right)$. Tìm $t$ để $S(t)=10$.
\shortans{$3$}
\loigiai{
\textbf{Cách 1.} Ta có $S(t)=\displaystyle\int\limits_1^t|2 x+1| \mathrm{d}x=\displaystyle\int\limits_1^t(2 x+1) \mathrm{d}x$.\\
Suy ra $S(t)=\left.\left(x^2+x\right)\right|_1^t=t^2+t-2$.\\
Do đó $S(t)=10 \Leftrightarrow t^2+t-2=10 \Leftrightarrow t^2+t-12=0 \Leftrightarrow\hoac{&t=3 \\ &t=-4\text{ (L)}}.$\\
Vậy $t=3$.\\
\textbf{Cách 2}. Hình phẳng đã cho là hình thang có đáy nhỏ bằng $y(1)=3$, đáy lớn bằng $y(t)=2 t+1$ và chiều cao bằng $t-1$.
Ta có \[\dfrac{(3+2t+1)(t-1)}{2}=10 \Leftrightarrow 2 t^2+2 t-24=0 \Leftrightarrow\hoac{t=3 \\ t=-4}.\]\\
Vì $t>1$ nên $t=3$.
}
\end{ex}

\begin{ex}%[2D4V3-1]
Gọi $S$ là diện tích hình phẳng giới hạn bởi các đường $m y=x^2$, $m x=y^2(m>0)$. Tìm giá trị của $m$ để $S=3$.

\shortans{$3$}
\loigiai{
Tọa độ giao điểm của hai đồ thị hàm số là nghiệm của hệ phương trình $\heva{my=x^2 &\quad(1)\\mx=y^2 &\quad(2)}$
Thế $(1)$ vào $(2)$ ta được $m x=\left(\dfrac{x^2}{m}\right)^2 \Leftrightarrow m^3 x-x^4=0 \Leftrightarrow\hoac{&x=0\\&x=m>0.}$\\
Vì $y=\dfrac{x^2}{m}>0$ nên $m x=y^2 \text{ (với $y>0$) }  \Leftrightarrow y=\sqrt{m x}$\\
Khi đó diện tích hình phẳng cần tìm là
\begin{eqnarray*}
S&=&\displaystyle\int\limits_0^m\left|\sqrt{m x}-\dfrac{x^2}{m}\right| \mathrm{d}x=\left|\displaystyle\int\limits_0^m\left(\sqrt{m x}-\dfrac{x^2}{m}\right) \mathrm{d}x\right|\\
&=&\left|\left(\dfrac{2 \sqrt{m}}{3} \cdot x^{\dfrac{3}{2}}-\dfrac{x^3}{3 m}\right)\right|_0^m\\
&=& \left|\dfrac{1}{3} m^2 \right|=\dfrac{1}{3} m^2
\end{eqnarray*}
Yêu cầu bài toán $S=3 \Leftrightarrow \dfrac{1}{3} m^2=3 \Leftrightarrow m^2=9 \Leftrightarrow m=3$.

}
\end{ex}

\begin{ex}%[2D4V3-1]
Giá trị dương của tham số $m$ sao cho diện tích hình phẳng giới hạn bởi đồ thị của hàm số $y=2 x+3$ và các đường thẳng $y=0$, $x=0$, $x=m$ bằng 10 là?

\shortans{$2$}
\loigiai{
Vì $m>0$ nên $2 x+3>0, \forall x \in[0;m]$.
Diện tích hình phẳng giới hạn bởi đồ thị hàm số $y=2 x+3$ và các đường thẳng $y=0$, $x=0$, $x=m$ là
\[S=\displaystyle\int\limits_0^m(2 x+3) \cdot \mathrm{d}x=\left.\left(x^2+3 x\right)\right|_0^m=m^2+3m.\]
Theo giả thiết ta có
\[S=10 \Leftrightarrow m^2+3 m=10 \Leftrightarrow m^2+3 m-10=0 \Leftrightarrow \hoac{&m=2\\&m=-5}
\Leftrightarrow m=2 \text { do } m>0.\]
}
\end{ex}

\begin{ex}%[2D4V3-1]
Cho hàm số  $f(x)=\heva{&7-4x^3\text{ khi }  0 \leq x \leq 1\\&4-x^2 \text{ khi } x>1}$. Tính diện tích hình phẳng giới hạn bởi đồ thị hàm số $f(x)$ và các đường thẳng $x=0$, $x=3$, $y=0$.
\shortans{10}
\loigiai{

\begin{center}
\begin{tikzpicture}[scale=0.6, font=\footnotesize, line join=round, line cap=round, >=stealth]
% Vẽ trục
\draw[->] (-0.5,0) -- (4,0) node[right] {$x$};
\draw[->] (0,-5) -- (0,7) node[above] {$y$};

% Đánh dấu các điểm trên trục x
\foreach \x in {1,2,3}{
\draw[fill=black] (\x,0) circle(0.03) node[below]{$\x$};
}

% Đánh dấu các điểm trên trục y
\foreach \y in {-5,-4,-3,-2,-1,1,2,3,4,5,6,7}{
\draw[fill=black] (0,\y) circle(0.03) node[left]{$\y$};
}
% Đánh dấu gốc tọa độ
\draw[fill=black] (0,0) circle(0.03) node[below left] {$0$};



% Draw function curve
\draw[thick,domain=0:1,samples=100] plot (\x,{-4*(\x)*(\x)+7});
\draw[thick,domain=1:3,samples=100] plot (\x,{-(\x)*(\x)+4});
% Draw lines
\draw (1,0) -- (1,3);
\draw (3,0) -- (3,-5);
\draw[dashed] (0,3) -- (1,3);
\draw[dashed] (0,-5) -- (3,-5);
% Draw shaded area
\fill[pattern=north east lines, pattern color=black!50]
(0,0) -- plot[domain=0:1,samples=100] (\x,{-4*(\x)*(\x)+7}) -- (1,0) -- cycle;
\fill[pattern=north east lines, pattern color=black!50]
(1,0) -- plot[domain=1:2,samples=100] (\x,{-(\x)*(\x)+4}) -- (2,0) -- cycle;
\fill[pattern=north east lines, pattern color=black!50]
(2,0) -- plot[domain=2:3,samples=100] (\x,{-(\x)*(\x)+4}) -- (3,0) -- cycle;
\end{tikzpicture}
\end{center}

\begin{eqnarray*}
S&=&\displaystyle\int\limits_0^1\left(7-4 x^3\right) \mathrm{d}x+\displaystyle\int\limits_1^2\left(4-x^2\right) \mathrm{d}x+\displaystyle\int\limits_2^3\left(x^2-4\right) \mathrm{d}x \\ & =&\left.\left(7 x-x^4\right)\right|_0 ^1+\left.\left(4 x-\dfrac{x^3}{3}\right)\right|_1 ^2+\left.\left(\dfrac{x^3}{3}-4 x\right)\right|_2 ^3
\\&=&6+4-\dfrac{7}{3}-3-\dfrac{8}{3}+8=10 .
\end{eqnarray*}
}
\end{ex}

\begin{ex} %[2D4V3-1]
Cho $(H)$ là hình phẳng được giới hạn bởi các đường có phương trình $y=\frac{10}{3}x-{x^2}$, $y=\heva{&-x\, \, &\text{khi} \, x\le 1 \\ &x-2\, \, &\text{khi} \, x>1 }$. Diện tích của $( H )$ bằng bao nhiêu.
\shortans[1]{$6{,}5$}
\loigiai{
Xét $y=\heva{&-x\, \, &\text{khi} \, x\le 1 \\ &x-2\, \, &\text{khi} \, x>1 }$ \\
$\underset{x\to {1^+}}{\mathop{\lim\limits_{n \to +\infty}}}y=\underset{x\to {1^+}}{\mathop{\lim\limits_{n \to +\infty}}}( x-2 )=-1$ và $y(1)=-1
\Rightarrow \underset{x\to {1^+}}{\mathop{\lim\limits_{n \to +\infty}}}y=y( 1 )$.\\
Do đó hàm số liên tục tại $ x=1$.\\
Xét phương trình hoành độ giao điểm $\dfrac{10}{3}x-x^2=-x$ khi $x\le 1$ có nghiệm $x=0$.\\
Xét phương trình hoành độ giao điểm $\dfrac{10}{3}x-x^2=x-2$ khi $x>1$ có nghiệm $x=3$.\\
Diện tích hình phẳng cần tính là\\
\begin{eqnarray*}
S&=&\int\limits_0^1{( \frac{10}{3}x-{x^2}+x )\mathrm{\,d}x}+\int\limits_1^3{( \frac{10}{3}x-{x^2}-x+2 )\mathrm{\,d}x}.\\
&=&\int\limits_0^1{( \frac{13}{3}x-{x^2} )\mathrm{\,d}x}+\int\limits_1^3{( \frac{7}{3}x-{x^2}+2 )\mathrm{\,d}x}\\
&=&\left( \frac{13}{6}{x^2}-\frac{{x^3}}{3}\right) \bigg|_0^1+\left( \frac{7}{6}{x^2}-\frac{{x^3}}{3}+2x\right) \bigg|_1^3\\
&=&\frac{13}{2}=6{,}5.
\end{eqnarray*}}
\end{ex}

\begin{ex}%[Dự án 2025 - đề cấu trúc mới, Hung Doan]%[2D4V3-1]
Cho hàm số $y=f(x)$. Hàm số có đồ thị hàm số $y=f'(x)$ như hình vẽ dưới đây.
\begin{center}
\begin{tikzpicture}[scale=1, font=\footnotesize, line join=round, line cap=round, >=stealth]
\def\xt{-2.75} \def\xp{4.75} \def\yt{2.5} \def\yd{-3}
\draw[->](\xt, 0)--(\xp, 0) node[below]{$x$};
\draw[->](0,\yd)--(0,\yt) node[left]{$y$};
\node at (0, 0) [below left]{$O$};
\node at (-1.5, 1.4) {$y=f'(x)$};
\foreach \x in {1,4}
\draw[shift ={ (\x,0)}]node[above]{$\x$} (0pt,2pt) --(0pt,-2pt);
\draw[shift ={ (-2,0)}]node[above right]{$-2$} (0pt,2pt) --(0pt,-2pt);
\clip(\xt+0.1,\yd+0.1) rectangle(4,\yt-0.1);
\draw plot[smooth,tension=0.7] coordinates{(-2.4,2) (-1.2,-1.8) (1,0) (3,-2.5) (4,0)};
\end{tikzpicture}
\end{center}
Biết diện tích hình phẳng giới hạn bởi trục $Ox$ và đồ thị hàm số $y=f'(x)$ trên đoạn $[-2;1]$ và $[1;4]$ lần lượt bằng $9$ và $12$. Cho biết $f(1)=3$. Tính giá trị biểu thức $P=f(-2)+f(4)$.
\shortans{$3$}
\loigiai{
Ta có $\displaystyle\int \limits_{-2}^1 |f'(x)|\mathrm{d}x=9\Leftrightarrow \displaystyle\int \limits_{-2}^1 f'(x)\mathrm{d}x=-9\Rightarrow f(1)-f(-2)=-9$.\\
Mà $f(1)=3 \Rightarrow f(-2)=12$.\\
Ta có $\displaystyle\int \limits_1^4 |f'(x)|\mathrm{d}x=12\Leftrightarrow \displaystyle\int \limits_1^4 f'(x)\mathrm{d}x=-12\Rightarrow f(4)-f(1)=-12$.\\
Mà $f(1)=3 \Rightarrow f(4)=-9$.\\
Vậy $P=f(-2)+f(4)=3$.
}
\end{ex}

\begin{ex}%[Vovanle]%[2D4V3-1]
\immini{Cho hàm số $y=ax^4+bx^2+c$ có đồ thị $(C)$, biết rằng $(C)$ đi qua điểm $A(-1;0)$, tiếp tuyến $d$ tại $A$ của $(C)$, cắt $(C)$ tại hai điểm có hoành độ lần lượt là $0$ và $2$. Diện tích hình phẳng giới hạn bởi $d$, đồ thị $(C)$ và hai đường thẳng $x=0$; $x=2$ có diện tích bằng $\dfrac{28}{5}$ (phần gạch sọc trong hình vẽ).

Tính diện tích hình phẳng giới hạn bởi $(C)$, trục hoành và hai đường thẳng $x=-1$; $x=0$.
}{
\begin{tikzpicture}[line join=round,line cap=round, font=\footnotesize,scale=0.5,>=stealth]
\draw[-stealth](-2,0)--(2.5,0)node[below]{$x$};
\draw[-stealth](0,-0.7)--(0,7)node[right]{$y$};
\fill[pattern=north east lines]plot[domain=0:2](\x,{(\x)^4-3*(\x)^2+2})--cycle;
\draw[smooth,samples=300,domain=-2.02:2.02] plot(\x,{(\x)^4-3*(\x)^2+2});
\draw[smooth,samples=300,domain=-1.5:2.3] plot(\x,{2*(\x+1)});
\fill (0,0) circle(1pt)node[below right]{$O$}(-1,0)circle(1pt)+(0.1,0) node[below]{$-1$}(2,0)circle(1pt)node[below]{$2$};
\draw[dashed] (2,0)--(2,6);
\end{tikzpicture}
}
\shortans{$0{,}2$}
\loigiai{
Ta có $y'=4ax^3+2bx$ $\Rightarrow d\colon y=\left(-4a-2b\right)\left(x+1\right)$.
Phương trình hoành độ giao điểm của $d$ và $(C)$ là $\left(-4a-2b\right)\left(x+1\right)=ax^4+bx^2+c.\hfill(1)$\\
Phương trình $(1)$ phải cho $2$ nghiệm là $x=0$, $x=2$.
\[\Rightarrow\heva{&-4a-2b=c\\&-12a-6b=16a+4b+c}
\Leftrightarrow \heva{&-4a-2b-c=0&(2)\\&28a+10b+c=0&(3).}\]
Mặt khác, diện tích phần gạch sọc là
\allowdisplaybreaks
\begin{eqnarray*}
&&\dfrac{28}{5}=\displaystyle\int\limits_0^2{\left[\left(-4a-2b \right)\left( x+1 \right)-ax^4-bx^2-c\right]\mathrm{\,d}x}\\
&\Leftrightarrow& \dfrac{28}{5}=4\left(-4a-2b\right)-\dfrac{32}{5}a-\dfrac{8}3b-2c\\
&\Leftrightarrow& \dfrac{112}{5}a+\dfrac{32}3b+2c=-\dfrac{28}{5}\qquad(4)
\end{eqnarray*}
Giải hệ 3 phương trình $(2)$, $(3)$ và $(4)$ ta được $a=1$, $b=-3$, $c=2$.\\
Khi đó, $(C)\colon y=x^4-3x^2+2$, $d\colon y=2\left(x+1\right)$.\\
Diện tích cần tìm là
\[S=\displaystyle\int\limits_{-1}^0\left[x^4-3x^2+2-2\left(x+1\right)\right]\mathrm{\,d}x=\displaystyle\int\limits_{-1}^0\left(x^4-3x^2-2x\right)\mathrm{\,d}x=\dfrac1{5}=0{,}2.\]
}
\end{ex}

\begin{ex}%[Vovanle]%[2D4V3-1]
\immini{Cho hai đường tròn $\left(O_1;5\right)$ và $\left(O_2;3\right)$ cắt nhau tại hai điểm $A$, $B$ sao cho $AB$ là một đường kính của đường tròn $\left(O_2\right)$. Gọi $(D)$ là hình thẳng được giới hạn bởi hai đường tròn (phần ở ngoài đường tròn lớn, được gạch chéo như hình vẽ). Một vật trang trí có dạng một khối tròn xoay được tạo thành khi quay miền $(D)$ quanh trục $O_1O_2$. Thể tích của khối tròn xoay được tạo thành có $V=\dfrac{a\pi}{b}$ ($\dfrac{a}{b}$ là phân số tối giản) thì $a^2+b^3$ bằng bao nhiêu?
}{
\begin{tikzpicture}[line join=round,line cap=round, font=\footnotesize,scale=1,>=stealth]
\def \r{0.3}
\path
(-4*\r,0) coordinate (O_1)
(0,0) coordinate (O_2)
(90:3*\r) coordinate (A)
(-90:3*\r) coordinate (B)
;
\pgfmathsetmacro\g{atan (3/4)}
\fill[pattern=north east lines]plot[domain=-\g:\g](5*\r*cos \x-4*\r,5*\r*sin \x)--plot[domain=90:-90](3*\r*cos \x,3*\r*sin \x);
\draw (O_1) circle(5*\r);
\draw (O_2) circle(3*\r);
\draw (2*\r,0) coordinate (D)node[above]{$(D)$};
\draw (A)--(B)(-9*\r,0)--(3*\r,0);
\foreach \x/\g in {O_1/90,O_2/140,A/60,B/-60}\fill[black] (\x) circle (1pt)+(\g:.3)node{$\x$};
\end{tikzpicture}
}
\shortans{$1627$}
\loigiai{
\immini{Chọn hệ tọa độ $Oxy$ với \\
$O_2\equiv O$, $O_2C\equiv Ox$, $O_2A\equiv Oy$.\\
Đoạn $O_1O_2=\sqrt{O_1A^2-O_2A^2}=\sqrt{5^2-3^2}=4$.\\
Suy ra $\left(O_1\right):{{\left( x+4 \right)}^2}+y^2=25$.\\
Kí hiệu $\left(H_1\right)$ là hình phẳng giới hạn bởi các đường $\left(O_1\right)\colon \left(x+4\right)^2+y^2=25$, $Oy\colon x=0$, $x\geq 0$.\\
Kí hiệu $\left(H_2\right)$ là hình phẳng giới hạn bởi các đường $\left(O_2\right)\colon x^2+y^2=9$, $Oy\colon x=0$, $x\geq 0$.
}{
\begin{tikzpicture}[line join=round,line cap=round, font=\footnotesize,scale=1,>=stealth]
\def \r{0.4}
\path
(-4*\r,0) coordinate (O_1)
(0,0) coordinate (O_2)
(90:3*\r) coordinate (A)
(-90:3*\r) coordinate (B)
;
\pgfmathsetmacro\g{atan (3/4)}
\fill[pattern=north east lines]plot[domain=-\g:\g](5*\r*cos \x-4*\r,5*\r*sin \x)--plot[domain=90:-90](3*\r*cos \x,3*\r*sin \x);
\draw[-stealth](-9.5*\r,0)--(4*\r,0)node[below]{$x$};
\draw[-stealth](0,-5.5*\r)--(0,5.5*\r)node[right]{$y$};
\draw (O_1) circle(5*\r);
\draw (O_2) circle(3*\r);
\draw (2*\r,0) coordinate (D)node[above]{$(D)$};
\foreach \x/\g in {O_1/90,O_2/140,A/60,B/-60}\fill[black] (\x) circle (1pt)+(\g:.3)node{$\x$};
\end{tikzpicture}
}
\noindent Khi đó thể tích $V$ cần tìm chính bằng thể tích $V_2$ của khối tròn xoay thu được khi quay hình $\left(H_2\right)$ xung quanh trục $Ox$ trừ đi thể tích $V_1$ của khối tròn xoay thu được khi quay hình $\left(H_1\right)$ xung quanh trục $Ox$.\\
Ta có $V_2=\dfrac{1}{2}\cdot \dfrac{4}{3}\pi r^3=\dfrac{2}{3}\pi {\cdot 3^3}=18\pi$.\\
Lại có $V_1=\pi\displaystyle\int\limits_0^1y^2\mathrm{\,d}x=\pi\displaystyle\int\limits_0^1\left[25-\left(x+4\right)^2\right]\mathrm{\,d}x=\left.\pi \left[25x-\dfrac{\left(x+4\right)^3}{3}\right]\right|_0^1
=\dfrac{14\pi}{3}$.\\
Do đó $V=V_2-V_1=18\pi-\dfrac{14\pi}{3}=\dfrac{40\pi}{3}$.\\
Vậy $a^2+b^3=1627$.
}
\end{ex}

\begin{ex}%[Dự án 2025 - đề cấu trúc mới, Nguyễn Kiều Nhã Tú]%[2D4V3-1]
Cho hàm số $f(x)=ax^4+bx^2+c$ ($a\ne 0$) có bảng biến thiên như sau
\begin{center}
\begin{tikzpicture}
\tkzTabInit[lgt=1.5,espcl=2.5]
{$x$/1,$f’(x)$/1,$f(x)$/2}
{$-\infty$,$-1$,$0$,$1$,$+\infty$}
\tkzTabLine{ ,+,z,-,z,+,z,-, }
\tkzTabVar{-/$-\infty$,+/$2$,-/$1$,+/$2$,-/$-\infty$}
\end{tikzpicture}
\end{center}
Tính diện tích hình phẳng giới hạn bởi đồ thị $(C) \colon y=f(x)$ và đồ thị $(P)\colon y=\dfrac{1}{2}x^2$. Kết quả làm tròn đến hàng phần trăm.
\shortans{$3{,}39$}
\loigiai{
\begin{center}
\begin{tikzpicture}
\def\a{-1}
\def\b{2}
\def\c{1}
\def\f(#1){\a*((#1)^4)+\b*((#1)^2)+\c}
\def\xmin{-1.7}
\def\xmax{1.7}
\def\ymin{-1}
\def\ymax{2}
\draw[->] (\xmin-1,0)--(\xmax+1,0) node[below] { $x$};
\draw[->] (0,\ymin-1)--(0,\ymax+1) node[left] { $y$};
\draw (0,0) node [below left] {$O$};
\draw[smooth,samples=200] plot[domain=\xmin:\xmax]  (\x,{\f(\x)});
\draw[smooth,samples=200] plot[domain=\xmin-0.3:\xmax+0.3]  (\x,{0.5*(\x)^2}) node [above right]{$(P): y=\dfrac{1}{2}x^2$};
\draw (\xmax,{\f(\xmax)}) node [above right]{$(C): y=x^4-2x^2+1$};
\draw[dashed] (-1,0) node[below]{$-1$}--(-1,2)--(0,2) node[above right]{$2$}--(1,2)--(1,0)node[below]{$1$};
\end{tikzpicture}
\end{center}
Ta có $y=f(x)=ax^4+bx^2+c$ với $a \ne 0$. Khi đó $y'=f'(x)=4ax^3+2bx$.\\
Từ bảng biến thiên ta có
$\heva{&y'(1)=0\\&y(1)=2\\&y(0)=1} \Leftrightarrow \heva{&4a+2b=0\\&a+b+c=2\\&c=1} \Leftrightarrow \heva{&a=-1\\&b=2\\&c=1.}$\\
Vậy $y=f(x)=-x^4+2x^2+1$.\\
Phương trình hoành độ giao điểm của $(C)$ và $(P)$ là $-x^4+2x^2+1=\dfrac{x^2}{2}\Leftrightarrow x=\pm \sqrt{2}$.\\
Diện tích hình phẳng giới hạn bởi đồ thị $(C)$ và $(P)$ là \[S=\displaystyle\int\limits_{-\sqrt{2}}^{\sqrt{2}}\left(-x^4+\dfrac{3}{2}x^2+1\right) \mathrm{\,d}x=\dfrac{12\sqrt{2}}{5}\approx 3{,}39.\]
}
\end{ex}

\begin{ex}%[Dự án 2025 - đề cấu trúc mới, Nguyễn Kiều Nhã Tú]%[2D4V3-1]
Cho parabol $(P_1)\colon y=\dfrac{x^2}{2}$, $(P_2)\colon y=-\dfrac{x^2}{2}$ và đường tròn $(C)$ có tâm là gốc tọa độ, bán kính bằng $\sqrt{8}$. Tính diện tích hình phẳng giới hạn bởi $(P_1)$, $(P_2)$ và $(C)$ (kết quả làm tròn đến hàng phần chục).
\begin{center}
\begin{tikzpicture}
\def\a{0.5}
\def\d{-0.5}
\def\b{0}
\def\c{0}
\def\f(#1){\a*((#1)^2)+\b*(#1)+\c}
\def\g(#1){\d*((#1)^2)+\b*(#1)+\c}
\def\xmin{-2.84}
\def\xmax{2.84}
\def\ymin{-3}
\def\ymax{3}
\begin{scope}
\clip (0,0) circle (2.82 cm);
\fill[gray,smooth] (-2.82,0) --
plot[domain=-2.82:2.82] (\x,{\f(\x)})
-- (2.82,0) -- cycle;
\fill[gray,smooth] (-2.82,0) --
plot[domain=-2.82:2.82] (\x,{\g(\x)})
-- (2.82,0) -- cycle;
\end{scope}
\draw[->] (\xmin-1,0)--(\xmax+1,0) node[below] {$x$};
\draw[->] (0,\ymin-1)--(0,\ymax+1) node[left] {$y$};
\draw (0,0) node [below left] {$O$};
\foreach \x in {-3,-2,-1,1,2,3}		\draw (\x,0.1)--(\x,-0.1) node [below] { $\x$};
\foreach \y in {-3,-2,-1,1,2,3}		\draw (0.1,\y)--(-0.1,\y) node [left] { $\y$};
\clip (\xmin,\ymin) rectangle (\xmax,\ymax);
\draw[smooth,samples=200] plot[domain=\xmin:\xmax] (\x,{\f(\x)});
\draw[smooth,samples=200] plot[domain=\xmin:\xmax] (\x,{\g(\x)});
\draw (0,0) circle (2.82 cm);
\end{tikzpicture}
\end{center}
\shortans{$9{,}9$}
\loigiai{
Phương trình hoành độ giao điểm của $(P_1)$ và $(C)$ là $\sqrt{8-x^2}=\dfrac{x^2}{2} \Leftrightarrow x= \pm 2$.\\
Phương trình hoành độ giao điểm của $(P_2)$ và $(C)$ là $-\sqrt{8-x^2}=-\dfrac{x^2}{2} \Leftrightarrow x= \pm 2$.\\
Diện tích hình phẳng giới hạn bởi $(P_1)$, $(P_2)$ và $(C)$ là \[S=4\left(\displaystyle\int\limits_1^2 \dfrac{x^2}{2}\mathrm{\,d}x+\displaystyle\int\limits_2^3 \sqrt{8-x^2} \mathrm{\,d}x\right) \approx 9{,}9.\]
}
\end{ex}

\begin{ex}%[HSA-ĐỢT 2-Nguyễn Mạnh Trung Anh]%[2D4V3-1]
\immini[thm]{
Tính thể tích của một chiếc cốc thủy tinh được tạo thành bằng cách quay hình thang vuông $OABC$ với $O A=6$ cm, $AB=3$ cm và $O C=2$ cm, quanh trục đối xứng của nó (làm tròn đến hàng đơn vị).
\par\shortans{119}

}{
\rotatebox{90}{\begin{tikzpicture}[scale=0.8, font=\footnotesize, line join=round, line cap=round, >=stealth]
\fill[yellow!30](0,0)--(0,1)--(3,2)--(3,0)--cycle;
%\draw[->] (3,0) -- (4.5,0) node[right] {$x$};
%\draw[->] (0,1) -- (0,2.5) node[above] {$y$};
\draw (0,0) ellipse (0.3cm and 1cm);
\draw (3,0) ellipse (0.5cm and 2cm);
\draw (0,1)--(3,2)(0,-1)--(3,-2);
\draw[dashed] (0,-1)--(0,1);
\draw (3,2)--(3,0);
\draw[dashed] (-0.3,0)--(3,0);
%\draw[dashed] (0,0)--(0,1) node[pos=0.5,left]{$r$};
\fill (0,0) circle (1pt) node[below right]{$O$};
\fill (3,0) circle (1pt) node[below]{$A$};
\fill (3,2) circle (1pt) node[above]{$B$};
\fill (0,1) circle (1pt) node[above left]{$C$};
%\fill (-3,0) circle (1pt) node[above left]{$D$};
%\fill (0.1,0.5) circle (0pt) node[left]{$r$};
\fill (0.5,-2) circle (0pt) ;%node[below]{Hình $4.28$};
\end{tikzpicture}}
}
\loigiai{
\immini{
Đặt $OC=r, AB=R, OA=h$.\\
Ta tìm được phương trình đường thẳng $CB$ đi qua hai điểm $\left(0;r\right)$ và $\left(h; R\right)$ là $y=\dfrac{R-r}{h}x+r$.\\
Khi đó thể tích của khối tròn xoay sinh ra là
\begin{eqnarray*}
V &=& \pi \displaystyle\int_{0}^h\left( \dfrac{R-r}{h}x+r\right)^2 \mathrm{\,d}x \\
&=& \pi \displaystyle\int_{0}^h \left[ \dfrac{\left( R-r \right)^2}{h^2}x^2+2\dfrac{\left( R-r \right)r}{h}x+r^2 \right] \mathrm{\,d}x \\
&=&\pi \left. \left[ \dfrac{\left( R-r \right)^2}{h^2}\cdot \dfrac{x^3}{3}+2\dfrac{\left( R-r \right)r}{h} \cdot \dfrac{x^2}{2}+r^2x \right] \right|_0^h \\
&=&\pi \left[ \dfrac{\left( R-r \right)^2}{h^2}\cdot \dfrac{h^3}{3}+2\dfrac{\left( R-r \right)r}{h}\cdot \dfrac{h^2}{2}+r^2h \right] \\
&=&\pi \left[ \dfrac{\left( R-r \right)^2h}{3}+\left( R-r \right)rh+r^2h \right] \approx 119 (\mathrm{cm^3}).
\end{eqnarray*}
Vậy thể tích của chiếc cốc thủy tinh là $119\, \mathrm{cm^3}$.
}
{
\begin{tikzpicture}[scale=1, font=\footnotesize, line join=round, line cap=round, >=stealth]
\fill[yellow!30](0,0)--(0,1)--(3,2)--(3,0)--cycle;
\draw[->] (-1,0) -- (4.5,0) node[right] {$x$};
\draw[->] (0,-1.5) -- (0,2.5) node[above] {$y$};
\draw (0,0) ellipse (0.3cm and 1cm);
\draw (3,0) ellipse (0.5cm and 2cm);
\draw (0,1)--(3,2)(0,-1)--(3,-2);
\draw[dashed] (0,-1)--(0,1);
\draw (3,2)--(3,0);
\draw (3,1) node [right]{$R$} (1.5,0) node [above]{$h$};
%\draw[dashed] (0,0)--(0,1) node[pos=0.5,left]{$r$};
\fill (0,0) circle (1pt) node[below right]{$O$};
\fill (3,0) circle (1pt) node[below]{$A$};
\fill (3,2) circle (1pt) node[above]{$B$};
\fill (0,1) circle (1pt) node[above left]{$C$};
%\fill (-3,0) circle (1pt) node[above left]{$D$};
\fill (0.1,0.5) circle (0pt) node[left]{$r$};
\fill (0.5,-2) circle (0pt) ;%node[below]{Hình $4.28$};
\end{tikzpicture}
}
}
\end{ex}

\begin{ex}%[2D4V3-1]
Biết $F(x)$, $G(x)$ là hai nguyên hàm của $f(x)$ trên $\mathbb{R}$ và $\displaystyle \int\limits_{0}^{7} f(x) \mathrm{\,d}x = F(7)-G(0)+3m$ $(m>0)$. Gọi $S$ là diện tích hình phẳng giới hạn bởi các đường $y=F(x)$, $y=G(x)$, $x=0$ và $x=7$. Khi $S=105$ thì giá trị của $m$ bằng bao nhiêu?
\shortans{$5$}
\loigiai{
Ta có $G(x)= F(x)+C$.\\
Theo giả thiết $\displaystyle \int\limits_{0}^{7} f(x) \mathrm{\,d}x = F(7)-G(0)+3m$, $(m>0)$ nên
\begin{eqnarray*}
\heva{& F(7) - F(0) = F(7) - G(0)+3m\\& G(7)-G(0)=  F(7) - G(0)+3m} \Rightarrow \heva{& G(0)-F(0)=3m\\& G(7)-F(7)= 3m} \Rightarrow G(x)-F(x)=3m.
\end{eqnarray*}
Khi đó $S= \displaystyle \int\limits_{0}^{7} \left| G(x)-F(x) \right| \mathrm{\,d}x = \int\limits_{0}^{7} \left|3m \right| \mathrm{\,d} = \int\limits_{0}^{7} 3m \mathrm{\,d}x=21m$.\\
Theo giả thiết thì $21m=105 \Leftrightarrow m = 5$.
}
\end{ex}

\begin{ex}%[2D4V3-1]
Tính diện tích hình phẳng giới hạn bởi đồ thị của hai hàm số $y=\dfrac{x^2}{12}$ và $y=\sqrt{4-\dfrac{x^2}{4}}$ (\textit{làm tròn kết quả đến chữ số thập phân thứ hai}).
\shortans{9,53}
\loigiai{
Điều kiện xác định $4-\dfrac{x^2}{4}\geq0\Leftrightarrow-4\leq x\leq4$.\\

Xét phương trình hoành độ giao điểm
\begin{eqnarray*}
& &\dfrac{x^2}{12}=\sqrt{4-\dfrac{x^2}{4}}\\
&\Leftrightarrow&\dfrac{x^4}{12^2}+\dfrac{x^2}{4}-4=0\\
&\Leftrightarrow&x=\pm2\sqrt{3} \quad \text{(thoả mãn điều kiện xác định).}
\end{eqnarray*}

Diện tích hình phẳng giới cần tính là \[S=\int_{-2\sqrt{3}}^{2\sqrt{3}} \left|\dfrac{x^2}{12}-\sqrt{4-\dfrac{x^2}{4}}\right|dx\approx9,53.\]
}
\end{ex}

\begin{ex}%[HSA - đợt 2, Nguyen The Tuan Vu]%[2D4V3-1]
Cho hình phẳng $(H)$ giới hạn bởi đồ thị hàm số $y=2x^2$ và đường thẳng $y=mx$ với $m\neq 0$. Có bao nhiêu giá trị nguyên dương của tham số $m$ để diện tích hình phẳng $(H)$ nhỏ hơn $10$?
\shortans{$6$}
\loigiai{
Phương trình hoành độ giao điểm của đồ thị hàm số $y=2x^2$ và đường thẳng $y=mx$ là \[2x^2=mx\Leftrightarrow 2x^2-mx=0\Leftrightarrow x(2x-m)=0\Leftrightarrow\hoac{& x=0 \\ & x=\dfrac{m}{2}.}\]\\
Vì $m>0$ nên ta có đồ thị hàm số $y=2x^2$ và đường thẳng $y=mx$ là
\begin{center}
\begin{tikzpicture}[scale=1,>=stealth, font=\footnotesize, line join=round, line cap=round]
\def\xmin{-2} \def\xmax{2}
\def\ymin{-1} \def\ymax{3}
\draw[->] (\xmin,0)--(\xmax,0) node [below]{$x$};
\draw[->] (0,\ymin)--(0,\ymax) node [left]{$y$};
\node at (0,0) [below left]{$O$};
\clip (\xmin+0.1,\ymin+0.1) rectangle (\xmax-0.5,\ymax-0.1);
\draw[smooth,samples=300] plot(\x,{2*(\x)^2});
\draw[smooth,samples=300] plot(\x,{2*(\x)});
\end{tikzpicture}
\end{center}
$\Rightarrow$ Diện tích của hình phẳng $(H)$ là \[S=\displaystyle\int\limits_{0}^{\tfrac{m}{2}} \left(mx-2x^2\right) \mathrm{\,d}x=\left(\dfrac{mx^2}{2}-\dfrac{2x^3}{3}\right)\Big|^{\tfrac{m}{2}}_0=\dfrac{m}{2}\cdot\left(\dfrac{m}{2}\right)^2-\dfrac{2}{3}\cdot\left(\dfrac{m}{2}\right)^3=\dfrac{m^3}{24}.\]
Vì $S<10$ nên $\dfrac{m^3}{24}<10\Leftrightarrow m^3<240\Leftrightarrow m<2\sqrt[3]{30}$.\\
Vì $m$ nguyên dương nên $m\in\{1;2;3;4;5;6\}$.\\
$\Rightarrow$ Có $6$ giá trị nguyên dương của tham số $m$.
}
\end{ex}

\begin{ex}%[12-EX-3-2026, Bùi Quang Phú]%[2D4V3-1]
Cho hàm số $y=f(x)$. Đồ thị hàm số $y=f'(x)$ là đường cong trong hình dưới. Biết rằng diện tích của các phần hình phẳng $A$ và $B$ lần lượt là $S_A=4$ và $S_B=10$. Tính giá trị của $f(3)$, biết giá trị của $f(0)=2$.
\begin{center}
\begin{tikzpicture}[scale=1, font=\footnotesize, line join=round, line cap=round, >=stealth]
\def \xmin{-1}
\def \xmax{4}
\def \ymin{-2.2}
\def \ymax{1.2}
%%Tự động
\draw[->] (\xmin,0)--(\xmax,0) node[below left] {$x$};
\draw[->] (0,\ymin)--(0,\ymax) node[below left] {$y$};
\draw[fill=black] (0,0) circle(1pt) node [below left] {$O$};
\draw[fill=black]
(1,0) circle(1pt) node[shift=(50:0.2)]{$1$}
(3,0) circle(1pt) node[shift=(50:0.2)]{$3$}
;
%%Tự động
\begin{scope}[xshift=1cm]
\clip (\xmin-0.4,\ymin+0.01) rectangle (\xmax-0.01,\ymax-0.01);
\draw[samples=350,domain=\xmin-1:\xmax-0.01,smooth,variable=\x] plot (\x,{(\x+1)*(\x)*(\x-2)});
\draw[fill=gray!60,opacity=0.5] (-1,0)plot[domain=-1:0] (\x,{(\x+1)*(\x)*(\x-2)})--(0,0)--cycle;
\draw[fill=gray!60,opacity=0.5] (0,0)plot[domain=0:2] (\x,{(\x+1)*(\x)*(\x-2)})--(2,0)--cycle;
\end{scope}
\path
(0.5,0.3) node{$(A)$}
(2.1,-1) node{$(B)$}
;
\end{tikzpicture}
\end{center}
\shortans[oly]{$-4$}
\loigiai{
Ta có $S_A=\displaystyle\int\limits_{0}^{1} f'(x) \mathrm{\,d}x=f(x)\Big|_{0}^{1}=f(1)-f(0)=4$. Suy ra $f(1)=4+f(0)=4+2=6$.\\
Mặt khác $S_B=-\displaystyle\int\limits_{1}^{3} f'(x) \mathrm{\,d}x=-f(x)\Big|_{1}^{3}=f(1)-f(3)=10$. Suy ra $f(3)=f(1)-10=6-10=-4$.
}
\end{ex}

\begin{ex}%[12-EX3-2026. Trần Hoà]%[2D4V3-1]
\immini{Cho hàm số $y=f(x)$ liên tục và có đồ thị trên đoạn $[-2;5]$ gồm một phần của đường parabol $(P)$ và một phần của đường thẳng $(d)$ như hình vẽ sau. Tính tích phân $I=\displaystyle\int\limits_{-2}^{5}f(x)\mathrm{\,d}x$. (Làm tròn kết quả đến hàng phần mười)
\shortans[]{$15{,}2$}}{\begin{tikzpicture}[scale=1,font=\footnotesize,line join = round, line cap = round, >= stealth]
\draw (-2,3)--(1,0);
\draw (1,0) parabola bend ((3,4)(5,0);%hoanh-tang-dan
\def\a{-2.5} \def\b{6} \def\c{-.5} \def\d{5}
\coordinate (O) at (0,0);
\draw[->] (\a,0)--(\b,0)node[below]{$x$};
\draw[->] (0,\c)--(0,\d)node[right]{$y$};
\foreach \p/\g in {-2/-90,1/-90,3/-90,5/-90} \draw[fill] (\p,0) circle(.5pt)node [shift={(\g:.3)}] {$\p$};
\foreach \p/\g in {3/145,4/180} \draw[fill] (0,\p) circle(.5pt)node [shift={(\g:.3)}] {$\p$};
\draw[dashed] (-2,0)|-(0,3) (3,0)|-(0,4);
\draw[fill] (O) circle (.5 pt) node[above right]{$O$} ;
\end{tikzpicture}}
\loigiai{
Dựa vào đồ thị\\
Trên đoạn $[-2; 1]$, đồ thị là đoạn thẳng đi qua $(-2; 3)$ và $(1; 0)$.\\
Phương trình đường thẳng: $\dfrac{x-1}{-2-1} = \dfrac{y-0}{3-0} \Rightarrow y = -(x-1)\Rightarrow y = -x+1$.\\
Trên đoạn $[1; 5]$, đồ thị là parabol có đỉnh $I(3; 4)$ và đi qua $(1; 0), (5; 0)$.\\
Do đồ thị đi qua $(1; 0)$, $(5; 0)$ nên phương trình có dạng $y = a(x-1)(x-5)$.\\
Đỉnh $I(3;4) \Rightarrow 4 = a(3-1)(3-5)\Leftrightarrow 4=-4a \Rightarrow a=-1$.\\
Vậy $(P)\colon y = -x^2+6x-5$.\\
Suy ra
$I = \displaystyle\int\limits_{-2}^1 (-x+1) \mathrm{\,d}x + \displaystyle\int\limits_1^{5} (-x^2+6x-5) \mathrm{\,d}x = \dfrac{9}{2} + \dfrac{32}{3}=\dfrac{91}{6} \approx 15{,}2$. (làm tròn đến hàng phần mười)
}
\end{ex}

\begin{ex}%[12-EX-3-2025, Trần Văn Hưng]%[2D4V3-1]
Biết $F(x)$ và $G(x)$ là hai nguyên hàm của hàm số $f(x)$ trên $\mathbb{R}$ và thoả mãn $\displaystyle\int\limits_0^4f(x)\mathrm{\,d}x=F(4)-G(0)+2m$, với $m>0$. Gọi $S$ là diện tích hình phẳng giới hạn bởi các đường $y=F(x)$, $y=G(x)$, $x=0$ và $x=4$. Khi $S=8$ thì $m$ bằng?
\par
\shortans{1}
\loigiai{
Theo đề ta có $\displaystyle\int\limits_0^4f(x)\mathrm{\,d}x=F(4)-G(0)+2m\Rightarrow F(x)\bigg|_0^4=F(4)-G(0)+2m$.\\
$\Rightarrow F(4)-F(0)=F(4)-G(0)+2m\Rightarrow G(0)-F(0)=2m$. $(1)$\\
Mặt khác, do $F(x)$ và $G(x)$ là hai nguyên hàm của hàm số $f(x)$ trên $\mathbb{R}$ nên ta có $G(x)-F(x)=C$ (không đổi) với mọi $x\in \mathbb{R}$. $(2)$\\
Từ $(1)$ và $(2)$ suy ra $G(x)-F(x)=2m>0$, với mọi $x\in \mathbb{R}$.\\
Khi đó ta có $S=\displaystyle\int\limits_0^4\left|G(x)-F(x)\right|\mathrm{\,d}x=\displaystyle\int\limits_0^42m \mathrm{\,d}x=2mx\bigg|_0^4=8m$.\\
Theo đề ta có $8m=8\Leftrightarrow m=1$.
}
\end{ex}

\begin{ex}%[12-EX-3-2025, Thái Bảo]%[2D4V3-1]
\immini{
Một gia đình muốn làm cái cổng (như hình vẽ). Phần phía trên cổng có hình dạng là một parabol với $IH=2{,}5$ m, phần phía dưới là một hình chữ nhật có kích thước $AD=4\,\text{m}$, $AB=6\,\text{m}$. Giả sử giá để làm phần cổng được tô màu (phần hình chữ nhật $ABCD$) là $900\,000$ đ/m$^2$ và giá để làm phần cổng phía trên là $1\,300\,000$ đ/m$^2$. Tính số tiền gia đình đó phải trả là bao nhiêu triệu đồng.
\par\shortans{$34{,}6$}
}{
\begin{tikzpicture}[font=\footnotesize, line join=round, line cap=round, >=stealth, scale=.8]
\path
(0,0) coordinate (A)
(6,0) coordinate (B)
(6,-4) coordinate (C)
(0,-4) coordinate (D);
\draw (A)--(B)--(C)--(D)--cycle (3,0)--(3,2.5) node [midway, right] {$2{,}5$ m};
\draw [fill= black] (3,2.5) circle (2pt) node [above] {$I$};
\draw [fill= black] (3,0) circle (2pt) node [below] {$H$};
\node at (0,-2) [right] {$4$ m};
\node at (3,-4) [below] {$6$ m};
\draw[domain=0:6, samples=100]
plot (\x, {-0.28*\x*\x + 1.67*\x});
\foreach \t/\g in {A/135,B/45,C/-45,D/-135}{
\draw [fill=black] (\t) circle (2pt) node [shift={(\g:7pt)}] {$\t$};
}
\end{tikzpicture}

}
\loigiai{
Diện tích hình chữ nhật $ABCD$ là
\[
S_1=AB\cdot AD=6\cdot4=24\ (\text{m}^2).
\]
Chọn hệ trục tọa độ sao cho $H$ là gốc tọa độ, trục $Ox$ trùng với $AB$, trục $Oy$ hướng lên trên.\\
Khi đó parabol có đỉnh $I(0;2{,}5)$ và cắt trục $Ox$ tại $A(-3;0)$ và $B(3;0)$.\\
Gọi phương trình parabol có dạng $y=ax^2+2{,}5$.\\
Thay tọa độ điểm $A(-3;0)$ vào, ta được $0=9a+2{,}5 \Rightarrow a=-\dfrac{2{,}5}{9}$.\\
Suy ra phương trình parabol là $y=-\dfrac{2{,}5}{9}x^2+2{,}5$.\\
Diện tích phần cổng phía trên là
\begin{eqnarray*}
S_2&=&\displaystyle\int\limits_{-3}^{3}\left(-\dfrac{2{,}5}{9}x^2+2{,}5\right)\mathrm{\,d}x\\
&=&2\displaystyle\int\limits_0^{3}\left(-\dfrac{2{,}5}{9}x^2+2{,}5\right)\mathrm{\,d}x.\\
&=&2\left.\left(2{,}5x-\dfrac{2{,}5}{27}x^3\right)\right|_0^{3}\\
&=&2\cdot(7{,}5-2{,}5)\\
&=&10\ (\text{m}^2).
\end{eqnarray*}
Chi phí làm phần hình chữ nhật là
$24\cdot900\,000=21\,600\,000$ (đồng).\\
Chi phí làm phần phía trên là $10\cdot1\,300\,000=13\,000\,000$ (đồng).\\
Tổng số tiền gia đình phải trả là $21\,600\,000+13\,000\,000=34\,600\,000$ (đồng).\\
Vậy số tiền gia đình đó phải trả là $34{,}6$ triệu đồng.
}
\end{ex}

\begin{ex}%[Đề kiểm tra giữa kì 2 - THPT Thanh Miện, Hải Dương, Nguyễn Anh Tuấn, Dự án 12EX-3-2026]%[2D4V3-1]
\immini{Cho hình thang cong $(H)$ giới hạn bởi các đường $y=\mathrm{e}^x$, $y=0$, $x=0$, $x=\ln 4$. Đường thẳng \break  $x=k\,(0<k<\ln 4)$ chia $(H)$ thành hai phần có diện tích là $S_1$ và $S_2$ như hình vẽ bên. Tìm $k$ để $S_1=2 S_2$. (kết quả làm tròn đến hàng phần chục)}
{\begin{tikzpicture}[
scale=0.9,
>=stealth,
every node/.style={font=\footnotesize}
]
\draw[domain=-1:1.4, smooth=200]
plot (\x,{exp(\x)});
\fill[gray!25]
plot[domain=0:0.7] (\x,{exp(\x)}) --
(0.7,0) -- (0,0) -- cycle;
\node at (0.35,0.9) {$S_1$};
\fill[gray!40]
plot[domain=0.7:1.386] (\x,{exp(\x)}) --
(1.386,0) -- (0.7,0) -- cycle;
\node at (1.05,1.7) {$S_2$};
\draw[->] (-1,0) -- (2.2,0) node[right] {$x$};
\draw[->] (0,-0.3) -- (0,4.6) node[above] {$y$};
\draw[dashed] (0.7,0) -- (0.7,4.5);
\draw[dashed] (1.386,0) -- (1.386,4.5);
\foreach \x/\y in{0/1,0/0}\draw[fill=black](\x,\y)circle(1.25pt);
\foreach \x/\y/\goc/\z in{1.386/0/-90/\ln 4,0/1/150/1,0.7/0/-90/k,0/0/-140/O}
\path(\x,\y)node[shift={(\goc:7pt)}]{\scriptsize$\z$};
\end{tikzpicture}}
\shortans[oly]{1,1}
\loigiai{
Theo giả thiết, ta có công thức tính diện tích hai miền phẳng $S_1$ và $S_2$ như sau
\begin{itemize}
\item $S_1 = \displaystyle\int\limits_{0}^{k} \mathrm{e}^x \,\mathrm{\,d}x = \mathrm{e}^x \left|_{0}^{k} = \mathrm{e}^k - \mathrm{e}^0 = \mathrm{e}^k-1 \right.$.
\item $S_2 = \displaystyle\int\limits_{k}^{\ln 4} \mathrm{e}^x \,\mathrm{\,d}x = \mathrm{e}^x \left|_{k}^{\ln 4}=\mathrm{e}^{\ln 4} - \mathrm{e}^k=4-\mathrm{e}^k\right.$.
\end{itemize}
Theo giả thiết $S_1=2S_2$, ta có phương trình
\[ \mathrm{e}^k-1=2(4 - \mathrm{e}^k) \Leftrightarrow \mathrm{e}^k-1=8-2\mathrm{e}^k \Leftrightarrow 3\mathrm{e}^k=9 \Leftrightarrow \mathrm{e}^k=3 \Leftrightarrow k=\ln 3\approx 1{,}1. \]
Vậy $k \approx 1{,}1$.
}
\end{ex}

\begin{ex}%[2D4V3-1]%[Dự án C - đợt 3 - NH24-25-PNL-Deso13-Lê Minh Thien Anh]
\immini{Cho hàm số bậc hai trên bậc nhất $f(x) = \dfrac{ax^2+bx+c}{4x-d}$ và parabol $(P)$ có đồ thị như hình vẽ.
Hàm số $y=f(x)$ có đường tiệm cận đứng $x=1$ và đường tiệm cận xiên cắt parabol tại điểm $A(3;1)$. Ký hiệu các diện tích hình phẳng $S_1, S_2$ lần lượt là các phần diện tích được tô đậm và không tô màu như hình vẽ. Biết rằng giá trị $S_1 + S_2 = \dfrac{32}{3}$ và $f(4) = -\dfrac{7}{4}$. Tính giá trị của $S_2$ (Kết quả làm tròn đến chữ số hàng phần chục).}
{
\begin{tikzpicture}[scale=0.95, font=\footnotesize, line join=round,
line cap=round, >=stealth]
\def \xmin{-2}
\def \xmax{6}
\def \ymin{-2}
\def \ymax{2}
\fill[cyan!50] (-1,-5/3)--(3,1)--plot[domain=3:2.9] (\x,{-1/3*(\x)^2+4/3*(\x)}) -- plot[domain=2.9:1.19] (\x,{2/3*(\x)-1+1/(4*(\x)-4)})--plot[domain=1.19:-1] (\x,{-1/3*(\x)^2+4/3*(\x)});
\draw[->] (\xmin,0)--(\xmax,0) node[shift=(-110:0.2)] {$x$};
\draw[->] (0,\ymin)--(0,\ymax+0.5) node[shift=(-150:0.2)] {$y$};
\fill (0,0) circle(1pt) node[shift=(135:0.25)]{$O$};
\begin{scope}
\clip (\xmin,\ymin) rectangle (\xmax,\ymax);
\draw[smooth,samples=100,domain=\xmin:\xmax,thick] plot(\x,{-1/3*(\x)^2+4/3*(\x)});
\draw[smooth,samples=100,domain=1.1:\xmax,thick] plot(\x,{2/3*(\x)-1+1/(4*(\x)-4)});
\draw[smooth,samples=100,domain=\xmin:\xmax,thick,dashed] plot(\x,{2/3*(\x)-1});
\draw[dashed] (1,\ymin)--(1,\ymax) (-1,-5/3) circle(1pt) node[above left] {$B$} (3,1) circle(1pt) node[right] {$A(3;1)$};
\end{scope}
\draw[fill=white] (0.1,-.6) rectangle (0.5,-.2) node[pos=0.5] {$S_1$} ;
\draw[fill=white] (1.7,0.7) rectangle (2.1,1.1) node[pos=0.5] {$S_2$} ;
\draw (1.1,1.8) node[right] {$y=f(x)$};
\draw (1,0) node[above left] {$1$};
\end{tikzpicture}
}
\par \shortans[0]{0,94}
\loigiai{
Gọi điểm $B(x_2, y_2)$.\\
Sử dụng công thức tính diện tích hình phẳng giới hạn bởi đường thẳng và parabol, ta có
\[ \dfrac{(x_1 - x_2)^3}{6} = S_1 + S_2 = \dfrac{32}{3} \Leftrightarrow (x_1 - x_2)^3 = 64 \Leftrightarrow x_1 - x_2 = 4 \Leftrightarrow \hoac{& x_1 = 3 \\& x_2 = -1.} \]
Phương trình parabol $(P)\colon y = ax(x-4)$ đi qua điểm $A(1;3)$ nên suy ra $a = -\dfrac{1}{3}$.\\
Khi đó phương trình parabol $(P)$ là $y = -\dfrac{1}{3}x(x-4)$ suy ra $B\left(-1; -\dfrac{5}{3}\right)$.\\
Phương trình đường thẳng $AB$ cũng là đường tiệm cận xiên của hàm số $y=f(x)$ có phương trình là $y = \dfrac{2}{3}x - 1$.\\
Đồ thị hàm số có đường tiệm cận đứng là đường thẳng $x=1$ nên suy ra $d=4$.\\
Ta có $f(x) = \dfrac{ax^2+bx+c}{4x-d} = \dfrac{\left(\dfrac{2}{3}x-1\right)(4x-4)+k}{(4x-4)}$ mà $f(4)=-\dfrac{7}{4}$ nên $k=1$.\\
Vậy $f(x) = \dfrac{ax^2+bx+c}{4x-d} = \dfrac{\left(\tfrac{2}{3}x-1\right)(4x-4)+1}{(4x-4)} = \left(\dfrac{2}{3}x-1\right) + \dfrac{1}{(4x-4)}$.\\
Phương trình hoành độ giao điểm giữa $y=f(x)$ và parabol là
\[ \left(\dfrac{2}{3}x-1\right) + \dfrac{1}{(4x-4)} = -\dfrac{1}{3}x^2 + \dfrac{4}{3}x \Leftrightarrow -4x^3 + 12x^2 + 4x - 15 = 0 \Leftrightarrow \hoac{& x \approx 2{,}89868083208 = a \\& x \approx 1,18919299296 = b \\& x \approx -1{,}08787382504 \text{ (loại)}.} \]
Vậy $S_2 = \displaystyle\int\limits_{a}^{b} \left[ \left(\dfrac{2}{3}x-1\right) + \dfrac{1}{(4x-4)} - \left(-\dfrac{1}{3}x^2 + \dfrac{4}{3}x\right) \right] \mathrm{\,d}x \approx 0{,}94$.
}
\end{ex}

\begin{ex}%[2D4V3-1]
\immini{Một bước tường hình chữ nhật $ABCD$ có kích thước $6$ m $\times 4$ m, được bạn An trang trí bằng cách vẽ hai đồ thị $f(x) = a^x$, $g(x) = \log_b{x}$ đối xứng nhau qua đường thẳng $d \colon y= x$ và chia thành ba phần (tham khảo hình vẽ). Phần $H_1$ được sơn mà xanh da trời, phần $H_2$ được sơn màu vàng, phần $H_3$ được sơn màu xanh lá cây. Biết rằng mỗi hộp sơn các màu chỉ sơn được $3$ m$^2$ tường, đồng thời giá của hộp sơn màu xanh da trời là $100\,000$ đồng/hộp, hộp sơn màu vàng là $140\,000$ đồng/hộp, hộp sơn xanh lá cây là $130\,000$ đồng/hộp. Tính giá tiền mà bạn An mua để sơn bức tường này? (đơn vị là triệu đồng và của cửa hàng sơn chi bán số nguyên của hộp).}{
\begin{tikzpicture}[line join = round, line cap=round,>=stealth,font=\footnotesize,scale=1]
\path
(0,0) coordinate (O)
(-3,-2) coordinate (A)
(-3,2) coordinate (D)
(3,2) coordinate (C)
(3,-2) coordinate (B)
;
\draw
(A) -- (B) -- (C) -- (D) -- cycle
;
\begin{scope}
\clip
(-3,-2) rectangle (3,2);
\fill[yellow!30] (A) -- (B) -- (C) -- (D) -- (A)
;
\fill[blue!30] plot (\x,{3^(1/2*\x)});
\fill[green!30,domain=0.2:3] plot (\x,{ln(\x)/ln(sqrt(3))}) -- (3,-2);
\draw[samples=100,domain=-3:3,smooth,variable=\x] plot (\x,{3^(1/2*\x)});
\draw[samples=100,domain=0.1:3,smooth,variable=\x] plot (\x,{ln(\x)/ln(sqrt(3))});

\end{scope}
\fill
(O) circle(1pt) node[above left]{$O$}
(A) circle(1pt) node[below]{$A$}
(B) circle(1pt) node[below]{$B$}
(C) circle(1pt) node[above]{$C$}
(D) circle(1pt) node[above]{$D$}
(-2,1.5) node{$H_1$}
(2,-1.5) node{$H_3$}
(-2,-1.5) node{$H_2$}
(-3,0) circle(1pt) node[above left]{$-3$}
(3,0) circle(1pt) node[above right]{$3$}
(0,-2) circle(1pt) node[below left]{$-2$}
(0,2) circle(1pt) node[above left]{$2$}
;
\draw[->] (0,-3) -- (0,3) node[right]{$y$};
\draw[->] (-4,0) -- (4,0) node[below]{$x$};
\end{tikzpicture}
}
\par\shortans[0]{1{,}15}
\loigiai{
Từ hình vẽ ta thấy đồ thị hàm số $g(x) = \log_{b}x$ đi qua điểm $\left(3; 2 \right)$.\\
Nên $2 = \log_{b}3 \Rightarrow b^2 = 3 \Rightarrow b = \sqrt{3}$ (do $b > 0$).\\
Khi đó $g(x) = \log_{\sqrt{3}}x$.\\
Mặt khác đồ thị hàm số $f(x)$ và $g(x)$ đối xứng nhau qua $y = x$ nên $f(x) = \left(\sqrt{3} \right)^x$.\\
Hoành độ giao điểm của đồ thị hàm số $y = f(x)$ và $y = 2$ là nghiệm của phương trình\\ $f(x) = 2 \Leftrightarrow \left(\sqrt{3} \right)^2 = 2 \Leftrightarrow x = \log_{\sqrt{3}}2$.\\
Vậy diện tích phần $H_1$ là $S_1 = \displaystyle\int\limits_{-3}^{\log_{\sqrt{3}}2} \left(2 - \left(\sqrt{3} \right)^x \right)\,\mathrm{d}x \approx 5{,}23$ m$^2$.\\
Suy ra ta cần dùng $2$ hộp sơn màu xanh da trời.\\
Diện tích phần $H_3$ là $S_3 = \displaystyle\int\limits_{\frac{1}{3}}^3 \left(\log_{\sqrt{3}}x + 2 \right) \,\mathrm{d}x \approx 7{,}15$ m$^2$.\\
Suy ra ta cần dùng $3$ hộp sơn màu xanh lá cây.\\
Diện tích phần $H_2$ là $S_2 = 24 - S_1 - S_2 \approx 11{,}62$ m$^2$.\\
Suy ra ta cần dùng $4$ hộp sơn màu vàng.\\
Vậy tổng số tiền là $100\,000 \cdot 2 + 130\,000 \cdot 3 + 140\,000 \cdot 4 = 1\,150\,000$ đồng = $1{,}15$ triệu đồng.
}
\end{ex}

\begin{ex}%[2D4V3-1]%[ĐỀ ÔN THI TỐT NGHIỆP 2025 - SỐ 2- Hector Tran]%[THPT Lê Quý Đôn- Tỉnh Đồng Nai]
\immini[thm]{
Cho hai hàm số $f(x)=ax^3+bx^2+cx-2$ và $g(x)=dx^2+ex+2$ ($a$, $b$, $c$, $d$, $e\in \mathbb{R}$). Biết rằng đồ thị của hàm số $y=f(x)$ và $y=g(x)$ cắt nhau tại ba điểm có hoành độ lần lượt là $-2$; $-1$; $1$ (tham khảo hình vẽ). Tính diện tích hình phẳng giới hạn bởi hai đồ thị đã cho \textit{(làm tròn đến chữ số thập phân thứ nhất)}?
}{
\begin{tikzpicture}[scale=0.7, font=\footnotesize, line join=round, line cap=round,>=stealth,y=0.8cm]
\tikzset{every node/.style={scale=0.8}}
\fill[pattern=crosshatch,pattern color=gray!70] plot[smooth,samples=100,domain=-2:1] (\x,{-1*(\x)^2+2})--plot[smooth,samples=100,domain=1:-2] (\x,{2*(\x)^3+3*(\x)^2-2*(\x)-2});
\draw[->] (-2.5,0)--(2,0)node[below] {$x$};
\draw[->] (0,-2.5) -- (0,2.5)node[left] {$y$};
\draw[fill=black] (0,0) circle (1pt)node[above left,circle] {$O$};
\foreach \x/\g in {-2/90,-1/-90,1/-90}{
\draw[shift={(\x,0)}] (0pt,1pt)--(0pt,-1pt);
\path (\x,0) node[shift={(\g:0.3)}] {\footnotesize $\x$};}
\draw[smooth,samples=100,domain=-2.1:1.2]
plot(\x,{-1*(\x)^2+2})node[right]{$f(x)$};
\draw[smooth,samples=100,domain=-2.05:1.1]
plot(\x,{2*(\x)^3+3*(\x)^2-2*(\x)-2})node[right]{$g(x)$};
\draw[dashed] (-2,0) -- (-2,-2) (-1,0) -- (-1,1) (1,0) -- (1,1);
\end{tikzpicture}
}
\par
\shortans[oly]{6{,}2}
\loigiai{
Phương trình hoành độ giao điểm của đồ thị $f(x)$ và $g(x)$ là
\[ax^3+bx^2+cx-2=dx^2+3x+2\Leftrightarrow ax^3+(b-d)x^2+(c-e)x-4=0.\quad (\ast)\]
Vì đồ thị của hai hàm số cắt nhau tại ba điểm suy ra phương trình $(\ast)$ có ba nghiệm $x=-2$; $x=-1$; $x=1$ nên ta có
\[ax^3+(b-d)x^2+(c-e)x-4=a(x+2)(x+1)(x-1).\]
Khi đó $-4=-2a\Rightarrow a=2$.\\
Vậy diện tích hình phẳng cần tìm là
$\displaystyle\int \limits_{-2}^1\left|2(x+2)(x+1)(x-1)\right|\,\mathrm{d}x=\dfrac{37}{6}\approx 6{,}2$.
}
\end{ex}

\begin{ex}%[Hiếu Mai]%[2D4V3-1]
\immini[thm]
{
Cho đồ thị hàm số $y=\cos x$ và hình phẳng được tô màu như hình vẽ. Tính diện tích hình phẳng đó (viết kết quả dưới dạng số thập phân và làm tròn đến hàng phần mười).
}
{
\begin{tikzpicture}[scale=1,>=stealth, font=\footnotesize, line join=round, line cap=round]
\fill[black, opacity=0.4]plot[smooth, samples=100, domain=1:3](\x,{cos(\x r)})--plot[smooth, samples=100, domain=3:1](\x,{\x});
\draw[->](-1,0)--(4,0)node[below]{$x$};
\draw[->](0,-1)--(0,3.5)node[right]{$y$};
\draw[smooth, samples=100, domain=-1:3.5] plot(\x,{cos(\x r)})node[right]{$y=\cos x$};
\draw[dashed](1,0)--(1,1);
\draw(3,-1)--(3,3);
\draw[smooth, samples=100, domain=-1:3.5] plot(\x,{\x})node[right]{$y=x$};
\path
(0,0)node[below right]{$O$}
(0,1.15)node[left]{$1$}
(1,0)node[below]{$1$}
(3,0)node[below right]{$3$}
;
\end{tikzpicture}
}
\shortans[oly]{$4{,}7$}
\loigiai
{
Hình phẳng đã cho được giới hạn bởi các đồ thị hàm số $y=\cos x$, $y=x$ và hai đường thẳng $x=1$, $x=3$. \\
Khi đó, diện tích hình phẳng là
$S = \displaystyle\int\limits_{1}^{3} |x-\cos x|\mathrm{,d}x$. \\
Từ đồ thị ta có $x\geqslant\cos x$, $\forall x\in[1;3]$. \\
Do đó $S=\displaystyle\int\limits_{1}^{3} (x-\cos x)\mathrm{\,d}x=\left(\dfrac{x^2}{2}-\sin x\right)\bigg|_{1}^{3}=4-\sin3+\sin1 \approx4{,}7$.
}
\end{ex}

\begin{ex}%[Hiếu Mai]%[2D4V3-1]
\immini[thm]
{
Cho $g(x)=\displaystyle\int\limits_{0}^{x} f(t)\mathrm{\,d}t$ $(0\leqslant x\leqslant7)$ trong đó $f(t)$ là hàm số có đồ thị như hình vẽ. Tính $g(3)$.
}
{
\begin{tikzpicture}[scale=0.8,>=stealth, font=\footnotesize, line join=round, line cap=round]
\draw[color=gray!50,dashed](-1,-3)grid(8,5);
\draw[->](-1,0)--(8,0)node[below]{$t$};
\draw[->](0,-3)--(0,5)node[left]{$f$};
\path
(0,0)coordinate(O)node[below left]{$O$}
(0,2)coordinate(A)
(1,2)coordinate(B)
(1,0)coordinate(C)
(2,0)coordinate(D)
(2,4)coordinate(E)
(3,0)coordinate(F)
(5,-2)coordinate(G)
(6,-2)coordinate(H)
(7,0)coordinate(I)
(0,1)node[left]{$1$}
(1,0)node[below]{$1$}
(5,0)node[above]{$5$}
;
\draw[very thick, smooth, samples=300](A)--(B)--(E)--(F)--(G)--(H)--(I);
\end{tikzpicture}
}
\shortans[oly]{$7$}
\loigiai
{
\immini
{
Gọi các điểm $A(0;2)$, $B(1;2)$, $C(1;0)$, $D(2;0)$, $E(2;4)$, $F(3;0)$. \\
Ta có
\[\begin{aligned}
g(3)
&=\displaystyle\int\limits_{0}^{3} f(t)\mathrm{\,d}t=S_{OABC}+S_{BCDE}+S_{DEF} \\
&=OA\cdot OC+\dfrac{(BC+DE)\cdot CD}{2}+\dfrac{DE\cdot DF}{2} \\
&=2\cdot1+\dfrac{(2+4)\cdot1}{2}+\dfrac{4\cdot1}{2} \\
&=7.
\end{aligned}\]
}
{
\begin{tikzpicture}[scale=0.8,>=stealth, font=\footnotesize, line join=round, line cap=round]
\draw[color=gray!50,dashed](-1,-3)grid(8,5);
\draw[->](-1,0)--(8,0)node[below]{$t$};
\draw[->](0,-3)--(0,5)node[left]{$f$};
\path
(0,0)coordinate(O)
(0,2)coordinate(A)
(1,2)coordinate(B)
(1,0)coordinate(C)
(2,0)coordinate(D)
(2,4)coordinate(E)
(3,0)coordinate(F)
(5,-2)coordinate(G)
(6,-2)coordinate(H)
(7,0)coordinate(I)
;
\draw[thick, blue, smooth, samples=300](A)--(B)--(E)--(F)--(G)--(H)--(I);
\foreach \p/\q in {O/-135, A/-180, B/120, C/-90, D/-90, E/90, F/-90} \fill[black] (\p) circle (2pt) ($(\p)+(\q:3.5mm)$) node{$\p$};
\end{tikzpicture}
}
}
\end{ex}

\begin{ex}%[2D4V3-1]
Gọi $S$ là phần hình phẳng gạch chéo được tạo bởi hai đồ thị hàm số $y=\sin x+1$ và $y=\cos x+1$ như hình vẽ dưới. Biết thể tích của vật thể được tạo thành khi quay $S$ quanh trục Ox là $V=a\sqrt {b} \pi +c\pi^2$, với $a,b,c\in \mathbb{N}$, $b$ là số nguyên tố. Tính giá trị của biểu thức $P=a+b+c$.
\begin{center}
\begin{tikzpicture}[scale=0.8,>=stealth, font=\footnotesize, line join=round, line cap=round]
\def\xmin{-3} \def\xmax{6} \def\ymin{-1} \def\ymax{2.8}
\draw[->] (\xmin,0)--(\xmax,0) node [below]{$x$};
\draw[->] (0,\ymin)--(0,\ymax) node [right]{$y$};
\node at (0,0) [below right]{$O$};
\pgfmathsetmacro{\m}{pi/4}
\pgfmathsetmacro{\n}{5*pi/4}
\clip (\xmin+0.1,\ymin+0.1) rectangle (\xmax-0.5,\ymax-0.1);
\draw[smooth,samples=400,domain=\xmin:\xmax] plot(\x,{cos(\x r)+1});
\draw[smooth,samples=400,domain=\xmin:\xmax] plot(\x,{sin(\x r)+1});
\draw [draw=none, pattern = north west lines, pattern color=gray!100, samples=100]
plot[domain=\m:\n] (\x, {cos(\x r)+1})--plot[domain=\n:\m] (\x, {sin(\x r)+1})--cycle;
\node at (0,2.3) [left]{$2$};
\node at (0,1.3) [left]{$1$};
\node at (-0.5*pi,0) [below]{$-\frac{\pi}{2}$};
\node at (0.5*pi,0) [below]{$\frac{\pi}{2}$};
\node at (pi-0.1,0) [below]{$\pi$};
\node at (1.5*pi,0) [below]{$\frac{3\pi}{2}$};
\end{tikzpicture}
\end{center}
\shortans{6}
\loigiai{
Xét phương trình hoành độ giao điểm giữa hai đồ thị hàm số\\
\[\sin x+1=\cos x+1\Leftrightarrow \sin x-\cos x=0\Leftrightarrow \sqrt {2} \sin\left(x-\dfrac{\pi}{4}\right)=0\Leftrightarrow x=\dfrac{\pi}{4}+k\pi (k\in \mathbb{Z}).\]
Khi đó, phần hình phẳng màu xanh được giới hạn bởi $x=\dfrac{\pi}{4}$ và $x=\dfrac{5\pi}{4}$.\\
Khi $x\in \left[\dfrac{\pi}{4};\dfrac{5\pi}{4}\right]$, ta có $\sin x+1\geq \cos x+1\geq 0$.\\
Thể tích của vật thể tròn xoay khi đó bằng\\
$V=\pi \displaystyle\int\limits_\frac{\pi }{4}^\frac{5\pi }{4} \left| (\sin x+1)^{2}-(\cos x+1)^{2} \right| \mathrm{\,d}x=\pi \displaystyle\int\limits_\frac{\pi }{4}^\frac{5\pi }{4} \left( \sin^{2}x-\cos^{2}x+2\sin x-2\cos x \right)\mathrm{\,d}x$.\\
$V=\pi \cdot \displaystyle\int\limits_{\tfrac{\pi}{4}}^{\tfrac{5\pi}{4}} \left(-\cos 2x+2\sin x-2\cos x\right)\mathrm{\,d}x =\pi \left(\dfrac{-\sin2x}{2}-2\cos x-2\sin x\right)\Bigg\rvert_{\tfrac{5\pi}{4}}^{\tfrac{\pi}{4}}=4\pi \sqrt {2} $.\\
Khi đó, $a=4,b=2,c=0\Rightarrow P=a+b+c=6$.}
\end{ex}

\begin{ex}%[HSA-L1-2526, Huỳnh Xuân Tín]%[2D4V3-1]
\immini[thm]{Hình phẳng $(H)$ được giới hạn bởi đồ thị $(C)$ của hàm số đa thức bậc ba và parabol $(P)$ có trục đối
xứng vuông góc với trục hoành. Phần tô đậm như hình vẽ có diện tích bằng $\dfrac{a}{b}$ $(a, b \in \mathbb{Z};(a, b)=1)$.
Giá trị của biểu thức $a-3b$ bằng bao nhiêu?}{\begin{tikzpicture}[smooth,samples=300,scale=0.8,>=stealth,font=\footnotesize]
\def\xmin{-1.5} \def\xmax{3}  \def\ymin{-3}  \def\ymax{3}
%\draw[color=gray!50,dashed] (\xmin,\ymin) grid (\xmax,\ymax);
\draw[->] (\xmin,0)--(\xmax,0) node [below]{$x$};
\draw[->] (0,\ymin)--(0,\ymax) node [left]{$y$};
\clip (\xmin,\ymin) rectangle (\xmax,\ymax);
%\draw[->] (-2.8,0)--(3,0) node[below]{$x$};
%\draw[->] (0,-2.5)--(0,3) node[right]{$y$};
%\draw[pattern = north east lines,smooth] plot[domain=-1:1](\x,{(\x)^3-3*(\x)^2+2})--cycle;
\draw[name path=done,line width=0.8, red, domain=-1.5:3,smooth]
plot (\x,{(\x)^3-3*(\x)^2+2});
\draw[name path=dtwo,line width=0.8, blue, domain=-1.5:3,smooth]
plot (\x,{-(\x)^2+\x});
\fill[pattern = north east lines,smooth,draw=none] plot[domain=-1:2](\x,{(\x)^3-3*(\x)^2+2})--plot[domain=2:-1](\x,{-(\x)^2+(\x)})--cycle;
\draw[-][dashed] (-1,0)|- (0,-2) -| (2,0);
\foreach \p in{-1,1,2}
\draw[fill=black] (\p,0) circle(1pt) node[above]{$\p$} (0,-2) circle(1pt) node[below left]{$-2$} (0,2) circle(1pt) node[above right]{$2$} (2,-2) circle (1pt);
\draw[] (3,2) node[below left]{$(C)$} (2.3,-3) node[above left]{$(P)$};
\end{tikzpicture}}
\par \shortans{-5}
\loigiai{
Giả sử $(C)\colon y=f(x)$, $(P)\colon y=g(x)$.\\
Dựa vào đồ thị, $(C)$ và $(P)$ cắt nhau tại ba điểm có hoành độ lần lượt là $-1$, $1$, $2$.\\
Suy ra $f(x)-g(x)=a\cdot (x+1)(x-1)(x-2)$.\\
Với $x=0\Rightarrow f(0)-g(0)=2a=2\Rightarrow a=1$.\\
Vậy diện tích phần giới hạn bởi hai đồ thị $(C)$ và $(P)$ là
\[S=\displaystyle\int \limits_{-1}^2 \left|f(x)-g(x)\right|\mathrm{d} x=\displaystyle\int\limits_{-1}^2\left|(x+1)(x-1)(x-2)\right|\mathrm{d}x=\dfrac{37}{12}.\]
$\Rightarrow a=37$, $b=12\Rightarrow a-3b=1$.
}
\end{ex}

\begin{ex}%[HSA-L1-2526, Đào Trung Kiên]%[2D4V3-1]
Hình phẳng $(H)$ được giới hạn bởi đồ thị $(C)$ của hàm số đa thức bậc ba và parabol $(P)$ có trục đối xứng vuông góc với trục hoành. Phần tô đậm như hình vẽ có diện tích bằng $\dfrac{a}{b}$ ($a, b \in \mathbb{Z}; (a,b)=1$). Giá trị của biểu thức $a - 3b$ bằng bao nhiêu?
\begin{center}
\begin{tikzpicture}[scale=1,>=stealth, font=\footnotesize, line join=round, line cap=round, x=0.8 cm, y=0.5 cm]
\tikzset{declare function={
f(\x)=(\x)^3-3*(\x)^2+2;
g(\x)=-1*(\x)^2+(\x);}}
\def\xmin{-2.5}
\def\xmax{4}
\def\ymin{-3.5}
\def\ymax{4}
\draw[->] (\xmin,0) -- (0,0) node[below right]{$O$} --  (\xmax,0) node[below]{\footnotesize $x$};
\draw[->] (0,\ymin) -- (0,\ymax) node[left]{$y$};
\foreach \i in {(-1,0), (0,0), (1,0), (2,0), (0,-2)} \fill \i circle(1pt);
\begin{scope}
\clip (\xmin,\ymin) rectangle (\xmax,\ymax);
\fill [pattern = north east lines] plot[domain = -1:1] (\x,{f(\x)}) -- plot[domain = 1:-1] (\x,{g(\x)});
\fill [pattern = north west lines] plot[domain = 1:2] (\x,{f(\x)}) -- plot[domain = 2:1] (\x,{g(\x)});
\draw[smooth, samples=100, domain=\xmin: \xmax] plot (\x,{f(\x)})
plot (\x,{g(\x)});
\end{scope}
\draw[dashed](-1,0) node[above]{$-1$} --(-1,-2)--(2,-2)--(2,0)node[above]{$2$}
(1,0) node[above]{$1$}
(0,-2)node[below right]{$-2$};
\end{tikzpicture}
\end{center}
\shortans{1}

\loigiai{
Gọi $(C): y = ax^3 + bx^2 + cx + d \quad (a \ne 0)$.

Do $(C)$ cắt trục $Oy$ tại điểm có tung độ bằng 2 nên $d = 2$.

Vì $(C)$ đi qua 3 điểm $A(-1; -2)$, $B(1; 0)$ và $C(2; -2)$ nên ta được hệ phương trình:
\[
\heva{&	-a + b - c = -4 \\&a + b + c = -2 \\&4a + 2b + c = -2}
\Rightarrow
\heva{&a = 1\\&b = -3\\&c = 0}
\]

Do đó $(C): y = x^3 - 3x^2 + 2.$

Gọi $(P): y = mx^2 + nx + r \quad (m \ne 0)$.

Do $(P)$ đi qua 3 điểm $A(-1; -2)$, $O(0; 0)$ và $C(2; -2)$ nên ta được:
\[
\heva{&m - n + r = -2 \\&r = 0 \\&4m + 2n + r = -2}
\Rightarrow
\heva{&m = -1 \\&n = 1 \\&r = 0}
\]

Do đó $(P): y = -x^2 + x.$

Vậy
\[
S_{(H)} =\displaystyle \int_{-1}^{2} \left| x^3 - 2x^2 - x + 2 \right| \, \mathrm{d}x = \dfrac{37}{12}.
\]

Suy ra $a = 37$, $b = 12$ nên $a - 3b = 1$.
}
\end{ex}

\begin{ex}%[HSA-L1-2526]%[2D4V3-1]
Cho hàm số $y=f(x)$ có đồ thị như hình vẽ dưới. Tính giá trị của tích phân $I=\displaystyle\int\limits_{-6}^4\left[f(x)+x^2-x\right] \mathrm{\,d}x$. (nhập đáp án vào ô trống, kết quả làm tròn đến hàng đơn vị)
\begin{center}
\begin{tikzpicture}[line join=round, line cap=round,scale=1,transform shape]
\clip (-6.5,-3.3) rectangle (4.5,4.3);
\draw[gray!50] (-6.5,-3) grid (4.5,4);
\def\xmin{-6.5} \def\xmax{4.5}
\def\ymin{-3.3} \def\ymax{4.3}
\draw[->] (\xmin,0)--(\xmax,0) node [above left]{$x$};
\draw[->] (0,\ymin)--(0,\ymax) node [below right]{$y$};
\node at (0,0) [below left]{$O$};
\foreach \x in{-6,-5,-4,-3,-2,-1,1,2,3,4} \draw[fill=black] (\x,0) circle (.5pt) node[below]{$\x$};
\foreach \y in{-3,-2,-1,1,2,3,4}\draw[fill=black] (0,\y) circle (.5pt) node[left]{$\y$};
\draw[red, smooth,samples=100,domain=\xmin:0] plot(\x,{.5*(\x)+2})--(1,1) arc (180:0:1 cm)--++(1.5,0);
\end{tikzpicture}
\end{center}
\par\shortans{112}
\loigiai{
\begin{center}
\begin{tikzpicture}[line join=round, line cap=round,scale=1,transform shape]
\clip (-6.5,-3.3) rectangle (4.5,4.3);
\draw[gray!50] (-6.5,-3) grid (4.5,4);
\def\xmin{-6.5} \def\xmax{4.5}
\def\ymin{-3.3} \def\ymax{4.3}

\fill[red!50] (-6,0)--(-6,-1)-- plot[smooth,samples=100,domain=-6:-4](\x,{.5*(\x)+2});
\fill[blue!50] plot[smooth,samples=100,domain=-4:0](\x,{.5*(\x)+2})--(0,0);
\fill[green!50] (0,0)--(0,2)--(1,1)--(1,0);
\fill[yellow!50] (1,1) arc (180:0:1 cm)--++(0,-1)--++(-2,0);
\fill[purple!50] (3,1) --(4,1)--(4,0)--(3,0);
\draw[smooth,samples=100,domain=\xmin:0] plot(\x,{.5*(\x)+2})--(1,1) arc (180:0:1 cm)--++(1.5,0);
\draw[->] (\xmin,0)--(\xmax,0) node [above left]{$x$};
\draw[->] (0,\ymin)--(0,\ymax) node [below right]{$y$};
\node at (0,0) [below left]{$O$};
\foreach \x in{-6,-5,-4,-3,-2,-1,1,2,3,4} \draw[fill=black] (\x,0) circle (.5pt) node[below]{$\x$};
\foreach \y in{-3,-2,-1,1,2,3,4}\draw[fill=black] (0,\y) circle (.5pt) node[left]{$\y$};
\node at (-5.5,-.4) {$S_1$};
\node at (-1,.7) {$S_2$};
\node at (.5,.7) {$S_3$};
\node at (2,.7) {$S_4$};
\node at (3.5,.5) {$S_5$};
\end{tikzpicture}

\end{center}
Xét các diện tích $S_1$, $S_2$, $S_3$, $S_4$, $S_5$ được ký hiệu như hình.\\
$S_1=\dfrac{1\cdot 2}{2}=1$,\\
$S_2=\dfrac{2\cdot 4}{2}=4$,\\
$S_3=\dfrac{1\cdot(1+2)}{2}=\dfrac{3}{2}$,\\
$S_4=\dfrac{1}{2} \cdot \pi \cdot 1^2+1\cdot 2=2+\dfrac{\pi}{2}$,\\
$S_5=1\cdot 1=1$.\\
$\Rightarrow \displaystyle\int\limits_{-6}^4f(x) \mathrm{\,d}x=-S_1+S_2+S_3+S_4+S_5=-1+4+\dfrac{3}{2}+2+\dfrac{\pi}{2}+1=\dfrac{\pi+15}{2}$.\\
Lại có $\displaystyle\int\limits_{-6}^4\left(x^2-x\right) \mathrm{\,d}x=\left.\left(\dfrac{x^3}{3}-\dfrac{x^2}{2}\right)\right|_{-6} ^4=\dfrac{310}{3}$
$\Rightarrow I=\dfrac{\pi+15}{2}+\dfrac{310}{3} \approx 112{,}4$.
}
\end{ex}

\begin{ex}%[2D4V3-2]
Trong chương trình nông thôn mới, tại một xã Y có xây một đoạn đường hầm bằng bê tông như hình vẽ. Tính thể tích khối bê tông để đổ đủ đoạn đường hầm. (Đường cong trong hình vẽ là các đường Parabol).
\begin{center}
\begin{tikzpicture}[scale=1, font=\footnotesize, line join=round, line cap=round, >=stealth]
\path
(0,0) coordinate (O)
(0,1) coordinate (B)
(-1,0) coordinate (A)
(1,0) coordinate (C)
($(B)+(0,0.5)$) coordinate (E)
($(A)+(-0.5,0)$) coordinate (D)
($(C)+(0.5,0)$) coordinate (F)
;
\foreach \i in {O,A,B,C,D,E,F}
\path ($(\i)+(-6,-2)$) coordinate (\i');
\draw (E)--(E') (F)--(F')
(D')--(F') node[midway,below]{$19$};
\draw[dashed] %(O)--(O')
(A)--(A') (B)--(B') (D)--(D') (C)--(C')
(D)--(F)
;
\draw[<->] ($(A')-(0,0.2)$)--($(D')-(0,0.2)$) node[midway,below]{$0{,}5$};
\draw[<->] ($(F')-(0,0.2)$)--($(C')-(0,0.2)$) node[midway,below]{$0{,}5$};
\draw[<->] ($(A')-(0,0.1)$)--($(C')-(0,0.1)$) node[midway,below]{$19$};
\draw[<->] ($(O)+(0.1,0)$)--($(B)+(0.1,0)$) node[midway,right]{$2$};
\draw[<->] (B)--(E) node[midway,right]{$0{,}5$};
\begin{scope}
\draw[samples=350,domain=-1.5:1.5,smooth,variable=\x] plot (\x,{(-2/3)*(\x)^2+1.5});
\draw[samples=350,domain=-1:1,smooth,variable=\x] plot (\x,{-(\x)^2+1});
\draw[samples=350,domain=-7.5:-4.5,smooth,variable=\x] plot (\x,{(-2/3)*(\x+6)^2-0.5});
\draw[samples=350,domain=-7:-5,smooth,variable=\x] plot (\x,{-(\x+6)^2-1});
\end{scope}
\end{tikzpicture}
\end{center}
\shortans[]{$40$}
\loigiai{
Chọn hệ trục $Oxy$ như hình vẽ.\\
\begin{center}
\begin{tikzpicture}[scale=1, font=\footnotesize, line join=round, line cap=round, >=stealth]
\draw[->] (0,0)--(1.5,0) node[below]{$x$};
\draw[->] (0,0)--(0,2) node[left]{$y$};
\draw[fill=black] (0,0) circle(1pt) node[below left]{$O$};
\path
(0,0) coordinate (O)
(0,1) coordinate (B)
(-1,0) coordinate (A) node[above right]{$B$}
(1,0) coordinate (C) node[above right]{$A$}
($(B)+(0,0.5)$) coordinate (E)
($(A)+(-0.5,0)$) coordinate (D)
($(C)+(0.5,0)$) coordinate (F)
;
\foreach \i in {O,A,B,C,D,E,F}
\path ($(\i)+(-6,-2)$) coordinate (\i');
\draw (E)--(E') (F)--(F')
(D')--(F');
\draw[dashed] %(O)--(O')
(A)--(A') (B)--(B') (D)--(D') (C)--(C')
(D)--(F)
;
\draw[<->] ($(A')-(0,0.2)$)--($(D')-(0,0.2)$) node[midway,below]{$0{,}5$};
\draw[<->] ($(F')-(0,0.2)$)--($(C')-(0,0.2)$) node[midway,below]{$0{,}5$};
%\draw[<->] ($(O)+(0.1,0)$)--($(B)+(0.1,0)$) node[midway,right]{$2$};
%	\draw[<->] (B)--(E) node[midway,right]{$0{,}5$};
\begin{scope}
\draw[samples=350,domain=-1.5:1.5,smooth,variable=\x] plot (\x,{(-2/3)*(\x)^2+1.5});
\draw[samples=350,domain=-1:1,smooth,variable=\x] plot (\x,{-(\x)^2+1});
\draw[samples=350,domain=-7.5:-4.5,smooth,variable=\x] plot (\x,{(-2/3)*(\x+6)^2-0.5});
\draw[samples=350,domain=-7:-5,smooth,variable=\x] plot (\x,{-(\x+6)^2-1});
\end{scope}
\end{tikzpicture}
\end{center}
Gọi $(P_1)\colon y=a_1 x^2+b_1$ là Parabol đi qua hai điểm $A\left( \dfrac{19}{2};0 \right)$, $B(0;2)$.\\
Nên ta có hệ phương trình sau:
\[\heva{&0=a\left( \dfrac{19}{2} \right)^2+2\\&2=b}
\Leftrightarrow \heva{&a_1=-\dfrac{8}{361}\\&b_1=2}
\Rightarrow (P_1)\colon y=-\dfrac{8}{361}{x^2}+2.\]
Gọi $(P_2)\colon y=a_2x^2+b_2$ là Parabol đi qua hai điểm $C\left( 10;0 \right)$, $D\left( 0;\dfrac{5}{2} \right)$.\\
Nên ta có hệ phương trình sau:
\[\heva{&0=a_2\left( 10 \right)^2+\dfrac{5}{2}\\&\dfrac{5}{2}=b_2}
\Leftrightarrow \heva{&a_2=-\dfrac{1}{40}\\&b_2=\dfrac{5}{2}}
\Rightarrow (P_2)\colon y=-\dfrac{1}{40} x^2+\dfrac{5}{2}.\]
Ta có thể tích của bê tông là
\[ V=5\cdot 2\left[ \int_0^{10} \left( -\dfrac{1}{40}{x^2}+\dfrac{5}{2} \right)\mathrm{\,d}x- \int_0^{\tfrac{19}{2}} \left( -\dfrac{8}{361}{x^2}+2 \right)\mathrm{\,d}x \right]
=40~\text{m}^3.\]
}
\end{ex}

\begin{ex}%[2D4V3-2]%[Tổ 21 - Đợt 17 - Chương 4 - - Cánh Diều - Đề 2]
\immini{Cho một cái cổng hình Parabol như hình vẽ sau. Chiều cao $G H=4$m, chiều rộng $A B=4$m, $A C=B D=0{,}9$m. Chủ nhà làm hai cánh cổng khi đóng lại là hình chữ nhật $C D E F$ tô đậm có giá là $1\,200\,000$ đồng/m$^{2}$, còn các phần để trắng thì trang trí hoa có giá là $900\,000$ đồng/m$^{2}$. Hỏi tổng số tiền để làm hai phần nói trên cần ít nhất bao nhiêu? (triệu đồng)}{
\begin{tikzpicture}[scale=1, font=\footnotesize, line join=round, line cap=round, >=stealth]
\def\xt{-1} \def\xp{5} \def\yt{5} \def\yd{-1}
\draw[fill=cyan,opacity=0.5] (0.9,0)--(0.9,2.79)--(3.1,2.79)--(3.1,0)--(0.9,0);
\clip(\xt+0.1,\yd+0.1) rectangle(\xp+1,\yt-0.1);
\draw[smooth,samples=300,domain=0:4] plot(\x,{4*(\x)-(\x)^2});
\fill (0.9,2.79) node[ left]{$F$} circle(1pt);
\fill (3.1,2.79) node[right]{$E$} circle(1pt);
\fill (0.9,0) node[above left]{$C$} circle(1pt);
\fill (3.1,0) node[above right]{$D$} circle(1pt);
\fill (0,0) node[below right]{$A$} circle(1pt);
\fill (4,0) node[below right]{$B$} circle(1pt);
\fill (2,0) node[below]{$H$} circle(1pt);
\fill (2,4) node[above]{$G$} circle(1pt);
\draw[dashed] (2,4)--(2,0);
\draw(0,0)--(4,0);
\end{tikzpicture}}
\shortans{$12$}
\loigiai
{
\immini{Gắn hệ trục tọa độ $O x y$ sao cho $A B$ trùng $O x$, $A$ trùng với gốc $O$. Khi đó parabol có đỉnh $G(2 ; 4)$ và đi qua gốc tọa độ.\\
Giả sử phương trình của parabol có dạng $y=a x^{2}+b x+c$, $(a \neq 0)$.\\
Vì parabol có đỉnh là $G(2 ; 4)$ và đi qua điểm $O(0 ; 0)$ nên ta có
\[\heva{&c=0\\&-\dfrac{b}{2 a}=2\\&a \cdot 2^{2}+b \cdot 2+c=4}\Leftrightarrow\heva{&a=-1 \\&b=4\\&c=0.}\]
}{
\begin{tikzpicture}[scale=1.1, font=\footnotesize, line join=round, line cap=round, >=stealth]
\def\xt{-1} \def\xp{5} \def\yt{5} \def\yd{-1}
\draw[fill=cyan,opacity=0.5] (0.9,0)--(0.9,2.79)--(3.1,2.79)--(3.1,0)--(0.9,0);
\draw[->](\xt, 0)--(\xp, 0) node[below]{$x$};
\draw[->](0,\yd)--(0,\yt) node[left]{$y$};
\node at (0, 0) [below left]{$O$};
\clip(\xt+0.1,\yd+0.1) rectangle(\xp+1,\yt-0.1);
\draw[smooth,samples=300] plot(\x,{4*(\x)-(\x)^2});
\fill (0.9,2.79) node[ left]{$F$} circle(1pt);
\fill (3.1,2.79) node[right]{$E$} circle(1pt);
\fill (3.1,0) node[below]{$3{,}1$} circle(1pt);
\fill (0.9,0) node[below]{$0{,}9$} circle(1pt);
\fill (0.9,0) node[above left]{$C$} circle(1pt);
\fill (3.1,0) node[above right]{$D$} circle(1pt);
\fill (0,0) node[below right]{$A$} circle(1pt);
\fill (4,0) node[below right]{$B$} circle(1pt);
\fill (4,0) node[above right]{$4$} circle(1pt);
\fill (0,4) node[left]{$4$} circle(1pt);
\fill (2,0) node[below]{$H$} circle(1pt);
\fill (2,4) node[above]{$G$} circle(1pt);
\fill (4,4) node[above]{$f(x)=-x^2+4x$};
\draw[dashed] (0,4)--(2,4)--(2,0);
\end{tikzpicture}}
\noindent 		Suy ra phương trình parabol là $y=f(x)=-x^{2}+4 x$.\\
Diện tích của cả cổng là $S=\displaystyle\int_{0}^{4}\left(-x^{2}+4 x\right) \mathrm{d} x=\left.\left(-\dfrac{x^{3}}{3}+2 x^{2}\right)\right|_{0} ^{4}=\dfrac{32}{3}\mathrm{~m}^{2}$.\\
Mặt khác chiều cao $C F=D E=f(0{,}9)=2{,}79\mathrm{~m}$; $C D=4-2\cdot0{,}9=2{,}2 \mathrm{~m}$.\\
Diện tích hai cánh cổng là $S_{\text {CDEF }}=C D \cdot C F=6{,}138 \mathrm{~m}^{2}$.\\
Diện tích phần để trắng là $S_{1}=S-S_{C D E F}=\dfrac{32}{3}-6{,}138=\dfrac{6\,793}{1\,500} \mathrm{~m}^{2}$.\\
Vậy tổng số tiền để làm cổng là $6{,}138 \cdot 1\,200\,000+\dfrac{6\,793}{1\,500} \cdot 900\,000=11\,441\,400$ đồng.\\
Do đó cần ít nhất $12$ triệu đồng.
}
\end{ex}

\begin{ex}%[2D4V3-2]%[Tổ 21 - Đợt 17 - Chương 4 - - Cánh Diều - Đề 2]%[Dương Hồng]
\immini{Bác Bình muốn làm cửa rào sắt có hình dạng và kích thước như hình vẽ bên, biết đường cong phía trên là một Parabol. Giá $1  \mathrm{m}^{2}$ của rào sắt là $700\,000$ đồng. Hỏi bác Bình phải trả bao nhiêu tiền để làm cái cửa sắt như vậy (nghìn đồng).}{
\definecolor{apricot}{rgb}{0.98, 0.81, 0.69}
\begin{tikzpicture}[line join=round, line cap=round,scale=1.2,transform shape]

\tikzset{cong/.pic={
\def\N{
(-1.8,-1.85)--(-1.8,.25)
..controls +(30:.4) and +(175:.4) ..(0,.75)--(0,-1.85)--cycle
(-1.8,0)--(0,0)
(-1.8,-1.3)--(0,-1.3)
(-1.2,-1.2) rectangle (-.6,-.15);
;
}
\draw[fill=black!30] (-1.8,-1.8) rectangle (0,-1.3);
\draw[line width=.4mm,black]\N;
\draw[fill=black] (0,.75)--(0,-1.85)--(-.1,-1.85)--(-.1,.75)--cycle
;
%\fill[ecru!80] \N;
\draw[pattern=vertical lines,pattern color=black](-1.8,-1.3)rectangle(-1.3,0);
\draw[pattern=vertical lines,pattern color=black](-.5,-1.3)rectangle(0,0);
\draw[pattern=vertical lines,pattern color=black](-1.2,-.12)rectangle(-.6,0);
\draw[pattern=vertical lines,pattern color=black](-1.2,-1.3)rectangle(-.6,-1.2);
%Cột
\draw[fill=apricot,apricot,line width=.2mm](-2.3,-1.9)rectangle(-1.82,1.28);
%\draw[fill=apricot](-3.5,-1.9)rectangle(-2.3,-.5);
%======================viền cong
\def\V{
(-1.2,.2)
..controls +(-110:.4) and +(-35:.4) ..(-1.1,.4)
..controls +(150:.2) and +(100:.2) ..(-1.3,.1)
(-.6,.2)
..controls +(-110:.4) and +(-35:.4) ..(-.5,.4)
..controls +(150:.2) and +(100:.2) ..(-.7,.1)
(-.85,.2)
..controls +(-170:.2) and +(155:.2) ..(-.8,.3)
..controls +(-20:.2) and +(-55:.24) ..(-1.05,.1)
(-1.45,.2)
..controls +(-170:.2) and +(155:.2) ..(-1.4,.3)
..controls +(-20:.2) and +(-55:.6) ..(-1.8,.25)
(-.25,.2)
..controls +(-170:.2) and +(155:.2) ..(-.18,.3)
..controls +(-20:.2) and +(-55:.24) ..(-.45,.1)
(-.25,.5)
..controls +(-170:.2) and +(155:.2) ..(-.18,.7)
..controls +(-20:.2) and +(-65:.55) ..(-.45,.5)
(-.55,.52)
..controls +(-170:.2) and +(155:.2) ..(-.44,.65)
..controls +(-20:.2) and +(-65:.2) ..(-.75,.5)
(-.7,.36)
..controls +(-170:.2) and +(155:.2) ..(-.85,.55)
..controls +(-20:.2) and +(-55:.25) ..(-1.15,.5)
;
}
\draw[line width=.4mm,black]\V;

}}
%========================hoa, lá\tikzset{hoa/.pic={
\tikzset{hoa/.pic={
\draw (0,0) ..controls +(40:1) and +(170:1) .. (1.6,-2.4)
(0,0) ..controls +(40:1) and +(170:1) .. (1.6,2.6)
(0,0) ..controls +(40:1) and +(30:1) .. (-1.4,-2.6)
(0,0) ..controls +(120:1) and +(80:1) .. (-1.4,2.6);
\def\H{
(0,0)
..controls +(0:2) and +(60:2) .. (0,0)
..controls +(60:2) and +(120:2) .. (0,0)
..controls +(120:2) and +(180:2) .. (0,0)
..controls +(180:2) and +(240:2) .. (0,0)
..controls +(240:2) and +(300:2) .. (0,0)
..controls +(300:2) and +(360:2) .. (0,0)
;}
\fill[black] \H;
\draw \H;

}}
\tikzset{la/.pic={
\def\L{
(1.6,.95)
..controls +(30:.5) and +(160:1) ..  (2.4,1.3)
..controls +(20:0) and +(-160:0.3) ..  (2.9,1.3)
..controls +(-120:1.4) and +(-120:1) ..  (1.6,.95)
;}
\draw \L;
\fill[black!80] \L;
%Vẽ gân lá
\draw (1.35,1.2)
..controls +(-40:.1) and +(-160:1) ..  (2.5,.9) (2.1,1.15)--(1.7,.84)
..controls +(-70:.3) and +(170:0) ..  (2,.5)
(2.2,.82)--(2.5,1.1) (2.3,.78)--(2.5,.8);
}}
%==================Mái nhà
\definecolor{burntsienna}{rgb}{0.91, 0.45, 0.32}
\tikzset{mai/.pic={
\def\mainha{
(.5,2)
foreach \n in {1,2,...,30} { -- ++ (0,0) -- ++ (0,1) -- ++ (1,0) -- ++ (0,-1) } -- cycle

(.25,1)
foreach \n in {1,2,...,30} { -- ++ (0,0) -- ++ (0,1) -- ++ (1,0) -- ++ (0,-1) } -- cycle

(0,0)
foreach \n in {1,2,...,30} { -- ++ (0,0) -- ++ (0,1) -- ++ (1,0) -- ++ (0,-1) } -- cycle
;
}
\clip (-1.5,-3)--(26.5,-3)--(23.5,3)--(1.5,3)--cycle;
\draw[fill=burntsienna] \mainha;

}}



\path (-2.5,1.3)pic[scale=.2]{mai}
(0,0)pic[scale=1]{cong} (-.9,-.7)pic[scale=.2]{hoa} (-1,-1.7)pic[scale=.2]{la}
(-.75,-1.7)pic[rotate=-20,scale=.2]{la}
(-.4,-1.8)pic[rotate=140,scale=.2]{la}

(-1.4,-2)pic[rotate=35,scale=.2]{la}
(-.95,-1.9)pic[rotate=125,scale=.3]{la}
%-----------

(0,0)pic[scale=1,xscale=-1]{cong}
(.9,-.7)pic[scale=.2]{hoa}
(.89,-1.7)pic[rotate=-5,scale=.2,xscale=-1]{la}

(.72,-1.77)pic[rotate=10,scale=.2,xscale=-1]{la}

(.3,-1.85)pic[rotate=230,scale=.2,xscale=-1]{la}

(1.4,-2)pic[rotate=-35,scale=.2,xscale=-1]{la}
(.9,-1.8)pic[rotate=225,scale=.3,xscale=-1]{la}
;

\draw[dashed] (0,.77)--(2.6,.77);
\draw[<->,red] (2.6,-1.9)--(2.6,.77);
\draw[<->,red] (-1.8,-2.1)--(1.8,-2.1);
\draw[<->,red] (1.9,-1.9)--(1.9,.3);
\node at (0,-2.3) {\tiny $5$ m};
\node at (2.9,-.5) {\tiny $2$ m};
\node at (2.2,-.5) {\tiny $1{,}5$ m};
\end{tikzpicture}}
\shortans{$6417$}
\loigiai{
Ta chọn hệ trục tọa độ như hình vẽ.
\begin{center}
\begin{tikzpicture}[scale=0.7, font=\footnotesize, line join=round, line cap=round, >=stealth]
\def\xt{-4} \def\xp{4} \def\yt{3} \def\yd{-1}
\draw[->](\xt, 0)--(\xp, 0) node[below]{$x$};
\draw[->](0,\yd)--(0,\yt) node[left]{$y$};
\fill (0,0) node[below left]{$O$} circle(1.5pt);
\clip(\xt+0.1,\yd+0.1) rectangle(\xp-0.1,\yt-0.1);
\draw[color=blue, smooth, samples=300] plot(\x,{-0.08*(\x)^2+2});
\draw[color=red, smooth, samples=300] plot(-2.5,\x);
\draw[color=red, smooth, samples=300] plot(2.5,\x);
% Tô màu giữa hai đồ thị
\fill[fill=orange, opacity=0.5]
(-2.5,0) -- plot[samples=300, domain=-2.5:2.5] (\x,{-0.08*(\x)^2+2}) -- (2.5,0) -- cycle;
\fill (2.5,0) node[below left]{$2{,}5$} circle(1.5pt);
\fill (-2.5,0) node[below left]{$-2{,}5$} circle(1.5pt);
\fill (0,2) node[above left]{$2$} circle(1.5pt);
\fill (0,2) node[above right]{$C$} circle(1.5pt);
\fill (-2.5,1.5) node[above right]{$A$} circle(1.5pt);
\fill (2.5,1.5) node[above right]{$B$} circle(1.5pt);
\end{tikzpicture}
\end{center}
Trong đó $A(-2{,}5 ; 1{,}5)$, $B(2{,}5 ; 1{,}5)$, $C(0 ; 2)$.\\
Giả sử đường cong phía trên là một Parabol có dạng $y=a x^{2}+b x+c$, với $a$; $b$; $c \in \mathbb{Q}$.
Do Parabol đi qua các điểm $A(-2{,}5 ; 1{,}5)$, $B(2{,}5 ; 1{,}5)$, $C(0 ; 2)$ nên ta có hệ phương trình
\[
\heva{&a(-2{,}5)^{2} + b(-2{,}5) + c = 1{,}5\\
&a(2{,}5 )^{2} + b(2{,}5 ) + c = 1{,}5 \\
&	c = 2
}
\Leftrightarrow
\heva{&a=-\dfrac{2}{25} \\
&b=0\\
&	c=2.}
\]
Khi đó phương trình Parabol là $y=-\dfrac{2}{25} x^{2}+2$.\\
Diện tích $S$ của cửa rào sắt là diện tích phần hình phẳng giới hạn bởi đồ thị hàm số $y=-\dfrac{2}{25} x^{2}+2$, trục hoành và hai đường thẳng $x=-2{,}5$; $x=2{,}5$.\\
Ta có $S=\displaystyle\int_{-2{,}5}^{2{,}5}\left(-\dfrac{2}{25} x^{2}+2\right) \mathrm{d} x=\left.\left(-\dfrac{2}{25} \dfrac{x^{3}}{3}+2 x\right)\right|_{-2{,}5} ^{2{,}5}=\dfrac{55}{6}$.\\
Vậy bác Bình phải trả số tiền để làm cái cửa sắt là
\[S \times 700\,000=\dfrac{55}{6} \times 700\,000 \approx 6\,416 \,667 \text{( đồng).}\]
Hay $6\,417$ nghìn đồng.
}
\end{ex}

\begin{ex}%[12-MH-1-MH2025]%[MH-2025, Mã/Tên TV biên soạn]%[2D4V3-2]
Bạn Hải nhận thiết kế logo hình con mắt (phần gạch sọc như hình vẽ) cho một cơ sở y tế: Logo là hình phẳng giới hạn bởi hai parabol $y=f(x)$ và $y=g(x)$ như hình vẽ (đơn vị trên mỗi trục toạ độ là decimét).
\begin{center}
\begin{tikzpicture}[font=\footnotesize,line join=round, line cap=round, >=stealth,scale=0.8]
\def \xmin{-5}\def \xmax{5}\def \ymin{-2}\def \ymax{3.2}
\draw[->] (\xmin,0)--(\xmax,0) node[shift=(-110:0.2)] {$x$};
\draw[->] (0,\ymin)--(0,\ymax) node[shift=(-30:0.2)] {$y$};
\fill (0,0) circle(1pt) node[shift=(-135:0.25)]{$O$}
(-2,0) circle(1pt) node[shift=(-130:0.23)]{$-2$}
(2,0) circle(1pt) node[shift=(-90:0.2)]{$2$}
(0,-1) circle(1pt) node[shift=(-150:0.28)]{$-1$}
(0,1) circle(1pt) node[shift=(180:0.2)]{$1$}
(0,2) circle(1pt) node[shift=(135:0.22)]{$2$}
(2,1) circle(1pt);
\draw[dashed] (2,0)|-(0,1);
\clip (\xmin,\ymin) rectangle (\xmax,\ymax);
\fill[pattern=north west lines,opacity=.6] plot[domain={-sqrt(6)}:{sqrt(6)}](\x,{(\x)^(2)/4-1})--plot[domain={sqrt(6)}:{-sqrt(6)}](\x,{(-(\x)^(2))/4+2});
\draw[smooth,samples=100,domain=\xmin:\xmax] plot(\x,{(\x)^(2)/4-1});
\draw[smooth,samples=100,domain=\xmin:\xmax] plot(\x,{(-(\x)^(2))/4+2});
\node[rotate=-60] at (-3.3,2.4) {$y=f(x)$};
\node[rotate=60] at (-3.9,-1.2) {$y=g(x)$};
\end{tikzpicture}
\end{center}
Bạn Hải cần tính diện tích của logo để báo giá cho cơ sở y tế đó trước khi kí hợp đồng. Diện tích của logo là bao nhiêu decimét vuông (làm tròn kết quả đến hàng phần mười)?
\shortans{$9{,}8$}
\loigiai{
Gọi parabol $y=f(x)$ có dạng $f(x)=ax^2+bx+c$. \\
Parabol $y=f(x)$ nhận $Oy$ làm trục đối xứng nên ta có $\dfrac{-b}{2 a}=0 \Leftrightarrow b=0$. \\
Lại có đồ thị hàm số $y=f(x)$ đi qua điểm $(0;-1)$ và điểm $(2;0)$ nên $\heva{&a=\dfrac{1}{4}\\& c=-1.}$\\
Suy ra $y=f(x)=\dfrac{1}{4} x^2-1$.\\
Tương tự, ta cũng có $y=g(x)=-\dfrac{1}{4} x^2+2$.\\
Phương trình hoành độ giao điểm của hai parabol là
\begin{align*}
\dfrac{1}{4} x^2-1=-\dfrac{1}{4} x^2+2 \Leftrightarrow \hoac{&x=\sqrt{6} \\&x=-\sqrt{6}.}
\end{align*}
Vậy, diện tích của logo là
\begin{align*}
S & =\displaystyle\int\limits_{-\sqrt{6}}^{\sqrt{6}}\left[\left(-\dfrac{1}{4} x^2+2\right)-\left(\dfrac{1}{4} x^2-1\right)\right] \mathrm{d} x \\
& =\displaystyle\int\limits_{-\sqrt{6}}^{\sqrt{6}}\left(3-\dfrac{1}{2} x^2\right) \mathrm{d} x=\left.\left(3 x-\dfrac{x^3}{6}\right)\right|_{-\sqrt{6}} ^{\sqrt{6}}=4 \sqrt{6} \approx 9{,}8\left(\mathrm{dm}^2\right) .
\end{align*}
}
\end{ex}

\begin{ex}%[2D4V3-2]
\immini
{
Một gia đình muốn làm cánh cổng (như hình vẽ). Phần phía trên cổng có hình dạng là một parabol với $IH=2{,}5$ m, phần phía dưới là một hình chữ nhật kích thước cạnh là $AD=4$ m, $A B=6$ m. Giả sử giá để làm phần cổng được tô màu là $1\,000\,000 \text{~đồng} / \mathrm{m}^2$ và giá để làm phần cổng phía trên là $1\,200\,000 \text{~đồng} / \mathrm{m}^2$. Số tiền tổng cộng gia đình cần trả là bao nhiêu triệu đồng?
}
{
\begin{tikzpicture}[scale=0.7, font=\footnotesize, line join=round, line cap=round, >=stealth]
\draw (-3,-2.5) parabola bend (0,0) (3,-2.5);
\fill (0,0)coordinate (I)  circle (1pt) node[above]{$I$};
\fill (0,-2.5)coordinate (H)  circle (1pt) node[above right]{$H$};
\fill (-3,-2.5)coordinate (A)  circle (1pt) node[left]{$A$};
\fill (3,-2.5)coordinate (B)  circle (1pt) node[right]{$B$};
\fill (-3,-6.5)coordinate (D)  circle (1pt) node[left]{$D$};
\fill (3,-6.5)coordinate (C)  circle (1pt) node[right]{$C$};
\draw (I)--(H) node[midway, right]{$2{,}5$ m};
\draw (A)--(D) node[midway, left]{$4$ m};
\draw (D)--(C) node[midway, below]{$6$ m};
\draw[pattern = north east lines] (-3,-2.5) rectangle (3,-6.5);
\end{tikzpicture}
}
\shortans{$36$}
\loigiai{
\begin{center}
\begin{tikzpicture}[scale=0.8, font=\footnotesize, line join=round, line cap=round, >=stealth]
\draw (-3,-2.5) parabola bend (0,0) (3,-2.5);
\draw (-11,-2.5) parabola bend (-8,0) (-5,-2.5);
\draw[->] (-8,-5)--(-8,1)node[right]{$y$};
\draw[->] (-12,0)--(-4,0)node[above]{$x$};
\fill (-8,0) circle (1pt) node[above right]{$O$};
\fill (-11,-2.5) circle (1.5pt) node[left]{$A$};
\fill (-5,-2.5) circle (1.5pt) node[right]{$B$};
\draw[dashed] (-11,0)--(-11,-2.5)--(-8,-2.5);
\draw[dashed] (-5,0)--(-5,-2.5)--(-8,-2.5);
\fill (0,0)coordinate (I)  circle (1pt) node[above]{$I$};
\fill (0,-2.5)coordinate (H)  circle (1pt) node[above right]{$H$};
\fill (-3,-2.5)coordinate (A)  circle (1pt) node[left]{$A$};
\fill (3,-2.5)coordinate (B)  circle (1pt) node[right]{$B$};
\fill (-3,-6.5)coordinate (D)  circle (1pt) node[left]{$D$};
\fill (3,-6.5)coordinate (C)  circle (1pt) node[right]{$C$};
\draw (I)--(H) node[midway, right]{$2.5$ m};
\draw (A)--(D) node[midway, left]{$4$ m};
\draw (D)--(C) node[midway, below]{$6$ m};
\draw[pattern = north east lines] (-3,-2.5) rectangle (3,-6.5);
\fill (-11,0) circle (1pt) node[above]{$-3$};
\fill (-5,0) circle (1pt) node[above]{$3$};
\fill (-8,-2.5) circle (1pt) node[below left]{$-2{,}5$};
\end{tikzpicture}
\end{center}
Diện tích hình chữ nhật $ABCD$ là $S_{ABCD}=24 \text{~m}^2$.\\
Chọn hệ trục toạ độ như hình vẽ. Khi đó Parabol đi qua $3$ điểm $A$, $O$, $B$ có dạng $(P)\colon y=ax^2$.\\
Mặt khác $B(3;-2{,}5)\in (P)\Rightarrow 9a=-2{,}5\Leftrightarrow a=-\dfrac{5}{18}$.\\
Do đó $(P)\colon y=-\dfrac{5x^2}{18}$.\\
Diện tích cổng hình Parabol là \[S=\displaystyle\int\limits_{-3}^3 \left[ \left( -\dfrac{5x^2}{18}\right)-(-2{,}5)\right]  \mathrm{\,d}x=\left( -\dfrac{5x^3}{54}+\dfrac{5x}{2}\right) \bigg|_{-3}^3=10.\]
Vậy tổng số tiền làm cổng là $T=24\cdot 1\,000\,000+10\cdot 1\,200\,000=36$ triệu đồng.
}
\end{ex}

\begin{ex}%[2D4V3-2]
\immini{Bạn An cần mua một chiếc gương có viền là đường parabol bậc $2$ (xem hình vẽ). Biết rằng đoạn $AB=60$ cm, $OH=30$ cm. Diện tích của chiếc gương bạn An mua bằng bao nhiêu cm$^2$?
\shortans{$1200$}
}{\begin{tikzpicture}[scale=1,font=\footnotesize,line join = round, line cap = round, >= stealth]
%\tkzDefPoints{0/0/A, 5/0/x, 1.5/1/y, 0/5/z}
\draw[thick] (-3,0) parabola bend (0,3) (3,0);
\draw[thick,<->] (-3,0)--(3,0);
\draw[thick,<->] (0,0)--(0,3);
\draw (0,3) node[above]{H} (-3,0) node[below]{A} (3,0) node[below]{B} (0,0) node[above left]{O};
\draw (0,0) node[below]{$60$ cm};
\draw[](0,1.5) node[left]{$30$ cm};
%\clip (-3.5,-1) rectangle (3.5,3.5);
\end{tikzpicture}}
\loigiai{
\immini{Xét hệ trục tọa độ $Oxy$, sao cho gốc $O$ trùng với điểm $O$, trục $Ox$ trùng với tia $OB$, trục $Oy$ trùng với tia $OH$.\\
Khi đó $B(30;0)$, $H(0;30)$.\\
Phương trình Parabol đi qua $A$, $H$, $B$ có dạng\\ $y=ax^2+b\quad(P)$.
}{\begin{tikzpicture}[scale=.85,font=\footnotesize,line join = round, line cap = round, >= stealth]
%\tkzDefPoints{0/0/A, 5/0/x, 1.5/1/y, 0/5/z}
\draw[thick] (-3,0) parabola bend (0,3) (3,0);
\draw[thick,->] (-3.2,0)--(3.5,0) node[below]{$x$};
\draw[thick,->] (0,-0.5)--(0,3.5) node[right]{$y$} ;
\draw (0,3) node[above left]{H} (-3,0) node[below]{A} (3,0) node[below]{B} (0,0) node[below left]{O};
%\draw (0,0) node[below]{$60$ cm};
%\draw[](0,1.5) node[left]{$30$ cm};
%\clip (-3.5,-1) rectangle (3.5,3.5);
\end{tikzpicture}}
\noindent $H\in (P)\Rightarrow 30=b.\quad(1)$\\
$B\in (P)\Rightarrow 0=a\cdot 30^2+30\Rightarrow a =-\dfrac{1}{30}.\quad(2)$\\
Từ $(1)$, $(2)\Rightarrow (P)\colon y=-\dfrac{1}{30}x^2+30$.\\
Diện tích của chiếc gương là
\[\displaystyle\int\limits_{-30}^{30}\left(-\dfrac{1}{30}x^2+30\right)\mathrm{\,d}x=\left.\left(-\dfrac{x^3}{90}+30x\right)\right|_{-30}^{30}=1200\, \text{(cm$^2$)}.\]
}
\end{ex}

\begin{ex}%[2D4V3-2]
\immini{
Một mảnh vườn hình elip có trục lớn bằng $100$ m, trục nhỏ bằng $80$ m được chia thành $2$ phần bởi một đoạn thẳng nối hai đỉnh liên tiếp của elip. Phần nhỏ hơn trồng cây con và phần lớn hơn trồng rau. Hỏi diện tích trồng rau của mảnh vườn là bao nhiêu? (Kết quả làm tròn đến hàng đơn vị).}
{\begin{tikzpicture}[>=stealth,scale=0.6]
\draw (0,0) ellipse (5cm and 4cm);
\draw(-6,0) -- (6,0);
\draw(0,-5) -- (0,5);
%			\draw (0,0) node[below left]{\footnotesize $O$};
\draw (0,4)--(5,0);
\draw (5,0) node[below right]{\footnotesize $A$};
\draw (0,4) node[above left]{\footnotesize $B$};
\fill (0,0) circle (2pt);
%\fill (0,-3) circle (2pt) node[below right] {\footnotesize $-\dfrac{3}{2}$};
%\fill (0,3) circle (2pt) node[above right] {\footnotesize $\dfrac{3}{2}$};
\end{tikzpicture}}
\shortans{$5712$}
\loigiai{\immini{
Chọn hệ trục tọa độ $Oxy$ như hình vẽ.\\
Theo giả thiết phương trình của Ellip là\\
$\dfrac{x^2}{2500}+\dfrac{y^2}{1600}=1\Leftrightarrow y=\dfrac{4}{5}\sqrt{2500-x^2}$ m$^2$.\\
Diện tích của cả khu vườn là
\[S=4\displaystyle \int\limits_{0}^{50}\dfrac{4}{5}\sqrt{2500-x^2}\mathrm{\,d}x=2000\pi\,\text{m}^2.\]
Diện tích phần trồng cây con là
\[\begin{aligned}[t]
S_1&=\displaystyle \int\limits_{0}^{50}\dfrac{4}{5}\sqrt{2500-x^2}\mathrm{\,d}x-S_{OAB}\\
&=500\pi-\dfrac{1}{2}\cdot 40\cdot 50=500\pi-1000\,\text{m}^2.
\end{aligned}\]
Diện tích phần trồng rau là\\
\[S_2=S-S_1=3\cdot500\pi+1000\approx 5712\,\text{m}^2.\]}{\begin{tikzpicture}[>=stealth,scale=0.6]
\draw (0,0) ellipse (5cm and 4cm);
\draw[->] (-6,0) -- (6,0) node[above] {$x$};
\draw[->] (0,-5) -- (0,5) node[right] {$y$};
\draw (0,0) node[below left]{\footnotesize $O$};
\draw (0,4)--(5,0);
\draw (5,0) node[below right]{\footnotesize $A$};
\draw (0,4) node[above left]{\footnotesize $B$};
%				\fill (0,0) circle (2pt);
%\fill (0,-3) circle (2pt) node[below right] {\footnotesize $-\dfrac{3}{2}$};
%\fill (0,3) circle (2pt) node[above right] {\footnotesize $\dfrac{3}{2}$};
\end{tikzpicture}}
}
\end{ex}

\begin{ex}%[2D4V3-2]
\immini{Một viên gạch hoa hình vuông cạnh $40$ cm. Người ta đã dùng bốn đường parabol có chung đỉnh tại tâm của viên gạch để tạo ra bốn cánh hoa (phần tô đậm như hình vẽ). Diện tích của mỗi cánh hoa đó bằng bao nhiêu? (Kết quả làm tròn đến hàng đơn vị).
\shortans[0]{$133$}}{\begin{tikzpicture}[scale=1.2,>=stealth, line join = round, line cap = round,font=\footnotesize]
\draw (-1,1) rectangle (1,-1);
\foreach \i in {0,90,180,270}
\draw [fill=gray, smooth, samples=100,rotate=\i] plot [domain=0:1] (\x, {(\x)^2})-- plot [domain=1:0] (\x, {sqrt(\x)});
\draw [<-] (-1,-1.5)--(0,-1.5) node[below] {$40$ cm};
\draw [->] (0,-1.5)--(1,-1.5);
\end{tikzpicture}}
\loigiai{
\immini{Đặt hình vuông vào hệ trục tọa độ $Oxy$ như hình vẽ, cạnh của hình vuông ta có thể coi có độ dài bằng $2$ trên hệ trục tọa độ.\\
Khi đó diện tích một cánh hoa chính là diện tích hình phẳng giới hạn bởi hai đồ thị $y=x^2$ và $y=\sqrt{x}$. Ta có
\[S_1=\displaystyle\int\limits_0^1 (\sqrt{x}-x^2)\mathrm{\,d}x=\left. \dfrac{2}{3}x\sqrt{x}-\dfrac{x^3}{3}\right|_0^1=\dfrac{1}{3}.\]
Vậy thực tế diện tích của mỗi cánh hoa là
\[S=20\cdot 20\cdot S_1=\dfrac{400}{3}\approx 133 \,\text{cm}^2.\] }{\begin{tikzpicture}[scale=1.2,>=stealth, line join = round, line cap = round,font=\footnotesize]
\draw[->] (-1.5,0)--(0,0)%
node[below right]{$O$}--(1.5,0) node[below]{$x$};
\draw[->] (0,-1.5) --(0,1.5) node[right]{$y$};
\draw (-1,1) rectangle (1,-1);
\foreach \i in {-1,1}
\foreach \j in {-1,1}
\draw[fill=black]  (\i,0) circle (0.5 pt) node [below] {\footnotesize $\i$};
\draw [fill=gray, smooth, samples=100] plot [domain=0:1] (\x, {(\x)^2})-- plot [domain=1:0] (\x, {sqrt(\x)});
\end{tikzpicture}}
}
\end{ex}

\begin{ex}%[2D4V3-2]
\immini{Một chiếc cổng có hình dạng là một Parabol có khoảng cách giữa hai chân cổng là $AB=8$ m. Người ta treo một tấm phông hình chữ nhật có hai đỉnh $M,\,N$ nằm trên Parabol và hai đỉnh $P,\,Q$ nằm trên mặt đất (như hình vẽ). Ở phần phía ngoài phông (phần không tô đen) người ta mua hoa để trang trí, biết $MN=4$ m, $MQ=6$ m. Hỏi diện tích trang trí hoa chiếc cổng bằng bao nhiêu? (Kết quả làm tròn đến hàng phần mười).
\shortans[0]{$18{,}7$}}
{
\begin{tikzpicture}[font=\footnotesize, line join=round, line cap=round, >=stealth]
\begin{scope}[xscale=0.5,yscale=0.5]
\draw[samples=100,domain=-4:4]plot(\x,{-0.5*(\x)^2+8});
\path (-2,6)coordinate(M)(2,6)coordinate(N)(-2,0)coordinate(Q)(2,0)coordinate(P)(-4,0)coordinate(A)(4,0)coordinate(B);
\fill[gray!60,draw=black](M)--(N)--(P)--(Q)--cycle;
\draw(A)--(B);
\foreach \d/\g in{A/180,B/0,M/180,N/0,P/-90,Q/-90}
\draw(\d)circle(0.7pt)node[shift={(\g:0.35)}]{$\d$};
\end{scope}
\end{tikzpicture}
}
\loigiai{
\immini{Chọn hệ trục tọa độ như hình vẽ.\\
Parabol đối xứng qua $Oy$ nên có dạng $(P)\colon y=ax^2+c$.\\
Vì $(P)$ đi qua $B(4;0)$ và $N(2;6)$ nên $(P)\colon y=-\dfrac{1}{2}x+8$.\\
Diện tích hình phẳng giới hạn bởi $(P)$ và $Ox$ là
\[S=2\displaystyle\int\limits_{0}^{4} {\left(-\dfrac{1}{2}x^2+8\right)}\mathrm{\,d}x=\dfrac{128}{3}\mathrm{m}^2.\]
Diện tích phần trang trí hoa là \[S=S_1-S_{MNPQ}=\dfrac{128}{3}-24=\dfrac{56}{3}\approx18{,}7\,\text{m}^2.\]
}
{\begin{tikzpicture}[font=\footnotesize, line join=round, line cap=round, >=stealth]
\begin{scope}[xscale=0.5,yscale=0.5]
\draw[samples=100,domain=-4:4]plot(\x,{-0.5*(\x)^2+8});
\path (-2,6)coordinate(M)(2,6)coordinate(N)(-2,0)coordinate(Q)(2,0)coordinate(P)(-4,0)coordinate(A)(4,0)coordinate(B);
\fill[gray!60,draw=black](M)--(N)--(P)--(Q)--cycle;
\draw(A)--(B) (0,0) node[below left]{$O$};
\foreach \d/\g in{A/135,B/45,M/135,N/0,P/45,Q/135}
\draw(\d)circle(0.7pt)node[shift={(\g:0.35)}]{$\d$};
\draw [->](-5,0)--(5,0)node[below]{$x$};
\draw [->](0,-1)--(0,9)node[left]{$y$};
\foreach \x in{-4,-2,2,4}
\draw(\x,0)circle(0.7pt)node[below]{$\x$};
\end{scope}
\end{tikzpicture}}
}
\end{ex}

\begin{ex}%[2D4V3-2]
\immini{Một biển quảng cáo có dạng hình elip với bốn đỉnh $A_1$, $A_2$, $B_1$, $B_2$ như hình vẽ bên. Người ta cần sơn màu phần tô đậm như trong hình vẽ. Hỏi diện tích để sơn gần bằng bao nhiêu? Biết $A_1A_2=8$ m, $B_2=6$ m và tứ giác $MNPQ$ là hình chữ nhật có $MQ=3$ m? (Kết quả làm tròn đến hàng phần muời).
\shortans[0]{$35{,}5$}
}{
\begin{tikzpicture}[>=stealth,scale= 1.2,line width=0.6pt,xscale=1,yscale=1]
\def\R{2}
\def\r{1.36689}
\tkzDefPoints{0/0/O,-2/0/A_1,2/0/A_2}
\coordinate (N) at ($({\R*cos(40)},{\r*sin(40)})$);
\coordinate (P) at ($({\R*cos(-40)},{\r*sin(-40)})$);
\coordinate (B_2) at ($({\R*cos(90)},{\r*sin(90)})$);
\coordinate (B_1) at ($({\R*cos(-90)},{\r*sin(-90)})$);
\coordinate (M) at ($({\R*cos(140)},{\r*sin(140)})$);
\coordinate (Q) at ($({\R*cos(-140)},{\r*sin(-140)})$);
\draw[line width=0.6pt] (A_2) arc (0:360:\R cm and \r cm);
\fill[color=gray!10!black,shading=axis,opacity=0.2]
(Q) arc (220:320:\R cm and \r cm)--(N) arc (40:140:\R cm and \r cm)--(Q);
\draw[dashed,line width=0.6pt] (Q)--(P) (M)--(N);
\draw[line width=0.6pt] (M)--(Q) (N)--(P);
\foreach \x/\g in {A_1/-135,B_1/-135,A_2/-45,B_2/135,M/180,N/0,P/0,Q/180}
\fill[black] 	(\x) circle (1pt)
($(\g:3mm)+(\x)$) node {$\x$};
\end{tikzpicture}}
\loigiai{
\begin{center}
\begin{tikzpicture}[>=stealth,scale= 1.2,line width=0.6pt,xscale=1,yscale=1]
\draw[->] (-2.5,0) -- (3,0) node[below] { $x$};
\draw[->] (0,-2) -- (0,2) node[below right] {$y$};
\draw (0,0) node[below right] {\footnotesize $O$};
\def\R{2}
\def\r{1.36689}
\tkzDefPoints{0/0/O,-2/0/A_1,2/0/A_2}
\coordinate (N) at ($({\R*cos(40)},{\r*sin(40)})$);
\coordinate (P) at ($({\R*cos(-40)},{\r*sin(-40)})$);
\coordinate (B_2) at ($({\R*cos(90)},{\r*sin(90)})$);
\coordinate (B_1) at ($({\R*cos(-90)},{\r*sin(-90)})$);
\coordinate (M) at ($({\R*cos(140)},{\r*sin(140)})$);
\coordinate (Q) at ($({\R*cos(-140)},{\r*sin(-140)})$);
\draw[line width=0.6pt] (A_2) arc (0:360:\R cm and \r cm);
\fill[color=blue,opacity=0.2]
(Q) arc (220:320:\R cm and \r cm)--(N) arc (40:140:\R cm and \r cm)--(Q);
\draw[dashed,line width=0.6pt] (Q)--(P) (M)--(N);
\draw[line width=0.6pt] (M)--(Q) (N)--(P);
\foreach \x/\g in {A_1/-135,B_1/-135,A_2/-45,B_2/135,M/180,N/0,P/0,Q/180}
\fill[black] 	(\x) circle (1pt)
($(\g:3mm)+(\x)$) node {$\x$};
\end{tikzpicture}
\end{center}
Giả sử phương trình chính tắc của elip $(E)\colon \dfrac{x^2}{a^2}+\dfrac{y^2}{b^2}=1,\,(a>0,b>0)$.\\
Theo giả thiết ta có $\heva{&A_1A_2=8\\&B_1B_2=6}\Leftrightarrow \heva{&2a=8\\&2b=6}\Leftrightarrow \heva{&a=4\\&b=3.}$\\
$\Rightarrow (E)\colon \dfrac{x^2}{16}+\dfrac{y^2}{9}=1\Rightarrow y=\pm \dfrac{3}{4}\sqrt{16-x^2}$.\\
Diện tích phần sơn màu là $\displaystyle S=2\int_{-2\sqrt{3}}^{2\sqrt{3}}\dfrac{3}{4}\sqrt{16-x^2}\mathrm{\,d}x =8\pi+6\sqrt{3}\approx 35{,}5$ m$^2$.\\
}
\end{ex}

\begin{ex}%[2D4V3-2]
Một người thiết kế một vật trang trí hình vuông $ OABC $ có độ dài cạnh bằng $4$ được chia thành hai phần bởi đường cong parabol có đỉnh là $O$ và đi qua $B$. Gọi $ S_1 $, $ S_2 $ lần lượt là diện tích của phần không bị gạch và phần bị gạch (tham khảo hình vẽ bên dưới). Tỉ số $ \dfrac{S_1}{S_2} $ bằng
\shortans[0]{$2$}
\begin{center}
\begin{tikzpicture}
[scale=0.7,line cap=round,line join=round,>=triangle 45]
%			\draw[->] (-0.5,0) -- (5,0) node[right] {$x$};
%			\draw[->,color=black] (0,-0.5) -- (0,5)  node[right] {$y$};
\draw[->] (3.5,2) -- (5,1.5) node[right] {$S_2$};

\fill (0,0) node[below right] {\normalsize $O$} circle (2pt);
\node[right] at (1.2,2) {$ S_1 $};
\fill (2,0) node [below right]  {$ 4 $} (4,0)node [above right] {$ C $} circle (2pt);
\fill (0,2) node [left] {$ 4 $} (0,4) node [above right] {$ A $}circle (2pt);
\fill (4,4) node [right] {$ B $} circle (2pt);

\clip(-0.5,-0.5) rectangle (4.9,4.9);
\draw[smooth,samples=100]
plot[domain=0:4] (\x,{0.25*(\x)^2});
\draw (4,0)--(4,4)--(0,4);
\draw (0,4)--(0,0)--(4,0);
\fill[pattern=north west lines,smooth]
plot[domain=0:4] (\x,{0.25*(\x)^2})--(4,{0})-- cycle;
\end{tikzpicture}
\end{center}
\loigiai{
\immini{
Chọn hệ trục tọa độ như hình vẽ.\\
Gọi $y=ax^2$ là đường cong parabol đi qua đỉnh $O$ và điểm $B$ theo đề.\\
Vì điểm $B(4;4)$ thuộc parabol nên ta có $4=a\cdot 4^2\Rightarrow a=\dfrac{1}{4}$.\\
Vậy đường cong parabol cần tìm là $y=\dfrac{1}{4}x^2$.\\
Diện tích của hình vuông $ S=4^2=16 $.\\
Diện tích phần bị gạch là $ S_2=\displaystyle\int\limits_{0}^{4} \dfrac{1}{4}x^2\mathrm{\,d}x=\dfrac{16}{3} $. \\
Vậy $ \dfrac{S_1}{S_2}= \dfrac{S-S_2}{S_2}=\dfrac{16-\dfrac{16}{3}}{\dfrac{16}{3}}=2 $.
}
{\begin{tikzpicture}
[scale=0.7,line cap=round,line join=round,>=triangle 45,x=1cm,y=1cm]
\draw[->] (-0.5,0) -- (5,0) node[right] {$x$};
\draw[->,color=black] (0,-0.5) -- (0,5)  node[right] {$y$};
\draw[->] (3.5,2) -- (5,1.5) node[right] {$S_2$};

\fill (0cm,0cm) node[below right] {\normalsize $O$} circle (2pt);
\node[right] at (1.2,2) {$ S_1 $};
\fill (4cm,0cm) node [below right] {$ 4 $} node [above right] {$ C $} circle (2pt);
\fill (0cm,4cm) node [left] {$ 4 $} node [above right] {$ A $}circle (2pt);
\fill (4cm,4cm) node [right] {$ B $} circle (2pt);

\clip(-0.5,-0.5) rectangle (4.9,4.9);
\draw[smooth,samples=100]
plot[domain=0:4] (\x,{0.25*(\x)^2});
\draw (4,0)--(4,4)--(0,4);
%			\draw (0,4)--(0,0)--(4,0);
\fill[pattern=north west lines,smooth]
plot[domain=0:4] (\x,{0.25*(\x)^2})--(4,{0})-- cycle;
\end{tikzpicture}}
}
\end{ex}

\begin{ex}%[2D4V3-2]
\immini[thm]{	Cho hàm số $ y=f(x)=ax^4 + bx^2 + c$ có đồ thị $(C)$, biết rằng $(C)$ đi qua điểm $ A\left(-1;0\right)$, tiếp tuyến $ d$ tại $ A$ của $(C)$ cắt $(C)$ tại hai điểm có hoành độ lần lượt là $ 0$ và $ 2$. Khi diện tích hình phẳng giới hạn bởi $d$, đồ thị $(C)$ và hai đường thẳng $x=0$; $ x=2$ có diện tích bằng $\dfrac{28}{5}$ (phần gạch sọc) . Tính $\displaystyle\int\limits_{-1}^0 f(x)\mathrm{\,d}x$ (làm tròn kết quả đến hàng phần trăm).
}{
\begin{tikzpicture}[yscale=.5,>=stealth, font=\footnotesize, line join=round, line cap=round]
\def\a{1} \def\b{-3} \def\c{2} % Hệ số
\def\xmin{-3} \def\xmax{3}
\def\ymin{-1} \def\ymax{7}
%	\draw[color=gray!50,dashed] (\xmin,\ymin) grid (\xmax,\ymax);
\draw[->] (\xmin,0)--(\xmax,0) node [below]{$x$};
\draw[->] (0,\ymin)--(0,\ymax) node [left]{$y$};
\node at (0,0) [below left]{$O$};
\clip (\xmin+0.1,\ymin+0.1) rectangle (\xmax-0.5,\ymax-0.1);
\draw[smooth,samples=300] plot(\x,{\a*(\x)^4+\b*(\x)^2+\c});
\draw[smooth,samples=300] plot(\x,{2*(\x+1)});
\fill[pattern=north east lines]  plot[domain= 0:2](\x,{2*(\x+1)})--plot[domain= 0:2](\x,{\a*(\x)^4+\b*(\x)^2+\c});
\draw[dashed] (2,6)--(2,0) node[below]{$2$};
\node at (-1,0) [below]{$-1$};
\end{tikzpicture}
}
\shortans{$1{,}20$}
\loigiai{
Vì $d$ là tiếp tuyến với $(C)$ tại điểm $ A\left(-1;0\right)$.\\
Ta có $y'=4a{x^3}+2bx$ $\Rightarrow $ $ d\colon y=\left(-4a-2b\right)\left(x+1\right)$.\\
Phương trình hoành độ giao điểm của $ d$ và $(C)$ là: $\left(-4a-2b\right)\left(x+1\right)=a{x^4}+b{x^2}+c \quad (1)$.\\
Phương trình $(1)$ phải cho $ 2$ nghiệm là $ x=0$, $ x=2$.\\
$\Rightarrow \[\left\{
\begin{aligned}
&-4a-2b=c\\
&-12a-6b=16a+4b+c\\
\end{aligned}
\right.\]\Rightarrow\left\{
\begin{aligned}
&-4a-2b-c=0 & (2)\\
& 28a+10b+c=0 & (3)\\
\end{aligned}
\right.$.\\
Mặt khác, diện tích phần tô màu là $\dfrac{28}{5}=\displaystyle\int\limits_0^2\left[\left(-4a-2b\right)\left(x+1\right)-a{x^4}-b{x^2}-c\right]\mathrm{\,d}x$\\
$\Rightarrow\dfrac{28}{5}=4\left(-4a-2b\right)-\dfrac{32}{5}a-\dfrac{8}{3}b-2c$ $\Rightarrow\dfrac{112}{5}a+\dfrac{32}{3}b+2c=-\dfrac{28}{5} \quad (4)$.\\
Giải hệ 3 phương trình $(2)$, $(3)$ và $(4)$ ta được $ a=1$, $ b=-3$, $ c=2$.\\
Khi đó $ y=f(x)=x^4-3x^2+2$, $ d\colon y=2\left(x+1\right)$.\\
Vậy $\displaystyle\int\limits_{-1}^0\left(x^4-3x^2+2\right)\text{d}x=\dfrac{6}{5}=1{,}2$.}
\end{ex}

\begin{ex}%[EX-Ôn Tập TN 2025, Viết Tường]%[2D4V3-2]
\immini{
Bạn Hải nhận thiết kế logo hình con mắt (phần được tô đậm) cho một cơ sở y tế. Logo là hình phẳng giới hạn bởi hai parabol $y=f(x)$ và $y=g(x)$ như hình vẽ (đơn vị trên mỗi trục toạ độ là decimét). Bạn Hải cần tính diện tích của logo để báo giá cho cơ sở y tế trước khi kí hợp đồng. Diện tích của logo là bao nhiêu decimét vuông? (Làm tròn kết quả đến hàng phần mười).
}{
\begin{tikzpicture}[font=\footnotesize ,line cap=round,line join=round,scale=.8,>=stealth]
\fill[gray,smooth]
plot [domain=-2.45:2.45] (\x,{-0.25*(\x)^2+2})
-- plot[domain=-2.45:2.45] (\x,{0.25*(\x)^2-1})
-- cycle;
\draw[->] (-4,0)--(4,0) node[below left] {$x$};
\draw[->] (0,-2)--(0,3) node[below left] {$y$};
\draw (0,0) node [below left] {$O$};
\foreach \x in {-2,2}
\draw[thin] (\x,1pt)--(\x,-1pt) node[below] {$\x$};
\draw[thin] (-3,1pt)--(-3,-1pt) node[below left] {$-3$};
\draw[thin] (3,1pt)--(3,-1pt) node[below right] {$3$};
\foreach \y in {1,2}
\draw[thin] (1pt,\y)--(-1pt,\y) node[above left] {$\y$};
\draw[thin] (1pt,-1)--(-1pt,-1) node[below left] {$-1$};
\begin{scope}
\clip (-4,-2) rectangle (4, 3);
\draw[domain= -3.5:3.5 ,smooth,variable=\x] plot (\x,{-0.25*(\x)^2+2});
\draw[domain= -3.5:3.5 ,smooth,variable=\x] plot (\x,{0.25*(\x)^2-1});
\end{scope}
\draw[dashed,thin](2,0)--(2,1)--(0,1);
\end{tikzpicture}
}
\shortans{$9{,}8$}
\loigiai
{
Gọi parabol $y=f(x)$ có dạng $f(x)=ax^2+bx+c$.\\
Parabol $y=f(x)$ nhận $Oy$ làm trục đối xứng nên ta có $\dfrac{-b}{2a}=0\Leftrightarrow b=0$.\\
Lại có đồ thị hàm số $y=f(x)$ đi qua điểm $(0;-1)$ và điểm $(2;0)$ nên $a=\dfrac{1}{4}$, $c=-1$.\\
Vậy $f(x)=\dfrac{1}{4}x^2-1$.\\
Tương tự, ta cũng có $y=g(x)=-\dfrac{1}{4}x^2+2$.\\
Phương trình hoành độ giao điểm của $f(x)$ và $g(x)$ là
\[\dfrac{1}{4}x^2-1=-\dfrac{1}{4}x^2+2\Leftrightarrow x=\sqrt{6}\text{ hoặc }x=-\sqrt{6}.\]
Khi đó diện tích của logo là
\begin{eqnarray*}
S&=&\displaystyle\int\limits_{-\sqrt{6}}^{\sqrt{6}}\left[\left(-\dfrac{1}{4}x^2+2\right)-\left(\dfrac{1}{4}x^2-1\right)\right]\mathrm{\,d}x\\
&=&\displaystyle\int\limits_{-\sqrt{6}}^{\sqrt{6}}\left(3-\dfrac{1}{2}x^2\right)\mathrm{\,d}x=\left(3x-\dfrac{x^3}{6}\right)\Bigg|_{-\sqrt{6}}^{\sqrt{6}}=4\sqrt{6}\approx 9{,}8\text{ dm}^2.
\end{eqnarray*}
}
\end{ex}

\begin{ex}%[Cánh diều, mức độ 3]%[BG12-4in1, Đỗ Vũ Minh Thắng]%[2D4V3-2]
Người ta dự định lắp kính cho cửa của một vòm có dạng parabol.
\begin{center}
\begin{tikzpicture}[scale=1, font=\footnotesize, line join=round, line cap=round, >=stealth]
\def\xt{-6.5}
\def\xp{6.5}
\def\yd{-0.5}
\def\yt{7.5}
\begin{scope}[scale=0.5]
\draw[purple, line width=1pt] plot[smooth, samples=100, domain=-6:6] (\x,{-0.2*(\x+6)*(\x-6)}) -- (-6,0);
\draw[<->, blue] (-6,-1)--(6,-1) node[midway, below]{$70$ m};
\clip (-6,0) -- plot[smooth, samples=100, domain=-6:6] (\x,{-0.2*(\x+6)*(\x-6)}) -- (-6,0);
\foreach \i in {0,1,2,...,11}
{\fill[cyan!40!white, opacity=0.1] (\i-6,0) rectangle (\i-5,2)
(\i-6,2) rectangle (\i-5,4)
(\i-6,4) rectangle (\i-5,7.5);
\draw[lightgray, line width=1] (\i-6,0) rectangle (\i-5,2)
(\i-6,2) rectangle (\i-5,4)
(\i-6,4) rectangle (\i-5,7.5);}
\draw[<->, blue] (0,0)--(0,7.2) node[sloped, midway, shift={(150:0.3)}]{$21$ m};
\end{scope}
\end{tikzpicture}
\end{center}
Hãy tính diện tích mặt kính (theo đơn vị m$^2$) cần lắp vào, biết rằng vòm cửa cao $21~\text{m}$ và rộng $70~\text{m}$.
\shortans{$ 980 $}
\loigiai{
\begin{center}
\begin{tikzpicture}[scale=1, font=\footnotesize, line join=round, line cap=round, >=stealth]
\def\xt{-7}
\def\xp{7}
\def\yd{-1}
\def\yt{9}
\begin{scope}[scale=0.5]
\draw[purple, line width=1pt] plot[smooth, samples=100, domain=-6:6] (\x,{-0.2*(\x+6)*(\x-6)}) -- (-6,0);
\clip (-6,0) -- plot[smooth, samples=100, domain=-6:6] (\x,{-0.2*(\x+6)*(\x-6)}) -- (-6,0);
\foreach \i in {0,1,2,...,11}
{\fill[cyan!40!white, opacity=0.1] (\i-6,0) rectangle (\i-5,2)
(\i-6,2) rectangle (\i-5,4)
(\i-6,4) rectangle (\i-5,7.5);
\draw[lightgray, line width=1] (\i-6,0) rectangle (\i-5,2)
(\i-6,2) rectangle (\i-5,4)
(\i-6,4) rectangle (\i-5,7.5);}
\end{scope}
\begin{scope}[scale=0.5]
\draw[->, thick] (\xt,0)--(0,0) node[below left]{$O$}--(\xp,0) node[below]{$x$};
\draw[->, thick] (0,\yd)--(0,\yt) node[left]{$y$};
\fill (-6,0) node[below]{$B$} circle(2pt)
(6,0) node[below]{$C$} circle(2pt)
(0,7.2) node[above right]{$A$} circle(2pt);
\end{scope}
\end{tikzpicture}
\end{center}

\noindent Đặt vào hệ trục tọa độ, đỉnh cửa trùng với điểm $A \in Oy$, hai điểm $B$, $C$ thuộc trục $Ox$.\\
Do cửa có chiều cao là $35\text{m}$ và đáy có chiều dài là $70\text{m}$ nên $A(0;21)$, $B(-35;0)$ và $C(35;0)$.\\
Do đồ thị Parabol cắt trục hoành tại $B(-35;0)$ và $C(35;0)$ nên parabol có phương trình là\\ $y=a(x-35)(x+35)$.\\
Mà parabol qua điểm $A(0;21)$ nên $y=\dfrac{-3}{175}(x-35)(x+35)=\dfrac{-3}{175}(x^2-1225)$.\\ Khi đó, diện tích kính của cánh cửa là $S=\displaystyle\int\limits_{-35}^{35} \dfrac{-3}{175}(x^2-1225) \mathrm{\,d}x=980~\text{m}^2$.
}
\end{ex}

\begin{ex}%[CTST, mức độ 3]%[BG12-4in1, Đỗ Vũ Minh Thắng]%[2D4V3-2]
\immini{
Mặt cắt của một cửa hầm có dạng là hình phẳng giới hạn bởi một parabol và đường thẳng nằm ngang như hình bên. Tính diện tích của cửa hầm (theo đơn vị m$^2$).
}
{
\begin{tikzpicture}
\fill[pattern=north east lines,pattern color=black](-2.5,0)--(-2.5,4.5)--(2.5,4.5)--(2.5,0)--cycle;
\draw plot[domain=-2:2](\x,{-\x*(\x)+4})--(2,0);
\draw[<->](-2,-0.5)--(2,-0.5);
\draw[<->](3,0)--(3,4);
\draw[dashed] (-2,0)--(-2,-0.5) (2,-0.5)--(2,0)
(0,4)--(3,4) (2,0)--(3,0);
\path (0,-0.5)node[above]{$6$ m} (3,2)node[right]{$6$ m};
\draw (-2,0)--(2,0);
\fill[white](-2,0)--plot[domain=-2:2](\x,{-\x*(\x)+4})--(2,0);
\end{tikzpicture}
}
\shortans{$ 24 $}
\loigiai{
\immini{Chọn hệ trục tọa độ như hình vẽ.\\
Gọi đồ thị hàm số của parabol có dạng $y=ax^2+bx+c \,\,(a\neq 0)$.\\
Vì đồ thị hàm số có trục đối xứng là trục $Oy$ nên ta có: $-\dfrac{b}{2a}=0\Rightarrow b=0$.\\
Vì đồ thị hàm số đi qua điểm $(0;6)$ nên ta có: $c=6$.\\
Vì đồ thị hàm số đi qua điểm $(3;0)$ nên ta có: $a\cdot3^2+6=0\Rightarrow a=-\dfrac{2}{3}$.\\
Vậy $y=-\dfrac{2}{3}x^2+6$.}{
\begin{tikzpicture}[scale=0.7]
%\fill[yellow!30](-2,0)--plot[domain=-2:2](\x,{-\x*(\x)+4})--(2,0);
\draw[->](-3,0)--(3,0)node[right]{$x$};
\draw[->](0,-0.5)--(0,4.5)node[above]{$y$};
%\fill[black!10](-2.5,0)--(-2.5,4.5)--(2.5,4.5)--(2.5,0)--cycle;
\draw plot[domain=-2:2](\x,{-\x*(\x)+4})--(2,0);
\path (-2,0)node[below]{$-3$}
(2,0)node[below]{$3$}
(0,0)node[below left]{$O$}
(0,4)node[below right]{$6$};
\end{tikzpicture}
}
Diện tích của cửa hầm là diện tích hình phẳng giới hạn bởi đồ thị hàm số $y=-\dfrac{2}{3}x^2+6$, trục hoành và hai đường thẳng $x=-3, x=3$. Do đó
\begin{align*}
S=\displaystyle\int\limits_{-3}^{3} \left|-\dfrac{2}{3}x^2+6\right|\mathrm{\,d}x
=\left|\displaystyle\int\limits_{-3}^{3} (-\dfrac{2}{3}x^2+6)\mathrm{\,d}x\right|
=\left|\left(-\dfrac{2x^3}{9}+6x\right)\bigg|_{-3}^{3}\right|
=\left|24\right|
=24~\text{m}^2.
\end{align*}
}
\end{ex}

\begin{ex}%[Cánh diều, mức độ 3]%[BG12-4in1, Đỗ Vũ Minh Thắng]%[2D4V3-2]
Hình dưới đây minh họa mặt đứng của một con kênh đặt trong hệ trục tọa độ $Oxy$. Đáy của con kênh là một đường cong cho bởi phương trình $y=f(x)=\dfrac{3}{100} \left(-\dfrac{1}{3}x^3+5x^2 \right)$.
\begin{center}
\begin{tikzpicture}[scale=1, font=\footnotesize, line join=round, line cap=round, >=stealth]
\def\xt{-9}
\def\xp{15}
\def\yd{-2}
\def\yt{8}
\begin{scope}[scale=0.25]
\fill[pattern=north east lines,pattern color=black] (\xt+1,\yd+0.5)--(\xt+1,5) -- (-5,5) node[above, text=black]{$A$}-- plot[domain=-5:10, smooth, samples=150] (\x, {-0.01*(\x)^3+0.15*(\x)^2}) node[above, text=black]{$B$} -- (\xp-1,5) -- (\xp-1,\yd+0.5) -- cycle;
\draw (\xt+1,\yd+0.5)--(\xt+1,5) -- (-5,5)-- plot[domain=-5:10, smooth, samples=150] (\x, {-0.01*(\x)^3+0.15*(\x)^2}) -- (\xp-1,5) -- (\xp-1,\yd+0.5) -- cycle;
\draw[white] (-5,5) -- plot[domain=-5:10] (\x, {-0.01*(\x)^3+0.15*(\x)^2}) --cycle;
\draw (-5,5) -- plot[domain=-5:10] (\x, {-0.01*(\x)^3+0.15*(\x)^2}) --cycle;

\draw[->](\xt,0) -- (0,0)node[below left, fill=white, thick, rounded corners=3pt, inner sep=2pt]{$O$} -- (\xp,0)node[below]{$x$};
\draw[->](0,\yd)--(0,\yt)node[left]{$y$};
\fill (-5,0) circle(3pt) (10,0) circle(3pt) (0,5) node[above left]{$5$};
\draw[dashed] (-5,0) node[below, fill=white, thick, rounded corners=3pt, inner sep=2pt]{$-5$} -- (-5,5)
(10,0) node[below, fill=white, thick, rounded corners=3pt, inner sep=2pt]{$10$} -- (10,5);
\path [shift={(90:0.5)}, decoration={text along path, text={$y=f(x)$}}, decorate] plot[domain=3:10] (\x, {-0.01*(\x)^3+0.15*(\x)^2});
\end{scope}
\end{tikzpicture}
\end{center}
Biết đơn vị trên mỗi trục tọa độ là mét, diện tích mặt cắt đứng của mức nước trong con kênh là một phân số tối giản có dạng $ \dfrac{a}{b} $ m$^2$, trong đó $ a, b \in \mathbb{N^*} $. Tìm $ a $.
\shortans{$ 675 $}
\loigiai{
\noindent Từ hình vẽ, mặt cắt đứng của mức nước trong con kênh là hình phẳng tô màu giới hạn bởi đồ thị hàm số $f(x)$, $y=5$ và hai đường thẳng $x=-5$, $x=10$.\\
\noindent Khi đó diện tích mặt cắt đứng của mức nước trong con kênh là
\[S=\displaystyle\int\limits_{-5}^{10} \left[ 5-\dfrac{3}{100} \left(-\dfrac{1}{3}x^3+5x^2 \right) \right] \mathrm{\,d}x=\dfrac{675}{16}~\text{m}^2.\]
}
\end{ex}

\begin{ex}%[BG-12-4in1, Phạm Đức]%[2D4V3-2]
\immini
{
Người ta xây dựng một đường hầm hình parabol đi qua núi có chiều cao $OI=9$ m, chiều rộng $AB=10$ m. Tính diện tích cửa đường hầm.
}
{
\begin{tikzpicture}[scale=0.4,font=\footnotesize,>=stealth]
\draw (0,0) parabola bend (5,9) (10,0)--cycle (5,9)--(5,0) node[midway,sloped, above]{$9$ m};
\fill (0,0) circle(2pt) node[below left]{$A$} (10,0) circle(2pt) node[below right]{$B$} (5,0) circle(2pt) node[below]{$O$} (5,9) circle(2pt) node[above]{$I$};
\draw[<->,yshift=-0.9cm] (0,0)--(10,0) node[midway,below]{$10$ m};
\end{tikzpicture}
}
\shortans{$60$}
\loigiai
{
Gắn hệ trục tọa độ $Oxy$ vào hình vẽ sao cho $Ox$ trùng với tia $OB$, $Oy$ trùng với tia $OI$. Khi đó $A(-5;0)$, $B(5;0)$, $I(0;9)$.\\
Giả sử parabol có phương trình $f(x)=ax^2+bx+c$ ($a\neq 0$).\\
Vì parabol đi qua ba điểm $A(-5;0)$, $B(5;0)$, $I(0;9)$ nên
$\left\{\begin{aligned}&25a-5b+c=0 \\&25a+5b+c=0 \\&c=9\end{aligned}\right. \Leftrightarrow \left\{\begin{aligned}&a=-\dfrac{9}{25} \\&b=0 \\&c=9.\end{aligned}\right.$\\
Vậy $f(x)=-\dfrac{9}{25}x^2+9$.\\
Diện tích cửa đường hầm là
\[S=\displaystyle\int\limits_{-5}^5\left|-\dfrac{9}{25}x^2+9\right|\mathrm{\,d}x = \displaystyle\int\limits_{-5}^5\left(-\dfrac{9}{25}x^2+9\right)\mathrm{\,d}x = \left(-\dfrac{3}{25}x^3+9x\right)\bigg|_{-5}^5 = 30+30=60.\]
Vậy diện tích cửa đường hầm bằng $60\text{\,m}^2$.
}

\end{ex}

\begin{ex}%[2D4V3-2]
\immini{
Sàn của một viện bảo tàng mỹ thuật được lát bằng những viên gạch hình vuông cạnh $40$ cm  như hình bên. Biết rằng người thiết kế đã sử dụng các đường cong có phương trình $4x^2=y^4$  và  $4y^2 =x^4$ để tạo hoa văn cho viên gạch. Tính diện tích (đơn vị cm$^2$) phần được tô đậm (làm tròn kết quả đến hàng đơn vị).
}
{
\begin{tikzpicture}
\draw (2,2)--(2,-2)--(-2,-2)--(-2,2)--cycle;
\draw[fill=cyan, smooth, samples=200] plot[domain=-2:2,variable=\y]({0.5*(\y)^2},{\y})--plot[domain=2:0,variable=\x]({\x},{0.5*(\x)*(\x)})--plot[domain=0:2,variable=\x]({\x},{(-0.5)*(\x)*(\x)});
\draw[fill=cyan, smooth, samples=200] plot[domain=-2:2,variable=\y]({-0.5*(\y)^2},{\y})--plot[domain=-2:0,variable=\x]({\x},{0.5*(\x)*(\x)})--plot[domain=0:-2,variable=\x]({\x},{(-0.5)*(\x)*(\x)});
\end{tikzpicture}
}
\shortans{$533$}
\loigiai{
\immini{
Do tính đối xứng nên diện tích cần tìm bằng $4$ lần diện tích phần tô đậm của hình vẽ bên. Ta chỉ xét đồ thị trong góc phần tư thứ nhất do đó
\begin{align*}
&4x^2=y^4 \Leftrightarrow y = \sqrt{2x} \\
&4y^2 =x^4 \Leftrightarrow  y = \dfrac{1}{2} x^2.
\end{align*}
}
{
\begin{tikzpicture}[scale=1,>=stealth]
\draw[->] (-1,0)--(3,0) node[below] {$x$};
\draw[->] (0,-1)--(0,3) node[left] {$y$};
%\draw[fill] (1,0) circle (1pt) node[below] {$1$};
\draw[fill] (0,0) circle node[below left=-2pt] {$O$};
\draw[dashed] (2,0)--(2,2)--(0,2);
\draw[fill] (2,0) circle (1pt) node[below] {$2$};
\draw[fill] (0,2) circle (1pt) node[left] {$2$};
\draw[fill=cyan, smooth, samples=200] plot[domain=0:2,variable=\y]({0.5*(\y)^2},{\y})--plot[domain=2:0,variable=\x]({\x},{0.5*(\x)*(\x)});
\end{tikzpicture}
}
\noindent Một đơn vị trong hệ tọa độ $Oxy$ bằng $10$ cm, do đó diện tích của phần tô đậm ban đầu là
\[
S =4\cdot 10^2 \int \limits_0^2 \left ( \sqrt{2x}-\dfrac{1}{2}x ^2\right) \mathrm{d}x = 400 \left ( \dfrac{2\sqrt{2}}{3} \sqrt{x^3}-\dfrac{1}{6}x^3 \right ) \Bigg|_0^2= \dfrac{1600}{3}\approx 533 \ \mathrm{(cm^2)}.
\]
}
\end{ex}

\begin{ex}%[Dự án 2025 - đề cấu trúc mới, Hung Doan]%[2D4V3-2]
Một viên gạch hoa hình vuông cạnh $40$ cm. Người thiết kế đã sử dụng bốn đường parabol có chung đỉnh tại tâm viên gạch để tạo ra bốn cánh hoa (được tô màu sẫm như hình vẽ bên).
\begin{center}
\begin{tikzpicture}[scale=1, font=\footnotesize, line join=round, line cap=round, >=stealth]
\draw (2,2)--(2,-2)--(-2,-2)--(-2,2)--(2,2);
\clip(-2,-2) rectangle(2,2);
\draw[smooth,samples=300] plot(\x,{0.5*(\x)^2});
\draw[smooth,samples=300] plot(\x,{-0.5*(\x)^2});
\draw[samples=200,domain=0:2,smooth] plot (\x,{sqrt(2*(\x))});
\draw[samples=200,domain=0:2,smooth] plot (\x,{-sqrt(2*(\x))});
\draw[samples=200,domain=-2:0,smooth] plot (\x,{sqrt(-2*(\x))});
\draw[samples=200,domain=-2:0,smooth] plot (\x,{-sqrt(-2*(\x))});
\fill[gray,smooth] (-2,0)--plot[domain=-2:0](\x,{-sqrt(-2*(\x))})--(0,0)--cycle;
\fill[gray,smooth] (-2,0)--plot[domain=-2:0](\x,{sqrt(-2*(\x))})--(0,0)--cycle;
\fill[gray,smooth] (0,0)--plot[domain=0:2](\x,{-sqrt(2*(\x))})--(2,0)--cycle;
\fill[gray,smooth] (0,0)--plot[domain=0:2](\x,{sqrt(2*(\x))})--(2,0)--cycle;
\fill[white,smooth] (-2,0)--plot[domain=-2:2](\x,{-0.5*(\x)^2})--(2,0)--cycle;
\fill[white,smooth] (-2,0.01)--plot[domain=-2:2](\x,{0.5*(\x)^2})--(2,0.01)--cycle;
\end{tikzpicture}
\end{center}
Diện tích mỗi cánh hoa của viên gạch bằng bằng $\dfrac{a}{b}$ (cm$^2$), với $\dfrac{a}{b}$ là phân số tối giản thì $a$ bằng bao nhiêu?
\shortans{$400$}
\loigiai{
\begin{center}
\begin{tikzpicture}[scale=1, font=\footnotesize, line join=round, line cap=round, >=stealth]
\begin{scope}
\clip(-2,-2) rectangle(2,2);
\draw[smooth,samples=300] plot(\x,{0.5*(\x)^2});
\draw[smooth,samples=300] plot(\x,{-0.5*(\x)^2});
\draw[samples=200,domain=0:2,smooth] plot (\x,{sqrt(2*(\x))});
\draw[samples=200,domain=0:2,smooth] plot (\x,{-sqrt(2*(\x))});
\draw[samples=200,domain=-2:0,smooth] plot (\x,{sqrt(-2*(\x))});
\draw[samples=200,domain=-2:0,smooth] plot (\x,{-sqrt(-2*(\x))});
\fill[gray,smooth] (-2,0)--plot[domain=-2:0](\x,{-sqrt(-2*(\x))})--(0,0)--cycle;
\fill[gray,smooth] (-2,0)--plot[domain=-2:0](\x,{sqrt(-2*(\x))})--(0,0)--cycle;
\fill[gray,smooth] (0,0)--plot[domain=0:2](\x,{-sqrt(2*(\x))})--(2,0)--cycle;
\fill[gray,smooth] (0,0)--plot[domain=0:2](\x,{sqrt(2*(\x))})--(2,0)--cycle;
\fill[white,smooth] (-2,0)--plot[domain=-2:2](\x,{-0.5*(\x)^2})--(2,0)--cycle;
\fill[white,smooth] (-2,0.01)--plot[domain=-2:2](\x,{0.5*(\x)^2})--(2,0.01)--cycle;
\end{scope}
\draw[->](-2.75, 0)--(3, 0) node[below]{$x$};
\draw[->](0,-2.5)--(0,2.75) node[left]{$y$};
\node at (0, 0) [below left]{$O$};
\draw[shift ={ (-2,0)}]node[below left]{$-2$} (0pt,2pt) --(0pt,-2pt);
\draw[shift ={ (2,0)}]node[below right]{$2$} (0pt,2pt) --(0pt,-2pt);
\draw[shift ={ (0,2)}]node[above right]{$2$} (2pt,0pt) --(-2pt,0pt);
\draw[shift ={ (0,-2)}]node[below left]{$-2$} (2pt,0pt) --(-2pt,0pt);
\draw[smooth] (2,2)--(2,-2)--(-2,-2)--(-2,2)--(2,2);
\end{tikzpicture}
\end{center}
Chọn hệ tọa độ như hình vẽ ($1$ đơn vị trên trục bằng $10$ cm=$1$ dm), các cánh hoa tạo bởi các đường parabol có phương trình $y=\dfrac{x^2}{2}$, $y=-\dfrac{x^2}{2}$, $x=-\dfrac{y^2}{2}$, $x=\dfrac{y^2}{2}$.\\
Diện tích một cánh hoa (nằm trong góc phần tư thứ nhất) bằng diện tích hình phẳng giới hạn bởi hai đồ thị hàm số $y=\dfrac{x^2}{2}$, $y=\sqrt{2x}$ và hai đường thẳng $x=0$; $x=2$.\\
Do đó diện tích một cánh hoa bằng
\begin{align*}
\displaystyle\int \limits_0^2 \left(\sqrt{2x}-\dfrac{x^2}{2}\right)\mathrm{d}x &=\left(\dfrac{2\sqrt{2}}{3}\sqrt{(2x)^3}-\dfrac{x^3}{6}\right)\bigg|\bigg|_0^2\\
&=\dfrac{4}{3}(\text{dm}^2)=\dfrac{400}{3}\,(\text{cm}^2).
\end{align*}
Suy ra $a=400$.
}
\end{ex}

\begin{ex}%[Vovanle]%[2D4V3-2]
\immini{Một khuôn viên dạng nửa hình tròn có đường kính bằng $4\sqrt{5}$ (m). Trên đó người thiết kế hai phần để trồng hoa có dạng của một cánh hoa hình parabol có đỉnh trùng với tâm nửa hình tròn và hai đầu mút của cánh hoa nằm trên nửa đường
}{
\begin{tikzpicture}[line join=round,line cap=round, font=\footnotesize,scale=0.5,>=stealth]
\path
({-2*sqrt (5)},0) coordinate (A)
({2*sqrt (5)},0) coordinate (B)
(2,4) coordinate (M)
(-2,4) coordinate (N)
(-2,0) coordinate (C)
(2,0) coordinate (D)
($(M)!0.5!(D)$) coordinate (G)node[right]{$4$ m}
;
\fill[pattern=north east lines]plot[domain=-2:2](\x,{sqrt (20-(\x)^2)})--plot[domain=2:-2](\x,{(\x)^2});
\draw (A) arc(180:0:{2*sqrt (5)});
\draw plot[domain=-2:2](\x,{(\x)^2});
\draw[dashed] (N)--(C)(M)--(D)(M)--(N);
\path ($(M)!0.5!(N)$) coordinate (H)node[below]{$4$ m};
\draw (A)--(B);
\end{tikzpicture}
}
\noindent tròn (phần gạch sọc), cách nhau một khoảng bằng $4\,\mathrm{m}$, phần còn lại của khuôn viên (phần không gạch sọc) dành để trang trí cỏ nhân tạo. Biết các kích thước cho như hình vẽ và kinh phí cỏ nhân tạo là $100\,000$ đồng/m$^2$. Hỏi cần bao nhiêu tiền để trang trí cỏ trên phần đất đó? (Số tiền được làm tròn đến hàng nghìn).
\shortans{$1948$}
\loigiai{
\immini{Đặt hệ trục tọa độ như hình vẽ. Khi đó phương trình nửa đường tròn là
\[y=\sqrt{R^2-x^2}=\sqrt{\left(2\sqrt{5}\right)^2-x^2}=\sqrt{20-x^2}.\]
Phương trình parabol $(P)$ có đỉnh là gốc $O$ sẽ có dạng $y=ax^2$. Mặt khác $(P)$ qua điểm $M(2;4)$.
}{
\begin{tikzpicture}[line join=round,line cap=round, font=\footnotesize,scale=0.6,>=stealth]
\path
({-2*sqrt (5)},0) coordinate (A)
({-2*sqrt (5)},0) coordinate (B)
(2,4) coordinate (M)node[above right]{$M(2,4)$}
(-2,4) coordinate (N)
(-2,0) coordinate (C)
(2,0) coordinate (D)
;
\fill[pattern=north east lines]plot[domain=-2:2](\x,{sqrt (20-(\x)^2)})--plot[domain=2:-2](\x,{(\x)^2});
\draw[-stealth](-5,0)--(5,0)node[below]{$x$};
\draw[-stealth](0,-0.7)--(0,5)node[right]{$y$};
\fill (0,0) circle(1pt)node[below left]{$O$}(-2,0)circle(1pt) node[below]{$-2$}(2,0)circle(1pt)node[below]{$2$}(0,4)circle(1pt)node[below left]{$4$};
\draw (A) arc(180:0:{2*sqrt (5)});
\draw plot[domain=-2:2](\x,{(\x)^2});
\draw[dashed] (N)--(C)(M)--(D)(M)--(N);
\end{tikzpicture}
}
\noindent Do đó $4=a\cdot(-2)^2\Rightarrow a=1$.\\
Phần diện tích của hình phẳng giới hạn bởi $(P)$ và nửa đường tròn.(phần gạch sọc).\\
Ta có công thức $S_1=\displaystyle\int\limits_{-2}^2{\left(\sqrt{20-x^2}-x^2 \right)\mathrm{\,d}x}\approx 11{,}94\,\mathrm{m}^2$.\\
Vậy phần diện tích trồng cỏ là $S_{\text{cỏ}}=\dfrac{1}{2}{{S}_{\text{htron}}}-S_1=\dfrac{1}{2}\cdot 20\cdot \pi-11{,}94\approx 19{,}476\,\mathrm{m}^2$.\\
Số tiền cần có là $S_{\text{cỏ}}\times 100000\approx 1947592\text{ (đồng)}\approx 1948$ (nghìn đồng).
}
\end{ex}

\begin{ex}%[Dự án 2025 - đề cấu trúc mới, Nguyễn Kiều Nhã Tú]%[2D4V3-2]
\immini[thm]{Ông Bình muốn làm một cổng sắt có hình dạng và kích thước giống hình vẽ, biết đường cong phía trên là một parabol. Nếu giá một mét vuông cổng sắt là $1\,200\,000$ đồng thì ông Bình phải trả bao nhiêu tiền để làm cổng sắt như vậy (số tiền phải trả tính theo đơn vị là triệu đồng)?
}{
\begin{tikzpicture}[line join = round, line cap = round,>=stealth,font=\footnotesize,scale=1]
\def \rong{3}
\def \cao{1.5}
\def \so{9}
\pgfmathsetmacro{\bk}{3/(2*\so+3)}
\draw[line width=1pt,double,shift={(90:.5)}] (-\rong,0) to[bend left]coordinate[pos=0.5](A)
(\rong,0);
\draw[,double,line width=1pt] (-\rong,0) to[bend left] (\rong,0);
\draw[line width=3pt] (\rong,.5) --++(0,-\cao-0.5)--++(-2*\rong,0)--++(0,\cao+0.5)
(0,-\cao)--(A)
;
\foreach \i in {1,...,\so}% Tạo nội dung lặp
{
\pgfmathsetmacro{\tile}{\i/(2*\so+2)}
\path (-\rong,0) to[bend left] coordinate[pos=\tile](A\i) (\rong,0);
\path (\rong,0) to[bend right] coordinate[pos=\tile](B\i) (-\rong,0);
\coordinate (C\i) at ($(-\rong,-\cao)!\tile!(\rong,-\cao)$);
\coordinate (D\i) at ($(\rong,-\cao)!\tile!(-\rong,-\cao)$);
\coordinate (E\i) at ($(C\i)+(0,\cao/2)$);
\coordinate (F\i) at ($(D\i)+(0,\cao/2)$);
\coordinate (G\i) at ($(A\i)+(0,-0.5)$);
\coordinate (H\i) at ($(B\i)+(0,-0.5)$);
\draw (E\i) arc(0:360:\bk) (F\i) arc(180:-180:\bk) (G\i) arc(0:360:\bk) (H\i) arc(-180:180:\bk);
\draw (A\i) arc(270:630:{.1cm} and {.2cm})
(B\i) arc(270:630:{.1cm} and {.2cm})
;
\fill
(A\i) circle (1.5pt)
(B\i) circle (1.5pt)
(C\i) circle (1.5pt)
(D\i) circle (1.5pt)
;
\draw[->,-latex] (C\i)--(A\i)--++(0,0.5);
\draw[->,-latex] (D\i)--(B\i)--++(0,0.5);
}
\draw
(E9) arc (180:540:{\bk})
(F9) arc (0:360:{\bk})
([shift={(0:2*\bk)}]E9) --(-\rong,-\cao/2) ([shift={(180:2*\bk)}]F9)--(\rong,-\cao/2)
;
\draw[white] (A)--(0,-\cao) (\rong,.5) --++(0,-\cao-0.5)--++(-2*\rong,0)--++(0,\cao+0.5)
(0,-\cao)--(A);
\draw[dashed] (\rong,-\cao)--++(0,-0.5)(-\rong,-\cao)--++(0,-0.5) (\rong,-\cao)--++(1.5,0) (\rong,0.5)--++(0.5,0) (A)--++(\rong+1.5,0);
\draw[<->] (\rong,-\cao-0.3)--++(-2*\rong,0) node[fill=white,midway]{$5$ m};
\draw[<->] (\rong+0.2,-\cao)--++(0,\cao+0.5) node[midway,right]{$1{,5}$ m};
\draw[<->] ($(A)+(\rong+1.3,0)$)--(\rong+1.3,-\cao) node[midway,right]{$2$ m};
\end{tikzpicture}
}
\shortans{$11$}
\loigiai{
Chọn hệ trục tọa độ $Oxy$ như hình vẽ.
\begin{center}
\begin{tikzpicture}
\def\a{-0.08}
\def\b{0}
\def\c{0.5}
\def\f(#1){\a*((#1)^2)+\b*(#1)+\c}
\pgfmathsetmacro\ct{-\b/(2*\a)}
\def\xmin{-4.5}
\def\xmax{5}
\def\ymin{-1.5}
\def\ymax{1.5}
\draw[->] (\xmin-1,0)--(\xmax+1,0) node[below] { $x$};
\draw[->] (0,\ymin-1)--(0,\ymax+1) node[left] { $y$};
\draw (0,0) node [below left] {$O$};
\foreach \x in {-2.5}		\draw (\x,0.1)--(\x,-0.1) node [above left] {$\x$};
\foreach \x in {2.5}		\draw (\x,0.1)--(\x,-0.1) node [above right] {$\x$};
\foreach \y in {-1.5,0.5} \draw (0.1,\y)--(-0.1,\y) node [above left] {$\y$};
%	\clip (\xmin,\ymin) rectangle (\xmax,\ymax);
\draw[smooth,samples=200] plot[domain=\xmin:\xmax] (\x,{\f(\x)});
\draw (-2.5,0)--(-2.5,-1.5)--(2.5,-1.5)--(2.5,0);
\end{tikzpicture}
\end{center}
Diện tích cổng gồm diện tích hình chữ nhật và diện tích phần giới hạn bởi parabol $(P)$ và trục hoành.\\
Từ tọa độ $3$ điểm thuộc parabol $(P)$ là $(0;0{,}5)$, $(-2{,}5;0)$ và $(2{,}5;0)$, ta tìm được phương trình của parabol $(P)$ là $y=f(x)=-\dfrac{2}{25}x^2+\dfrac{1}{2}$.\\
Diện tích của cổng $S=S_{(P)}+S_{\text{hcn}}=\displaystyle\int\limits_{-2{,}5}^{2{,}5}\left(-\dfrac{2}{25}x^2+\dfrac{1}{2}\right)\mathrm{\,d}x+5\cdot 1{,}5=\dfrac{5}{3}+\dfrac{15}{2}=\dfrac{55}{6}$ m$^2$.\\
Giá một mét vuông cổng sắt là $1\,200\,000$ đồng.\\
Vậy số tiền ông Bình phải trả để làm cổng sắt là $\dfrac{55}{6}\cdot 1\,200\,000=11\,000\,000$ đồng.
}
\end{ex}

\begin{ex}%[2D4V3-2]
\immini[thm]{Một khu vườn hình bán nguyệt có bán kính $R=4$ m, ở giữa khu vườn người ta muốn tạo một cái bể cá có đường biên là một parabol có phương trình $ y=\dfrac{\sqrt{3}}{4}{x^2}+\sqrt{3}$ (như hình vẽ), phần còn lại sẽ trồng hoa. Biết chi phí xây bể cá là $400\,000$ đồng/m$^2$, chi phí xây khu trồng hoa là $200\,000$ đồng/m$^2$.
}
{
\begin{tikzpicture}[scale=0.7,>=stealth, font=\footnotesize, line join=round, line cap=round]
\pgfmathsetmacro{\bb}{sqrt(3)/4}
\pgfmathsetmacro{\cc}{sqrt(3)}
\def\a{1} \def\b{-4} \def\c{3} % Hệ số
\def\xmin{-4.5} \def\xmax{4.5}
\def\ymin{-0.5} \def\ymax{4.5}
\clip (\xmin+0.1,\ymin+0.1) rectangle (\xmax-0.5,\ymax-0.1);
\draw[smooth,samples=300,domain=-2:2] plot(\x,{\bb*(\x)^2+\cc});
\draw[smooth,samples=300,domain=-4:4] plot(\x,{sqrt(16-1*(\x)^2)});
\draw(-4,0)--(4,0);
\fill[pattern=north east lines,opacity=0.8] plot[domain=-2:2](\x,{sqrt(16-(\x)^2)})--plot[domain=-2:2](\x,{\bb*(\x)^2+\cc});
\end{tikzpicture}}\noindent
Khi đó chi phí xây dựng toàn bộ khu vườn hết bao nhiêu triệu đồng (làm tròn đến hàng phần trăm)?

\shortans{6,24}
\loigiai{
Phương trình hoành độ giao điểm
\allowdisplaybreaks
\begin{eqnarray*}
&&\sqrt{16-x^2}=\dfrac{\sqrt{3}}{4}{x^2}+\sqrt{3}\\
&\Leftrightarrow& 16-x^2=\left(\dfrac{\sqrt{3}}{4}{x^2}+\sqrt{3}\right)^2\\
&\Leftrightarrow&\dfrac{3}{16}{x^4}+\dfrac{5}{2}{x^2}-13=0\Rightarrow\hoac{
&x=2\\
&x=-2.}
\end{eqnarray*}
Diện tích bể cá $ S=\displaystyle\int\limits_{-2}^2\left(\sqrt{16-x^2}-\dfrac{\sqrt{3}}{4}{x^2}-\sqrt{3}\right)\mathrm{d}x=\dfrac{8\pi-4\sqrt{3}}{3}$ (m$^2$).\\
Diện tích trồng hoa $8\pi-\dfrac{8\pi-4\sqrt{3}}{3}=\dfrac{16\pi+4\sqrt{3}}{3}$ (m$^2$).\\
Chi phí xây dựng $\dfrac{8\pi-4\sqrt{3}}{3}\cdot 400\, 000+\dfrac{16\pi+4\sqrt{3}}{3}\cdot200\, 000\approx 6\,240\, 184$ đồng.}
\end{ex}

\begin{ex}%[Dự án C đề thi thử THPT QG 2025]%[Đỗ Chí Tâm]%[2D4V3-2]
\immini[thm]{Một công ty có ý định thiết kế một logo hình vuông có độ dài nửa đường chéo bằng $4$. Biểu tượng $4$ chiếc lá (được tô màu) được tạo thành bởi các đường cong đối xứng với nhau qua tâm của hình vuông và qua các đường chéo. Một trong số các đường cong ở nửa bên phải của logo là một phần của đồ thị hàm số bậc ba dạng $y = ax^3 + bx^2 - x$ với hệ số $a<0$. Để kỷ niệm ngày thành lập $2/3$, công ty thiết kế để tỉ số diện tích được tô màu so với phần không được tô màu bằng $\dfrac{2}{3}$. Tính $2a + 2b$.
}
{\begin{tikzpicture}[scale=0.7, font=\footnotesize, line join=round, line cap=round,>=stealth]
\foreach \i in {0,1,2,3}{
\fill[cyan,smooth,rotate=90*\i] plot [domain=0:3](\x,{-(\x)*(\x-3)^2/8})--plot [domain=3:0](\x,{(\x)*(\x-3)^2/8})--cycle;
\draw (0,0)--(90*\i:3)--(90*\i+90:3);
}
\end{tikzpicture}
}
\par
\shortans[0]{0{,}8}
\loigiai
{
Ta có nửa đường chéo hình vuông có độ dài là $4$, cạnh hình vuông sẽ là $4\sqrt{2}$ và diện tích hình vuông là $32$, khi đó ta có được diện tích phần tô màu là $\dfrac{64}{5}$.\\
Gọi $f(x) = ax^3 + bx^2 - x$ là hàm số bậc ba biểu diễn đường cong trên logo.\\
Ta có $x = 4$ là nghiệm của phương trình nên $64a + 16b - 4 = 0 \Leftrightarrow 4a+b=1$.\quad $(1)$\\
Do đó phương trình $f(x) = 0$ sẽ có các nghiệm là $0$, $4$.\\
Khi đó diện tích hình phẳng $S$ giới hạn bởi các đường $y = f(x)$, trục $Ox$, đường thẳng $x = 0$, $x = 4$ là
\allowdisplaybreaks
\begin{eqnarray*}
S&= & \displaystyle \int \limits_{0}^{4} \left|ax^3+bx^2-x\right|\mathrm{\,d}x = -\displaystyle \int \limits_{0}^{4} \left(ax^3+bx^2-x\right)\mathrm{\,d}x\\
&= & -\left(\dfrac{ax^4}{4}+\dfrac{bx^3}{3}-\dfrac{x^2}{2}\right)\Bigg|_0^4 = -64a-\dfrac{64}{3}b+8.
\end{eqnarray*}
Mà $S = \dfrac{1}{8}\cdot \dfrac{64}{5} = \dfrac{8}{5}$ nên $-64a-\dfrac{64}{3}b+8 = \dfrac{8}{5} \Leftrightarrow 64a+\dfrac{64}{3}b=\dfrac{32}{5}$.\quad $(2)$\\
Từ $(1)$ và $(2)$, ta có $\heva{&4a+b=1\\&64a+\dfrac{64}{3}b=\dfrac{32}{5}} \Leftrightarrow \heva{&a = -\dfrac{1}{20}\\&b=\dfrac{9}{20}} \Rightarrow 2a+2b = \dfrac{4}{5}$.
}
\end{ex}

\begin{ex}%[Dự án C THPTQG 2025]%[Võ Hoàng Nghĩa]%[2D4V3-2]
Một ô tô bắt đầu chuyển động nhanh dần đều với vận tốc $v_1(t)=2t$ m/s. Đi được $12$ giây, người lái xe gặp chướng ngại vật và phanh gấp, ô tô tiếp tục chuyển động chậm dần đều với gia tốc $a=-12$ m/s$^2$. Tính quãng đường $s$ m đi được của ô tô từ lúc bắt đầu chuyển động đến khi dừng hẳn?
\par
\shortans[0]{168}
\loigiai{
\textbf{Giai đoạn 1:} Xe bắt đầu chuyển động đến khi gặp chướng ngại vật thì quãng đường xe đi được là
\[S_1=\displaystyle\int\limits_0^{12}v_1(t)\,\mathrm{d}t=\displaystyle\int\limits_0^{12}2t\,\mathrm{d}t=t^2\bigg|_0^{12}=144.\]
\textbf{Giai đoạn 2:} Xe gặp chướng ngại vật đến khi dừng hẳn.\\
Ô tô chuyển động chậm dần đều với vận tốc $v_2(t)=\displaystyle\int a(t)\,\mathrm{d}t=\displaystyle\int(-12)\,\mathrm{d}t=-12t+C$.\\
Vận tốc của xe khi gặp chướng ngại vật là
\[v_2(0)=v_1(12)=2\cdot12=24\Rightarrow -12\cdot0+C=24\Rightarrow C=24\Rightarrow v_2(t)=-12t+24.\]
Thời gian khi xe gặp chướng ngại vật đến khi dừng hẳn là nghiệm của phương trình
\[-12t+24=0\Leftrightarrow t=2.\]
Khi đó quãng đường xe đi được từ lúc đạp phanh đến khi dừng hẳn là
\[S_1=\displaystyle\int\limits_0^2v_2(t)\,\mathrm{d}t=\displaystyle\int\limits_0^2(-12t+24)\,\mathrm{d}t=\left(-6t^2+24t\right)\bigg|_0^2=24\text{ (m)}.\]
Vậy tổng quãng đường xe đi được là $S=S_1+S_2=144+24=168$ (m).
}
\end{ex}

\begin{ex}%[TT-THPT-Trần Đại Nghĩa-Đề 1-Đồng Nai-NH24-25-TN]%[2D4V3-2]
\immini[thm]{
Ông An xây dựng một sân bóng đá mini hình chữ nhật có chiều rộng $30$\,(m) và chiều dài $50$\,(m). Để giảm bớt kinh phí cho việc trồng cỏ nhân tạo, ông An chia sân bóng ra làm hai phần (tô màu và không tô màu) như hình vẽ.
\begin{itemize}
\item Phần tô màu gồm hai miền diện tích bằng nhau và đường cong $AIB$ là một parabol có đỉnh $I$.
\item Phần tô màu được trồng cỏ nhân tạo với giá $130$ nghìn đồng/m$^2$ và phần còn lại được trồng cỏ nhân tạo với giá $90$ nghìn đồng/m$^2$.
\end{itemize}
Hỏi số tiền ông An phải trả để trồng cỏ nhân tạo cho sân bóng là bao nhiêu triệu đồng?

}{
\begin{tikzpicture}[scale=0.8, font=\footnotesize, line join=round, line cap=round,>=stealth]
\def\a{1.5} \def\b{5} \def\c{1.5}
\path
(180:\a) coordinate(A)
(0:\a) coordinate(B)
($(A)+(90:\b)$) coordinate(D)
($(B)+(D)-(A)$) coordinate(C)
($(A)!1/5!(D)$) coordinate(M)
($(A)!4/5!(D)$) coordinate(N)
($(B)!1/5!(C)$) coordinate(P)
($(B)!4/5!(C)$) coordinate(Q)
($(M)!1/2!(P)$) coordinate(I)
($(N)!1/2!(Q)$) coordinate(J);
\draw (A)--(B)--(C)--(D)--cycle;
\draw[dashed] (M)--(P) (N)--(Q);
\begin{scope}
\clip (-1.5,0) rectangle (1.5,5);
\draw[samples=200,domain=-1.5:1.5,smooth,variable=\x] plot (\x,{(4/9)*(\x)^2+4});
\fill[pattern=north west lines,opacity=0.8,draw,thin] (-1.5,5)--plot[domain=-1.5:1.5](\x,{(4/9)*(\x)^2+4})--(1.5,5)--cycle;
\draw[samples=200,domain=-1.5:1.5,smooth,variable=\x] plot (\x,{-(4/9)*(\x)^2+1});
\fill[pattern=north west lines,opacity=0.8,draw,thin] (-1.5,0)--plot[domain=-1.5:1.5](\x,{-(4/9)*(\x)^2+1})--(1.5,0)--cycle;
\end{scope}
\foreach \x/\y  in{A/-150,B/-50,I/90}{\draw[fill=black](\x) circle(1pt)++(\y:0.35)node{$\x$};}
\draw[|<->|,transform canvas={yshift=-9pt}] (A)--(B) node[midway,below]{$30$ m};
\draw[|<->|,transform canvas={xshift=-9pt}] (A)--(D) node[midway,above,rotate=90]{$50$ m};
\draw[|<->|,transform canvas={xshift=9pt}] (C)--(Q) node[midway,below,rotate=90]{$10$ m};
\draw[|<->|,transform canvas={yshift=-9pt}] (N)--(J) node[midway,below]{$15$ m};
\end{tikzpicture}
}
\par
\shortans[0]{151}
\loigiai{
\immini{
Diện tích sân bóng là $S=50\cdot 30=1\,500$ m$^2$.\\
Đặt hệ trục toạ độ như hình vẽ.\\
Phương trình đường parabol đi qua $A$, $I$, $B$ là $y=-\dfrac{2}{45}x^2+10$.\\
Diện tích hình phẳng giới hạn bởi đường cong $AIB$ và đường thẳng $AB$ là
\[S=\displaystyle\displaystyle\displaystyle\int\limits_{-1{,}5}^{1{,}5}\left(-\dfrac{2}{45}x^2+10\right)\,\mathrm{d}x=200\;\left(\mathrm{m}^2\right).\]
Số tiền trải cỏ nhân tạo cho sân bóng là
\begin{eqnarray*}
T&=&2\cdot 200\cdot 130\,000+(1\,500-2\cdot 200)\cdot 90\,000\\
&=&151\,000\,000\,(\text{đồng})=151\,(\text{triệu đồng}).
\end{eqnarray*}
}
{\begin{tikzpicture}[scale=1, font=\footnotesize, line join=round, line cap=round,>=stealth]
\def\a{1.5} \def\b{5} \def\c{1.5}
\path
(0:0) coordinate(O)
(180:\a) coordinate(A)
(0:\a) coordinate(B)
($(A)+(90:\b)$) coordinate(D)
($(B)+(D)-(A)$) coordinate(C)
($(A)!1/5!(D)$) coordinate(M)
($(A)!4/5!(D)$) coordinate(N)
($(B)!1/5!(C)$) coordinate(P)
($(B)!4/5!(C)$) coordinate(Q)
($(M)!1/2!(P)$) coordinate(I)
($(N)!1/2!(Q)$) coordinate(J);
\draw (A)--(B)--(C)--(D)--cycle;
\draw[dashed] (M)--(P) (N)--(Q);
\begin{scope}
\clip (-1.5,0) rectangle (1.5,5);
\draw[samples=200,domain=-1.5:1.5,smooth,variable=\x] plot (\x,{(4/9)*(\x)^2+4});
\fill[pattern=north west lines,opacity=0.8,draw,thin] (-1.5,5)--plot[domain=-1.5:1.5](\x,{(4/9)*(\x)^2+4})--(1.5,5)--cycle;
\draw[samples=200,domain=-1.5:1.5,smooth,variable=\x] plot (\x,{-(4/9)*(\x)^2+1});
\fill[pattern=north west lines,opacity=0.8,draw,thin] (-1.5,0)--plot[domain=-1.5:1.5](\x,{-(4/9)*(\x)^2+1})--(1.5,0)--cycle;
\end{scope}
\draw[->]($(B)!1.3!(A)$)--($(A)!1.3!(B)$) node[below]{$x$};
\draw[->](-90:0.5)--(90:5.5) node[left]{$y$};
\foreach \x/\y  in{A/150,B/30,I/30,O/-45}{\draw[fill=black](\x) circle(1pt)++(\y:0.35)node{$\x$};}
\path
(-1.5,0) node[below]{$-15$}
(1.5,0) node[below]{$15$}
(0,1) node[above left]{$10$};
\end{tikzpicture}}
}
\end{ex}

\begin{ex}%[12-GHK2-2025, Tran Tony]%[2D4V3-2]
\immini{
Lô gô gắn tại các Showroom của hãng ô tô VINFAST là một hình tròn như hình vẽ. Phần tô đậm nằm giữa parabol đỉnh $I$ và đường gấp khúc $AJB$ được dát bạc với chi phí $10$ triệu đồng $\mathrm{m}^2$. Phần còn lại được phủ sơn với chi phí $2$ triệu đồng $\mathrm{m}^2$. Biết $AB=2\mathrm{~m}$, $IA=IB=\sqrt{5}\mathrm{~m}$ và $JA=JB=\dfrac{\sqrt{13}}{2}\mathrm{~m}$. Tính tổng số tiền (triệu đồng) để dát bạc và phủ sơn của lô gô nói trên (làm tròn đến hàng phần mười)?
\shortans[]{19,2}
}{
\begin{tikzpicture}[scale=1.3, font=\footnotesize, line join=round, line cap=round, >=stealth]
\path
(0,0) coordinate (I)
(-1,2) coordinate (A)
(1,2) coordinate (B)
(0,0.5) coordinate (J)
(0,1.25) coordinate (K)
;
\draw (0,1.8) node{\textbf{VINFAST}};
\draw (A)--(J)--(B);
\draw[samples=200,domain=-1:1,smooth,variable=\x] plot (\x,{2*(\x)^2});
\path let \p1=(K),\p2=(A),\n1={veclen(\x2-\x1,\y2-\y1)} in \pgfextra{\xdef\Temp{\n1}};
\draw[thick] (K) circle (\Temp);
\fill[gray]plot[samples=200,domain=-1:0,smooth,variable=\x] (\x,{-1.5*(\x)+0.5})--plot[samples=200,domain=0:-1,smooth,variable=\x] (\x,{2*(\x)^2})--cycle;
\draw[dashed]plot[samples=200,domain=-1:0,smooth,variable=\x] (\x,{-1.5*(\x)+0.5})--(0,0);
\draw plot[samples=200,domain=-1:0,smooth,variable=\x] (\x,{-1.5*(\x)+0.5});
\draw plot[samples=200,domain=-1:0,smooth,variable=\x] (\x,{2*(\x)^2});
\fill[gray]plot[samples=200,domain=0:1,smooth,variable=\x] (\x,{1.5*(\x)+0.5})--plot[samples=200,domain=1:0,smooth,variable=\x] (\x,{2*(\x)^2})--cycle;
\draw (0,0)--plot[samples=200,domain=0:1,smooth,variable=\x] (\x,{1.5*(\x)+0.5});
\draw plot[samples=200,domain=0:1,smooth,variable=\x] (\x,{1.5*(\x)+0.5});
\draw plot[samples=200,domain=0:1,smooth,variable=\x] (\x,{2*(\x)^2});
\foreach \x/\g in {A/180,B/0,I/-45,J/90}
\fill[black](\x) circle (1pt)	($(\x)+(\g:3mm)$) node{$\x$};
\end{tikzpicture}
}
\loigiai{
\immini{Chọn gốc tọa độ tại điểm $I(0; 0)$, tia $Iy$ trùng với tia $IJ$, ta có $A(-1; 2)$, $B(1; 2)$ và $J\left(0; \dfrac{1}{2}\right)$.\\
Phương trình parabol qua các điểm $I(0; 0)$, $A(-1; 2)$, $B(1; 2)$ là $y=2x^2$.\\
Phương trình đường tròn đi qua ba điểm $I(0; 0)$, $A(-1; 2)$, $B(1; 2)$ là \[x^2+\left(y-\frac{5}{4}\right)^2=\dfrac{25}{16}.\]
(có thể tính $R=\dfrac{AB \cdot IA \cdot IB}{4 S_{\triangle ABI}}=\dfrac{5}{4}$).\\
Diện tích cả hình tròn bằng $S=\pi R^2=\dfrac{25\pi}{16}$.
}{
\begin{tikzpicture}[scale=1.3, font=\footnotesize, line join=round, line cap=round, >=stealth]
\path
(0,0) coordinate (I)
(-1,2) coordinate (A)
(1,2) coordinate (B)
(0,0.5) coordinate (J)
(0,1.25) coordinate (K)
;
\draw[->] (-2,0)--(2,0) node[below] {$x$};
\draw[->] (0,-0.6)--(0,3) node[left] {$y$};
\draw (A)--(J)--(B);
\draw[samples=200,domain=-1:1,smooth,variable=\x] plot (\x,{2*(\x)^2});
\path let \p1=(K),\p2=(A),\n1={veclen(\x2-\x1,\y2-\y1)} in \pgfextra{\xdef\Temp{\n1}};
\draw[thick] (K) circle (\Temp);
\fill[gray]plot[samples=200,domain=-1:0,smooth,variable=\x] (\x,{-1.5*(\x)+0.5})--plot[samples=200,domain=0:-1,smooth,variable=\x] (\x,{2*(\x)^2})--cycle;
\draw[dashed]plot[samples=200,domain=-1:0,smooth,variable=\x] (\x,{-1.5*(\x)+0.5})--(0,0);
\draw plot[samples=200,domain=-1:0,smooth,variable=\x] (\x,{-1.5*(\x)+0.5});
\draw plot[samples=200,domain=-1:0,smooth,variable=\x] (\x,{2*(\x)^2});
\fill[gray]plot[samples=200,domain=0:1,smooth,variable=\x] (\x,{1.5*(\x)+0.5})--plot[samples=200,domain=1:0,smooth,variable=\x] (\x,{2*(\x)^2})--cycle;
\draw (0,0)--plot[samples=200,domain=0:1,smooth,variable=\x] (\x,{1.5*(\x)+0.5});
\draw plot[samples=200,domain=0:1,smooth,variable=\x] (\x,{1.5*(\x)+0.5});
\draw plot[samples=200,domain=0:1,smooth,variable=\x] (\x,{2*(\x)^2});
\foreach \x/\g in {A/180,B/0,I/-45,J/70}
\fill[black](\x) circle (1pt)	($(\x)+(\g:3mm)$) node{$\x$};
\end{tikzpicture}
}
\noindent
Đoạn thẳng $JA$ thuộc đường thẳng có phương trình $y=-\dfrac{3}{2}x+\dfrac{1}{2}$.\\
Vậy diện tích phần tô đậm bằng
\[S_1=2\displaystyle\int\limits_{-1}^0\left|2x^2-\left(-\dfrac{3}{2}x+\dfrac{1}{2}\right)\right| \mathrm{\,d}x=\dfrac{7}{6}.\]
Diện tích phần không tô đậm bằng $S_2=S-S_1=\dfrac{25}{16}\pi-\dfrac{7}{6}$.\\
Chi phí bằng $T=S_1 \cdot 10+2 S_2=\dfrac{7}{6}\cdot 10+\left(\dfrac{25 \pi}{16}-\dfrac{7}{6}\right) \cdot 2 \approx 19{,}2$ triệu đồng.
}
\end{ex}

\begin{ex}%[EX-3-2026, GHK2-12 THPT Nguyễn Quốc TRinh, Hà Nội]%[Huỳnh Xuân Tín]%[2D4V3-2]
\immini[thm]{Bạn Hải nhận thiết kế logo hình con mắt (phần được tô đậm) cho một cơ sở y tế. Logo là hình phẳng giới hạn bởi hai parabol $y=f(x)$ và $y=g(x)$ như hình vẽ (đơn vị trên mỗi trục tọa độ là decimet). Bạn Hải cần tính diện tích của logo để báo giá cho cơ sở y tế đó trước khi kí hợp đồng. Diện tích của logo là bao nhiêu decimet vuông (làm tròn kết quả đến hàng phần mười)?
\shortans{9,8}
}
{
\begin{tikzpicture}[scale=0.85,>=stealth, font=\footnotesize, line join=round, line cap=round]
\def\xmin{-4} \def\xmax{4}
\def\ymin{-2} \def\ymax{3}
\draw[->] (\xmin,0)--(\xmax,0) node [below]{$x$};
\draw[->] (0,\ymin)--(0,\ymax) node [left]{$y$};
\draw[fill=black] (0,0)node[below left]{$O$} circle(1pt);
\draw[smooth,samples=300,domain=-3.5:3.5] plot(\x,{(-0.25)*(\x)^2+2})node[above right]{$y=g(x)$};
\draw[smooth,samples=300,domain=-3.5:3.5] plot(\x,{0.25*(\x)^2-1})node[above right]{$y=f(x)$};
\foreach \d/\g in{-3/90,-2/-120,2/-90,3/90}
\draw[fill=black](\d,0)circle(1pt)node[shift={(\g:0.3)}]{$\d$};
\foreach \d/\g in{-1/-135,1/180,2/135}
\draw[fill=black](0,\d)circle(1pt)node[shift={(\g:0.3)}]{$\d$};
\draw[dashed] (2,0)--(2,1)--(0,1);
\fill[pattern=north east lines,opacity=0.4] plot[domain=-2.45:2.45](\x,{(-0.25)*(\x)^2+2})--plot[domain=+2.45:-2.45](\x,{0.25*(\x)^2-1});
\end{tikzpicture}
}
\loigiai{
Gọi phương trình parabol $f(x)=ax^2+bx+c$. Vì parabol có đỉnh nằm trên trục tung nên $b=0$. \\
Mặt khác đồ thị hàm số $f(x)$ đi qua $(0;-1)$ và $(2;0)$ nên ta có hệ phương trình
\[\heva{&c=-1\\&4a+c=0}\Leftrightarrow \heva{&a=\dfrac{1}{4}\\&c=-1.}\]
Vậy $f(x)=\dfrac{1}{4}x^2-1$.  \\
Tương tự, gọi phương trình parabol $g(x)=mx^2+nx+p$. Vì parabol có đỉnh nằm trên trục tung nên $n=0$. \\
Mặt khác đồ thị hàm số $g(x)$ đi qua $(0;2)$ và $(2;1)$ nên ta có hệ phương trình
\[\heva{&c=2\\&4a+c=1}\Leftrightarrow \heva{&a=-\dfrac{1}{4}\\&c=2.}\]
Vậy $g(x)=-\dfrac{1}{4}x^2+2$.  \\
Ta có $f(x)=g(x)\Leftrightarrow \dfrac{1}{4}x^2-1=-\dfrac{1}{4}x^2+2\Leftrightarrow x=\pm\sqrt{6}$. \\
Diện tích hình phẳng giới hạn bởi đồ thị hàm số $y=f(x)$ và $y=g(x)$ là
\[S=\int\limits_{-\sqrt{6}}^{\sqrt{6}}\left|\dfrac{1}{4}x^2-1-\left(\-\dfrac{1}{4}x^2+2\right)\right|\mathrm{\,d}x\approx 9{,}8.\]
Vậy diện tích của logo gần bằng $9{,}8$ dm$^2$.
}
\end{ex}

\begin{ex}%[12-EX-2-2025, Diamond]%[2D4V3-2]
\immini{
Cho hàm số $y=f(x)$. Đồ thị hàm số $y=f'(x)$ là đường cong trong hình bên. Biết rằng diện tích của các phần hình phẳng $A$ và $B$ lần lượt là $S_A=4$ và $S_B=10$. Tính giá trị của $f(3)$, biết giá trị của $f(0)=2$.
}{
\begin{tikzpicture}[scale=1, font=\footnotesize, line join=round, line cap=round, >=stealth]
\draw[->] (-.5,0)--(3.5,0)node[below]{$x$};
\draw[->] (0,-2.2)--(0,1.5)node[left]{$y$};
\fill[pattern=dots] (0,0)--plot[domain=0:1](\x,{(\x)*(\x-1)*(\x-3)})--cycle;
\fill[pattern=dots] (1,0)--plot[domain=1:3](\x,{(\x)*(\x-1)*(\x-3)})--cycle;
\draw[samples=100, domain=-.2:3.2] plot(\x,{(\x)*(\x-1)*(\x-3)})node[left]{$y=f'(x)$};
\foreach \i/\j in {1/60, 3/-45}
\draw[fill=black] (\i,0)node[shift=(\j:.32)]{$\i$}circle(1pt);
\draw[fill=black] (0,0)node[below left]{$O$}circle(1pt) (.5,.5)node[below, fill=white, rounded corners]{$A$} (2,-1)node[below, fill=white, rounded corners]{$B$};
\end{tikzpicture}
}
\shortans{$-4$}
\loigiai{
Từ hình vẽ ta có $\displaystyle\int\limits_0^1f'(x)\mathrm{\,d}x=4\Leftrightarrow f(x)\bigg|_0^1=4\Leftrightarrow f(1)-f(0)=4$.\hfill(1)\\
và $\displaystyle\int\limits_1^3f'(x)\mathrm{\,d}x=-10\Leftrightarrow f(x)\bigg|_1^3=-10\Leftrightarrow f(3)-f(1)=-10$.\hfill(2)\\
Cộng (1) và (2) ta được $f(3)-f(0)=-6\Leftrightarrow f(3)=-6+2=-4$.
}
\end{ex}

\begin{ex}%[Đề GHK2 Toán 12, THPT Minh Hà, Hà Nội, 2024-2025]%[Nguyễn Trung Kiên, dự án 12-EX3-2026]%[2D4V3-2]
\immini
{Một chiếc cổng có hình dạng là một parabol có khoảng cách giữa hai chân cổng là $8$ m. Người ta treo một tấm phông hình chữ nhật có hai đỉnh $M$, $N$ nằm trên parabol và hai đỉnh $P$, $Q$ nằm trên mặt đất. Ở phần phía ngoài phông người ta mua hoa để trang trí với chi phí $200\,000$ đồng/m$^2$, biết $MN=4$ m, $MQ=6$ m. Tính số tiền (đơn vị: triệu đồng) để mua hoa trang trí (làm tròn đến hàng phần mười).}
{\begin{tikzpicture}[scale=1, font=\footnotesize, line join=round, line cap=round, >=stealth, x=0.55cm, y=0.55cm]
\tikzset{declare function={f(\x)=-0.5*(\x-4)*(\x+4);}}
\draw[thick, blue] plot[smooth, samples=150, domain=-4:4] (\x,{f(\x)});
\draw[fill=cyan] (-2,0) coordinate(Q) -- (-2,6) coordinate(M) node[midway, sloped, above]{$6$\,m}
--(2,6) coordinate(N) node[midway, above]{$4$\,m}
--(2,0) coordinate(P) --cycle;
\path (-4,0) coordinate(A) (4,0) coordinate(B);
\begin{scope}
\clip (-5,0) rectangle (5,-0.5);
\draw (-5,0)--(5,0);
\foreach \i in {-5.5,-5.4,...,5} \draw[brown] (\i,0)--++(-70:1);
\end{scope}
\foreach \i/\j in {A/150, B/30, Q/150, P/30, M/135, N/45} \fill (\i) node[shift={(\j:3mm)}]{$\i$} circle(1pt);
\end{tikzpicture}}
\shortans{3,7}
\loigiai{
Chọn hệ trục tọa độ $Oxy$ sao cho gốc $O$ là trung điểm của $PQ$ (trên mặt đất), trục $Oy$ là trục đối xứng của cổng.
\begin{center}
\begin{tikzpicture}[scale=1, font=\footnotesize, line join=round, line cap=round, >=stealth, x=0.55cm, y=0.55cm]
\tikzset{declare function={f(\x)=-0.5*(\x-4)*(\x+4);}}
\draw[thick, blue] plot[smooth, samples=150, domain=-4:4] (\x,{f(\x)});
\draw[fill=cyan] (-2,0) coordinate(Q) -- (-2,6) coordinate(M)
--(2,6) coordinate(N)
--(2,0) coordinate(P) --cycle;
\path (-4,0) coordinate(A) (4,0) coordinate(B);
\draw[->] (-5,0)--(5.5,0) node[below]{$x$};
\draw[->] (0,-1)--(0,9) node[left]{$y$};
\foreach \i/\j in {A/150, B/30, Q/150, P/30, M/135, N/45} \fill (\i) node[shift={(\j:3mm)}]{$\i$} circle(1pt);
\path (A) node[below]{$-4$}
(B) node[below]{$4$}
(Q) node[below]{$-2$}
(P) node[below]{$2$}
(0,6) node[above left]{$6$};
\fill (0,0) node[below left]{$O$} circle(1pt);
\end{tikzpicture}
\end{center}
Cổng có dạng là parabol $(P)\colon y=ax^2+c$ đi qua hai điểm $A(-4,0)$ và $M(-2,6)$.\\
Do đó, ta có $\heva{& a(-4)^2+c=0\\& a(-2)^2+c=6} \Leftrightarrow \heva{& a=-\dfrac{1}{2}\\& c=8}$,
suy ra $(P)\colon y=-\dfrac{1}{2}x^2+8$.\\
Diện tích của toàn bộ phần cổng là
\[S=\displaystyle\int\limits_{-4}^{4} \left(-\dfrac{1}{2}x^2+8\right) \mathrm{\,d}x = \left(-\dfrac{x^3}{6}+8x\right)\Bigg|_{-4}^4 = \dfrac{128}{3}\, \left(\text{m}^2\right).\]
Diện tích tấm phông hình chữ nhật $MNPQ$ là $S_1 = 4\cdot 6 =24\, \left(\text{m}^2\right)$.\\
Diện tích cần trang trí hoa: $S_2 = S-S_1=\dfrac{128}{3}-24 =\dfrac{56}{3}\, \left(\text{m}^2\right)$.\\
Vậy chi phí mua hoa trang trí là $T=S_2\cdot 0{,}2 = \dfrac{56}{15} \approx 3{,}7$ (triệu đồng).
}
\end{ex}

\begin{ex}%[12-EX-3-2025, Nguyễn Chiến]%[2D4V3-2]
Ông An muốn thiết kế một cái cổng có kích thước như hình vẽ. Vòm cổng có hình dạng một parabol có đỉnh $I(0;2)$ và đi qua điểm $B\left(\dfrac{5}{2};\dfrac{3}{2}\right)$. Biết đơn giá vật liệu làm cổng là $1{,}3$ triệu/m$^2$. Hỏi ông An phải chi phí hết bao nhiêu triệu đồng để hoàn thành cổng (làm tròn kết quả đến hàng phần chục)?
\shortans{$11{,}9$}
\begin{center}
\begin{tikzpicture}[>=stealth, scale=1, font=\footnotesize]
\def \xmin{-3}
\def \xmax{3.5}
\def \ymin{-0.5}
\def \ymax{3}
\clip (\xmin,\ymin) rectangle (\xmax+0.1,\ymax);
\draw[->] (\xmin,0)--(\xmax,0) node[below left=0 and -0.15] {$x$};
\draw[->] (0,\ymin)--(0,\ymax) node[below right=-0.12 and 0] {$y$};
\draw[smooth,samples=200,domain=-2.5:2.5] plot (\x,{(-2/25)*(\x)^2+2});
\draw[fill=black] (0,0) node[below left=-0.1] {$O$} circle (1pt);
\draw[fill=black] (2.5,0) node[below] {$\tfrac{5}{2}$} circle (1pt);
\draw[fill=black] (-2.5,0) node[below] {$\tfrac{-5}{2}$} circle (1pt);
\draw[fill=black] (-2.5,1.5) node[above] {$A$};
\draw[fill=black] (2.5,1.5) node[above] {$B$};
\draw[fill=black] (2.5,0) node[above right=-0.1] {$C$};
\draw[fill=black] (-2.5,0) node[above left=-0.1] {$D$};
\draw[fill=black] (0,2) node[above right] {$2$} circle (1pt);
\draw[fill=black] (0,1.5) node[below right] {$\tfrac{3}{2}$} circle (1pt);
\draw[dashed] (-2.5,0)--(-2.5,1.5) (2.5,0)--(2.5,1.5);
\draw[dashed] (-2.5,1.5)--(2.5,1.5);
\end{tikzpicture}

\end{center}
\loigiai{
Giả sử phương trình parabol là $(P) \colon y=ax^2+bx+c$.\\
Vì parabol có đỉnh $I(0;2)$ và đi qua $B\left(\dfrac{5}{2};\dfrac{3}{2}\right)$ nên\\
\[\heva{&-\dfrac{b}{2a}=0\\&a\cdot 0^2+b\cdot 0+c=2\\&a \cdot \left(\dfrac{5}{2}\right)^2+b \cdot \dfrac{5}{2}+c=\dfrac{3}{2}} \Leftrightarrow \heva{&a=-\dfrac{2}{25}\\&b=0\\&c=2.}\]
Do đó $(P) \colon y=-\dfrac{2}{25}x^2+2$.\\
Diện tích phần hình chữ nhật $ABCD$ là $S_1=5\cdot\dfrac{3}{2}=\dfrac{15}{2}$.\\
Phương trình cạnh $AB$ là $y=\dfrac{3}{2}$.\\
Diện tích phần vòm là $S_2=\displaystyle\int \limits_{-\frac{5}{2}}^{\frac{5}{2}}\left(-\dfrac{2}{25}x^2+2-\dfrac{3}{2}\right)\mathrm{\,d}x=\dfrac{5}{3}$.\\
Suy ra $S=S_1+S_2=\dfrac{55}{6}$ m$^2$.\\
Chi phí là $C=1{,}3\cdot\dfrac{55}{6}=\dfrac{715}{60}\approx  11{,}9$ triệu đồng.
}
\end{ex}

\begin{ex}%[12-EX3-2026. Trần Hoà]%[2D4V3-2]
\immini{Bác Bình muốn làm một cái cửa bằng inox hình parabol có chiều cao từ mặt đất đến đỉnh là $3$ mét, chiều rộng tiếp giáp với mặt đất là $3$ mét. Biết rằng giá vật liệu và tiền công mỗi mét vuông là $1\,700\,000$ đồng. Vậy bác Bình phải trả bao nhiêu tiền để làm cái cửa đó (đơn vị triệu đồng)? (Làm tròn kết quả đến hàng phần mười).
\shortans[]{$10{,}2$}}{\begin{tikzpicture}[scale=1,font=\footnotesize,line join = round, line cap = round, >= stealth]
\draw (0,0) parabola bend ((3/2,3)(3,0);%hoanh-tang-dan
\def\a{-.5} \def\b{3.5} \def\c{-.5} \def\d{3.5}
\coordinate (O) at (0,0);
\draw (0,0)--(3,0);
\fill[pattern=dots] (0,0) parabola bend ((3/2,3)(3,0)--cycle;
\draw[<->] ($(0,0)+(-90:0.5)$)--($(3,0)+(-90:0.5)$) node[pos=0.5,above]{$3$m};
\draw[<->] ($(3,0)+(0:0.5)$)--($(3,3)+(0:0.5)$) node[pos=0.5,left]{$3$m};
\end{tikzpicture}}



\loigiai{
\immini{Chọn hệ trục toạ độ như hình vẽ.\\
Parabol cắt trục hoành tia $0$ và $3$ nên phương trình có dạng
$y=ax(x-3)$.\\
Lại có  đỉnh $I\left(\dfrac{3}{2};3\right)$ nên $3=-\dfrac{9}{4}a\Rightarrow a=-\dfrac{4}{3}$.\\
Vậy $(P)\colon y=-\dfrac{4}{3}x^2+4x$.\\
Diện tích cửa là $S=\displaystyle\int\limits_{0}^{3}\left(-\dfrac{4}{3}x^2+4x\right) \mathrm{\,d}x = \left(-\dfrac{4}{9}x^3+2x^2\right)\bigg|_0^3=6$ (m$^2$).\\
Số tiền để làm cửa là $6 \cdot 1\,700\,000 = 10\,200\,000$ đồng $= 10{,}2$ triệu đồng.}{\begin{tikzpicture}[scale=1,font=\footnotesize,line join = round, line cap = round, >= stealth]
\draw (0,0) parabola bend ((3/2,3)(3,0);%hoanh-tang-dan
\def\a{-.5} \def\b{3.5} \def\c{-.5} \def\d{3.5}
\coordinate (O) at (0,0);
\draw[->] (\a,0)--(\b,0)node[below]{$x$};
\draw[->] (0,\c)--(0,\d)node[right]{$y$};
\foreach \p/\g in {3/-90} \draw[fill] (\p,0) circle(.5pt)node [shift={(\g:.3)}] {$\p$};
\foreach \p/\g in {3/180} \draw[fill] (0,\p) circle(.5pt)node [shift={(\g:.3)}] {$\p$};
\draw[dashed] (3/2,0)|-(0,3) ;
\draw[fill] (O) circle (.5 pt) node[above right]{$O$} ;
\end{tikzpicture}}
}
\end{ex}

\begin{ex}%[12-EX-2-2026, Lương Như Quỳnh]%[2D4V3-2]
Nhà bác An có một bồn hoa dạng hình tròn bán kính bằng $6$ m. Bác An chia bồn hoa thành các phần như hình vẽ dưới đây và bác định bố trí như sau: Phần diện tích bên trong hình vuông $ABC D$ để trồng hoa. Phần diện tích kéo dài từ bốn cạnh của hình vuông đến đường tròn (phần gạch chéo) dùng để trồng cỏ. Ở bốn góc còn lại, mỗi góc trồng một cây cau voi. Biết $AB=4$ m; giá trồng hoa là $200$ nghìn đồng m$^2$; giá trồng cỏ là 100 nghìn đồng m$^2$; giá mỗi cây cau voi là $500$ nghìn đồng. Hỏi số tiền để bác An thực hiện việc trang trí bồn hoa như trên là bao nhiêu (đơn vị là triệu đồng, làm tròn kết quả đến hàng phần mười)?
\begin{center}
\begin{tikzpicture}[>=stealth,line join=round,line cap=round,font=\footnotesize,scale=0.5]
\def\R{6}
\path
(0,0) coordinate (O);
\draw[domain=-2:2,samples=250,smooth,pattern = north east lines] plot(\x,{sqrt(36-(\x)^2)})--plot[domain=2:-2](\x,{-sqrt(36-(\x)^2)})--cycle;
\draw[rotate=90,domain=-2:2,samples=250,smooth,pattern = north east lines] plot(\x,{sqrt(36-(\x)^2)})--plot[domain=2:-2](\x,{-sqrt(36-(\x)^2)})--cycle;
\draw (0,0) circle (\R cm);
\draw[fill=green!30] (-2,-2) rectangle (2,2);
\end{tikzpicture}
\end{center}
\shortans[]{11{,}4}
\loigiai{
Đặt hệ trục tọa độ $Oxy$ như hình vẽ, tâm bồn hoa trùng gốc tọa độ $O$.
\begin{center}
\begin{tikzpicture}[>=stealth,line join=round,line cap=round,scale=0.5]%font=\footnotesize,
\def\R{6}
\path
(0,0) coordinate (O);
\draw[domain=-2:2,samples=250,smooth,pattern = north east lines] plot(\x,{sqrt(36-(\x)^2)})--plot[domain=2:-2](\x,{-sqrt(36-(\x)^2)})--cycle;
\draw[rotate=90,domain=-2:2,samples=250,smooth,pattern = north east lines] plot(\x,{sqrt(36-(\x)^2)})--plot[domain=2:-2](\x,{-sqrt(36-(\x)^2)})--cycle;
\draw (0,0) circle (\R cm);
\draw[fill=green!30] (-2,-2) rectangle (2,2);
\def\xmin{-7.5} \def\xmax{7.5}
\def\ymin{-7.5} \def\ymax{7.5}
\draw[->] (\xmin,0)--(\xmax,0) node [above left]{$x$};
\draw[->] (0,\ymin)--(0,\ymax) node [below right]{$y$};
\node at (0,0) [below right]{$O$};
%\foreach \x in{-5,...,-3,-1,1,3,4,5} \draw[fill=black] (\x,0) circle (.5pt) node[below]{\tiny $\x$};
\draw[fill=black] (6,0) circle (.5pt) node[below right]{$6$};
\draw[fill=black] (-6,0) circle (.5pt) node[below left]{$-6$};
\draw[fill=black] (2,0) circle (.5pt) node[below right]{$2$};
\draw[fill=black] (-2,0) circle (.5pt) node[below left]{$-2$};
%\foreach \y in{-5,...,-3,-1,1,3,4,5}\draw[fill=black] (0,\y) circle (.5pt) node[left]{\tiny $\y$};
\draw[fill=black] (0,6) circle (.5pt) node[above left]{$6$};
\draw[fill=black] (0,-6) circle (.5pt) node[below left]{$-6$};
\draw[fill=black] (0,2) circle (.5pt) node[above left]{$2$};
\draw[fill=black] (0,-2) circle (.5pt) node[below left]{$-2$};
\end{tikzpicture}
\end{center}
Phương trình đường tròn tâm $O$, bán kính $6$ là
$x^2+y^2=36$.\\
Hình vuông có tâm $O$, cạnh $4$.
\begin{itemize}
\item Diện tích trồng hoa là $S_{\text{hoa}}=16$ m$^2$.\\
Tiền trồng hoa là $16\cdot 200=3200$ nghìn đồng $=3{,}2$  triệu đồng.
\item Diện tích trồng cỏ là $4\displaystyle\int\limits_{-2}^2\left(\sqrt{36-x^2}-2\right)\mathrm{\,d}x\approx 62$ m$^2$.\\
Tiền trồng cỏ là $62\cdot 100=6\,200$ nghìn đồng $=6{,}2$ triệu đồng.
\item Tiền trồng $4$ cây cau voi là $
4\cdot 500=2000 \text{ nghìn đồng}=2$ triệu đồng.
\end{itemize}
Vậy số tiền để bác An thực hiện trang trí bồn hoa là
$3{,}2+6{,}2+2{,}0=11{,}4$ triệu đồng.
}
\end{ex}

\begin{ex}%[Đề thi GHK2, THPT Nguyễn Hữu Huân, TPHCM năm học 2024-2025]%[BCTuan, dự án 12EX3-26]%[2D4V3-2]
Một bức tường hình chữ nhật $ABCD$ có $AB=6$ m, $BC=4$ m được bạn An trang trí bằng cách vẽ một parabol có trục đối xứng là trục đối xứng của hình chữ nhật và đỉnh của parabol cách $CD$ một đoạn là $1$ m, các nhánh parabol chia hình chữ nhật thành ba phần $H_1$, $H_2$, $H_3$ như hình vẽ bên dưới ($HK=2$ m).
\begin{center}
\begin{tikzpicture}[line join=round, line cap=round,>=stealth,font=\footnotesize,scale=1]
\path
(-3,4) coordinate (A)
(3,4) coordinate (B)
(-3,0) coordinate (D)
(3,0) coordinate (C)
({sqrt(5)},4) coordinate (I)
({-sqrt(5)},4) coordinate (P)
(1,0) coordinate (K)
(-1,0) coordinate (H);
\begin{scope}
\clip (-3,0) rectangle (3,4);
\fill[yellow!30] (-3,0) rectangle (3,4);
\fill[cyan!30]
(A) -- (D) -- (H)
-- plot[domain=-1:-3, samples=100] (\x,{(\x)^2-1})
-- (P) -- cycle;
\fill[green!50]
(B) -- (C) -- (K)
-- plot[domain=1:3, samples=100] (\x,{(\x)^2-1})
-- (I) -- cycle;
\end{scope}
\begin{scope}
\clip (-3,0) rectangle (3,4);
\draw[samples=200,domain=-3:3,smooth,variable=\x] plot (\x,{1*(\x)^2-1});
\end{scope}
\draw (A)--(B)node[above,midway,sloped]{$6$\,m}--(C)node[right,midway]{$4$\,m}--(K)--(H)node[below,midway,sloped]{$2$\,m}--(D)--(A)
(-2.2,2)node{$H_1$}
(0,2)node{$H_2$}
(2.2,2)node{$H_3$}
;
\foreach \i/\g in {A/135,B/45,C/-45,D/-135,K/-90,H/-90}{\draw[fill=black](\i) circle (1pt) ($(\i)+(\g:3mm)$) node[scale=1]{$\i$};}
\end{tikzpicture}
\end{center}
Phần $H_1$ được sơn màu xanh da trời, phần $H_2$ được sơn màu vàng, phần $H_3$ được sơn màu xanh lá cây. Biết rằng mỗi hộp sơn các màu chỉ sơn được $3$ m$^2$ tường, đồng thời giá của hộp sơn màu xanh da trời là $120\,000$ đồng/hộp, hộp sơn màu vàng là $130\,000$ đồng/hộp, hộp sơn màu xanh lá cây là $110\,000$ đồng/hộp. Tính giá tiền bạn An phải trả để sơn bức tường này. Biết rằng cửa hàng sơn chỉ bán số nguyên hộp (đơn vị là triệu đồng và làm tròn kết quả đến hàng phần trăm).
\shortans[oly]{$1{,}11$}
\loigiai{
\begin{center}
\begin{tikzpicture}[line join=round, line cap=round,>=stealth,font=\footnotesize,scale=1]
\path
(0,0) coordinate (O)
(-3,4) coordinate (A)
(3,4) coordinate (B)
(-3,0) coordinate (D)
(3,0) coordinate (C)
(0,-1) coordinate (E)
({sqrt(5)},4) coordinate (I)
({sqrt(5)},0) coordinate (J)
({-sqrt(5)},4) coordinate (P)
(1,0) coordinate (K)
(-1,0) coordinate (H);
\begin{scope}
\clip (-3,0) rectangle (3,4);
\fill[yellow!30] (-3,0) rectangle (3,4);
\fill[cyan!30]
(A) -- (D) -- (H)
-- plot[domain=-1:-3, samples=100] (\x,{(\x)^2-1})
-- (P) -- cycle;
\fill[green!50]
(B) -- (C) -- (K)
-- plot[domain=1:3, samples=100] (\x,{(\x)^2-1})
-- (I) -- cycle;
\end{scope}
\begin{scope}
\clip (-3,-2) rectangle (3,4);
\draw[samples=200,domain=-3:3,smooth,variable=\x] plot (\x,{1*(\x)^2-1});
\end{scope}
\draw[->] (-4,0)--(4,0) node[below left] {$x$};
\draw[->] (0,-1.5)--(0,5.1) node[below left] {$y$};
\draw (A)--(B)--(C)--(K)--(H)--(D)--(A)
;
\draw[dashed] (I)--(J);
\fill (P)circle (1pt);
\foreach \i/\g in {A/135,B/45,C/-45,D/-135,K/-80,H/-120,I/90,J/-90,E/-120,O/-45}{\draw[fill=black](\i) circle (1pt) ($(\i)+(\g:3mm)$) node[scale=1]{$\i$};}
\end{tikzpicture}
\end{center}
Ta thiết lập hệ trục tọa độ $Oxy$ với $E(0;-1)$ là đỉnh của parabol như hình vẽ.\\
Gọi parabol cần tìm có dạng $y=ax^2+c$ với ($a\ne 0$).\\
Do parabol đi qua điểm $K(1;0)$ và có đỉnh $E(0;-1)$ nên suy ra $a=1$ và $c=-1$.\\
Do đó parabol là $y=x^2-1$.\\
Gọi giao điểm của parabol với đường thẳng $AB$ là điểm $I$ $ (IB<IA)$.\\
Xét phương trình hoành độ giao điểm $x^2-1=4\Leftrightarrow \hoac{& x=\sqrt{5} \\ & x=-\sqrt{5}.}$\\
Suy ra điểm $I\left(\sqrt{5};4\right)$.\\
Gọi hình chiếu vuông góc của điểm $I$ lên $DC$ là $J$.\\
Ta tính diện tích $S_0$ của hình phẳng giới hạn bởi các đường thẳng $x=\sqrt{5}$, $y=0$ và parabol.\\
\[S_0=\displaystyle\int\limits_1^{\sqrt{5}} \left( x^2-1\right) \mathrm{\,d}x=\left( \dfrac{x^3}{3}-x\right) \bigg|_1^{\sqrt{5}}=\dfrac{2\sqrt{5}+2}{3}\ (\text{m}^2).\]
Suy ra diện tích hình $H_3$ và $H_1$ là
\[S_{H_3}=S_{H_1}=S_0+S_{IJCB}=\dfrac{2\sqrt{5}+2}{3}+4\cdot \left(3-\sqrt{5} \right)\approx 5{,}21\ (\text{m}^2).\]
Vậy với mỗi diện tích $H_1$ và $H_3$ bạn An cần mua $2$ hộp sơn mỗi loại.\\
Diện tích hình $H_2$ là
\[S_{H_2}=S_{ABCD}-2S_{H_3}\approx 6\cdot 4-2\cdot 5{,}21=13{,}58\ (\text{m}^2).\]
Vậy với diện tích hình $H_2$ bạn An cần mua $5$ hộp sơn màu vàng.\\
Tổng số tiền bạn An cần để mua sơn là
\[T=2\cdot 120\,000 + 2\cdot 110\,000 + 5\cdot 130\,000=1\,110\,000\,\,(\text{đồng})=1{,}11\,\ (\text{triệu đồng}).\]
}
\end{ex}

\begin{ex}%[Giữa học kì 2, Trần Cao Vân, Khánh Hòa, 2025]%[12-EX-3-2026, Thành Đức Trung]%[2D4V3-2]
Một khuôn viên dạng nửa hình tròn, trên đó người thiết kế phần để trồng hoa có dạng của một cánh hoa hình parabol có đỉnh trùng với tâm và có trục đối xứng vuông góc với đường kính của nửa hình tròn, hai đầu mút của cánh hoa nằm trên nửa đường tròn (phần tô màu) và cách nhau một khoảng bằng $4$ (m).
\immini
{
Phần còn lại của khuôn viên (phần không tô màu) dành để trồng cỏ Nhật Bản. Biết các kích thước cho như hình vẽ, chi phí để trồng hoa và cỏ Nhật Bản tương ứng là $150.000$ đồng/m$^2$ và $100.000$ đồng/m$^2$. Hỏi cần bao nhiêu tiền (đơn vị: triệu đồng) để trồng hoa và trồng cỏ Nhật Bản trong khuôn viên đó \textit{(làm tròn kết qủa đến hàng phần trăm)}?
}
{
\begin{tikzpicture}[scale=0.7, font=\footnotesize, line join=round, line cap=round, >=stealth]
\draw (4.47,0) arc [ start angle=0, end angle=180, x radius=4.47,y radius=4.47];
\draw[samples=100,domain=-2:2] plot (\x,{(\x)^2});
\fill[pattern=north west lines] plot[smooth,samples=200,domain=-2:2] (\x,{sqrt(20-(\x)^2)})--plot[smooth,samples=200,domain=2:-2] (\x,{(\x)^2});
\draw (-4.47,0)--(4.47,0);
\draw[dashed] (-2,0)--(-2,4) node[pos=0.5,above,rotate=90]{$4$ m};
\draw[dashed] (-2,4)--(2,4) node[pos=0.5,rotate=0,circle,fill=white,inner sep=0.5]{$ 4$ m};
\draw[dashed] (2,4)--(2,0) node[pos=0.5,above,rotate=270]{$4$ m};
\foreach \x in {-4.47,-2,0,2,4.47} \draw(\x,0) circle (1pt);
\end{tikzpicture}
}
\shortans[]{$3{,}74$}
\loigiai
{
\begin{center}
\begin{tikzpicture}[scale=0.8, font=\footnotesize, line join=round, line cap=round, >=stealth]
\draw[->] (-5,0)--(5,0) node[below]{$x$};
\draw[->] (0,-0.5)--(0,5) node[left]{$y$};
\draw[fill=white] (0,0)coordinate (O) circle(1pt)node[below left]{$O$};
\draw (4.47,0) arc [ start angle=0, end angle=180, x radius=4.47,y radius=4.47];
\draw[samples=100,domain=-2:2] plot (\x,{(\x)^2});
\fill[pattern=north west lines] plot[smooth,samples=200,domain=-2:2] (\x,{sqrt(20-(\x)^2)})--plot[smooth,samples=200,domain=2:-2] (\x,{(\x)^2});
\draw (-4.47,0)--(4.47,0);
%\draw[dashed] (-2,0)--(-2,4)--(2,4)--(2,0);
\draw[dashed] (-2,0)--(-2,4) node[pos=0.5,above,rotate=90]{$ 4$ m};
\draw[dashed] (-2,4)--(2,4) node[pos=0.5,rotate=0,circle,fill=white,inner sep=0.5]{$ 4$ m};
\draw[dashed] (2,4)--(2,0) node[pos=0.5,above,rotate=270]{$ 4$ m};
\foreach \x in {-2,2} \draw(\x,0) circle (1pt) node[below]{$\x$};
\end{tikzpicture}
\end{center}
Phương trình của nửa đường tròn $(C)$ là $x^2+y^2=20$, $y\geqslant0$, suy ra $y=\sqrt{20-x^2}$.\\
Parabol $(P)$ có đỉnh $O(0;0)$ và đi qua điểm $(2;4)$ nên có phương trình $y=x^2$.\\
Diện tích phần tô màu là $S_1=\displaystyle\int\limits_{-2}^2 \left[\sqrt{20-x^2}-x^2\right]\mathrm{\,d}x\approx11{,}9396$ (m$^2$). \\
Diện tích phần không tô màu là $S_2=\dfrac{1}{2}\cdot\pi\cdot\left(2\sqrt{5}\right)^2-S_1\approx10\pi-11{,}9396$ (m$^2$). \\
Số tiền để trồng hoa và trồng cỏ Nhật Bản trong khuôn viên đó là
\[150\,000\cdot11{,}9\,396+100\,000\cdot\left(10\pi -11{,}9396\right)\approx3\,738\,573\,\left(\text{đồng}\right).\]
Vậy cần $3{,}74$ triệu đồng để làm khuôn viên.
}
\end{ex}

\begin{ex}%[GK2, Thpt Trần Quốc Tuấn PY, 2025]%[Bùi Anh Tuấn, dự án EX3]%[2D4V3-2]
Một công ty quảng cáo $X$ muốn làm bức tranh trang trí hình $MNEF$ ở chính giữa của một bức tường hình chữ nhật $ABCD$ có chiều cao $BC=6$ m, chiều dài  $CD=12$ m (hình vẽ bên dưới). Cho biết $MNEF$ là hình chữ nhật có $MN=4$ m, cung $EIF$ có hình dạng là một phần của cung parabol có đỉnh $I$ là trung điểm của cạnh $AB$ và đi qua $2$ điểm $C$, $D$. Kinh phí làm bức tranh là $600\, 000$ đồng/m$^2$.
\begin{center}
\begin{tikzpicture}[scale=0.5, font=\footnotesize, line join=round, line cap=round, >=stealth]
\draw (-6,0)--(6,0);
\draw [domain=-6:6, smooth,samples=100] plot (\x, {(-(1/6)*(\x)^2+6)}) node [above right] {};
\fill [pattern=north east lines,pattern color=black,domain=-2:2,smooth,samples=100, variable=\x] (-2,0) -- plot (\x,{(-(1/6)*(\x)^2+6)}) -- (2,0);
\draw (-6,0)--(-6,6)--(6,6)--(6,0) (-2,0)--(-2,16/3) (2,16/3)--(2,0);
\draw
(0,6) node [above] {$I$} circle(1pt)
(-6,6) node [above left] {$A$} circle(1pt)
(6,6) node [above] {$B$} circle(1pt)
(6,0) node [below] {$C$} circle(1pt)
(-6,0) node [below] {$D$} circle(1pt)
(-2,0) node [below] {$M$} circle(1pt)
(2,0) node [below] {$N$} circle(1pt)
(2,16/3) node [right] {$E$} circle(1pt)
(-2,16/3) node [left] {$F$} circle(1pt)
(0,0) node [below] {$4$ m}
(6,3) node [right] {$6$ m};
\draw[<-, dashed] (-6,7)--(0,7)node [above] {$12$ m};
\draw[->,dashed] (0,7)--(6,7);
\end{tikzpicture}
\end{center}
Hỏi công ty $X$ cần bao nhiêu triệu đồng để làm bức tranh đó (kết quả làm tròn đến hàng phần mười)?
\shortans[]{$13{,}9$}
\loigiai{
Chọn hệ trục tọa độ $Oxy$ như hình vẽ, gốc tọa độ $O$ là trung điểm của $CD$ và trục $Oy$ trùng với đường thẳng $OI$.\\
Khi đó $I(0;6)$, $C(6;0)$, $D(-6;0)$, $M(-2;0)$, $N(2;0)$.\\
Gọi $(P)\colon y=ax^2+bx+c, \, a\neq 0$ là parabol có đỉnh là $I$ và đi qua $C$, $D$.\\
Suy ra $\heva{&c=6\\&36a+6b+c=0\\&36a-6b+c=0}\Leftrightarrow \heva{&a=-\dfrac{1}{6}\\&b=0\\&c=6}\Rightarrow y=-\dfrac{1}{6}x^2+6$.
\begin{center}
\begin{tikzpicture}[scale=0.5, font=\footnotesize, line join=round, line cap=round, >=stealth]
\draw [->] (-7,0)--(7,0) node [below] {$x$};
\draw [->] (0,-1) -- (0,8) node [left] {$y$};
\draw [domain=-6:6, smooth,samples=100] plot (\x, {(-(1/6)*(\x)^2+6)}) node [above right] {};
\fill [pattern=north east lines,pattern color=black,domain=-2:2,smooth,samples=100, variable=\x] (-2,0) -- plot (\x,{(-(1/6)*(\x)^2+6)}) -- (2,0);
\draw (-6,0)--(-6,6)--(6,6)--(6,0) (-2,0)--(-2,16/3) (2,16/3)--(2,0);
\draw  (0,0) node [below left] {$O$} circle(1pt)
(0,6) node [above left] {$I$} circle(1pt)
(-6,6) node [above left] {$A$} circle(1pt)
(6,6) node [above] {$B$} circle(1pt)
(6,0) node [below] {$C$} circle(1pt)
(-6,0) node [below] {$D$} circle(1pt)
(-2,0) node [below] {$M$} circle(1pt)
(2,0) node [below] {$N$} circle(1pt)
(2,16/3) node [right] {$E$} circle(1pt)
(-2,16/3) node [left] {$F$} circle(1pt)
;
\end{tikzpicture}
\end{center}
Diện tích bức tranh là
\[S=\displaystyle\int\limits_{-2}^2 \left|-\dfrac{1}{6}x^2+6\right|\mathrm{\,d}x=\displaystyle\int\limits_{-2}^2 \left(-\dfrac{1}{6}x^2+6\right)\mathrm{\,d}x=\left(-\dfrac{x^3}{18}+6x\right)\Bigg|_{-2}^2=\dfrac{208}{9}\,\text{m}^2.\]
Vì kinh phí $600\,000$ đồng/m$^2=0{,}6$ triệu đồng/m$^2$ nên số tiền cần để trang trí bức tranh là
\begin{center}
$\dfrac{208}{9}\cdot 0{,}6\approx 13{,}9$ triệu đồng.
\end{center}
}
\end{ex}

\begin{ex}%[12-EX-2-2025, Đình Nguyên]%[2D4V3-2]
Mặt cắt của một hầm có dạng là hình phẳng giới hạn bởi một parabol và đường thẳng nằm ngang như hình vẽ. Tính diện tích của cửa hầm (đơn vị là $m^2$, kết quả làm tròn đến hàng phần mười).
\begin{center}
\begin{tikzpicture}[scale=0.5, font=\footnotesize, line join=round, line cap=round, >=stealth]
\draw [fill=gray!30] (-6,0) rectangle (6,7);
\draw [fill=white] plot[domain=-5:5] (\x, {-0.26*\x*\x + 6.5}) -- cycle;
\draw [dashed] (-5,0) -- (-5,-1) (5,0) -- (5,-1);
\draw [<->] (-5,-0.8) -- (5,-0.8) node[midway, fill=white] {$10$ m};
\draw [dashed] (0,6.5) -- (7,6.5) (0,0) -- (7,0);
\draw [<->] (6.5,0) -- (6.5,6.5) node[midway, fill=white, rotate=90] {$6{,}5$ m};
\end{tikzpicture}
\end{center}
\shortans{43,3}
\loigiai{
Chọn hệ trục tọa độ $Oxy$ sao cho đỉnh parabol là $I(0; 6{,}5)$ và hai giao điểm với trục hoành là $(-5;0), (5;0)$.\\
Phương trình parabol có dạng $y = ax^2 + 6{,}5$. Thay $x=5, y=0 \Rightarrow a = -0{,}26$.\\
Diện tích cửa hầm là $S = \displaystyle\int\limits_{-5}^5 \left(-0{,}26x^2 + 6{,}5\right) \mathrm{\,d}x = \dfrac{130}{3} \approx 43{,}3$ m$^2$.
}
\end{ex}

\begin{ex}%[12-EX-3-2026]%[Hoàng Thanh Phương]%[2D4V3-2]
\immini{Một vật chuyển động trong $4$ giờ với vận tốc $v(t)$ phụ thuộc vào thời gian $t$ có đồ thị vận tốc như hình vẽ bên. Trong khoảng thời gian $1$ giờ kể từ khi bắt đầu chuyển động, đồ thị đó là một phần của đường parabol có đỉnh $I(2;10)$ và trục đối xứng song song với trục tung. Khoảng thời gian còn lại vật chuyển động chậm dần đều. Tính quãng đường $S$ mà vật đi được trong $4$ giờ đó (kết quả làm tròn đến hàng phần mười).
}
{
\begin{tikzpicture}[scale=0.7,>=stealth, font=\footnotesize, line join=round, line cap=round,xscale=0.7,yscale=0.7]
\def\a{-1.5} \def\b{6} \def\c{4} % Hệ số
\def\xmin{-1} \def\xmax{5}
\def\ymin{-1} \def\ymax{12}
\draw[->] (\xmin,0)--(\xmax,0) node [below]{$t$};
\draw[->] (0,\ymin)--(0,\ymax) node [left]{$v$};
\node at (0,0) [below left]{$O$};
\draw[smooth,samples=300,domain=0:1] plot(\x,{\a*(\x)^2+\b*(\x)+\c});
\draw[smooth,samples=300,domain=1:3,dashed] plot(\x,{\a*(\x)^2+\b*(\x)+\c});
\draw[dashed] (1,0)--(1,8.5) (2,0)--(2,10)--(0,10) (3,0)--(3,8.5)--(1,8.5) (4,0)--(4,4)--(0,4)
;
\draw (1,8.5)--(4,4);
\foreach \d/\g in{1,2,3,4}
\draw[fill=black](\d,0)circle(1pt)node[shift={(-90:0.35)}]{$\d$};
\foreach \d/\g in{4,10}
\draw[fill=black](0,\d)circle(1pt)node[shift={(180:0.35)}]{$\d$};
\end{tikzpicture}
}
\shortans{25,3}
\loigiai{
Gọi phương trình vận tốc của vật là $v=at^2+bt+c$, với $a\neq 0$. \\
Vì đồ thị hàm số đi qua các điểm $(0;4)$ và $(2;10)$, đồng thời nhận $t=2$ là trục đối xứng nên ta có hệ phương trình
\[\heva{&c=4\\&4a+2b+c=10\\&\dfrac{-b}{2a}=2}\Leftrightarrow \heva{&a=-\dfrac{3}{2}\\&b=6\\&c=4.}\]
Vậy phương trình vận tốc của vật là $v(t)=-\dfrac{3}{2}t^2+6t+4$. \\
Ta có $v(1)=\dfrac{17}{2}$. \\
Quãng đường chuyển động của vật trong $4$ giờ đó là
\[S=\displaystyle\int\limits_{0}^{4} v(t)\mathrm{\,d}t=\displaystyle\int\limits_{0}^{1} v(t)\mathrm{\,d}t+\displaystyle\int\limits_{1}^{4} v(t)\mathrm{\,d}t=\displaystyle\int\limits_{0}^{1} \left(-\dfrac{3}{2}t^2+6t+4\right)\mathrm{\,d}t+\dfrac{\left(\dfrac{17}{2}+4\right)\cdot 3}{2}=25{,}25.\]
Vậy quãng đường vật đi được trong $4$ giờ là $25{,}25$.
}
\end{ex}

\begin{ex}%[Đề kiểm tra giữa kì 2 - THPT Thị Xã Quảng Trị, Đăng Tạ, Dự án 12EX-3-2025]%[2D4V3-2]
Ông An muốn làm một cái cổng hình Parabol như hình vẽ bên dưới
\begin{center}
\begin{tikzpicture}[scale=1, font=\footnotesize, line join=round, line cap=round, >=stealth]
% Vẽ Parabol y = -x^2 + 4
\draw[ domain=-2:2, samples=100] plot (\x, {-\x*\x + 4});

% Xác định các điểm
\coordinate (H) at (0,0);
\coordinate (G) at (0,4);
\coordinate (A) at (-2,0);
\coordinate (B) at (2,0);
\coordinate (C) at (-1.1,0);
\coordinate (D) at (1.1,0);
\coordinate (E) at (1.1, 2.79);
\coordinate (F) at (-1.1, 2.79);

% Tô đậm hình chữ nhật CDEF
\fill[gray!30] (C) -- (D) -- (E) -- (F) -- cycle;
\draw (C) -- (D) -- (E) -- (F) -- cycle;

% Vẽ đường thẳng AB và đường cao GH
\draw (A) -- (B);
\draw[dashed] (G) -- (H);

% Kí hiệu các điểm
\fill (A) circle (1pt) node[below left] {$A$};
\fill (B) circle (1pt) node[below right] {$B$};
\fill (C) circle (1pt) node[below] {$C$};
\fill (D) circle (1pt) node[below] {$D$};
\fill (E) circle (1pt) node[above right] {$E$};
\fill (F) circle (1pt) node[above left] {$F$};
\fill (G) circle (1pt) node[above] {$G$};
\fill (H) circle (1pt) node[below] {$H$};
\end{tikzpicture}
\end{center}
Chiều cao $GH = 4$ m, chiều rộng $AB = 4$ m, $AC = BD = 0{,}9$ m. Ông An làm hai cánh cổng khi đóng lại là hình chữ nhật $CDEF$ tô đậm có giá là $1\,200\,000$ đồng/m$^2$, còn các phần để trống làm xiên hoa có giá là $900\,000$ đồng/m$^2$. Hỏi để làm hai phần nói trên ông An phải trả bao nhiêu triệu đồng (làm tròn đến hàng phần mười)?
\shortans[0]{11{,}4}
\loigiai{
\begin{center}
\begin{tikzpicture}[scale=1, font=\footnotesize, line join=round, line cap=round, >=stealth]
% Vẽ hệ trục tọa độ (ẩn để giống hình vẽ đề bài)
\draw[->] (-3,0) -- (3,0) node[below] {$x$};
\draw[->] (0,-0.5) -- (0,5) node[left] {$y$};

% Vẽ Parabol y = -x^2 + 4
\draw[ domain=-2:2, samples=100] plot (\x, {-\x*\x + 4});

% Xác định các điểm
\coordinate (H) at (0,0);
\coordinate (G) at (0,4);
\coordinate (A) at (-2,0);
\coordinate (B) at (2,0);
\coordinate (C) at (-1.1,0);
\coordinate (D) at (1.1,0);
\coordinate (E) at (1.1, 2.79);
\coordinate (F) at (-1.1, 2.79);

% Tô đậm hình chữ nhật CDEF
\fill[gray!30] (C) -- (D) -- (E) -- (F) -- cycle;
\draw (C) -- (D) -- (E) -- (F) -- cycle;

% Vẽ đường thẳng AB và đường cao GH
\draw (A) -- (B);
\draw[dashed] (G) -- (H);

% Kí hiệu các điểm
\fill (A) circle (1pt) node[below left] {$A$};
\fill (B) circle (1pt) node[below right] {$B$};
\fill (C) circle (1pt) node[below] {$C$};
\fill (D) circle (1pt) node[below] {$D$};
\fill (E) circle (1pt) node[above right] {$E$};
\fill (F) circle (1pt) node[above left] {$F$};
\fill (G) circle (1pt) node[above right] {$G$};
\fill (H) circle (1pt) node[below right] {$H$};
\end{tikzpicture}
\end{center}
Chọn hệ trục tọa độ $Oxy$ sao cho $H \equiv O(0;0)$, trục $Oy$ trùng với $HG$, trục $Ox$ trùng với $AB$.
\begin{itemize}
\item Vì $AB = 4$ nên $A(-2;0)$ và $B(2;0)$.
\item Chiều cao $GH = 4$ nên đỉnh Parabol là $G(0;4)$.
\item Phương trình Parabol $(P)$ có dạng: $y = ax^2 + 4$ do $(P)$ nhận trục $Oy$ làm trục đối xứng và đi qua điểm $G(0;4)$.
\item $(P)$ đi qua $B(2;0) \Rightarrow 0 = a\cdot 2^2 + 4 \Rightarrow a = -1$.
\end{itemize}
Vậy $(P) \colon  y = -x^2 + 4$.\\
Diện tích của cổng Parabol (giới hạn bởi $(P)$ và trục hoành) là
\[S = \displaystyle\int\limits_{-2}^{2} (-x^2 + 4) \mathrm{\,d}x = \left. \left( -\dfrac{x^3}{3} + 4x \right) \right|_{-2}^{2} = \frac{32}{3} \approx 10,667 \text{ m}^2.\]
Tính diện tích hình chữ nhật $CDEF$
\begin{itemize}
\item Ta có $HD = HB - BD = 2 - 0{,}9 = 1{,}1$ m. Suy ra $D(1,1; 0)$.
\item Điểm $E$ thuộc $(P)$ có hoành độ $x_E = 1{,}1 \Rightarrow y_E = -1{,}1^2 + 4 = 2{,}79$ m.
\item Diện tích hình chữ nhật $CDEF$ là $S_{CDEF} = CD \cdot DE = (2 \cdot  1{,}1) \cdot 2{,}79 = 6{,}138$ m$^2$.
\end{itemize}
Diện tích phần làm xiên hoa (phần còn lại) là
\[S_{\text{hoa}} = S - S_{CDEF} = \dfrac{32}{3} - 6{,}138 \approx 4{,}529 \text{ m}^2.\]
Tổng chi phí ông An phải trả là
\[T = 6,138 \cdot 1\,200\,000 + \left( \dfrac{32}{3} - 6,138 \right) \cdot 900\,000 = 11\,441\,400 \text{ đồng}.\]
Làm tròn đến hàng phần mười theo đơn vị triệu đồng, ta được $11{,}4$ triệu đồng.
}
\end{ex}

\begin{ex}%[2D4V3-2]%[Dự án C đợt 3 - KSCL LeThanhTong-Võ Hoàng Nghĩa]
Một nhóm các kĩ sư muốn xây dựng một cây cầu vòm dàn thép với giá đỡ dưới bằng thép cao cấp có hình dáng là một parabol nối từ $2$ cột trụ $A$ và $B$ nằm bên dưới cây cầu. Biết hai cột trụ cách nhau $400$ m, khoảng cách từ trụ $A$ đến cây cầu là $50$ m và $AB$ song song với mặt đường.
\begin{center}
\begin{tikzpicture}[>=stealth]
\def\vongtrencau{(12,-1.7) .. controls +(135:2.5) and +(0:2.6) ..
(6,2) coordinate (A) .. controls +(180:2.6) and +(45:2.5) ..
(0,-1.7)}
\def\vongduoicau{ (0,-2) .. controls +(42:2.4) and +(180:2.5) ..
(6,1.75) .. controls +(0:2.5) and +(138:2.4) ..
(12,-2)}
\begin{scope}
\clip \vongduoicau -- (12.2,0)--(0,0)--cycle;
\foreach \x in {0.25,0.5,...,12}{
\draw (\x,-2)--+(0,4);
}
\end{scope}
\fill[yellow!65] \vongduoicau--\vongtrencau--cycle;

%Hai bờ
\path[fill=brown!65!white] (-1,0) rectangle +(14,-0.25);
\path[fill=brown!75!black] (-2,0) -- ++(2,0)--+(0,-3)--cycle
(12,0) -- ++(2,0)--+(-2,-3)--cycle;

%Dòng nước
\path[fill=cyan!35] (0,-2.1) rectangle ++(12,-0.9);
\begin{scope}[shift={(-0.5,-1.75)}]
\draw[thick,gray!45,scale=1.5] (0.79,-0.43) -- (2.4,-0.43) -- (2.38,-0.47)
(2.77,-0.43) -- (5.98,-0.43) (6.56,-0.43) -- (8.17,-0.43)
(8.2,-0.5) -- (7.32,-0.5) (7.2,-0.51) -- (5.92,-0.51)
(5.69,-0.48) -- (3.25,-0.48) (3.13,-0.48) -- (1.89,-0.48)
(1.64,-0.5) -- (0.45,-0.5) (0.44,-0.6) -- (1.67,-0.6)
(2.51,-0.62) -- (4.44,-0.62) (4.85,-0.58) -- (8.2,-0.58);
\end{scope}

\path[draw=black,thick] \vongtrencau;
\path[draw=gray,thick] \vongduoicau;

\draw[dashed] (A)--(0,2) node[left]{$30$};

\draw[->] (-2.5,0)--(14.5,0) node[below]{$x$};
\draw[->] (0,-3.5)--(0,3) node[left]{$y$};

%Ghi chú
\path (0,-2) node[anchor = east,text=yellow]{$-30$}
(135:9pt) node{$O$};
\end{tikzpicture}
\end{center}
Gắn hệ trục $Oxy$ vào cây cầu với đơn vị trục tọa độ là $10$ m. Giá đỡ dưới bằng thép là đường cong parabol tạo với hai trục tọa độ các hình phẳng có diện tích $S_1$, $S_2$ như hình vẽ bên. Biết rằng $S_2-2S_1=\dfrac{2200}{21}$. Điểm cao nhất của giá đỡ dưới bằng thép cao cấp cách mặt đường cây cầu bao nhiêu mét? (làm tròn đến hàng phần mười)

\par
\shortans[0]{$64{,}3$}
\loigiai{
\immini{
Parabol có dạng $(P)\colon y=ax^2+bx+c$.\\
Vì $(P)$ đi qua $(0;-5)$ nên $f(0)=-5 \Rightarrow c=-5$.\\
Và $(P)$ đi qua $(40;-5)$ nên $f(40)=-5$.
Suy ra $1600a+40b+c=-5 \Rightarrow 1600a+40b=0 $. \hfill (1)
}{
\begin{tikzpicture}[scale=0.8, font=\footnotesize, line join=round, line cap=round,>=stealth, xscale=0.2, yscale=0.2]
%Gán số liệu.
\def\xmin{-5};\def\ymin{-10};\def\xmax{50};\def\ymax{10};
%Gán tọa độ.
\coordinate (O) at (0,0);
%Trục Oxy.
\draw[->] (\xmin,0)--(\xmax,0) node[below]{$x$};
\draw[->] (0,\ymin)--(0,\ymax) node[left]{$y$};
\fill (O) node[below left]{$O$} circle(1pt);
%Giới hạn đồ thị.
\clip ({\xmin-0.1},{\ymin-0.1}) rectangle ({\xmax+0.1},{\ymax+0.1});
%Vẽ đồ thị.
\draw[thick,samples=100] plot[domain=-5:50](\x,{-1/35*(\x)^2+8/7*\x-5});
\draw(40,-5) node[right]{$y=f(x)$};
\draw [dashed] (40,-5)--(40,0);
\draw [dashed] (0,-5)--(40,-5) node[midway,sloped,below]{$40$};
\fill[fill=black](0,-5) node[left]{$-5$} circle(5pt);
\fill[fill=black](40,-5) node[below left]{$B$} circle(5pt);
\fill[fill=black](40,0) node[above]{$C$} circle(5pt);
\draw(0,-5) node [below right]{$A$};
\draw(1.8,-1.5) node{$S_1$};
\draw(38.2,-1.5) node{$S_1$};
\draw(20,2.5) node{$S_2$};
\draw(20,-2.5) node{$S_3$};
\end{tikzpicture}
}
\noindent
Ta có
\allowdisplaybreaks
\begin{eqnarray*}
S_2-2S_1 &=& S_2-(S_{OABC}-S_3)=(S_2+S_3)-S_{OABC}\\
&=&\displaystyle\int\limits_0^{40} \left( ax^2+bx+c+5\right) \,\mathrm{d}x-5\cdot40\\
&=& \dfrac{64000}{3}a+800b+40c.
\end{eqnarray*}
Suy ra $\dfrac{64000}{3}a+800b +40c= \dfrac{6400}{21}$. \hfill $(2)$\\
Từ $(1)$ và $(2)$ ta giải được $a=-\dfrac{1}{35}$, $b=\dfrac{8}{7}$. \\
Suy ra $(P)\colon y=-\dfrac{1}{35}x^2+\dfrac{8}{7}x-5$.\\
Do đó $(P)$ có đỉnh $I\left( 20;\dfrac{45}{7}\right) $.\\
Vậy điểm cao nhất của giá đỡ dưới bằng thép cao cấp cách mặt đường cây cầu $10\cdot \dfrac{45}{7} \approx 64{,}3$ m.

}
\end{ex}

\begin{ex}%[Bồ Văn Hậu]%[Dự án Đề ôn thi tốt nghiệp THPT]%[2D4V3-2]
\immini{Người ta thiết kế một mẫu gạch lát nền nhà có dạng hình vuông cạnh $4$ dm. Bốn góc viên gạch màu trắng, phần ở giữa màu xanh (Hình 5). Đường viền của phần màu xanh bao gồm bốn đoạn thẳng nằm trên các cạnh hình vuông và bốn đường cong có tính chất: Tích khoảng cách từ một điểm bất kì thuộc đường cong đó đến hai trục đối xứng của viên gạch (hai đường thẳng đi qua tâm viên gạch và lần lượt song song với hai cạnh vuông góc) bằng $2$ dm$^2$.
Hãy cho biết phần màu xanh có diện tích bằng bao nhiêu decimét vuông (\textit{làm tròn kết quả đến hàng phần mười})?
\par\shortans[0]{13{,}5}}{\begin{tikzpicture}[line join=round, line cap=round,>=stealth,thick,font=\scriptsize]
\tikzset{every node/.style={scale=0.9}}
\fill[gray] (0,0)--(0,2)--plot[domain=1:2] (\x,{(0*(\x)+2)/(1*(\x)+0)})--(2,0)--(2,1)--plot[domain=2:1] (\x,{(0*(\x)+-2)/(1*(\x)+0)})--(1,-2)--(0,-2)--cycle;
\fill[xscale=-1,gray] (0,0)--(0,2)--plot[domain=1:2] (\x,{(0*(\x)+2)/(1*(\x)+0)})--(2,0)--(2,1)--plot[domain=2:1] (\x,{(0*(\x)+-2)/(1*(\x)+0)})--(1,-2)--(0,-2)--cycle;
\draw (-2,-2)rectangle (2,2);

\end{tikzpicture}}

\loigiai{
Gắn trục toạ độ $Oxy$ vào viên gạch sao cho hai trục trùng với hai đường đối xứng, gốc $O$ ở tâm hình vuông như hình dưới. Giả sử toạ độ một điểm trên đường viền cong là $( x;y )$. \\
Theo giả thiết, ta có $| xy |=2$. Suy ra $ y=\dfrac {2}{x}$ hoặc $ y=-\dfrac {2}{x}$.\\ Ứng với hình dưới, ta có các đường viền cong $AK$, $DE$ là một phần của đồ thị hàm số $ y=-\dfrac {2}{x}$; các đường viền cong $BC,GH$ là một phần của đồ thị hàm số $ y=\dfrac {2}{x}$.\\
\begin{center}
\begin{tikzpicture}[line join=round, line cap=round,>=stealth,thick,font=\scriptsize]
\tikzset{every node/.style={scale=0.9}}
\path (-1,2) coordinate (A)
(1,2) coordinate (B)
(2,1) coordinate (C)
(2,-1) coordinate (D)
(1,-2) coordinate (E)
(-1,-2) coordinate (G)
(-2,-1) coordinate (H)
(-2,1) coordinate (K)
;
\fill[gray] (0,0)--(0,2)--plot[domain=1:2] (\x,{(0*(\x)+2)/(1*(\x)+0)})--(2,0)--(2,1)--plot[domain=2:1] (\x,{(0*(\x)+-2)/(1*(\x)+0)})--(1,-2)--(0,-2)--cycle;
\fill[xscale=-1,gray] (0,0)--(0,2)--plot[domain=1:2] (\x,{(0*(\x)+2)/(1*(\x)+0)})--(2,0)--(2,1)--plot[domain=2:1] (\x,{(0*(\x)+-2)/(1*(\x)+0)})--(1,-2)--(0,-2)--cycle;
\draw[->] (-3.1,0)--(3.1,0) node[below left] {$x$};
\draw[->] (0,-3.1)--(0,3.1) node[below left] {$y$};
\draw (0,0) node [above left] {$O$};
\foreach \x/\nx in {-2/-2,-1/-1,1/1,2/2}
\draw[thin] (\x,1pt)--(\x,-1pt) node [shift={(-60:0.3)}] {$\nx$};
\foreach \y/\ny in {-2/-2,-1/-1,2/2}
\draw[thin] (1pt,\y)--(-1pt,\y) node [shift={(145:0.3)}] {$\ny$};
\draw[dashed,thin] (A)--(G)
(B)--(E) (C)--(K) (D)--(H)
;
\begin{scope}
\clip (-2,-2) rectangle (2,2);
\draw[samples=200,domain=-2:-0.01,smooth,variable=\x] plot (\x,{(0*(\x)+2)/(1*(\x)+0)});
\draw[samples=200,domain=0.01:2,smooth,variable=\x] plot (\x,{(0*(\x)+2)/(1*(\x)+0)});
\draw[dashed,thin] (-2,0/1)--(2,0/1);
\draw[samples=200,domain=-2:-0.01,smooth,variable=\x] plot (\x,{(0*(\x)+-2)/(1*(\x)+0)});
\draw[samples=200,domain=0.01:2,smooth,variable=\x] plot (\x,{(0*(\x)+-2)/(1*(\x)+0)});
\draw[dashed,thin] (-2,0/1)--(2,0/1);
\end{scope}
\draw (-2,-2)rectangle (2,2);
\foreach \t/\g in {A/90,B/90,C/0,D/0,E/-90,G/-90,H/180,K/180}{
\path (\t) node[shift={(\g:7pt)},font=\scriptsize]{$ \t $};
}
\end{tikzpicture}
\end{center}
Khi đó, diện tích phần màu xanh bằng
\allowdisplaybreaks
\begin{eqnarray*}
& &\displaystyle\int\limits_{-2}^{-1} \left| -\dfrac {2}{x}-\dfrac {2}{x} \right|\,\mathrm{d}x +\displaystyle\int\limits_1^2 \left| \dfrac {2}{x}-\dfrac {-2}{x} \right|\,\mathrm{d}x +S_{ABEG}\\
&=&\displaystyle\int\limits_{-2}^{-1} \left( -\dfrac {2}{x}-\dfrac {2}{x} \right)\,\mathrm{d}x +\displaystyle\int\limits_1^2 \left( \dfrac {2}{x}-\dfrac {-2}{x} \right)\,\mathrm{d}x +4{,}2\\
&=&\left.-4\ln |x|\right|_{-2}^{-1} + \left.4\ln |x|\right|_1^2+8\approx 13{,}5 \,(\text{dm}^2).
\end{eqnarray*}
}
\end{ex}

\begin{ex}%[2D4V3-2]%[Dự án C - đợt 3 - NH24-25-PNL-Deso13-Lê Minh Thien Anh]
Một phần của bề mặt phía trên của các gợn sóng của nước biển có hình dạng đồ thị hàm số bậc ba khi gắn hệ trục tọa độ $Oxy$. Biết rằng đồ thị hàm số bậc ba có các điểm cực trị lần lượt là $M(-2;1)$ và $N(0;1{,}2)$. Đơn vị trên hệ trục tọa độ là mét. Hai vị trí $A$ và $B$ có hoành độ lần lượt là $-3$ và $1$ nằm trên đường cong của các gợn sóng đó. Tính độ dài đường cong $AB$ (Kết quả làm tròn đến hàng phần trăm). Biết rằng độ dài đường cong có phương trình $y=f(x)$ từ điểm $C(c;f(c))$ tới điểm $D(d;f(d))$ với $c<d$ được tính bởi công thức $T = \displaystyle\int\limits_{c}^{d} \sqrt{1+(f'(x))^2} \mathrm{\,d}x$.

\par \shortans[0]{4,07}
\loigiai{
Gọi hàm số bậc ba là $y = f(x) = ax^3 + bx^2 + cx + d$ có $f'(x) = 3ax^2 + 2bx + c$.\\
Do $M(-2;1)$ và $N(0;1{,}2)$ là hai điểm cực trị nên $\heva{& f'(-2)=0 \\& f'(0)=0 }$ và $\heva{& f(-2)=1 \\ &f(0)=1{,}2.}$\\
Thay vào tìm được $a = -0{,}05$; $b = -0{,}15$; $c = 0$; $d = 1{,}2$ suy ra $f'(x) = -0{,}15x^2 - 0{,}3x$.\\
Vậy độ dài đường cong $AB$ là $l_{AB} = \displaystyle\int\limits_{-3}^{1} \sqrt{1+(-0{,}15x^2 - 0{,}3x)^2} \mathrm{\,d}x \approx 4{,}07$ mét.
}
\end{ex}

\begin{ex}%[2D4V3-2]
Trong lễ kỷ niệm 30 năm thành lập trường THPT Chu Văn An – Biên Hòa (1994 - 2024), nhà trường có cho dựng 1 cổng chào hình parabol với điểm cao nhất là $6$ m và chiều rộng giữa hai chân cổng chào là $8$ m. Hỏi diện tích của cổng chào là bao nhiêu mét vuông?
\par\shortans[0]{32}
\loigiai{
\immini{Đặt hệ quy chiếu $Oxy$, vào 1 điểm chân cổng chào.\\
Khi đó, ta cần xác định hàm số $y=ax^2+bx+c$, biết parabol có đỉnh $I(4;6)$ và đi qua gốc tọa độ $O(0;0)$.\\
Thay tọa độ điểm $O$ vào hàm số, suy ra $c=0$.\\
Ta có hệ phương trình $\heva{&16a+4b=6 \\ &8a+b=0} \Leftrightarrow \heva{&a=\dfrac{-3}{8} \\ &b=3}$\\
Suy ra hàm số $y = \dfrac{-3}{8}x^2+3x$\\
Diện tích của cổng chào: $S = \displaystyle \int_0^8\left| \dfrac{-3}{8}x^2+3x \, dx \right|  = 32$ (m$^2$).}
{\begin{tikzpicture}[line join=round, line cap=round,>=stealth,thick,scale=0.6]
\tikzset{every node/.style={scale=0.9}}
\draw[->] (-1.1,0)--(9.1,0) node[below left] {$x$};
\draw[->] (0,-1.1)--(0,7.1) node[below left] {$y$};
\draw (0,0) node [below left] {$O$};
\foreach \x/\nx in {4/4,8/8}
\draw[thin] (\x,1pt)--(\x,-1pt) node [below] {$\nx$};
\foreach \y/\ny in {6/6}
\draw[thin] (1pt,\y)--(-1pt,\y) node [left] {$\ny$};
\begin{scope}
\clip (-1,-1) rectangle (9,7);
\draw[samples=200,domain=0:8,smooth,variable=\x] plot (\x,{-0.375*(\x)^2+3*(\x)+0});
\end{scope}
\draw[dashed] (4,0)--(4,6)--(0,6);
\foreach \p in {(4,6),(8,0),(0,0)}
\fill \p circle (1.75pt);
\end{tikzpicture}}
}
\end{ex}

\begin{ex}%[2D4V3-2]
\immini{Bác Bình muốn làm cửa rào sắt hình parabol có chiều cao từ mặt đất đến đỉnh là $3$ (m), chiều rộng tiếp giáp với mặt đất là $2$ (m) (được mô phỏng bởi hình vẽ bên dưới). Giá $1$ (m$^2$) vật tư và công làm là $1{,}4$ (triệu đồng), các chi phí khác không đáng kể. Vậy để làm cửa rào sắt đó thì số tiền Bác Bình phải trả là bao nhiêu triệu đồng (\textit{làm tròn kết quả đến hàng phần chục})?}{\begin{tikzpicture}[line join=round, line cap=round,>=stealth,thick,font=\scriptsize]
\tikzset{every node/.style={scale=1}}
\path (-1,0) coordinate (A)
(1,0) coordinate (B)
(0,3) coordinate (C)
;
\begin{scope}
\clip (-1.5,-1) rectangle (1.5,3.5);
\fill[blue,thin,pattern=north east lines](-1,0) --plot[samples=200,domain=-1:1,smooth,variable=\x] (\x,{-3*(\x)^2+3})--(1,0);
\draw plot[samples=200,domain=-1:1,smooth,variable=\x] (\x,{-3*(\x)^2+3})
(-1,0)--(1,0)
;
\end{scope}

\path (A)--(B)--([turn]-90:0.35) coordinate (xt)
($(A)-(B)+(xt)$) coordinate (yt);
\draw[>=stealth,|<->|] (xt)--(yt) node[fill=white,inner sep=0pt,midway,sloped]{$2\ \mathrm{m}$};
\path ($(C)+(0:1.2)$) coordinate (xt)
($(B)+(0:0.2)$) coordinate (yt);
\draw[dash pattern=on 2pt off 2pt] (C)--(xt) (B)--(yt);
\draw[>=stealth,|<->|] (xt)--(yt) node[fill=white,inner sep=0pt,midway,sloped]{$3\ \mathrm{m}$};
\end{tikzpicture}}
\par\shortans[0]{$5{,}6$}
\loigiai{\immini{Do tính đối xứng của parabol nên chọn hệ trục tọa độ $Oxy$ sao cho parabol có đỉnh $I\in Oy$.\\
Gọi phương trình parabol $(P)\colon y=ax^2+bx+c$ $(a\ne 0)$.\\
Ta có $\heva{& I(0;3)\in (P) \\ & A(-1;0)\in (P) \\ & B(1;0)\in (P) }\Rightarrow \heva{& 3=c \\ & a-b+c=0 \\ & a+b+c=0} \Leftrightarrow \heva{& a=-3 \\ & b=0 \\ & c=3.}$\\
Do đó $(P)\colon y=-3x^2+3$.\\
Diện tích cửa rào hình parabol là
\allowdisplaybreaks
\begin{eqnarray*}
&S=&\displaystyle\int\limits_{-1}^1 (-3x^2+3 )\mathrm{\,d}x\\
&=&2\displaystyle\int\limits_0^1 (-3x^2+3)\mathrm{\,d}x\\
&=&2\cdot\left. \left(-x^3+3x\right)\right|_0^1 \\
&=&4\, \left(\text{m}^2\right).
\end{eqnarray*}
Vậy số tiền Bác Bình phải trả là $4\cdot 1{,}4=5{,}6$ (triệu đồng).}{\begin{tikzpicture}[line join=round, line cap=round,>=stealth,thick,font=\scriptsize]
\tikzset{every node/.style={scale=0.9}}
\draw[->] (-2,0)--(2.5,0) node[below left] {$x$};
\draw[->] (0,-1.1)--(0,4) node[below left] {$y$};
\draw (0,0) node [above left] {$O$};
\begin{scope}
\clip (-1.5,-1) rectangle (1.5,3.5);
\fill[blue,thin,pattern=north east lines](-1,0) --plot[samples=200,domain=-1:1,smooth,variable=\x] (\x,{-3*(\x)^2+3})--(1,0);
\draw plot[samples=200,domain=-1:1,smooth,variable=\x] (\x,{-3*(\x)^2+3});
\end{scope}
\draw (1,0)node[below]{$B(1;0)$}
(-1,0)node[below]{$A(-1;0)$}
(0,3)node[above right]{$I(0;3)$}
;
\end{tikzpicture}}
}
\end{ex}

\begin{ex}%[2D4V3-2]%[Dự Án C - Đợt 2 - Lê Thánh Tông - TP HCM]%[Nguyễn Văn Sang]
Một nhóm các kỹ sư muốn xây dựng một cây cầu vòm dàn thép với giá đỡ dưới bằng thép cao cấp có hình dáng là một đường cong Parabol nối từ $2$ cột trụ $A$ và $B$ nằm bên dưới cây cầu, biết hai cột trụ cách nhau $400$ m, khoảng cách từ trụ A đến cây cầu là $50$ m và $AB$ song song với mặt đường.
\begin{center}
\begin{tikzpicture}[>=stealth]
\def\vongtrencau{(12,-1.7) .. controls +(135:2.5) and +(0:2.6) ..
(6,2) coordinate (A) .. controls +(180:2.6) and +(45:2.5) ..
(0,-1.7)}
\def\vongduoicau{ (0,-2) .. controls +(42:2.4) and +(180:2.5) ..
(6,1.75) .. controls +(0:2.5) and +(138:2.4) ..
(12,-2)}
\begin{scope}
\clip \vongduoicau -- (12.2,0)--(0,0)--cycle;
\foreach \x in {0.25,0.5,...,12}{
\draw (\x,-2)--+(0,4);
}
\end{scope}
\fill[yellow!65] \vongduoicau--\vongtrencau--cycle;

%Hai bờ
\path[fill=brown!65!white] (-1,0) rectangle +(14,-0.25);
\path[fill=brown!75!black] (-2,0) -- ++(2,0)--+(0,-3)--cycle
(12,0) -- ++(2,0)--+(-2,-3)--cycle;

%Dòng nước
\path[fill=cyan!35] (0,-2.1) rectangle ++(12,-0.9);
\begin{scope}[shift={(-0.5,-1.75)}]
\draw[thick,gray!45,scale=1.5] (0.79,-0.43) -- (2.4,-0.43) -- (2.38,-0.47)
(2.77,-0.43) -- (5.98,-0.43) (6.56,-0.43) -- (8.17,-0.43)
(8.2,-0.5) -- (7.32,-0.5) (7.2,-0.51) -- (5.92,-0.51)
(5.69,-0.48) -- (3.25,-0.48) (3.13,-0.48) -- (1.89,-0.48)
(1.64,-0.5) -- (0.45,-0.5) (0.44,-0.6) -- (1.67,-0.6)
(2.51,-0.62) -- (4.44,-0.62) (4.85,-0.58) -- (8.2,-0.58);
\end{scope}

\path[draw=black,thick] \vongtrencau;
\path[draw=gray,thick] \vongduoicau;

\draw[dashed] (A)--(0,2) node[left]{$30$};

\draw[->] (-2.5,0)--(14.5,0) node[below]{$x$};
\draw[->] (0,-3.5)--(0,3) node[left]{$y$};

%Ghi chú
\path (0,-2) node[anchor = east,text=yellow]{$-30$}
(135:9pt) node{$O$};
\end{tikzpicture}
\end{center}
Gắn hệ trục tọa độ Oxy vào cây cầu với đơn vị trục tọa độ là $10$ m. Giá đỡ dưới bằng thép là đường cong Parabol tạo với $2$ trục tọa độ các hình phẳng có diện tích $S_1$, $S_2$ như hình vẽ bên, biết rằng
\[
S_2 - 2S_1 = \dfrac{2\,200}{21}.
\]
Điểm cao nhất của giá đỡ dưới bằng thép cao cấp cách mặt đường cây cầu bao nhiêu mét (\textit{làm tròn đến hàng phần mười}).
\begin{center}
\begin{tikzpicture}[scale=1, font=\footnotesize, line join=round, line cap=round, >=stealth]
% Vẽ hệ trục tọa độ
\draw[->] (-1,0) -- (9,0) node[right] {$x$};
\draw[->] (0,-3) -- (0,2.5) node[above] {$y$};
\fill (0,0) circle (1pt) node[below left]{$O$};
\fill (0,-1) circle (1pt) node[left]{$-5$};
\fill (0,-1) circle (1pt) node[below right]{$A$};
\fill (6.54,-1) circle (1pt) node[below]{$B$};
\node at (3.25,.7) {$S_2$};
\node at (0.24,-.35) {$S_1$};
%vẽ nối AB node 40m ở giữa
\draw[dashed,<->] (0,-1) -- (6.54,-1) node[midway,above]{$40$};
\draw[dashed] (6.54,-1) -- (6.54,0);
% Vẽ đồ thị parabol
\draw[domain=-1:7.5,smooth,thick,blue] plot (\x,{(-0.19)*(\x)^2 + 1.25*(\x) - 1});
\end{tikzpicture}
\end{center}
\par
\shortans{64{,}3}
\loigiai{
\begin{center}
\begin{tikzpicture}[scale=1, font=\footnotesize, line join=round, line cap=round, >=stealth]
% Vẽ hệ trục tọa độ
\draw[->] (-1,0) -- (9,0) node[right] {$x$};
\draw[->] (0,-3) -- (0,2.5) node[above] {$y$};
\fill (0,0) circle (1pt) node[below left]{$O$};
\fill (0,-1) circle (1pt) node[left]{$-5$};
\fill (0,-1) circle (1pt) node[below right]{$A$};
\fill (6.54,-1) circle (1pt) node[below]{$B$};
\fill (6.54,0) circle (1pt) node[above]{$C$};
\node at (3.25,.7) {$S_2$};
\node at (3.25,-.35) {$S_3$};
\node at (0.24,-.35) {$S_1$};
\node at (6.25,-.35) {$S_1$};
%vẽ nối AB node 40m ở giữa
\draw[dashed,<->] (0,-1) -- (6.54,-1) node[midway,above]{$40$};
\draw[dashed] (6.54,-1) -- (6.54,0);
% Vẽ đồ thị parabol
\draw[domain=-1:7.5,smooth,thick,blue] plot (\x,{(-0.19)*(\x)^2 + 1.25*(\x) - 1});
\end{tikzpicture}
\end{center}
Parabol là $y = f(x) = ax^2 + bx + c \quad (a \neq 0)$.\\
Parabol đi qua $(0; -5)$, $(40; -5)$ nên
\begin{itemize}
\item  $f(0) = -5 \Rightarrow c = -5$. \hfill $(1)$
\item  $f(40) = -5 \Rightarrow 1\,600a + 40b + c = -5$. \hfill $(2)$
\end{itemize}
Ta có
\begin{eqnarray*}
S_2 - 2S_1 &=& S_2 - (S_{OABC} - S_3)\\
&=& S_2 - S_{OABC} + S_3\\
&=& (S_2 + S_3) - S_{OABC}\\
&=&  \displaystyle\int\limits_{0}^{40} (ax^2 + bx + c + 5) \,\mathrm{d}x -5\cdot40\\
&=& \left( a \cdot \dfrac{x^3}{3} + b \cdot \dfrac{x^2}{2} + (c+5) x \right) \Bigg|_0^{40} - 200\\
&=& \left( a \cdot \dfrac{40^3}{3} + b \cdot \dfrac{40^2}{2} + (c+5) \cdot 40 \right) - 200\\
&=& \dfrac{40^3}{3} a + \dfrac{40^2}{2} b + 40c - 200.
\end{eqnarray*}
Từ giả thiết suy ra
$\dfrac{40^3}{3} a + \dfrac{40^2}{2} b + 40c = \dfrac{2\,200}{21}$. \hfill $(3)$\\
Từ $(1)$, $(2)$ và $(3)$ suy ra
$a = \dfrac{-1}{35}$, $b = \dfrac{8}{7}$, $c= -5
$.\\
Vậy parabol $y = \dfrac{-1}{35} x^2 + \dfrac{8}{7} x - 5$.\\
Gọi $I$ là đỉnh của parabol.\\
Hoành độ của đỉnh parabol là $x_I = \dfrac{-b}{2a} = \dfrac{-\dfrac{8}{7}}{2 \cdot \left(\dfrac{-1}{35}\right)} = 20$.\\
Tung độ của đỉnh parabol là $y_I=-\dfrac{1}{35} \cdot 20^2+\dfrac{8}{7} \cdot 20-5=\dfrac{45}{7}$.\\
Điểm cao nhất của giá đỡ dưới bằng thép cao cấp cách mặt đường cây cầu là
\[10 \cdot \dfrac{45}{7} \approx 64{,}3\ (\text{m}).\]
}
\end{ex}

\begin{ex}%[2D4V3-3]
\immini{
Từ hình chữ nhật $ABCD$ có chiều dài $AB=10$ cm và chiều rộng $BC=5$ cm. Người ta cắt bỏ miền $(R)$ được giới hạn bởi cạnh $CD$ của hình chữ nhật và hai nửa đường parabol có chung đỉnh là trung điểm của cạnh $AB$, chúng lần lượt đi qua hai đầu mút $C$, $D$ của hình chữ nhật đó (phần tô đậm như hình vẽ).
}
{
\begin{tikzpicture}[scale=.6, font=\footnotesize, line join=round, line cap=round, >=stealth]
\tikzset{label style/.style={font=\footnotesize}}
%Nhập giới hạn đồ thị và hàm số cần vẽ
%\pgfmathsetmacro{\so}{sqrt{5}}
\def \hamso{sin((2*(\x))*180/pi)}
%\def \tiemcanxien{\x+1}
%Tự động
%	\draw[->] (-5.5,0)--(6,0) node[below left] {$x$};
%	\draw[->] (0,-0.5)--(0,6) node[below left] {$y$};
%	\draw[fill=black] (0,0) circle(1pt) node [below right] {$O$};
%Vẽ các điểm trên 2 hệ trục
\draw
(-5,0) coordinate (A) circle(1pt) node[above left]{$A$} %node[below]{$-5$}
(5,0) coordinate (B) circle(1pt) node[above right]{$B$} %node[below]{$5$}
(-5,5) coordinate (D) circle(1pt) node[above left]{$D$}
(5,5) coordinate (C) circle(1pt) node[above right]{$C$}
%(0,5) circle(1pt) node[above right]{$5$}
;
\draw
(A)--(B)--(C)--(D)--cycle;
%Tự động
\begin{scope}
\draw[samples=350,domain=0:5,smooth,variable=\x] plot (\x,{sqrt(5*\x)});
\draw[samples=350,domain=-5:0,smooth,variable=\x] plot (\x,{sqrt{-5*(\x)}});
\end{scope}
\end{tikzpicture}
}
\noindent Phần còn lại cho quay quanh trục $AB$ để tạo nên một đồ vật làm trang trí, thể tích của vật trang trí đó bằng bao nhiêu? (làm tròn đến hàng đơn vị).
\shortans[]{393}
\loigiai{
\begin{center}
\begin{tikzpicture}[scale=.6, font=\footnotesize, line join=round, line cap=round, >=stealth]
\tikzset{label style/.style={font=\footnotesize}}
%Nhập giới hạn đồ thị và hàm số cần vẽ
%\pgfmathsetmacro{\so}{sqrt{5}}
\def \hamso{sin((2*(\x))*180/pi)}
%\def \tiemcanxien{\x+1}
%Tự động
\draw[->] (-5.5,0)--(6,0) node[below left] {$x$};
\draw[->] (0,-0.5)--(0,6) node[below left] {$y$};
\draw[fill=black] (0,0) circle(1pt) node [below right] {$O$};
%Vẽ các điểm trên 2 hệ trục
\draw
(-5,0) coordinate (A) circle(1pt) node[above left]{$A$} node[below]{$-5$}
(5,0) coordinate (B) circle(1pt) node[above right]{$B$} node[below]{$5$}
(-5,5) coordinate (D) circle(1pt) node[above left]{$D$}
(5,5) coordinate (C) circle(1pt) node[above right]{$C$}
(0,5) circle(1pt) node[above right]{$5$}
;
\draw
(A)--(B)--(C)--(D)--cycle;
%Tự động
\begin{scope}
\draw[line width=1.5,color=red,samples=350,domain=0:5,smooth,variable=\x] plot (\x,{sqrt(5*\x)});
\draw[samples=350,domain=-5:0,smooth,variable=\x] plot (\x,{sqrt{-5*(\x)}});
\end{scope}
\end{tikzpicture}
\end{center}
Khi quay hình chữ nhật $ABCD$ quanh trục $AB$ ta được khối trụ có bán kính đáy $r=BC=5$ cm, chiều cao $h=AB=10$ cm.\\
Do đó thể tích khối trụ này có thể tích $V_1=\pi \cdot 5^2\cdot 10=250\pi$ (cm$^3$).\\
Mặt khác, chọn hệ trục tọa độ như hình vẽ thì $C(5;5)$ và parabol bên phải trục $Ox$ có dạng $(P_1)\colon y^2=2px$.\\
Ta có $C\in (P_1)\Leftrightarrow p=\dfrac{5}{2}
\Rightarrow (P_1)\colon y^2=5x$ hay $y=\sqrt{5x}$.\\
Khi đó miền $(R)$ khi quay quanh trục $Ox$ có thể tích $V_2=2\pi \displaystyle \int\limits_0^5 \left[ {5^2}-5x \right]\,\mathrm{\,d}x=2\pi \left. \left( 25x-\dfrac{5}{2}{x^2} \right) \right|_0^5=125\pi $.\\
Vậy thể tích phần còn lại là $V=V_1-V_2=125\pi$ (cm$^3$).\\
}
\end{ex}

\begin{ex}%[2D4V3-3] %[Tổ 21 - Đợt 17 - Chương 4 - - Cánh Diều - Đề 2]
\immini{Cho hình $(H)$ giới hạn bởi đồ thị hàm số $y=\dfrac{\sqrt{3}}{9} x^{3}$, cung tròn có phương trình $y=\sqrt{4-x^{2}}$
(với $0 \leq x \leq 2$ ) và trục hoành (phần tô đậm trong hình vẽ).
Biết thể tích của khối tròn xoay tạo thành khi quay $(H)$ quanh trục hoành là $V=\left(-\dfrac{a}{b} \sqrt{3}+\dfrac{c}{d}\right) \pi$, trong đó $a$, $b$, $c$, $d \in\mathbb{N}^{*}$ và $\dfrac{a}{b}$, $\dfrac{c}{d}$ là các phân số tối giản. Tính $P=a+b+c+d$.}{
\begin{tikzpicture}[scale=0.7, font=\footnotesize, line join=round, line cap=round, >=stealth]
\def\xt{-3} \def\xp{3} \def\yt{3} \def\yd{-2}
\draw[->](\xt, 0)--(\xp, 0) node[below]{$x$};
\draw[->](0,\yd)--(0,\yt) node[left]{$y$};
\fill (0,0) node[below left]{$O$} circle(1.5pt);
\clip(\xt+0.1,\yd+0.1) rectangle(\xp+1,\yt-0.1);
\draw[color=blue, smooth, samples=300, domain=0:2] plot(\x,{sqrt(4-(\x)^2)});
\draw[color=red, smooth, samples=300] plot(\x,{(sqrt(3))/9*(\x)^3});
% Tô màu giữa hai đồ thị
\fill[fill=orange, opacity=0.5]
(0,0) -- plot[samples=300, domain=0:1.732] (\x,{(sqrt(3))/9*(\x)^3}) -- (1.732,0) -- cycle;
\fill[fill=orange, opacity=0.5]
(1.732,0) -- plot[samples=300, domain=1.732:2] (\x,{sqrt(4-(\x)^2)}) -- (2,0) -- cycle;
\fill (2,0) node[below]{$2$} circle(1.5pt);
\fill (0,2) node[left]{$2$} circle(1.5pt);
\end{tikzpicture}	}
\shortans{$46$}
\loigiai
{
Phương trình hoành độ giao điểm $\dfrac{\sqrt{3}}{9} x^{3}=\sqrt{4-x^{2}} \Leftrightarrow x=\sqrt{3}$.
\begin{align*}
& V=\pi\left[\displaystyle\int_{0}^{\sqrt{3}}\left(\dfrac{\sqrt{3}}{9} x^{3}\right)^{2} \mathrm{d} x+\displaystyle\int_{\sqrt{3}}^{2}\left(\sqrt{4-x^{2}}\right)^{2} \mathrm{d} x\right] \\
& =\pi\left[\displaystyle\int_{0}^{\sqrt{3}} \dfrac{1}{27} x^{6} \mathrm{d} x+\displaystyle\int_{\sqrt{3}}^{2}\left(4-x^{2}\right) \mathrm{d} x\right] \\
& =\pi\left[\left.\dfrac{1}{27} \cdot \dfrac{x^{7}}{7}\right|_{0} ^{\sqrt{3}}+\left.\left(4 x-\dfrac{x^{3}}{3}\right)\right|_{\sqrt{3}} ^{2}\right] \\
& =\left(-\dfrac{20}{7} \sqrt{3}+\dfrac{16}{3}\right) \pi.
\end{align*}
Suy ra $a=20$, $b=7$, $c=16$, $d=3$. \\
Nên $P=a+b+c+d=46$.
}
\end{ex}

\begin{ex}%[2D4V3-3] %[Tổ 21 - Đợt 17 - Chương 4 - - Cánh Diều - Đề 2]%[Lưu Công Chinh]
Cho hình phẳng $(H)$ giới hạn bởi đồ thị hai hàm số $y=\sqrt{3 x+1}$ và $y=|x-1|$. Biết thể tích của khối tròn xoay tạo thành khi quay $(H)$ quanh trục hoành là $V=\dfrac{a \pi}{b}$, trong đó $a$, $b \in \mathbb{N}^{*}$ và $\dfrac{a}{b}$ là các phân số tối giản. Tính $P=a+b$.
\shortans{$131$}
\loigiai
{Phương trình hoành độ giao điểm \[\sqrt{3 x+1}=|x-1| \Leftrightarrow 3 x+1=(x-1)^{2} \Leftrightarrow x^{2}-5 x=0 \Leftrightarrow
\hoac{&x=0 \\& x=5.}
\]
Thể tích khối tròn xoay là
\begin{align*}
V&=\pi \displaystyle\int_{0}^{5}\left|(\sqrt{3 x+1})^{2}-\left| x-1\right|^{2}\right|\mathrm{d}x\\
&=\pi \displaystyle\int_{0}^{5}\left| x^{2}-5 x \right| \mathrm{d}x\\
&=\dfrac{125 \pi}{6}.
\end{align*}
Suy ra $a=125$, $b=6$, nên $P=125+6=131$.
}
\end{ex}

\begin{ex}%[12-MH-1-MH2025]%[MH-2025, Mã/Tên TV biên soạn]%[2D4V3-3]
\immini[thm]{Một chiếc bát thuỷ tinh có bề dày của phần xung quanh là một khối tròn xoay, khi xoay hình phẳng $D$ quanh một đường thẳng $a$ bất kì nào đó mà khi gắn hệ trục tọa độ $Oxy$ (đơn vị trên trục là decimét) vào hình phẳng $D$ tại một vị trí thích hợp, thì đường thẳng $a$ sẽ trùng với trục $Ox$. Khi đó, hình phẳng $D$ được giới hạn bởi các đồ thị hàm số $y=x+\dfrac{1}{x}$, $y=x$ và hai đường thẳng $x=1$, $x=4$. Thể tích của bề dày chiếc bát thuỷ tinh đó bằng bao nhiêu decimét khối (làm tròn kết quả đến hàng phần mười)?}
{\begin{tikzpicture}[font=\footnotesize,line join=round, line cap=round, >=stealth,scale=0.6]
\def \xmin{-1}\def \xmax{5.5}\def \ymin{-5.2}\def \ymax{5.5}
\draw[->] (\xmin,0)--(\xmax,0) node[shift=(-110:0.2)] {$x$};
\draw[->] (0,\ymin)--(0,\ymax) node[shift=(-150:0.2)] {$y$};
\fill (0,0) circle(1pt) node[shift=(135:0.25)]{$O$}
(1,0) circle(1pt) node[shift=(-65:0.23)]{$1$}
(2,0) circle(1pt) node[shift=(-90:0.2)]{$2$}
(3,0) circle(1pt) node[shift=(-90:0.2)]{$3$}
(4,0) circle(1pt) node[shift=(-65:0.23)]{$4$}
(0,1) circle(1pt) node[shift=(180:0.2)]{$1$}
(0,2) circle(1pt) node[shift=(180:0.2)]{$2$};
\clip (\xmin,\ymin) rectangle (\xmax,\ymax);
\fill[gray!30!] plot[domain=1:4](\x,{\x+1/\x})--plot[domain=4:1](\x,{\x});
\draw[smooth,samples=100,domain=0.01:\xmax] plot(\x,{\x+1/(\x)});
\draw[smooth,samples=100,domain=\xmin:\xmax] plot(\x,{(\x)});
\draw[dashed,smooth,samples=100,domain=1:4] plot(\x,{-\x});
\draw[dashed,smooth,samples=100,domain=1:4] plot(\x,{-\x-(1/(\x))});
\draw (1,0)--(1,2) (4,0)--(4,4.25);
\draw[dashed] (1,0)--(1,-2) (4,0)--(4,-4.25);
\draw[dashed] (1,1)arc(90:270:0.25 and 1)(1,2)arc(90:270:0.5 and 2);
\draw[dashed] (1,1)arc(90:-270:0.25 and 1)(1,2)arc(90:-270:0.5 and 2);
\draw[dashed] (4,4)arc(90:270:0.45 and 4)(4,4.25)arc(90:270:0.6 and 4.25);
\draw[dashed] (4,4)arc(90:-270:0.45 and 4)(4,4.25)arc(90:-270:0.6 and 4.25);
\end{tikzpicture}}
\shortans[]{$21{,}2$}
\loigiai{
Gọi $V_1$ là thể tích của khối tròn xoay được tạo thành khi cho hình phẳng giới hạn bởi đồ thị hàm số $y=x+\dfrac{1}{x}$, trục hoành và hai đường thẳng $x=1$, $x=4$ quay quanh trục $Ox$.\\
Khi đó \[V_1=\pi\displaystyle \int \limits_1^4\left(x+\dfrac{1}{x}\right)^2\mathrm{\,d}x=\dfrac{111\pi}{4}\, (\text{dm}^3).\]
Gọi $V_2$ là thể tích của khối tròn xoay được tạo thành khi cho hình phẳng giới hạn bởi đồ thị hàm số $y=x$, trục hoành và hai đường thẳng $x=1$, $x=4$ quay quanh trục $Ox$.\\
Khi đó \[V_2=\pi\displaystyle \int \limits_1^4 x^2\mathrm{\,d}x=21\pi \, (\text{dm}^3).\]
Vậy thể tích của bề dày chiếc bát thuỷ tinh là
\begin{align*}
V=V_1-V_2=\dfrac{111\pi}{4}-21\pi=\dfrac{27\pi}{4}\approx 21{,}2\, (\text{dm}^3).
\end{align*}
}
\end{ex}

\begin{ex}%[2D4V3-3]
Cho hình $(H)$ giới hạn bởi đồ thị hàm số $y=\dfrac{\sqrt{3}}{9}x^3$, cung tròn có phương trình $y=\sqrt{4-x^2}$ (với $0\le x\le 2)$ và trục hoành (phần tô đậm trong hình vẽ).
\begin{center}
\begin{tikzpicture}[x=1cm,y=1cm]
\draw[->] (-3,0)--(3,0)
node[below]{$x$};
\draw[->] (0,-3) --(0,3)
node[left]{$y$};
\begin{scope}
\clip (0,0) circle (2);
\fill[blue,smooth] (-2,0) --
plot[domain=0:2] (\x,{0.1924500897*(\x)^3})
-- (2,0) -- cycle;
\end{scope}
\draw plot[domain=-2.3:2.3,smooth]
(\x,{0.1924500897*(\x)^3});
\begin{scope}
\clip (2,3.5) rectangle (0,0);
\draw (0,0) circle(2);
\end{scope}
\path
(0,0) node[above left]{$O$}
(2,0) node[below right]{$2$}
(-2,0) node[below left]{$-2$}
(0,2) node[above left]{$2$}
(0,-2) node[below left]{$-2$};
\end{tikzpicture}
\end{center}
Biết thể tích của khối tròn xoay tạo thành khi quay $(H)$ quanh trục hoành là $V=\left(-\dfrac{a}{b}\sqrt{3}+\dfrac{c}{d}\right)\pi $, trong đó $a, b, c, d\in {\mathbb{N}}^{*}$ và $\dfrac{a}{b}$, $\dfrac{c}{d}$ là các phân số tối giản. Tính $P=a+b+c+d$.
\shortans{$46$}
\loigiai{
Phương trình hoành độ giao điểm là \[\dfrac{\sqrt{3}}{9}x^3=\sqrt{4-x^2}\Leftrightarrow x=\sqrt{3}.\]
Khi đó thể tích của khối tròn xoay được tạo thành là
\begin{eqnarray*}
&V&=\pi \left[\displaystyle \int\limits_0^{\sqrt{3}}\left(\dfrac{\sqrt{3}}{9}{x^3}\right)^2\mathrm{\,d}x+\displaystyle \int\limits_{\sqrt{3}}^{2}\left(\sqrt{4-x^2}\right)^2\mathrm{\,d}x\right]\\
& &=\pi \left[\displaystyle \int\limits_0^{\sqrt{3}}{\dfrac{1}{27}x^6\mathrm{\,d}x}+\displaystyle \int\limits_{\sqrt{3}}^2{\left(4-x^2\right)\mathrm{\,d}x}\right]\\
& &=\pi \left[\left. \dfrac{1}{27}\cdot\dfrac{x^7}{7}\right|_0^{\sqrt{3}}+\left. \left(4x-\dfrac{x^3}{3}\right)\right|_{\sqrt{3}}^2\right]\\
& &=\left(-\dfrac{20\sqrt{3}}{7}+\dfrac{16}{3}\right)\pi\\
& &\Rightarrow a=20, b=7, c=16, d=3\\
& &\Rightarrow P=a+b+c+d=46.
\end{eqnarray*}
}
\end{ex}

\begin{ex}%[2D4V3-3]
\immini[thm]{Cho hình phẳng $(H)$ là tam giác cong $OAB$ trong hình vẽ bên. Tính thể hình tròn xoay sinh ra bởi $(H)$ khi quay $(H)$ quanh trục $Ox$ (Kết quả viết dưới dạng số thập phân và làm tròn đến hàng phần trăm).
}{
\begin{tikzpicture}[line join=round, line cap=round,>=stealth,thick]
\tikzset{every node/.style={scale=0.9}}
\draw[dashed, step=1, gray!50,very thin] (-.5,-0.9) grid (4.5,4.5);
\draw[->] (-1.6,0)--(5.1,0) node[below] {$x$};
\draw[->] (0,-1.8)--(0,4.5) node[right] {$y$};
\draw (0,0) node [below left] {$O$};
\foreach \x/\nx in {-1/-1,1/1,2/2,3/3,4/4}
\draw[thin] (\x,1pt)--(\x,-1pt) node [below] {$\nx$};
\foreach \y/\ny in {-1/-1,1/1,2/2,3/3,4/4}
\draw[thin] (1pt,\y)--(-1pt,\y) node [left] {$\ny$};
\begin{scope}
\clip (-1.5,-1.5) rectangle (4.5,4.5);
\draw plot[samples=200,domain=-1.2:1.7,smooth,variable=\x] (\x,{(1*(\x)^3});
\path (1,1)--(2,8) node[pos=0.4,above, sloped]{$y=x^3$};
\draw plot[samples=200,domain=-0.3:4.3,smooth,variable=\x] (\x,{1*(\x)^2+-4*(\x)+4});
\path (3,1)--(4,4) node[pos=0.55,above, sloped]{$y=x^2-4x+4$};
\fill[black](1,1) circle (1.5pt) node[right]{$A$};
\fill[black](2,0) circle (1.5pt) node[above]{$B$};
\end{scope}
\end{tikzpicture}
}
\shortans{$1{,}08$}
\loigiai{
Ta có $V=\pi\displaystyle\int_0^1{\left(x^3\right)^2\mathrm{\,d}x}+\pi\displaystyle\int_1^2{\left(x^2-4x+4\right)^2\mathrm{\,d}x} \approx 1{,}08$.
}
\end{ex}

\begin{ex}%[2D4V3-3]
\immini[thm]{Gọi $V$ là thể tích khối tròn xoay tạo thành khi quay hình phẳng giới hạn bởi các đường $y=\sqrt{x}$, $y=0$ và $x=4$ quanh trục $Ox$. Đường thẳng $x=a$, $\left(0<a<4\right)$ cắt đồ thị hàm số $y=\sqrt{x}$ tại $M$ (hình vẽ). Gọi $V_1$ là thể tích khối tròn xoay tạo thành khi quay tam giác $OMH$ quanh trục $Ox$. Biết rằng $V=2V_1$. Tìm $a$.
}{
\begin{tikzpicture}[line join=round, line cap=round,>=stealth,thick]
\tikzset{every node/.style={scale=0.9}}
\draw[->] (-0.6,0)--(5.1,0) node[below left] {$x$};
\draw[->] (0,-0.4)--(0,2.5) node[below left] {$y$};
\draw (0,0) node [below left] {$O$};
\foreach \x/\nx in {3/a,4/4}
\draw[thin] (\x,1pt)--(\x,-1pt) node [below] {$\nx$};
\begin{scope}
\clip (-1.0,-1.0) rectangle (4.5,2.5);
\draw plot[samples=200,domain=0:4.5,smooth,variable=\x] (\x,{sqrt((\x))});
\path (0,0)--(4,1) node[pos=0.45,above=0.8cm, sloped]{$y=\sqrt x$};
\fill[black](3,1.732) circle (1.5pt) node[above right]{$M$} (4,0) node[above right]{$H$};
\draw (3,2)--(3,-0.1);
\draw[pattern=north east lines](0,0)--(3,1.732)--(4,0);
\end{scope}
\end{tikzpicture}
}
\shortans{$3$}
\loigiai{
\immini{
Ta có $V=\pi\displaystyle\int\limits_0^4x\mathrm{\,d}x=\pi\dfrac{x^2}{2}\Bigg|_0^4=8\pi$.\\
Mà $V=2V_1\Rightarrow{V_1}=4\pi$.\\
Gọi $K$ là hình chiếu của $M$ trên $Ox$.\\
Suy ra $OK=a$, $KH=4-a$, $MK=\sqrt a$.\\
Khi xoay tam giác $OMH$ quanh $Ox$ ta được khối
}{
\begin{tikzpicture}[line join=round, line cap=round,>=stealth,thick]
\tikzset{every node/.style={scale=0.9}}
\draw[->] (-0.6,0)--(5.1,0) node[below left] {$x$};
\draw[->] (0,-0.4)--(0,2.5) node[below left] {$y$};
\draw (0,0) node [below left] {$O$};
\foreach \x/\nx in {3/a,4/4}
\draw[thin] (\x,1pt)--(\x,-1pt) node [below] {$\nx$};
\begin{scope}
\clip (-1.0,-1.0) rectangle (4.5,2.5);
\draw plot[samples=200,domain=0:4.5,smooth,variable=\x] (\x,{sqrt((\x))});
\path (0,0)--(4,1) node[pos=0.45,above=0.8cm, sloped]{$y=\sqrt x$};
\fill[black](3,1.732) circle (1.5pt) node[above right]{$M$} (4,0) node[above right]{$H$} (3,0)node[below right]{$K$};
\draw (3,2)--(3,-0.1);
\draw[pattern=north east lines](0,0)--(3,1.732)--(4,0);
\end{scope}
\end{tikzpicture}
}\hspace{-0.77cm}
tròn xoay là sự lắp ghép của hai khối nón sinh bởi các tam giác $OMK$, $MHK$, hai khối nón đó có cùng mặt đáy và có tổng chiều cao là $OH=4$ nên thể tích của khối tròn xoay đó là $V_1=\dfrac{1}{3} \cdot \pi \cdot 4 \cdot \left(\sqrt a\right)^2=\dfrac{4\pi a}{3}$, từ đó suy ra $a=3$.
}
\end{ex}

\begin{ex}%[Dự án EX-Ôn Tập TN 2025, Duong Xuan Loi]%[2D4V3-3]
\immini{Một chiếc bát thuỷ tinh có bề dày của phần xung quanh là một khối tròn xoay, khi xoay hình phẳng $D$ quanh một đường thẳng $a$ bất kì nào đó mà khi gắn hệ trục tọa độ $Oxy$ (đơn vị trên trục là decimét) vào hình phẳng $D$ tại một vị trí thích hợp, thì đường thẳng $a$ sẽ trùng với trục $Ox$. Khi đó, hình phẳng $D$ được giới hạn bởi các đồ thị hàm số $y=x+\dfrac{1}{x}$, $y=x$ và hai đường thẳng $x=1$, $x=4$. Thể tích của bề dày chiếc bát thuỷ tinh đó bằng bao nhiêu decimét khối? (làm tròn kết quả đến hàng phần mười).}
{\begin{tikzpicture}[x=1cm,y=1cm,scale=0.6,>=stealth]
\def\hsf(#1){(#1)+1/(#1)}
\def\hsg(#1){(#1)}
\fill[gray!30!](1,0)--plot[domain=1:4](\x,{\hsf(\x)})--plot[domain=4:1](\x,{\hsg(\x)});
\draw[->] (-0.88,0)--(5.5,0)node[below right]{$x$};
\draw[->] (0,-5.2)--(0,5.5)node[left]{$y$};
\fill (0,0)node[below left]{$O$};
\foreach \a in {2,3} \fill (\a,0) circle (1.2pt)node[below]{$\a$};
\foreach \b in {1,2} \fill (0,\b) circle (1.2pt)node[left]{$\b$};
\draw (1,0)circle(1.2pt)node[above right]{\small{1}};
\draw (4,0)circle(1.2pt)node[above right]{\small{4}};
\draw[black,samples=150,smooth,domain=0.2:5.5] plot(\x,{(\x)+(1/(\x))});
\draw[black,samples=150,smooth,domain=-0.3:5.5] plot(\x,{(\x)});
\draw[dashed,black,samples=150,smooth,domain=-0.3:5.5] plot(\x,{-1*(\x)});
\draw[dashed,black,samples=150,smooth,domain=1:4] plot(\x,{-1*(\x)-(1/(\x))});
\draw (1,0)--(1,2) (4,0)--(4,4.25);
\draw[dashed] (1,0)--(1,-2) (4,0)--(4,-4.25);
\draw[dashed] (1,1)arc(90:270:0.25 and 1)(1,2)arc(90:270:0.5 and 2);
\draw[dashed] (1,1)arc(90:-270:0.25 and 1)(1,2)arc(90:-270:0.5 and 2);
\draw[dashed] (4,4)arc(90:270:0.45 and 4)(4,4.25)arc(90:270:0.6 and 4.25);
\draw[dashed] (4,4)arc(90:-270:0.45 and 4)(4,4.25)arc(90:-270:0.6 and 4.25);
\end{tikzpicture}}
\shortans{$21{,}2$}
\loigiai{
Gọi $V_1$ là thể tích của khối tròn xoay được tạo thành khi cho hình phẳng giới hạn bởi đồ thị hàm số $y=x+\dfrac{1}{x}$, trục hoành và hai đường thẳng $x=1$, $x=4$ quay quanh trục $Ox$.\\
Khi đó \[V_1=\pi\displaystyle \int \limits_1^4\left(x+\dfrac{1}{x}\right)^2\mathrm{\,d}x=\dfrac{111\pi}{4}\, (\text{dm}^3).\]
Gọi $V_2$ là thể tích của khối tròn xoay được tạo thành khi cho hình phẳng giới hạn bởi đồ thị hàm số $y=x$, trục hoành và hai đường thẳng $x=1$, $x=4$ quay quanh trục $Ox$.\\
Khi đó \[V_2=\pi\displaystyle \int \limits_1^4 x^2\mathrm{\,d}x=21\pi \, (\text{dm}^3).\]
Vậy thể tích của bề dày chiếc bát thuỷ tinh đó là \[V=V_1-V_2=\dfrac{111\pi}{4}-21\pi=\dfrac{27\pi}{4}\approx 21{,}2\, (\text{dm}^3).\]
}
\end{ex}

\begin{ex} %[2D4V3-3]
Cho tam giác vuông $OAB$ có cạnh $OA$ nằm trên trục $Ox$ và $\widehat{AOB}=\alpha \, (0<\alpha <\frac{\pi }{2})$ và $B(a;b)$ với $a$, $b$ là các số thực thỏa ${a^2}+{b^2}=1$. Gọi $\beta $ là khối tròn xoay sinh ra khi quay miền tam giác $OAB$ xung quanh trục $Ox$.
\begin{center}
\begin{tikzpicture}[scale=1,line join=round, line cap=round,>=stealth,declare function={rt=.03;xmin=-2.5;xmax=8;ymin=-2.5;ymax=3;a=2.5;b=a/3;}]
\path (0,0) coordinate (O)
(6,0)coordinate (A)++(90:a) coordinate (B)
;
\draw[->] (0,ymin)--(0,ymax) node[right] {$y$};
\draw[->] (0,0)--(-150:2) node[above left] {$z$};
\draw[fill=black] (0,0) circle (rt);
\fill[red!15] (O)--(B)--(A)--cycle;
\draw[->]  (A)--(xmax,0) node[below] {$x$};
\draw (A) ellipse ({b} and {a})
(-1,0)--(0,0) node[above left]{$O$} (O)--(B)--(A) (O)--($(A)+(-90:a)$)
pic["$\alpha$",angle radius=15mm] {angle = A--O--B}
pic[draw,angle radius=7mm] {angle = A--O--B}
;
\draw[dashed](O)--(A);
\foreach \p/\g in {A/-90,B/90}
\fill (\p) circle(1pt)+(\g:.3) node{$\p$};
\end{tikzpicture}
\end{center}
Tính giá trị $\tan \alpha $ khi thể tích của khối $\beta $ đạt giá trị lớn nhất (Làm tròn kết quả đến chữ số thập phân thứ 2).
\shortans[1]{$1{,}41$}
\loigiai{
Do $OB$ đi qua gốc tọa độ nên ta đặt $OB\colon y=kx$ với $ k$ là số thực dương.\\
Do $OB$ đi qua $B(a;b) \Rightarrow OB\colon y=\frac{b}{a}x$ và $\tan \alpha =\frac{b}{a}$.\\
Khi đó, $V=\pi \int\limits_0^a{{\left(  \dfrac{b}{a}x\right) ^2}\mathrm{\,d}x}=
\pi \int\limits_0^a{\dfrac{{b^2}}{{a^2}}{x^2}\mathrm{\,d}x}=
\dfrac{\pi {b^2}{x^3}}{3{a^2}} \bigg|_0^a=
\dfrac{\pi }{3}{b^2}a$.\\
Áp dụng bất đẳng thức Am – Gm:\\
\[{V^2}=\dfrac{{{\pi }^2}}{9}{b^4}{a^2}=\dfrac{{{\pi }^2}}{18}{b^2}{b^2}2{a^2}\le \dfrac{{{\pi }^2}}{18}\dfrac{{{( {b^2}+{b^2}+2{a^2} )}^3}}{27}=\dfrac{4{{\pi }^2}}{243}\\
\Rightarrow V\le \frac{2\pi \sqrt {3}}{27}.\]
Đẳng thức xảy ra khi và chỉ khi ${b^2}=2{a^2} \Rightarrow \dfrac{b}{a}=\tan \alpha =\sqrt {2} \approx 1{,}41$.}
\end{ex}

\begin{ex}%[2D4V3-3]
Một quả trứng khủng long đồ chơi bằng nhựa có thiết diện qua trục lớn là một đường elip. Biết độ dài mỗi trục là $12$ cm và $8$ cm.
Bên trong quả trứng người ta cần thiết kế một chiếc hộp hình trụ để đựng các đồ chơi trẻ con như bóng đèn xanh đỏ, kẹo $\ldots$
Hỏi khối trụ như thế có thể tích tối đa bao nhiêu cm$^3$ (làm tròn đến hàng đơn vị)?
\begin{center}
\begin{tikzpicture}[line join=round, line cap=round,>=stealth,font=\footnotesize,scale=0.9]
\path (0,0) coordinate (O);
\draw (O) ellipse (3 cm and 2 cm);
\path ($(O) + (40:3 and 2)$) coordinate (A)
($(O) + (140:3 and 2)$) coordinate (B)
($(O) + (220:3 and 2)$) coordinate (C)
($(O) + (-40:3 and 2)$) coordinate (D)			;
\draw[dashed] (A)--(B) (C)--(D)
(A) to [bend right=20](D)
(B) to [bend right=20](C)
;
\draw (A) to [bend left=20](D)
(B) to [bend left=20](C)
;
\foreach \t/\g in {A/90,B/90,C/180,D/0}{
\draw[fill=black] (\t) circle (1pt) ;
}
\end{tikzpicture}
\end{center}
\shortans{232}
\loigiai{
\begin{center}
\begin{tikzpicture}[line join=round, line cap=round,>=stealth,font=\footnotesize,scale=0.9]
\path (0,0) coordinate (O);
\draw (O) ellipse (3 cm and 2 cm);
\path ($(O) + (40:3 and 2)$) coordinate (A)
($(O) + (140:3 and 2)$) coordinate (B)
($(O) + (220:3 and 2)$) coordinate (C)
($(O) + (-40:3 and 2)$) coordinate (D)	;
\draw[->] (-4,0)--(4,0)node[below]{$x$};
\draw[->] (0,-3)--(0,3)node[left]{$y$};
\draw[dashed] (A)--(B) (C)--(D)
(A) to [bend right=20](D)
(B) to [bend right=20](C)
;
\draw (A) to [bend left=20](D)
(B) to [bend left=20](C)
;
\foreach \t/\g in {A/90,B/90,C/180,D/0}{
\draw[fill=black] (\t) circle (1pt);
}
\foreach \t/\g in {O/230,A/90}{
\draw[fill=black] (\t) circle (1pt) node[shift={(\g:7pt)},font=\scriptsize]{$ \t $};
}
\end{tikzpicture}
\end{center}
Gắn elip lên hệ trục tọa độ $Oxy$ như hình vẽ, elip có \[2a=12\Rightarrow a=6; 2b=8\Rightarrow b=4.\]
Phương trình chính tắc elip $(E)\colon\dfrac{x^2}{36}+\dfrac{y^2}{16}=1$.\\
Đặt chiều cao và bán kính đáy hình trụ nội tiếp elip là $h$, $r$ thì điểm tiếp xúc $A\left(\dfrac{h}{2}; r\right)$ với $0< h < 12$; $0< r < 4$.\\
Điểm $A$ thuộc $(E)\colon\dfrac{x^2}{36}+\dfrac{y^2}{16}=1\Rightarrow \dfrac{\left(\dfrac{h}{2} \right)^2}{36}+\dfrac{r^2}{16}=1
\Rightarrow \dfrac{h^2}{144}+\dfrac{r^2}{16}=1\Rightarrow r^2=16\left(1-\dfrac{h^2}{144}\right)$.\\
Thể tích khối trụ là $V=\pi r^2h=\pi \cdot16\left(1-\dfrac{h^2}{144} \right)\cdot h=\pi \left(16h-\dfrac{h^3}{9} \right)$.\\
Ta có $V'=\pi \left(16-\dfrac{h^2}{3} \right)$;
$V'=0\Rightarrow 16-\dfrac{h^2}{3}=0\Rightarrow h^2=48\Rightarrow h=4\sqrt{3} \in (0; 12)$.\\
Bảng biến thiên
\begin{center}
\begin{tikzpicture}[font=\normalsize,t style/.style={style=solid}]
\tkzTabInit[nocadre=false, lgt=1.2, espcl=2.5, deltacl=0.5]%,help]
{$h$/0.7, $V'$/0.7, $V$/2}
{$0$, $4\sqrt{3}$, $12$}
\tkzTabLine{ , +, 0, -}
\path
($(N13)!0.1!(N12)$) node (A1){}
($(N23)!0.8!(N22)$) node (A2){}
($(N33)!0.1!(N32)$) node (A3){};
\foreach \x/\y in {A1/A2, A2/A3}{
\draw[-stealth] (\x)--(\y);
}
\end{tikzpicture}
\end{center}
Giá trị lớn nhất của thể tích khối trụ là \[V_{\max}=V\left(4\sqrt{3} \right)=\pi \left(16\cdot 4\sqrt{3}-\dfrac{\left(4\sqrt{3} \right)^3}{9} \right)\approx 232\ (\mathrm{cm}^3).\]
}
\end{ex}

\begin{ex}%[Hiếu Mai]%[2D4V3-3]
\immini[thm]
{
Cho khối tròn xoay như hình vẽ. Tính thể thích khối tròn xoay được tạo thành bởi hình phẳng cho ở hình bên khi quay quanh trục $Ox$ (viết kết quả dưới dạng số thập phân và làm tròn kết quả đến hàng phần mười).
}
{
\begin{tikzpicture}[scale=0.95, font=\footnotesize, line join=round, line cap=round, >=stealth]
\draw[->,teal, line width = 1pt, color=black] (-1,0) -- (5,0);
\draw(5,0) node [anchor=south west] {$x$};
\draw[->,teal, line width = 1pt, color=black](0,-2) -- (0,3);
\draw(0,3) node [anchor=west] {$y$};
\draw[smooth,samples=100,domain=0:4] plot(\x,{sqrt((\x))}) node[above]{$y=\sqrt{x}$};
\draw[smooth,dashed,samples=100,domain=0:4] plot(\x,{-sqrt((\x))});
\draw(4,-2) arc (-90:90:0.5 cm and 2 cm);
\draw[dashed] (4,2) arc (90:270:0.5 cm and 2 cm);
\draw[dotted](4,-2)--(4,2);
\draw (0,0) circle (1pt) node[below left] {$O$} (4,0) circle (1pt) node[below left] {$1$};
\end{tikzpicture}
}
\shortans[oly]{$1{,}6$}
\loigiai
{
Hình phẳng đã cho được giới hạn bởi đồ thị hàm số $y=\sqrt{x}$, trục hoành và các đường thẳng $x=0 $, $x=1$. \\
Khi quay hình phẳng đó quanh trục $Ox$ ta được khối tròn xoay có thể tích là
\[V=\pi\cdot\displaystyle\int\limits_{0}^{1} \left(\sqrt{x}\right)^2\mathrm{\,d}x=\pi\cdot\displaystyle\int\limits_{0}^{1} x\mathrm{\,d}x=\dfrac{\pi}{2}\approx1{,}6.\]
}
\end{ex}

\begin{ex}%[Nguyễn Thái Hoàng]%[Đợt 1 - HSA - Đề 22]%[2D4V3-3]
Cho hai hình cầu như hình vẽ. Biết rằng bán kính của hai hình cầu đều bằng $5$, khoản cách giữa tâm của hai hình cầu bằng $6$. Thể tích phần chung giữa hai hình cầu là $V_c=\dfrac{a}{b}\pi$, với $\dfrac{a}{b}$ là phân số tối giản $a$, $b>0$. Tính giá trị biểu thức $T=ab$.
\begin{center}
\begin{tikzpicture}[scale=.4, font=\footnotesize, line join=round, line cap=round, >=stealth]
\fill[red!40] (-3,0) circle (5);
\fill[red!40] (3,0) circle (5);
\draw (-3,0) circle[radius=5cm];
\draw (3,0) circle[radius=5cm];
\end{tikzpicture}
\end{center}
\shortans[]{312}
\loigiai{
\begin{center}
\begin{tikzpicture}[scale=.4, font=\footnotesize, line join=round, line cap=round, >=stealth]
\def \xmin{-10}\def \xmax{10}\def \ymin{-6}\def \ymax{6}
\draw[->] (\xmin,0)--(\xmax,0) node[shift=(-110:0.2)] {$x$};
\draw[->] (0,\ymin)--(0,\ymax) node[shift=(-30:0.2)] {$y$};
\draw (-3,0) circle[radius=5cm];
\draw (3,0) circle[radius=5cm];
\fill (0,0) circle(1pt) node[shift=(135:0.25)]{$O$}
(-3,0) circle(1pt) node[shift=(90:0.2)]{$O_1$}
(3,0) circle(1pt) node[shift=(45:0.2)]{$O_2$}
(2,0) circle(1pt) node[shift=(-45:0.3)]{$2$}
;
%	\fill[green] (0,0)--(0,4)--(2,0)--cycle;
\end{tikzpicture}
\end{center}
Phương trình phần đường tròn $O_1(-3,0)$ bán kính $R_1=5$ phía trên trục $Ox$ có phương trình là
\[(x+3)^2+y^2=25\Rightarrow y=\sqrt{25-(x+3)^2}.\]
Phương trình đường tròn $O_2(3,0)$ bán kính $R_2=5$ có phương trình là
\[(x-3)^2 + y^2 = 25.\]
Thể tích phần giao của hai hình cầu là
\allowdisplaybreaks
\begin{eqnarray*}
V_c&=&2\pi\displaystyle\int\limits_0^2\left(\sqrt{25-(x+3)^2}\right)^2\mathrm{\,d}x\\&=&2\pi \displaystyle\int\limits_0^2\left(25-(x+3)^2\right)\mathrm{\,d}x\\&=&2\pi \left[25x - \dfrac{(x+3)^3}{3}\right]\Bigg|_0^2\\&=&\dfrac{104}{3}\pi.
\end{eqnarray*}
Suy ra $a = 104$, $b=3$.\\
Khi đó $T=ab=312$.
}
\end{ex}

\begin{ex}%[HSA-L1-2526]%[2D4V3-3]
Gọi $S$ là phần hình phẳng màu xanh được tạo bởi hai đồ thị hàm số $y=\sin x+1$ và $y=\cos x+1$ như hình vẽ dưới. Biết thể tích của vật thể được tạo thành khi quay $S$ quanh trục Ox là $V=a \sqrt{b} \pi+c \pi^2$, với $a$, $b$, $c \in \mathbb{N}$, $b$ là số nguyên tố. Tính giá trị của biểu thức $P=a+b+c$. (nhập đáp án vào ô trống)
\begin{center}
\begin{tikzpicture}[>=stealth,x=1cm,y=1cm,scale=1]
\draw[->] (-3,0) -- (6,0) node[below left] {$x$};
\draw[->] (0,-1) -- (0,3) node[below left] {$y$};
\fill[gray!60] plot[red,samples=100,domain=pi/4:1.25*pi] (\x,{sin((\x)*180/pi)+1})--plot[red,samples=100,domain=1.25*pi:pi/4] (\x,{cos((\x)*180/pi)+1});
\draw[samples=100,domain=-3:6] plot(\x,{sin((\x)*180/pi)+1});
\draw[samples=100,domain=-3:6] plot(\x,{cos((\x)*180/pi)+1});
\node at (0,0) [below left]{$O$};
\foreach \x in{-1.57,1.57,3.17,4.71}\draw[fill=black] (\x,0) circle (.5pt);
\foreach \y in{2}\draw[fill=black] (0,\y) circle (.5pt);
\node at (-1.57,0) [below]{$-\dfrac{\pi}{2}$};
\node at (1.57,0) [below]{$\dfrac{\pi}{2}$};
\node at (3.14,0) [below]{$\pi$};
\node at (4.71,0) [below]{$\dfrac{3\pi}{2}$};
\draw (0,2) node[above right] {$2$};
\end{tikzpicture}
\end{center}
\par\shortans{6}
\loigiai{
Xét phương trình hoành độ giao điểm giữa hai đồ thị hàm số
\[\sin x+1=\cos x+1\Leftrightarrow \sin x-\cos x=0\Leftrightarrow \sqrt{2} \sin \left(x-\dfrac{\pi}{4}\right)=0\Leftrightarrow x=\dfrac{\pi}{4}+k \pi \quad(k \in \mathbb{Z}).\]
Khi đó, phần hình phẳng màu xanh được giới hạn bởi $x=\dfrac{\pi}{4}$ và $x=\dfrac{5\pi}{4}$.\\
Khi $x \in\left[\dfrac{\pi}{4}; \dfrac{5\pi}{4}\right]$, ta có $\sin x+1\geq \cos x+1\geq 0$. Thể tích của vật thể tròn xoay khi đó bằng
\begin{align*}
V&=\pi \displaystyle\int\limits_{\tfrac{\pi}{4}}^{\tfrac{5\pi}{4}}\left|(\sin x+1)^2-(\cos x+1)^2\right|\mathrm{\,d}x
\\
&=\pi \displaystyle\int\limits_{\tfrac{\pi}{4}}^{\tfrac{5\pi}{4}}\left(\sin ^2x-\cos ^2x+2\sin x-2\cos x\right) \mathrm{\,d}x
\\
&=\pi \cdot \displaystyle\int\limits_{\tfrac{\pi}{4}}^{\tfrac{5\pi}{4}}(-\cos 2x+2\sin x-2\cos x) \mathrm{\,d}x=\pi \cdot\left(\dfrac{-\sin 2x}{2}-2\cos x-2\sin x\right)\bigg|_{\tfrac{\pi}{4}} ^{\tfrac{5\pi}{4}}=4\pi \sqrt{2}.
\end{align*}
Khi đó, $a=4$, $b=2$, $c=0\Rightarrow P=a+b+c=6$.
}
\end{ex}

\begin{ex}%[2D4V3-4]
\immini{Cho vật thể đáy là hình tròn có bán kính bằng $1$ (tham khảo hình vẽ). Khi cắt vật thể bằng mặt phẳng vuông góc với trục $Ox$ tại điểm có hoành độ $x\ (-1\leq x\leq 1)$ thì được thiết diện là một tam giác đều. Tính thể tích vật thể đó (làm tròn đến chữ số đầu tiên sau dấu phẩy).
}{
\begin{tikzpicture}[line cap=round,line join=round,scale=.7,font=\footnotesize,>=stealth]
\def\h{4}
\def\a{3}
\def\x{.65}
\pgfmathsetmacro\b{\a/3}
\pgfmathsetmacro{\g}{asin(\b/\h)}
\pgfmathsetmacro\xo{\a *cos(\g)}
\pgfmathsetmacro\yo{\b *sin(\g)}
\fill[draw,thick,top color=gray!15,bottom color=gray,smooth] (90:\h)coordinate(S)--(-\xo,\yo)coordinate(A) arc (180-\g:360+\g:{\a} and {\b})--(\x,{-(\x)^2+\h}) plot[domain=\x:0](\x,{-(\x)^2+\h});
\begin{scope}
\clip (90:\h)--(-\xo,\yo) arc (180-\g:360+\g:{\a} and {\b})--(\x,{-(\x)^2+\h}) plot[domain=\x:0](\x,{-(\x)^2+\h});
\draw[smooth,thick] plot[domain=0:\a](\x,{-(\x)^2+\h});
\foreach \x in {.1,.2,...,\a}\draw[opacity=.5] (\x,{-(\x)^2+\h})--++(-128:2*\h);
\foreach \y in {-10,...,30}\draw[smooth,shift={($($(S)!\y/25!(A)$)-(S)$)},opacity=.5] plot[domain=0:\a](\x,{-(\x)^2+\h});
\foreach \x in {.65,.75,...,\a}\draw[opacity=.5] (\x,{-(\x)^2+\h})--++(-55:\h);
\end{scope}
\end{tikzpicture}}
\shortans{$2{,}3$}
\loigiai{
\immini{
Ta có $\triangle OAB$ đều có cạnh $AB=2AH=2\sqrt{1-x^2}.$\\
Diện tích thiết diện là $S(x)=\dfrac{AB^2\cdot\sqrt{3}}{4}=\sqrt{3}\left(1-x^2\right)$.\\
Khi đó thể tích vật thể là
\allowdisplaybreaks
\begin{eqnarray*}
V&=& \displaystyle\int\limits_{-1}^1 S(x) \mathrm{\,d}x=\sqrt{3}\displaystyle\int\limits_{-1}^1 (1-x^2) \mathrm{\,d}x\\
&=&\sqrt{3}\left(x-\dfrac{x^3}{3}\right)\bigg|_{-1}^1=\dfrac{4\sqrt{3}}{3}\approx 2{,}3.
\end{eqnarray*}
}{\begin{tikzpicture}[scale=1.2, font=\footnotesize,line join=round, line cap=round, >=stealth];
\draw[->] (-2,0) -- (2.25,0) node[below] {\scriptsize $x$};
\draw[->] (0,-2) -- (0,2.25) node[left] {\scriptsize $y$};
\draw (-0.3,0) node[below] {$O$};
\draw  (60:1.6) node[right]{$A$};
\draw  (-60:1.6) node[right]{$B$};
\draw (0:1.5) arc (0:360:1.5);
\draw  (0:0)--(48:1.5) (0:0)--(-48:1.5) (-48:1.5)--(48:1.5);
\draw[fill=black]  (48:1.5) circle (1.2pt) (-48:1.5) circle (1.2pt) (1,0) circle (1.2pt);
\draw (1.2,0) node[below] {$x$};
\draw (1.2,0) node[above] {$H$};
\draw (0.5,1) node[below] {$1$};
\end{tikzpicture}}

}
\end{ex}

\begin{ex}%[2D4V3-4]%[Dự Án EX-TF-TLN Toán 12 - Đợt 2 - Quan Ón]
Một vật thể nằm giữa hai mặt phẳng $x = 0$ và $x = \dfrac{\pi}{2}$; biết rằng mặt cắt của vật thể cắt bởi mặt phẳng vuông góc với trục $Ox$ tại điểm có hoành độ $x$ $\left(0 \leq x \leq \dfrac{\pi}{2}\right)$ là tam giác đều có cạnh là $2\sqrt{\cos x + \sin x}$. Thể tích của vật thể trên có dạng $a\sqrt{b}$. Hãy tính $5a + b^2$.
\shortans{$19$}
\loigiai{
Diện tích của mặt cắt là $S(x) = \dfrac{\left(2\sqrt{\cos x + \sin x} \right)^2\cdot \sqrt{3}}{4} = \sqrt{3}\left(\cos x + \sin x\right)$.
Vậy thể tích của vật thể đã cho là
\[ V = \int_{0}^{\frac{\pi}{2}}S(x)\,\mathrm{d}x = \int_{0}^{\frac{\pi}{2}}\sqrt{3}\left(\cos x + \sin x\right)\,\mathrm{d}x = 2\sqrt{3}.\]
Do đó $a = 2$ và $b = 3$.\\
Vậy $5a + b^2 = 5\cdot 2 + 3^2 = 19$.
}
\end{ex}

\begin{ex}%[Vovanle]%[2D4V3-4]
Một téc nước hình trụ, đang chứa nước được đặt nằm ngang, có chiều dài $3$ m và đường kính đáy $1$ m. Hiện tại mặt nước trong téc cách phía trên đỉnh của téc $0{,}25$ m (xem hình vẽ).
\begin{center}
\begin{tikzpicture}[line join=round,line cap=round, font=\footnotesize,scale=1,>=stealth]
\def \x{0.5}
\def \y{1.5}
\def \z{6}
\path
(80:{\x} and {\y}) coordinate (A)
(80:{\x} and {\y})+(\z,0) coordinate (B)
(-80:{\x} and {\y}) coordinate (C)
(-80:{\x} and {\y})+(\z,0) coordinate (D)
(10:{\x} and {\y}) coordinate (E)
(200:{\x} and {\y}) coordinate (F)
($(B)!0.5!(D)$) coordinate (K)
($(A)!0.5!(C)$) coordinate (H)
(40:{\x} and {\y}) coordinate (M)
(160:{\x} and {\y}) coordinate (N)
(40:{\x} and {\y})+(\z,0) coordinate (U)
(160:{\x} and {\y})+(\z,0) coordinate (V)
(C)+(0,-0.7) coordinate (I)
(D)+(0,-0.7) coordinate (J)
(B)+(2,0) coordinate (P)
(D)+(2,0) coordinate (Q)
(intersection of B--D and U--V) coordinate (T)
(B)+(0.6,0) coordinate (G)
(T)+(0.6,0) coordinate (R)
($(P)!0.5!(Q)$) coordinate (m)node[right]{$1$ m}
($(I)!0.5!(J)$) coordinate (X)node[above]{$3$ m}
($(G)!0.5!(R)$) coordinate (Z)node[right]{$0{,}25$ m}
;
\fill[blue!20] plot [domain=-90:-200] (0.5*cos \x,{1.5*sin \x})--(M)--(U)--plot [domain=40:-90] (\z+0.5*cos \x,{1.5*sin \x})--cycle;
\draw  (A) arc (80:280:{\x} and {\y});
\draw[dashed] (C) arc (-80:80:{\x} and {\y});
\draw  (B) arc (80:280:{\x} and {\y});
\draw (D) arc (-80:80:{\x} and {\y});
\draw (A)--(B)(C)--(D)(U)--(V)(N)--(V);
\draw[<->](I)--(J);
\draw[<->](P)--(Q);
\draw[<->](G)--(R);
\draw[dashed](C)--(I)(D)--(J)(B)--(P)(C)--(Q)(M)--(N)(M)--(U)(T)--(R);

\end{tikzpicture}
\end{center}
Tính thể tích của nước trong téc (kết quả làm tròn đến hàng phần trăm)?
\shortans{$1{,}9$}
\loigiai{
\immini{Thế tích phần dầu còn lại sẽ bằng diện tích hình phẳng gạch sọc trong hình nhân với chiều dài của bồn (chiều cao của trụ).

Đường tròn có tâm $O(0;0)$, $R=0{,}5$ có phương trình là
\[x^2+y^2=0{,}25 \Leftrightarrow y=\pm \sqrt{0{,}25-x^2}.\]
Diện tích hình gạch sọc chính là diện tích hình phẳng giới hạn bởi các đường
\[y=\sqrt{0{,}25-x^2};\,y=-\sqrt{0{,}25-x^2};\,x=-0{,}5;\,x=0{,}25.\]
Do đó
\[V=Sh=3 \displaystyle\int_{-0{,}5}^{0{,}25}\left|\sqrt{0{,}25-x^2}-\left(-\sqrt{0{,}25-x^2}\right)\right|\mathrm{\,d}x \approx 1{,}896\mathrm{\,m}^3 \approx 1{,}9\mathrm{\,m}^3.\]
}{
\begin{tikzpicture}[line join=round,line cap=round, font=\footnotesize,scale=1,>=stealth]
\def \r{1.5}
\fill[pattern=north east lines]plot[domain=60:300](\r*cos \x,\r*sin \x)--cycle;
\draw[-stealth](-2,0)--(2.5,0)node[below]{$x$};
\draw[-stealth](0,-2)--(0,2)node[right]{$y$};
\fill (0,0) circle(1pt)node[below left]{$O$}(\r,0)circle(1pt)+(-0.1,0.2) node[right]{$0{,}5$}(-\r,0)circle(1pt)+(0.1,0.2) node[left]{$-0{,}5$}(0,-\r)circle(1pt)+(0.1,-0.2) node[left]{$-0{,}5$}(0,\r)circle(1pt)+(0.1,0.2) node[left]{$0{,}5$}(0.5*\r,0)circle(1pt)+(-0.1,-0.2) node[right]{$0{,}25$};
\draw (60:\r)--(-60:\r);
\draw (0,0) circle(\r);
\end{tikzpicture}
}
}
\end{ex}

\begin{ex}%[12-EX-2-2025, Diamond]%[2D4V3-4]
\immini{
Một chiếc lều mái vòm có hình dạng như hình bên. Nếu cắt lều bằng mặt phẳng song song với mặt đáy và cách mặt đáy một khoảng $x$ (mét) ($0 \le x \le 3$) thì được hình chữ nhật có các kích thước lần lượt là $x$ và $\sqrt{9-x^2}$. Tính thể tích cái lều (đơn vị m$^3$).
}{
\tdplotsetmaincoords{80}{110}
\begin{tikzpicture}[tdplot_main_coords, scale=2, font=\footnotesize, line cap=round, line join=round, >=stealth]
% --- 2. VẼ MẶT PARABOLOID ---
% Phương trình: z = 2 - r^2 (với r^2 = x^2 + y^2)
%Vẽ hình vuông ở độ cao h=1.
\path (1,0,1)coordinate(A) (-1,0,1)coordinate(C) (0,1,1)coordinate(B) (0,-1,1)coordinate(D) (1.414,0,0)coordinate(A') (-1.414,0,0)coordinate(C') (0,1.414,0)coordinate(B') (0,-1.414,0)coordinate(D');
\fill[yellow, opacity=0.5] (A)--(B)--(C)--(D)--cycle;
\draw (D)--(A)--(B) (D')--(A')--(B');
\draw[dashed] (B)--(C)--(D) (B')--(C')--(D') (0,0,2)--(0,0,0);
% Vẽ đường biên (profile) của parabol
% Dùng tham số hoá để vẽ các đường sinh
\draw[thick, dashed] plot[domain=-1.414:0, samples=50] (\x, 0, {2-\x*\x}); % Đường trên mặt xOz
\draw[thick] plot[domain=0:1.414, samples=50] (\x, 0, {2-\x*\x}); % Đường trên mặt xOz
\draw[thick] plot[domain=-1.414:1.414, samples=50] (0, \x, {2-\x*\x}); % Đường trên mặt yOz
\end{tikzpicture}
}
\shortans{$9$}
\loigiai{
Diện tích hình chữ nhật là $S(x)=x\sqrt{9-x^2}$.\\
Thể tích của cái lều là $V=\displaystyle\int\limits_0^3x\sqrt{9-x^2}\mathrm{\,d}x=9$ m$^3$.
}
\end{ex}

\begin{ex}%[12-EX-3-2026, Nguyên Quang Hiệp]%[2D4V3-4]
Một mảnh đất có hình dạng là hình thang cong có các thông số như hình vẽ dưới đây, biết phần đường cong là phần đồ thị của hàm số $y=a\sqrt{x}$ ($a>0$). Diện tích của mảnh đất đó là bao nhiêu? (kết quả làm tròn đến hàng phần chục).
\begin{center}
\begin{tikzpicture}[>=stealth,line join=round,line cap=round,font=\footnotesize,x=0.5cm,y=0.45cm]
\tikzset{declare function={
f(\x)=2*sqrt(\x);
}}
% Vẽ parabol y=2√x và y=-2√x
\begin{scope}[rotate=-90]
\draw[thick,blue,smooth,samples=100,domain=8:12] plot(\x,{f(\x)});
\draw[thick,blue,smooth,samples=100,domain=8:12] plot(\x,{-f(\x)});
% Vùng tô màu từ x=8 đến x=12
\fill[blue!15,opacity=0.6]
(8,0)
-- plot[domain=8:12,smooth,samples=50] (\x,{f(\x)})
-- (12,0)
-- plot[domain=12:8,smooth,samples=50] (\x,{-f(\x)})
-- cycle;
% Đường thẳng đứng x=8 và x=12
\draw[thick] (8,0)--(8,{f(8)})coordinate (B);
\draw[thick] (8,0)node[shift={(90:2mm)}]{$8$}--(8,{-f(8)})coordinate (A);
\draw[thick] (12,0)node[shift={(-90:2mm)}]{$12$}--(12,{f(12)})coordinate (C);
\draw[thick] (12,0)--(12,{-f(12)})coordinate (D);
\draw (8,0)coordinate (M)--(12,0)node[midway,right]{$h=5$};
\end{scope}
\foreach \point/\goc in {A/90,B/90,C/-90,D/-90}{
\draw[fill=black](\point)circle(.8pt)+(\goc:2mm)node[scale=.8]{$\point$};
}
\end{tikzpicture}
\end{center}
\shortans[oly]{50{,}7}
\loigiai{
\begin{center}
\begin{tikzpicture}[>=stealth,line join=round,line cap=round,font=\footnotesize,x=0.5cm,y=0.45cm]
\tikzset{declare function={
f(\x)=2*sqrt(\x);
}}
\draw[thick,blue,smooth,samples=100,domain=8:12] plot(\x,{f(\x)});
\draw[thick,blue,smooth,samples=100,domain=8:12] plot(\x,{-f(\x)});
% Vùng tô màu từ x=8 đến x=12
\fill[blue!15,opacity=0.6]
(8,0)
-- plot[domain=8:12,smooth,samples=50] (\x,{f(\x)})
-- (12,0)
-- plot[domain=12:8,smooth,samples=50] (\x,{-f(\x)})
-- cycle;
% Đường thẳng đứng x=8 và x=12
\draw[thick] (8,0)--(8,{f(8)})coordinate (B);
\draw[thick] (8,0)node[shift={(180:2mm)}]{$8$}--(8,{-f(8)})coordinate (A);
\draw[thick] (12,0)node[shift={(0:2mm)}]{$12$}--(12,{f(12)})coordinate (C);
\draw[thick] (12,0)--(12,{-f(12)})coordinate (D);
\draw (8,0)coordinate (M)--(12,0)node[midway,above]{$h=5$};
\foreach \point/\goc in {A/-90,B/90,C/90,D/-90}{
\draw[fill=black](\point)circle(.8pt)+(\goc:2mm)node[scale=.8]{$\point$};
}
\end{tikzpicture}	\qquad
\begin{tikzpicture}[>=stealth,line join=round,line cap=round,font=\footnotesize,x=0.5cm,y=0.45cm]
\tikzset{declare function={
f(\x)=2*sqrt(\x);
}}
% Vẽ parabol y=2√x và y=-2√x
\draw[thick,blue,smooth,samples=100,domain=0:13] plot(\x,{f(\x)});
\draw[thick,blue,smooth,samples=100,domain=0:13] plot(\x,{-f(\x)});
% Vùng tô màu từ x=8 đến x=12
\fill[blue!15,opacity=0.6]
(8,0)
-- plot[domain=8:12,smooth,samples=50] (\x,{f(\x)})
-- (12,0)
-- plot[domain=12:8,smooth,samples=50] (\x,{-f(\x)})
-- cycle;
% Trục tọa độ
\draw[->] (-0.5,0)--(14,0) node[above]{$x$};
\draw[->] (0,-8)--(0,8) node[left]{$y$};
\filldraw(0,0)circle(1.5pt)node[below left]{$O$};
% Đường thẳng đứng x=8 và x=12
\draw[dashed,thick] (8,0)--(8,{f(8)})coordinate (A);
\draw[dashed,thick] (8,0)--(8,{-f(8)})coordinate (D);
\draw[dashed,thick] (12,0)--(12,{f(12)})coordinate (B);
\draw[dashed,thick] (12,0)--(12,{-f(12)})coordinate (C);
\path (A)--(B)node[above=1mm,midway]{$y=2\sqrt{x}$} (C)--(D)node[below=1mm,midway]{$y=-2\sqrt{x}$};
\draw[|<->|] ($(A)+(180:2mm)$)--($(D)+(180:2mm)$)node[pos=.8,left]{$8$};
\draw[|<->|] ($(C)+(0:2mm)$)--($(B)+(0:2mm)$)node[pos=.8,right]{$12$};
\end{tikzpicture}
\end{center}
Chọn hệ trục tọa độ $Oxy$ sao cho $Ox$ là trục đối xứng nằm ngang của mảnh đất, chiều dương hướng từ đáy nhỏ sang đáy lớn. Gốc tọa độ $O$ được chọn sao cho phù hợp với hàm số $y=a\sqrt{x}$.\\
Theo hình vẽ và giả thiết
\begin{itemize}
\item Tại vị trí đáy nhỏ (kích thước $8$): Tung độ các điểm biên là $y = \pm 4$.
\item Tại vị trí đáy lớn (kích thước $12$): Tung độ các điểm biên là $y = \pm 6$.
\item Chiều cao $h=5$ là khoảng cách giữa hai đáy dọc theo trục $Ox$.
\end{itemize}
Gọi hoành độ của đáy nhỏ là $x_1$ và đáy lớn là $x_2$. Ta có $x_2 - x_1 = 5$.\\
Do biên của mảnh đất nằm trên đồ thị $y=a\sqrt{x}$ nên tại $x_1$, ta có $a\sqrt{x_1} = 4 \Rightarrow a^2 x_1 = 16$.\\
Tại $x_2$, ta có $a\sqrt{x_2} = 6 \Rightarrow a^2 x_2 = 36$.\\
Suy ra $a^2(x_2 - x_1) = 36 - 16 = 20$. \\
Mà $x_2 - x_1 = 5$ nên $a^2 \cdot 5 = 20 \Rightarrow a^2 = 4 \Rightarrow a = 2$ (vì $a>0$).\\
Khi đó $x_1=\dfrac{16}{4}=4$ và $x_2=\dfrac{36}{4}=9$.\\
Diện tích mảnh đất được giới hạn bởi $y=2\sqrt{x}$, $y=-2\sqrt{x}$, $x=4$, $x=9$ là
\[S = \displaystyle\int\limits_{4}^{9} \big[2\sqrt{x} - (-2\sqrt{x})\big] \mathrm{\,d}x = \displaystyle\int\limits_{4}^{9} 4\sqrt{x} \mathrm{\,d}x = 4 \cdot \dfrac{2}{3}x\sqrt{x} \bigg|_4^9 = \dfrac{8}{3} (27 - 8) = \dfrac{152}{3} \approx 50{,}7.\]
}
\end{ex}

\begin{ex}%[12-EX-3-2025, Phạm Hoài]%[2D4V3-4]
\immini{ Một khối bê tông có chiều cao $2$ m được đặt trên mặt đất. Nếu cắt khối bê tông bởi một mặt phẳng nằm ngang và cách mặt đất $x$ m ($0 \leq x \leq 2$) thì được mặt cắt là một hình chữ nhật có chiều dài $5$ m và chiều rộng bằng $(0{,}5)^x$ m (\textit{tham khảo hình vẽ bên}). Tính thể tích (đơn vị m$^3$) của khối bê tông đó (\textit{làm tròn kết quả đến hàng phần trăm}).}
{\begin{tikzpicture}[scale=0.7, font=\footnotesize, line join=round, line cap=round, >=stealth]
\path
(0,0) coordinate (A)
(-3,-1) coordinate (B)
(5,0) coordinate (D)
($(B)+(D)-(A)$) coordinate (C)
(1.5,4) coordinate (A')
(-1.5,3) coordinate (B')
(3.5,4) coordinate (D')
($(B')+(D')-(A')$) coordinate (C')
;

\draw (B)--(C)--(D);
\draw (A')--(B')--(C')--(D')--(A');
\path (A)..controls++(10:0.3)and++(-90:0.65)..(A')foreach \i in{0,1,2,...,10}{coordinate[pos=\i/15](A\i)};;
\path (B)..controls++(10:0.3)and++(-90:0.65)..(B')foreach \i in{0,1,2,...,10}{coordinate[pos=\i/15](B\i)};
\path (C)..controls++(170:0.3)and++(-90:0.65)..(C')foreach \i in{0,1,2,...,10}{coordinate[pos=\i/15](C\i)};
\path (D)..controls++(170:0.3)and++(-90:0.65)..(D')foreach \i in{0,1,2,...,10}{coordinate[pos=\i/15](D\i)};
\fill[green!30!yellow](A4)--(B4)--(C4)--(D4)--cycle;
\draw[dashed] (A)--(B)(D)--(A);
\draw[dashed] (A)..controls++(10:0.3)and++(-90:0.65)..(A');
\draw (B)..controls++(10:0.3)and++(-90:0.65)..(B');
\draw (C)..controls++(170:0.3)and++(-90:0.65)..(C');
\draw (D)..controls++(170:0.3)and++(-90:0.65)..(D');
\draw[dashed] (B4)--(A4)--(D4);
\draw (B4)--(C4)--(D4);
\path (B4)--(C4)node[pos=0.5,below]{$0{,}5^x$ m};
\path (2.7,-0.2)node[right]{$5$ m};
\draw[<->](0.5,-1)--(0.5,-0.35)node[pos=0.5,right]{$x$ m};
\draw[<->](-3.5,-1)--(-3.5,3)node[pos=0.5,left]{$2$ m};
\end{tikzpicture}
}
\shortans[]{$5{,}41$}
\loigiai{
Ta có diện tích mặt cắt khối bê tông là $S(x)=5 \cdot 0{,}5^x$.\\
Do đó thể tích của khối bê tông là $
V=\displaystyle\int\limits_0^25\cdot  0{,}5^x \mathrm{\,d}x=\dfrac{5\cdot 0{,}5^x}{\ln 0{,}5}\bigg|_0^2\approx 5{,}41$ m$^3$.
}
\end{ex}

\begin{ex}%[Giữa học kì 2, Trần Cao Vân, Khánh Hòa, 2025]%[12-EX-3-2026, Thành Đức Trung]%[2D4V3-4]
\immini
{
Chướng ngại vật “tường cong” trong một sân thi đấu X-Game là một khối bê tông có chiều cao từ mặt đất lên là $3{,}5$ m. Giao của mặt tường cong và mặt đất là đoạn thẳng $AB=2$ m. Thiết diện của khối tường cong cắt bởi mặt phẳng vuông góc với $AB$ tại $A$ là một hình tam giác vuông cong $ACE$ với $AC=4$ m, $CE=3{,}5$ m và cạnh cong $AE$ nằm trên một đường parabol có trục đối xứng
}{
\begin{tikzpicture}[scale=0.8,>=stealth, font=\footnotesize, line join=round, line cap=round]
%Vẽ gạch cong
\foreach \x in {0,.5,...,2.5}{
\begin{scope}[shift={(150:\x)}]
\draw (0,0) to[out=20,in=-120]
coordinate[pos=.1](X1)
coordinate[pos=.2](X2)
coordinate[pos=.3](X3)
coordinate[pos=.4](X4)
coordinate[pos=.5](X5)
coordinate[pos=.6](X6)
coordinate[pos=.7](X7)
coordinate[pos=.8](X8)
coordinate[pos=.9](X9)
coordinate[pos=1](X10)
++(4,3.5);
\draw (150:.25) to[out=20,in=-120]
coordinate[pos=.05](Y1)
coordinate[pos=.15](Y2)
coordinate[pos=.25](Y3)
coordinate[pos=.35](Y4)
coordinate[pos=.45](Y5)
coordinate[pos=.55](Y6)
coordinate[pos=.65](Y7)
coordinate[pos=.75](Y8)
coordinate[pos=.85](Y9)
coordinate[pos=.95](Y10)
++(4,3.5);
\foreach \i in {1,...,10}\draw (X\i)--++(150:.25) (Y\i)--++(150:.25);
\end{scope}
}
%======
\path (0,0)node[below left]{$A$}to[out=20,in=-120]coordinate[pos=.4](A)(4,3.5)node[above right]{$E$};
\draw[dashed] (4,0)--++(150:3)--++(-4,0) ($(4,0)+(150:3)$)--++(0,3.5);
\draw($(0,0)!(A)!(4,0)$) node[below]{$M$}-- node[right]{1m}(A)
(4,1.75) node[right]{3{,}5m} (1.2,0) node[below]{4m};
\draw (0,0)to[out=20,in=-120](4,3.5)--(4,0)node[below]{$C$}--cycle;
\draw (0,0)--++(150:3)node[below left]{$B$}node[midway,below]{$2$m} to[out=20,in=-120]++(4,3.5)--(4,3.5);
\draw[fill](0,0) circle (1pt) (4,3.5) circle (1pt);
\end{tikzpicture}
}\noindent vuông góc với mặt đất. Tại vị trí $M$ là trung điểm của $AC$ thì tường cong có độ cao $1$ m (xem hình minh họa bên). Tính thể tích bê tông cần sử dụng để tạo nên khối tường cong đó.
\shortans[]{$10$}
\loigiai
{
\immini
{
Chọn hệ trục $Oxy$ như hình vẽ sao cho $A\equiv O$. \\
Do đó cạnh cong $AE$ nằm trên parabol $(P)\colon y=ax^2+bx$ đi qua các điểm $(2; 1)$ và $\left(4;\dfrac{7}{2}\right)$ nên $(P)\colon y=\dfrac{3}{16}x^2+\dfrac{1}{8}x$.\\
Vì khi cắt tường cong bởi mặt phẳng vuông góc với $Ox$ tại điểm $x$ có hoành độ $0\leq x\leq 4$ ta được thiết diện là hình chữ nhật có diện tích $S(x)=2\left(\dfrac{3}{16}x^2+\dfrac{1}{8}x\right)$.\\
%Khi đó diện tích tam giác cong $ACE$ có diện tích
Vậy thể tích của khối tường cong
\[V=\displaystyle\int\limits_0^4 2\left(\dfrac{3}{16}x^2+\dfrac{1}{8}x\right)\mathrm{\,d}x=10\;\mathrm{m}^3.\]
%	Vậy thể tích khối bê tông cần sử dụng là $V=5\cdot 2=10\;\mathrm{m}^3.$
}{
\begin{tikzpicture}[scale=0.8,>=stealth, font=\footnotesize, line join=round, line cap=round]
\draw[latex-latex] (0,5)node[left]{$y$}|-(5,0)node[below]{$x$};
\draw[fill=black!20] (0,0)to[out=20,in=-120](4,3.5)--(4,0)node[below]{$4$}--cycle;
%Vẽ gạch cong
\foreach \x in {0,.5,...,2.5}{
\begin{scope}[shift={(150:\x)}]
\draw (0,0) to[out=20,in=-120]
coordinate[pos=.1](X1)
coordinate[pos=.2](X2)
coordinate[pos=.3](X3)
coordinate[pos=.4](X4)
coordinate[pos=.5](X5)
coordinate[pos=.6](X6)
coordinate[pos=.7](X7)
coordinate[pos=.8](X8)
coordinate[pos=.9](X9)
coordinate[pos=1](X10)
++(4,3.5);
\draw (150:.25) to[out=20,in=-120]
coordinate[pos=.05](Y1)
coordinate[pos=.15](Y2)
coordinate[pos=.25](Y3)
coordinate[pos=.35](Y4)
coordinate[pos=.45](Y5)
coordinate[pos=.55](Y6)
coordinate[pos=.65](Y7)
coordinate[pos=.75](Y8)
coordinate[pos=.85](Y9)
coordinate[pos=.95](Y10)
++(4,3.5);
\foreach \i in {1,...,10}\draw (X\i)--++(150:.25) (Y\i)--++(150:.25);
\end{scope}
}
%======
\path (0,0)node[below left]{$A$}to[out=20,in=-120]coordinate[pos=.4](A)(4,3.5)node[above right]{$E$};
\draw[dashed]($(0,0)!(A)!(4,0)$) node[below]{$2$}-- (A)--($(0,0)!(A)!(0,4)$)node[left]{$1$} (4,0)--++(150:3)--++(-4,0) (4,3.5)--(0,3.5)node[left]{$3{,}5$} ($(4,0)+(150:3)$)--++(0,3.5);
\draw (0,0)--++(150:3)node[below left]{$B$}node[midway,below]{$2$m} to[out=20,in=-120]++(4,3.5)--(4,3.5);
\draw[fill](0,0) circle (1pt) (4,3.5) circle (1pt) (A)circle (1pt);
\end{tikzpicture}
}
}
\end{ex}

\begin{ex}%[12-EX-3-2025, Thái Bảo]%[2D4V3-4]
\immini{
Người ta cần đổ một khối bê tông cao $2\,\text{m}$ để làm trụ chân cầu được đặt trên mặt đất phẳng như hình vẽ. Nếu cắt khối bê tông này bằng mặt phẳng nằm ngang, cách mặt đất $x$ (m) ($0\le x\le2$) thì được thiết diện là một hình chữ nhật có chiều dài $10\,\text{m}$, chiều rộng là $(0{,}5)^x\,\text{m}$. Mỗi mét khối bê tông có giá $5$ triệu đồng. Hãy tính số tiền cần đổ bê tông cho trụ chân cầu (làm tròn đến hàng phần chục).
\par\shortans{$54{,}1$}
}{
\begin{tikzpicture}[font=\footnotesize, line join=round, line cap=round,>=stealth,scale=.8]
\path (0,0) coordinate (A)
(2.46,0.58) coordinate (B)
(-3,1.9) coordinate (a)
($(B)+(a)$) coordinate (B')
($(A)+(a)$) coordinate (A')
($(B)+(0,4.6)$) coordinate (C)
($(C)+(a)$) coordinate (C')
(2.46,2.78) coordinate (D)
($(D)+(a)$) coordinate (D')
(1.46,2.51) coordinate (E)
($(E)+(a)$) coordinate (E')
(2,5) coordinate (F)
($(F)+(a)$) coordinate (F')
;
\fill [fill=cyan!50](D)--(D')--(E')--(E)--cycle;
\draw (A')--(A)--(B)--(C)--(F) (D)--(E)--(E') (C)--(C')--(F')--(F)--cycle;
\draw [dashed] (D)--(D')node [midway,shift={(0,0.5)}] {$10$ m}--(E') (A')--(B')--(B) (B')--(C');
\draw [domain=0:2, samples=100]
plot (\x, {0.36*\x^3+0.23*\x^2 +0.61*\x});
\draw[domain=0:2, samples=100, shift={(-3,1.9)}]
plot (\x, {0.36*\x^3 + 0.23*\x^2 + 0.61*\x});
\node at (0,1) {$(0{,}5)^x$ m};
\path (B)--(C)--([turn]-90:0.35) coordinate (xt)
($(B)-(C)+(xt)$) coordinate (yt);
\draw[dash pattern=on 2pt off 2pt] (B)--(yt) (C)--(xt);
\draw[>=stealth,|<->|] (xt)--(yt) node[midway,above right]{$2$ m};
\path (B)--(D)--([turn]-90:0.7) coordinate (zt)
($(B)-(D)+(zt)$) coordinate (tt);
\draw[dash pattern=on 2pt off 2pt] (B)--(tt) (D)--(zt);
\draw[>=stealth,|<->|] (zt)--(tt) node[midway,right]{$x$ m};
\end{tikzpicture}
}
\loigiai{
Tại độ cao $x$ (m), thiết diện ngang của khối bê tông là hình chữ nhật có diện tích
\[
S(x)=10\cdot(0{,}5)^x\ (\text{m}^2).
\]
Thể tích khối bê tông là
\begin{eqnarray*}
V&=&\displaystyle\int\limits_0^2 S(x)\mathrm{\,d}x\\
&=&10\displaystyle\int\limits_0^2(0{,}5)^x\mathrm{\,d}x\\
&=&10\left[\dfrac{(0{,}5)^x}{\ln(0{,}5)}\right]_0^2\\
&=&10\cdot\dfrac{(0{,}5)^2-1}{\ln(0{,}5)}\\
&\approx&10{,}82
\end{eqnarray*}
Giá tiền cần đổ bê tông là
\[
10{,}82\cdot5=54{,}1\ (\text{triệu đồng}).
\]
Vậy số tiền cần đổ bê tông cho trụ chân cầu là khoảng $54{,}1$ triệu đồng.
}
\end{ex}

\begin{ex}%[2D4V3-4]
Cắt một vật thể bởi hai mặt phẳng vuông góc với trục $Ox$ tại $x=1$; $x=3$. Khi cắt một vật thể bởi mặt phẳng vuông góc với trục $Ox$ tại điểm có hoành độ $x, (1\leq x\leq 3)$, mặt cắt là tam giác vuông có một góc $45^\circ$ và độ dài một cạnh góc vuông là $\sqrt{4-\dfrac{1}{2}x^2}$. Tính thể tích vật thể trên, kết quả làm tròn đến hàng phần trăm.
\shortans{1{,}83}
\loigiai{
Diện tích tam giác vuông cân là $S(x) = \dfrac{1}{2}\sqrt{4-\dfrac{1}{2}x^2}\cdot\sqrt{4-\dfrac{1}{2}x^2}=\dfrac{1}{2}\left(4-\dfrac{1}{2}x^2\right)$.\\
Thể tích vật thể là $V=\displaystyle\int\limits_1^3\dfrac{1}{2}\left( 4-\dfrac{1}{2}x^2\right)\mathrm{\,d}x=\dfrac{11}{6}\approx 1{,}83$.
}
\end{ex}

\begin{ex}%[Tổ 20 - Đợt 17 - Chương 4 - - CD - Đề 8]%[Huỳnh Xuân Tín]%[2D4V3-5]
\immini{Một chiếc ly có hình dạng tròn xoay và kích thước như hình vẽ, thiết diện dọc của cốc (đường sinh) là một đường Parabol.
Tính thể tích của cốc. (kết quả làm tròn đến hàng đơn vị).}{	\definecolor{almond}{rgb}{0.94, 0.87, 0.8}%màu ly
\definecolor{anti-flashwhite}{rgb}{0.95, 0.95, 0.96}%màu miệng ly
\begin{tikzpicture}[line join=round, line cap=round,scale=1,transform shape,line width=.3mm]

\tikzset{cai_ly/.pic={
\path
(1.85,4)coordinate (A)
(-1.85,4) coordinate (B)
(-3,3)coordinate (C)
(-3,-1.5)coordinate (D)
;
\draw (A)--(B) (C)--(D);
\foreach\p in {A,B,C,D}
{\draw[fill=black](\p) circle (1pt);}
\node at (-3.8,1) {$10$ cm};
\node at (0,4.5) {$8$ cm};
\draw (2,1)--(5,2.5) node[above] {Parabol};

\draw[fill=almond!50!black] (1.75,-4.93)  arc (0:360:1.75 cm and .54cm);
\draw[fill=almond] (1.75,-4.8)  arc (0:360:1.75 cm and .4cm);

\draw[fill=anti-flashwhite] (1.85,3)  arc (0:360:1.85 cm and .3cm);
\draw[fill=almond] (1.85,3)
..controls +(-95:1.3) and +(30:1.7) ..(.4,-1.6)
..controls +(-150:.1) and +(90:2) ..(.2,-4.2)
..controls +(-150:.1) and +(-30:.1) ..(-.2,-4.2)
..controls +(90:2) and +(-30:.1) ..(-.4,-1.6)
..controls +(150:1.7) and +(-85:1.3) ..(-1.85,3)
arc (-180:0:1.85 cm and .3cm);
;
\draw[fill=almond] (.2,-4.2)
..controls +(-150:.1) and +(-30:.1) ..(-.2,-4.2)
..controls +(180:.1) and +(60:.1) ..(-.4,-4.4)
..controls +(180:.3) and +(60:.1) ..(-.9,-4.8)
..controls +(-30:.4) and +(-150:.4) ..(.9,-4.8)
..controls +(120:.1) and +(0:.3) ..(.4,-4.4)
..controls +(150:.1) and +(-30:.1) ..(.2,-4.2)
;

}}

\path
(0,0)pic[scale=1]{cai_ly};
\end{tikzpicture}}
\shortans{$251$}
\loigiai
{Ta chọn hệ tọa độ như hình vẽ
\begin{center}
\begin{tikzpicture}[line cap=round,line join=round, >=stealth,scale=0.56]
\def\a{0.625}
\def\b{0}
\def\c{0}
\def\f(#1){\a*((#1)^2)+\b*(#1)+\c}
\pgfmathsetmacro\ct{-\b/(2*\a)}
\def\xmin{-3.5}
\def\xmax{3.5}
\def\ymin{-3}
\def\ymax{7}
\draw[->] (\xmin-1,0)--(\xmax+1,0) node[below] { $x$};
\draw[->] (0,-1)--(0,\ymax+1) node[left] { $y$};
\draw (0,0) node [below left] { $O$};
%	\foreach \x in {-3,-2,-1,1,2,3}
%	\draw (-3,0)--(-3,-0.1) node [below] { $\x$};
%	\foreach \y in {-3,-2,-1,1,2,3}		\draw (0.1,\y)--(-0.1,\y) node [left] { $\y$};
%	\clip (\xmin,\ymin) rectangle (\xmax,\ymax);
\draw[thick,smooth,samples=200] plot[domain=\xmin:\xmax] (\x,{\f(\x)});
\draw[dashed] (\ct,0)--(\ct,{\f(\ct)})--(0,{\f(\ct)});
\fill (\ct,{\f(\ct)}) circle (1pt);
%	\draw (\xmax,{\f(\xmax)}) node [above right]{${x^2}+2{x}-2$};
\end{tikzpicture}
\end{center}
Khi này Parabol có phương trình $y=\dfrac{5}{8} x^2 \Leftrightarrow x^2=\dfrac{8}{5} y$.\\
Thể tích tối đa cốc $V=\pi \displaystyle\int\limits_0^{10}\left(\dfrac{8}{5} y\right) \cdot\mathrm{\,d}y \approx 251$.
}
\end{ex}

\begin{ex}%[Tổ 20 - Đợt 17 - Chương 4 - - CD - Đề 8]%[Huỳnh Xuân Tín]%[2D4V3-5]
Một chiếc đồng hồ cát như hình vẽ, gồm hai phần đối xứng nhau qua mặt nằm ngang và đặt rong một hình trụ. Thiết diện thẳng đứng qua trục của nó là hai parabol chung đỉnh và đối xứng nhau qua mặt nằm ngang. Ban đầu lượng cát dồn hết ở phần trên của đồng hồ thì chiều ao $h$ của mực cát bằng $\dfrac{3}{4}$ chiều cao của nửa trên đồng hồ (xem hình).




Cát chảy từ trên xuống dưới với lưu lượng không đổi $2{,}9$ cm$^3$/phút. Khi chiều cao của cát còn $4$ cm thì bề mặt trên cùng của cát tạo thành một đường tròn chu vi $8 \pi$ (cm) (xem hình). Biết sau $30$ phút thì cát chảy hết xuống phần bên dưới của đồng hồ. Tính chiều cao của khối trụ bên ngoài. (Kết quả làm tròn đến hàng đơn vị)
\shortans{}
\loigiai
{Xét thiết diện chứa trục theo phương thẳng đứng của đồng hồ cát là parabol. Gọi $(P)$ là đường Parabol phía trên. Chọn hệ trục $O x y$ như hình vẽ.\\
Đường tròn thiết diện có chu vi bằng $8 \pi=2 \pi R \Leftrightarrow R=4$.\\
Do $(P)$ có đỉnh là $O(0 ; 0)$ và đi qua điểm có tọa độ $(4 ; 4)$ nên phương trình $(P)\colon y=\dfrac{1}{4} x^2$.\\
Thể tích phần cát ban đầu chính bằng thể tích khối tròn xoay sinh ra khi quay nhánh phải của $(P)$ quay quanh trục $O y$ và bằng lượng cát đã chảy trong thời gian $30$ phút.
\[
V=\pi \displaystyle\int\limits_0^h(2 \sqrt{y})^2 \mathrm{\,d} y=2 \pi h^2.		\]
Ta có lượng cát chảy trong 30 phút là $2{,}9\cdot30=87\left(\mathrm{~cm}^3\right)$.\\
Vậy $V=2 \pi h^2=87 \Leftrightarrow h \approx 3{,}7$ (cm).\\
Khi đó chiều cao hình trụ bên ngoài là $l=2 h \cdot \dfrac{4}{3} \approx 10$ (cm).
}
\end{ex}

\begin{ex}%[ST12-dot1-Trần xuân Hòa]%[2D4V3-5]
Du khách ghé thăm Bình Định không thể bỏ qua địa danh Tháp Bánh Ít nổi tiếng. Tháp có hai cửa, mỗi cửa có hình dáng là một cung Parabol nằm cùng một trục (hướng Đông - Tây). Hai cửa cách nhau $ 8 $ mét, có chiều cao $ 4 $ mét, lối đi rộng $ 1 $ mét thông hai cửa với nhau. Hãy tính thể tích phần không gian lối đi giới hạn giữa hai cửa. (Kết quả làm tròn đến chữ số thập phân thứ nhất).
\shortans{$21{,}3$}
\loigiai{
\immini{
Xét một cửa của Tháp Bánh Ít, đặt hệ trục $ Oxy $ như hình vẽ. Dễ dàng tìm được phương trình của parabol là $ y=-16x^2+4 $.\\
Diện tích của parabol là $ S=\displaystyle \int \limits^{0,5}_{-0,5} \left (-16x^2+4 \right )\mathrm{\,d}x =\dfrac{8}{3}$.}{
\begin{tikzpicture}[scale=1, line join=round, line cap=round,>=stealth]
\draw[->,line width = 1pt] (-2,-2.25)--(-2,-2.25)--(2,-2.25) node[below]{$x$};
\draw[->,line width = 1pt] (0,-3) --(0,1) node[right]{$y$};
\draw [<->] (1.5,-2.25)--(-1.5,-2.25);
\draw [<->] (0,0)--(0,-2.25);
\draw (0,-1.125) node[right]{$4 \mathrm{m}$};
\draw (0,-2.25) node[below right]{$1 \mathrm{m}$} ;
\draw (0,-2.25) node[above left]{$O$} ;
\draw (0,0) parabola (1.5,-2.25) ;
\draw (0,0) parabola (-1.5,-2.25) ;
\end{tikzpicture}
}
\noindent Thể tích phần không gian lối đi giới hạn giữa hai cửa là $ 8\cdot \dfrac{8}{3}=\dfrac{64}{3}\approx 21{,}3$ m$^3$.
}
\end{ex}

\begin{ex}%[12-MH-2-MH2025]%[MH-2025, Chu Hà]%[2D4V3-5]
Nghệ thuật làm gốm có lịch sử phát triển lâu đời và vẫn còn tồn tại cho đến ngày nay. Giả sử một bình gốm có mặt trong của bình là một mặt tròn xoay sinh ra khi cho phần đồ thị của hàm số $y=\dfrac{1}{175} x^2+\dfrac{3}{35} x+5$ $\left(0 \leq x \leq 30\right)$ ($x$, $y$ tính theo cm) quay tròn quanh bệ gốm có trục trùng với trục hoành $O x$. Hỏi để hoàn thành bình gốm đó ta cần sử dụng bao nhiêu cm$^3$ đất sét, biết rằng bình gốm đó có độ dày không đổi là $1$ cm.
\shortans[]{$94{,}2$}
\loigiai{
Thể tích phần bên trong của bình gốm là
\[V_1=\pi \displaystyle \int\limits_0^{30} \left(\dfrac{1}{175} x^2+\dfrac{3}{35} x+5\right)^2  \mathrm{\,d}x=240\pi \, (\text{cm}^3).\]
Thể tích của bình gốm bao gồm cả độ dày là
\[V_2=\pi \displaystyle \int\limits_0^{30} \left(\dfrac{1}{175} x^2+\dfrac{3}{35} x+5+1\right)^2  \mathrm{\,d}x=270\pi \, (\text{cm}^3).\]
Thể tích cần dùng để hoàn thành bình gốm là
\[V=V_2-V_1=270 \pi-240\pi=30 \pi =94{,}2 \, (\text{cm}^3).\]
}
\end{ex}

\begin{ex}%[KNTT, mức độ 3]%[BG12-4in1, Đỗ Vũ Minh Thắng]%[2D4V3-5]
Nghệ thuật làm gốm có lịch sử phát triển lâu đời và vẫn còn tồn tại cho đến ngày nay. Giả sử một bình gốm có mặt trong của bình là một mặt tròn xoay sinh ra khi cho phần đồ thị của hàm số $y=\dfrac{1}{175} x^2+\dfrac{3}{35} x+5$ $\left(0 \leq x \leq 30\right)$ ($x$, $y$ tính theo cm) quay tròn quanh bệ gốm có trục trùng với trục hoành $O x$. Để hoàn thành bình gốm đó ta cần sử dụng $ x $ cm$^3$ đất sét. Biết rằng bình gốm đó có độ dày không đổi là $1$ cm. Tính $ \dfrac{x}{\pi} $.
\shortans{$ 510 $}
\loigiai{
Thể tích phần bên trong của bình gốm là
\[V_1=\pi \displaystyle \int\limits_0^{30} \left(\dfrac{1}{175} x^2+\dfrac{3}{35} x+5\right)^2  \mathrm{\,d}x=\dfrac{101586}{49}\pi \, (\text{cm}^3).\]
Thể tích của bình gốm bao gồm cả độ dày là
\[V_2=\pi \displaystyle \int\limits_0^{30} \left(\dfrac{1}{175} x^2+\dfrac{3}{35} x+5+1\right)^2  \mathrm{\,d}x=\dfrac{126576}{49}\pi \, (\text{cm}^3).\]
Thể tích cần dùng để hoàn thành bình gốm là
\[V=V_2-V_1=510 \pi \, (\text{cm}^3).\]
}
\end{ex}

\begin{ex}%[BG-12-4in1, Phạm Đức]%[2D4V3-5]
\immini
{
Một thùng rượu vang có dạng hình tròn xoay có hai đáy là hai hình tròn bằng nhau, khoảng cách giữa hai đáy bằng $80$ (cm). Đường sinh của mặt xung quanh thùng là một phần đường tròn có bán kính bằng $60$ (cm) (tham khảo hình vẽ). Hỏi thùng đó có thể đựng được bao nhiêu lít rượu? (Kết quả làm tròn đến hàng đơn vị).
}
{
\begin{tikzpicture}[scale=0.5, font=\footnotesize,line join=round, line cap=round, >=stealth]
\def\a{2}
\def\b{0.25}
\def\h{4}
\path
(0,0) coordinate (M)
($(M)+(2*\a,0)$) coordinate (N)
($(M)!0.5!(N)$)coordinate (O)
($(M)+(0,\h)$) coordinate (M')
($(N)+(0,\h)$) coordinate (N')
($(O)+(0,\h)$) coordinate (O')
;
\draw[dashed,thin] (M) arc (180:0:\a cm and \b cm);
\draw (O') ellipse (\a cm and \b cm)	(M) arc (-180:0:\a cm and \b cm);
\draw (M) arc(-135:-225:2.828);
\draw[dashed] (N') arc(45:135:2.828);
\draw (N) arc(-45:45:2.828);
\draw[dashed] (N) arc(-45:-135:2.828);
\end{tikzpicture}
}
\shortans{$771$}
\loigiai{
\immini
{
Gắn hệ trục tọa độ $Oxy$ như hình vẽ.\\
Ta có $80$ cm $= 8$ dm, $60$ cm $=6$ dm.\\
Ta lại có phương trình đường tròn
\[x^2+y^2=36\Leftrightarrow y^2=36-x^2.\]
Thể tích của thùng rượu là
\[V=\pi\displaystyle\int\limits_{-4}^{4}(36-x^2)\mathrm{\,d}x=\dfrac{736}{3}\pi\approx 771 \text{ dm}^3=771 \text{ lít}.\]
}
{
\begin{tikzpicture}[>=stealth,line join=round,line cap=round,font=\footnotesize,scale=0.5]
\pgfmathsetmacro{\g}{acos(4/6)}
\coordinate (A) at ($(0,0)+(180-\g:6)$);
\coordinate (B) at ($(0,0)+(\g:6)$);
\draw[dashed](-4,0)--(A)(4,0)--(B);
\draw (B) arc (\g:180-\g:6cm)($(0,0)+(10+\g:6)$)node[above,rotate=-20,yshift=5pt]{$y=\sqrt{36-x^2}$};
\draw[pattern = north east lines,draw=none] (4,0)--(B)arc (\g:180-\g:6cm)--(-4,0)--cycle;
\draw[->] (-4.5,0)--(5,0)node[above left]{\scriptsize $x$};
\draw[->] (0,-0.5)--(0,7)node[left]{\scriptsize $y$};
\foreach \x/\y/\z/\g in{-4/0/-4/-90,4/0/4/-90,0/0/O/-135,0/6/6/45}\draw[fill=black](\x,\y)node[shift={(\g:7pt)}]{$\z$}circle(1pt);
\end{tikzpicture}
}
}
\end{ex}

\begin{ex}%[Cánh Diều, mức độ 3]%[BG12-4in1, Đỗ Vũ Minh Thắng]%[2D4V3-5]
Xét chiếc chén (cái ly) trong bộ ấm chén uống trà, bạn Dương ước lượng được rằng chiếc chén được tạo thành khi cho hình phẳng giới hạn bởi đồ thị hàm số
$f(x)=0{,}14x^3-0{,}87x^2+1{,}92x+0{,}85$, trục hoành và hai đường thẳng $x=0$, $x=3$ quay quanh trục $Ox$ (đơn vị trên mỗi trục toạ độ là centimét). Tính thể tích của chiếc chén  (làm tròn đến hàng đơn vị của centimét khối).
\begin{center}
\begin{tikzpicture}[scale=0.8, font=\footnotesize, line join=round, line cap=round, >=stealth]
\tikzset{chen/.pic={\def\a{0.14*3^3-0.87*3^2+1.92*3+0.85}
\draw[brown, thick, fill=brown!40, fill opacity=0.6] (0,0)-- plot[smooth, samples=150, domain=0:3] (\x,{0.14*(\x)^3-0.87*(\x)^2+1.92*(\x)+0.85}) -- (3,\a) arc (90:270:{0.45*\a} and {\a}) -- plot[smooth, samples=150, fill=brown!40, domain=3:0] (\x,{-0.14*(\x)^3+0.87*(\x)^2-1.92*(\x)-0.85})--cycle;
\draw[brown, line width=1pt] (3,0) ellipse ({0.45*\a} and {\a});}}
\draw (-4.5,-2) pic[scale=0.8, rotate=90]{chen};
\draw (0,0) pic[scale=0.8]{chen};
\draw[->] (-1,0)--(0,0)node[below left]{$O$}--(5,0)node[below]{$x$};
\draw[->] (0,-3)--(0,3.5)node[left]{$y$};
\draw[blue, fill opacity=0.5, fill=cyan] (0,0) -- plot[smooth, samples=150, domain=0:3] (\x,{0.14*(\x)^3-0.87*(\x)^2+1.92*(\x)+0.85})--(3,0)--cycle;
\end{tikzpicture}
\end{center}
\shortans{$42$}
\loigiai
{
Thể tích khối tròn xoay tạo thành khi cho hình phẳng giới hạn bởi đồ thị hàm số
\begin{align*}
f(x)=0{,}14x^3-0{,}87x^2+1{,}92x+0{,}85,
\end{align*}
trục hoành và hai đường thẳng $x=0$, $x=3$ quay quanh trục $Ox$ là
\begin{eqnarray*}
V&=&\pi \displaystyle\int\limits_{0}^{3} \left(0{,}14x^3-0{,}87x^2+1{,}92x+0{,}85\right)^2 \mathrm{\,d}x\\
&=&\pi \displaystyle\int\limits_{0}^{3} \left(0{,}019x^6-0{,}2436x^5+1{,}2945x^4-3{,}1028x^3+2{,}2074x^2+3{,}264x+0{,}7225\right) \mathrm{\,d}x\\
&=&\pi \left(0{,}0028x^7-0{,}0406x^6+0{,}2589x^5-0{,}7757x^4+0{,}7358x^3+1{,}632x^2+0{,}7225x\right)\bigg|_{0}^3\\
&=&13{,}3293\pi\approx 42 \text{ cm}^3.
\end{eqnarray*}
Vậy thể tích của chiếc chén khoảng $42$ cm$^3$.
}
\end{ex}

\begin{ex}%[Dự án 2025 - đề cấu trúc mới, Hung Doan]%[2D4V3-5]
\immini{Một cốc có hình dạng tròn xoay và kích thước như hình vẽ, thiết diện dọc của mặt bên trong cốc (bổ dọc cốc thành $2$ phần bằng nhau) là một đường Parabol. Tính thể tích tối đa mà cốc có thể chứa được (kết quả làm tròn đến chữ số hàng đơn vị).}{
\definecolor{almond}{rgb}{0.94, 0.87, 0.8}%màu ly
\definecolor{anti-flashwhite}{rgb}{0.95, 0.95, 0.96}%màu miệng ly
\begin{tikzpicture}[line join=round, line cap=round,scale=0.5,transform shape,line width=.3mm]

\tikzset{co_vat/.pic={
\path
(1.85,4)coordinate (A)
(-1.85,4) coordinate (B)
(-3,3)coordinate (C)
(-3,-1.5)coordinate (D)

;
\draw (A)--(B) (C)--(D);
\foreach\p in {A,B,C,D}
{\draw[fill=black](\p) circle (1pt);}
\node at (-3.8,1) {$10$ cm};
\node at (0,4.5) {$8$ cm};
\draw (2,1)--(5,2.5) node[above] {Parabol};

\draw[fill=almond!50!black] (1.75,-4.93)  arc (0:360:1.75 cm and .54cm);
\draw[fill=almond] (1.75,-4.8)  arc (0:360:1.75 cm and .4cm);

\draw[fill=anti-flashwhite] (1.85,3)  arc (0:360:1.85 cm and .3cm);
\draw[fill=almond] (1.85,3)
..controls +(-95:1.3) and +(30:1.7) ..(.4,-1.6)
..controls +(-150:.1) and +(90:2) ..(.2,-4.2)
..controls +(-150:.1) and +(-30:.1) ..(-.2,-4.2)
..controls +(90:2) and +(-30:.1) ..(-.4,-1.6)
..controls +(150:1.7) and +(-85:1.3) ..(-1.85,3)
arc (-180:0:1.85 cm and .3cm);
;
\draw[fill=almond] (.2,-4.2)
..controls +(-150:.1) and +(-30:.1) ..(-.2,-4.2)
..controls +(180:.1) and +(60:.1) ..(-.4,-4.4)
..controls +(180:.3) and +(60:.1) ..(-.9,-4.8)
..controls +(-30:.4) and +(-150:.4) ..(.9,-4.8)
..controls +(120:.1) and +(0:.3) ..(.4,-4.4)
..controls +(150:.1) and +(-30:.1) ..(.2,-4.2)
;

}}

\path
(0,0)pic[scale=1]{co_vat};
\end{tikzpicture}
}
\shortans{$251$}
\loigiai{
\immini{Parabol có phương trình $y=\dfrac{5}{8}x^2\Leftrightarrow x^2=\dfrac{8}{5}y$.\\
Thể tích tối đa cốc $V=\pi \displaystyle\int \limits_0^{10} \left(\dfrac{8}{5}y\right)\cdot \mathrm{\,d}y\approx 251$.}{
\begin{tikzpicture}[scale=0.7, font=\footnotesize, line join=round, line cap=round, >=stealth]
\draw[->, line width=0.8pt](-3.5, 0)--(3.75, 0) node[below]{$x$};
\draw[->, line width=0.8pt](0,-0.5)--(0,6.5) node[left]{$y$};
\node at (0, 0) [below left]{$O$};
\fill (3,0) node[below]{$4$} circle(1pt);
\fill (-3,0) node[below]{$-4$} circle(1pt);
\fill (0,5.6125) node[below left]{$10$} circle(1pt);
\clip(-3,-0.1) rectangle(3,5.75);
\draw[ smooth,samples=300] plot(\x,{5/8*(\x)^2});
\draw[dashed] (-3,0)--(-3,5.6125)--(3,5.6125)--(3,0);
\end{tikzpicture}
}
}
\end{ex}

\begin{ex}%[Đề thi tốt nghiệp THPT 2025]%[Nguyễn Vương Hiển]%[2D4V3-5]
Để đặt được một vật trang trí trên mặt bàn, người ta thiết kế một chân đế như sau. Lấy một khối gỗ có dạng khối chóp cụt tứ giác đều với độ dài hai cạnh đáy lần lượt bằng $7{,}4$ cm và $10{,}4$ cm, bề dày của khối gỗ bằng $1{,}5$ cm. Sau đó khoét bỏ đi một phần của khối gỗ sao cho phần đó có dạng vật thể $H$, ở đó $H$ nhận được bằng cách cắt khối cầu bán kính $5{,}8$ cm bởi một mặt phẳng cắt mà mặt cắt là hình tròn bán kính $3{,}5$ cm (xem hình dưới).\\
\centerline{
\begin{tikzpicture}[line join=round, line cap=round,>=stealth,thick]
% Hình trái
\coordinate (A) at (0,0);
\coordinate (B) at (-1.5,-1);
\coordinate (D) at (4,0);
\coordinate (C) at ($(B)+(D)-(A)$);
\coordinate (O) at ($(A)!0.5!(C)$);
\coordinate (I) at ($(O)+(0,3.5)$);
\coordinate (A') at ([scale around={0.71:(I)}] A);
\coordinate (B') at ([scale around={0.71:(I)}] B);
\coordinate (C') at ([scale around={0.71:(I)}] C);
\coordinate (D') at ([scale around={0.71:(I)}] D);
\coordinate (O') at ($(A')!0.5!(C')$);
\draw(B)--(C)--(D)--(D')--(C')--(B')--(A')--(D') (C)--(C') (B)--(B');
\draw[dashed,thin](A')--(A)--(B) (A)--(D);
\draw[dashed,<->](O)--(O')node[right,pos=0.3]{\tiny $1{,}5$ cm};
\coordinate (vt) at (0,0.2);
\coordinate (A1) at ($(A')+(vt)$);
\coordinate (D1) at ($(D')+(vt)$);
\draw[|<->|](A1)--(D1)node[above,sloped,pos=0.5]{\tiny $7{,}4$ cm};
\coordinate (vt1) at (0,0.1);
\coordinate (B1) at ($(B)-(vt)$);
\coordinate (C1) at ($(C)-(vt)$);
\draw[|<->|](B1)--(C1)node[below,sloped,pos=0.5]{\tiny $10{,}4$ cm};
\end{tikzpicture}\qquad
\begin{tikzpicture}[line join=round, line cap=round,>=stealth,thick,scale=0.8]
% Hình giữa
\def\bk{3}
\def\bkba{0.15}
\def\bkbb{0.8}
\def\ga{-50}
\def\gb{180-\ga}
\def\gc{0-\gb}
\coordinate (O) at (0,0);
\coordinate (M) at ($(O) + (\ga:\bk)$);
\coordinate (N) at ($(O) + (\gb:\bk)$);
\coordinate (I) at ($(M)!0.5!(N)$);
\draw[name path=c1,dashed] (M) arc (\ga:\gb:\bk);
\path let \p1=(I),\p2=(M),\n1={scalar(veclen(\x2-\x1,\y2-\y1)/1cm)} in \pgfextra{\xdef\bke{\n1}};
\fill[gray](M) arc (\ga:\gc:\bk)--(N) arc (180:0:\bke cm and \bkba cm)--cycle;
\draw[name path=c1] (M) arc (\ga:\gc:\bk);
\draw[name path=ne1,dashed] (M) arc (0:180:\bke cm and \bkba cm) (M)--(N);
\draw[name path=ne1] (M) arc (0:-180:\bke cm and \bkba cm);
\draw[name path=ela,dashed] (O) ellipse (\bk cm and \bkbb cm) ($(O)-(\bk,0)$)--($(O)+(\bk,0)$);
\coordinate (P) at ($(O) + (250:0.9*\bk)$);
\coordinate (Q) at ($(O) + (235:1.1*\bk)$);
\path(Q) node [left] {\Large\bf $H$};
\draw[->](Q)--(P);
\end{tikzpicture}\qquad
\begin{tikzpicture}[line join=round, line cap=round,>=stealth,thick,rotate=90,scale=0.8]
% Hình phải
\def\bk{3}
\def\bkba{0.15}
\def\bkbb{0.8}
\def\ga{-50}
\def\gb{180-\ga}
\def\gc{0-\gb}
\coordinate (O) at (0,0);
\coordinate (M) at ($(O) + (\ga:\bk)$);
\coordinate (N) at ($(O) + (\gb:\bk)$);
\coordinate (I) at ($(M)!0.5!(N)$);
\coordinate (J) at ($(O)+(M)-(I)$);
\draw(M) arc (\ga:\gb:\bk);
\path let \p1=(I),\p2=(M),\n1={scalar(veclen(\x2-\x1,\y2-\y1)/1cm)} in \pgfextra{\xdef\bke{\n1}};
\fill[gray](M) arc (\ga:-90:\bk)--(I)--cycle;
\draw(M) arc (\ga:\gc:\bk);
\draw(M) arc (0:180:\bke cm and \bkba cm);
\draw[dashed](M) arc (0:-180:\bke cm and \bkba cm);
\coordinate (P) at ($(O) + (250:0.9*\bk)$);
\coordinate (Q) at ($(O) + (235:1.1*\bk)$);
\path(Q) node [below] {\Large\bf $H$};
\coordinate (A) at ($(O) + (180:\bk)$);
\coordinate (B) at ($(O) + (0:\bk)$);
\coordinate (C) at ($(O) + (-90:\bk)$);
\coordinate (D) at ($(O) + (90:\bk)$);
\draw(B) arc (0:180:\bk cm and \bkbb cm);
\draw[dashed](B) arc (0:-180:\bk cm and \bkbb cm);
\draw[dashed](J)--(M)--(N) (C)--(D) (A)--(B);
\draw[->](Q)--(P);
\draw[->](C)--($(C)-(0,0.5)$)node[below]{$x$};
\draw[->](B)--($(B)+(0.5,0)$)node[right]{$y$};
\draw(A)--($(A)-(0.5,0)$) (D)--($(D)+(0,0.5)$);
\path(O) node[below left]{$O$};
\path(J) node[left=-0.05]{\tiny $3{,}5$};
\path($(B)+(0.15,0)$)node[left]{\tiny $5{,}8$};
\end{tikzpicture}
}\\
Thể tích của khối chân đế bằng bao nhiêu centimét khối (không làm tròn kết quả các phép tính trung gian, chỉ làm tròn kết quả cuối cùng đến hàng phần mười)?
\shortans[oly]{$96{,}5$}
\loigiai{
Gọi $V_2$ là thể tích khối chóp cụt, $V_1$ là thể tích khối chỏm cầu.\\
Ta có
\begin{itemize}
\item 	Diện tích đáy lớn là $10{,}4^2 = 108{,}16$ cm$^2$.
\item  Diện tích đáy nhỏ là $7{,}4^2 = 54{,}76$ cm$^2$.
\item  	Chiều cao là $1{,}5$ cm.
\end{itemize}
Thể tích của khối chóp cụt là $V_2 = \dfrac{1{,}5}{3} \left(108{,}16 + 54{,}76 + \sqrt{108{,}16 \cdot 54{,}76} \right) =119{,}94$ cm$^3$.\\
Phương trình đường tròn $(C)$ tâm $O$, bán kính $R = 5{,}8$ là
\[(C)\colon x^2 + y^2 = 5{,}8^2 \Rightarrow y = \sqrt{5{,}8^2 - x^2} \quad \text{(nửa đường tròn trên).}\]
Hoành độ giao điểm của hình tròn bán kính $3{,}5$ trong hình và  $(C)$ là
\[\sqrt{5{,}8^2 - x^2} = 3{,}5 \Rightarrow x^2 =\dfrac{2139}{100} \Rightarrow x = \dfrac{\sqrt{2139}}{10}.\]
Thể tích của khối chỏm cầu là
\[V_1 = \pi\displaystyle \int\limits_{\tfrac{\sqrt{2139}}{10}}^{5{,}8} \left(5{,}8^2 - x^2 \right)\mathrm{\,d}x.\]
Vậy thể tích của khối chân đế là $V= V_2 - V_1 \approx 96{,}5$ cm$^3$.
}
\end{ex}

\begin{ex}%[12EX-PT-ĐỀ SỐ 2-2025, Nguyễn Anh Tuấn]%[2D4V3-5]
Một đồ lưu niệm bằng thuỷ tinh có chiều cao bằng $14$ cm, được thiết kế gồm hai phần, phần dưới là một khối lập phương cạnh bằng $8$ cm và phần trên là một phần của khối cầu có đường kính bằng $8$ cm (tham khảo hình vẽ). Thể tích của đồ lưu niệm đó là bao nhiêu (làm tròn kết quả đến hàng đơn vị của cm$^3$)?
\shortans{738}
\begin{center}
\begin{tikzpicture}[scale=.7,font=\footnotesize, line join=round, line cap=round, >=stealth]
\coordinate (A) at (0,0);
\coordinate (B) at (-2,-2);
\coordinate (D) at (5,0);
\coordinate (C) at ($(B)+(D)-(A)$);
\coordinate (A') at (0,5);
\coordinate (B') at ($(B)+(A')-(A)$);
\coordinate (C') at ($(C)+(A')-(A)$);
\coordinate (D') at ($(D)+(A')-(A)$);
\coordinate (O) at ($(A')!0.5!(C')$);
\coordinate (x) at ($(A')!0.3!(B')$);
\coordinate (y) at ($(D')!0.22!(A')$);
\coordinate (I) at ($(O)+(0,2)$);
\coordinate (K) at ($(I)+(2.53,0)$);
\coordinate (M) at ($(O)+(-1.8,0)$);
\coordinate (N) at ($(O)+(1.8,0)$);
\fill[gray!20] (B')--(B)--(C)--(C') (C')--(D')--(D)--(C) (A')--(B')--(C')--(D');
\fill[gray, opacity=0.7] (N) arc (-45:225:2.55cm);
\fill[gray, opacity=0.7] (N) arc (0:-180:1.8cm and 0.7cm);
\draw (B')--(B)--(C)--(D)--(D')--(y) (x)--(B')--(C')--(C) (C')--(D');
\draw (N) arc (0:-180:1.8cm and 0.7cm);
\draw [dashed] (N) arc (0:180:1.8cm and 0.6cm);
\draw (K) arc (0:-180:2.53cm and 0.6cm);
\begin{scope}[shift={(10,0)}] % Shift the second figure to the right
\coordinate (A) at (0,0);
\coordinate (B) at (-2,-2);
\coordinate (D) at (5,0);
\coordinate (C) at ($(B)+(D)-(A)$);
\coordinate (A') at (0,5);
\coordinate (B') at ($(B)+(A')-(A)$);
\coordinate (C') at ($(C)+(A')-(A)$);
\coordinate (D') at ($(D)+(A')-(A)$);
\coordinate (O) at ($(A')!0.5!(C')$);
\coordinate (x) at ($(A')!0.3!(B')$);
\coordinate (y) at ($(D')!0.22!(A')$);
\coordinate (I) at ($(O)+(0,2)$);
\coordinate (K) at ($(I)+(2.53,0)$);
\coordinate (M) at ($(O)+(-1.8,0)$);
\coordinate (N) at ($(O)+(1.8,0)$);
\coordinate (t) at ($(K)+(2,2)$);
\coordinate (z) at ($(K)+(2,-7)$);
\draw (N) arc (0:-180:1.8cm and 0.7cm);
\draw [dashed] (N) arc (0:180:1.8cm and 0.6cm);
\draw (K) arc (0:-180:2.53cm and 0.6cm);
\draw [dashed] (K) arc (0:180:2.53cm and 0.7cm);
\draw (N) arc (-45:225:2.55cm);
\draw[dashed] (B)--(A)node[sloped, midway, above]{$8\text{ cm}$}--(D) (A)--(A') (x)--(A')--(y);
\draw (B')--(B)--(C)node[sloped, midway, below]{$8\text{ cm}$}--(D)--(D')--(y) (x)--(B')--(C')--(C) (C')--(D');
\draw[<->] (t)--(z)node[sloped, midway, above]{$14\text{ cm}$};
\end{scope}
\end{tikzpicture}
\end{center}
\loigiai{
\begin{itemize}
\item Thể tích khối lập phương $V_1=8^3=512$ cm$^3$.
\item Chỏm cầu (bên trong khối lập phương) bán kính $R=4$ cm và chiều cao  $h=8-6=2$ cm.\\
$V_2=\pi\cdot 2^2\cdot \left(4-\dfrac{2}{3}\right)$ cm$^3$.
\item Vậy thể tích đồ lưu niệm là $V=V_1+\dfrac{4}{3}\cdot\pi\cdot 4^3-V_2\approx 738$ cm$^3$.
\end{itemize}
}
\end{ex}

\begin{ex}%[2D4V3-5]%[Dự án Tex hoá đề Tốt nghiệp 2025 + Phát triển 05 đề]%[Viết Tường]
Để đặt được một vật trang trí trên mặt bàn, người ta thiết kế một chân đế như sau. Lấy một khối gỗ có dạng khối chóp cụt tứ giác đều với độ dài hai cạnh đáy lần lượt bằng $7$ cm và $10$ cm, bề dày của khối gỗ bằng $1$ cm. Sau đó khoét bỏ đi một phần của khối gỗ sao cho phần đó có dạng vật thể $H$, ở đó $H$ nhận được bằng cách cắt khối cầu bán kính $5$ cm bởi một mặt phẳng cắt mà mặt cắt là hình tròn bán kính $3$ cm (xem hình dưới).\\
\centerline{
\begin{tikzpicture}[line join=round, line cap=round,>=stealth,thick]
% Hình trái
\coordinate (A) at (0,0);
\coordinate (B) at (-1.5,-1);
\coordinate (D) at (4,0);
\coordinate (C) at ($(B)+(D)-(A)$);
\coordinate (O) at ($(A)!0.5!(C)$);
\coordinate (I) at ($(O)+(0,3.5)$);
\coordinate (A') at ([scale around={0.71:(I)}] A);
\coordinate (B') at ([scale around={0.71:(I)}] B);
\coordinate (C') at ([scale around={0.71:(I)}] C);
\coordinate (D') at ([scale around={0.71:(I)}] D);
\coordinate (O') at ($(A')!0.5!(C')$);
\draw(B)--(C)--(D)--(D')--(C')--(B')--(A')--(D') (C)--(C') (B)--(B');
\draw[dashed,thin](A')--(A)--(B) (A)--(D);
\draw[dashed,<->](O)--(O')node[right,pos=0.3]{\tiny $1$ cm};
\coordinate (vt) at (0,0.2);
\coordinate (A1) at ($(A')+(vt)$);
\coordinate (D1) at ($(D')+(vt)$);
\draw[|<->|](A1)--(D1)node[above,sloped,pos=0.5]{\tiny $7$ cm};
\coordinate (vt1) at (0,0.1);
\coordinate (B1) at ($(B)-(vt)$);
\coordinate (C1) at ($(C)-(vt)$);
\draw[|<->|](B1)--(C1)node[below,sloped,pos=0.5]{\tiny $10$ cm};
\end{tikzpicture}\qquad
\begin{tikzpicture}[line join=round, line cap=round,>=stealth,thick,scale=0.8]
% Hình giữa
\def\bk{3}
\def\bkba{0.15}
\def\bkbb{0.8}
\def\ga{-50}
\def\gb{180-\ga}
\def\gc{0-\gb}
\coordinate (O) at (0,0);
\coordinate (M) at ($(O) + (\ga:\bk)$);
\coordinate (N) at ($(O) + (\gb:\bk)$);
\coordinate (I) at ($(M)!0.5!(N)$);
\draw[name path=c1,dashed] (M) arc (\ga:\gb:\bk);
\path let \p1=(I),\p2=(M),\n1={scalar(veclen(\x2-\x1,\y2-\y1)/1cm)} in \pgfextra{\xdef\bke{\n1}};
\fill[gray](M) arc (\ga:\gc:\bk)--(N) arc (180:0:\bke cm and \bkba cm)--cycle;
\draw[name path=c1] (M) arc (\ga:\gc:\bk);
\draw[name path=ne1,dashed] (M) arc (0:180:\bke cm and \bkba cm) (M)--(N);
\draw[name path=ne1] (M) arc (0:-180:\bke cm and \bkba cm);
\draw[name path=ela,dashed] (O) ellipse (\bk cm and \bkbb cm) ($(O)-(\bk,0)$)--($(O)+(\bk,0)$);
\coordinate (P) at ($(O) + (250:0.9*\bk)$);
\coordinate (Q) at ($(O) + (235:1.1*\bk)$);
\path(Q) node [left] {\Large\bf $H$};
\draw[->](Q)--(P);
\path (0,0) node{\hypersetup{hidelinks}\href{yLRPIuU}{ }};
\end{tikzpicture}\qquad
\begin{tikzpicture}[line join=round, line cap=round,>=stealth,thick,rotate=90,scale=0.8]
% Hình phải
\def\bk{3}
\def\bkba{0.15}
\def\bkbb{0.8}
\def\ga{-50}
\def\gb{180-\ga}
\def\gc{0-\gb}
\coordinate (O) at (0,0);
\coordinate (M) at ($(O) + (\ga:\bk)$);
\coordinate (N) at ($(O) + (\gb:\bk)$);
\coordinate (I) at ($(M)!0.5!(N)$);
\coordinate (J) at ($(O)+(M)-(I)$);
\draw(M) arc (\ga:\gb:\bk);
\path let \p1=(I),\p2=(M),\n1={scalar(veclen(\x2-\x1,\y2-\y1)/1cm)} in \pgfextra{\xdef\bke{\n1}};
\fill[gray](M) arc (\ga:-90:\bk)--(I)--cycle;
\draw(M) arc (\ga:\gc:\bk);
\draw(M) arc (0:180:\bke cm and \bkba cm);
\draw[dashed](M) arc (0:-180:\bke cm and \bkba cm);
\coordinate (P) at ($(O) + (250:0.9*\bk)$);
\coordinate (Q) at ($(O) + (235:1.1*\bk)$);
\path(Q) node [below] {\Large\bf $H$};
\coordinate (A) at ($(O) + (180:\bk)$);
\coordinate (B) at ($(O) + (0:\bk)$);
\coordinate (C) at ($(O) + (-90:\bk)$);
\coordinate (D) at ($(O) + (90:\bk)$);
\draw(B) arc (0:180:\bk cm and \bkbb cm);
\draw[dashed](B) arc (0:-180:\bk cm and \bkbb cm);
\draw[dashed](J)--(M)--(N) (C)--(D) (A)--(B);
\draw[->](Q)--(P);
\draw[->](C)--($(C)-(0,0.5)$)node[below]{$x$};
\draw[->](B)--($(B)+(0.5,0)$)node[right]{$y$};
\draw(A)--($(A)-(0.5,0)$) (D)--($(D)+(0,0.5)$);
\path(O) node[below left]{$O$};
\path(J) node[left=-0.05]{\tiny $3$};
\path($(B)+(0.15,0)$)node[left]{\tiny $5$};
\end{tikzpicture}
}\\
Thể tích của khối chân đế bằng bao nhiêu centimét khối (không làm tròn kết quả các phép tính trung gian, chỉ làm tròn kết quả cuối cùng đến hàng phần mười)?
\shortans{58,3}
\loigiai
{
Gọi $V_2$ là thể tích khối chóp cụt, $V_1$ là thể tích khối chỏm cầu.\\
Ta có
\begin{itemize}
\item 	Diện tích đáy lớn là $10^2 = 100$ cm$^2$.
\item  Diện tích đáy nhỏ là $7^2 = 49$ cm$^2$.
\item  	Chiều cao là $1$ cm.
\end{itemize}
Thể tích của khối chóp cụt là $V_2 = \dfrac{1}{3} \left(100 + 49 + \sqrt{100 \cdot 49} \right) = 73$ cm$^3$.\\
Phương trình đường tròn $(C)$ tâm $O$, bán kính $R = 5$ là
\[(C)\colon x^2 + y^2 = 5^2 \Rightarrow y = \sqrt{5^2 - x^2} \quad \text{(nửa đường tròn trên).}\]
Hoành độ giao điểm của hình tròn bán kính $3$ trong hình và  $(C)$ là
\[\sqrt{5^2 - x^2} = 3 \Rightarrow x=4.\]
Thể tích của khối chỏm cầu là
\[V_1 = \pi\displaystyle \int\limits_{4}^{5} \left(5^2 - x^2 \right)\mathrm{\,d}x.\]
Vậy thể tích của khối chân đế là $V= V_2 - V_1 \approx 58{,}3$ cm$^3$.
}
\end{ex}

\begin{ex}%[Dự án Tex hoá đề Tốt nghiệp 2025 + Phát triển 05 đề]%[Đình Nguyên]%[2D4V3-5]
Một xưởng mộc muốn chế tạo một chân nến bằng gỗ theo thiết kế đặc biệt. Người ta bắt đầu với một khối gỗ có dạng hình chóp cụt đều, trong đó hai đáy là hình vuông có cạnh lần lượt là $6$ cm và $9$ cm, chiều cao khối gỗ là $2$ cm.
Để đặt được cây nến tròn, người thợ khoét bỏ đi một phần gỗ có hình dạng như phần vật thể $H$ trong hình vẽ dưới. Phần bị khoét chính là một phần của hình cầu bán kính $R = 4{,}2$ cm, được cắt bởi một mặt phẳng song song với mặt đáy tạo thành mặt cắt là hình tròn bán kính $r = 2{,}8$ cm.
\begin{center}
\begin{tikzpicture}[line join=round, line cap=round,>=stealth,font=\footnotesize]
% Hình trái
\coordinate (A) at (0,0);
\coordinate (B) at (-1.5,-1);
\coordinate (D) at (4,0);
\coordinate (C) at ($(B)+(D)-(A)$);
\coordinate (O) at ($(A)!0.5!(C)$);
\coordinate (I) at ($(O)+(0,3.5)$);
\coordinate (A') at ([scale around={0.71:(I)}] A);
\coordinate (B') at ([scale around={0.71:(I)}] B);
\coordinate (C') at ([scale around={0.71:(I)}] C);
\coordinate (D') at ([scale around={0.71:(I)}] D);
\coordinate (O') at ($(A')!0.5!(C')$);
\draw(B)--(C)--(D)--(D')--(C')--(B')--(A')--(D') (C)--(C') (B)--(B');
\draw[dashed,thin](A')--(A)--(B) (A)--(D);
\draw[dashed,<->](O)--(O')node[right,pos=0.3]{\tiny $2$ cm};
\coordinate (vt) at (0,0.2);
\coordinate (A1) at ($(A')+(vt)$);
\coordinate (D1) at ($(D')+(vt)$);
\draw[|<->|](A1)--(D1)node[above,sloped,pos=0.5]{\tiny $6$ cm};
\coordinate (vt1) at (0,0.1);
\coordinate (B1) at ($(B)-(vt)$);
\coordinate (C1) at ($(C)-(vt)$);
\draw[|<->|](B1)--(C1)node[below,sloped,pos=0.5]{\tiny $9$ cm};
\end{tikzpicture}\qquad
\begin{tikzpicture}[line join=round, line cap=round,>=stealth,thick,scale=0.8,font=\footnotesize]
% Hình giữa
\def\bk{3}
\def\bkba{0.15}
\def\bkbb{0.8}
\def\ga{-50}
\def\gb{180-\ga}
\def\gc{0-\gb}
\coordinate (O) at (0,0);
\coordinate (M) at ($(O) + (\ga:\bk)$);
\coordinate (N) at ($(O) + (\gb:\bk)$);
\coordinate (I) at ($(M)!0.5!(N)$);
\draw[name path=c1,dashed] (M) arc (\ga:\gb:\bk);
\path let \p1=(I),\p2=(M),\n1={scalar(veclen(\x2-\x1,\y2-\y1)/1cm)} in \pgfextra{\xdef\bke{\n1}};
\fill[gray](M) arc (\ga:\gc:\bk)--(N) arc (180:0:\bke cm and \bkba cm)--cycle;
\draw[name path=c1] (M) arc (\ga:\gc:\bk);
\draw[name path=ne1,dashed] (M) arc (0:180:\bke cm and \bkba cm) (M)--(N);
\draw[name path=ne1] (M) arc (0:-180:\bke cm and \bkba cm);
\draw[name path=ela,dashed] (O) ellipse (\bk cm and \bkbb cm) ($(O)-(\bk,0)$)--($(O)+(\bk,0)$);
\coordinate (P) at ($(O) + (250:0.9*\bk)$);
\coordinate (Q) at ($(O) + (235:1.1*\bk)$);
\path(Q) node [left] {\Large\bf $H$};
\draw[->](Q)--(P);
\end{tikzpicture}\qquad
\begin{tikzpicture}[line join=round, line cap=round,>=stealth,thick,rotate=90,scale=0.8,font=\footnotesize]
% Hình phải
\def\bk{3}
\def\bkba{0.15}
\def\bkbb{0.8}
\def\ga{-50}
\def\gb{180-\ga}
\def\gc{0-\gb}
\coordinate (O) at (0,0);
\coordinate (M) at ($(O) + (\ga:\bk)$);
\coordinate (N) at ($(O) + (\gb:\bk)$);
\coordinate (I) at ($(M)!0.5!(N)$);
\coordinate (J) at ($(O)+(M)-(I)$);
\draw(M) arc (\ga:\gb:\bk);
\path let \p1=(I),\p2=(M),\n1={scalar(veclen(\x2-\x1,\y2-\y1)/1cm)} in \pgfextra{\xdef\bke{\n1}};
\fill[gray](M) arc (\ga:-90:\bk)--(I)--cycle;
\draw(M) arc (\ga:\gc:\bk);
\draw(M) arc (0:180:\bke cm and \bkba cm);
\draw[dashed](M) arc (0:-180:\bke cm and \bkba cm);
\coordinate (P) at ($(O) + (250:0.9*\bk)$);
\coordinate (Q) at ($(O) + (235:1.1*\bk)$);
\path(Q) node [below] {\Large\bf $H$};
\coordinate (A) at ($(O) + (180:\bk)$);
\coordinate (B) at ($(O) + (0:\bk)$);
\coordinate (C) at ($(O) + (-90:\bk)$);
\coordinate (D) at ($(O) + (90:\bk)$);
\draw(B) arc (0:180:\bk cm and \bkbb cm);
\draw[dashed](B) arc (0:-180:\bk cm and \bkbb cm);
\draw[dashed](J)--(M)--(N) (C)--(D) (A)--(B);
\draw[->](Q)--(P);
\draw[->](C)--($(C)-(0,0.5)$)node[below]{$x$};
\draw[->](B)--($(B)+(0.5,0)$)node[right]{$y$};
\draw(A)--($(A)-(0.5,0)$) (D)--($(D)+(0,0.5)$);
\path(O) node[below left]{$O$};
\path(J) node[left=-0.05]{\tiny $2{,}8$};
\path($(B)+(0.15,0)$)node[left]{\tiny $4{,}2$};
\end{tikzpicture}
\end{center}
Thể tích của khối chân đế bằng bao nhiêu centimét khối (không làm tròn kết quả các phép tính trung gian, chỉ làm tròn kết quả cuối cùng đến hàng đơn vị)?
\shortans[oly]{$100$}
\loigiai{
Gọi $V_2$ là thể tích khối chóp cụt, $V_1$ là thể tích khối chỏm cầu.\\
Ta có
\begin{itemize}
\item 	Diện tích đáy lớn là $9^2 = 81$ cm$^2$.
\item  Diện tích đáy nhỏ là $6^2=36$ cm$^2$.
\item  	Chiều cao là $2$ cm.
\end{itemize}
Thể tích của khối chóp cụt là $V_2 = \dfrac{1}{3}\cdot 2\cdot \left(81+36+\sqrt{81\cdot 36} \right) =114$ cm$^3$.\\
Phương trình đường tròn $(C)$ tâm $O$, bán kính $R = 4{,}2$ là
\[(C)\colon x^2 + y^2 = 4{,}2^2 \Rightarrow y = \sqrt{4{,}2^2 - x^2} \quad \text{(nửa đường tròn trên).}\]
Hoành độ giao điểm của hình tròn bán kính $3{,}5$ trong hình và  $(C)$ là
\[\sqrt{4{,}2^2 - x^2} = 2{,}8 \Rightarrow x^2 =\dfrac{49}{5} \Rightarrow x = \dfrac{7\sqrt5}{5}.\]
Thể tích của khối chỏm cầu là
\[V_1 = \pi\displaystyle \int\limits_{\tfrac{7\sqrt5}{5}}^{4{,}2} \left(4{,}2^2 - x^2 \right)\mathrm{\,d}x.\]
Vậy thể tích của khối chân đế là $V= V_2 - V_1 \approx 100$ cm$^3$.
}
\end{ex}

\begin{ex}%[2-GHK2-1-ChuyenViThanh-HauGiang-25]%[12-EX-3-2026, Tran Quoc]%[2D4V3-5]
Để chuẩn bị cho buổi dã ngoại, nhóm du lịch dự định dựng một cái lều trại có dạng như hình vẽ sau:
\begin{center}
\begin{tikzpicture}[font=\footnotesize, line join=round, line cap=round, >=stealth, scale=0.8]
\foreach \x/\y/\pos in {-2/0/A, 2/0/B, 0/3/S} \path ($(\x,\y)$) coordinate (\pos) ($(\x,\y)+(4.5,0.8)$) coordinate (\pos');
\draw[smooth,samples=100,domain=-2:2] plot(\x,{-0.75*(\x)^2+3});
\draw[dashed, smooth, samples=100, domain=-2:0] plot({\x+4.5},{-0.75*(\x)^2+3.8});
\draw[smooth, samples=100, domain=0:2] plot({\x+4.5},{-0.75*(\x)^2+3.8});
\draw[dashed] (A) -- (A') -- (B');
\draw (A) -- (B) -- (B') (S)--(S');
\end{tikzpicture}
\end{center}
Biết rằng mặt trước và mặt sau của trại là hai parabol bằng nhau, nằm trên hai mặt phẳng song song với nhau và cùng vuông góc với mặt nền. Nền của lều trại là một hình chữ nhật có kích thước chiều rộng là $4$\,m (lối vào lều), chiều dài là $6$\,m, đỉnh parabol cách nền $3$\,m. Tính thể tích phần không gian bên trong lều trại.
\par\shortans{48}
\loigiai{
\begin{center}
\begin{minipage}[]{0.68\textwidth}
\begin{center}
\begin{tikzpicture}[font=\footnotesize, line join=round, line cap=round, >=stealth, scale=0.8]
\foreach \x/\y/\pos in {-2/0/A, 2/0/B, 0/3/S} \path ($(\x,\y)$) coordinate (\pos) ($(\x,\y)+(4.5,0.8)$) coordinate (\pos');
\path
($(A)!1/2!(B)$) coordinate (O)
($(A')!1/2!(B')$) coordinate (O')
($(O)!1.5!(S)$) coordinate (y)
($(O)!1.42!(O')$) coordinate (x')
($(O)!1.8!(O')$) coordinate (x);
\draw[smooth,samples=100,domain=-2:2] plot(\x,{-0.75*(\x)^2+3});
\draw[dashed, smooth, samples=100, domain=-2:0] plot({\x+4.5},{-0.75*(\x)^2+3.8});
\draw[smooth, samples=100, domain=0:2] plot({\x+4.5},{-0.75*(\x)^2+3.8});
\draw[dashed] (A) -- (A') -- (B') (O)--(x');
\draw (A) -- (B) -- (B') (S)--(S');
\draw[->] (x')--(x) node[shift=(-110:0.2)] {$x$};
\draw[->] (O)--(y) node[shift=(-150:0.2)] {$y$};
\fill (O) circle(1pt) node[shift=(-90:0.2)]{$O$};
\fill (S) circle(1pt) node[shift=(150:0.25)]{$3$};
\fill (O') circle(1pt) node[shift=(90:0.2)]{$6$};
\end{tikzpicture}\\
Hình $1$.
\end{center}
\end{minipage}
\begin{minipage}[]{0.28\textwidth}
\begin{center}
\begin{tikzpicture}[font=\footnotesize, line join=round, line cap=round, >=stealth, scale=0.8]
\def \xmin{-3}\def \xmax{3}\def \ymin{-1}\def \ymax{4}
\draw[->] (\xmin,0)--(\xmax,0) node[shift=(-110:0.2)] {$x$};
\draw[->] (0,\ymin)--(0,\ymax) node[shift=(-150:0.2)] {$y$};
\fill (0,0) circle(1pt) node[shift=(-135:0.25)]{$O$};
\fill (-2,0) circle(1pt) node[shift=(-90:0.2)]{$-2$};
\fill (2,0) circle(1pt) node[shift=(-90:0.2)]{$2$};
\fill (0,3) circle(1pt) node[shift=(150:0.25)]{$3$};
\fill[pattern=north west lines,opacity=.8] plot[domain=-2:2](\x,{-0.75*(\x)^2+3})--cycle;
\draw[smooth,samples=100,domain=-2:2] plot(\x,{-0.75*(\x)^2+3});
\node at (1.8,2) {$(P)$};
\end{tikzpicture}\\
Hình $2$.
\end{center}
\end{minipage}
\end{center}
Chọn hệ tục $Oxy$ như hình $1$.\\
Khi đó, mỗi mặt cắt vuông góc với trục $Ox$ tại điểm có hoành độ $x$ sẽ cắt phần không gian bên trong lều trại là miền gạch sọc giới hạn bởi parabol $(P)$ và trục hoành $Ox$ như hình $2$.\\
Trong hình $2$, parabol $(P)\colon f(x)=ax^2+bx+c$ đi qua $3$ điểm $A(-2;0)$, $B(0;3)$ và $C(2;0)$ nên ta có hệ phương trình
\[ \heva{&f(-2)=0\\&f(0)=3\\&f(2)=0} \Leftrightarrow \heva{&4a-2b+c=0\\&0a+0b+c=3\\&4a+2b+c=0}\Leftrightarrow \heva{&a=-\dfrac{3}{4}\\&b=0\\&c=3.} \]
Do đó $(P)\colon y=-\dfrac{3}{4}x^2+3$.\\
Suy ra diện tích mặt cắt là
\[S = \int\limits_{-2}^{2} \left(-\dfrac{3}{4}x^2 + 3\right) \mathrm{\, d}x = \left(-\dfrac{x^3}{4} + 3x\right) \bigg|_{-2}^{2} = 8\,\left(\text{m}^2\right).\]
Vậy thể tích phần không gian bên trong lều trại là
\[V=\displaystyle \int\limits_{0}^{6} S \mathrm{\, d} x =6S=6\cdot 8=48\,\left(\text{m}^3\right).\]
}
\end{ex}

\begin{ex}%[12-EX-3-2025, Quan Ón]%[2D4V3-5]
Một khuôn viên dạng nửa hình tròn, trên đó người thiết kế phần để trồng hoa Cẩm Tú Cầu (\textit{phần được gạch chéo trên hình vẽ})
\begin{center}
\begin{tikzpicture}[scale=0.8, font=\footnotesize, line join=round, line cap=round, >=stealth]
\draw[pattern=north west lines, pattern color=gray!30] (-2,4) plot[domain=-2:2](\x, {(\x)^2})--(2,4) arc (65:105:4.78cm)--cycle;
\draw [domain=-2:2, samples=100] %
plot (\x, {(\x)^2});
\draw (4.47,0) arc (0:180:4.47cm)--(-4.47,0)--(4.47,0)--cycle;

\end{tikzpicture}
\end{center}
Biết rằng phần gạch chéo là hình phẳng giới hạn bởi parabol $y=x^2$ và nửa trên của đường tròn $y=\sqrt{20-x^2}$. Tính số tiền (nghìn đồng) tối thiểu để trồng xong vườn hoa Cẩm Tú Cầu biết rằng để trồng mỗi $m^2$ hoa cần ít nhất là $250$ nghìn đồng. (\textit{Số tiền được làm tròn đến hàng đơn vị})
\shortans[]{$2985$}
\loigiai{
Phương trình hoành độ giao điểm của $y=x^2$ và $y=\sqrt{20-x^2}$ là
\[x^2=\sqrt{20-x^2}\Leftrightarrow x^4+x^2-20=0\Leftrightarrow \hoac{&x^2=4\text{ (nhận)}\\&x^2=-5\text{ (loại)}}\Rightarrow x=\pm 2.\]
Diện tích phần gạch chéo là
\[S=\int\limits_{-2}^{2}\left(\sqrt{20-x^2}-x^2\right)\mathrm{\,d}x\approx 11{,}9396.\]
Số tiền tối thiểu để trồng xong vườn hoa Cẩm Tú Cầu là
\[250\cdot 11{,}9396 \approx 2985\text{ (nghìn đồng).}\]
}
\end{ex}

\begin{ex}%[12-EX3-2026. Trần Hoà]%[2D4V3-5]
\immini{Một thùng chứa rượu bằng gỗ là một hình tròn xoay như hình bên có hai mặt đáy là hai hình tròn bằng nhau, khoảng cách giữa hai đáy là $6$ dm. Đường cong mặt bên trong của thùng là một phần của đường elip có độ dài trục lớn bằng $10$ dm, trục bé bằng $8$ dm. Hỏi thùng gỗ này đựng được tối đa bao nhiêu lít rượu? (Làm tròn kết quả đến hàng đơn vị).
\shortans[]{$265$}}{\begin{tikzpicture}[scale=1, font=\footnotesize, line join=round, line cap=round, >=stealth, yscale=0.5, xscale=0.4]
\def\bl{8}
\def\bn{3}
\pgfmathsetmacro{\a}{\bl*cos(60)}
\pgfmathsetmacro{\b}{\bn*sin(60)}
\pgfmathsetmacro{\c}{\bl*cos(65)}
\pgfmathsetmacro{\d}{\bn*sin(65)}
\pgfmathsetmacro{\e}{\bl*cos(70)}
\pgfmathsetmacro{\f}{\bn*sin(70)}
\begin{scope}
\draw (-\a,\b) arc (120:60:{\bl} and {\bn});
\draw (-\a,-\b) arc (240:300: {\bl} and {\bn});
\draw[fill=brown!90!black] (\a,0) ellipse ({0.4*\b} and {\b});
\draw[fill=brown!60] (\a,0) ellipse ({0.32*\b} and {0.8*\b});
\draw[fill=brown!75!black,shift={(-0.1,0)}] (\a,0) ellipse ({0.28*\b} and {0.7*\b});
\draw (-\a,\b) arc (90:270:{0.4*\b} and {\b});
\fill[red!0.5!brown] (-\a,-\b) arc (270:90:{0.4*\b} and {\b}) --(-\a,\b) arc (120:60:{\bl} and {\bn})-- (\a,\b) arc (90:270:{0.4*\b} and {\b}) -- (\a,-\b) arc (300:240:{\bl} and {\bn});
\draw[brown!60!black] ({0.4*\b*cos(135)-\a},{\b*sin(135)}) arc (120:60:8 and 2.5)
({0.4*\b*cos(225)-\a},{\b*sin(225)}) arc (240:300:8 and 2.5)
({-0.4*\b-\a},0)--({-0.4*\b+\a},0);
\draw[line width=1pt,gray] (-\c,\d) arc (90:270:{0.4*\d} and {\d});
\draw[line width=1pt,gray] (-\e,\f) arc (90:270:{0.4*\f} and {\f});
\draw[line width=1pt,gray] (\c,\d) arc (90:270:{0.4*\d} and {\d});
\draw[line width=1pt,gray] (\e,\f) arc (90:270:{0.4*\f} and {\f});
\end{scope}
\end{tikzpicture}}

\loigiai{
\immini{Chọn hệ trục tọa độ $Oxy$ với gốc $O$ là tâm elip, trục lớn nằm trên $Ox$.\\
Phương trình elip có dạng $\dfrac{x^2}{a^2}+\dfrac{y^2}{b^2}=1$. ($0<b<a$).\\
Do trục lớn bằng $10$ nên $2a=10 \Rightarrow a=5$.\\
Trục nhỏ bằng $8$ nên $2b=8 \Rightarrow b=4$.\\
Phương trình elip là $\dfrac{x^2}{25} + \dfrac{y^2}{16} = 1$.\\
Suy ra $y^2 = 16 \left( 1 - \dfrac{x^2}{25} \right)$.\\
}{\begin{tikzpicture}[scale=1,font=\footnotesize,line join = round, line cap = round, >= stealth]
\def\a{-3.5} \def\b{3.5} \def\c{-.5} \def\d{4.5}
\draw [domain=-3:3, samples=100] %
plot (\x, {4*sqrt(1-(\x/5)^2)});
\fill[blue,opacity=0.25] (-3,0)--plot[domain=-3:3, samples=100] (\x, {4*sqrt(1-(\x/5)^2)})--(3,0)--cycle;
\draw (-3,0)--(-3,3.2) (3,0)--(3,3.2);
\coordinate (O) at (0,0);
\draw[->] (\a,0)--(\b,0)node[below]{$x$};
\draw[->] (0,\c)--(0,\d)node[right]{$y$};
\foreach \p/\g in {-3/-90,3/-90} \draw[fill] (\p,0) circle(.5pt)node [shift={(\g:.3)}] {$\p$};
\foreach \p/\g in {4/60} \draw[fill] (0,\p) circle(.5pt)node [shift={(\g:.3)}] {$\p$};
%\draw[dashed] (3/2,0)|-(0,3) ;
\draw[fill] (O) circle (.5 pt) node[below right]{$O$} ;
\end{tikzpicture}}
\noindent Thùng rượu là vật thể tròn xoay tạo bởi hình phẳng giới hạn bởi cung elip, trục $Ox$ và hai đường thẳng $x=-3, x=3$ (do chiều cao thùng là $6$).\\
Thể tích thùng gỗ là
{\allowdisplaybreaks
\begin{eqnarray*}
V &=& \pi \displaystyle\int\limits_{-3}^3 y^2 \mathrm{\,d}x = \pi \displaystyle\int\limits_{-3}^3 16 \left( 1 - \dfrac{x^2}{25} \right) \mathrm{\,d}x\\
&=& 32\pi \displaystyle\int\limits_{0}^3 \left( 1 - \dfrac{x^2}{25} \right) \mathrm{\,d}x = 32\pi \left( x - \dfrac{x^3}{75} \right) \Big|_0^3 = 32\pi \left( 3 - \dfrac{27}{75} \right) \approx 265.
\end{eqnarray*}}
\noindent Vậy thùng gỗ này đựng được tối đa $265$ lít rượu.}
\end{ex}

\begin{ex}%[12-EX-3-2025, Trần Văn Hưng]%[2D4V3-5]
	Một hộ gia đình sản xuất cơ khí nhỏ mỗi ngày sản xuất được $x$ sản phẩm $(0\leq x\leq 20)$. Chi phí biên để sản xuất $x$ sản phẩm, tính bằng nghìn đồng, cho bởi hàm số $C'(x)=3x^2-4x+10$.\\
	Biết rằng chi phí cố định ban đầu để sản xuất là $500$ nghìn đồng. Giả sử cơ sở này bán hết sản phẩm mỗi ngày với giá $270$ nghìn đồng/sản phẩm. Tính lợi nhuận tối đa mà gia đình đó thu được khi sản xuất và bán sản phẩm?
	\shortans{1300}
	\loigiai{ Chi phí để sản xuất $x$ sản phẩm bằng $C(x)=\displaystyle\int C'(x)\mathrm{\,d}x=x^3-2x^2+10x+C$\\
		Mà chi phí cố định ban đầu để sản xuất là $500$ nghìn đồng nên $C(0)=500\Rightarrow C=500\Rightarrow C(x)=x^3-2x^2+10x+500$.\\
		Khi bán $x$ sản phẩm, số tiền thu được là $270x$ (nghìn đồng).\\
		Do đó lợi nhuận thu được là $T(x)=-x^3+2x^2+260x-500$ (nghìn đồng).\\
		Ta có $T'(x)=-3x^2+4x+260=0\Leftrightarrow \hoac{&x=10\\&x=-\dfrac{26}{3}\,\text{(loại)}.}$
		\begin{center}
			\begin{tikzpicture}		
				\tkzTabInit%[nocadre=false,lgt=1.2,espcl=2.5,deltacl=0.6]
				{$x$ /0.6, $T'$ /0.6, $T$ /2.5}
				{$0$,$10$,$20$}
				\tkzTabLine{,+,$0$,-,}
				\tkzTabVar{-/$ $,+/$1\,300$,-/$ $}	
		\end{tikzpicture}
		\end{center}
	}
\end{ex}

\begin{ex}%[12-EX-2-2025, Đình Nguyên]%[2D4V3-5]
\immini{Hình dưới mô phỏng phần bên trong của một chậu cây có dạng khối tròn xoay tạo thành khi quay một phần của đồ thị hàm số $y=\sqrt{x}+\dfrac{3}{2}$ với $0 \le x \le 4$ quanh trục hoành. Biết đơn vị trên các trục $Ox, Oy$ là decimet, thể tích phần bên trong (dung tích) của chậu cây là bao nhiêu lít (kết quả làm tròn đến hàng đơn vị)?
\shortans{104}
}
{
\begin{tikzpicture}[scale=0.6, font=\footnotesize, line join=round, line cap=round, >=stealth]
\fill[gray!20] plot[domain=0:4] (\x, {sqrt(\x)+1.5}) -- plot[domain=4:0] (\x, {-sqrt(\x)-1.5}) -- cycle;
\draw[->] (-2.5,0) -- (5.5,0) node[below] {$x$};
\draw[->] (0,-4.5) -- (0,4.5) node[left] {$y$};
\node at (0,0) [below left] {$O$};
\foreach \x in {-2,-1,1,2,3,4,5} \draw (\x,1pt) -- (\x,-1pt) node[below] {$\x$};
\foreach \y in {-4,-3,-2,-1,1,2,3,4} \draw (1pt,\y) -- (-1pt,\y) node[left] {$\y$};
\draw[domain=0:4, samples=100] plot (\x, {sqrt(\x)+1.5});
\draw[domain=0:4, samples=100] plot (\x, {-sqrt(\x)-1.5});
\draw (4,0) ellipse (0.4 and 3.5);
\draw[dashed] (0,0) ellipse (0.2 and 1.5);
\end{tikzpicture}
}
\loigiai{
Thể tích khối tròn xoay tạo thành khi quay đồ thị hàm số $y=f(x)$ quanh trục hoành được tính bởi công thức $V = \pi \displaystyle\int\limits_a^b [f(x)]^2 \mathrm{\,d}x$.\\
Áp dụng với $f(x) = \sqrt{x} + \dfrac{3}{2}$ trên đoạn $[0; 4]$, ta có
\begin{eqnarray*}
V &=& \pi \displaystyle\int\limits_0^4 \left(\sqrt{x} + \dfrac{3}{2}\right)^2 \mathrm{\,d}x = \pi \displaystyle\int\limits_0^4 \left(x + 3\sqrt{x} + \dfrac{9}{4}\right) \mathrm{\,d}x\\
&=& \pi \left. \left( \dfrac{x^2}{2} + 2x\sqrt{x} + \dfrac{9}{4}x \right) \right|_0^4\\
&=& \pi \left[ \left( \dfrac{4^2}{2} + 2 \cdot 4 \cdot \sqrt{4} + \dfrac{9}{4} \cdot 4 \right) - 0 \right]\\
&=& \pi (8 + 16 + 9) = 33\pi \approx 103{,}67 \text{ dm}^3.
\end{eqnarray*}
Vì $1 \text{ dm}^3 = 1$ lít nên dung tích của chậu cây khoảng $104$ lít.
}
\end{ex}

\begin{ex}%[Đề kiểm tra giữa kì 2 - THPT Thanh Miện, Hải Dương, Nguyễn Anh Tuấn, Dự án 12EX-3-2026]%[2D4V3-5]
Một bình cắm hoa dạng khối tròn xoay với đáy bình và miệng bình có đường kính lần lượt là $2$ và $4$. Mặt xung quanh của bình là một phần của mặt tròn xoay khi quay đường cong $y=\sqrt{x-1}$ quay quanh trục $Ox$. Thể tích của bình cắm hoa đó bằng bao nhiêu? (kết quả làm tròn đến hàng phần chục)
\shortans[oly]{23,6}
\begin{center}

\begin{tikzpicture}[
scale=0.9,
>=stealth,
every node/.style={font=\footnotesize}
]
% Trục tọa độ
\draw[->] (-0.5,0) -- (6,0) node[above] {$x$};
\draw[->] (0,-0.5) -- (0,3) node[right] {$y$};
\node[below left] at (0,0) {$O$}; % Thêm gốc tọa độ O

% Vẽ đồ thị hàm số y = sqrt(x-1)
% Thêm samples để đường cong mượt hơn
\draw[domain=1:5.6, smooth,samples=300]
plot (\x, {sqrt(\x-1)});

% Đường dóng (Dashed lines)
\draw[dashed] (2,0) |- (0,1); % Cách viết nhanh cho đường dóng vuông góc
\draw[dashed] (5,0) |- (0,2);

% Vẽ các điểm chấm tròn
\foreach \p in {(1,0), (2,0), (5,0), (0,1), (0,2)}
\fill \p circle (1.5pt);

% Gán nhãn số trên trục
\node[below] at (1,0) {$1$};
\node[below] at (2,0) {$2$};
\node[below] at (5,0) {$5$};
\node[left] at (0,1) {$1$};
\node[left] at (0,2) {$2$};

\end{tikzpicture}
\end{center}
\loigiai{
Thể tích của bình hoa là thể tích khối tròn xoay khi quay hình phẳng giới hạn bởi $y=\sqrt{x-1}$, trục hoành, $x=2$, $x=5$ quanh trục $Ox$.
\[ V = \pi \displaystyle\int\limits_{2}^{5} \left(\sqrt{x-1}\right)^2 \,\mathrm{d}x = \pi \displaystyle\int\limits_{2}^{5} (x-1) \,\mathrm{d}x. \]
\[ V = \pi \left( \dfrac{x^2}{2} - x \right) \bigg|_{2}^{5} = \pi \left[ \left(\dfrac{25}{2} - 5\right) - \left(\dfrac{4}{2} - 2\right) \right] = \pi \left( \dfrac{15}{2} - 0 \right) = \dfrac{15\pi}{2}\approx 23{,}6. \]
Vậy thể tích của bình cắm hoa xấp xỉ $23{,}6$ (đơn vị thể tích).
}
\end{ex}

\begin{ex}%[2D4V3-5]%[Quang Vinh NT]
Một chiếc bát thủy tinh có bề dày của phần xung quanh là một khối tròn xoay, khi xoay hình phẳng $D$ quanh một đường thẳng $a$ bất kì nào đó mà khi gắn hệ trục tọa độ $Oxy$ (đơn vị trên trục là decimet) vào hình phẳng $D$ tại một vị trí thích hợp, thì đường thẳng $a$ sẽ trùng với trục $Ox$. Khi đó hình phẳng $D$ được giới hạn bởi các đồ thị hàm số $y = x + \dfrac{1}{x}$, $y = x$ và hai đường thẳng $x=1$, $x=4$ như hình dưới đây.
\begin{center}
\begin{tikzpicture}[line join=round, line cap=round,>=stealth,scale=0.7]
\tikzset{every node/.style={scale=0.9}}
\draw[->] (-1.1,0)--(6.1,0) node[below left] {$x$};
\draw[->] (0,-5.1)--(0,5.1) node[below left] {$y$};
\draw (0,0) node [below left] {$O$};
\foreach \x/\nx in {1/1,2/2,3/3,4/4}
\draw[thin] (\x,1pt)--(\x,-1pt) node [below] {$\nx$};
\foreach \y/\ny in {1/1,2/2}
\draw[thin] (1pt,\y)--(-1pt,\y) node [left] {$\ny$};
\draw[dashed,thin] (0.01,-5)--(0.01,5);
\draw(4,0)--(4,4);
\begin{scope}
\clip (-1,-5) rectangle (6,5);
\draw[fill=gray] (1,0)--plot[domain=1:4] (\x,{(1*((\x)^2) +0*(\x)+1)/(1*(\x)+0)})--(4,4)--(1,1)--cycle;
\draw[samples=200,domain=0.01:6,smooth,variable=\x] plot (\x,{(1*((\x)^2) +0*(\x)+1)/(1*(\x)+0)});
\draw[dashed,samples=200,domain=1:6,smooth,variable=\x] plot (\x,{(-1*((\x)^2)-0*(\x)-1)/(1*(\x)+0)});
\draw (0,0)--(6.1,6.1); %Tiệm cận xiên
\draw[dashed,samples=200,domain=0:360,smooth,variable=\t]
plot ({1 + 0.5*cos(\t)}, {2*sin(\t)});
\draw[dashed,samples=200,domain=0:360,smooth,variable=\t]
plot ({1 + 0.25*cos(\t)}, {1*sin(\t)});
\draw[dashed,samples=200,domain=0:360,smooth,variable=\t]
plot ({4 + 0.5*cos(\t)}, {4.25*sin(\t)});
\draw[dashed,samples=200,domain=0:360,smooth,variable=\t]
plot ({4 + 0.4*cos(\t)}, {4*sin(\t)});
\draw[dashed] (0,0)--(6.1,-6.1); %Tiệm cận xiên
\end{scope}
\end{tikzpicture}
\end{center}
Thể tích của bề dày chiếc bát thủy tinh đó bằng bao nhiêu decimet khối? (làm tròn kết quả đến hàng phần mười).
\par
\shortans{21,2}
\loigiai{
Hình phẳng $D$ giới hạn bởi các đường $y = f(x) = x + \dfrac{1}{x}$, $y = g(x) = x $, $x=1$, $x=4$.\\
Khi quay hình phẳng $D$ quanh trục $Ox$, ta được khối tròn xoay là phần bề dày của chiếc bát.\\
Thể tích của khối tròn xoay
\begin{center}
$V = \pi \displaystyle \int_{1}^{4} \left[ f^2(x) - g^2(x) \right] \mathrm{\,d}x= \pi \displaystyle\int_{1}^{4} \left(2 + \dfrac{1}{x^2}\right) \mathrm{\,d}x= \pi \left[ 2x - \dfrac{1}{x} \right]\bigg|_{1}^{4}= \dfrac{27\pi}{4}\approx21{,}2$(dm$^3$).
\end{center}
}
\end{ex}

\begin{ex}%[Nguyễn Hữu Chung Kiên]%[MathTex_TN THPT 2026]%[2D4V3-5]
\immini
{Kiến trúc sư thiết kế một khu sinh hoạt cộng đồng có dạng hình chữ nhật với chiều rộng và chiều dài lần lượt là $60$ m và $80$ m. Trong đó, phần được tô màu đậm là sân chơi, phần còn lại để trồng hoa. Mỗi phần trồng hoa có đường biên cong là một phần của parabol với đỉnh thuộc một trục đối xứng của hình chữ nhật và khoảng cách từ đỉnh đó đến trung điểm cạnh tương ứng của hình chữ nhật bằng $20$ m (xem hình minh họa).
Diện tích của phần sân chơi là bao nhiêu mét vuông?
}
{\begin{tikzpicture}[scale=.6, font=\footnotesize, line join=round, line cap=round, >=stealth]
\coordinate (A) at (0,0);
\coordinate (B) at (6,0);
\coordinate (C) at (6,8);
\coordinate (D) at (0,8);
\coordinate (E) at (3,6);
\coordinate (F) at (3,2);
\draw[fill=green!30!gray] (A) -- (B) -- (C) -- (D) -- cycle;
\draw[fill=white]  (A) parabola bend (F) (B);
\draw[fill=white] (D) parabola bend (E) (C);
\draw (A)--(B) node[below,midway] {$60$ m};
\draw (C)--(B) node[right,midway] {$80$ m};
\draw[<->,dashed] (E)--(3,8) node[right,midway] {$20$ m};
\draw[<->,dashed] (F)--(3,0) node[right,midway] {$20$ m};
\end{tikzpicture}}
\par \shortans[0]{3200}
\loigiai{
\immini
{Diện tích của phần sân chơi bằng diện tích của hình chữ nhật trừ đi diện tích của hai phần trồng hoa.\\
Diện tích của hình chữ nhật là $60 \cdot 80=4\,800 \mathrm{~m}^2$.\\
Diện tích của mỗi phần trồng hoa bằng diện tích của một parabol.\\
Chọn hệ trục $Oxy$ như hình vẽ.\\
Phương trình parabol bên dưới có dạng $y=a x^2+b$ (vì đỉnh parabol thuộc trục tung).\\
Ta có $\heva{
&y(0)=20 \\
&y(30)=0} \Leftrightarrow \heva{&b=20 \\
&a\cdot30^2+20=0} \Rightarrow\heva{&a=-\dfrac{1}{45} \\
&b=20.}$\\
Phương trình Parabol dưới là $y=-\dfrac{1}{45} x^2+20$.}
{\begin{tikzpicture}[scale=.6, font=\footnotesize, line join=round, line cap=round, >=stealth]
% Định nghĩa toạ độ 6 điểm:
\coordinate (A) at (0,0);
\coordinate (B) at (6,0);
\coordinate (C) at (6,8);
\coordinate (D) at (0,8);
\coordinate (E) at (3,6);
\coordinate (F) at (3,2);
\draw[fill=green!30!gray] (A) -- (B) -- (C) -- (D) -- cycle;
\draw[fill=white]  (A) parabola bend (F) (B);
\draw[fill=white] (D) parabola bend (E) (C);
\fill (3,0) circle (1pt) node[below] {$O$};
\draw (C)--(B) node[right,midway] {$80$ m};
\draw[dashed] (E)--(3,8) node[right,midway] {$20$ m};
\draw[dashed] (F)--(3,0) node[right,midway] {$20$ m};
% % Vẽ trục toạ độ (x và y), chiều dài tuỳ ý
\draw[->] (3,0) -- (7,0) node[right] {$x$};
\draw[->] (3,0) -- (3,9) node[above] {$y$};

\end{tikzpicture}}
\noindent Diện tích của mỗi phần trồng hoa bằng
\[
\displaystyle\int\limits_{-30}^{30}\left(-\dfrac{1}{45} x^2+20\right) \,\mathrm{d}x=800 \mathrm{~m}^2.
\]
Diện tích của phần sân chơi là $4\,800-2 \cdot 800=3\,200 \mathrm{~m}^2$.}
\end{ex}

\begin{ex}%[12-TT-PNL-Deso12-PNL]%[pb Lương Pho]%[2D4V3-5]
\immini[thm]{Người ta cần thiết kế đoạn một cây cầu cong nối giữa hai điểm $A$ và $B$ trên hai bờ sông sao cho cây cầu này có hình dạng thẩm mỹ và ổn định. Để mô phỏng hình dáng của cây cầu, người ta sử dụng một hàm bậc ba. Trong hệ tọa độ $Oxy$ (đơn vị: feet). Bằng cách đo đạc tại thực địa, ta xác định được đường cong của cây cầu đi qua hai điểm có toạ độ lần lượt là $A(-1\,200;80)$,  $B(1\,200;60)$ và điểm giữa cây cầu (tại $x=0$) có độ cao là $200$ feet. Tiếp tuyến tại điểm giữa cây cầu (tại $x=0$) cắt trục hoành tại điểm có hoành độ là $-1\,000$ feet (Hình vẽ bên).
Tính tổng chi phí xây dựng cây cầu dựa trên diện tích bề mặt của cây cầu (Đơn vị: triệu USD), biết chi phí xây dựng là $ 100 $ USD cho mỗi feet vuông của mặt cầu (Diện tích bề mặt của cây cầu được xác định là diện tích của phần hình phẳng được giới hạn bởi đồ thị của hàm số $y=f(x)$, trục hoành và các đường thẳng $x=-1\,200$; $x=1\,200$).
}{
\begin{tikzpicture}[scale=.6,line cap=round,line join=round,>=stealth,font=\footnotesize]
\definecolor{peru}{rgb}{0.80,0.52,0.25}
\definecolor{Maroon}{rgb}{0.50,0.00,0.00}
\definecolor{DodgerBlue}{rgb}{0.12,0.56,1.00}
\definecolor{SeaGreen}{rgb}{0.18,0.55,0.34}
\definecolor{Cyan}{rgb}{0.00,1.00,1.00}
\definecolor{DeepSkyBlue}{rgb}{0.00,0.75,1.00}
\tikzset{chancau/.pic={
\draw[fill=gray] (2.80,-3.23)-- (2.80,-3.59)-- (3.12,-3.84)-- (3.13,-3.04) -- cycle;
\draw[fill=gray!50] (3.12,-3.84)-- (3.47,-3.82)-- (3.48,-2.79)-- (3.13,-3.04) -- cycle;
\draw[fill=gray] (3.47,-3.82)-- (3.71,-3.64)-- (3.71,-2.63)-- (3.48,-2.79) -- cycle;
}}
\tikzset{khicau/.pic={
\definecolor{royal blue}{rgb}{0.25,0.41,0.88}
\definecolor{dark turquoise}{rgb}{0,0.81,0.82}
\definecolor{crimson}{rgb}{0.86,0.08,0.24}
\definecolor{khaki}{rgb}{0.94,0.90,0.55}
\definecolor{ballblue}{rgb}{0.13,0.67,0.8}
\def\c{(0.49,11.54)..controls++(268:0.54) and++(141:0.74)..(1.6,8.72)..controls++(93:1.34) and++(252:0.9)..(1.92,11.75)..controls++(290:1.59) and++(127:0.44)..(3.43,8.52)..controls++(88:0.47) and++(253:1.29)..(4.06,11.9)..controls++(289:1.31) and++(126:1.63)..(5.71,8.42)..controls++(61:1.72) and++(259:1.12)..(7.03,11.83)..controls++(284:1.55) and++(89:1.07)..(7.49,8.62)..controls++(21:0.29) and++(255:1.86)..(9.18,11.51)..controls++(291:1.14) and++(87:0.93)..(9.46,8.89)}
%Chỗ ngồi
\def\e{(3.15,2.66) -- (3.15,0.72)..controls++(288:0.22) and++(180:0.07).. (3.43,0.44) -- (6.60,0.44)..controls++(4:0.12) and++(265:0.12).. (6.88,0.66) -- (6.89,2.64) -- cycle}
\draw[fill=brown]
\e;
\draw[fill=brown,pattern=fivepointed stars,pattern color=ballblue]
\e;
\draw[fill=red!80!blue] (3.02,3.08) -- (7.11,3.05) -- (7.12,2.65) -- (3.01,2.66) -- cycle;
\fill[gray!70!blue] (3.13,2.64)-- (6.86,2.64)-- (6.87,2.52)-- (3.12,2.52)--cycle;
%4 dây khéo
\draw[fill=brown] (3.09,3.08) -- (3.17,3.08) -- (3.17,5.96) -- (3.09,6.03) -- cycle;
\draw[fill=brown] (3.22,5.93) -- (3.62,3.08) -- (3.72,3.08) -- (3.33,5.84) -- cycle;
\draw[fill=brown] (6.86,5.87) -- (6.50,3.05) -- (6.60,3.05) -- (6.96,5.95) -- cycle;
\draw[fill=brown] (6.99,5.98) -- (6.99,3.05) -- (7.07,3.05) -- (7.07,6.04) -- cycle;
%Phần quả kinh khi cầu
\draw[fill=crimson]
(4.1,5.17)..controls++(145:5.84) and++(244:0.35)..(0.49,11.54)..controls++(268:0.54) and++(141:0.74)..(1.6,8.72)..controls++(93:1.34) and++(252:0.9)..(1.92,11.75)..controls++(290:1.59) and++(127:0.44)..(3.43,8.52)..controls++(88:0.47) and++(253:1.29)..(4.06,11.9)..controls++(289:1.31) and++(126:1.63)..(5.71,8.42)..controls++(61:1.72) and++(259:1.12)..(7.03,11.83)..controls++(284:1.55) and++(89:1.07)..(7.49,8.62)..controls++(21:0.29) and++(255:1.86)..(9.18,11.51)..controls++(291:1.14) and++(87:0.93)..(9.46,8.89)..controls++(246:1.14) and++(31:2.5)..(5.88,5.17)..controls++(185:0.8) and++(353:0.68)..(4.11,5.18);
\draw[fill=dark turquoise]
\c..controls++(71:1.3) and++(321:1.36)..(8.87,12.53)..controls++(131:2.5) and++(13:1.04)..(3.58,13.94)..controls++(189:1.64) and++(72:1.2)..(0.49,11.54);
\draw[fill=khaki]
(4.1,5.17)..controls++(128:2.52) and++(296:0.89)..(1.61,8.66)..controls++(97:0.58) and++(251:0.94)..(1.92,11.69)..controls++(269:2.73) and++(117:2.5)..(4.3,5.16)--cycle
(4.54,5.15)..controls++(109:1.99) and++(284:0.5)..(3.47,8.47)..controls++(89:1.41) and++(255:1.34)..(4.04,11.87)..controls++(269:2.87) and++(97:2.57)..(4.77,5.16)--cycle
(5.24,5.15)..controls++(81:1.66) and++(265:0.58)..(5.74,8.46)..controls++(63:1.13) and++(260:1.5)..(7.02,11.79)..controls++(271:2.33) and++(71:2.77)..(5.48,5.16)--cycle
(5.72,5.18)..controls++(60:2.18) and++(251:0.51)..(7.49,8.59)..controls++(39:1.03) and++(254:1.43)..(9.18,11.46)..controls++(270:2.54) and++(54:2.22)..(5.91,5.19);
\draw[fill=blue!80]
(1.59,8.67)..controls++(134:0.91) and++(273:1.07)..(0.5,11.58)..controls++(65:0.43) and++(227:0.63)..(1.23,12.8)..controls++(255:0.46) and++(116:3.02)..(1.59,8.67)
(3.45,8.49)..controls++(121:1.73) and++(294:0.53)..(1.93,11.72)..controls++(83:1) and++(183:1.42)..(3.61,13.92)..controls++(237:1.11) and++(104:3.26)..(3.45,8.49)
(5.71,8.46)..controls++(131:1.12) and++(281:0.53)..(4.06,11.87)..controls++(93:1.13) and++(240:0.32)..(4.33,14.03)..controls++(21:0.51) and++(163:0.24)..(5.73,14.02)..controls++(292:0.44) and++(79:2.4)..(5.71,8.46)
(7.49,8.61)..controls++(92:1.08) and++(290:0.56)..(7.04,11.79)..controls++(94:0.85) and++(306:0.88)..(6.41,13.9)..controls++(360:2.13) and++(67:2.32)..(7.5,8.6)
(9.45,8.93)..controls++(86:1.28) and++(296:0.66)..(9.16,11.5)..controls++(86:0.29) and++(301:0.43)..(8.8,12.64)..controls++(318:0.67) and++(66:1.98)..(9.45,8.93);
%Bóng mờ
\draw[opacity=0.5,fill=white]
(3.51,5.9)..controls++(135:2.77) and++(247:2.6)..(1.18,12.68)..controls++(259:1.61) and++(129:4.74)..(3.49,5.92)
(3.84,5.97)..controls++(123:4.05) and++(222:2.36)..(2.75,13.74)..controls++(257:1.09) and++(119:5.39)..(3.84,5.97)
(4.32,5.65)..controls++(113:5.66) and++(226:1.27)..(3.65,13.95)..controls++(240:2.09) and++(108:4.75)..(4.32,5.65)
(4.71,5.74)..controls++(104:4.41) and++(241:1.33)..(4.34,14.08)..controls++(263:1.86) and++(100:4.12)..(4.71,5.74)
(5.28,5.64)..controls++(81:4.84) and++(285:1.06)..(5.46,14.11)..controls++(354:1.05) and++(77:5.67)..(5.28,5.6)
(5.72,5.7)..controls++(72:5.24) and++(286:1.18)..(6.42,13.91)..controls++(329:1.31) and++(68:4.59)..(5.72,5.7)
(5.94,5.53)..controls++(62:3.44) and++(294:2.79)..(7.36,13.67)..controls++(312:0.77) and++(57:6.84)..(5.94,5.53)
(6.64,5.98)..controls++(54:5.04) and++(273:1.49)..(8.83,12.57)..controls++(295:2.82) and++(50:2.95)..(6.61,5.98);
\draw[fill=purple]
(4.,5.2)..controls++(260:0.25) and++(177:0.68)..(5.09,5)..controls++(3.3:0.99) and++(259:0.19)..(5.94,5.19)..controls++(188:0.37) and++(354:0.4)..(4.,5.2);
}}
%Nền
\clip (-3.6,-4.3) rectangle (6.5,10);
\fill[DeepSkyBlue!80] (-3.6,-4.3) rectangle (6.5,10);
\fill[top color=Cyan,bottom color=white] (-4,3)..controls++(5:4.09) and++(190:2.4)..(6.7,3)--(6.7,10)--(-4,10)--cycle;
\fill[blue!70] (-2.34,1.56)..controls++(335:1.84) and++(127:1.86)..(2.07,-0.69)--(2.59,0.34)..controls++(166:5.69) and++(327:0.76)..(0.36,2.02)--cycle;
%đất
\draw[left color=green!30,right color=green!60]
(-3.63,-3.57)..controls++(26:10.12) and++(196:3.02)..(3.28,1.45)--(6.5,1.63)--(6.5,-4.3)--(-3.62,-4.3)--cycle;
\draw[fill=SeaGreen]
(-3.61,1.37)..controls++(6:3.43) and++(188:2.85)..(0.47,2.07)--(-3.7,2.76)--cycle;
\pic at (0.,-0.1) {chancau};
\pic at (2,1.3) {chancau};
\pic[scale=0.15] at (3,6) {khicau};
\tikzset{cay/.pic={
\draw[fill=SeaGreen] (2.68,-1.18)..controls++(195:0.36) and++(156:0.52)..(2.63,-0.49)..controls++(141:0.57) and++(76:0.97)..(3,-0.47)..controls++(46:0.9) and++(39:0.45)..(3.03,-1.09)--cycle;
}}
\pic at (-3,-1) {cay};
\pic at (-1.2,0.8) {cay};
\pic at (-1,0.6) {cay};
\pic at (1,0.6) {cay};
\pic at (0.7,0.6) {cay};
\pic at (3.4,2.5) {cay};
\draw[fill=black!30]
(-2.42,2.75)..controls++(354:1.94) and++(260:0.66)..(-1.67,4.42)..controls++(76:0.45) and++(244:0.19)..(-1.39,5.07)..controls++(265:1.42) and++(164:0.76)..(-0.22,2.62)..controls++(173:1.06) and++(348:0.64)..(-2.42,2.75);
\draw[fill=gray]
(6.5,0.59)..controls++(269:1.96) and++(35:5.44)..(1.1,-4.5)--(0.5,-4.5)..controls++(37:3.6) and++(262:1.55)..(6.5,0.59);
\draw[fill=peru]
(0,-4.48)..controls++(29:10.24) and++(359:12.37)..(-3.6,2.77)--(-3.6,2.47)..controls++(358:17.58) and++(27:5.21)..(-3.62,-4.41);
\draw[fill=Maroon]
(0,-4.48)..controls++(29:10.24) and++(359:12.37)..(-3.6,2.77)--(-3.6,2.81)..controls++(360:17.89) and++(36:3.66)..(0.58,-4.48)
(-3.6,2.4)..controls++(358:17.58) and++(27:5.21)..(-3.62,-4.3)--(-3.42,-4.41)..controls++(27:5.21) and++(358:17.58)..(-3.6,2.47);
\draw[fill=orange!30]
(4.6,-0.33)..controls++(358:0.84) and++(357:1.81)..(0.68,7.42)..controls++(177:2.66) and++(17:1.92)..(-3.6,2.29)--(-3.3,2.02)..controls++(30:1.72) and++(170:2.27)..(0.63,7.02)..controls++(356:1.47) and++(87:1.35)..(4.22,0.99)..controls++(134:1.96) and++(353:1.34)..(-2.58,2.4)--(-2.7,2.3)..controls++(358:5.69) and++(97:1.67)..(4.6,-0.33);
\draw[fill=gray]
(-3.3,2.03)..controls++(32:2.23) and++(175:2.51)..(0.63,7.04)..controls++(187:2.16) and++(40:2.04)..(-3.04,2)--cycle;
\draw[fill=orange!70]
(4.78,-0.25)..controls++(93:1.14) and++(354:1.91)..(0.68,7.4)..controls++(352:2.84) and++(7:0.35)..(6.44,0.84)..controls++(142:0.2) and++(345:0.29)..(5.63,1.68)..controls++(56:0.77) and++(0:0.28)..(2.06,6.58)..controls++(307:3.69) and++(129:2.65)..(5.23,0.96)..controls++(313:0.43) and++(41:0.64)..(4.78,-0.25);
\end{tikzpicture}
}
\par
\shortans[0]{37,6}
\loigiai{
Gọi $y=f(x)=ax^3+bx^2+cx+d$ ($a\ne 0$) là hàm số có đồ thị mô ta hình dáng cây cầu.\\
Ta có $f'(x)=3ax^2+2bx+c$.\\
Do đồ thị của hàm số đi qua các điểm $A(-1\,200;80)$, $B(1\,200;60)$ nên ta có
\[\heva{&{f(-1\,200)=80}\\
&{f(1\,200)=60}}\Rightarrow\heva{&
{a{(-1\,200)^3}+b{(-1\,200)^2}+c(-1\,200)+d=80}\\
&{a\cdot 1\,200^3+b\cdot 1\,200^2+c\cdot 1\,200+d=60}.
} \qquad (*)\]
Điểm giữa cây cầu (tại $x=0$) có độ cao là $200$ feet nên \[f(0)=200\Rightarrow d=200.\]
Tại $x=0$ ta có $f(0)=200$, $f'(0)=c$.\\
Tiếp tuyến tại điểm giữa cây cầu (tại $x=0$ ) có dạng
$y=cx+200$.\\
Tiếp tuyến này cắt trục hoành tại điểm có hoành độ là $-1\,000$ feet. Do đó
\[0=c\cdot (-1\,000)+200\Leftrightarrow c=0{,}2.\]
Khi đó hệ $(*)$ được viết lại
\begin{eqnarray*}
& &\heva{
&a(-1\,200)^3+b(-1\,200)^2+0{,}2\cdot(-1\,200)+200=80\\
&a\cdot 1\,200^3+b\cdot 1\,200^2+0{,}2\cdot1\,200+200=60}\\
&\Leftrightarrow &
\heva{
&a(-1\,200)^3+b(-1\,200)^2=120\\
&a\cdot 1\,200^3+b\cdot 1\,200^2=-380
}\\
&\Leftrightarrow &\heva{
&a=-\dfrac{1}{6\,912\,000}\\
&b=-\dfrac{13}{144\,000}.}
\end{eqnarray*}
Vậy $y=f(x)=-\dfrac{1}{6\,912\,000}{x^3}-\dfrac{13}{144\,000}{x^2}+0{,}2x+200$.\\
Diện tích bề mặt cây cầu là \[S=\displaystyle\int\limits_{-1\,200}^{1\,200}\left|-\dfrac{1}{6\,912\,000} x^3-\dfrac{13}{144\,000} x^2+0{,}2 x+200\right| \mathrm{d} x=376\,000\,\left(\text{feet}^2\right).\]
Chi phí cần để xây cầu là $C=376\,000\cdot100=37\,600\,000$ (USD) $=37{,}6$ (triệu USD).
}
\end{ex}

\begin{ex}%[Đề thi thử trường THPT Hoàng Diệu Đồng Nai- đề 02, Nguyễn Hữu Đức]%[2D4V3-5]
Để chuẩn bị cho hội trại do Đoàn trường tổ chức, lớp $12$B dự định dựng một cái lều trại có dạng hình parabol. Nền của lều trại là một hình chữ nhật có kích thước bề ngang $3$ mét, chiều dài $6$ mét, đỉnh trại cách nền $3$ mét. Tính thể tích phần không gian bên trong lều trại.
\par\shortans[0]{36}
\loigiai{
\immini{
Chọn hệ trục tọa độ $Oxy$ sao cho đỉnh lều thuộc trục $Oy$ và nền lều thuộc trục $Ox$.
Khi đó, parabol $(P)$ có đỉnh $I(0;3)$ và đi qua hai điểm $A\left(-1{,}5; 0\right)$, $B\left(1{,}5; 0\right)$ (vì bề ngang là $3$ m).
Phương trình $(P)$ có dạng $y = ax^2 + c$ (do trục đối xứng là $Oy$).
Ta có:
\[ \heva{& c = 3 \\ & a \cdot (1{,}5)^2 + c = 0} \Leftrightarrow \heva{& c = 3 \\ & 2{,}25a = -3} \Leftrightarrow \heva{& c = 3 \\ & a = -\dfrac{4}{3}.} \]
Vậy $(P): y = -\dfrac{4}{3}x^2 + 3$.
Diện tích mặt cắt ngang (đáy parabol) là:
\[ S = \displaystyle\int\limits_{-1{,}5}^{1{,}5} \left(-\dfrac{4}{3}x^2 + 3\right) \,\mathrm{d}x = 6\; \text{m}^2. \]
}{
\begin{tikzpicture}[>=stealth,scale=1]
\draw[->] (-2.5,0) -- (2.5,0) node[below] {$x$};
\draw[->] (0,-1) -- (0,4) node[right] {$y$};
\draw[fill=black] (0,0) circle (1pt) node[below left]{$O$};
\draw[thick, smooth, domain=-1.6:1.6] plot (\x, {-4/3*(\x)^2 + 3});
\draw[dashed] (-1.5,0) -- (-1.5,0) node[below left] {$-1{,}5$};
\draw[dashed] (1.5,0) -- (1.5,0) node[below left] {$1{,}5$};
\draw (0,3) node[above right] {$3$} circle (1pt);
\fill (-1.5,0) circle (1pt) (1.5,0) circle (1pt) (0,3) circle (1pt);
\end{tikzpicture}
}
\noindent
\textit{Chú ý:} Ta có thể dùng công thức nhanh diện tích hình phẳng giới hạn bởi parabol và trục hoành: $S = \dfrac{2}{3} \cdot \text{ngang} \cdot \text{cao} = \dfrac{2}{3} \cdot 3 \cdot 3 = 6\;\text{m}^2$.\\
Thể tích lều trại (chiều dài $h=6$ m) là:
\[ V = S \cdot h = 6 \cdot 6 = 36\; \text{m}^3. \]
}
\end{ex}

